% Auto-generated from data/segments.json + layout.json (+ transforms) — do not edit by hand.
% volume: 06Di01
% mode: printing
% segments: 1230
% transform_rules: 0
% toc_marks: 292

% Volume layout defaults (layout.json ``layout``).
\csromanlayoutapply{1}{1.5}{21.6pt}{5pt}{30pt}{65pt}{2.5em}%

\setcsromanpitaka{สุตฺตนฺตปิฏก}
\setcsromannikaya{ทีฆนิกาย}
\csromanheader{2}{\textbf{ทีฆนิกาย}}
\csromanmatikaanchor{mtk.0}
\csromantocmarknopage{chapter}{สีลกฺขนฺธวคฺคปาฬิ}{mtk.0}
\bookmark[dest={mtk.0},level=0]{สีลกฺขนฺธวคฺคปาฬิ}
\gambhira{\textbf{สีลกฺขนฺธวคฺคปาฬิ}}
\namakkaram{นโม ตสฺส ภควโต อรหโต สมฺมาสมฺพุทฺธสฺส.}
\csromanmatikaanchor{mtk.1}
\csromanmatikaanchor{mtk.2}
\csromantocmarknopage{chapter}{1. พฺรหฺมชาลสุตฺต}{mtk.1}
\bookmark[dest={mtk.1},level=0]{1. พฺรหฺมชาลสุตฺต}
\csromantocmarkref{section}{ปริพฺพาชกกถา}{mtk.2}
\bookmark[dest={mtk.2},level=1]{ปริพฺพาชกกถา}
\setcsromanheadcha{1. พฺรหฺมชาลสุตฺต}
\setcsromanheadhclear{}
\csromanheader{1}{\textbf{ปริพฺพาชกกถา}}
\proseitem{1}{\csromanfolio{1}เอวํ เม สุตํ, เอกํ สมยํ ภควา อนฺตรา จ ราชคหํ อนฺตรา จ นาฬนฺทํ อทฺธานมคฺคปฺปฏิปนฺโน โหติ มหตา ภิกฺขุสํเฆน สทฺธิํ ปญฺจมตฺเตหิ ภิกฺขุสเตหิ, สุปฺปิโยปิ โข ปริพฺพาชโก อนฺตรา จ ราชคหํ อนฺตรา จ นาฬนฺทํ อทฺธานมคฺคปฺปฏิปนฺโน โหติ สทฺธิํ อนฺเตวาสินา พฺรหฺมทตฺเตน มาณเวน, ตตฺร สุทํ สุปฺปิโย ปริพฺพาชโก อเนกปริยาเยน พุทฺธสฺส อวณฺณํ ภาสติ, ธมฺมสฺส อวณฺณํ ภาสติ, สํฆสฺส อวณฺณํ ภาสติ, สุปฺปิยสฺส ปน ปริพฺพาชกสฺส อนฺเตวาสี พฺรหฺมทตฺโต มาณโว อเนกปริยาเยน พุทฺธสฺส วณฺณํ ภาสติ, ธมฺมสฺส วณฺณํ ภาสติ, สํฆสฺส วณฺณํ ภาสติ, อิติห เต อุโภ อาจริยนฺเตวาสี อญฺญมญฺญสฺส อุชุวิปจฺจนีกวาทา ภควนฺตํ ปิฏฺฐิโต ปิฏฺฐิโต อนุพนฺธา\footnote{อนุพทฺธา (กสี, อิ)} โหนฺติ ภิกฺขุสํฆญฺจ.}

\proseitem{2}{อถ โข ภควา อมฺพลฏฺฐิกายํ ราชาคารเก เอกรตฺติวาสํ อุปคจฺฉิ\footnote{อุปคญฺฉิ (สี, สฺยา, กํ, อิ)} สทฺธิํ ภิกฺขุสํเฆน, สุปฺปิโยปิ โข ปริพฺพาชโก อมฺพลฏฺฐิกายํ ราชาคารเก เอกรตฺติวาสํ อุปคจฺฉิ2 สทฺธิํ อนฺเตวาสินา พฺรหฺมทตฺเตน มาณเวน.\csromanspacer{} ตตฺรปิ สุทํ สุปฺปิโย ปริพฺพาชโก อเนกปริยาเยน \csromanfolio{2}พุทฺธสฺส อวณฺณํ ภาสติ, ธมฺมสฺส อวณฺณํ ภาสติ, สํฆสฺส อวณฺณํ ภาสติ, สุปฺปิยสฺส ปน ปริพฺพาชกสฺส อนฺเตวาสี พฺรหฺมทตฺโต มาณโว อเนกปริยาเยน พุทฺธสฺส วณฺณํ ภาสติ, ธมฺมสฺส วณฺณํ ภาสติ, สํฆสฺส วณฺณํ ภาสติ, อิติห เต อุโภ อาจริยนฺเตวาสี อญฺญมญฺญสฺส อุชุวิปจฺจนีกวาทา วิหรนฺติ.}

\proseitem{3}{อถ โข สมฺพหุลานํ ภิกฺขูนํ รตฺติยา ปจฺจูสสมยํ ปจฺจุฏฺฐิตานํ มณฺฑลมาเฬ สนฺนิสินฺนานํ สนฺนิปติตานํ อยํ สํขิยธมฺโม อุทปาทิ “อจฺฉริยํ อาวุโส, อพฺภุตํ อาวุโส, ยาวญฺจิทํ เตน ภควตา ชานตา ปสฺสตา อรหตา สมฺมาสมฺพุทฺเธน สตฺตานํ นานาธิมุตฺติกตา สุปฺปฏิวิทิตา, อยํ หิ สุปฺปิโย ปริพฺพาชโก อเนกปริยาเยน พุทฺธสฺส อวณฺณํ ภาสติ, ธมฺมสฺส อวณฺณํ ภาสติ, สํฆสฺส อวณฺณํ ภาสติ, สุปฺปิยสฺส ปน ปริพฺพาชกสฺส อนฺเตวาสี พฺรหฺมทตฺโต มาณโว อเนกปริยาเยน พุทฺธสฺส วณฺณํ ภาสติ, ธมฺมสฺส วณฺณํ ภาสติ, สํฆสฺส วณฺณํ ภาสติ, อิติหเม อุโภ อาจริยนฺเตวาสี อญฺญมญฺญสฺส อุชุวิปจฺจนีกวาทา ภควนฺตํ ปิฏฺฐิโต ปิฏฺฐิโต อนุพนฺธา โหนฺติ ภิกฺขุสํฆญฺจา”ติ.}

\proseitem{4}{อถ โข ภควา เตสํ ภิกฺขูนํ อิมํ สํขิยธมฺมํ วิทิตฺวา เยน มณฺฑลมาโฬ เตนุปสงฺกมิ, อุปสงฺกมิตฺวา ปญฺญตฺเต อาสเน นิสีทิ, นิสชฺช โข ภควา ภิกฺขู อามนฺเตสิ “กายนุตฺถ ภิกฺขเว เอตรหิ กถาย สนฺนิสินฺนา สนฺนิปติตา, กา จ ปน โว อนฺตรากถา วิปฺปกตา”ติ.\csromanspacer{} เอวํ วุตฺเต เต ภิกฺขู ภควนฺตํ เอตทโวจุํ “อิธ ภนฺเต อมฺหากํ รตฺติยา ปจฺจูสสมยํ ปจฺจุฏฺฐิตานํ มณฺฑลมาเฬ สนฺนิสินฺนานํ สนฺนิปติตานํ อยํ สํขิยธมฺโม อุทปาทิ ‘อจฺฉริยํ อาวุโส, อพฺภุตํ อาวุโส, ยาวญฺจิทํ เตน ภควตา ชานตา ปสฺสตา อรหตา สมฺมาสมฺพุทฺเธน สตฺตานํ นานาธิมุตฺติกตา สุปฺปฏิวิทิตา, อยํ หิ สุปฺปิโย ปริพฺพาชโก อเนกปริยาเยน พุทฺธสฺส อวณฺณํ ภาสติ, ธมฺมสฺส อวณฺณํ ภาสติ, สํฆสฺส อวณฺณํ ภาสติ, สุปฺปิยสฺส ปน ปริพฺพาชกสฺส อนฺเตวาสี พฺรหฺมทตฺโต มาณโว อเนกปริยาเยน พุทฺธสฺส วณฺณํ ภาสติ, ธมฺมสฺส วณฺณํ ภาสติ, สํฆสฺส วณฺณํ ภาสติ, อิติหเม อุโภ อาจริยนฺเตวาสี อญฺญมญฺญสฺส อุชุวิปจฺจนีกวาทา ภควนฺตํ ปิฏฺฐิโต \csromanfolio{3}ปิฏฺฐิโต อนุพนฺธา โหนฺติ ภิกฺขุสํฆญฺจา’ติ, อยํ โข โน ภนฺเต อนฺตรากถา วิปฺปกตา, อถ ภควา อนุปฺปตฺโต”ติ.}

\proseitem{5}{มมํ วา ภิกฺขเว ปเร อวณฺณํ ภาเสยฺยุํ, ธมฺมสฺส วา อวณฺณํ ภาเสยฺยุํ, สํฆสฺส วา อวณฺณํ ภาเสยฺยุํ, ตตฺร ตุเมฺหหิ น อาฆาโต น อปฺปจฺจโย น เจตโส อนภิรทฺธิ กรณียา.\csromanspacer{} มมํ วา ภิกฺขเว ปเร อวณฺณํ ภาเสยฺยุํ, ธมฺมสฺส วา อวณฺณํ ภาเสยฺยุํ, สํฆสฺส วา อวณฺณํ ภาเสยฺยุํ, ตตฺร เจ ตุเมฺห อสฺสถ กุปิตา วา อนตฺตมนา วา, ตุมฺหํเยวสฺส เตน อนฺตราโย.\csromanspacer{} มมํ วา ภิกฺขเว ปเร อวณฺณํ ภาเสยฺยุํ, ธมฺมสฺส วา อวณฺณํ ภาเสยฺยุํ, สํฆสฺส วา อวณฺณํ ภาเสยฺยุํ, ตตฺร เจ ตุเมฺห อสฺสถ กุปิตา วา อนตฺตมนา วา, อปิ นุ ตุเมฺห ปเรสํ สุภาสิตํ ทุพฺภาสิตํ อาชาเนยฺยาถาติ.\csromanspacer{} โนเหตํ ภนฺเต.\csromanspacer{} มมํ วา ภิกฺขเว ปเร อวณฺณํ ภาเสยฺยุํ, ธมฺมสฺส วา อวณฺณํ ภาเสยฺยุํ, สํฆสฺส วา อวณฺณํ ภาเสยฺยุํ, ตตฺร ตุเมฺหหิ อภูตํ อภูตโต นิพฺเพเฐตพฺพํ “อิติเปตํ อภูตํ, อิติเปตํ อตจฺฉํ, นตฺถิ เจตํ อเมฺหสุ, น จ ปเนตํ อเมฺหสุ สํวิชฺชตี”ติ.}

\proseitem{6}{มมํ วา ภิกฺขเว ปเร วณฺณํ ภาเสยฺยุํ, ธมฺมสฺส วา วณฺณํ ภาเสยฺยุํ, สํฆสฺส วา วณฺณํ ภาเสยฺยุํ, ตตฺร ตุเมฺหหิ น อานนฺโท น โสมนสฺสํ น เจตโส อุปฺปิลาวิตตฺตํ กรณียํ.\csromanspacer{} มมํ วา ภิกฺขเว ปเร วณฺณํ ภาเสยฺยุํ, ธมฺมสฺส วา วณฺณํ ภาเสยฺยุํ, สํฆสฺส วา วณฺณํ ภาเสยฺยุํ, ตตฺร เจ ตุเมฺห อสฺสถ อานนฺทิโน สุมนา อุปฺปิลาวิตา, ตุมฺหํเยวสฺส เตน อนฺตราโย.\csromanspacer{} มมํ วา ภิกฺขเว ปเร วณฺณํ ภาเสยฺยุํ, ธมฺมสฺส วา วณฺณํ ภาเสยฺยุํ, สํฆสฺส วา วณฺณํ ภาเสยฺยุํ, ตตฺร ตุเมฺหหิ ภูตํ ภูตโต ปฏิชานิตพฺพํ “อิติเปตํ ภูตํ, อิติเปตํ ตจฺฉํ, อตฺถิ เจตํ อเมฺหสุ, สํวิชฺชติ จ ปเนตํ อเมฺหสู”ติ.}

\csromanmatikaanchor{mtk.3}
\csromantocmarkref{section}{จูฬสีล}{mtk.3}
\bookmark[dest={mtk.3},level=1]{จูฬสีล}
\csromanheader{1}{\textbf{จูฬสีล}}
\proseitem{7}{อปฺปมตฺตกํ โข ปเนตํ ภิกฺขเว โอรมตฺตกํ สีลมตฺตกํ, เยน ปุถุชฺชโน ตถาคตสฺส วณฺณํ วทมาโน วเทยฺย, กตมญฺจ ตํ ภิกฺขเว อปฺปมตฺตกํ โอรมตฺตกํ สีลมตฺตกํ, เยน ปุถุชฺชโน}

\proseitem{8}{\csromanfolio{4}“ปาณาติปาตํ ปหาย ปาณาติปาตา ปฏิวิรโต สมโณ โคตโม นิหิตทณฺโฑ นิหิตสตฺโถ ลชฺชี ทยาปนฺโน สพฺพปาณภูต\-หิตา\-นุกมฺปี วิหรตี”ติ, อิติ วา หิ ภิกฺขเว ปุถุชฺชโน ตถาคตสฺส วณฺณํ วทมาโน}

\prose{“อทินฺนาทานํ ปหาย อทินฺนาทานา ปฏิวิรโต สมโณ โคตโม ทินฺนาทายี ทินฺนปาฏิกงฺขี อเถเนน สุจิภูเตน อตฺตนา วิหรตี”ติ, อิติ วา หิ ภิกฺขเว ปุถุชฺชโน ตถาคตสฺส วณฺณํ วทมาโน วเทยฺย.}

\prose{“อพฺรหฺมจริยํ ปหาย พฺรหฺมจารี สมโณ โคตโม อาราจารี\footnote{อนาจารี (ก)} วิรโต\footnote{ปฏิวิรโต (กตฺถจิ)} เมถุนา คามธมฺมา”ติ, อิติ วา หิ ภิกฺขเว ปุถุชฺชโน ตถาคตสฺส}

\proseitem{9}{“มุสาวาทํ ปหาย มุสาวาทา ปฏิวิรโต สมโณ โคตโม สจฺจวาที สจฺจสนฺโธ เถโต\footnote{เฐโต (สฺยา, กํ)} ปจฺจยิโก อวิสํวาทโก โลกสฺสา”ติ, อิติ วา หิ ภิกฺขเว ปุถุชฺชโน ตถาคตสฺส วณฺณํ วทมาโน วเทยฺย.}

\prose{“ปิสุณํ วาจํ ปหาย ปิสุณาย วาจาย ปฏิวิรโต สมโณ โคตโม, อิโต สุตฺวา น อมุตฺร อกฺขาตา อิเมสํ เภทาย, อมุตฺร วา สุตฺวา น อิเมสํ อกฺขาตา อมูสํ เภทาย, อิติ ภินฺนานํ วา สนฺธาตา, สหิตานํ วา อนุปฺปทาตา, สมคฺคาราโม สมคฺครโต สมคฺคนนฺที สมคฺคกรณิํ วาจํ ภาสิตา”ติ, อิติ วา หิ ภิกฺขเว ปุถุชฺชโน ตถาคตสฺส วณฺณํ}

\prose{“ผรุสํ วาจํ ปหาย ผรุสาย วาจาย ปฏิวิรโต สมโณ โคตโม, ยา สา วาจา เนลา กณฺณสุขา เปมนียา หทยงฺคมา โปรี พหุชนกนฺตา พหุชนมนาปา, ตถารูปิํ วาจํ ภาสิตา”ติ, อิติ วา หิ ภิกฺขเว ปุถุชฺชโน ตถาคตสฺส วณฺณํ วทมาโน วเทยฺย.}

\prose{“สมฺผปฺปลาปํ ปหาย สมฺผปฺปลาปา ปฏิวิรโต สมโณ โคตโม กาลวาที ภูตวาที อตฺถวาที ธมฺมวาที วินยวาที, นิธานวติํ วาจํ ภาสิตา กาเลน สาปเทสํ ปริยนฺตวติํ อตฺถสํหิตนฺ”ติ, อิติ วา หิ ภิกฺขเว ปุถุชฺชโน ตถาคตสฺส วณฺณํ วทมาโน วเทยฺย.}

\proseitem{10}{\csromanfolio{5}“พีชคาม\-ภูตคาม\-สมารมฺภา\footnote{สมารพฺภา (สี, ก)} ปฏิวิรโต สมโณ โคตโม”ติ, อิติ วา หิ ภิกฺขเว ฯเปฯ.}

\prose{เอกภตฺติโก สมโณ โคตโม รตฺตูปรโต วิรโต\footnote{ปฏิวิรโต (กตฺถจิ)} วิกาลโภชนา.}

\prose{นจฺจ\-คีต\-วาทิต\-วิสูก\-ทสฺสนา\footnote{นจฺจคีตวาทิตวิสุกทสฺสนา (ก)} ปฏิวิรโต สมโณ โคตโม.}

\prose{มาลาคนฺธวิเลปนธารณมณฺฑนวิภูสนฏฺฐานา ปฏิวิรโต สมโณ โคตโม.}

\prose{อุจฺจาสยนมหาสยนา ปฏิวิรโต สมโณ โคตโม.}

\prose{ชาตรูป\-รชต\-ปฏิคฺคหณา ปฏิวิรโต สมโณ โคตโม.}

\prose{อามกธญฺญปฏิคฺคหณา ปฏิวิรโต สมโณ โคตโม.}

\prose{อามกมํสปฏิคฺคหณา ปฏิวิรโต สมโณ โคตโม.}

\prose{อิตฺถิ\-กุมาริก\-ปฏิคฺคหณา ปฏิวิรโต สมโณ โคตโม.}

\prose{ทาสิทาสปฏิคฺคหณา ปฏิวิรโต สมโณ โคตโม.}

\prose{อเชฬกปฏิคฺคหณา ปฏิวิรโต สมโณ โคตโม.}

\prose{กุกฺกุฏสูกรปฏิคฺคหณา ปฏิวิรโต สมโณ โคตโม.}

\prose{หตฺถิควสฺสวฬวปฏิคฺคหณา ปฏิวิรโต สมโณ โคตโม.}

\prose{เขตฺตวตฺถุปฏิคฺคหณา ปฏิวิรโต สมโณ โคตโม.}

\prose{ทูเตยฺย\-ปหิณ\-คมนา\-นุโยคา ปฏิวิรโต สมโณ โคตโม.}

\prose{กยวิกฺกยา ปฏิวิรโต สมโณ โคตโม.}

\prose{ตุลา\-กูฏ\-กํส\-กูฏ\-มาน\-กูฏา ปฏิวิรโต สมโณ โคตโม.}

\prose{อุกฺโกฏน\-วญฺจน\-นิกติ\-สาจิโยคา\footnote{สาวิโยคา (สฺยา, กํ, ก)} ปฏิวิรโต สมโณ โคตโม.}

\prosewithcloser{เฉทนวธพนฺธนวิปราโมส-อาโลปสหสาการา ปฏิวิรโต สมโณ โคตโมติ, อิติ วา หิ ภิกฺขเว ปุถุชฺชโน ตถาคตสฺส วณฺณํ วทมาโน}{จูฬสีลํ นิฏฺฐิตํ.}
\csromanmatikaanchor{mtk.4}
\csromantocmarkref{section}{มชฺฌิมสีล}{mtk.4}
\bookmark[dest={mtk.4},level=1]{มชฺฌิมสีล}
\csromanheader{1}{\textbf{มชฺฌิมสีล}}
\proseitem{11}{\csromanfolio{6}ยถา วา ปเนเก โภนฺโต สมณพฺราหฺมณา สทฺธาเทยฺยานิ โภชนานิ ภุญฺชิตฺวา เต เอวรูปํ พีชคาม\-ภูตคาม\-สมารมฺภํ อนุยุตฺตา วิหรนฺติ, เสยฺยถิทํ\footnote{เสยฺยถีทํ (สี, สฺยา)}, มูลพีชํ ขนฺธพีชํ ผฬุพีชํ อคฺคพีชํ พีชพีชเมว ปญฺจมํ\footnote{ปญฺจมํ อิติ วา (สี, สฺยา, ก)}, อิติ เอวรูปา พีชคาม\-ภูตคาม\-สมารมฺภา ปฏิวิรโต สมโณ โคตโมติ, อิติ วา หิ ภิกฺขเว ปุถุชฺชโน ตถาคตสฺส วณฺณํ วทมาโน}

\proseitem{12}{ยถา วา ปเนเก โภนฺโต สมณพฺราหฺมณา สทฺธาเทยฺยานิ โภชนานิ ภุญฺชิตฺวา เต เอวรูปํ สนฺนิธิการปริโภคํ อนุยุตฺตา วิหรนฺติ, เสยฺยถิทํ, อนฺนสนฺนิธิํ ปานสนฺนิธิํ วตฺถสนฺนิธิํ ยานสนฺนิธิํ สยนสนฺนิธิํ คนฺธสนฺนิธิํ อามิสสนฺนิธิํ อิติ วา, อิติ เอวรูปา สนฺนิธิการปริโภคา ปฏิวิรโต สมโณ โคตโมติ, อิติ วา หิ ภิกฺขเว ปุถุชฺชโน ตถาคตสฺส วณฺณํ วทมาโน วเทยฺย.}

\proseitem{13}{ยถา วา ปเนเก โภนฺโต สมณพฺราหฺมณา สทฺธาเทยฺยานิ โภชนานิ ภุญฺชิตฺวา เต เอวรูปํ วิสูกทสฺสนํ อนุยุตฺตา วิหรนฺติ, เสยฺยถิทํ, นจฺจํ คีตํ วาทิตํ เปกฺขํ อกฺขานํ ปาณิสฺสรํ เวตาฬํ กุมฺภถูณํ\footnote{กุมฺภถูนํ (สฺยา, ก), กุมฺภถุณํ (สี)} โสภนกํ\footnote{โสภนฆรกํ (สี), โสภ นครกํ (สฺยา, กํ, อิ)} จณฺฑาลํ วํสํ โธวนํ หตฺถิยุทฺธํ อสฺสยุทฺธํ มหิํสยุทฺธํ\footnote{มหิสยุทฺธํ (สี, สฺยา, กํ, อิ)} อุสภยุทฺธํ อชยุทฺธํ เมณฺฑยุทฺธํ กุกฺกุฏยุทฺธํ วฏฺฏกยุทฺธํ ทณฺฑยุทฺธํ มุฏฺฐิยุทฺธํ นิพฺพุทฺธํ อุยฺโยธิกํ พลคฺคํ เสนาพฺยูหํ อนีกทสฺสนํ อิติ วา, อิติ เอวรูปา วิสูกทสฺสนา ปฏิวิรโต สมโณ โคตโมติ, อิติ วา หิ ภิกฺขเว ปุถุชฺชโน ตถาคตสฺส วณฺณํ วทมาโน}

\proseitem{14}{ยถา วา ปเนเก โภนฺโต สมณพฺราหฺมณา สทฺธาเทยฺยานิ โภชนานิ ภุญฺชิตฺวา เต เอวรูปํ ชูตปฺปมาทฏฺฐานานุโยคํ อนุยุตฺตา วิหรนฺติ, เสยฺยถิทํ, อฏฺฐปทํ ทสปทํ อากาสํ ปริหารปถํ สนฺติกํ ขลิกํ ฆฏิกํ สลากหตฺถํ อกฺขํ ปงฺคจีรํ วงฺกกํ โมกฺขจิกํ จิงฺคุลิกํ\footnote{จิงฺคุลกํ (กสี)} ปตฺตาฬฺหกํ รถกํ ธนุกํ อกฺขริกํ มเนสิกํ ยถาวชฺชํ \csromanfolio{7}อิติ วา, อิติ เอวรูปา ชูตปฺปมาทฏฺฐานานุโยคา ปฏิวิรโต สมโณ โคตโมติ, อิติ วา หิ ภิกฺขเว ปุถุชฺชโน ตถาคตสฺส วณฺณํ วทมาโน วเทยฺย.}

\proseitem{15}{ยถา วา ปเนเก โภนฺโต สมณพฺราหฺมณา สทฺธาเทยฺยานิ โภชนานิ ภุญฺชิตฺวา เต เอวรูปํ อุจฺจาสยนมหาสยนํ อนุยุตฺตา วิหรนฺติ, เสยฺยถิทํ, อาสนฺทิํ ปลฺลงฺกํ โคนกํ จิตฺตกํ ปฏิกํ ปฏลิกํ ตูลิกํ วิกติกํ อุทฺทโลมิํ เอกนฺตโลมิํ กฏฺฏิสฺสํ โกเสยฺยํ กุตฺตกํ หตฺถตฺถรํ อสฺสตฺถรํ รถตฺถรํ\footnote{หตฺถตฺถรณํ อสฺสตฺถรณํ รถตฺถรณํ (สี, ก, อิ) 2. มาลาวิเลปนํ (สี, สฺยา, กํ, อิ)} อชินปฺปเวณิํ กทลิ\-มิค\-ปวร\-ปจฺจตฺถรณํ ส-อุตฺตรจฺฉทํ อุภโตโลหิตกูปธานํ อิติ วา, อิติ เอวรูปา อุจฺจาสยนมหาสยนา ปฏิวิรโต สมโณ โคตโมติ, อิติ วา หิ ภิกฺขเว ปุถุชฺชโน ตถาคตสฺส}

\proseitem{16}{ยถา วา ปเนเก โภนฺโต สมณพฺราหฺมณา สทฺธาเทยฺยานิ โภชนานิ ภุญฺชิตฺวา เต เอวรูปํ มณฺฑน\-วิภูสนฏฺฐานา\-นุโยคํ อนุยุตฺตา วิหรนฺติ, เสยฺยถิทํ, อุจฺฉาทนํ ปริมทฺทนํ นฺหาปนํ สมฺพาหนํ อาทาสํ อญฺชนํ มาลาคนฺธวิเลปนํ2 มุขจุณฺณํ มุขเลปนํ หตฺถพนฺธํ สิขาพนฺธํ ทณฺฑํ นาฬิกํ อสิํ\footnote{ขคฺคํ (สี, อิ), อสิํ ขคฺคํ (สฺยา, กํ) 4. อิตฺถิกถํ ปุริสกถํ (สฺยา, กํ, ก)} ฉตฺตํ จิตฺรุปาหนํ อุณฺหีสํ มณิํ วาลพีชนิํ โอทาตานิ วตฺถานิ ทีฆทสานิ อิติ วา, อิติ เอวรูปา มณฺฑน\-วิภูสนฏฺฐานา\-นุโยคา ปฏิวิรโต สมโณ โคตโมติ, อิติ วา หิ ภิกฺขเว ปุถุชฺชโน ตถาคตสฺส วณฺณํ วทมาโน วเทยฺย.}

\proseitem{17}{ยถา วา ปเนเก โภนฺโต สมณพฺราหฺมณา สทฺธาเทยฺยานิ โภชนานิ ภุญฺชิตฺวา เต เอวรูปํ ติรจฺฉานกถํ อนุยุตฺตา วิหรนฺติ, เสยฺยถิทํ, ราชกถํ โจรกถํ มหามตฺตกถํ เสนากถํ ภยกถํ ยุทฺธกถํ อนฺนกถํ ปานกถํ วตฺถกถํ สยนกถํ มาลากถํ คนฺธกถํ ญาติกถํ ยานกถํ คามกถํ นิคมกถํ นครกถํ ชนปทกถํ อิตฺถิกถํ4 สูรกถํ วิสิขากถํ กุมฺภฏฺฐานกถํ ปุพฺพเปตกถํ นานตฺตกถํ โลกกฺขายิกํ สมุทฺทกฺขายิกํ อิติภวาภวกถํ อิติ วา, อิติ เอวรูปาย ติรจฺฉานกถาย ปฏิวิรโต สมโณ โคตโมติ, อิติ วา หิ ภิกฺขเว ปุถุชฺชโน ตถาคตสฺส วณฺณํ วทมาโน วเทยฺย.}

\proseitem{18}{\csromanfolio{8}ยถา วา ปเนเก โภนฺโต สมณพฺราหฺมณา สทฺธาเทยฺยานิ โภชนานิ ภุญฺชิตฺวา เต เอวรูปํ วิคฺคาหิกกถํ อนุยุตฺตา วิหรนฺติ, เสยฺยถิทํ, น ตฺวํ อิมํ ธมฺมวินยํ อาชานาสิ, อหํ อิมํ ธมฺมวินยํ อาชานามิ, กิํ ตฺวํ อิมํ ธมฺมวินยํ อาชานิสฺสสิ, มิจฺฉา ปฏิปนฺโน ตฺวมสิ, อหมสฺมิ สมฺมา ปฏิปนฺโน, สหิตํ เม, อสหิตํ เต, ปุเรวจนียํ ปจฺฉา อวจ, ปจฺฉาวจนียํ ปุเร อวจ, อธิจิณฺณํ เต วิปราวตฺตํ, อาโรปิโต เต วาโท, นิคฺคหิโต ตฺวมสิ, จร วาทปฺปโมกฺขาย, นิพฺเพเฐหิ วา สเจ ปโหสีติ อิติ วา, อิติ เอวรูปาย วิคฺคาหิกกถาย ปฏิวิรโต สมโณ โคตโมติ, อิติ วา หิ ภิกฺขเว ปุถุชฺชโน ตถาคตสฺส วณฺณํ วทมาโน}

\proseitem{19}{ยถา วา ปเนเก โภนฺโต สมณพฺราหฺมณา สทฺธาเทยฺยานิ โภชนานิ ภุญฺชิตฺวา เต เอวรูปํ ทูเตยฺย\-ปหิณ\-คมนา\-นุโยคํ อนุยุตฺตา วิหรนฺติ, เสยฺยถิทํ, รญฺญํ, ราชมหามตฺตานํ, ขตฺติยานํ, พฺราหฺมณานํ, คหปติกานํ, กุมารานํ “อิธ คจฺฉ, อมุตฺราคจฺฉ, อิทํ หร, อมุตฺร อิทํ อาหรา”ติ อิติ วา, อิติ เอวรูปา ทูเตยฺย\-ปหิณ\-คมนา\-นุโยคา ปฏิวิรโต สมโณ โคตโมติ, อิติ วา หิ ภิกฺขเว ปุถุชฺชโน ตถาคตสฺส วณฺณํ}

\proseitemwithcloserruled{20}{ยถา วา ปเนเก โภนฺโต สมณพฺราหฺมณา สทฺธาเทยฺยานิ โภชนานิ ภุญฺชิตฺวา เต กุหกา จ โหนฺติ, ลปกา จ เนมิตฺติกา จ นิปฺเปสิกา จ, ลาเภน ลาภํ นิชิคีสิตาโร จ\footnote{ลาเภน ลาภํ นิชิคิํสิตาโร (สี, สฺยา), ลาเภน จ ลาภํ นิชิคีสิตาโร (อิ)} อิติ\footnote{อิติ วา, อิติ (สฺยา, กํ, ก)} เอวรูปา กุหนลปนา ปฏิวิรโต สมโณ โคตโมติ, อิติ วา หิ ภิกฺขเว ปุถุชฺชโน ตถาคตสฺส วณฺณํ วทมาโน}{มชฺฌิมสีลํ นิฏฺฐิตํ.}
\csromanmatikaanchor{mtk.5}
\csromantocmarkref{section}{มหาสีล}{mtk.5}
\bookmark[dest={mtk.5},level=1]{มหาสีล}
\csromanheader{1}{\textbf{มหาสีล}}
\proseitem{21}{ยถา วา ปเนเก โภนฺโต สมณพฺราหฺมณา สทฺธาเทยฺยานิ โภชนานิ ภุญฺชิตฺวา เต เอวรูปาย ติรจฺฉานวิชฺชาย มิจฺฉาชีเวน ชีวิตํ \csromanfolio{9}กปฺเปนฺติ, เสยฺยถิทํ, องฺคํ นิมิตฺตํ อุปฺปาตํ สุปินํ ลกฺขณํ มูสิกจฺฉินฺนํ อคฺคิโหมํ ทพฺพิโหมํ ถุสโหมํ กณโหมํ ตณฺฑุลโหมํ สปฺปิโหมํ เตลโหมํ มุขโหมํ โลหิตโหมํ องฺควิชฺชา วตฺถุวิชฺชา ขตฺตวิชฺชา\footnote{เขตฺตวิชฺชา (พหูสุ)} สิววิชฺชา ภูตวิชฺชา ภูริวิชฺชา อหิวิชฺชา วิสวิชฺชา วิจฺฉิกวิชฺชา มูสิกวิชฺชา สกุณวิชฺชา วายสวิชฺชา ปกฺกชฺฌานํ สรปริตฺตาณํ มิคจกฺกํ อิติ วา, อิติ เอวรูปาย ติรจฺฉานวิชฺชาย มิจฺฉาชีวา ปฏิวิรโต สมโณ โคตโมติ, อิติ วา หิ ภิกฺขเว ปุถุชฺชโน}

\proseitem{22}{ยถา วา ปเนเก โภนฺโต สมณพฺราหฺมณา สทฺธาเทยฺยานิ โภชนานิ ภุญฺชิตฺวา เต เอวรูปาย ติรจฺฉานวิชฺชาย มิจฺฉาชีเวน ชีวิตํ กปฺเปนฺติ, เสยฺยถิทํ, มณิลกฺขณํ วตฺถลกฺขณํ ทณฺฑลกฺขณํ สตฺถลกฺขณํ อสิลกฺขณํ อุสุลกฺขณํ ธนุลกฺขณํ อาวุธลกฺขณํ อิตฺถิลกฺขณํ ปุริสลกฺขณํ กุมารลกฺขณํ กุมาริลกฺขณํ ทาสลกฺขณํ ทาสิลกฺขณํ หตฺถิลกฺขณํ อสฺสลกฺขณํ มหิํสลกฺขณํ\footnote{มหิสลกฺขณํ (สี, สฺยา, กํ, อิ)} อุสภลกฺขณํ โคลกฺขณํ อชลกฺขณํ เมณฺฑลกฺขณํ กุกฺกุฏลกฺขณํ วฏฺฏกลกฺขณํ โคธาลกฺขณํ กณฺณิกาลกฺขณํ กจฺฉปลกฺขณํ มิคลกฺขณํ อิติ วา, อิติ เอวรูปาย ติรจฺฉานวิชฺชาย มิจฺฉาชีวา ปฏิวิรโต สมโณ โคตโมติ, อิติ วา หิ ภิกฺขเว ปุถุชฺชโน ตถาคตสฺส วณฺณํ วทมาโน วเทยฺย.}

\proseitem{23}{ยถา วา ปเนเก โภนฺโต สมณพฺราหฺมณา สทฺธาเทยฺยานิ โภชนานิ ภุญฺชิตฺวา เต เอวรูปาย ติรจฺฉานวิชฺชาย มิจฺฉาชีเวน ชีวิตํ กปฺเปนฺติ, เสยฺยถิทํ, รญฺญํ นิยฺยานํ ภวิสฺสติ, รญฺญํ อนิยฺยานํ ภวิสฺสติ, อพฺภนฺตรานํ รญฺญํ อุปยานํ ภวิสฺสติ, พาหิรานํ รญฺญํ อปยานํ ภวิสฺสติ, พาหิรานํ รญฺญํ อุปยานํ ภวิสฺสติ, อพฺภนฺตรานํ รญฺญํ อปยานํ ภวิสฺสติ, อพฺภนฺตรานํ รญฺญํ ชโย ภวิสฺสติ, พาหิรานํ รญฺญํ ปราชโย ภวิสฺสติ, พาหิรานํ รญฺญํ ชโย ภวิสฺสติ, อพฺภานฺตรานํ รญฺญํ ปราชโย ภวิสฺสติ, อิติ อิมสฺส ชโย ภวิสฺสติ, อิมสฺส ปราชโย ภวิสฺสติ อิติ วา, อิติ เอวรูปาย ติรจฺฉานวิชฺชาย มิจฺฉาชีวา ปฏิวิรโต สมโณ โคตโมติ, อิติ วา หิ ภิกฺขเว ปุถุชฺชโน ตถาคตสฺส}

\proseitem{24}{\csromanfolio{10}ยถา วา ปเนเก โภนฺโต สมณพฺราหฺมณา สทฺธาเทยฺยานิ โภชนานิ ภุญฺชิตฺวา เต เอวรูปาย ติรจฺฉานวิชฺชาย มิจฺฉาชีเวน ชีวิตํ กปฺเปนฺติ, เสยฺยถิทํ, จนฺทคฺคาโห ภวิสฺสติ, สูริยคฺคาโห\footnote{สุริยคฺคาโห (สี, สฺยา, กํ, อิ)} ภวิสฺสติ, นกฺขตฺตคฺคาโห ภวิสฺสติ, จนฺทิมสูริยานํ ปถคมนํ ภวิสฺสติ, จนฺทิมสูริยานํ อุปฺปถคมนํ ภวิสฺสติ, นกฺขตฺตานํ ปถคมนํ ภวิสฺสติ, นกฺขตฺตานํ อุปฺปถคมนํ ภวิสฺสติ, อุกฺกาปาโต ภวิสฺสติ, ทิสาฑาโห ภวิสฺสติ, ภูมิจาโล ภวิสฺสติ, เทวทุทฺรภิ\footnote{เทวทุนฺทุภิ (สฺยา, กํ, อิ)} ภวิสฺสติ, จนฺทิมสูริยนกฺขตฺตานํ อุคฺคมนํ โอคมนํ สํกิเลสํ โวทานํ ภวิสฺสติ, เอวํวิปาโก จนฺทคฺคาโห ภวิสฺสติ, เอวํวิปาโก สูริยคฺคาโห ภวิสฺสติ, เอวํวิปาโก นกฺขตฺตคฺคาโห ภวิสฺสติ, เอวํวิปากํ จนฺทิมสูริยานํ ปถคมนํ ภวิสฺสติ, เอวํวิปากํ จนฺทิมสูริยานํ อุปฺปถคมนํ ภวิสฺสติ, วิปากํ นกฺขตฺตานํ ปถคมนํ ภวิสฺสติ, เอวํวิปากํ นกฺขตฺตานํ อุปฺปถคมนํ ภวิสฺสติ, เอวํวิปาโก อุกฺกาปาโต ภวิสฺสติ, เอวํวิปาโก ทิสาฑาโห ภวิสฺสติ, เอวํวิปาโก ภูมิจาโล ภวิสฺสติ, เอวํวิปาโก เทวทุทฺรภิ ภวิสฺสติ, เอวํวิปากํ จนฺทิมสูริยนกฺขตฺตานํ อุคฺคมนํ โอคมนํ สํกิเลสํ โวทานํ ภวิสฺสติ อิติ วา, อิติ เอวรูปาย ติรจฺฉานวิชฺชาย มิจฺฉาชีวา ปฏิวิรโต สมโณ โคตโมติ, อิติ วา หิ ภิกฺขเว ปุถุชฺชโน ตถาคตสฺส วณฺณํ วทมาโน วเทยฺย.}

\proseitem{25}{ยถา วา ปเนเก โภนฺโต สมณพฺราหฺมณา สทฺธาเทยฺยานิ โภชนานิ ภุญฺชิตฺวา เต เอวรูปาย ติรจฺฉานวิชฺชาย มิจฺฉาชีเวน ชีวิตํ กปฺเปนฺติ, เสยฺยถิทํ, สุวุฏฺฐิกา ภวิสฺสติ, ทุพฺพุฏฺฐิกา ภวิสฺสติ, สุภิกฺขํ ภวิสฺสติ, ทุพฺภิกฺขํ ภวิสฺสติ, เขมํ ภวิสฺสติ, ภยํ ภวิสฺสติ, โรโค ภวิสฺสติ, อาโรคฺยํ ภวิสฺสติ, มุทฺทา, คณนา, สงฺขานํ, กาเวยฺยํ, โลกายตํ อิติ วา, อิติ เอวรูปาย ติรจฺฉานวิชฺชาย มิจฺฉาชีวา ปฏิวิรโต สมโณ โคตโมติ, อิติ วา หิ ภิกฺขเว ปุถุชฺชโน ตถาคตสฺส วณฺณํ}

\proseitem{26}{ยถา วา ปเนเก โภนฺโต สมณพฺราหฺมณา สทฺธาเทยฺยานิ โภชนานิ ภุญฺชิตฺวา เต เอวรูปาย ติรจฺฉานวิชฺชาย มิจฺฉาชีเวน ชีวิตํ \csromanfolio{11}กปฺเปนฺติ, เสยฺยถิทํ, อาวาหนํ วิวาหนํ สํวรณํ วิวรณํ สํกิรณํ วิกิรณํ สุภคกรณํ ทุพฺภคกรณํ วิรุทฺธคพฺภกรณํ ชิวฺหานิพนฺธนํ หนุสํหนนํ หตฺถาภิชปฺปนํ หนุชปฺปนํ กณฺณชปฺปนํ อาทาสปญฺหํ กุมาริกปญฺหํ เทวปญฺหํ อาทิจฺจุปฏฺฐานํ มหตุปฏฺฐานํ อพฺภุชฺชลนํ สิริวฺหายนํ อิติ วา, อิติ เอวรูปาย ติรจฺฉานวิชฺชาย มิจฺฉาชีวา ปฏิวิรโต สมโณ โคตโมติ, อิติ วา หิ ภิกฺขเว ปุถุชฺชโน}

\proseitem{27}{ยถา วา ปเนเก โภนฺโต สมณพฺราหฺมณา สทฺธาเทยฺยานิ โภชนานิ ภุญฺชิตฺวา เต เอวรูปาย ติรจฺฉานวิชฺชาย มิจฺฉาชีเวน ชีวิตํ กปฺเปนฺติ, เสยฺยถิทํ, สนฺติกมฺมํ ปณิธิกมฺมํ ภูตกมฺมํ ภูริกมฺมํ วสฺสกมฺมํ โวสฺสกมฺมํ วตฺถุกมฺมํ วตฺถุปริกมฺมํ อาจมนํ นฺหาปนํ ชุหนํ วมนํ วิเรจนํ อุทฺธํวิเรจนํ อโธวิเรจนํ สีสวิเรจนํ กณฺณเตลํ เนตฺตตปฺปนํ นตฺถุกมฺมํ อญฺชนํ ปจฺจญฺชนํ สาลากิยํ สลฺลกตฺติยํ ทารกติกิจฺฉา, มูลเภสชฺชานํ อนุปฺปทานํ, โอสธีนํ ปฏิโมกฺโข อิติ วา, อิติ เอวรูปาย ติรจฺฉานวิชฺชาย มิจฺฉาชีวา ปฏิวิรโต สมโณ โคตโมติ, อิติ วา หิ ภิกฺขเว ปุถุชฺชโน ตถาคตสฺส วณฺณํ วทมาโน วเทยฺย.}

\prosewithcloserruled{อิทํ โข ภิกฺขเว อปฺปมตฺตกํ โอรมตฺตกํ สีลมตฺตกํ, เยน ปุถุชฺชโน ตถาคตสฺส วณฺณํ วทมาโน วเทยฺย.}{มหาสีลํ นิฏฺฐิตํ.}
\csromanmatikaanchor{mtk.6}
\csromantocmarkref{section}{ปุพฺพนฺตกปฺปิก}{mtk.6}
\bookmark[dest={mtk.6},level=1]{ปุพฺพนฺตกปฺปิก}
\csromanheader{1}{\textbf{ปุพฺพนฺตกปฺปิก}}
\proseitem{28}{อตฺถิ ภิกฺขเว อญฺเญว ธมฺมา คมฺภีรา ทุทฺทสา ทุรนุโพธา สนฺตา ปณีตา อตกฺกาวจรา นิปุณา ปณฺฑิตเวทนียา, เย ตถาคโต สยํ อภิญฺญา สจฺฉิกตฺวา ปเวเทติ, เยหิ ตถาคตสฺส ยถาภุจฺจํ วณฺณํ สมฺมา วทมานา วเทยฺยุํ.\csromanspacer{} กตเม จ เต ภิกฺขเว ธมฺมา คมฺภีรา ทุทฺทสา ทุรนุโพธา สนฺตา ปณีตา อตกฺกาวจรา นิปุณา ปณฺฑิตเวทนียา, เย ตถาคโต สยํ อภิญฺญา สจฺฉิกตฺวา ปเวเทติ, เยหิ ตถาคตสฺส ยถาภุจฺจํ วณฺณํ สมฺมา วทมานา วเทยฺยุํ.}

\proseitem{29}{\csromanfolio{12}สนฺติ ภิกฺขเว เอเก สมณพฺราหฺมณา ปุพฺพนฺตกปฺปิกา ปุพฺพนฺตานุทิฏฺฐิโน ปุพฺพนฺตํ อารพฺภ อเนกวิหิตานิ อธิมุตฺติปทานิ\footnote{อธิวุตฺติปทานิ (สี, อิ)} อภิวทนฺติ อฏฺฐารสหิ วตฺถูหิ, เต จ โภนฺโต สมณพฺราหฺมณา กิมาคมฺม กิมารพฺภ ปุพฺพนฺตกปฺปิกา ปุพฺพนฺตานุทิฏฺฐิโน ปุพฺพนฺตํ อารพฺภ อเนกวิหิตานิ อธิมุตฺติปทานิ อภิวทนฺติ อฏฺฐารสหิ วตฺถูหิ?}

\csromanmatikaanchor{mtk.7}
\csromantocmarkref{section}{สสฺสตวาท}{mtk.7}
\bookmark[dest={mtk.7},level=1]{สสฺสตวาท}
\csromanheader{1}{\textbf{สสฺสตวาท}}
\proseitem{30}{สนฺติ ภิกฺขเว เอเก สมณพฺราหฺมณา สสฺสตวาทา สสฺสตํ อตฺตานญฺจ โลกญฺจ ปญฺญเปนฺติ จตูหิ วตฺถูหิ, เต จ โภนฺโต สมณพฺราหฺมณา กิมาคมฺม กิมารพฺภ สสฺสตวาทา สสฺสตํ อตฺตานญฺจ โลกญฺจ ปญฺญเปนฺติ จตูหิ วตฺถูหิ?}

\proseitem{31}{อิธ ภิกฺขเว เอกจฺโจ สมโณ วา พฺราหฺมโณ วา อาตปฺปมนฺวาย ปธานมนฺวาย อนุโยคมนฺวาย อปฺปมาทมนฺวาย สมฺมา\-มนสิการม\-นฺวาย ตถารูปํ เจโตสมาธิํ ผุสติ, ยถาสมาหิเต จิตฺเต ( )\footnote{(ปริสุทฺเธ ปริโยทาเต อนงฺคเณ วิคตูปกฺกิเลเส) (สฺยา, ก)} อเนกวิหิตํ ปุพฺเพนิวาสํ อนุสฺสรติ.\csromanspacer{} เสยฺยถิทํ, เอกํปิ ชาติํ เทฺวปิ ชาติโย ติสฺโสปิ ชาติโย จตสฺโสปิ ชาติโย ปญฺจปิ ชาติโย ทสปิ ชาติโย วีสํปิ ชาติโย ติํสํปิ ชาติโย จตฺตาลีสํปิ ชาติโย ปญฺญาสํปิ ชาติโย ชาติสตํปิ ชาติสหสฺสํปิ ชาติสตสหสฺสํปิ อเนกานิปิ ชาติสตานิ อเนกานิปิ ชาติสหสฺสานิ อเนกานิปิ ชาติสตสหสฺสานิ, “อมุตฺราสิํ เอวํนาโม เอวํโคตฺโต เอวํวณฺโณ เอวมาหาโร เอวํสุขทุกฺขปฺปฏิสํเวที เอวมายุปริยนฺโต, โส ตโต จุโต อมุตฺร อุทปาทิํ, ตตฺราปาสิํ เอวํนาโม เอวํโคตฺโต เอวํวณฺโณ เอวมาหาโร เอวํสุขทุกฺขปฺปฏิสํเวที เอวมายุปริยนฺโต, โส ตโต จุโต อิธูปปนฺโน”ติ, อิติ สาการํ ส-อุทฺเทสํ อเนกวิหิตํ ปุพฺเพนิวาสํ อนุสฺสรติ.}

\prose{โส เอวมาห “สสฺสโต อตฺตา จ โลโก จ วญฺโฌ กูฏฏฺโฐ เอสิกฏฺฐายิฏฺฐิโต, เต จ สตฺตา สนฺธาวนฺติ สํสรนฺติ จวนฺติ อุปปชฺชนฺติ, อตฺถิเตฺวว สสฺสติสมํ.\csromanspacer{} ตํ กิสฺส เหตุ, อหํ หิ อาตปฺปมนฺวาย \csromanfolio{13}ปธานมนฺวาย อนุโยคมนฺวาย อปฺปมาทมนฺวาย สมฺมา\-มนสิการม\-นฺวาย ตถารูปํ เจโตสมาธิํ ผุสามิ, ยถาสมาหิเต จิตฺเต อเนกวิหิตํ ปุพฺเพนิวาสํ อนุสฺสรามิ.\csromanspacer{} เสยฺยถิทํ, เอกํปิ ชาติํ เทฺวปิ ชาติโย ติสฺโสปิ ชาติโย จตสฺโสปิ ชาติโย ปญฺจปิ ชาติโย ทสปิ ชาติโย วีสํปิ ชาติโย ติํสํปิ ชาติโย จตฺตาลีสํปิ ชาติโย ปญฺญาสํปิ ชาติโย ชาติสตํปิ ชาติสหสฺสํปิ ชาติสตสหสฺสํปิ อเนกานิปิ ชาติสตานิ อเนกานิปิ ชาติสหสฺสานิ อเนกานิปิ ชาติสตสหสฺสานิ, “อมุตฺราสิํ เอวํนาโม เอวํโคตฺโต เอวํวณฺโณ เอวมาหาโร เอวํสุขทุกฺขปฺปฏิสํเวที เอวมายุปริยนฺโต, โส ตโต จุโต อมุตฺร อุทปาทิํ, ตตฺราปาสิํ เอวํนาโม เอวํโคตฺโต เอวํวณฺโณ เอวมาหาโร เอวํสุขทุกฺขปฺปฏิสํเวที เอวมายุปริยนฺโต, โส ตโต จุโต อิธูปปนฺโน”ติ, อิติ สาการํ ส-อุทฺเทสํ อเนกวิหิตํ ปุพฺเพนิวาสํ อนุสฺสรามิ.}

\prose{อิมินามหํ เอตํ ชานามิ “ยถา สสฺสโต อตฺตา จ โลโก จ วญฺโฌ กูฏฏฺโฐ เอสิกฏฺฐายิฏฺฐิโต, เต จ สตฺตา สนฺธาวนฺติ สํสรนฺติ จวนฺติ อุปปชฺชนฺติ, อตฺถิเตฺวว สสฺสติสมนฺ”ติ.\csromanspacer{} อิทํ ภิกฺขเว ปฐมํ ฐานํ, ยํ อาคมฺม ยํ อารพฺภ เอเก สมณพฺราหฺมณา สสฺสตวาทา สสฺสตํ อตฺตานญฺจ โลกญฺจ ปญฺญเปนฺติ. (1)}

\proseitem{32}{ทุติเย จ โภนฺโต สมณพฺราหฺมณา กิมาคมฺม กิมารพฺภ สสฺสตวาทา สสฺสตํ อตฺตานญฺจ โลกญฺจ ปญฺญเปนฺติ? อิธ ภิกฺขเว เอกจฺโจ สมโณ วา พฺราหฺมโณ วา อาตปฺปมนฺวาย ปธานมนฺวาย อนุโยคมนฺวาย อปฺปมาทมนฺวาย สมฺมา\-มนสิการม\-นฺวาย ตถารูปํ เจโตสมาธิํ ผุสติ, ยถาสมาหิเต จิตฺเต อเนกวิหิตํ ปุพฺเพนิวาสํ อนุสฺสรติ.\csromanspacer{} เสยฺยถิทํ, เอกํปิ สํวฏฺฏวิวฏฺฏํ เทฺวปิ สํวฏฺฏวิวฏฺฏานิ ตีณิปิ สํวฏฺฏวิวฏฺฏานิ จตฺตาริปิ สํวฏฺฏวิวฏฺฏานิ ปญฺจปิ สํวฏฺฏวิวฏฺฏานิ ทสปิ สํวฏฺฏวิวฏฺฏานิ, “อมุตฺราสิํ เอวํนาโม เอวํโคตฺโต เอวํวณฺโณ เอวมาหาโร เอวํสุขทุกฺขปฺปฏิสํเวที เอวมายุปริยนฺโต, โส ตโต จุโต อมุตฺร อุทปาทิํ, ตตฺราปาสิํ เอวํนาโม เอวํโคตฺโต เอวํวณฺโณ เอวมาหาโร เอวํสุขทุกฺขปฺปฏิสํเวที เอวมายุปริยนฺโต, โส ตโต จุโต อิธูปปนฺโน”ติ, อิติ สาการํ ส-อุทฺเทสํ อเนกวิหิตํ ปุพฺเพนิวาสํ อนุสฺสรติ.}

\prose{\csromanfolio{14}โส เอวมาห “สสฺสโต อตฺตา จ โลโก จ วญฺโฌ กูฏฏฺโฐ เอสิกฏฺฐายิฏฺฐิโต, เต จ สตฺตา สนฺธาวนฺติ สํสรนฺติ จวนฺติ อุปปชฺชนฺติ, อตฺถิเตฺวว สสฺสติสมํ, ตํ กิสฺส เหตุ, อหํ หิ อาตปฺปมนฺวาย ปธานมนฺวาย อนุโยคมนฺวาย อปฺปมาทมนฺวาย สมฺมา\-มนสิการม\-นฺวาย ตถารูปํ เจโตสมาธิํผุสามิ, ยถาสมาหิเต จิตฺเต อเนกวิหิตํ ปุพฺเพนิวาสํ อนุสฺสารามิ.\csromanspacer{} เสยฺยถิทํ, เอกํปิ สํวฏฺฏวิวฏฺฏํ เทฺวปิ สํวฏฺฏวิวฏฺฏานิ ตีณิปิ สํวฏฺฏวิวฏฺฏานิ จตฺตาริปิ สํวฏฺฏวิวฏฺฏานิ ปญฺจปิ สํวฏฺฏวิวฏฺฏานิ ทสปิ สํวฏฺฏวิวฏฺฏานิ, “อมุตฺราสิํ เอวํนาโม เอวํโคตฺโต เอวํวณฺโณ เอวมาหาโร เอวํสุขทุกฺขปฺปฏิสํเวที เอวมายุปริยนฺโต, โส ตโต จุโต อมุตฺร อุทปาทิํ, ตตฺราปาสิํ เอวํนาโม เอวํโคตฺโต เอวํวณฺโณ เอวมาหาโร เอวํสุขทุกฺขปฺปฏิสํเวที เอวมายุปริยนฺโต, โส ตโต จุโต อิธูปปนฺโน”ติ, อิติ สาการํ ส-อุทฺเทสํ อเนกวิหิตํ ปุพฺเพนิวาสํ อนุสฺสรามิ.}

\prose{อิมินามหํ เอตํ ชานามิ “ยถา สสฺสโต อตฺตา จ โลโก จ วญฺโฌ กูฏฏฺโฐ เอสิกฏฺฐายิฏฺฐิโต, เต จ สตฺตา สนฺธาวนฺติ สํสรนฺติ จวนฺติ อุปปชฺชนฺติ, อตฺถิเตฺวว สสฺสติสมนฺ”ติ.\csromanspacer{} อิทํ ภิกฺขเว ทุติยํ ฐานํ, ยํ อาคมฺม ยํ อารพฺภ เอเก สมณพฺราหฺมณา สสฺสตวาทา สสฺสตํ อตฺตานญฺจ โลกญฺจ ปญฺญเปนฺติ. (2)}

\proseitem{33}{ตติเย จ โภนฺโต สมณพฺราหฺมณา กิมาคมฺม กิมารพฺภ สสฺสตวาทา สสฺสตํ อตฺตานญฺจ โลกญฺจ ปญฺญเปนฺติ? อิธ ภิกฺขเว เอกจฺโจ สมโณ วา พฺราหฺมโณ วา อาตปฺปมนฺวาย ปธานมนฺวาย อนุโยคมนฺวาย อปฺปมาทมนฺวาย สมฺมา\-มนสิการม\-นฺวาย ตถารูปํ เจโตสมาธิํ ผุสติ, ยถาสมาหิเต จิตฺเต อเนกวิหิตํ ปุพฺเพนิวาสํ อนุสฺสรติ.\csromanspacer{} เสยฺยถิทํ, ทสปิ สํวฏฺฏวิวฏฺฏานิ วีสํปิ สํวฏฺฏวิวฏฺฏานิ ติํสํปิ สํวฏฺฏวิวฏฺฏานิ จตฺตาลีสํปิ สํวฏฺฏวิวฏฺฏานิ, “อมุตฺราสิํ เอวํนาโม เอวํโคตฺโต เอวํวณฺโณ เอวมาหาโร เอวํสุขทุกฺขปฺปฏิสํเวที เอวมายุปริยนฺโต, โส ตโต จุโต อมุตฺร อุทปาทิํ, ตตฺราปาสิํ เอวํนาโม เอวํโคตฺโต เอวํวณฺโณ เอวมาหาโร เอวํสุขทุกฺขปฺปฏิสํเวที เอวมายุปริยนฺโต, โส ตโต จุโต อิธูปปนฺโน”ติ, อิติ สาการํ ส-อุทฺเทสํ อเนกวิหิตํ ปุพฺเพนิวาสํ อนุสฺสรติ.}

\prose{\csromanfolio{15}โส เอวมาห “สสฺสโต อตฺตา จ โลโก จ วญฺโฌ กูฏฏฺโฐ เอสิกฏฺฐายิฏฺฐิโต, เต จ สตฺตา สนฺธาวนฺติ สํสรนฺติ จวนฺติ อุปปชฺชนฺติ, อตฺถิเตฺวว สสฺสติสมํ.\csromanspacer{} ตํ กิสฺส เหตุ, อหํ หิ อาตปฺปมนฺวาย ปธานมนฺวาย อนุโยคมนฺวาย อปฺปมาทมนฺวาย สมฺมา\-มนสิการม\-นฺวาย ตถารูปํ เจโตสมาธิํ ผุสามิ, ยถาสมาหิเต จิตฺเต อเนกวิหิตํ ปุพฺเพนิวาสํ อนุสฺสรามิ.\csromanspacer{} เสยฺยถิทํ, ทสปิ สํวฏฺฏวิวฏฺฏานิ วีสํปิ สํวฏฺฏวิวฏฺฏานิ ติํสํปิ สํวฏฺฏวิวฏฺฏานิ จตฺตาลีสํปิ สํวฏฺฏวิวฏฺฏานิ, “อมุตฺราสิํ เอวํนาโม เอวํโคตฺโต เอวํวณฺโณ เอวมาหาโร เอวํสุขทุกฺขปฺปฏิสํเวที เอวมายุปริยนฺโต, โส ตโต จุโต อมุตฺร อุทปาทิํ, ตตฺราปาสิํ เอวํนาโม เอวํโคตฺโต เอวํวณฺโณ เอวมาหาโร เอวํสุขทุกฺขปฺปฏิสํเวที เอวมายุปริยนฺโต, โส ตโต จุโต อิธูปปนฺโน”ติ, อิติ สาการํ ส-อุทฺเทสํ อเนกวิหิตํ ปุพฺเพนิวาสํ อนุสฺสรามิ.}

\prose{อิมินามหํ เอตํ ชานามิ “ยถา สสฺสโต อตฺตา จ โลโก จ วญฺโฌ กูฏฏฺโฐ เอสิกฏฺฐายิฏฺฐิโต, เต จ สตฺตา สนฺธาวนฺติ สํสรนฺติ จวนฺติ อุปปชฺชนฺติ, อตฺถิเตฺวว สสฺสติสมนฺ”ติ, อิทํ ภิกฺขเว ตติยํ ฐานํ, ยํ อาคมฺม ยํ อารพฺภ เอเก สมณพฺราหฺมณา สสฺสตวาทา สสฺสตํ อตฺตานญฺจ โลกญฺจ ปญฺญเปนฺติ. (3)}

\proseitem{34}{จตุตฺเถ จ โภนฺโต สมณพฺราหฺมณา กิมาคมฺม กิมารพฺภ สสฺสตวาทา สสฺสตํ อตฺตานญฺจ โลกญฺจ ปญฺญเปนฺติ? อิธ ภิกฺขเว เอกจฺโจ สมโณ วา พฺราหฺมโณ วา ตกฺกี โหติ วีมํสี, โส ตกฺกปริยาหตํ วีมํสานุจริตํ สยํ ปฏิภานํ เอวมาห “สสฺสโต อตฺตา จ โลโก จ วญฺโฌ กูฏฏฺโฐ เอสิกฏฺฐายิฏฺฐิโต, เต จ สตฺตา สนฺธาวนฺติ สํสรนฺติ จวนฺติ อุปปชฺชนฺติ, อตฺถิเตฺวว สสฺสติสมนฺ”ติ, อิทํ ภิกฺขเว จตุตฺถํ ฐานํ, ยํ อาคมฺม ยํ อารพฺภ เอเก สมณพฺราหฺมณา สสฺสตวาทา สสฺสตํ อตฺตานญฺจ โลกญฺจ ปญฺญเปนฺติ. (4)}

\proseitem{35}{อิเมหิ โข เต ภิกฺขเว สมณพฺราหฺมณา สสฺสตวาทา สสฺสตํ อตฺตานญฺจ โลกญฺจ ปญฺญเปนฺติ จตูหิ วตฺถูหิ, เย หิ เกจิ ภิกฺขเว สมณา วา พฺราหฺมณา วา สสฺสตวาทา สสฺสตํ อตฺตานญฺจ โลกญฺจ ปญฺญเปนฺติ, สพฺเพ เต อิเมเหว จตูหิ วตฺถูหิ, เอเตสํ วา อญฺญตเรน, นตฺถิ อิโต พหิทฺธา.}

\proseitem{36}{\csromanfolio{16}ตยิทํ ภิกฺขเว ตถาคโต ปชานาติ “อิเม ทิฏฺฐิฏฺฐานา เอวํคหิตา เอวํปรามฏฺฐา เอวํคติกา ภวนฺติ เอวํ-อภิสมฺปรายา”ติ, ตญฺจ ตถาคโต ปชานาติ, ตโต จ อุตฺตริตรํ ปชานาติ, ตญฺจ ปชานนํ\footnote{ปชานํ (?) ที 3. 23 ปิฏฺเฐ ปาฬิอฏฺฐกถา ปสฺสิตพฺพา.} น ปรามสติ, อปรามสโต จสฺส ปจฺจตฺตญฺเญว นิพฺพุติ วิทิตา.\csromanspacer{} เวทนานํ สมุทยญฺจ อตฺถงฺคมญฺจ อสฺสาทญฺจ อาทีนวญฺจ นิสฺสรณญฺจ ยถาภูตํ วิทิตฺวา อนุปาทา วิมุตฺโต ภิกฺขเว ตถาคโต.}

\proseitem{37}{อิเม โข เต ภิกฺขเว ธมฺมา คมฺภีรา ทุทฺทสา ทุรนุโพธา สนฺตา ปณีตา อตกฺกาวจรา นิปุณา ปณฺฑิตเวทนียา, เย ตถาคโต สยํ อภิญฺญา สจฺฉิกตฺวา ปเวเทติ, เยหิ ตถาคตสฺส ยถาภุจฺจํ วณฺณํ สมฺมา วทมานา วเทยฺยุํ.}

\csromanheader{2}{ปฐมภาณวาโร.}
\csromansectionrule
\csromanmatikaanchor{mtk.8}
\csromantocmarkref{section}{เอกจฺจสสฺสตวาท}{mtk.8}
\bookmark[dest={mtk.8},level=1]{เอกจฺจสสฺสตวาท}
\csromanheader{1}{\textbf{เอกจฺจสสฺสตวาท}}
\proseitem{38}{สนฺติ ภิกฺขเว เอเก สมณพฺราหฺมณา เอกจฺจสสฺสติกา เอกจฺจ- อสสฺสติกา เอกจฺจํ สสฺสตํ เอกจฺจํ อสสฺสตํ อตฺตานญฺจ โลกญฺจ ปญฺญเปนฺติ จตูหิ วตฺถูหิ.\csromanspacer{} เต จ โภนฺโต สมณพฺราหฺมณา กิมาคมฺม กิมารพฺภ เอกจฺจสสฺสติกา เอกจฺจ-อสสฺสติกา เอกจฺจํ สสฺสตํ เอกจฺจํ อสสฺสตํ อตฺตานญฺจ โลกญฺจ ปญฺญเปนฺติ จตูหิ วตฺถูหิ?}

\proseitem{39}{โหติ โข โส ภิกฺขเว สมโย, ยํ กทาจิ กรหจิ ทีฆสฺส อทฺธุโน อจฺจเยน อยํ โลโก สํวฏฺฏติ, สํวฏฺฏมาเน โลเก เยภุยฺเยน สตฺตา อาภสฺสรสํวตฺตนิกา โหนฺติ, เต ตตฺถ โหนฺติ มโนมยา ปีติภกฺขา สยํปภา อนฺตลิกฺขจรา สุภฏฺฐายิโน, จิรํ ทีฆมทฺธานํ ติฏฺฐนฺติ.}

\proseitem{40}{โหติ โข โส ภิกฺขเว สมโย, ยํ กทาจิ กรหจิ ทีฆสฺส อทฺธุโน อจฺจเยน อยํ โลโก วิวฏฺฏติ, วิวฏฺฏมาเน โลเก สุญฺญํ พฺรหฺมวิมานํ ปาตุภวติ.\csromanspacer{} อถ โข อญฺญตโร สตฺโต อายุกฺขยา วา ปุญฺญกฺขยา วา อาภสฺสรกายา จวิตฺวา สุญฺญํ พฺรหฺมวิมานํ อุปปชฺชติ, \csromanfolio{17}โส ตตฺถ โหติ มโนมโย ปีติภกฺโข สยํปโภ อนฺตลิกฺขจโร สุภฏฺฐายี, จิรํ ทีฆมทฺธานํ ติฏฺฐติ.}

\proseitem{41}{ตสฺส ตตฺถ เอกกสฺส ทีฆรตฺตํ นิวุสิตตฺตา อนภิรติ ปริตสฺสนา อุปฺปชฺชติ “อโห วต อญฺเญปิ สตฺตา อิตฺถตฺตํ อาคจฺเฉยฺยุนฺ”ติ.\csromanspacer{} อถ อญฺเญปิ สตฺตา อายุกฺขยา วา ปุญฺญกฺขยา วา อาภสฺสรกายา จวิตฺวา พฺรหฺมวิมานํ อุปปชฺชนฺติ ตสฺส สตฺตสฺส สหพฺยตํ, เตปิ ตตฺถ โหนฺติ มโนมยา ปีติภกฺขา สยํปภา อนฺตลิกฺขจรา สุภฏฺฐายิโน, จิรํ ทีฆมทฺธานํ ติฏฺฐนฺติ.}

\proseitem{42}{ตตฺร ภิกฺขเว โย โส สตฺโต ปฐมํ อุปปนฺโน, ตสฺส เอวํ โหติ “อหมสฺมิ พฺรหฺมา มหาพฺรหฺมา อภิภู อนภิภูโต อญฺญทตฺถุทโส วสวตฺตี อิสฺสโร กตฺตา นิมฺมาตา เสฏฺโฐ สชิตา\footnote{สชฺชิตา (สฺยา, กํ)} วสี ปิตา ภูตภพฺยานํ, มยา อิเม สตฺตา นิมฺมิตา, ตํ กิสฺส เหตุ, มมํ หิ ปุพฺเพ เอตทโหสิ ‘อโห วต อญฺเญปิ สตฺตา อิตฺถตฺตํ อาคจฺเฉยฺยุนฺ’ติ, อิติ มม จ มโนปณิธิ, อิเม จ สตฺตา อิตฺถตฺตํ อาคตา”ติ.}

\prose{เยปิ เต สตฺตา ปจฺฉา อุปปนฺนา, เตสมฺปิ เอวํ โหติ “อยํ โข ภวํ พฺรหฺมา มหาพฺรหฺมา อภิภู อนภิภูโต อญฺญทตฺถุทโส วสวตฺตี อิสฺสโร กตฺตา นิมฺมาตา เสฏฺโฐ สชิตา วสี ปิตา ภูตภพฺยานํ, อิมินา มยํ โภตา พฺรหฺมุนา นิมฺมิตา, ตํ กิสฺส เหตุ, อิมํ หิ มยํ อทฺทสาม อิธ ปฐมํ อุปปนฺนํ, มยํ ปนมฺห ปจฺฉา อุปปนฺนา”ติ.}

\proseitem{43}{ตตฺร ภิกฺขเว โย โส สตฺโต ปฐมํ อุปปนฺโน, โส ทีฆายุกตโร จ โหติ วณฺณวนฺตตโร จ มเหสกฺขตโร จ.\csromanspacer{} เย ปน เต สตฺตา ปจฺฉา อุปปนฺนา, เต อปฺปายุกตรา จ โหนฺติ ทุพฺพณฺณตรา จ อปฺเปสกฺขตรา จ.}

\proseitem{44}{ฐานํ โข ปเนตํ ภิกฺขเว วิชฺชติ, ยํ อญฺญตโร สตฺโต ตมฺหา กายา จวิตฺวา อิตฺถตฺตํ อาคจฺฉติ, อิตฺถตฺตํ อาคโต สมาโน อคารสฺมา อนคาริยํ ปพฺพชติ, อคารสฺมา อนคาริยํ ปพฺพชิโต \csromanfolio{18}สมาโน อาตปฺปมนฺวาย ปธานมนฺวาย อนุโยคมนฺวาย อปฺปมาทมนฺวาย สมฺมา\-มนสิการม\-นฺวาย ตถารูปํ เจโตสมาธิํ ผุสติ, ยถาสมาหิเต จิตฺเต ตํ ปุพฺเพนิวาสํ อนุสฺสรติ, ตโต ปรํ นานุสฺสรติ.}

\prose{โส เอวมาห “โย โข โส ภวํ พฺรหฺมา มหาพฺรหฺมา อภิภู อนภิภูโต อญฺญทตฺถุทโส วสวตฺตี อิสฺสโร กตฺตา นิมฺมาตา เสฏฺโฐ สชิตา วสี ปิตา ภูตภพฺยานํ, เยน มยํ โภตา พฺรหฺมุนา นิมฺมิตา, โส นิจฺโจ ธุโว สสฺสโต อวิปริณามธมฺโม สสฺสติสมํ ตเถว ฐสฺสติ.\csromanspacer{} เย ปน มยํ อหุมฺหา เตน โภตา พฺรหฺมุนา นิมฺมิตา, เต มยํ อนิจฺจา อทฺธุวา อปฺปายุกา จวนธมฺมา อิตฺถตฺตํ อาคตา”ติ.\csromanspacer{} อิทํ ภิกฺขเว ปฐมํ ฐานํ, ยํ อาคมฺม ยํ อารพฺภ เอเก สมณพฺราหฺมณา เอกจฺจสสฺสติกา เอกจฺจ-อสสฺสติกา เอกจฺจํ สสฺสตํ เอกจฺจํ อสสฺสตํ อตฺตานญฺจ โลกญฺจ ปญฺญเปนฺติ. (1=5)}

\proseitem{45}{ทุติเย จ โภนฺโต สมณพฺราหฺมณา กิมาคมฺม กิมารพฺภ เอกจฺจสสฺสติกา เอกจฺจ-อสสฺสติกา เอกจฺจํ สสฺสตํ เอกจฺจํ อสสฺสตํ อตฺตานญฺจ โลกญฺจ ปญฺญเปนฺติ? สนฺติ ภิกฺขเว ขิฑฺฑาปโทสิกา นาม เทวา, เต อติเวลํ หสฺส\-ขิฑฺฑา\-รติธมฺม\-สมาปนฺนา\footnote{หสขิฑฺฑารติธมฺมสมาปนฺนา (ก)} วิหรนฺติ, เตสํ อติเวลํ หสฺส\-ขิฑฺฑา\-รติธมฺม\-สมาปนฺนานํ วิหรตํ สติ สมฺมุสฺสติ\footnote{ปมุสฺสติ (สี, สฺยา)}, สติยา สมฺโมสา เต เทวา ตมฺหา กายา จวนฺติ.}

\proseitem{46}{ฐานํ โข ปเนตํ ภิกฺขเว วิชฺชติ, ยํ อญฺญตโร สตฺโต ตมฺหา กายา จวิตฺวา อิตฺถตฺตํ อาคจฺฉติ, อิตฺถตฺตํ อาคโต สมาโน อคารสฺมา อนคาริยํ ปพฺพชติ, อคารสฺมา อนคาริยํ ปพฺพชิโต สมาโน อาตปฺปมนฺวาย ปธานมนฺวาย อนุโยคมนฺวาย อปฺปมาทมนฺวาย สมฺมา\-มนสิการม\-นฺวาย ตถารูปํ เจโตสมาธิํ ผุสติ, ยถาสมาหิเต จิตฺเต ตํ ปุพฺเพนิวาสํ อนุสฺสรติ, ตโต ปรํ นานุสฺสรติ.}

\prose{โส เอวมาห “เย โข เต โภนฺโต เทวา น ขิฑฺฑาปโทสิกา, เต น อติเวลํ หสฺส\-ขิฑฺฑา\-รติธมฺม\-สมาปนฺนา วิหรนฺติ, เตสํ น อติเวลํ หสฺส\-ขิฑฺฑา\-รติธมฺม\-สมาปนฺนานํ วิหรตํ สติ น สมฺมุสฺสติ, \csromanfolio{19}สติยา อสมฺโมสา เต เทวา ตมฺหา กายา น จวนฺติ, นิจฺจา ธุวา สสฺสตา อวิปริณามธมฺมา สสฺสติสมํ ตเถว ฐสฺสนฺติ.\csromanspacer{} เย ปน มยํ อหุมฺหา ขิฑฺฑาปโทสิกา, เต มยํ อติเวลํ หสฺส\-ขิฑฺฑา\-รติธมฺม\-สมาปนฺนา วิหริมฺหา, เตสํ โน อติเวลํ หสฺส\-ขิฑฺฑา\-รติธมฺม\-สมาปนฺนานํ วิหรตํ สติ สมฺมุสฺสติ, สติยา สมฺโมสา เอวํ มยํ ตมฺหา กายา จุตา อนิจฺจา อทฺธุวา อปฺปายุกา จวนธมฺมา อิตฺถตฺตํ อาคตา”ติ.\csromanspacer{} อิทํ ภิกฺขเว ทุติยํ ฐานํ, ยํ อาคมฺม ยํ อารพฺภ เอเก สมณพฺราหฺมณา เอกจฺจสสฺสติกา เอกจฺจ- อสสฺสติกา เอกจฺจํ สสฺสตํ เอกจฺจํ อสสฺสตํ อตฺตานญฺจ โลกญฺจ ปญฺญเปนฺติ. (2=6)}

\proseitem{47}{ตติเย จ โภนฺโต สมณพฺราหฺมณา กิมาคมฺม กิมารพฺภ เอกจฺจสสฺสติกา เอกจฺจ-อสสฺสติกา เอกจฺจํ สสฺสตํ เอกจฺจํ อสสฺสตํ อตฺตานญฺจ โลกญฺจ ปญฺญเปนฺติ? สนฺติ ภิกฺขเว มโนปโทสิกา นาม เทวา, เต อติเวลํ อญฺญมญฺญํ อุปนิชฺฌายนฺติ, เต อติเวลํ อญฺญมญฺญํ อุปนิชฺฌายนฺตา อญฺญมญฺญมฺหิ จิตฺตานิ ปทูเสนฺติ, เต อญฺญมญฺญํ ปทุฏฺฐจิตฺตา กิลนฺตกายา กิลนฺตจิตฺตา เต เทวา ตมฺหา กายา จวนฺติ.}

\proseitem{48}{ฐานํ โข ปเนตํ ภิกฺขเว วิชฺชติ, ยํ อญฺญตโร สตฺโต ตมฺหา กายา จวิตฺวา อิตฺถตฺตํ อาคจฺฉติ, อิตฺถตฺตํ อาคโต สมาโน อคารสฺมา อนคาริยํ ปพฺพชติ, อคารสฺมา อนคาริยํ ปพฺพชิโต สมาโน อาตปฺปมนฺวาย ปธานมนฺวาย อนุโยคมนฺวาย อปฺปมาทมนฺวาย สมฺมา\-มนสิการม\-นฺวาย ตถารูปํ เจโตสมาธิํ ผุสติ, ยถาสมาหิเต จิตฺเต ตํ ปุพฺเพนิวาสํ อนุสฺสรติ, ตโต ปรํ นานุสฺสรติ.}

\prose{โส เอวมาห “เย โข เต โภนฺโต เทวา น มโนปโทสิกา, เต นาติเวลํ อญฺญมญฺญํ อุปนิชฺฌายนฺติ, เต นาติเวลํ อญฺญมญฺญํ อุปนิชฺฌายนฺตา อญฺญมญฺญมฺหิ จิตฺตานิ นปฺปทูเสนฺติ, เต อญฺญมญฺญํ อปฺปทุฏฺฐจิตฺตา อกิลนฺตกายา อกิลนฺตจิตฺตา เต เทวา ตมฺหา กายา น จวนฺติ, นิจฺจา ธุวา สสฺสตา อวิปริณามธมฺมา สสฺสติสมํ ตเถว ฐสฺสนฺติ.\csromanspacer{} เย ปน มยํ อหุมฺหา มโนปโทสิกา, เต มยํ อติเวลํ อญฺญมญฺญํ อุปนิชฺฌายิมฺหา, เต มยํ อติเวลํ อญฺญมญฺญํ อุปนิชฺฌายนฺตา อญฺญมญฺญมฺหิ จิตฺตานิ ปทูสิมฺหา, เต มยํ อญฺญมญฺญํ ปทุฏฺฐจิตฺตา กิลนฺตกายา กิลนฺตจิตฺตา, เอวํ มยํ ตมฺหา กายา จุตา อนิจฺจา อทฺธุวา อปฺปายุกา จวนธมฺมา อิตฺถตฺตํ \csromanfolio{20}อาคตา”ติ.\csromanspacer{} อิทํ ภิกฺขเว ตติยํ ฐานํ, ยํ อาคมฺม ยํ อารพฺภ เอเก สมณพฺราหฺมณา เอกจฺจสสฺสติกา เอกจฺจ-อสสฺสติกา เอกจฺจํ สสฺสตํ เอกจฺจํ อสสฺสตํ อตฺตานญฺจ โลกญฺจ ปญฺญเปนฺติ. (3=7)}

\proseitem{49}{จตุตฺเถ จ โภนฺโต สมณพฺราหฺมณา กิมาคมฺม กิมารพฺภ เอกจฺจสสฺสติกา เอกจฺจ-อสสฺสติกา เอกจฺจํ สสฺสตํ เอกจฺจํ อสสฺสตํ อตฺตานญฺจ โลกญฺจ ปญฺญเปนฺติ? อิธ ภิกฺขเว เอกจฺโจ สมโณ วา พฺราหฺมโณ วา ตกฺกี โหติ วีมํสี, โส ตกฺกปริยาหตํ วีมํสานุจริตํ สยํปฏิภานํ เอวมาห “ยํ โข อิทํ วุจฺจติ จกฺขุํ อิติปิ โสตํ อิติปิ ฆานํ อิติปิ ชิวฺหา อิติปิ กาโย อิติปิ, อยํ อตฺตา อนิจฺโจ อทฺธุโว อสสฺสโต วิปริณามธมฺโม, ยญฺจ โข อิทํ วุจฺจติ จิตฺตนฺติ วา มโนติ วา วิญฺญาณนฺติ วา, อยํ อตฺตา นิจฺโจ ธุโว สสฺสโต อวิปริณามธมฺโม สสฺสติสมํ ตเถว ฐสฺสตี”ติ.\csromanspacer{} อิทํ ภิกฺขเว จตุตฺถํ ฐานํ, ยํ อาคมฺม ยํ อารพฺภ เอเก สมณพฺราหฺมณา เอกจฺจสสฺสติกา เอกจฺจ- อสสฺสติกา เอกจฺจํ สสฺสตํ เอกจฺจํ อสสฺสตํ อตฺตานญฺจ โลกญฺจ ปญฺญเปนฺติ. (4=8)}

\proseitem{50}{อิเมหิ โข เต ภิกฺขเว สมณพฺราหฺมณา เอกจฺจสสฺสติกา เอกจฺจ- อสสฺสติกา เอกจฺจํ สสฺสตํ เอกจฺจํ อสสฺสตํ อตฺตานญฺจ โลกญฺจ ปญฺญเปนฺติ จตูหิ วตฺถูหิ.\csromanspacer{} เย หิ เกจิ ภิกฺขเว สมณา วา พฺราหฺมณา วา เอกจฺจสสฺสติกา เอกจฺจ-อสสฺสติกา เอกจฺจํ สสฺสตํ เอกจฺจํ อสสฺสตํ อตฺตานญฺจ โลกญฺจ ปญฺญเปนฺติ, สพฺเพ เต อิเมเหว จตูหิ วตฺถูหิ, เอเตสํ วา อญฺญตเรน, นตฺถิ อิโต พหิทฺธา.}

\proseitem{51}{ตยิทํ ภิกฺขเว ตถาคโต ปชานาติ “อิเม ทิฏฺฐิฏฺฐานา เอวํคหิตา เอวํปรามฏฺฐา เอวํคติกา ภวนฺติ เอวํ-อภิสมฺปรายา”ติ, ตญฺจ ตถาคโต ปชานาติ, ตโต จ อุตฺตริตรํ ปชานาติ, ตญฺจ ปชานนํ น ปรามสติ, อปรามสโต จสฺส ปจฺจตฺตญฺเญว นิพฺพุติ วิทิตา, เวทนานํ สมุทยญฺจ อตฺถงฺคมญฺจ อสฺสาทญฺจ อาทีนวญฺจ นิสฺสรณญฺจ ยถาภูตํ วิทิตฺวา อนุปาทาวิมุตฺโต ภิกฺขเว ตถาคโต.}

\csromanmatikaanchor{mtk.9}
\csromantocmarkref{section}{อนฺตานนฺตวท}{mtk.9}
\bookmark[dest={mtk.9},level=1]{อนฺตานนฺตวท}
\proseitem{52}{อิเม โข เต ภิกฺขเว ธมฺมา คมฺภีรา ทุทฺทสา ทุรนุโพธา สนฺตา ปณีตา อตกฺกาวจรา นิปุณา ปณฺฑิตเวทนียา, เย ตถาคโต \csromanfolio{21}สยํ อภิญฺญา สจฺฉิกตฺวา ปเวเทติ, เยหิ ตถาคตสฺส ยถาภุจฺจํ วณฺณํ สมฺมา วทมานา วเทยฺยุํ.}

\csromanheader{2}{\textbf{อนฺตานนฺตวาท}}
\proseitem{53}{สนฺติ ภิกฺขเว เอเก สมณพฺราหฺมณา อนฺตานนฺติกา อนฺตานนฺตํ โลกสฺส ปญฺญเปนฺติ จตูหิ วตฺถูหิ.\csromanspacer{} เต จ โภนฺโต สมณพฺราหฺมณา กิมาคมฺม กิมารพฺภ อนฺตานนฺติกา อนฺตานนฺตํ โลกสฺส ปญฺญเปนฺติ จตูหิ วตฺถูหิ?}

\proseitem{54}{อิธ ภิกฺขเว เอกจฺโจ สมโณ วา พฺราหฺมโณ วา อาตปฺปมนฺวาย ปธานมนฺวาย อนุโยคมนฺวาย อปฺปมาทมนฺวาย สมฺมา\-มนสิการม\-นฺวาย ตถารูปํ เจโตสมาธิํ ผุสติ, ยถาสมาหิเต จิตฺเต อนฺตสญฺญี โลกสฺมิํ วิหรติ.}

\prose{โส เอวมาห “อนฺตวา อยํ โลโก ปริวฏุโม, ตํ กิสฺส เหตุ, อหํ หิ อาตปฺปมนฺวาย ปธานมนฺวาย อนุโยคมนฺวาย อปฺปมาทมนฺวาย สมฺมา\-มนสิการม\-นฺวาย ตถารูปํ เจโตสมาธิํ ผุสามิ, ยถาสมาหิเต จิตฺเต อนฺตสญฺญี โลกสฺมิํ วิหรามิ, อิมินามหํ เอตํ ชานามิ ยถา อนฺตวา อยํ โลโก ปริวฏุโม”ติ.\csromanspacer{} อิทํ ภิกฺขเว ปฐมํ ฐานํ, ยํ อาคมฺม ยํ อารพฺภ เอเก สมณพฺราหฺมณา อนฺตานนฺติกา อนฺตานนฺตํ โลกสฺส ปญฺญเปนฺติ. (1=9)}

\proseitem{55}{ทุติเย จ โภนฺโต สมณพฺราหฺมณา กิมาคมฺม กิมารพฺภ อนฺตานนฺติกา อนฺตานนฺตํ โลกสฺส ปญฺญเปนฺติ? อิธ ภิกฺขเว เอกจฺโจ สมโณ วา พฺราหฺมโณ วา อาตปฺปมนฺวาย ปธานมนฺวาย อนุโยคมนฺวาย อปฺปมาทมนฺวาย สมฺมา\-มนสิการม\-นฺวาย ตถารูปํ เจโตสมาธิํ ผุสติ, ยถาสมาหิเต จิตฺเต อนนฺตสญฺญี โลกสฺมิํ วิหรติ.}

\prose{โส เอวมาห “อนนฺโต อยํ โลโก อปริยนฺโต.\csromanspacer{} เย เต สมณพฺราหฺมณา เอวมาหํสุ ‘อนฺตวา อยํ โลโก ปริวฏุโม’ติ, เตสํ มุสา.\csromanspacer{} อนนฺโต อยํ โลโก อปริยนฺโต, ตํ กิสฺส เหตุ, อหํ หิ อาตปฺปมนฺวาย ปธานมนฺวาย อนุโยคมนฺวาย อปฺปมาทมนฺวาย สมฺมา\-มนสิการม\-นฺวาย ตถารูปํ เจโตสมาธิํ \csromanfolio{22}ผุสามิ, ยถาสมาหิเต จิตฺเต อนนฺตสญฺญี โลกสฺมิํ วิหรามิ, อิมินามหํ เอตํ ชานามิ ยถา อนนฺโต อยํ โลโก อปริยนฺโต”ติ.\csromanspacer{} อิทํ ภิกฺขเว ทุติยํ ฐานํ, ยํ อาคมฺม ยํ อารพฺภ เอเก สมณพฺราหฺมณา อนฺตานนฺติกา อนฺตานนฺตํ โลกสฺส ปญฺญเปนฺติ. (2=10)}

\proseitem{56}{ตติเย จ โภนฺโต สมณพฺราหฺมณา กิมาคมฺม กิมารพฺภ อนฺตานนฺติกา อนฺตานนฺตํ โลกสฺส ปญฺญเปนฺติ? อิธ ภิกฺขเว เอกจฺโจ สมโณ วา พฺราหฺมโณ วา อาตปฺปมนฺวาย ปธานมนฺวาย อนุโยคมนฺวาย อปฺปมาทมนฺวาย สมฺมา\-มนสิการม\-นฺวาย ตถารูปํ เจโตสมาธิํ ผุสติ, ยถาสมาหิเต จิตฺเต อุทฺธมโธ อนฺตสญฺญี โลกสฺมิํ วิหรติ ติริยํ อนนฺตสญฺญี.}

\prose{โส เอวมาห “อนฺตวา จ อยํ โลโก อนนฺโต จ.\csromanspacer{} เย เต สมณพฺราหฺมณา เอวมาหํสุ ‘อนฺตวา อยํ โลโก ปริวฏุโม’ติ, เตสํ มุสา.\csromanspacer{} เยปิ เต สมณพฺราหฺมณา เอวมาหํสุ ‘อนนฺโต อยํ โลโก อปริยนฺโต’ติ, เตสมฺปิ มุสา.\csromanspacer{} อนฺตวา จ อยํ โลโก อนนฺโต จ, ตํ กิสฺส เหตุ, อหํ หิ อาตปฺปมนฺวาย ปธานมนฺวาย อนุโยคมนฺวาย อปฺปมาทมนฺวาย สมฺมา\-มนสิการม\-นฺวาย ตถารูปํ เจโตสมาธิํ ผุสามิ, ยถาสมาหิเต จิตฺเต อุทฺธมโธ อนฺตสญฺญี โลกสฺมิํ วิหรามิ ติริยํ อนนฺตสญฺญี, อิมินามหํ เอตํ ชานามิ ยถา อนฺตวา จ อยํ โลโก อนนฺโต จา”ติ.\csromanspacer{} อิทํ ภิกฺขเว ตติยํ ฐานํ, ยํ อาคมฺม ยํ อารพฺภ เอเก สมณพฺราหฺมณา อนฺตานนฺติกา อนฺตานนฺตํ โลกสฺส ปญฺญเปนฺติ. (3=11)}

\proseitem{57}{จตุตฺเถ จ โภนฺโต สมณพฺราหฺมณา กิมาคมฺม กิมารพฺภ อนฺตานนฺติกา อนฺตานนฺตํ โลกสฺส ปญฺญเปนฺติ? อิธ ภิกฺขเว เอกจฺโจ สมโณ วา พฺราหฺมโณ วา ตกฺกี โหติ วีมํสี, โส ตกฺกปริยาหตํ วีมํสานุจริตํ สยํปฏิภานํ เอวมาห “เนวายํ โลโก อนฺตวา, น ปนานนฺโต.\csromanspacer{} เย เต สมณพฺราหฺมณา เอวมาหํสุ ‘อนฺตวา อยํ โลโก ปริวฏุโม’ติ, เตสํ มุสา.\csromanspacer{} เยปิ เต สมณพฺราหฺมณา เอวมาหํสุ ‘อนนฺโต อยํ โลโก อปริยนฺโต’ติ, เตสมฺปิ มุสา.\csromanspacer{} เยปิ เต สมณพฺราหฺมณา เอวมาหํสุ ‘อนฺตวา จ อยํ โลโก อนนฺโต จา’ติ, เตสมฺปิ มุสา.\csromanspacer{} เนวายํ โลโก \csromanfolio{23}อนฺตวา, น ปนานนฺโต”ติ.\csromanspacer{} อิทํ ภิกฺขเว จตุตฺถํ ฐานํ, ยํ อาคมฺม ยํ อารพฺภ เอเก สมณพฺราหฺมณา อนฺตานนฺติกา อนฺตานนฺตํ โลกสฺส ปญฺญเปนฺติ. (4=12)}

\proseitem{58}{อิเมหิ โข เต ภิกฺขเว สมณพฺราหฺมณา อนฺตานนฺติกา อนฺตานนฺตํ โลกสฺส ปญฺญเปนฺติ จตูหิ วตฺถูหิ.\csromanspacer{} เย หิ เกจิ ภิกฺขเว สมณา วา พฺราหฺมณา วา อนฺตานนฺติกา อนฺตานนฺตํ โลกสฺส ปญฺญเปนฺติ, สพฺเพ เต อิเมเหว จตูหิ วตฺถูหิ, เอเตสํ วา อญฺญตเรน, นตฺถิ อิโต พหิทฺธา.}

\proseitem{59}{ตยิทํ ภิกฺขเว ตถาคโต ปชานาติ “อิเม ทิฏฺฐิฏฺฐานา เอวํคหิตา เอวํปรามฏฺฐา เอวํคติกา ภวนฺติ เอวํอภิสมฺปรายา”ติ, ตญฺจ ตถาคโต ปชานาติ, ตโต จ อุตฺตริตรํ ปชานาติ, ตญฺจ ปชานนํ น ปรามสติ, อปรามสโต จสฺส ปจฺจตฺตญฺเญว นิพฺพุติ วิทิตา, เวทนานํ สมุทยญฺจ อตฺถงฺคมญฺจ อสฺสาทญฺจ อาทีนวญฺจ นิสฺสรณญฺจ ยถาภูตํ วิทิตฺวา อนุปาทาวิมุตฺโต ภิกฺขเว ตถาคโต.}

\proseitem{60}{อิเม โข เต ภิกฺขเว ธมฺมา คมฺภีรา ทุทฺทสา ทุรนุโพธา สนฺตา ปณีตา อตกฺกาวจรา นิปุณา ปณฺฑิตเวทนียา, เย ตถาคโต สยํ อภิญฺญา สจฺฉิกตฺวา ปเวเทติ, เยหิ ตถาคตสฺส ยถาภุจฺจํ วณฺณํ สมฺมา วทมานา วเทยฺยุํ.}

\csromanmatikaanchor{mtk.10}
\csromantocmarkref{section}{อมราวิกฺเขปวาท}{mtk.10}
\bookmark[dest={mtk.10},level=1]{อมราวิกฺเขปวาท}
\csromanheader{1}{\textbf{อมราวิกฺเขปวาท}}
\proseitem{61}{สนฺติ ภิกฺขเว เอเก สมณพฺราหฺมณา อมราวิกฺเขปิกา ตตฺถ ตตฺถ ปญฺหํ ปุฏฺฐา สมานา วาจาวิกฺเขปํ อาปชฺชนฺติ อมราวิกฺเขปํ จตูหิ วตฺถูหิ.\csromanspacer{} เต จ โภนฺโต สมณพฺราหฺมณา กิมาคมฺม กิมารพฺภ อมราวิกฺเขปิกา ตตฺถ ตตฺถ ปญฺหํ ปุฏฺฐา สมานา วาจาวิกฺเขปํ อาปชฺชนฺติ อมราวิกฺเขปํ จตูหิ วตฺถูหิ?}

\proseitem{62}{อิธ ภิกฺขเว เอกจฺโจ สมโณ วา พฺราหฺมโณ วา “อิทํกุสลนฺ”ติ ยถาภูตํ นปฺปชานาติ, “อิทํอกุสลนฺ”ติ ยถาภูตํ นปฺปชานาติ, ตสฺส เอวํ โหติ “อหํ โข ‘อิทํ กุสลนฺ’ติ ยถาภูตํ นปฺปชานามิ, ‘อิทํ อกุสลนฺ’ติ ยถาภูตํ นปฺปชานามิ, อหญฺเจ โข ปน \csromanfolio{24}‘อิทํ กุสลนฺ’ติ ยถาภูตํ อปฺปชานนฺโต ‘อิทํ อกุสลนฺ’ติ ยถาภูตํ อปฺปชานนฺโต ‘อิทํ กุสลนฺ’ติ วา พฺยากเรยฺยํ, ‘อิทํ อกุสลนฺ’ติ วา พฺยากเรยฺยํ, ตํ มมสฺส มุสา.\csromanspacer{} ยํ มมสฺส มุสา, โส มมสฺส วิฆาโต.\csromanspacer{} โย มมสฺส วิฆาโต, โส มมสฺส อนฺตราโย”ติ.\csromanspacer{} อิติ โส มุสาวาทภยา มุสาวาทปริเชคุจฺฉา เนวิทํ กุสลนฺติ พฺยากโรติ, น ปนิทํ อกุสลนฺติ พฺยากโรติ, ตตฺถ ตตฺถ ปญฺหํ ปุฏฺโฐ สมาโน วาจาวิกฺเขปํ อาปชฺชติ อมราวิกฺเขปํ “เอวนฺติปิ เม โน, ตถาติปิ เม โน, อญฺญถาติปิ เม โน, โนติปิ เม โน, โน โนติปิ เม โน”ติ.\csromanspacer{} อิทํ ภิกฺขเว ปฐมํ ฐานํ, ยํ อาคมฺม ยํ อารพฺภ เอเก สมณพฺราหฺมณา อมราวิกฺเขปิกา ตตฺถ ตตฺถ ปญฺหํ ปุฏฺฐา สมานา วาจาวิกฺเขปํ อาปชฺชนฺติ อมราวิกฺเขปํ. (1=13)}

\proseitem{63}{ทุติเย จ โภนฺโต สมณพฺราหฺมณา กิมาคมฺม กิมารพฺภ อมราวิกฺเขปิกา ตตฺถ ตตฺถ ปญฺหํ ปุฏฺฐา สมานา วาจาวิกฺเขปํ อาปชฺชนฺติ อมราวิกฺเขปํ? อิธ ภิกฺขเว เอกจฺโจ สมโณ วา พฺราหฺมโณ วา “อิทํ กุสลนฺ”ติ ยถาภูตํ นปฺปชานาติ, “อิทํ อกุสลนฺ”ติ ยถาภูตํ นปฺปชานาติ.\csromanspacer{} ตสฺส เอวํ โหติ “อหํ โข ‘อิทํ กุสลนฺ’ติ ยถาภูตํ นปฺปชานามิ, ‘อิทํ อกุสลนฺ’ติ ยถาภูตํ นปฺปชานามิ, อหญฺเจ โข ปน ‘อิทํ กุสลนฺ’ติ ยถาภูตํ อปฺปชานนฺโต ‘อิทํ อกุสลนฺ’ติ ยถาภูตํ อปฺปชานนฺโต ‘อิทํ กุสลนฺ’ติ วา พฺยากเรยฺยํ, ‘อิทํ อกุสลนฺ’ติ วา พฺยากเรยฺยํ, ตตฺถ เม อสฺส ฉนฺโท วา ราโค วา โทโส วา ปฏิโฆ วา.\csromanspacer{} ยตฺถ\footnote{โย (?)} เม อสฺส ฉนฺโท วา ราโค วา โทโส วา ปฏิโฆ วา, ตํ มมสฺส อุปาทานํ.\csromanspacer{} ยํ มมสฺส อุปาทานํ, โส มมสฺส วิฆาโต.\csromanspacer{} โย มมสฺส วิฆาโต, โส มมสฺส อนฺตราโย”ติ.\csromanspacer{} อิติ โส อุปาทานภยา อุปาทานปริเชคุจฺฉา เนวิทํ กุสลนฺติ พฺยากโรติ, น ปนิทํ อกุสลนฺติ พฺยากโรติ, ตตฺถ ตตฺถ ปญฺหํ ปุฏฺโฐ สมาโน วาจาวิกฺเขปํ อาปชฺชติ อมราวิกฺเขปํ “เอวนฺติปิ เม โน, ตถาติปิ เม โน, อญฺญถาติปิ เม โน, โนติปิ เม โน, โน โนติปิ เม โน”ติ.\csromanspacer{} อิทํ ภิกฺขเว ทุติยํ ฐานํ, ยํ อาคมฺม ยํ อารพฺภ เอเก สมณพฺราหฺมณา \csromanfolio{25}อมราวิกฺเขปิกา ตตฺถ ตตฺถ ปญฺหํ ปุฏฺฐา สมานา วาจาวิกฺเขปํ อาปชฺชนฺติ อมราวิกฺเขปํ. (2=14)}

\proseitem{64}{ตติเย จ โภนฺโต สมณพฺราหฺมณา กิมาคมฺม กิมารพฺภ อมราวิกฺเขปิกา ตตฺถ ตตฺถ ปญฺหํ ปุฏฺฐา สมานา วาจาวิกฺเขปํ อาปชฺชนฺติ อมราวิกฺเขปํ? อิธ ภิกฺขเว เอกจฺโจ สมโณ วา พฺราหฺมโณ วา “อิทํกุสลนฺ”ติ ยถาภูตํ นปฺปชานาติ, “อิทํอกุสลนฺ”ติ ยถาภูตํ นปฺปชานาติ.\csromanspacer{} ตสฺส เอวํ โหติ “อหํ โข ‘อิทํ กุสลนฺ’ติ ยถาภูตํ นปฺปชานามิ, ‘อิทํ อกุสลนฺ’ติ ยถาภูตํ นปฺปชานามิ, อหญฺเจ โข ปน ‘อิทํ กุสลนฺ’ติ ยถาภูตํ อปฺปชานนฺโต ‘อิทํ อกุสลนฺ’ติ ยถาภูตํ อปฺปชานนฺโต ‘อิทํ กุสลนฺ’ติ วา พฺยากเรยฺยํ, ‘อิทํ อกุสลนฺ’ติ วา พฺยากเรยฺยํ.\csromanspacer{} สนฺติ หิ โข สมณพฺราหฺมณา ปณฺฑิตา นิปุณา กตปรปฺปวาทา วาลเวธิรูปา, เต ภินฺทนฺตา\footnote{โวภินฺทนฺตา (สี, อิ)} มญฺเญ จรนฺติ ปญฺญาคเตน ทิฏฺฐิคตานิ, เต มํ ตตฺถ สมนุยุญฺเชยฺยุํ สมนุคาเหยฺยุํ สมนุภาเสยฺยุํ.\csromanspacer{} เย มํ ตตฺถ สมนุยุญฺเชยฺยุํ สมนุคาเหยฺยุํ สมนุภาเสยฺยุํ, เตสาหํ น สมฺปาเยยฺยํ.\csromanspacer{} เยสาหํ น สมฺปาเยยฺยํ, โส มมสฺส วิฆาโต, โย มมสฺส วิฆาโต, โส มมสฺส อนฺตราโย”ติ, อิติ โส อนุโยคภยา อนุโยคปริเชคุจฺฉา เนวิทํ กุสลนฺติ พฺยากโรติ, น ปนิทํ อกุสลนฺติ พฺยากโรติ, ตตฺถ ตตฺถ ปญฺหํ ปุฏฺโฐ สมาโน วาจาวิกฺเขปํ อาปชฺชติ อมราวิกฺเขปํ “เอวนฺติปิ เม โน, ตถาติปิ เม โน, อญฺญถาติปิ เม โน, โนติปิ เม โน, โน โนติปิ เม โน”ติ.\csromanspacer{} อิทํ ภิกฺขเว ตติยํ ฐานํ, ยํ อาคมฺม ยํ อารพฺภ เอเก สมณพฺราหฺมณา อมราวิกฺเขปิกา ตตฺถ ตตฺถ ปญฺหํ ปุฏฺฐา สมานา วาจาวิกฺเขปํ อาปชฺชนฺติ อมราวิกฺเขปํ. (3=15)}

\proseitem{65}{จตุตฺเถ จ โภนฺโต สมณพฺราหฺมณา กิมาคมฺม กิมารพฺภ อมราวิกฺเขปิกา ตตฺถ ตตฺถ ปญฺหํ ปุฏฺฐา สมานา วาจาวิกฺเขปํ อาปชฺชนฺติ อมราวิกฺเขปํ? อิธ ภิกฺขเว เอกจฺโจ สมโณ วา พฺราหฺมโณ วา มนฺโท โหติ โมมูโห, โส มนฺทตฺตา โมมูหตฺตา ตตฺถ ตตฺถ ปญฺหํ ปุฏฺโฐ สมาโน วาจาวิกฺเขปํ อาปชฺชติ อมราวิกฺเขปํ “อตฺถิ ปโร โลโก”ติ อิติ เจ มํ ปุจฺฉสิ, “อตฺถิ ปโร โลโก”ติ อิติ เจ \csromanfolio{26}เม อสฺส, “อตฺถิ ปโร โลโก”ติ อิติ เต นํ พฺยากเรยฺยํ, เอวนฺติปิ เม โน, ตถาติปิ เม โน, อญฺญถาติปิ เม โน, โนติปิ เม โน, โน โนติปิ เม โนติ.\csromanspacer{} นตฺถิ ปโร โลโก ฯเปฯ.\csromanspacer{} อตฺถิ จ นตฺถิ จ ปโร โลโก ฯเปฯ.\csromanspacer{} เนวตฺถิ น นตฺถิ ปโร โลโก ฯเปฯ.\csromanspacer{} อตฺถิ สตฺตา โอปปาติกา ฯเปฯ.\csromanspacer{} นตฺถิ สตฺตา โอปปาติกา ฯเปฯ.\csromanspacer{} อตฺถิ จ นตฺถิ จ สตฺตา โอปปาติกา ฯเปฯ.\csromanspacer{} เนวตฺถิ น นตฺถิ สตฺตา โอปปาติกา ฯเปฯ.\csromanspacer{} อตฺถิ สุกตทุกฺกฏานํ\footnote{สุกฏทุกฺกฏานํ (สี, สฺยา, กํ)} กมฺมานํ ผลํ วิปาโก ฯเปฯ.\csromanspacer{} นตฺถิ สุกตทุกฺกฏานํ กมฺมานํ ผลํ วิปาโก ฯเปฯ.\csromanspacer{} อตฺถิ จ นตฺถิ จ สุกตทุกฺกฏานํ กมฺมานํ ผลํ วิปาโก ฯเปฯ.\csromanspacer{} เนวตฺถิ น นตฺถิ สุกตทุกฺกฏานํ กมฺมานํ ผลํ วิปาโก ฯเปฯ.\csromanspacer{} โหติ ตถาคโต ปรํ มรณา ฯเปฯ.\csromanspacer{} น โหติ ตถาคโต ปรํ มรณา ฯเปฯ.\csromanspacer{} โหติ จ น จ โหติ\footnote{น โหติ จ (สี, ก)} ตถาคโต ปรํ มรณา ฯเปฯ.\csromanspacer{} เนว โหติ น น โหติ ตถาคโต ปรํ มรณาติ อิติ เจ มํ ปุจฺฉสิ, “เนว โหติ น น โหติ ตถาคโต ปรํ มรณา”ติ อิติ เจ เม อสฺส, “เนว โหติ น น โหติ ตถาคโต ปรํ มรณา”ติ อิติ เต นํ พฺยากเรยฺยํ, เอวนฺติปิ เม โน, ตถาติปิ เม โน, อญฺญถาติปิ เม โน, โนติปิ เม โน, โน โนติปิ เม โนติ.\csromanspacer{} อิทํ ภิกฺขเว จตุตฺถํ ฐานํ, ยํ อาคมฺม ยํ อารพฺภ เอเก สมณพฺราหฺมณา อมราวิกฺเขปิกา ตตฺถ ตตฺถ ปญฺหํ ปุฏฺฐา สมานา วาจาวิกฺเขปํ อาปชฺชนฺติ อมราวิกฺเขปํ. (4=16)}

\proseitem{66}{อิเมหิ โข เต ภิกฺขเว สมณพฺราหฺมณา อมราวิกฺเขปิกา ตตฺถ ตตฺถ ปญฺหํ ปุฏฺฐา สมานา วาจาวิกฺเขปํ อาปชฺชนฺติ อมราวิกฺเขปํ จตูหิ วตฺถูหิ.\csromanspacer{} เย หิ เกจิ ภิกฺขเว สมณา วา พฺราหฺมณา วา อมราวิกฺเขปิกา ตตฺถ ตตฺถ ปญฺหํ ปุฏฺฐา สมานา วาจาวิกฺเขปํ อาปชฺชนฺติ อมราวิกฺเขปํ, สพฺเพ เต อิเมเหว จตูหิ วตฺถูหิ เอเตสํ วา อญฺญตเรน, นตฺถิ อิโต พหิทฺธา ฯเปฯ เยหิ ตถาคตสฺส ยถาภุจฺจํ วณฺณํ สมฺมา วทมานา วเทยฺยุํ.}

\csromanmatikaanchor{mtk.11}
\csromantocmarkref{section}{อธิจฺจสมุปฺปนฺนวาท}{mtk.11}
\bookmark[dest={mtk.11},level=1]{อธิจฺจสมุปฺปนฺนวาท}
\csromanheader{1}{\textbf{อธิจฺจสมุปฺปนฺนวาท}}
\proseitem{67}{สนฺติ ภิกฺขเว เอเก สมณพฺราหฺมณา อธิจฺจสมุปฺปนฺนิกา อธิจฺจสมุปฺปนฺนํ อตฺตานญฺจ โลกญฺจ ปญฺญเปนฺติ ทฺวีหิ วตฺถูหิ.\csromanspacer{} เต จ โภนฺโต \csromanfolio{27}สมณพฺราหฺมณา กิมาคมฺม กิมารพฺภ อธิจฺจสมุปฺปนฺนิกา อธิจฺจสมุปฺปนฺนํ อตฺตานญฺจ โลกญฺจ ปญฺญเปนฺติ ทฺวีหิ วตฺถูหิ?}

\proseitem{68}{สนฺติ ภิกฺขเว อสญฺญสตฺตา นาม เทวา, สญฺญุปฺปาทา จ ปน เต เทวา ตมฺหา กายา จวนฺติ.\csromanspacer{} ฐานํ โข ปเนตํ ภิกฺขเว วิชฺชติ, ยํ อญฺญตโร สตฺโต ตมฺหา กายา จวิตฺวา อิตฺถตฺตํ อาคจฺฉติ, อิตฺถตฺตํ อาคโต สมาโน อคารสฺมา อนคาริยํ ปพฺพชติ, อคารสฺมา อนคาริยํ ปพฺพชิโต สมาโน อาตปฺปมนฺวาย ปธานมนฺวาย อนุโยคมนฺวาย อปฺปมาทมนฺวาย สมฺมา\-มนสิการม\-นฺวาย ตถารูปํ เจโตสมาธิํ ผุสติ, ยถาสมาหิเต จิตฺเต สญฺญุปฺปาทํ อนุสฺสรติ, ตโต ปรํ นานุสฺสรติ.\csromanspacer{} โส เอวมาห “อธิจฺจสมุปฺปนฺโน อตฺตา จ โลโก จ, ตํ กิสฺส เหตุ, อหํ หิ ปุพฺเพ นาโหสิํ, โสมฺหิ เอตรหิ อหุตฺวา สนฺตตาย ปริณโต”ติ.\csromanspacer{} อิทํ ภิกฺขเว ปฐมํ ฐานํ, ยํ อาคมฺม ยํ อารพฺภ เอเก สมณพฺราหฺมณา อธิจฺจสมุปฺปนฺนิกา อธิจฺจสมุปฺปนฺนํ อตฺตานญฺจ โลกญฺจ ปญฺญเปนฺติ. (1=17)}

\proseitem{69}{ทุติเย จ โภนฺโต สมณพฺราหฺมณา กิมาคมฺม กิมารพฺภ อธิจฺจสมุปฺปนฺนิกา อธิจฺจสมุปฺปนฺนํ อตฺตานญฺจ โลกญฺจ ปญฺญเปนฺติ? อิธ ภิกฺขเว เอกจฺโจ สมโณ วา พฺราหฺมโณ วา ตกฺกี โหติ วีมํสี, โส ตกฺกปริยาหตํ วีมํสานุจริตํ สยํปฏิภานํ เอวมาห “อธิจฺจสมุปฺปนฺโน อตฺตา จ โลโก จา”ติ.\csromanspacer{} อิทํ ภิกฺขเว ทุติยํ ฐานํ, ยํ อาคมฺม ยํ อารพฺภ เอเก สมณพฺราหฺมณา อธิจฺจสมุปฺปนฺนิกา อธิจฺจสมุปฺปนฺนํ อตฺตานญฺจ โลกญฺจ ปญฺญเปนฺติ. (2=18)}

\proseitem{70}{อิเมหิ โข เต ภิกฺขเว สมณพฺราหฺมณา อธิจฺจสมุปฺปนฺนิกา อธิจฺจสมุปฺปนฺนํ อตฺตานญฺจ โลกญฺจ ปญฺญเปนฺติ ทฺวีหิ วตฺถูหิ, เย หิ เกจิ ภิกฺขเว สมณา วา พฺราหฺมณา วา อธิจฺจสมุปฺปนฺนิกา อธิจฺจสมุปฺปนฺนํ อตฺตานญฺจ โลกญฺจ ปญฺญเปนฺติ, สพฺเพ เต อิเมเหว ทฺวีหิ วตฺถูหิ, เอเตสํ วา อญฺญตเรน, นตฺถิ อิโต พหิทฺธา ฯเปฯ เยหิ ตถาคตสฺส ยถาภุจฺจํ วณฺณํ สมฺมา วทมานา วเทยฺยุํ.}

\proseitem{71}{อิเมหิ โข เต ภิกฺขเว สมณพฺราหฺมณา ปุพฺพนฺตกปฺปิกา ปุพฺพนฺตานุทิฏฺฐิโน ปุพฺพนฺตํ อารพฺภ อเนกวิหิตานิ อธิมุตฺติปทานิ อภิวทนฺติ \csromanfolio{28}อฏฺฐารสหิ วตฺถูหิ.\csromanspacer{} เย หิ เกจิ ภิกฺขเว สมณา วา พฺราหฺมณา วา ปุพฺพนฺตกปฺปิกา ปุพฺพนฺตานุทิฏฺฐิโน ปุพฺพนฺตมารพฺภ อเนกวิหิตานิ อธิมุตฺติปทานิ อภิวทนฺติ, สพฺเพ เต อิเมเหว อฏฺฐารสหิ วตฺถูหิ, เอเตสํ วา อญฺญตเรน, นตฺถิ อิโต พหิทฺธา.}

\proseitem{72}{ตยิทํ ภิกฺขเว ตถาคโต ปชานาติ “อิเม ทิฏฺฐิฏฺฐานา เอวํคหิตา เอวํปรามฏฺฐา เอวํคติกา ภวนฺติ เอวํ-อภิสมฺปรายา”ติ, ตญฺจ ตถาคโต ปชานาติ, ตโต จ อุตฺตริตรํ ปชานาติ, ตญฺจ ปชานนํ น ปรามสติ, อปรามสโต จสฺส ปจฺจตฺตญฺเญว นิพฺพุติ วิทิตา, เวทนานํ สมุทยญฺจ อตฺถงฺคมญฺจ อสฺสาทญฺจ อาทีนวญฺจ นิสฺสรณญฺจ ยถาภูตํ วิทิตฺวา อนุปาทาวิมุตฺโต ภิกฺขเว ตถาคโต.}

\proseitem{73}{อิเม โข เต ภิกฺขเว ธมฺมา คมฺภีรา ทุทฺทสา ทุรนุโพธา สนฺตา ปณีตา อตกฺกาวจรา นิปุณา ปณฺฑิตเวทนียา, เย ตถาคโต สยํ อภิญฺญา สจฺฉิกตฺวา ปเวเทติ, เยหิ ตถาคตสฺส ยถาภุจฺจํ วณฺณํ สมฺมา วทมานา วเทยฺยุํ.}

\csromanheader{2}{ทุติยภาณวาโร.}
\csromansectionrule
\csromanmatikaanchor{mtk.12}
\csromantocmarkref{section}{อปรนฺตกปฺปิก}{mtk.12}
\bookmark[dest={mtk.12},level=1]{อปรนฺตกปฺปิก}
\csromanheader{1}{\textbf{อปรนฺตกปฺปิก}}
\proseitem{74}{สนฺติ ภิกฺขเว เอเก สมณพฺราหฺมณา อปรนฺตกปฺปิกา อปรนฺตานุทิฏฺฐิโน อปรนฺตํ อารพฺภ อเนกวิหิตานิ อธิมุตฺติปทานิ อภิวทนฺติ จตุจตฺตารีสาย\footnote{จตุจตฺตาลีสาย (สฺยา, กํ)} วตฺถูหิ.\csromanspacer{} เต จ โภนฺโต สมณพฺราหฺมณา กิมาคมฺม กิมารพฺภ อปรนฺตกปฺปิกา อปรนฺตานุทิฏฺฐิโน อปรนฺตํ อารพฺภ อเนกวิหิตานิ อธิมุตฺติปทานิ อภิวทนฺติ จตุจตฺตารีสาย วตฺถูหิ?}

\csromanmatikaanchor{mtk.13}
\csromantocmarkref{section}{สญฺญีวาท}{mtk.13}
\bookmark[dest={mtk.13},level=1]{สญฺญีวาท}
\csromanheader{1}{\textbf{สญฺญีวาท}}
\proseitem{75}{สนฺติ ภิกฺขเว เอเก สมณพฺราหฺมณา อุทฺธมาฆาตนิกา สญฺญีวาทา อุทฺธมาฆาตนํ สญฺญิํ อตฺตานํ ปญฺญเปนฺติ โสฬสหิ วตฺถูหิ. \csromanfolio{29}เต จ โภนฺโต สมณพฺราหฺมณา กิมาคมฺม กิมารพฺภ อุทฺธมาฆาตนิกา สญฺญีวาทา อุทฺธมาฆาตนํ สญฺญิํ อตฺตานํ ปญฺญเปนฺติ โสฬสหิ วตฺถูหิ?}

\proseitem{76}{รูปี อตฺตา โหติ อโรโค ปรํ มรณา, สญฺญีติ นํ ปญฺญเปนฺติ.\csromanspacer{} อรูปี อตฺตา โหติ อโรโค ปรํ มรณา, สญฺญีติ นํ ปญฺญเปนฺติ.\csromanspacer{} รูปี จ อรูปี จ อตฺตา โหติ ฯเปฯ.\csromanspacer{} เนวรูปี นารูปี อตฺตา โหติ.\csromanspacer{} อนฺตวา อตฺตา โหติ.\csromanspacer{} อนนฺตวา อตฺตา โหติ.\csromanspacer{} อนฺตวา จ อนนฺตวา จ อตฺตา โหติ.\csromanspacer{} เนวนฺตวา นานนฺตวา อตฺตา โหติ.\csromanspacer{} เอกตฺตสญฺญี อตฺตา โหติ.\csromanspacer{} นานตฺตสญฺญี อตฺตา โหติ.\csromanspacer{} ปริตฺตสญฺญี อตฺตา โหติ.\csromanspacer{} อปฺปมาณสญฺญี อตฺตา โหติ.\csromanspacer{} เอกนฺตสุขี อตฺตา โหติ.\csromanspacer{} เอกนฺตทุกฺขี อตฺตา โหติ.\csromanspacer{} สุขทุกฺขี อตฺตา โหติ.\csromanspacer{} อทุกฺขมสุขี อตฺตา โหติ อโรโค ปรํ มรณา, สญฺญีติ นํ ปญฺญเปนฺติ. (16=34)}

\proseitem{77}{อิเมหิ โข เต ภิกฺขเว สมณพฺราหฺมณา อุทฺธมาฆาตนิกา สญฺญีวาทา อุทฺธมาฆาตนํ สญฺญิํ อตฺตานํ ปญฺญเปนฺติ โสฬสหิ วตฺถูหิ.\csromanspacer{} เย หิ เกจิ ภิกฺขเว สมณา วา พฺราหฺมณา วา อุทฺธมาฆาตนิกา สญฺญีวาทา อุทฺธมาฆาตนํ สญฺญิํ อตฺตานํ ปญฺญเปนฺติ, สพฺเพ เต อิเมเหว โสฬสหิ วตฺถูหิ, เอเตสํ วา อญฺญตเรน, นตฺถิ อิโต พหิทฺธา ฯเปฯ เยหิ ตถาคตสฺส ยถาภุจฺจํ วณฺณํ สมฺมา วทมานา วเทยฺยุํ.}

\csromanmatikaanchor{mtk.14}
\csromantocmarkref{section}{อสญฺญีวาท}{mtk.14}
\bookmark[dest={mtk.14},level=1]{อสญฺญีวาท}
\csromanheader{1}{\textbf{อสญฺญีวาท}}
\proseitem{78}{สนฺติ ภิกฺขเว เอเก สมณพฺราหฺมณา อุทฺธมาฆาตนิกา อสญฺญีวาทา อุทฺธมาฆาตนํ อสญฺญิํ อตฺตานํ ปญฺญเปนฺติ อฏฺฐหิ วตฺถูหิ.\csromanspacer{} เต จ โภนฺโต สมณพฺราหฺมณา กิมาคมฺม กิมารพฺภ อุทฺธมาฆาตนิกา อสญฺญีวาทา อุทฺธมาฆาตนํ อสญฺญิํ อตฺตานํ ปญฺญเปนฺติ อฏฺฐหิ วตฺถูหิ?}

\proseitem{79}{รูปี อตฺตา โหติ อโรโค ปรํ มรณา, อสญฺญีติ นํ ปญฺญเปนฺติ.\csromanspacer{} อรูปี อตฺตา โหติ อโรโค ปรํ มรณา, อสญฺญีติ นํ ปญฺญเปนฺติ.\csromanspacer{} รูปี จ อรูปี จ อตฺตา โหติ ฯเปฯ.\csromanspacer{} เนวรูปี นารูปี อตฺตา โหติ.\csromanspacer{} อนฺตวา อตฺตา โหติ.\csromanspacer{} อนนฺตวา อตฺตา โหติ.\csromanspacer{} อนฺตวา จ อนนฺตวา จ อตฺตา โหติ.\csromanspacer{} เนวนฺตวา นานนฺตวา อตฺตา โหติ อโรโค ปรํ มรณา, อสญฺญีติ นํ ปญฺญเปนฺติ. (8=24=42)}

\proseitem{80}{\csromanfolio{30}อิเมหิ โข เต ภิกฺขเว สมณพฺราหฺมณา อุทฺธมาฆาตนิกา อสญฺญีวาทา อุทฺธมาฆาตนํ อสญฺญิํ อตฺตานํ ปญฺญเปนฺติ อฏฺฐหิ วตฺถูหิ.\csromanspacer{} เย หิ เกจิ ภิกฺขเว สมณา วา พฺราหฺมณา วา อุทฺธมาฆาตนิกา อสญฺญีวาทา อุทฺธมาฆาตนํ อสญฺญิํ อตฺตานํ ปญฺญเปนฺติ, สพฺเพ เต อิเมเหว อฏฺฐหิ วตฺถูหิ, เอเตสํ วา อญฺญตเรน, นตฺถิ อิโต พหิทฺธา ฯเปฯ เยหิ ตถาคตสฺส ยถาภุจฺจํ วณฺณํ สมฺมา วทมานา วเทยฺยุํ.}

\csromanmatikaanchor{mtk.15}
\csromantocmarkref{section}{เนวสญฺญีนาสญฺญีวาท}{mtk.15}
\bookmark[dest={mtk.15},level=1]{เนวสญฺญีนาสญฺญีวาท}
\csromanheader{1}{\textbf{เนวสญฺญีนาสญฺญีวาท}}
\proseitem{81}{สนฺติ ภิกฺขเว เอเก สมณพฺราหฺมณา อุทฺธมาฆาตนิกา เนวสญฺญี\-นาสญฺญี\-วาทา อุทฺธมาฆาตนํ เนวสญฺญีนาสญฺญิํ อตฺตานํ ปญฺญเปนฺติ อฏฺฐหิ วตฺถูหิ.\csromanspacer{} เต จ โภนฺโต สมณพฺราหฺมณา กิมาคมฺม กิมารพฺภ อุทฺธมาฆาตนิกา เนวสญฺญี\-นาสญฺญี\-วาทา อุทฺธมาฆาตนํ เนวสญฺญีนาสญฺญิํ อตฺตานํ ปญฺญเปนฺติ อฏฺฐหิ วตฺถูหิ?}

\proseitem{82}{รูปี อตฺตา โหติ อโรโค ปรํ มรณา, เนวสญฺญีนาสญฺญีติ นํ ปญฺญเปนฺติ.\csromanspacer{} อรูปี อตฺตา โหติ ฯเปฯ.\csromanspacer{} รูปี จ อรูปี จ อตฺตา โหติ.\csromanspacer{} เนวรูปี นารูปี อตฺตา โหติ.\csromanspacer{} อนฺตวา อตฺตา โหติ.\csromanspacer{} อนนฺตวา อตฺตา โหติ.\csromanspacer{} อนฺตวา จ อนนฺตวา จ อตฺตา โหติ.\csromanspacer{} เนวนฺตวา นานนฺตวา อตฺตา โหติ อโรโค ปรํ มรณา, เนวสญฺญีนาสญฺญีติ นํ ปญฺญเปนฺติ. (8=32=50)}

\proseitem{83}{อิเมหิ โข เต ภิกฺขเว สมณพฺราหฺมณา อุทฺธมาฆาตนิกา เนวสญฺญี\-นาสญฺญี\-วาทา อุทฺธมาฆาตนํ เนวสญฺญีนาสญฺญิํ อตฺตานํ ปญฺญเปนฺติ อฏฺฐหิ วตฺถูหิ.\csromanspacer{} เย หิ เกจิ ภิกฺขเว สมณา วา พฺราหฺมณา วา อุทฺธมาฆาตนิกา เนวสญฺญี\-นาสญฺญี\-วาทา อุทฺธมาฆาตนํ เนวสญฺญีนาสญฺญิํ อตฺตานํ ปญฺญเปนฺติ, สพฺเพ เต อิเมเหว อฏฺฐหิ วตฺถูหิ ฯเปฯ เยหิ ตถาคตสฺส ยถาภุจฺจํ วณฺณํ สมฺมา วทมานา วเทยฺยุํ.}

\csromanmatikaanchor{mtk.16}
\csromantocmarkref{section}{อุจฺเฉทวาท}{mtk.16}
\bookmark[dest={mtk.16},level=1]{อุจฺเฉทวาท}
\csromanheader{1}{\textbf{อุจฺเฉทวาท}}
\proseitem{84}{สนฺติ ภิกฺขเว เอเก สมณพฺราหฺมณา อุจฺเฉทวาทา สโต สตฺตสฺส อุจฺเฉทํ วินาสํ วิภวํ ปญฺญเปนฺติ สตฺตหิ วตฺถูหิ.\csromanspacer{} เต จ \csromanfolio{31}โภนฺโต สมณพฺราหฺมณา กิมาคมฺม กิมารพฺภ อุจฺเฉทวาทา สโต สตฺตสฺส อุจฺเฉทํ วินาสํ วิภวํ ปญฺญเปนฺติ สตฺตหิ วตฺถูหิ?}

\proseitem{85}{อิธ ภิกฺขเว เอกจฺโจ สมโณ วา พฺราหฺมโณ วา เอวํวาที โหติ เอวํทิฏฺฐิ\footnote{เอวํทิฏฺฐี (ก, อิ)} “ยโต โข โภ อยํ อตฺตา รูปี จาตุมหาภูติโก มาตาเปตฺติกสมฺภโว กายสฺส เภทา อุจฺฉิชฺชติ วินสฺสติ, น โหติ ปรํ มรณา, เอตฺตาวตา โข โภ อยํ อตฺตา สมฺมา สมุจฺฉินฺโน โหตี”ติ.\csromanspacer{} อิตฺเถเก สโต สตฺตสฺส อุจฺเฉทํ วินาสํ วิภวํ ปญฺญเปนฺติ. (1=33=51)}

\proseitem{86}{ตมญฺโญ เอวมาห “อตฺถิ โข โภ เอโส อตฺตา, ยํ ตฺวํ วเทสิ, เนโส นตฺถีติ วทามิ, โน จ โข โภ อยํ อตฺตา เอตฺตาวตา สมฺมา สมุจฺฉินฺโน โหติ.\csromanspacer{} อตฺถิ โข โภ อญฺโญ อตฺตา ทิพฺโพ รูปี กามาวจโร กพฬีการาหารภกฺโข, ตํ ตฺวํ น ชานาสิ น ปสฺสสิ, ตมหํ ชานามิ ปสฺสามิ, โส โข โภ อตฺตา ยโต กายสฺส เภทา อุจฺฉิชฺชติ วินสฺสติ, น โหติ ปรํ มรณา, เอตฺตาวตา โข โภ อยํ อตฺตา สมฺมา สมุจฺฉินฺโน โหตี”ติ.\csromanspacer{} อิตฺเถเก สโต สตฺตสฺส อุจฺเฉทํ วินาสํ วิภวํ ปญฺญเปนฺติ. (2=34=52)}

\proseitem{87}{ตมญฺโญ เอวมาห “อตฺถิ โข โภ เอโส อตฺตา, ยํ ตฺวํ วเทสิ, เนโส นตฺถีติ วทามิ, โน จ โข โภ อยํ อตฺตา เอตฺตาวตา สมฺมา สมุจฺฉินฺโน โหติ.\csromanspacer{} อตฺถิ โข โภ อญฺโญ อตฺตา ทิพฺโพ รูปี มโนมโย สพฺพงฺคปจฺจงฺคี อหีนินฺทฺริโย, ตํ ตฺวํ น ชานาสิ น ปสฺสสิ, ตมหํ ชานามิ ปสฺสามิ, โส โข โภ อตฺตา ยโต กายสฺส เภทา อุจฺฉิชฺชติ วินสฺสติ, น โหติ ปรํ มรณา, เอตฺตาวตา โข โภ อยํ อตฺตา สมฺมา สมุจฺฉินฺโน โหตี”ติ.\csromanspacer{} อิตฺเถเก สโต สตฺตสฺส อุจฺเฉทํ วินาสํ วิภวํ ปญฺญเปนฺติ. (3=35=53)}

\proseitem{88}{ตมญฺโญ เอวมาห “อตฺถิ โข โภ เอโส อตฺตา, ยํ ตฺวํ วเทสิ, เนโส นตฺถีติ วทามิ, โน จ โข โภ อยํ อตฺตา \csromanfolio{32}เอตฺตาวตา สมฺมา สมุจฺฉินฺโน โหติ.\csromanspacer{} อตฺถิ โข โภ อญฺโญ อตฺตา สพฺพโส รูปสญฺญานํ สมติกฺกมา ปฏิฆสญฺญานํ อตฺถงฺคมา นานตฺตสญฺญานํ อมนสิการา ‘อนนฺโต อากาโส’ติ อากาสานญฺจายตนูปโค, ตํ ตฺวํ น ชานาสิ น ปสฺสสิ, ตมหํ ชานามิ ปสฺสามิ, โส โข โภ อตฺตา ยโต กายสฺส เภทา อุจฺฉิชฺชติ วินสฺสติ, น โหติ ปรํ มรณา, เอตฺตาวตา โข โภ อยํ อตฺตา สมฺมา สมุจฺฉินฺโน โหตี”ติ.\csromanspacer{} อิตฺเถเก สโต สตฺตสฺส อุจฺเฉทํ วินาสํ วิภวํ ปญฺญเปนฺติ. (4=36=54)}

\proseitem{89}{ตมญฺโญ เอวมาห “อตฺถิ โข โภ เอโส อตฺตา, ยํ ตฺวํ วเทสิ, เนโส นตฺถีติ วทามิ, โน จ โข โภ อยํ อตฺตา เอตฺตาวตา สมฺมา สมุจฺฉินฺโน โหติ.\csromanspacer{} อตฺถิ โข โภ อญฺโญ อตฺตา สพฺพโส อากาสานญฺจายตนํ สมติกฺกมฺม ‘อนนฺตํ วิญฺญาณนฺ’ติ วิญฺญาณญฺจายตนูปโค, ตํ ตฺวํ น ชานาสิ น ปสฺสสิ, ตมหํ ชานามิ ปสฺสามิ, โส โข โภ อตฺตา ยโต กายสฺส เภทา อุจฺฉิชฺชติ วินสฺสติ, น โหติ ปรํ มรณา, เอตฺตาวตา โข โภ อยํ อตฺตา สมฺมา สมุจฺฉินฺโน โหตี”ติ.\csromanspacer{} อิตฺเถเก สโต สตฺตสฺส อุจฺเฉทํ วินาสํ วิภวํ ปญฺญเปนฺติ. (5=37=55)}

\proseitem{90}{ตมญฺโญ เอวมาห “อตฺถิ โข โภ เอโส อตฺตา, ยํ ตฺวํ วเทสิ, เนโส นตฺถีติ วทามิ, โน จ โข โภ อยํ อตฺตา เอตฺตาวตา สมฺมา สมุจฺฉินฺโน โหติ.\csromanspacer{} อตฺถิ โข โภ อญฺโญ อตฺตา สพฺพโส วิญฺญาณญฺจายตนํ สมติกฺกมฺม ‘นตฺถิ กิญฺจี’ติ อากิญฺจญฺญายตนูปโค, ตํ ตฺวํ น ชานาสิ น ปสฺสสิ, ตมหํ ชานามิ ปสฺสามิ, โส โข โภ อตฺตา ยโต กายสฺส เภทา อุจฺฉิชฺชติ วินสฺสติ, น โหติ ปรํ มรณา, เอตฺตาวตา โข โภ อยํ อตฺตา สมฺมา สมุจฺฉินฺโน โหตี”ติ.\csromanspacer{} อิตฺเถเก สโต สตฺตสฺส อุจฺเฉทํ วินาสํ วิภวํ ปญฺญเปนฺติ. (6=38=56)}

\proseitem{91}{ตมญฺโญ เอวมาห “อตฺถิ โข โภ เอโส อตฺตา, ยํ ตฺวํ วเทสิ, เนโส นตฺถีติ วทามิ, โน จ โข โภ อยํ อตฺตา เอตฺตาวตา สมฺมา สมุจฺฉินฺโน โหติ.\csromanspacer{} อตฺถิ โข โภ อญฺโญ \csromanfolio{33}อตฺตา สพฺพโส อากิญฺจญฺญายตนํ สมติกฺกมฺม ‘สนฺตเมตํ ปณีตเมตนฺ’ติ เนวสญฺญานาสญฺญายตนูปโค, ตํ ตฺวํ น ชานาสิ น ปสฺสสิ, ตมหํ ชานามิ ปสฺสามิ, โส โข โภ อตฺตา ยโต กายสฺส เภทา อุจฺฉิชฺชติ วินสฺสติ, น โหติ ปรํ มรณา, เอตฺตาวตา โข โภ อยํ อตฺตา สมฺมา สมุจฺฉินฺโน โหตี”ติ.\csromanspacer{} อิตฺเถเก สโต สตฺตสฺส อุจฺเฉทํ วินาสํ วิภวํ ปญฺญเปนฺติ. (7=39=57)}

\proseitem{92}{อิเมหิ โข เต ภิกฺขเว สมณพฺราหฺมณา อุจฺเฉทวาทา สโต สตฺตสฺส อุจฺเฉทํ วินาสํ วิภวํ ปญฺญเปนฺติ สตฺตหิ วตฺถูหิ.\csromanspacer{} เย หิ เกจิ ภิกฺขเว สมณา วา พฺราหฺมณา วา อุจฺเฉทวาทา สโต สตฺตสฺส อุจฺเฉทํ วินาสํ วิภวํ ปญฺญเปนฺติ, สพฺเพ เต อิเมเหว สตฺตหิ วตฺถูหิ ฯเปฯ เยหิ ตถาคตสฺส ยถาภุจฺจํ วณฺณํ สมฺมา วทมานา วเทยฺยุํ.}

\csromanmatikaanchor{mtk.17}
\csromantocmarkref{section}{ทิฏฺฐธมฺมนิพฺพานวาท}{mtk.17}
\bookmark[dest={mtk.17},level=1]{ทิฏฺฐธมฺมนิพฺพานวาท}
\csromanheader{1}{\textbf{ทิฏฺฐธมฺมนิพฺพานวาท}}
\proseitem{93}{สนฺติ ภิกฺขเว เอเก สมณพฺราหฺมณา ทิฏฺฐธมฺมนิพฺพาน\-วาทา สโต สตฺตสฺส ปรมทิฏฺฐธมฺมนิพฺพานํ ปญฺญเปนฺติ ปญฺจหิ วตฺถูหิ.\csromanspacer{} เต จ โภนฺโต สมณพฺราหฺมณา กิมาคมฺม กิมารพฺภ ทิฏฺฐธมฺมนิพฺพาน\-วาทา สโต สตฺตสฺส ปรมทิฏฺฐธมฺมนิพฺพานํ ปญฺญเปนฺติ ปญฺจหิ วตฺถูหิ?}

\proseitem{94}{อิธ ภิกฺขเว เอกจฺโจ สมโณ วา พฺราหฺมโณ วา เอวํวาที โหติ เอวํทิฏฺฐิ “ยโต โข โภ อยํ อตฺตา ปญฺจหิ กามคุเณหิ สมปฺปิโต สมงฺคีภูโต ปริจาเรติ, เอตฺตาวตา โข โภ อยํ อตฺตา ปรมทิฏฺฐธมฺมนิพฺพานํ ปตฺโต โหตี”ติ.\csromanspacer{} อิตฺเถเก สโต สตฺตสฺส ปรมทิฏฺฐธมฺมนิพฺพานํ ปญฺญเปนฺติ. (1=40=58)}

\proseitem{95}{ตมญฺโญ เอวมาห “อตฺถิ โข โภ เอโส อตฺตา, ยํ ตฺวํ วเทสิ, เนโส นตฺถีติ วทามิ, โน จ โข โภ อยํ อตฺตา เอตฺตาวตา ปรมทิฏฺฐธมฺมนิพฺพานํ ปตฺโต โหติ, ตํ กิสฺส เหตุ, กามา หิ โภ อนิจฺจา ทุกฺขา วิปริณามธมฺมา, เตสํ วิปริณามญฺญถาภาวา อุปฺปชฺชนฺติ โสกปริเทวทุกฺขโทมนสฺสุปายาสา.\csromanspacer{} ยโต โข โภ อยํ อตฺตา วิวิจฺเจว กาเมหิ วิวิจฺจ อกุสเลหิ ธมฺเมหิ สวิตกฺกํ สวิจารํ วิเวกชํ ปีติสุขํ ปฐมํ ฌานํ อุปสมฺปชฺช วิหรติ, เอตฺตาวตา โข โภ \csromanfolio{34}อยํ อตฺตา ปรมทิฏฺฐธมฺมนิพฺพานํ ปตฺโต โหตี”ติ.\csromanspacer{} อิตฺเถเก สโต สตฺตสฺส ปรมทิฏฺฐธมฺมนิพฺพานํ ปญฺญเปนฺติ. (2=41=59)}

\proseitem{96}{ตมญฺโญ เอวมาห “อตฺถิ โข โภ เอโส อตฺตา, ยํ ตฺวํ วเทสิ, เนโส นตฺถีติ วาทามิ, โน จ โข โภ อยํ อตฺตา เอตฺตาวตา ปรมทิฏฺฐธมฺมนิพฺพานํ ปตฺโต โหติ, ตํ กิสฺส เหตุ, ยเทว ตตฺถ วิตกฺกิตํ วิจาริตํ, เอเตเนตํ โอฬาริกํ อกฺขายติ.\csromanspacer{} ยโต โข โภ อยํ อตฺตา วิตกฺกวิจารานํ วูปสมา อชฺฌตฺตํ สมฺปสาทนํ เจตโส เอโกทิภาวํ อวิตกฺกํ อวิจารํ สมาธิชํ ปีติสุขํ ทุติยํ ฌานํ อุปสมฺปชฺช วิหรติ, เอตฺตาวตา โข โภ อยํ อตฺตา ปรมทิฏฺฐธมฺมนิพฺพานํ ปตฺโต โหตี”ติ.\csromanspacer{} อิตฺเถเก สโต สตฺตสฺส ปรมทิฏฺฐธมฺมนิพฺพานํ ปญฺญเปนฺติ. (3=42=60)}

\proseitem{97}{ตมญฺโญ เอวมาห “อตฺถิ โข โภ เอโส อตฺตา, ยํ ตฺวํ วเทสิ, เนโส นตฺถีติ วทามิ, โน จ โข โภ อยํ อตฺตา เอตฺตาวตา ปรมทิฏฺฐธมฺมนิพฺพานํ ปตฺโต โหติ, ตํ กิสฺส เหตุ, ยเทว ตตฺถ ปีติคตํ เจตโส อุปฺปิลาวิตตฺตํ, เอเตเนตํ โอฬาริกํ อกฺขายติ.\csromanspacer{} ยโต โข โภ อยํ อตฺตา ปีติยา จ วิราคา อุเปกฺขโก จ วิหรติ, สโต จ สมฺปชาโน, สุขญฺจ กาเยน ปฏิสํเวเทติ, ยํ ตํ อริยา อาจิกฺขนฺติ ‘อุเปกฺขโก สติมา สุขวิหารี’ติ, ตติยํ ฌานํ อุปสมฺปชฺช วิหรติ, เอตฺตาวตา โข โภ อยํ อตฺตา ปรมทิฏฺฐธมฺมนิพฺพานํ ปตฺโต โหตี”ติ.\csromanspacer{} อิตฺเถเก สโต สตฺตสฺส ปรมทิฏฺฐธมฺมนิพฺพานํ ปญฺญเปนฺติ. (4=43=61)}

\proseitem{98}{ตมญฺโญ เอวมาห “อตฺถิ โข โภ เอโส อตฺตา, ยํ ตฺวํ วเทสิ, เนโส นตฺถีติ วทามิ, โน จ โข โภ อยํ อตฺตา เอตฺตาวตา ปรมทิฏฺฐธมฺมนิพฺพานํ ปตฺโต โหติ, ตํ กิสฺส เหตุ, ยเทว ตตฺถ สุขมิติ เจตโส อาโภโค, เอเตเนตํ โอฬาริกํ อกฺขายติ.\csromanspacer{} ยโต โข โภ อยํ อตฺตา สุขสฺส จ ปหานา ทุกฺขสฺส จ ปหานา ปุพฺเพว โสมนสฺสโทมนสฺสานํ อตฺถงฺคมา อทุกฺขมสุขํ อุเปกฺขาสติปาริสุทฺธิํ จตุตฺถํ ฌานํ อุปสมฺปชฺช วิหรติ, เอตฺตาวตา โข โภ อยํ อตฺตา ปรมทิฏฺฐธมฺมนิพฺพานํ ปตฺโต โหตี”ติ.\csromanspacer{} อิตฺเถเก สโต สตฺตสฺส ปรมทิฏฺฐธมฺมนิพฺพานํ ปญฺญเปนฺติ. (5=44=62)}

\proseitem{99}{\csromanfolio{35}อิเมหิ โข เต ภิกฺขเว สมณพฺราหฺมณา ทิฏฺฐธมฺมนิพฺพาน\-วาทา สโต สตฺตสฺส ปรมทิฏฺฐธมฺมนิพฺพานํ ปญฺญเปนฺติ ปญฺจหิ วตฺถูหิ.\csromanspacer{} เย หิ เกจิ ภิกฺขเว สมณา วา พฺราหฺมณา วา ทิฏฺฐธมฺมนิพฺพาน\-วาทา สโต สตฺตสฺส ปรมทิฏฺฐธมฺมนิพฺพานํ ปญฺญเปนฺติ, สพฺเพ เต อิเมเหว ปญฺจหิ วตฺถูหิ ฯเปฯ เยหิ ตถาคตสฺส ยถาภุจฺจํ วณฺณํ สมฺมา วทมานา วเทยฺยุํ.}

\proseitem{100}{อิเมหิ โข เต ภิกฺขเว สมณพฺราหฺมณา อปรนฺตกปฺปิกา อปรนฺตานุทิฏฺฐิโน อปรนฺตํ อารพฺภ อเนกวิหิตานิ อธิมุตฺติปทานิ อภิวทนฺติ จตุจตฺตารีสาย วตฺถูหิ.\csromanspacer{} เย หิ เกจิ ภิกฺขเว สมณา วา พฺราหฺมณา วา อปรนฺตกปฺปิกา อปรนฺตานุทิฏฺฐิโน อปรนฺตํ อารพฺภ อเนกวิหิตานิ อธิมุตฺติปทานิ อภิวทนฺติ, สพฺเพ เต อิเมเหว จตุจตฺตารีสาย วตฺถูหิ ฯเปฯ เยหิ ตถาคตสฺส ยถาภุจฺจํ วณฺณํ สมฺมา วทมานา วเทยฺยุํ.}

\proseitem{101}{อิเมหิ โข เต ภิกฺขเว สมณพฺราหฺมณา ปุพฺพนฺตกปฺปิกา จ อปรนฺตกปฺปิกา จ ปุพฺพนฺตา\-ปรนฺต\-กปฺปิกา จ ปุพฺพนฺตา\-ปรนฺตา\-นุทิฏฺฐิโน ปุพฺพนฺตาปรนฺตํ อารพฺภ อเนกวิหิตานิ อธิมุตฺติปทานิ อภิวทนฺติ ทฺวาสฏฺฐิยา วตฺถูหิ.}

\proseitem{102}{เย หิ เกจิ ภิกฺขเว สมณา วา พฺราหฺมณา วา ปุพฺพนฺตกปฺปิกา วา อปรนฺตกปฺปิกา วา ปุพฺพนฺตา\-ปรนฺต\-กปฺปิกา วา ปุพฺพนฺตา\-ปรนฺตา\-นุทิฏฺฐิโน ปุพฺพนฺตาปรนฺตํ อารพฺภ อเนกวิหิตานิ อธิมุตฺติปทานิ อภิวทนฺติ, สพฺเพ เต อิเมเหว ทฺวาสฏฺฐิยา วตฺถูหิ, เอเตสํ วา อญฺญตเรน, นตฺถิ อิโต พหิทฺธา.}

\proseitem{103}{ตยิทํ ภิกฺขเว ตถาคโต ปชานาติ “อิเม ทิฏฺฐิฏฺฐานา เอวํคหิตา เอวํปรามฏฺฐา เอวํคติกา ภวนฺติ เอวํ-อภิสมฺปรายา”ติ, ตญฺจ ตถาคโต ปชานาติ, ตโต จ อุตฺตริตรํ ปชานาติ, ตญฺจ ปชานนํ น ปรามสติ, อปรามสโต จสฺส ปจฺจตฺตญฺเญว นิพฺพุติ วิทิตา.\csromanspacer{} เวทนานํ สมุทยญฺจ อตฺถงฺคมญฺจ อสฺสาทญฺจ อาทีนวญฺจ นิสฺสรณญฺจ ยถาภูตํ วิทิตฺวา อนุปาทาวิมุตฺโต ภิกฺขเว ตถาคโต.}

\csromanmatikaanchor{mtk.18}
\csromantocmarkref{section}{ปริตสฺสิตวิปฺผนฺทิตวาท}{mtk.18}
\bookmark[dest={mtk.18},level=1]{ปริตสฺสิตวิปฺผนฺทิตวาท}
\proseitem{104}{\csromanfolio{36}อิเม โข เต ภิกฺขเว ธมฺมา คมฺภีรา ทุทฺทสา ทุรนุโพธา สนฺตา ปณีตา อตกฺกาวจรา นิปุณา ปณฺฑิตเวทนียา, เย ตถาคโต สยํ อภิญฺญา สจฺฉิกตฺวา ปเวเทติ, เยหิ ตถาคตสฺส ยถาภุจฺจํ วณฺณํ สมฺมา วทมานา วเทยฺยุํ.}

\csromanheader{2}{ปริตสฺสิตวิปฺผนฺทิต\-วาร}
\proseitem{105}{ตตฺร ภิกฺขเว เย เต สมณพฺราหฺมณา สสฺสตวาทา สสฺสตํ อตฺตานญฺจ โลกญฺจ ปญฺญเปนฺติ จตูหิ วตฺถูหิ, ตทปิ เตสํ ภวตํ สมณพฺราหฺมณานํ อชานตํ อปสฺสตํ เวทยิตํ ตณฺหาคตานํ ปริตสฺสิตวิปฺผนฺทิตเม\-ว.}

\proseitem{106}{ตตฺร ภิกฺขเว เย เต สมณพฺราหฺมณา เอกจฺจสสฺสติกา เอกจฺจ- อสสฺสติกา เอกจฺจํ สสฺสตํ เอกจฺจํ อสสฺสตํ อตฺตานญฺจ โลกญฺจ ปญฺญเปนฺติ จตูหิ วตฺถูหิ, ตทปิ เตสํ ภวตํ สมณพฺราหฺมณานํ อชานตํ อปสฺสตํ เวทยิตํ ตณฺหาคตานํ ปริตสฺสิตวิปฺผตเมว.}

\proseitem{107}{ตตฺร ภิกฺขเว เย เต สมณพฺราหฺมณา อนฺตานนฺติกา อนฺตานนฺตํ โลกสฺส ปญฺญเปนฺติ จตูหิ วตฺถูหิ, ตทปิ เตสํ ภวตํ สมณพฺราหฺมณานํ อชานตํ อปสฺสตํ เวทยิตํ ตณฺหาคตานํ ปริตสฺสิตวิปฺผนฺทิตเม\-ว.}

\proseitem{108}{ตตฺร ภิกฺขเว เย เต สมณพฺราหฺมณา อมราวิกฺเขปิกา ตตฺถ ตตฺถ ปญฺหํ ปุฏฺฐา สมานา วาจาวิกฺเขปํ อาปชฺชนฺติ อมราวิกฺเขปํ จตูหิ วตฺถูหิ, ตทปิ เตสํ ภวตํ สมณพฺราหฺมณานํ อชานตํ อปสฺสตํ เวทยิตํ ตณฺหาคตานํ ปริตสฺสิตวิปฺผนฺทิตเม\-ว.}

\proseitem{109}{ตตฺร ภิกฺขเว เย เต สมณพฺราหฺมณา อธิจฺจสมุปฺปนฺนิกา อธิจฺจสมุปฺปนฺนํ อตฺตานญฺจ โลกญฺจ ปญฺญเปนฺติ ทฺวีหิ วตฺถูหิ, ตทปิ เตสํ ภวตํ สมณพฺราหฺมณานํ อชานตํ อปสฺสตํ เวทยิตํ ตณฺหาคตานํ ปริตสฺสิตวิปฺผนฺทิตเม\-ว.}

\proseitem{110}{ตตฺร ภิกฺขเว เย เต สมณพฺราหฺมณา ปุพฺพนฺตกปฺปิกา ปุพฺพนฺตานุทิฏฺฐิโน ปุพฺพนฺตํ อารพฺภ อเนกวิหิตานิ อธิมุตฺติปทานิ อภิวทนฺติ \csromanfolio{37}อฏฺฐารสหิ วตฺถูหิ, ตทปิ เตสํ ภวตํ สมณพฺราหฺมณานํ อชานตํ อปสฺสตํ เวทยิตํ ตณฺหาคตานํ ปริตสฺสิตวิปฺผนฺทิตเม\-ว.}

\proseitem{111}{ตตฺร ภิกฺขเว เย เต สมณพฺราหฺมณา อุทฺธมาฆาตนิกา สญฺญีวาทา อุทฺธมาฆาตนํ สญฺญิํ อตฺตานํ ปญฺญเปนฺติ โสฬสหิ วตฺถูหิ, ตทปิ เตสํ ภวตํ สมณพฺราหฺมณานํ อชานตํ อปสฺสตํ เวทยิตํ ตณฺหาคตานํ ปริตสฺสิตวิปฺผนฺทิตเม\-ว.}

\proseitem{112}{ตตฺร ภิกฺขเว เย เต สมณพฺราหฺมณา อุทฺธมาฆาตนิกา อสญฺญีวาทา อุทฺธมาฆาตนํ อสญฺญิํ อตฺตานํ ปญฺญเปนฺติ อฏฺฐหิ วตฺถูหิ, ตทปิ เตสํ ภวตํ สมณพฺราหฺมณานํ อชานตํ อปสฺสตํ เวทยิตํ ตณฺหาคตานํ ปริตสฺสิตวิปฺผนฺทิตเม\-ว.}

\proseitem{113}{ตตฺร ภิกฺขเว เย เต สมณพฺราหฺมณา อุทฺธมาฆาตนิกา เนวสญฺญี\-นาสญฺญี\-วาทา อุทฺธมาฆาตนํ เนวสญฺญีนาสญฺญิํ อตฺตานํ ปญฺญเปนฺติ อฏฺฐหิ วตฺถูหิ, ตทปิ เตสํ ภวตํ สมณพฺราหฺมณานํ อชานตํ อปสฺสตํ เวทยิตํ ตณฺหาคตานํ ปริตสฺสิตวิปฺผนฺทิตเม\-ว.}

\proseitem{114}{ตตฺร ภิกฺขเว เย เต สมณพฺราหฺมณา อุจฺเฉทวาทา สโต สตฺตสฺส อุจฺเฉทํ วินาสํ วิภวํ ปญฺญเปนฺติ สตฺตหิ วตฺถูหิ, ตทปิ เตสํ ภวตํ สมณพฺราหฺมณานํ อชานตํ อปสฺสตํ เวทยิตํ ตณฺหาคตานํ ปริตสฺสิตวิปฺผนฺทิตเม\-ว.}

\proseitem{115}{ตตฺร ภิกฺขเว เย เต สมณพฺราหฺมณา ทิฏฺฐธมฺมนิพฺพาน\-วาทา สโต สตฺตสฺส ปรมทิฏฺฐธมฺมนิพฺพานํ ปญฺญเปนฺติ ปญฺจหิ วตฺถูหิ, ตทปิ เตสํ ภวตํ สมณพฺราหฺมณานํ อชานตํ อปสฺสตํ เวทยิตํ ตณฺหาคตานํ ปริตสฺสิตวิปฺผนฺทิตเม\-ว.}

\proseitem{116}{ตตฺร ภิกฺขเว เย เต สมณพฺราหฺมณา อปรนฺตกปฺปิกา อปรนฺตานุทิฏฺฐิโน อปรนฺตํ อารพฺภ อเนกวิหิตานิ อธิมุตฺติปทานิ อภิวทนฺติ จตุจตฺตารีสาย วตฺถูหิ, ตทปิ เตสํ ภวตํ สมณพฺราหฺมณานํ อชานตํ อปสฺสตํ เวทยิตํ ตณฺหาคตานํ ปริตสฺสิตวิปฺผนฺทิตเม\-ว.}

\csromanmatikaanchor{mtk.19}
\csromantocmarkref{section}{ผสฺสปจฺจยาวาท}{mtk.19}
\bookmark[dest={mtk.19},level=1]{ผสฺสปจฺจยาวาท}
\proseitem{117}{\csromanfolio{38}ตตฺร ภิกฺขเว เย เต สมณพฺราหฺมณา ปุพฺพนฺตกปฺปิกา จ อปรนฺตกปฺปิกา จ ปุพฺพนฺตา\-ปรนฺต\-กปฺปิกา จ ปุพฺพนฺตา\-ปรนฺตา\-นุทิฏฺฐิโน ปุพฺพนฺตาปรนฺตํ อารพฺภ อเนกวิหิตานิ อธิมุตฺติปทานิ อภิวทนฺติ ทฺวาสฏฺฐิยา วตฺถูหิ, ตทปิ เตสํ ภวตํ สมณพฺราหฺมณานํ อชานตํ อปสฺสตํ เวทยิตํ ตณฺหาคตานํ ปริตสฺสิตวิปฺผนฺทิตเม\-ว.}

\csromanheader{2}{\textbf{ผสฺสปจฺจยาวาร}}
\proseitem{118}{ตตฺร ภิกฺขเว เย เต สมณพฺราหฺมณา สสฺสตวาทา สสฺสตํ อตฺตานญฺจ โลกญฺจ ปญฺญเปนฺติ จตูหิ วตฺถูหิ, ตทปิ ผสฺสปจฺจยา.}

\proseitem{119}{ตตฺร ภิกฺขเว เย เต สมณพฺราหฺมณา เอกจฺจสสฺสติกา เอกจฺจ- อสสฺสติกา เอกจฺจํ สสฺสตํ เอกจฺจํ อสสฺสตํ อตฺตานญฺจ โลกญฺจ ปญฺญเปนฺติ จตูหิ วตฺถูหิ, ตทปิ ผสฺสปจฺจยา.}

\proseitem{120}{ตตฺร ภิกฺขเว เย เต สมณพฺราหฺมณา อนฺตานนฺติกา อนฺตานนฺตํ โลกสฺส ปญฺญเปนฺติ จตูหิ วตฺถูหิ, ตทปิ ผสฺสปจฺจยา.}

\proseitem{121}{ตตฺร ภิกฺขเว เย เต สมณพฺราหฺมณา อมราวิกฺเขปิกา ตตฺถ ตตฺถ ปญฺหํ ปุฏฺฐา สมานา วาจาวิกฺเขปํ อาปชฺชนฺติ อมราวิกฺเขปํ จตูหิ วตฺถูหิ, ตทปิ ผสฺสปจฺจยา.}

\proseitem{122}{ตตฺร ภิกฺขเว เย เต สมณพฺราหฺมณา อธิจฺจสมุปฺปนฺนิกา อธิจฺจสมุปฺปนฺนํ อตฺตานญฺจ โลกญฺจ ปญฺญเปนฺติ ทฺวีหิ วตฺถูหิ, ตทปิ ผสฺสปจฺจยา.}

\proseitem{123}{ตตฺร ภิกฺขเว เย เต สมณพฺราหฺมณา ปุพฺพนฺตกปฺปิกา ปุพฺพนฺตานุทิฏฺฐิโน ปุพฺพนฺตํ อารพฺภ อเนกวิหิตานิ อธิมุตฺติปทานิ อภิวทนฺติ อฏฺฐารสหิ วตฺถูหิ, ตทปิ ผสฺสปจฺจยา.}

\proseitem{124}{ตตฺร ภิกฺขเว เย เต สมณพฺราหฺมณา อุทฺธมาฆาตนิกา สญฺญีวาทา อุทฺธมาฆาตนํ สญฺญิํ อตฺตานํ ปญฺญเปนฺติ โสฬสหิ วตฺถูหิ, ตทปิ ผสฺสปจฺจยา.}

\proseitem{125}{\csromanfolio{39}ตตฺร ภิกฺขเว เย เต สมณพฺราหฺมณา อุทฺธมาฆาตนิกา อสญฺญีวาทา อุทฺธมาฆาตนํ อสญฺญิํ อตฺตานํ ปญฺญเปนฺติ อฏฺฐหิ วตฺถูหิ, ตทปิ ผสฺสปจฺจยา.}

\proseitem{126}{ตตฺร ภิกฺขเว เย เต สมณพฺราหฺมณา อุทฺธมาฆาตนิกา เนวสญฺญี\-นาสญฺญี\-วาทา อุทฺธมาฆาตนํ เนวสญฺญีนาสญฺญิํ อตฺตานํ ปญฺญเปนฺติ อฏฺฐหิ วตฺถูหิ, ตทปิ ผสฺสปจฺจยา.}

\proseitem{127}{ตตฺร ภิกฺขเว เย เต สมณพฺราหฺมณา อุจฺเฉทวาทา สโต สตฺตสฺส อุจฺเฉทํ วินาสํ วิภวํ ปญฺญเปนฺติ สตฺตหิ วตฺถูหิ, ตทปิ ผสฺสปจฺจยา.}

\proseitem{128}{ตตฺร ภิกฺขเว เย เต สมณพฺราหฺมณา ทิฏฺฐธมฺมนิพฺพาน\-วาทา สโต สตฺตสฺส ปรมทิฏฺฐธมฺมนิพฺพานํ ปญฺญเปนฺติ ปญฺจหิ วตฺถูหิ, ตทปิ ผสฺสปจฺจยา.}

\proseitem{129}{ตตฺร ภิกฺขเว เย เต สมณพฺราหฺมณา อปรนฺตกปฺปิกา อปรนฺตานุทิฏฺฐิโน อปรนฺตํ อารพฺภ อเนกวิหิตานิ อธิมุตฺติปทานิ อภิวทนฺติ จตุจตฺตารีสาย วตฺถูหิ, ตทปิ ผสฺสปจฺจยา.}

\proseitem{130}{ตตฺร ภิกฺขเว เย เต สมณพฺราหฺมณา ปุพฺพนฺตกปฺปิกา จ อปรนฺตกปฺปิกา จ ปุพฺพนฺตา\-ปรนฺต\-กปฺปิกา จ ปุพฺพนฺตา\-ปรนฺตา\-นุทิฏฺฐิโน ปุพฺพนฺตาปรนฺตํ อารพฺภ อเนกวิหิตานิ อธิมุตฺติปทานิ อภิวทนฺติ ทฺวาสฏฺฐิยา วตฺถูหิ, ตทปิ ผสฺสปจฺจยา.}

\csromanmatikaanchor{mtk.20}
\csromantocmarkref{section}{เนตํ ฐานํ วิชฺชติวาร}{mtk.20}
\bookmark[dest={mtk.20},level=1]{เนตํ ฐานํ วิชฺชติวาร}
\csromanheader{1}{\textbf{เนตํ ฐานํ วิชฺชติวาร}}
\proseitem{131}{ตตฺร ภิกฺขเว เย เต สมณพฺราหฺมณา สสฺสตวาทา สสฺสตํ อตฺตานญฺจ โลกญฺจ ปญฺญเปนฺติ จตูหิ วตฺถูหิ, เต วต อญฺญตฺร ผสฺสา ปฏิสํเวทิสฺสนฺตีติ เนตํ ฐานํ วิชฺชติ.}

\proseitem{132}{ตตฺร ภิกฺขเว เย เต สมณพฺราหฺมณา เอกจฺจสสฺสติกา เอกจฺจ- อสสฺสติกา เอกจฺจํ สสฺสตํ เอกจฺจํ อสสฺสตํ อตฺตานญฺจ โลกญฺจ ปญฺญเปนฺติ จตูหิ วตฺถูหิ, เต วต อญฺญตฺร ผสฺสา ปฏิสํเวทิสฺสนฺตีติ เนตํ ฐานํ วิชฺชติ.}

\proseitem{133}{\csromanfolio{40}ตตฺร ภิกฺขเว เย เต สมณพฺราหฺมณา อนฺตานนฺติกา อนฺตานนฺตํ โลกสฺส ปญฺญเปนฺติ จตูหิ วตฺถูหิ, เต วต อญฺญตฺร ผสฺสา ปฏิสํเวทิสฺสนฺตีติ เนตํ ฐานํ วิชฺชติ.}

\proseitem{134}{ตตฺร ภิกฺขเว เย เต สมณพฺราหฺมณา อมราวิกฺเขปิกา ตตฺถ ตตฺถ ปญฺหํ ปุฏฺฐา สมานา วาจาวิกฺเขปํ อาปชฺชนฺติ อมราวิกฺเขปํ จตูหิ วตฺถูหิ, เต วต อญฺญตฺร ผสฺสา ปฏิสํเวทิสฺสนฺตีติ เนตํ ฐานํ วิชฺชติ.}

\proseitem{135}{ตตฺร ภิกฺขเว เย เต สมณพฺราหฺมณา อธิจฺจสมุปฺปนฺนิกา อธิจฺจสมุปฺปนฺนํ อตฺตานญฺจ โลกญฺจ ปญฺญเปนฺติ ทฺวีหิ วตฺถูหิ, เต วต อญฺญตฺร ผสฺสา ปฏิสํเวทิสฺสนฺตีติ เนตํ ฐานํ วิชฺชติ.}

\proseitem{136}{ตตฺร ภิกฺขเว เย เต สมณพฺราหฺมณา ปุพฺพนฺตกปฺปิกา ปุพฺพนฺตานุทิฏฺฐิโน ปุพฺพนฺตํ อารพฺภ อเนกวิหิตานิ อธิมุตฺติปทานิ อภิวทนฺติ อฏฺฐารสหิ วตฺถูหิ, เต วต อญฺญตฺร ผสฺสา ปฏิสํเวทิสฺสนฺตีติ เนตํ ฐานํ วิชฺชติ.}

\proseitem{137}{ตตฺร ภิกฺขเว เย เต สมณพฺราหฺมณา อุทฺธมาฆาตนิกา สญฺญีวาทา อุทฺธมาฆาตนํ สญฺญิํ อตฺตานํ ปญฺญเปนฺติ โสฬสหิ วตฺถูหิ, เต วต อญฺญตฺร ผสฺสา ปฏิสํเวทิสฺสนฺตีติ เนตํ ฐานํ วิชฺชติ.}

\proseitem{138}{ตตฺร ภิกฺขเว เย เต สมณพฺราหฺมณา อุทฺธมาฆาตนิกา อสญฺญีวาทา อุทฺธมาฆาตนํ อสญฺญิํ อตฺตานํ ปญฺญเปนฺติ อฏฺฐหิ วตฺถูหิ, เต วต อญฺญตฺร ผสฺสา ปฏิสํเวทิสฺสนฺตีติ เนตํ ฐานํ วิชฺชติ.}

\proseitem{139}{ตตฺร ภิกฺขเว เย เต สมณพฺราหฺมณา อุทฺธมาฆาตนิกา เนวสญฺญี\-นาสญฺญี\-วาทา อุทฺธมาฆาตนํ เนวสญฺญีนาสญฺญิํ อตฺตานํ ปญฺญเปนฺติ อฏฺฐหิ วตฺถูหิ, เต วต อญฺญตฺร ผสฺสา ปฏิสํเวทิสฺสนฺตีติ เนตํ ฐานํ วิชฺชติ.}

\proseitem{140}{ตตฺร ภิกฺขเว เย เต สมณพฺราหฺมณา อุจฺเฉทวาทา สโต สตฺตสฺส อุจฺเฉทํ วินาสํ วิภวํ ปญฺญเปนฺติ สตฺตหิ วตฺถูหิ, เต วต อญฺญตฺร ผสฺสา ปฏิสํเวทิสฺสนฺตีติ เนตํ ฐานํ วิชฺชติ.}

\proseitem{141}{\csromanfolio{41}ตตฺร ภิกฺขเว เย เต สมณพฺราหฺมณา ทิฏฺฐธมฺมนิพฺพาน\-วาทา สโต สตฺตสฺส ปรมทิฏฺฐธมฺมนิพฺพานํ ปญฺญเปนฺติ ปญฺจหิ วตฺถูหิ, เต วต อญฺญตฺร ผสฺสา ปฏิสํเวทิสฺสนฺตีติ เนตํ ฐานํ วิชฺชติ.}

\proseitem{142}{ตตฺร ภิกฺขเว เย เต สมณพฺราหฺมณา อปรนฺตกปฺปิกา อปรนฺตานุทิฏฺฐิโน อปรนฺตํ อารพฺภ อเนกวิหิตานิ อธิมุตฺติปทานิ อภิวทนฺติ จตุจตฺตารีสาย วตฺถูหิ, เต วต อญฺญตฺร ผสฺสา ปฏิสํเวทิสฺสนฺตีติ เนตํ ฐานํ วิชฺชติ.}

\proseitem{143}{ตตฺร ภิกฺขเว เย เต สมณพฺราหฺมณา ปุพฺพนฺตกปฺปิกา จ อปรนฺตกปฺปิกา จ ปุพฺพนฺตา\-ปรนฺต\-กปฺปิกา จ ปุพฺพนฺตา\-ปรนฺตา\-นุทิฏฺฐิโน ปุพฺพนฺตาปรนฺตํ อารพฺภ อเนกวิหิตานิ อธิมุตฺติปทานิ อภิวทนฺติ ทฺวาสฏฺฐิยา วตฺถูหิ, เต วต อญฺญตฺร ผสฺสา ปฏิสํเวทิสฺสนฺตีติ เนตํ ฐานํ วิชฺชติ.}

\csromanmatikaanchor{mtk.21}
\csromantocmarkref{section}{ทิฏฺฐิคติกาธิฏฺฐานวฏฺฏกถา}{mtk.21}
\bookmark[dest={mtk.21},level=1]{ทิฏฺฐิคติกาธิฏฺฐานวฏฺฏกถา}
\csromanheader{1}{ทิฏฺฐิคติกา\-ธิฏฺฐาน\-วฏฺฏ\-กถา}
\proseitem{144}{ตตฺร ภิกฺขเว เย เต สมณพฺราหฺมณา สสฺสตวาทา สสฺสตํ อตฺตานญฺจ โลกญฺจ ปญฺญเปนฺติ จตูหิ วตฺถูหิ, เยปิ เต สมณพฺราหฺมณา เอกจฺจสสฺสติกา เอกจฺจ-อสสฺสติกา ฯเปฯ เยปิ เต สมณพฺราหฺมณา อนฺตานนฺติกา.\csromanspacer{} เยปิ เต สมณพฺราหฺมณา อมราวิกฺเขปิกา.\csromanspacer{} เยปิ เต สมณพฺราหฺมณา อธิจฺจสมุปฺปนฺนิกา.\csromanspacer{} เยปิ เต สมณพฺราหฺมณา ปุพฺพนฺตกปฺปิกา.\csromanspacer{} เยปิ เต สมณพฺราหฺมณา อุทฺธมาฆาตนิกา สญฺญีวาทา.\csromanspacer{} เยปิ เต สมณพฺราหฺมณา อุทฺธมาฆาตนิกา อสญฺญีวาทา.\csromanspacer{} เยปิ เต สมณพฺราหฺมณา อุทฺธมาฆาตนิกา เนวสญฺญี\-นาสญฺญี\-วาทา.\csromanspacer{} เยปิ เต สมณพฺราหฺมณา อุจฺเฉทวาทา.\csromanspacer{} เยปิ เต สมณพฺราหฺมณา ทิฏฺฐธมฺมนิพฺพาน\-วาทา.\csromanspacer{} เยปิ เต สมณพฺราหฺมณา อปรนฺตกปฺปิกา.\csromanspacer{} เยปิ เต สมณพฺราหฺมณา ปุพฺพนฺตกปฺปิกา จ อปรนฺตกปฺปิกา จ ปุพฺพนฺตา\-ปรนฺต\-กปฺปิกา จ ปุพฺพนฺตา\-ปรนฺตา\-นุทิฏฺฐิโน ปุพฺพนฺตาปรนฺตํ อารพฺภ อเนกวิหิตานิ อธิมุตฺติปทานิ อภิวทนฺติ ทฺวาสฏฺฐิยา วตฺถูหิ, สพฺเพ เต ฉหิ ผสฺสายตเนหิ ผุสฺส ผุสฺส ปฏิสํเวเทนฺติ, เตสํ เวทนาปจฺจยา ตณฺหา, ตณฺหาปจฺจยา อุปาทานํ, อุปาทานปจฺจยา ภโว, ภวปจฺจยา ชาติ, ชาติปจฺจยา ชรามรณํ โสกปริเทวทุกฺขโทมนสฺสุปายาสา สมฺภวนฺติ.}

\csromanmatikaanchor{mtk.22}
\csromantocmarkref{section}{วิวฏฺฏกถาทิ}{mtk.22}
\bookmark[dest={mtk.22},level=1]{วิวฏฺฏกถาทิ}
\csromanheader{1}{\textbf{วิวฏฺฏกถาทิ}}
\proseitem{145}{\csromanfolio{42}ยโต โข ภิกฺขเว ภิกฺขุ ฉนฺนํ ผสฺสายตนานํ สมุทยญฺจ อตฺถงฺคมญฺจ อสฺสาทญฺจ อาทีนวญฺจ นิสฺสรณญฺจ ยถาภูตํ ปชานาติ, อยํ อิเมหิ สพฺเพเหว อุตฺตริตรํ ปชานาติ.}

\proseitem{146}{เย หิ เกจิ ภิกฺขเว สมณา วา พฺราหฺมณา วา ปุพฺพนฺตกปฺปิกา วา อปรนฺตกปฺปิกา วา ปุพฺพนฺตา\-ปรนฺต\-กปฺปิกา วา ปุพฺพนฺตา\-ปรนฺตา\-นุทิฏฺฐิโน ปุพฺพนฺตาปรนฺตํ อารพฺภ อเนกวิหิตานิ อธิมุตฺติปทานิ อภิวทนฺติ, สพฺเพ เต อิเมเหว ทฺวาสฏฺฐิยา วตฺถูหิ อนฺโตชาลีกตา เอตฺถ สิตาว อุมฺมุชฺชมานา อุมฺมุชฺชนฺติ, เอตฺถ ปริยาปนฺนา อนฺโตชาลีกตาว อุมฺมุชฺชมานา อุมฺมุชฺชนฺติ.}

\prose{เสยฺยถาปิ ภิกฺขเว ทกฺโข เกวฏฺโฏ วา เกวฏฺฏนฺเตวาสี วา สุขุมจฺฉิเกน ชาเลน ปริตฺตํ อุทกทหํ\footnote{อุทกรหทํ (สี, สฺยา, อิ)} โอตฺถเรยฺย, ตสฺส เอวมสฺส “เย โข เกจิ อิมสฺมิํ อุทกทเห โอฬาริกา ปาณา, สพฺเพ เต อนฺโตชาลีกตา เอตฺถ สิตาว อุมฺมุชฺชมานา อุมฺมุชฺชนฺติ, เอตฺถ ปริยาปนฺนา อนฺโตชาลีกตาว อุมฺมุชฺชมานา อุมฺมุชฺชนฺตี”ติ, เอวเมว โข ภิกฺขเว เย หิ เกจิ สมณา วา พฺราหฺมณา วา ปุพฺพนฺตกปฺปิกา วา อปรนฺตกปฺปิกา วา ปุพฺพนฺตา\-ปรนฺต\-กปฺปิกา วา ปุพฺพนฺตา\-ปรนฺตา\-นุทิฏฺฐิโน ปุพฺพนฺตาปรนฺตํ อารพฺภ อเนกวิหิตานิ อธิมุตฺติปทานิ อภิวทนฺติ, สพฺเพ เต อิเมเหว ทฺวาสฏฺฐิยา วตฺถูหิ อนฺโตชาลีกตา เอตฺถ สิตาว อุมฺมุชฺชมานา อุมฺมุชฺชนฺติ, เอตฺถ ปริยาปนฺนา อนฺโตชาลีกตาว อุมฺมุชฺชมานา อุมฺมุชฺชนฺติ.}

\proseitem{147}{อุจฺฉินฺนภวเนตฺติโก ภิกฺขเว ตถาคตสฺส กาโย ติฏฺฐติ, ยาวสฺส กาโย ฐสฺสติ, ตาว นํ ทกฺขนฺติ เทวมนุสฺสา, กายสฺส เภทา อุทฺธํ ชีวิตปริยาทานา น นํ ทกฺขนฺติ เทวมนุสฺสา.}

\prose{เสยฺยถาปิ ภิกฺขเว อมฺพปิณฺฑิยา วณฺฏจฺฉินฺนาย ยานิ กานิจิ อมฺพานิ วณฺฏปฏิพนฺธานิ\footnote{วณฺฏูปนิพนฺธนานิ (สี, อิ), วณฺฑปฏิพทฺธานิ (ก)}, สพฺพานิ ตานิ ตทนฺวยานิ ภวนฺติ, เอวเมว โข ภิกฺขเว อุจฺฉินฺนภวเนตฺติโก ตถาคตสฺส กาโย ติฏฺฐติ, ยาวสฺส กาโย ฐสฺสติ, ตาว นํ ทกฺขนฺติ เทวมนุสฺสา, กายสฺส เภทา อุทฺธํ ชีวิตปริยาทานา น นํ ทกฺขนฺติ เทวมนุสฺสาติ.}

\proseitem{148}{\csromanfolio{43}เอวํ วุตฺเต อายสฺมา อานนฺโท ภควนฺตํ เอตทโวจ “อจฺฉริยํ ภนฺเต, อพฺภุตํ ภนฺเต, โกนาโม อยํ ภนฺเต ธมฺมปริยาโย”ติ.\csromanspacer{} ตสฺมาติห ตฺวํ อานนฺท อิมํ ธมฺมปริยายํ อตฺถชาลนฺติปิ นํ ธาเรหิ, ธมฺมชาลนฺติปิ นํ ธาเรหิ, พฺรหฺมชาลนฺติปิ นํ ธาเรหิ, ทิฏฺฐิชาลนฺติปิ นํ ธาเรหิ, อนุตฺตโร สงฺคามวิชโยติปิ นํ ธาเรหีติ, อิทมโวจ ภควา.}

\proseitemwithcloser{149}{อตฺตมนา เต ภิกฺขู ภควโต ภาสิตํ อภินนฺทุนฺติ, อิมสฺมิํ จ ปน เวยฺยากรณสฺมิํ ภญฺญมาเน ทสสหสฺสี\footnote{สหสฺสี (กตฺถจิ)} โลกธาตุ อกมฺปิตฺถาติ.}{พฺรหฺมชาลสุตฺตํ นิฏฺฐิตํ ปฐมํ.}
\csromanmatikaanchor{mtk.23}
\csromantocmarknopage{chapter}{2. สามญฺญผลสุตฺต}{mtk.23}
\bookmark[dest={mtk.23},level=0]{2. สามญฺญผลสุตฺต}
\setcsromanheadcha{2. สามญฺญผลสุตฺต}
\setcsromanheadhclear{}
\chapterheadpageread{2. \textbf{สามญฺญผลสุตฺต}}
\csromanmatikaanchor{mtk.24}
\csromantocmarkref{section}{ราชามจฺจกถา}{mtk.24}
\bookmark[dest={mtk.24},level=1]{ราชามจฺจกถา}
\csromanheader{1}{\textbf{ราชามจฺจกถา}}
\proseitem{150}{\csromanfolio{44}เอวํ เม สุตํ, เอกํ สมยํ ภควา ราชคเห วิหรติ ชีวกสฺส โกมารภจฺจสฺส อมฺพวเน มหตา ภิกฺขุสํเฆน สทฺธิํ อฑฺฒเตฬเสหิ ภิกฺขุสเตหิ.\csromanspacer{} เตน โข ปน สมเยน ราชา มาคโธ อชาตสตฺตุ เวเทหิปุตฺโต ตทหุโปสเถ ปนฺนรเส โกมุทิยา จาตุมาสินิยา ปุณฺณาย ปุณฺณมาย รตฺติยา ราชามจฺจปริวุโต อุปริปาสาทวรคโต นิสินฺโน โหติ.\csromanspacer{} อถ โข ราชา มาคโธ อชาตสตฺตุ เวเทหิปุตฺโต ตทหุโปสเถ อุทานํ อุทาเนสิ “รมณียา วต โภ โทสินา รตฺติ, อภิรูปา วต โภ โทสินา รตฺติ, ทสฺสนียา วต โภ โทสินา รตฺติ, ปาสาทิกา วต โภ โทสินา รตฺติ, ลกฺขญฺญา วต โภ โทสินา รตฺติ, กํ นุ ขฺวชฺช สมณํ วา พฺราหฺมณํ วา ปยิรุปาเสยฺยาม, ยํ โน ปยิรุปาสโต จิตฺตํ ปสีเทยฺยา”ติ.}

\proseitem{151}{เอวํ วุตฺเต อญฺญตโร ราชามจฺโจ ราชานํ มาคธํ อชาตสตฺตุํ เวเทหิปุตฺตํ เอตทโวจ “อยํ เทว ปูรโณ กสฺสโป สงฺฆี เจว คณี จ คณาจริโย จ ญาโต ยสสฺสี ติตฺถกโร สาธุสมฺมโต พหุชนสฺส รตฺตญฺญู จิรปพฺพชิโต อทฺธคโต วโย-อนุปฺปตฺโต, ตํ เทโว ปูรณํ กสฺสปํ ปยิรุปาสตุ, อปฺเปว นาม เทวสฺส ปูรณํ กสฺสปํ ปยิรุปาสโต จิตฺตํ ปสีเทยฺยา”ติ.\csromanspacer{} เอวํ วุตฺเต ราชา มาคโธ อชาตสตฺตุ เวเทหิปุตฺโต ตุณฺหี อโหสิ.}

\proseitem{152}{อญฺญตโรปิ โข ราชามจฺโจ ราชานํ มาคธํ อชาตสตฺตุํ เวเทหิปุตฺตํ เอตทโวจ “อยํ เทว มกฺขลิ โคสาโล สงฺฆี เจว คณี จ คณาจริโย จ ญาโต ยสสฺสี ติตฺถกโร สาธุสมฺมโต พหุชนสฺส รตฺตญฺญู จิรปพฺพชิโต อทฺธคโต วโย-อนุปฺปตฺโต, ตํ เทโว มกฺขลิํ โคสาลํ ปยิรุปาสตุ, อปฺเปว นาม เทวสฺส มกฺขลิํ โคสาลํ ปยิรุปาสโต จิตฺตํ ปสีเทยฺยา”ติ.\csromanspacer{} เอวํ วุตฺเต ราชา มาคโธ อชาตสตฺตุ เวเทหิปุตฺโต ตุณฺหี อโหสิ.}

\csromanheader{2}{2. สามญฺญาผลสุตฺต}
\proseitem{153}{\csromanfolio{45}อญฺญตโรปิ โข ราชามจฺโจ ราชานํ มาคธํ อชาตสตฺตุํ เวเทหิปุตฺตํ เอตทโวจ “อยํ เทว อชิโต เกสกมฺพโล สงฺฆี เจว คณี จ คณาจริโย จ ญาโต ยสสฺสี ติตฺถกโร สาธุสมฺมโต พหุชนสฺส รตฺตญฺญู จิรปพฺพชิโต อทฺธคโต วโย-อนุปฺปตฺโต, ตํ เทโว อชิตํ เกสกมฺพลํ ปยิรุปาสตุ, อปฺเปว นาม เทวสฺส อชิตํ เกสกมฺพลํ ปยิรุปาสตุ, อปฺเปว นาม เทวสฺส อชิตํ เกสกมฺพลํ ปยิรุปาสโต จิตฺตํ ปสีเทยฺยา”ติ.\csromanspacer{} เอวํ วุตฺเต ราชา มาคโธ อชาตสตฺตุ เวเทหิปุตฺโต ตุณฺหี อโหสิ.}

\proseitem{154}{อญฺญตโรปิ โข ราชามจฺโจ ราชานํ มาคธํ อชาตสตฺตุํ เวเทหิปุตฺตํ เอตทโวจ “อยํ เทว ปกุโธ\footnote{ปกุทฺโธ (สี)} กจฺจายโน สงฺฆี เจว คณี จ คณาจริโย จ ญาโต ยสสฺสี ติตฺถกโร สาธุสมฺมโต พหุชนสฺส รตฺตญฺญู จิรปพฺพชิโต อทฺธคโต วโย-อนุปฺปตฺโต, ตํ เทโว ปกุธํ กจฺจายนํ ปยิรุปาสตุ, อปฺเปว นาม เทวสฺส ปกุธํ กจฺจายนํ ปยิรุปาสโต จิตฺตํ ปสีเทยฺยา”ติ.\csromanspacer{} เอวํ วุตฺเต ราชา มาคโธ อชาตสตฺตุ เวเทหิปุตฺโต ตุณฺหี อโหสิ.}

\proseitem{155}{อญฺญตโรปิ โข ราชามจฺโจ ราชานํ มาคธํ อชาตสตฺตุํ เวเทหิปุตฺตํ เอตทโวจ “อยํ เทว สญฺจโย\footnote{สญฺชโย (สี, สฺยา)} เพลฏฺฐปุตฺโต\footnote{เพลฺลฏฺฐิปุตฺโต (สี), เวลฏฺฐปุตฺโต (สฺยา)} สงฺฆี เจว คณี จ คณาจริโย จ ญาโต ยสสฺสี ติตฺถกโร สาธุสมฺมโต พหุชนสฺส รตฺตญฺญู จิรปพฺพชิโต อทฺธคโต วโย-อนุปฺปตฺโต, ตํ เทโว สญฺจยํ เพลฏฺฐปุตฺตํ ปยิรุปาสตุ, อปฺเปว นาม เทวสฺส สญฺจยํ เพลฏฺฐปุตฺตํ ปยิรุปาสโต จิตฺตํ ปสีเทยฺยา”ติ.\csromanspacer{} เอวํ วุตฺเต ราชา มาคโธ อชาตสตฺตุ เวเทหิปุตฺโต ตุณฺหี อโหสิ.}

\proseitem{156}{อญฺญตโรปิ โข ราชามจฺโจ ราชานํ มาคธํ อชาตสตฺตุํ เวเทหิปุตฺตํ เอตทโวจ “อยํ เทว นิคณฺโฐ นาฏปุตฺโต\footnote{นาถปุตฺโต (สี), นาตปุตฺโต (อิ)} สงฺฆี เจว คณี จ คณาจริโย จ ญาโต ยสสฺสี ติตฺถกโร สาธุสมฺมโต พหุชนสฺส รตฺตญฺญู จิรปพฺพชิโต อทฺธคโต วโย-อนุปฺปตฺโต, ตํ เทโว นิคณฺฐํ นาฏปุตฺตํ ปยิรุปาสตุ, อปฺเปว นาม เทวสฺส \csromanfolio{46}นิคณฺฐํ นาฏปุตฺตํ ปยิรุปาสโต จิตฺตํ ปสีเทยฺยา”ติ.\csromanspacer{} เอวํ วุตฺเต ราชา มาคโธ อชาตสตฺตุ เวเทหิปุตฺโต ตุณฺหี อโหสิ.}

\csromanmatikaanchor{mtk.25}
\csromantocmarkref{section}{โกมารภจฺจชีวกกถา}{mtk.25}
\bookmark[dest={mtk.25},level=1]{โกมารภจฺจชีวกกถา}
\csromanheader{1}{\textbf{โกมารภจฺจชีวกกถา}}
\proseitem{157}{เตน โข ปน สมเยน ชีวโก โกมารภจฺโจ รญฺโญ มาคธสฺส อชาตสตฺตุสฺส เวเทหิปุตฺตสฺส อวิทูเร ตุณฺหีภูโต นิสินฺโน โหติ.\csromanspacer{} อถ โข ราชา มาคโธ อชาตสตฺตุ เวเทหิปุตฺโต ชีวกํ โกมารภจฺจํ เอตทโวจ “ตฺวํ ปน สมฺม ชีวก กิํ ตุณฺหี”ติ.\csromanspacer{} อยํ เทว ภควา อรหํ สมฺมาสมฺพุทฺโธ อมฺหากํ อมฺพวเน วิหรติ มหตา ภิกฺขุสํเฆน สทฺธิํ อฑฺฒเตฬเสหิ ภิกฺขุสเตหิ, ตํ โข ปน ภควนฺตํ\footnote{ภควนฺตํ โคตมํ (สี, ก, อิ)} เอวํ กลฺยาโณ กิตฺติสทฺโท อพฺภุคฺคโต “อิติปิ โส ภควา อรหํ สมฺมาสมฺพุทฺโธ วิชฺชาจรณสมฺปนฺโน สุคโต โลกวิทู อนุตฺตโร ปุริสทมฺมสารถิ สตฺถา เทวมนุสฺสานํ พุทฺโธ ภควา”ติ.\csromanspacer{} ตํ เทโว ภควนฺตํ ปยิรุปาสตุ, อปฺเปว นาม เทวสฺส ภควนฺตํ ปยิรุปาสโต จิตฺตํ ปสีเทยฺยาติ.}

\proseitem{158}{เตน หิ สมฺม ชีวก หตฺถิยานานิ กปฺปาเปหีติ. “เอวํ เทวา”ติ โข ชีวโก โกมารภจฺโจ รญฺโญ มาคธสฺส อชาตสตฺตุสฺส เวเทหิปุตฺตสฺส ปฏิสฺสุณิตฺวา ปญฺจมตฺตานิ หตฺถินิกาสตานิ กปฺปาเปตฺวา รญฺโญ จ อาโรหณียํ นาคํ, รญฺโญ มาคธสฺส อชาตสตฺตุสฺส เวเทหิปุตฺตสฺส ปฏิเวเทสิ “กปฺปิตานิ โข เต เทว หตฺถิยานานิ, ยสฺสทานิ กาลํ มญฺญสี”ติ.}

\proseitem{159}{อถ โข ราชา มาคโธ อชาตสตฺตุ เวเทหิปุตฺโต ปญฺจสุ หตฺถินิกาสเตสุ ปจฺเจกา อิตฺถิโย อาโรเปตฺวา อาโรหณียํ นาคํ อภิรุหิตฺวา อุกฺกาสุ ธาริยมานาสุ ราชคหมฺหา นิยฺยาสิ มหจฺจราชานุภาเวน, เยน ชีวกสฺส โกมารภจฺจสฺส อมฺพวนํ เตน ปายาสิ.}

\prose{อถ โข รญฺโญ มาคธสฺส อชาตสตฺตุสฺส เวเทหิปุตฺตสฺส อวิทูเร อมฺพวนสฺส อหุเทว ภยํ, อหุ ฉมฺภิตตฺตํ, อหุ โลมหํโส.\csromanspacer{} อถ โข ราชา มาคโธ อชาตสตฺตุ เวเทหิปุตฺโต ภีโต สํวิคฺโค}

\vspace{\tipitakaparskip}
\csromancenter{สามญฺญผลสุตฺต โลมหฏฺฏชาโต ชีวกํ โกมารภจฺจํ เอตทโวจ “กจฺจิ มํ สมฺม ชีวก น วญฺเจสิ, กจฺจิ มํ สมฺม ชีวก น ปลมฺเภสิ, กจฺจิ มํ สมฺม ชีวก น ปจฺจตฺถิกานํ เทสิ.\csromanspacer{} กถํ หิ นาม ตาว มหโต ภิกฺขุสํฆสฺส อฑฺฒเตฬสานํ ภิกฺขุสตานํ เนว ขิปิตสทฺโท ภวิสฺสติ, น อุกฺกาสิตสทฺโท, น นิคฺโฆโส”ติ.}
\prose{\csromanfolio{47}มา ภายิ มหาราช, มา ภายิ มหาราช, น ตํ เทว วญฺเจมิ, น ตํ เทว ปลมฺภามิ, น ตํ เทว ปจฺจตฺถิกานํ เทมิ, อภิกฺกม มหาราช, อภิกฺกม มหาราช, เอเต มณฺฑลมาเฬ ทีปา\footnote{ปทีปา (สี, สฺยา)} ฌายนฺตีติ.}

\csromanmatikaanchor{mtk.26}
\csromantocmarkref{section}{สามญฺญผลปุจฺฉา}{mtk.26}
\bookmark[dest={mtk.26},level=1]{สามญฺญผลปุจฺฉา}
\csromanheader{1}{\textbf{สามญฺญผลปุจฺฉา}}
\proseitem{160}{อถ โข ราชา มาคโธ อชาตสตฺตุ เวเทหิปุตฺโต ยาวติกา นาคสฺส ภูมิ นาเคน คนฺตฺวา นาคา ปจฺโจโรหิตฺวา ปตฺติโกว\footnote{ปทิโกว (สฺยา)} เยน มณฺฑลมาลสฺส ทฺวารํ เตนุปสงฺกมิ, อุปสงฺกมิตฺวา ชีวกํ โกมารภจฺจํ เอตทโวจ “กหํ ปน สมฺม ชีวก ภควา”ติ.\csromanspacer{} เอโส มหาราช ภควา, เอโส มหาราช ภควา มชฺฌิมํ ถมฺภํ นิสฺสาย ปุรตฺถาภิมุโข นิสินฺโน ปุรกฺขโต ภิกฺขุสํฆสฺสาติ.}

\proseitem{161}{อถ โข ราชา มาคโธ อชาตสตฺตุ เวเทหิปุตฺโต เยน ภควา เตนุปสงฺกมิ, อุปสงฺกมิตฺวา เอกมนฺตํ อฏฺฐาสิ.\csromanspacer{} เอกมนฺตํ ฐิโต โข ราชา มาคโธ อชาตสตฺตุ เวเทหิปุตฺโต ตุณฺหีภูตํ ภิกฺขุสํฆํ อนุวิโลเกตฺวา รหทมิว วิปฺปสนฺนํ อุทานํ อุทาเนสิ “อิมินา เม อุปสเมน อุทยภทฺโท\footnote{อุทายิภทฺโท (สี, อิ)} กุมาโร สมนฺนาคโต โหตุ, เยเนตรหิ อุปสเมน ภิกฺขุสํโฆ สมนฺนาคโต”ติ.\csromanspacer{} อคมา โข ตฺวํ มหาราช ยถาเปมนฺติ.\csromanspacer{} ปิโย เม ภนฺเต อุทยภทฺโท กุมาโร, อิมินา เม ภนฺเต อุปสเมน อุทยภทฺโท กุมาโร สมนฺนาคโต โหตุ, เยเนตรหิ อุปสเมน ภิกฺขุสํโฆ สมนฺนาคโตติ.}

\proseitem{162}{อถ โข ราชา มาคโธ อชาตสตฺตุ เวเทหิปุตฺโต ภควนฺตํ อภิวาเทตฺวา ภิกฺขุสํฆสฺส อญฺชลิํ ปณาเมตฺวา เอกมนฺตํ นิสีทิ. \csromanfolio{48}เอกมนฺตํ นิสินฺโน โข ราชา มาคโธ อชาตสตฺตุ เวเทหิปุตฺโต ภควนฺตํ เอตทโวจ “ปุจฺเฉยฺยามหํ ภนฺเต ภควนฺตํ กญฺจิเทว เทสํ\footnote{กิญฺจิเทว เทสํ เลสมตฺตํ (สฺยา, กํ, ก)} สเจ เม ภควา โอกาสํ กโรติ ปญฺหสฺส เวยฺยากรณายา”ติ.\csromanspacer{} ปุจฺฉ มหาราช, ยทากงฺขสีติ.}

\proseitem{163}{ยถา นุ โข อิมานิ ภนฺเต ปุถุสิปฺปายตนานิ, เสยฺยถิทํ, หตฺถาโรหา อสฺสาโรหา รถิกา ธนุคฺคหา เจลกา จลกา ปิณฺฑทายกา อุคฺคา ราชปุตฺตา ปกฺขนฺทิโน มหานาคา สูรา จมฺมโยธิโน ทาสิกปุตฺตา อาฬาริกา กปฺปกา นฺหาปกา\footnote{นหาปิกา (สี), นฺหาปิกา (สฺยา)} สูทา มาลการา รชกา เปสการา นฬการา กุมฺภการา คณกา มุทฺทิกา, ยานิ วา ปนญฺญานิปิ เอวํคตานิ ปุถุสิปฺปายตนานิ, เต ทิฏฺเฐว ธมฺเม สนฺทิฏฺฐิกํ สิปฺปผลํ อุปชีวนฺติ, เต เตน อตฺตานํ สุเขนฺติ ปีเณนฺติ\footnote{ปีเนนฺติ (กตฺถจิ)}, มาตาปิตโร สุเขนฺติ ปีเณนฺติ, ปุตฺตทารํ สุเขนฺติ ปีเณนฺติ, มิตฺตามจฺเจ สุเขนฺติ ปีเณนฺติ, สมณพฺราหฺมเณสุ\footnote{สมเณสุ พฺราหฺมเณสุ (ก)} อุทฺธคฺคิกํ ทกฺขิณํ ปติฏฺฐเปนฺติ โสวคฺคิกํ สุขวิปากํ สคฺคสํวตฺตนิกํ.\csromanspacer{} สกฺกา นุ โข ภนฺเต เอวเมว ทิฏฺเฐว ธมฺเม สนฺทิฏฺฐิกํ สามญฺญผลํ ปญฺญเปตุนฺติ.}

\proseitem{164}{อภิชานาติ โน ตฺวํ มหาราช อิมํ ปญฺหํ อญฺเญ สมณพฺราหฺมเณ ปุจฺฉิตาติ.\csromanspacer{} อภิชานามหํ ภนฺเต อิมํ ปญฺหํ อญฺเญ สมณพฺราหฺมเณ ปุจฺฉิตาติ.\csromanspacer{} ยถา กถํ ปน เต มหาราช พฺยากริํสุ, สเจ เต อครุ, ภาสสฺสูติ.\csromanspacer{} น โข เม ภนฺเต ครุ, ยตฺถสฺส ภควา นิสินฺโน ภควนฺตรูโป วาติ\footnote{จาติ (สี, ก)}.\csromanspacer{} เตน หิ มหาราช ภาสสฺสูติ.}

\csromanmatikaanchor{mtk.27}
\csromantocmarkref{section}{ปูรณกสฺสปวาท}{mtk.27}
\bookmark[dest={mtk.27},level=1]{ปูรณกสฺสปวาท}
\csromanheader{1}{\textbf{ปูรณกสฺสปวาท}}
\proseitem{165}{เอกมิทาหํ ภนฺเต สมยํ เยน ปูรโณ กสฺสโป เตนุปสงฺกมิํ, อุปสงฺกมิตฺวา ปูรเณน กสฺสเปน สทฺธิํ สมฺโมทิํ, สมฺโมทนียํ กถํ สารณียํ วีติสาเรตฺวา เอกมนฺตํ นิสีทิํ.\csromanspacer{} เอกมนฺตํ นิสินฺโน โข อหํ ภนฺเต ปูรณํ กสฺสปํ เอตทโวจํ “ยถา นุ โข อิมานิ โภ กสฺสป ปุถุสิปฺปายตนานิ, เสยฺยถิทํ หตฺถาโรหา อสฺสาโรหา รถิกา}

\vspace{\tipitakaparskip}
\csromancenter{สามญฺญผลสุตฺต ธนุคฺคหา เจลกา จลกา ปิณฺฑทายกา อุคฺคา ราชปุตฺตา ปกฺขนฺทิโน มหานาคา สูรา จมฺมโยธิโน ทาสิกปุตฺตา อาฬาริกา กปฺปกา นฺหาปกา สูทา มาลการา รชกา เปสการา นฬการา กุมฺภการา คณกา มุทฺทิกา, ยานิ วา ปนญฺญานิปิ เอวํคตานิ ปุถุสิปฺปายตนานิ, เต ทิฏฺเฐว ธมฺเม สนฺทิฏฺฐิกํ สิปฺปผลํ อุปชีวนฺติ, เต เตน อตฺตานํ สุเขนฺติ ปีเณนฺติ, มาตาปิตโร สุเขนฺติ ปีเณนฺติ, ปุตฺตทารํ สุเขนฺติ ปีเณนฺติ, มิตฺตามจฺเจ สุเขนฺติ ปีเณนฺติ, สมณพฺราหฺมเณสุ อุทฺธคฺคิกํ ทกฺขิณํ ปติฏฺฐเปนฺติ โสวคฺคิกํ สุขวิปากํ สคฺคสํวตฺตนิกํ.\csromanspacer{} สกฺกา นุ โข โภ กสฺสป เอวเมว ทิฏฺเฐว ธมฺเม สนฺทิฏฺฐิกํ สามญฺญผลํ ปญฺญเปตุนฺ”ติ.}
\proseitem{166}{\csromanfolio{49}เอวํ วุตฺเต ภนฺเต ปูรโณ กสฺสโป มํ เอตทโวจ “กโรโต โข มหาราช การยโต, ฉินฺทโต เฉทาปยโต, ปจโต ปาจาปยโต, โสจยโต โสจาปยโต, กิลมโต กิลมาปยโต, ผนฺทโต ผนฺทาปยโต, ปาณมติปาตาปยโต, อทินฺนํ อาทิยโต, สนฺธิํ ฉินฺทโต, นิลฺโลปํ หรโต, เอกาคาริกํ กโรโต, ปริปนฺเถ ติฏฺฐโต, ปรทารํ คจฺฉโต, มุสา ภณโต, กโรโต น กรียติ ปาปํ.\csromanspacer{} ขุรปริยนฺเตน เจปิ จกฺเกน โย อิมิสฺสา ปถวิยา ปาเณ เอกํ มํสขลํ เอกํ มํสปุญฺชํ กเรยฺย, นตฺถิ ตโตนิทานํ ปาปํ, นตฺถิ ปาปสฺส อาคโม.\csromanspacer{} ทกฺขิณญฺเจปิ คงฺคาย ตีรํ คจฺเฉยฺย หนนฺโต ฆาเตนฺโต ฉินฺทนฺโต เฉทาเปนฺโต ปจนฺโต ปาจาเปนฺโต, นตฺถิ ตโตนิทานํ ปาปํ, นตฺถิ ปาปสฺส อาคโม.\csromanspacer{} อุตฺตรญฺเจปิ คงฺคาย ตีรํ คจฺเฉยฺย ททนฺโต ทาเปนฺโต ยชนฺโต ยชาเปนฺโต, นตฺถิ ตโตนิทานํ ปุญฺญํ, นตฺถิ ปุญฺญสฺส อาคโม.\csromanspacer{} ทาเนน ทเมน สํยเมน สจฺจวชฺเชน นตฺถิ ปุญฺญํ, นตฺถิ ปุญฺญสฺส อาคโม”ติ.\csromanspacer{} อิตฺถํ โข เม ภนฺเต ปูรโณ กสฺสโป สนฺทิฏฺฐิกํ สามญฺญผลํ ปุฏฺโฐ สมาโน อกิริยํ พฺยากาสิ.}

\prose{เสยฺยถาปิ ภนฺเต อมฺพํ วา ปุฏฺโฐ ลพุชํ พฺยากเรยฺย, ลพุชํ วา ปุฏฺโฐ อมฺพํ พฺยากเรยฺย.\csromanspacer{} เอวเมว โข เม ภนฺเต ปูรโณ กสฺสโป สนฺทิฏฺฐิกํ สามญฺญผลํ ปุฏฺโฐ สมาโน อกิริยํ พฺยากาสิ, ตสฺส มยฺหํ ภนฺเต เอตทโหสิ “กถํ หิ นาม มาทิโส สมณํ วา พฺราหฺมณํ วา วิชิเต วสนฺตํ อปสาเทตพฺพํ มญฺเญยฺยา”ติ.\csromanspacer{} โส โข อหํ ภนฺเต ปูรณสฺส กสฺสปสฺส ภาสิตํ เนว อภินนฺทิํ นปฺปฏิกฺโกสิํ, อนภินนฺทิตฺวา \csromanfolio{50}อปฏิกฺโกสิตฺวา อนตฺตมโน อนตฺตมนวาจํ อนิจฺฉาเรตฺวา ตเมว วาจํ อนุคฺคณฺหนฺโต อนิกฺกุชฺชนฺโต\footnote{อนิกฺกุชฺเชนฺโต (สฺยา, กํ, ก)} อุฏฺฐายาสนา ปกฺกมิํ\footnote{ปกฺกามิํ (สี, สฺยา, กํ, อิ)}.}

\csromanmatikaanchor{mtk.28}
\csromantocmarkref{section}{มกฺขลิโคสาลวาท}{mtk.28}
\bookmark[dest={mtk.28},level=1]{มกฺขลิโคสาลวาท}
\csromanheader{1}{\textbf{มกฺขลิโคสาลวาท}}
\proseitem{167}{เอกมิทาหํ ภนฺเต สมยํ เยน มกฺขลิ โคสาโล เตนุปสงฺกมิํ, อุปสงฺกมิตฺวา มกฺขลินา โคสาเลน สทฺธิํ สมฺโมทิํ, สมฺโมทนียํ กถํ สารณียํ วีติสาเรตฺวา เอกมนฺตํ นิสีทิํ.\csromanspacer{} เอกมนฺตํ นิสินฺโน โข อหํ ภนฺเต มกฺขลิํ โคสาลํ เอตทโวจํ “ยตฺถา นุ โข อิมานิ โภ โคสาล ปุถุสิปฺปายตนานิ ฯเปฯ.\csromanspacer{} สกฺกา นุ โข โภ โคสาล เอวเมว ทิฏฺเฐว ธมฺเม สนฺทิฏฺฐิกํ สามญฺญผลํ ปญฺญเปตุนฺ”ติ.}

\proseitem{168}{เอวํ วุตฺเต ภนฺเต มกฺขลิ โคสาโล มํ เอตทโวจ “นตฺถิ มหาราช เหตุ นตฺถิ ปจฺจโย สตฺตานํ สํกิเลสาย, อเหตู\footnote{อเหตุ (กตฺถจิ)} อปจฺจยา สตฺตา สํกิลิสฺสนฺติ.\csromanspacer{} นตฺถิ เหตุ นตฺถิ ปจฺจโย สตฺตานํ วิสุทฺธิยา, อเหตู อปจฺจยา สตฺตา วิสุชฺฌนฺติ.\csromanspacer{} นตฺถิ อตฺตกาเร, นตฺถิ ปรกาเร, นตฺถิ ปุริสกาเร, นตฺถิ พลํ, นตฺถิ วีริยํ, นตฺถิ ปุริสถาโม, นตฺถิ ปุริสปรกฺกโม, สพฺเพ สตฺตา สพฺเพ ปาณา สพฺเพ ภูตา สพฺเพ ชีวา อวสา อพลา อวีริยา นิยติสํคติภาวปริณตา ฉเสฺววาภิชาตีสุ สุขทุกฺขํ\footnote{สุขญฺจ ทุกฺขญฺจ (สฺยา)} ปฏิสํเวเทนฺติ, จุทฺทส โข ปนิมานิ โยนิปมุขสตสหสฺสานิ สฏฺฐิ จ สตานิ ฉ จ สตานิ ปญฺจ จ กมฺมุโน สตานิ ปญฺจ จ กมฺมานิ ตีณิ จ กมฺมานิ กมฺเม จ อฑฺฒกมฺเม จ ทฺวฏฺฐิปฏิปทา ทฺวฏฺฐนฺตรกปฺปา ฉฬาภิชาติโย อฏฺฐ ปุริสภูมิโย เอกูนปญฺญาส อาชีวกสเต เอกูนปญฺญาส ปริพฺพาชกสเต เอกูนปญฺญาส นาคาวาสสเต วีเส อินฺทฺริยสเต ติํเส นิรยสเต ฉตฺติํส รโชธาตุโย สตฺต สญฺญีคพฺภา สตฺต อสญฺญีคพฺภา สตฺต นิคณฺฐิคพฺภา สตฺต เทวา สตฺต มานุสา สตฺต ปิสาจา สตฺต สรา สตฺต ปวุฏา\footnote{สปุฏา (ก), ปพุฏา (สี)} สตฺต ปวุฏสตานิ สตฺต ปปาตา สตฺต ปปาตสตานิ สตฺต สุปินา สตฺต สุปินสตานิ จุลฺลาสีติ มหากปฺปิโน\footnote{มหากปฺปุโน (กสี, อิ)} สตสหสฺสานิ, ยานิ พาเล จ ปณฺฑิเต จ สนฺธาวิตฺวา สํสริตฺวา ทุกฺขสฺสนฺตํ กริสฺสนฺติ, ตตฺถ}

\vspace{\tipitakaparskip}
\csromancenter{สามญฺญผลสุตฺต นตฺถิ ‘อิมินาหํ สีเลน วา วเตน วา ตเปน วา พฺรหฺมจริเยน วา อปริปกฺกํ วา กมฺมํ ปริปาเจสฺสามิ, ปริปกฺกํ วา กมฺมํ ผุสฺส ผุสฺส พฺยนฺติํ กริสฺสามี’ติ เหวํ นตฺถิ.\csromanspacer{} โทณมิเต สุขทุกฺเข, ปริยนฺตกเต สํสาเร นตฺถิ หายนวฑฺฒเน, นตฺถิ อุกฺกํสาวกํเส.\csromanspacer{} เสยฺยถาปิ นาม สุตฺตคุเฬ ขิตฺเต นิพฺเพฐิยมานเมว ปเลติ, เอวเมว พาเล จ ปณฺฑิเต จ สนฺธาวิตฺวา สํสริตฺวา ทุกฺขสฺสนฺตํ กริสฺสนฺตี”ติ.}
\proseitem{169}{\csromanfolio{51}อิตฺถํ โข เม ภนฺเต มกฺขลิ โคสาโล สนฺทิฏฺฐิกํ สามญฺญผลํ ปุฏฺโฐ สมาโน สํสารสุทฺธิํ พฺยากาสิ.\csromanspacer{} เสยฺยถาปิ ภนฺเต อมฺพํ วา ปุฏฺโฐ ลพุชํ พฺยกเรยฺย, ลพุชํ วา ปุฏฺโฐ อมฺพํ พฺยากเรยฺย.\csromanspacer{} เอวเมว โข เม ภนฺเต มกฺขลิ โคสาโล สนฺทิฏฺฐิกํ สามญฺญผลํ ปุฏฺโฐ สมาโน สํสารสุทฺธิํ พฺยากาสิ.\csromanspacer{} ตสฺส มยฺหํ ภนฺเต เอตทโหสิ “กถํ หิ นาม มาทิโส สมณํ วา พฺราหฺมณํ วา วิชิเต วสนฺตํ อปสาเทตพฺพํ มญฺเญยฺยา”ติ.\csromanspacer{} โส โข อหํ ภนฺเต มกฺขลิสฺส โคสาลสฺส ภาสิตํ เนว อภินนฺทิํ นปฺปฏิกฺโกสิํ, อนภินนฺทิตฺวา อปฺปฏิกฺโกสิตฺวา อนตฺตมโน อนตฺตมนวาจํ อนิจฺฉาเรตฺวา ตเมว วาจํ อนุคฺคณฺหนฺโต อนิกฺกุชฺชนฺโต อุฏฺฐายาสนา ปกฺกมิํ.}

\csromanmatikaanchor{mtk.29}
\csromantocmarkref{section}{อชิตเกสกมฺพลวาท}{mtk.29}
\bookmark[dest={mtk.29},level=1]{อชิตเกสกมฺพลวาท}
\csromanheader{1}{\textbf{อชิตเกสกมฺพลวาท}}
\proseitem{170}{เอกมิทาหํ ภนฺเต สมยํ เยน อชิโต เกสกมฺพโล เตนุปสงฺกมิํ, อุปสงฺกมิตฺวา อชิเตน เกสกมฺพเลน สทฺธิํ สมฺโมทิํ, สมฺโมทนียํ กถํ สารณียํ วีติสาเรตฺวา เอกมนฺตํ นิสีทิํ, เอกมนฺตํ นิสินฺโน โข อหํ ภนฺเต อชิตํ เกสกมฺพลํ เอตทโวจํ “ยถา นุ โข อิมานิ โภ อชิต ปุถุสิปฺปายตนานิ ฯเปฯ.\csromanspacer{} สกฺกา นุ โข โภ อชิต เอวเมว ทิฏฺเฐว ธมฺเม สนฺทิฏฺฐิกํ สามญฺญผลํ ปญฺญเปตุนฺ’ติ.}

\proseitem{171}{เอวํ วุตฺเต ภนฺเต อชิโต เกสกมฺพโล มํ เอตทโวจ “นตฺถิ มหาราช ทินฺนํ, นตฺถิ ยิฏฺฐํ, นตฺถิ หุตํ, นตฺถิ สุกตทุกฺกฏานํ กมฺมานํ ผลํ วิปาโก, นตฺถิ อยํ โลโก, นตฺถิ ปโร โลโก\footnote{ปรโลโก (สฺยา)}, นตฺถิ มาตา, นตฺถิ ปิตา, นตฺถิ สตฺตา โอปปาติกา, นตฺถิ โลเก สมณพฺราหฺมณา สมฺมคฺคตา\footnote{สมคฺคตา (ก), สมคฺคคตา (สฺยา)} สมฺมาปฏิปนฺนา, เย อิมญฺจ โลกํ ปรญฺจ โลกํ \csromanfolio{52}สยํ อภิญฺญา สจฺฉิกตฺวา ปเวเทนฺติ.\csromanspacer{} จาตุมหาภูติโก อยํ ปุริโส, ยทา กาลํ กโรติ, ปถวี ปถวิกายํ อนุเปติ อนุปคจฺฉติ, อาโป อาโปกายํ อนุเปติ อนุปคจฺฉติ, เตโช เตโชกายํ อนุเปติ อนุปคจฺฉติ, วาโย วาโยกายํ อนุเปติ อนุปคจฺฉติ, อากาสํ อินฺทฺริยานิ สงฺกมนฺติ, อาสนฺทิปญฺจมา ปุริสา มตํ อาทาย คจฺฉนฺติ, ยาวาฬาหนา ปทานิ ปญฺญายนฺติ, กาโปตกานิ อฏฺฐีนิ ภวนฺติ, ภสฺสนฺตา อาหุติโย, ทตฺตุปญฺญตฺตํ ยทิทํ ทานํ, เตสํ ตุจฺฉํ มุสา วิลาโป, เย เกจิ อตฺถิกวาทํ วทนฺติ, พาเล จ ปณฺฑิเต จ กายสฺส เภทา อุจฺฉิชฺชนฺติ วินสฺสนฺติ, น โหนฺติ ปรํ มรณา”ติ.}

\proseitem{172}{อิตฺถํ โข เม ภนฺเต อชิโต เกสกมฺพโล สนฺทิฏฺฐิกํ สามญฺญผลํ ปุฏฺโฐ สมาโน อุจฺเฉทํ พฺยากาสิ.\csromanspacer{} เสยฺยถาปิ ภนฺเต อมฺพํ วา ปุฏฺโฐ ลพุชํ พฺยากเรยฺย, ลพุชํ วา ปุฏฺโฐ อมฺพํ พฺยากเรยฺย.\csromanspacer{} เอวเมว โข เม ภนฺเต อชิโต เกสกมฺพโล สนฺทิฏฺฐิกํ สามญฺญผลํ ปุฏฺโฐ สมาโน อุจฺเฉทํ พฺยากาสิ.\csromanspacer{} ตสฺส มยฺหํ ภนฺเต เอตทโหสิ “กถํ หิ นาม มาทิโส สมณํ วา พฺราหฺมณํ วา วิชิเต วสนฺตํ อปสาเทตพฺพํ มญฺเญยฺยา”ติ.\csromanspacer{} โส โข อหํ ภนฺเต อชิตสฺส เกสกมฺพลสฺส ภาสิตํ เนว อภินนฺทิํ นปฺปฏิกฺโกสิํ, อนภินนฺทิตฺวา อปฺปฏิกฺโกสิตฺวา อนตฺตมโน อนตฺตมนวาจํ อนิจฺฉาเรตฺวา ตเมว วาจํ อนุคฺคณฺหนฺโต อนิกฺกุชฺชนฺโต อุฏฺฐายาสนา ปกฺกมิํ.}

\csromanmatikaanchor{mtk.30}
\csromantocmarkref{section}{ปกุธกจฺจายนวาท}{mtk.30}
\bookmark[dest={mtk.30},level=1]{ปกุธกจฺจายนวาท}
\csromanheader{1}{\textbf{ปกุธกจฺจายนวาท}}
\proseitem{173}{เอกมิทาหํ ภนฺเต สมยํ เยน ปกุโธ กจฺจายโน เตนุปสงฺกมิํ, อุปสงฺกมิตฺวา ปกุเธน กจฺจายเนน สทฺธิํ สมฺโมทิํ, สมฺโมทนียํ กถํ สารณียํ วีติสาเรตฺวา เอกมนฺตํ นิสีทิํ, เอกมนฺตํ นิสินฺโน โข อหํ ภนฺเต ปกุธํ กจฺจายนํ เอตทโวจํ “ยถา นุ โข อิมานิ โภ กจฺจายน ปุถุสิปฺปายตนานิ ฯเปฯ.\csromanspacer{} สกฺกา นุ โข โภ กจฺจายน เอวเมว ทิฏฺเฐว ธมฺเม สนฺทิฏฺฐิกํ สามญฺญผลํ ปญฺญเปตุนฺ”ติ.}

\proseitem{174}{เอวํ วุตฺเต ภนฺเต ปกุโธ กจฺจายโน มํ เอตทโวจ “สตฺติเม มหาราช กายา อกฏา อกฏวิธา อนิมฺมิตา อนิมฺมาตา วญฺฌา}

\vspace{\tipitakaparskip}
\csromancenter{สามญฺญผลสุตฺต กูฏฏฺฐา เอสิกฏฺฐายิฏฺฐิตา, เต น อิญฺชนฺติ, น วิปริณาเมนฺติ, น อญฺญมญฺญํ พฺยาพาเธนฺติ, นาลํ อญฺญมญฺญสฺส สุขาย วา ทุกฺขาย วา สุขทุกฺขาย วา.\csromanspacer{} กตเม สตฺต, ปถวิกาโย, อาโปกาโย, เตโชกาโย, วาโยกาโย, สุเข, ทุกฺเข, ชีเว สตฺตเม.\csromanspacer{} อิเม สตฺต กายา อกฏา อกฏวิธา อนิมฺมิตา อนิมฺมาตา วญฺฌา กูฏฏฺฐา เอสิกฏฺฐายิฏฺฐิตา, เต น อิญฺชนฺติ, น วิปริณาเมนฺติ, น อญฺญมญฺญํ พฺยาพาเธนฺติ, นาลํ อญฺญมญฺญสฺส สุขาย วา ทุกฺขาย วา สุขทุกฺขาย วา.\csromanspacer{} ตตฺถ นตฺถิ หนฺตา วา ฆาเตตา วา, โสตา วา สาเวตา วา, วิญฺญาตา วา วิญฺญาเปตา วา, โยปิ ติเณฺหน สตฺเถน สีสํ ฉินฺทติ, น โกจิ กิญฺจิ\footnote{กญฺจิ (กํ)} ชีวิตา โวโรเปติ, สตฺตนฺนํ เตฺวว\footnote{สตฺตนฺนํ เยว (สี, สฺยา, กํ, อิ)} กายานมนฺตเรน สตฺถํ วิวรมนุปตตี”ติ.}
\proseitem{175}{\csromanfolio{53}อิตฺถํ โข เม ภนฺเต ปกุโธ กจฺจายโน สนฺทิฏฺฐิกํ สามญฺญผลํ ปุฏฺโฐ สมาโน อญฺเญน อญฺญํ พฺยากาสิ.\csromanspacer{} เสยฺยถาปิ ภนฺเต อมฺพํ วา ปุฏฺโฐ ลพุชํ พฺยากเรยฺย, ลพุชํ วา ปุฏฺโฐ อมฺพํ พฺยากเรยฺย.\csromanspacer{} เอวเมว โข เม ภนฺเต ปกุโธ กจฺจายโน สนฺทิฏฺฐิกํ สามญฺญผลํ ปุฏฺโฐ สมาโน อญฺเญน อญฺญํ พฺยากาสิ.\csromanspacer{} ตสฺส มยฺหํ ภนฺเต เอตทโหสิ “กถํ หิ นาม มาทิโส สมณํ วา พฺราหฺมณํ วา วิชิเต วสนฺตํ อปสาเทตพฺพํ มญฺเญยฺยา”ติ.\csromanspacer{} โส โข อหํ ภนฺเต ปกุธสฺส กจฺจายนสฺส ภาสิตํ เนว อภินนฺทิํ นปฺปฏิกฺโกสิํ, อนภินนฺทิตฺวา อปฺปฏิกฺโกสิตฺวา อนตฺตมโน อนตฺตมนวาจํ อนิจฺฉาเรตฺวา ตเมว วาจํ อนุคฺคณฺหนฺโต อนิกฺกุชฺชนฺโต อุฏฺฐายาสนา ปกฺกมิํ.}

\csromanmatikaanchor{mtk.31}
\csromantocmarkref{section}{นิคณฺฐนาฏปุตฺตวาท}{mtk.31}
\bookmark[dest={mtk.31},level=1]{นิคณฺฐนาฏปุตฺตวาท}
\csromanheader{1}{\textbf{นิคณฺฐนาฏปุตฺตวาท}}
\proseitem{176}{เอกมิทาหํ ภนฺเต สมยํ เยน นิคณฺโฐ นาฏปุตฺโต เตนุปสงฺกมิํ, อุปสงฺกมิตฺวา นิคณฺเฐน นาฏปุตฺเตน สทฺธิํ สมฺโมทิํ, สมฺโมทนียํ กถํ สารณียํ วีติสาเรตฺวา เอกมนฺตํ นิสีทิํ, เอกมนฺตํ นิสินฺโน โข อหํ ภนฺเต นิคณฺฐํ นาฏปุตฺตํ เอตทโวจํ “ยถา นุ โข อิมานิ โภ อคฺคิเวสฺสน ปุถุสิปฺปายตนานิ ฯเปฯ.\csromanspacer{} สกฺกา นุ โข โภ อคฺคิเวสฺสน เอวเมว ทิฏฺเฐว ธมฺเม สนฺทิฏฺฐิกํ สามญฺญผลํ ปญฺญเปตุนฺ”ติ.}

\proseitem{177}{\csromanfolio{54}เอวํ วุตฺเต ภนฺเต นิคณฺโฐ นาฏปุตฺโต มํ เอตทโวจ “อิธ มหาราช นิคณฺโฐ จาตุยามสํวรสํวุโต โหติ.\csromanspacer{} กถญฺจ มหาราช นิคณฺโฐ จาตุยามสํวรสํวุโต โหติ.\csromanspacer{} อิธ มหาราช นิคณฺโฐ สพฺพวาริวาริโต จ โหติ, สพฺพวาริยุตฺโต จ, สพฺพวาริธุโต จ, สพฺพวาริผุโฏ จ.\csromanspacer{} เอวํ โข มหาราช นิคณฺโฐ จาตุยามสํวรสํวุโต โหติ, ยโต โข มหาราช นิคณฺโฐ เอวํ จาตุยามสํวรสํวุโต โหติ.\csromanspacer{} อยํ วุจฺจติ มหาราช นิคณฺโฐ\footnote{นิคณฺโฐ นาฏปุตฺโต (สฺยา, ก)} คตตฺโต จ ยตตฺโต จ ฐิตตฺโต จา”ติ.}

\proseitem{178}{อิตฺถํ โข เม ภนฺเต นิคณฺโฐ นาฏปุตฺโต สนฺทิฏฺฐิกํ สามญฺญผลํ ปุฏฺโฐ สมาโน จาตุยามสํวรํ พฺยากาสิ.\csromanspacer{} เสยฺยถาปิ ภนฺเต อมฺพํ วา ปุฏฺโฐ ลพุชํ พฺยากเรยฺย, ลพุชํ วา ปุฏฺโฐ อมฺพํ พฺยากเรยฺย.\csromanspacer{} เอวเมว โข เม ภนฺเต นิคณฺโฐ นาฏปุตฺโต สนฺทิฏฺฐิกํ สามญฺญผลํ ปุฏฺโฐ สมาโน จาตุยามสํวรํ พฺยากาสิ.\csromanspacer{} ตสฺส มยฺหํ ภนฺเต เอตทโหสิ “กถํ หิ นาม มาทิโส สมณํ วา พฺราหฺมณํ วา วิชิเต วสนฺตํ อปสาเทตพฺพํ มญฺเญยฺยา”ติ.\csromanspacer{} โส โข อหํ ภนฺเต นิคณฺฐสฺส นาฏปุตฺตสฺส ภาสิตํ เนว อภินนฺทิํ นปฺปฏิกฺโกสิํ, อนภินนฺทิตฺวา อปฺปฏิกฺโกสิตฺวา อนตฺตมโน อนตฺตมนวาจํ อนิจฺฉาเรตฺวา ตเมว วาจํ อนุคฺคณฺหนฺโต อนิกฺกุชฺชนฺโต อุฏฺฐายาสนา ปกฺกมิํ.}

\csromanmatikaanchor{mtk.32}
\csromantocmarkref{section}{สญฺจยเพลฏฺฐปุตฺตวาท}{mtk.32}
\bookmark[dest={mtk.32},level=1]{สญฺจยเพลฏฺฐปุตฺตวาท}
\csromanheader{1}{สญฺจย\-เพลฏฺฐปุตฺต\-วาท}
\proseitem{179}{เอกมิทาหํ ภนฺเต สมยํ เยน สญฺจโย เพลฏฺฐปุตฺโต เตนุปสงฺกมิํ, อุปสงฺกมิตฺวา สญฺจเยน เพลฏฺฐปุตฺเตน สทฺธิํ สมฺโมทิํ, สมฺโมทนียํ กถํ สารณียํ วีติสาเรตฺวา เอกมนฺตํ นิสีทิํ, เอกมนฺตํ นิสินฺโน โข อหํ ภนฺเต สญฺจยํ เพลฏฺฐปุตฺตํ เอตทโวจํ “ยถา นุ โข อิมานิ โภ สญฺจย ปุถุสิปฺปายตนานิ ฯเปฯ.\csromanspacer{} สกฺกา นุ โข โภ สญฺจย เอวเมว ทิฏฺเฐว ธมฺเม สนฺทิฏฺฐิกํ สามญฺญผลํ ปญฺญเปตุนฺ”ติ.}

\proseitem{180}{เอวํ วุตฺเต ภนฺเต สญฺจโย เพลฏฺฐปุตฺโต มํ เอตทโวจ “อตฺถิ ปโร โลโกติ อิติ เจ มํ ปุจฺฉสิ, อตฺถิ ปโร โลโกติ อิติ เจ เม อสฺส, อตฺถิ ปโร โลโกติ อิติ เต นํ พฺยากเรยฺยํ, เอวนฺติปิ เม โน, ตถาติปิ เม โน, อญฺญถาติปิ เม โน, โนติปิ}

\vspace{\tipitakaparskip}
\csromancenter{สามญฺญผลสุตฺต เม โน, โน โนติปิ เม โน.\csromanspacer{} นตฺถิ ปโร โลโก ฯเปฯ.\csromanspacer{} อตฺถิ จ นตฺถิ จ ปโร โลโก ฯเปฯ.\csromanspacer{} เนวตฺถิ น นตฺถิ ปโร โลโก ฯเปฯ.\csromanspacer{} อตฺถิ สตฺตา โอปปาติกา ฯเปฯ.\csromanspacer{} นตฺถิ สตฺตา โอปปาติกา ฯเปฯ.\csromanspacer{} อตฺถิ จ นตฺถิ จ สตฺตา โอปปาติกา ฯเปฯ.\csromanspacer{} เนวตฺถิ น นตฺถิ สตฺตา โอปปาติกา ฯเปฯ.\csromanspacer{} อตฺถิ สุกตทุกฺกฏานํ กมฺมานํ ผลํ วิปาโก ฯเปฯ.\csromanspacer{} นตฺถิ สุกตทุกฺกฏานํ กมฺมานํ ผลํ วิปาโก ฯเปฯ.\csromanspacer{} อตฺถิ จ นตฺถิ จ สุกตทุกฺกฏานํ กมฺมานํ ผลํ วิปาโก ฯเปฯ.\csromanspacer{} เนวตฺถิ น นตฺถิ สุกตทุกฺกฏานํ กมฺมานํ ผลํ วิปาโก ฯเปฯ.\csromanspacer{} โหติ ตถาคโต ปรํ มรณา ฯเปฯ.\csromanspacer{} น โหติ ตถาคโต ปรํ มรณา ฯเปฯ.\csromanspacer{} โหติ จ น จ โหติ ตถาคโต ปรํ มรณา ฯเปฯ.\csromanspacer{} เนว โหติ น น โหติ ตถาคโต ปรํ มรณาติ อิติ เจ มํ ปุจฺฉสิ, เนว โหติ น น โหติ ตถาคโต ปรํ มรณาติ อิติ เจ เม อสฺส, เนว โหติ น น โหติ ตถาคโต ปรํ มรณาติ อิติ เต นํ พฺยากเรยฺยํ, เอวนฺติปิ เม โน, ตถาติปิ เม โน, อญฺญถาติปิ เม โน, โนติปิ เม โน, โน โนติปิ เม โน”ติ.}
\proseitem{181}{\csromanfolio{55}อิตฺถํ โข เม ภนฺเต สญฺจโย เพลฏฺฐปุตฺโต สนฺทิฏฺฐิกํ สามญฺญผลํ ปุฏฺโฐ สมาโน วิกฺเขปํ พฺยากาสิ.\csromanspacer{} เสยฺยถาปิ ภนฺเต อมฺพํ วา ปุฏฺโฐ ลพุชํ พฺยากเรยฺย, ลพุชํ วา ปุฏฺโฐ อมฺพํ พฺยากเรยฺย.\csromanspacer{} เอวเมว โข เม ภนฺเต สญฺจโย เพลฏฺฐปุตฺโต สนฺทิฏฺฐิกํ สามญฺญผลํ ปุฏฺโฐ สมาโน วิกฺเขปํ พฺยากาสิ.\csromanspacer{} ตสฺส มยฺหํ ภนฺเต เอตทโหสิ “อยญฺจ อิเมสํ สมณพฺราหฺมณานํ สพฺพพาโล สพฺพมูโฬฺห, กถํ หิ นาม สนฺทิฏฺฐิกํ สามญฺญผลํ ปุฏฺโฐ สมาโน วิกฺเขปํ พฺยากริสฺสตี”ติ.\csromanspacer{} ตสฺส มยฺหํ ภนฺเต เอตทโหสิ “กถํ หิ นาม มาทิโส สมณํ วา พฺราหฺมณํ วา วิชิเต วสนฺตํ อปสาเทตพฺพํ มญฺเญยฺยา”ติ.\csromanspacer{} โส โข อหํ ภนฺเต สญฺจยสฺส เพลฏฺฐปุตฺตสฺส ภาสิตํ เนว อภินนฺทิํ นปฺปฏิกฺโกสิํ, อนภินนฺทิตฺวา อปฺปฏิกฺโกสิตฺวา อนตฺตมโน อนตฺตมนวาจํ อนิจฺฉาเรตฺวา ตเมว วาจํ อนุคฺคณฺหนฺโต อนิกฺกุชฺชนฺโต อุฏฺฐายาสนา ปกฺกมิํ.}

\csromanmatikaanchor{mtk.33}
\csromantocmarkref{section}{ปฐมสนฺทิฏฺฐิกสามญฺญผล}{mtk.33}
\bookmark[dest={mtk.33},level=1]{ปฐมสนฺทิฏฺฐิกสามญฺญผล}
\csromanheader{1}{ปฐม\-สนฺทิฏฺฐิก\-สามญฺญผล}
\proseitem{182}{โสหํ ภนฺเต ภควนฺตมฺปิ ปุจฺฉามิ “ยถา นุ โข อิมานิ ภนฺเต ปุถุสิปฺปายตนานิ.\csromanspacer{} เสยฺยถิทํ, หตฺถาโรหา อสฺสาโรหา รถิกา \csromanfolio{56}ธนุคฺคหา เจลกา จลกา ปิณฺฑทายกา อุคฺคา ราชปุตฺตา ปกฺขนฺทิโน มหานาคา สูรา จมฺมโยธิโน ทาสิกปุตฺตา อาฬาริกา กปฺปกา นฺหาปกา สูทา มาลการา รชกา เปสการา นฬการา กุมฺภการา คณกา มุทฺทิกา, ยานิ วา ปนญฺญานิปิ เอวํคตานิ ปุถุสิปฺปายตนานิ, เต ทิฏฺเฐว ธมฺเม สนฺทิฏฺฐิกํ สิปฺปผลํ อุปชีวนฺติ, เต เตน อตฺตานํ สุเขนฺติ ปีเณนฺติ, มาตาปิตโร สุเขนฺติ ปีเณนฺติ, ปุตฺตทารํ สุเขนฺติ ปีเณนฺติ, มิตฺตามจฺเจ สุเขนฺติ ปีเณนฺติ, สมณพฺราหฺมเณสุ อุทฺธคฺคิกํ ทกฺขิณํ ปติฏฺฐเปนฺติ โสวคฺคิกํ สุขวิปากํ สคฺคสํวตฺตนิกํ.\csromanspacer{} สกฺกา นุ โข ภนฺเต เอวเมว ทิฏฺเฐว ธมฺเม สนฺทิฏฺฐิกํ สามญฺญผลํ ปญฺญเปตุนฺ”ติ.}

\proseitem{183}{สกฺกา มหาราช, เตน หิ มหาราช ตญฺเญเวตฺถ ปฏิปุจฺฉิสฺสามิ, ยถา เต ขเมยฺย, ตถา ตํ พฺยากเรยฺยาสิ.\csromanspacer{} ตํ กิํ มญฺญสิ มหาราช.\csromanspacer{} อิธ เต อสฺส ปุริโส ทาโส กมฺมกาโร\footnote{กมฺมกโร (สี, สฺยา, กํ, อิ)} ปุพฺพุฏฺฐายี ปจฺฉานิปาตี กิํการปฏิสฺสาวี มนาปจารี ปิยวาที มุขุลฺโลกโก\footnote{มุขุลฺโลกิโก (สฺยา, กํ, ก)}.\csromanspacer{} ตสฺส เอวมสฺส “อจฺฉริยํ วต โภ อพฺภุตํ วต โภ ปุญฺญานํ คติ ปุญฺญานํ วิปาโก, อยํ หิ ราชา มาคโธ อชาตสตฺตุ เวเทหิปุตฺโต มนุสฺโส, อหมฺปิ มนุสฺโส, อยํ หิ ราชา มาคโธ อชาตสตฺตุ เวเทหิปุตฺโต ปญฺจหิ กามคุเณหิ สมปฺปิโต สมงฺคีภูโต ปริจาเรติ เทโว มญฺเญ, อหํ ปนมฺหิสฺส ทาโส กมฺมกาโร ปุพฺพุฏฺฐายี ปจฺฉานิปาตี กิํการปฏิสฺสาวี มนาปจารี ปิยวาที มุขุลฺโลกโก, โส วตสฺสาหํ ปุญฺญานิ กเรยฺยํ, ยํนูนาหํ เกสมสฺสุํ โอหาเรตฺวา กาสายานิ วตฺถานิ อจฺฉาเทตฺวา อคารสฺมา อนคาริยํ ปพฺพเชยฺยนฺ”ติ.\csromanspacer{} โส อปเรน สมเยน เกสมสฺสุํ โอหาเรตฺวา กาสายานิ วตฺถานิ อจฺฉาเทตฺวา อคารสฺมา อนคาริยํ ปพฺพเชยฺย, โส เอวํ ปพฺพชิโต สมาโน กาเยน สํวุโต วิหเรยฺย, วาจาย สํวุโต วิหเรยฺย, มนสา สํวุโต วิหเรยฺย, ฆาสจฺฉาทนปรมตาย สนฺตุฏฺโฐ, อภิรโต ปวิเวเก.\csromanspacer{} ตํ เจ เต ปุริสา เอวมาโรเจยฺยุํ “ยคฺเฆ เทว ชาเนยฺยาสิ, โย เต โส ปุริโส\footnote{โย เต ปุริโส (สี, ก)} ทาโส กมฺมกาโร ปุพฺพุฏฺฐายี ปจฺฉานิปาตี กิํ การปฏิสฺสาวี มนาปจารี ปิยวาที มุขุลฺโลกโก, โส เทว เกสมสฺสุํ โอหาเรตฺวา กาสายานิ วตฺถานิ อจฺฉาเทตฺวา อคารสฺมา อนคาริยํ ปพฺพชิโต, โส เอวํ}

\vspace{\tipitakaparskip}
\csromancenter{สามญฺญผลสุตฺต ปพฺพชิโต สมาโน กาเยน สํวุโต วิหรติ, วาจาย สํวุโต วิหรติ, มนสา สํวุโต วิหรติ, ฆาสจฺฉาทนปรมตาย สนฺตุฏฺโฐ, อภิรโต ปวิเวเก”ติ.\csromanspacer{} อปิ นุ ตฺวํ เอวํ วเทยฺยาสิ “เอตุ เม โภ โส ปุริโส, ปุนเทว โหตุ ทาโส กมฺมกาโร ปุพฺพุฏฺฐายี ปจฺฉานิปาตี กิํ การปฏิสฺสาวี มนาปจารี ปิยวาที มุขุลฺโลกโก”ติ.}
\proseitem{184}{\csromanfolio{57}โนเหตํ ภนฺเต, อถ โข นํ มยเมว อภิวาเทยฺยามปิ ปจฺจุฏฺเฐยฺยามปิ อาสเนนปิ นิมนฺเตยฺยาม อภินิมนฺเตยฺยามปิ นํ จีวรปิณฺฑปาตเสนาสนคิลานปฺปจฺจยเภสชฺชปริกฺขาเรหิ, ธมฺมิกมฺปิสฺส รกฺขาวรณคุตฺติํ สํวิทเหยฺยามาติ.}

\proseitem{185}{ตํ กิํ มญฺญสิ มหาราช.\csromanspacer{} ยทิ เอวํ สนฺเต โหติ วา สนฺทิฏฺฐิกํ สามญฺญผลํ โน วาติ.\csromanspacer{} อทฺธา โข ภนฺเต เอวํ สนฺเต โหติ สนฺทิฏฺฐิกํ สามญฺญผลนฺติ.\csromanspacer{} อิทํ โข เต มหาราช มยา ปฐมํ ทิฏฺเฐว ธมฺเม สนฺทิฏฺฐิกํ สามญฺญผลํ ปญฺญตฺตนฺติ.}

\csromanmatikaanchor{mtk.34}
\csromantocmarkref{section}{ทุติยสนฺทิฏฺฐิกสามญฺญผล}{mtk.34}
\bookmark[dest={mtk.34},level=1]{ทุติยสนฺทิฏฺฐิกสามญฺญผล}
\csromanheader{1}{ทุติย\-สนฺทิฏฺฐิก\-สามญฺญผล}
\proseitem{186}{สกฺกา ปน ภนฺเต อญฺญมฺปิ เอวเมว ทิฏฺเฐว ธมฺเม สนฺทิฏฺฐิกํ สามญฺญผลํ ปญฺญเปตุนฺติ.\csromanspacer{} สกฺกา มหาราช, เตน หิ มหาราช ตญฺเญเวตฺถ ปฏิปุจฺฉิสฺสามิ, ยถา เต ขเมยฺย, ตถา นํ พฺยากเรยฺยาสิ.\csromanspacer{} ตํ กิํ มญฺญสิ มหาราช.\csromanspacer{} อิธ เต อสฺส ปุริโส กสฺสโก คหปติโก กรการโก ราสิวฑฺฒโก.\csromanspacer{} ตสฺส เอวมสฺส “อจฺฉริยํ วต โภ อพฺภุตํ วต โภ ปุญฺญานํ คติ ปุญฺญานํ วิปาโก, อยํ หิ ราชา มาคโธ อชาตสตฺตุ เวเทหิปุตฺโต มนุสฺโส, อหมฺปิ มนุสฺโส, อยํ หิ ราชา มาคโธ อชาตสตฺตุ เวเทหิปุตฺโต ปญฺจหิ กามคุเณหิ สมปฺปิโต สมงฺคีภูโต ปริจาเรติ เทโว มญฺเญ, อหํ ปนมฺหิสฺส กสฺสโก คหปติโก กรการโก ราสิวฑฺฒโก, โส วตสฺสาหํ ปุญฺญานิ กเรยฺยํ, ยํนูนาหํ เกสมสฺสุํ โอหาเรตฺวา กาสายานิ วตฺตานิ อจฺฉาเทตฺวา อคารสฺมา อนคาริยํ ปพฺพเชยฺยนฺ”ติ.}

\csromanmatikaanchor{mtk.35}
\csromantocmarkref{section}{ปณีตตลสามญฺญผล}{mtk.35}
\bookmark[dest={mtk.35},level=1]{ปณีตตลสามญฺญผล}
\prose{โส อปเรน สมเยน อปฺปํ วา โภคกฺขนฺธํ ปหาย มหนฺตํ วา โภคกฺขนฺธํ ปหาย อปฺปํ วา ญาติปริวฏฺฏํ ปหาย มหนฺตํ วา ญาติปริวฏฺฏํ ปหาย เกสมสฺสุํ โอหาเรตฺวา กาสายานิ วตฺถานิ อจฺฉาเทตฺวา \csromanfolio{58}อคารสฺมา อนคาริยํ ปพฺพเชยฺย, โส เอวํ ปพฺพชิโต สมาโน กาเยน สํวุโต วิหเรยฺย, วาจาย สํวุโต วิหเรยฺย, มนสา สํวุโต วิหเรยฺย, ฆาสจฺฉาทนปรมตาย สนฺตุฏฺโฐ, อภิรโต ปวิเวเก.\csromanspacer{} ตํ เจ เต ปุริสา เอวมาโรเจยฺยุํ “ยคฺเฆ เทว ชาเนยฺยาสิ, โย เต โส ปุริโส\footnote{โย เต ปุริโส (สี)} กสฺสโก คหปติโก กรการโก ราสิวฑฺฒโก, โส เทว เกสมสฺสุํ โอหาเรตฺวา กาสายานิ วตฺถานิ อจฺฉาเทตฺวา อคารสฺมา อนคาริยํ ปพฺพชิโต, โส เอวํ ปพฺพชิโต สมาโน กาเยน สํวุโต วิหรติ, วาจาย สํวุโต วิหรติ, มนสา สํวุโต วิหรติ, ฆาสจฺฉาทนปรมตาย สนฺตุฏฺโฐ, อภิรโต ปวิเวเก”ติ.\csromanspacer{} อปิ นุ ตฺวํ เอวํ วเทยฺยาสิ “เอตุ เม โภ โส ปุริโส, ปุนเทว โหตุ กสฺสโก คหปติโก กรการโก ราสิวฑฺฒโก”ติ.}

\proseitem{187}{โนเหตํ ภนฺเต, อถ โข นํ มยเมว อภิวาเทยฺยามปิ ปจฺจุฏฺเฐยฺยามปิ อาสเนนปิ นิมนฺเตยฺยาม อภินิมนฺเตยฺยามปิ นํ จีวรปิณฺฑปาตเสนาสนคิลานปฺปจฺจยเภสชฺชปริกฺขาเรหิ, ธมฺมิกมฺปิสฺส รกฺขาวรณาคุตฺติํ สํวิทเหยฺยามาติ.}

\proseitem{188}{ตํ กิํ มญฺญสิ มหาราช.\csromanspacer{} ยทิ เอวํ สนฺเต โหติ วา สนฺทิฏฺฐิกํ สามญฺญผลํ โน วาติ.\csromanspacer{} อทฺธา โข ภนฺเต เอวํ สนฺเต โหติ สนฺทิฏฺฐิกํ สามญฺญผลนฺติ.\csromanspacer{} อิทํ โข เต มหาราช มยา ทุติยํ ทิฏฺเฐว ธมฺเม สนฺทิฏฺฐิกํ สามญฺญผลํ ปญฺญตฺตนฺติ.}

\csromanheader{2}{\textbf{ปณีตตรสามญฺญผล}}
\proseitem{189}{สกฺกา ปน ภนฺเต อญฺญมฺปิ ทิฏฺเฐว ธมฺเม สนฺทิฏฺฐิกํ สามญฺญผลํ ปญฺญเปตุํ อิเมหิ สนฺทิฏฺฐิเกหิ สามญฺญผเลหิ อภิกฺกนฺตตรญฺจ ปณีตตรญฺจาติ.\csromanspacer{} สกฺกา มหาราช, เตน หิ มหาราช สุโณหิ สาธุกํ มนสิ กโรหิ ภาสิสฺสามีติ. “เอวํ ภนฺเต”ติ โข ราชา มาคโธ อชาตสตฺตุ เวเทหิปุตฺโต ภควโต ปจฺจสฺโสสิ.}

\proseitem{190}{ภควา เอตทโวจ, อิธ มหาราช ตถาคโต โลเก อุปฺปชฺชติ อรหํ สมฺมาสมฺพุทฺโธ วิชฺชาจรณสมฺปนฺโน สุคโต โลกวิทู อนุตฺตโร ปุริสทมฺมสารถิ สตฺถา เทวมนุสฺสานํ พุทฺโธ ภควา, โส}

\vspace{\tipitakaparskip}
\csromancenter{สามญฺญผลสุตฺต อิมํ โลกํ สเทวกํ สมารกํ สพฺรหฺมกํ สสฺสมณพฺราหฺมณิํ ปชํ สเทวมนุสฺสํ สยํ อภิญฺญา สจฺฉิกตฺวา ปเวเทติ, โส ธมฺมํ เทเสติ อาทิกลฺยาณํ มชฺเฌกลฺยาณํ ปริโยสานกลฺยาณํ สาตฺถํ สพฺยญฺชนํ เกวลปริปุณฺณํ ปริสุทฺธํ พฺรหฺมจริยํ ปกาเสติ.}
\proseitem{191}{\csromanfolio{59}ตํ ธมฺมํ สุณาติ คหปติ วา คหปติปุตฺโต วา อญฺญตรสฺมิํ วา กุเล ปจฺจาชาโต, โส ตํ ธมฺมํ สุตฺวา ตถาคเต สทฺธํ ปฏิลภติ, โส เตน สทฺธาปฏิลาเภน สมนฺนาคโต อิติ ปฏิสญฺจิกฺขติ ‘สมฺพาโธ ฆราวาโส รโชปโถ, อพฺโภกาโส ปพฺพชฺชา, นยิทํ สุกรํ อคารํ อชฺฌาวสตา เอกนฺตปริปุณฺณํ เอกนฺตปริสุทฺธํ สงฺขลิขิตํ พฺรหฺมจริยํ จริตุํ, ยํนูนาหํ เกสมสฺสุํ โอหาเรตฺวา กาสายานิ วตฺถานิ อจฺฉาเทตฺวา อคารสฺมา อนคาริยํ ปพฺพเชยฺยนฺ’ติ.}

\proseitem{192}{โส อปเรน สมเยน อปฺปํ วา โภคกฺขนฺธํ ปหาย มหนฺตํ วา โภคกฺขนฺธํ ปหาย อปฺปํ วา ญาติปริวฏฺฏํ ปหาย มหนฺตํ วา ญาติปริวฏฺฏํ ปหาย เกสมสฺสุํ โอหาเรตฺวา กาสายานิ วตฺถานิ อจฺฉาเทตฺวา อคารสฺมา อนคาริยํ ปพฺพชติ.}

\proseitem{193}{โส เอวํ ปพฺพชิโต สมาโน ปาติโมกฺขสํวรสํวุโต วิหรติ อาจารโคจรสมฺปนฺโน, อณุมตฺเตสุ วชฺเชสุ ภยทสฺสาวี, สมาทาย สิกฺขติ สิกฺขาปเทสุ, กายกมฺมวจีกมฺเมน สมนฺนาคโต กุสเลน, ปริสุทฺธาชีโว สีลสมฺปนฺโน, อินฺทฺริเยสุ คุตฺตทฺวาโร\footnote{คุตฺตทฺวาโร, โภชเน มตฺตญฺญู (ก)} สติสมฺปชญฺเญน สมนฺนาคโต สนฺตุฏฺโฐ.}

\csromanmatikaanchor{mtk.36}
\csromantocmarkref{section}{จูฬสีล}{mtk.36}
\bookmark[dest={mtk.36},level=1]{จูฬสีล}
\csromanheader{1}{\textbf{จูฬสีล}}
\proseitem{194}{กถญฺจ มหาราช ภิกฺขุ สีลสมฺปนฺโน โหติ.\csromanspacer{} อิธ มหาราช ภิกฺขุ ปาณาติปาตํ ปหาย ปาณาติปาตา ปฏิวิรโต โหติ นิหิตทณฺโฑ นิหิตสตฺโถ ลชฺชี ทยาปนฺโน, สพฺพปาณภูต\-หิตา\-นุกมฺปี วิหรติ.\csromanspacer{} อิทมฺปิสฺส โหติ สีลสฺมิํ.}

\prose{\csromanfolio{60}อทินฺนาทานํ ปหาย อทินฺนาทานา ปฏิวิรโต โหติ ทินฺนาทายี ทินฺนปาฏิกงฺขี, อเถเนน สุจิภูเตน อตฺตนา วิหรติ.\csromanspacer{} อิทมฺปิสฺส โหติ สีลสฺมิํ.}

\prose{อพฺรหฺมจริยํ ปหาย พฺรหฺมจารี โหติ อาราจารี วิรโต เมถุนา คามธมฺมา.\csromanspacer{} อิทมฺปิสฺส โหติ สีลสฺมิํ.}

\prose{มุสาวาทํ ปหาย มุสาวาทา ปฏิวิรโต โหติ สจฺจวาที สจฺจสนฺโธ เถโต ปจฺจยิโก อวิสํวาทโก โลกสฺส.\csromanspacer{} อิทมฺปิสฺส โหติ สีลสฺมิํ.}

\prose{ปิสุณํ วาจํ ปหาย ปิสุณาย วาจาย ปฏิวิรโต โหติ, อิโต สุตฺวา น อมุตฺร อกฺขาตา อิเมสํ เภทาย, อมุตฺร วา สุตฺวา น อิเมสํ อกฺขาตา อมูสํ เภทาย, อิติ ภินฺนานํ วา สนฺธาตา, สหิตานํ วา อนุปฺปทาตา, สมคฺคาราโม สมคฺครโต สมคฺคนนฺที สมคฺคกรณิํ วาจํ ภาสิตา โหติ.\csromanspacer{} อิทมฺปิสฺส โหติ สีลสฺมิํ.}

\prose{ผรุสํ วาจํ ปหาย ผรุสาย วาจาย ปฏิวิรโต โหติ, ยา สา วาจา เนลา กณฺณสุขา เปมนียา หทยงฺคมา โปรี พหุชนกนฺตา พหุชนมนาปา, ตถารูปิํ วาจํ ภาสิตา โหติ.\csromanspacer{} อิทมฺปิสฺส โหติ สีลสฺมิํ.}

\prose{สมฺผปฺปลาปํ ปหาย สมฺผปฺปลาปา ปฏิวิรโต โหติ กาลวาที ภูตวาที อตฺถวาที ธมฺมวาที วินยวาที, นิธานวติํ วาจํ ภาสิตา โหติ กาเลน สาปเทสํ ปริยนฺตวติํ อตฺถสญฺหิตํ.\csromanspacer{} อิทมฺปิสฺส โหติ สีลสฺมิํ.}

\prose{พีชคาม\-ภูตคาม\-สมารมฺภา ปฏิวิรโต โหติ ฯเปฯ.\csromanspacer{} เอกภตฺติโก โหติ รตฺตูปรโต วิรโต วิกาลโภชนา.\csromanspacer{} นจฺจ\-คีต\-วาทิต\-วิสูก\-ทสฺสนา ปฏิวิรโต โหติ.\csromanspacer{} มาลาคนฺธวิเลปนธารณมณฺฑนวิภูสนฏฺฐานา ปฏิวิรโต โหติ.\csromanspacer{} อุจฺจาสยนมหาสยนา ปฏิวิรโต โหติ.\csromanspacer{} ชาตรูป\-รชต\-ปฏิคฺคหณา ปฏิวิรโต โหติ.\csromanspacer{} อามกธญฺญปฏิคฺคหณา ปฏิวิรโต โหติ.\csromanspacer{} อามกมํสปฏิคฺคหณา ปฏิวิรโต โหติ.\csromanspacer{} อิตฺถิ\-กุมาริก\-ปฏิคฺคหณา ปฏิวิรโต โหติ.\csromanspacer{} ทาสิทาสปฏิคฺคหณา ปฏิวิรโต โหติ.\csromanspacer{} อเชฬกปฏิคฺคหณา ปฏิวิรโต โหติ.}

\vspace{\tipitakaparskip}
\csromancenter{สามญฺญผลสุตฺต กุกฺกุฏสูกรปฏิคฺคหณา ปฏิวิรโต โหติ.\csromanspacer{} หตฺถิควสฺสวฬวปฏิคฺคหณา ปฏิวิรโต โหติ.\csromanspacer{} เขตฺตวตฺถุปฏิคฺคหณา ปฏิวิรโต โหติ.\csromanspacer{} ทูเตยฺย\-ปหิณ\-คมนา\-นุโยคา ปฏิวิรโต โหติ.\csromanspacer{} กยวิกฺกยา ปฏิวิรโต โหติ.\csromanspacer{} ตุลา\-กูฏ\-กํส\-กูฏ\-มาน\-กูฏา ปฏิวิรโต โหติ.\csromanspacer{} อุกฺโกฏน\-วญฺจน\-นิกติ\-สาจิโยคา ปฏิวิรโต โหติ.\csromanspacer{} เฉทนวธพนฺธนวิปราโมส-อาโลปสหสาการา ปฏิวิรโต โหติ.\csromanspacer{} อิทมฺปิสฺส โหติ สีลสฺมิํ.}
\nitthitamruled{จูฬสีลํ นิฏฺฐิตํ.}
\csromanmatikaanchor{mtk.37}
\csromantocmarkref{section}{มชฺฌิมสีล}{mtk.37}
\bookmark[dest={mtk.37},level=1]{มชฺฌิมสีล}
\csromanheader{1}{\textbf{มชฺฌิมสีล}}
\proseitem{195}{\csromanfolio{61}ยถา วา ปเนเก โภนฺโต สมณพฺราหฺมณา สทฺธาเทยฺยานิ โภชนานิ ภุญฺชิตฺวา เต เอวรูปํ พีชคาม\-ภูตคาม\-สมารมฺภํ อนุยุตฺตา วิหรนฺติ.\csromanspacer{} เสยฺยถิทํ, มูลพีชํ ขนฺธพีชํ ผฬุพีชํ อคฺคพีชํ พีชพีชเมว ปญฺจมํ, อิติ เอวรูปา พีชคาม\-ภูตคาม\-สมารมฺภา ปฏิวิรโต โหติ.\csromanspacer{} อิทมฺปิสฺส โหติ สีลสฺมิํ.}

\proseitem{196}{ยถา วา ปเนเก โภนฺโต สมณพฺราหฺมณา สทฺธาเทยฺยานิ โภชนานิ ภุญฺชิตฺวา เต เอวรูปํ สนฺนิธิการปริโภคํ อนุยุตฺตา วิหรนฺติ.\csromanspacer{} เสยฺยถิทํ, อนฺนสนฺนิธิํ ปานสนฺนิธิํ วตฺถสนฺนิธิํ ยานสนฺนิธิํ สยนสนฺนิธิํ คนฺธสนฺนิธิํ อามิสสนฺนิธิํ อิติ วา, อิติ เอวรูปา สนฺนิธิการปริโภคา ปฏิวิรโต โหติ.\csromanspacer{} อิทมฺปิสฺส โหติ สีลสฺมิํ.}

\proseitem{197}{ยถา วา ปเนเก โภนฺโต สมณพฺราหฺมณา สทฺธาเทยฺยานิ โภชนานิ ภุญฺชิตฺวา เต เอวรูปํ วิสูกทสฺสนํ อนุยุตฺตา วิหรนฺติ.\csromanspacer{} เสยฺยถิทํ, นจฺจํ คีตํ วาทิตํ เปกฺขํ อกฺขานํ ปาณิสฺสรํ เวตาฬํ กุมฺภถูณํ โสภนกํ จณฺฑาลํ วํสํ โธวนํ หตฺถิยุทฺธํ อสฺสยุทฺธํ มหิํสยุทฺธํ อุสภยุทฺธํ อชยุทฺธํ เมณฺฑยุทฺธํ กุกฺกุฏยุทฺธํ วฏฺฏกยุทฺธํ ทณฺฑยุทฺธํ มุฏฺฐิยุทฺธํ นิพฺพุทฺธํ อุยฺโยธิกํ พลคฺคํ เสนาพฺยูหํ อนีกทสฺสนํ อิติ วา, อิติ เอวรูปา วิสูกทสฺสนา ปฏิวิรโต โหติ.\csromanspacer{} อิทมฺปิสฺส โหติ สีลสฺมิํ.}

\proseitem{198}{\csromanfolio{62}ยถา วา ปเนเก โภนฺโต สมณพฺราหฺมณา สทฺธาเทยฺยานิ โภชนานิ ภุญฺชิตฺวา เต เอวรูปํ ชูตปฺปมาทฏฺฐานานุโยคํ อนุยุตฺตา วิหรนฺติ.\csromanspacer{} เสยฺยถิทํ, อฏฺฐปทํ ทสปทํ อากาสํ ปริหารปถํ สนฺติกํ ขลิกํ ฆฏิกํ สลากหตฺถํ อกฺขํ ปงฺคจีรํ วงฺกกํ โมกฺขจิกํ จิงฺคุลิกํ ปตฺตาฬฺหกํ รถกํ ธนุกํ อกฺขริกํ มเนสิกํ ยถาวชฺชํ อิติ วา, อิติ เอวรูปา ชูตปฺปมาทฏฺฐานานุโยคา ปฏิวิรโต โหติ.\csromanspacer{} อิทมฺปิสฺส โหติ สีลสฺมิํ.}

\proseitem{199}{ยถา วา ปเนเก โภนฺโต สมณพฺราหฺมณา สทฺธาเทยฺยานิ โภชนานิ ภุญฺชิตฺวา เต เอวรูปํ อุจฺจาสยนมหาสยนํ อนุยุตฺตา วิหรนฺติ.\csromanspacer{} เสยฺยถิทํ, อาสนฺทิํ ปลฺลงฺกํ โคนกํ จิตฺตกํ ปฏิกํ ปฏลิกํ ตูลิกํ วิกติกํ อุทฺทโลมิํ เอกนฺตโลมิํ กฏฺฏิสฺสํ โกเสยฺยํ กุตฺตกํ หตฺถตฺถรํ อสฺสตฺถรํ รถตฺถรํ อชินปฺปเวณิํ กทลิ\-มิค\-ปวร\-ปจฺจตฺถรณํ ส-อุตฺตรจฺฉทํ อุภโตโลหิตกูปธานํ อิติ วา, อิติ เอวรูปา อุจฺจาสยนมหาสยนา ปฏิวิรโต โหติ.\csromanspacer{} อิทมฺปิสฺส โหติ สีลสฺมิํ.}

\proseitem{200}{ยถา วา ปเนเก โภนฺโต สมณพฺราหฺมณา สทฺธาเทยฺยานิ โภชนานิ ภุญฺชิตฺวา เต เอวรูปํ มณฺฑน\-วิภูสนฏฺฐานา\-นุโยคํ อนุยุตฺตา วิหรนฺติ.\csromanspacer{} เสยฺยถิทํ, อุจฺฉาทนํ ปริมทฺทนํ นฺหาปนํ สมฺพาหนํ อาทาสํ อญฺชนํ มาลาคนฺธวิเลปนํ มุขจุณฺณํ มุขเลปนํ หตฺถพนฺธํ สิขาพนฺธํ ทณฺฑํ นาฬิกํ อสิํ\footnote{ขคฺคํ (สี, อิ), อสิํ ขคฺคํ (สฺยา, กํ), ขคฺคํ อสิํ (ก)} ฉตฺตํ จิตฺรุปาหนํ อุณฺหีสํ มณิํ วาลพีชนิํ โอทาตานิ วตฺถานิ ทีฆทสานิ อิติ วา, อิติ เอวรูปา มณฺฑน\-วิภูสนฏฺฐานา\-นุโยคา ปฏิวิรโต โหติ.\csromanspacer{} อิทมฺปิสฺส โหติ สีลสฺมิํ.}

\proseitem{201}{ยถา วา ปเนเก โภนฺโต สมณพฺราหฺมณา สทฺธาเทยฺยานิ โภชนานิ ภุญฺชิตฺวา เต เอวรูปํ ติรจฺฉานกถํ อนุยุตฺตา วิหรนฺติ.\csromanspacer{} เสยฺยถิทํ, ราชกถํ โจรกถํ มหามตฺตกถํ เสนากถํ ภยกถํ ยุทฺธกถํ อนฺนกถํ ปานกถํ วตฺถกถํ สยนกถํ มาลากถํ คนฺธกถํ ญาติกถํ ยานกถํ คามกถํ นิคมกถํ นครกถํ ชนปทกถํ อิตฺถิกถํ\footnote{อิตฺถิกถํ ปุริสกถํ กุมารกถํ กุมาริกถํ (ก)} สูรกถํ วิสิขากถํ กุมฺภฏฺฐานกถํ ปุพฺพเปตกถํ นานตฺตกถํ}

\vspace{\tipitakaparskip}
\csromancenter{สามญฺญผลสุตฺต โลกกฺขายิกํ สมุทฺทกฺขายิกํ อิติภวาภวกถํ อิติ วา, อิติ เอวรูปาย ติรจฺฉานกถาย ปฏิวิรโต โหติ.\csromanspacer{} อิทมฺปิสฺส โหติ สีลสฺมิํ.}
\proseitem{202}{\csromanfolio{63}ยถา วา ปเนเก โภนฺโต สมณพฺราหฺมณา สทฺธาเทยฺยานิ โภชนานิ ภุญฺชิตฺวา เต เอวรูปํ วิคฺคาหิกกถํ อนุยุตฺตา วิหรนฺติ.\csromanspacer{} เสยฺยถิทํ, น ตฺวํ อิมํ ธมฺมวินยํ อาชานาสิ, อหํ อิมํ ธมฺมวินยํ อาชานามิ, กิํ ตฺวํ อิมํ ธมฺมวินยํ อาชานิสฺสสิ, มิจฺฉา ปฏิปนฺโน ตฺวมสิ, อหมสฺมิ สมฺมา ปฏิปนฺโน, สหิตํ เม, อสหิตํ เต, ปุเร วจนียํ ปจฺฉา อวจ, ปจฺฉา วจนียํ ปุเร อวจ, อธิจิณฺณํ เต วิปราวตฺตํ, อาโรปิโต เต วาโท, นิคฺคหิโต ตฺวมสิ, จร วาทปฺปโมกฺขาย, นิพฺเพเฐหิ วา สเจ ปโหสีติ อิติ วา, อิติ เอวรูปาย วิคฺคาหิกกถาย ปฏิวิรโต โหติ.\csromanspacer{} อิทมฺปิสฺส โหติ สีลสฺมิํ.}

\proseitem{203}{ยถา วา ปเนเก โภนฺโต สมณพฺราหฺมณา สทฺธาเทยฺยานิ โภชนานิ ภุญฺชิตฺวา เต เอวรูปํ ทูเตยฺย\-ปหิณ\-คมนา\-นุโยคํ อนุยุตฺตา วิหรนฺติ.\csromanspacer{} เสยฺยถิทํ, รญฺญํ, ราชมหามตฺตานํ, ขตฺติยานํ, พฺราหฺมณานํ, คหปติกานํ, กุมารานํ “อิธ คจฺฉ, อมุตฺราคจฺฉ, อิทํ หร, อมุตฺร อิทํ อาหรา”ติ อิติ วา, อิติ เอวรูปา ทูเตยฺย\-ปหิณ\-คมนา\-นุโยคา ปฏิวิรโต โหติ.\csromanspacer{} อิทมฺปิสฺส โหติ สีลสฺมิํ.}

\proseitemwithcloserruled{204}{ยถา วา ปเนเก โภนฺโต สมณพฺราหฺมณา สทฺธาเทยฺยานิ โภชนานิ ภุญฺชิตฺวา เต กุหกา จ โหนฺติ ลปกา จ เนมิตฺติกา จ นิปฺเปสิกา จ ลาเภน ลาภํ นิชิคีสิตาโร จ, อิติ เอวรูปา กุหนลปนา ปฏิวิรโต โหติ.\csromanspacer{} อิทมฺปิสฺส โหติ สีลสฺมิํ.}{มชฺฌิมสีลํ นิฏฺฐิตํ.}
\csromanmatikaanchor{mtk.38}
\csromantocmarkref{section}{มหาสีล}{mtk.38}
\bookmark[dest={mtk.38},level=1]{มหาสีล}
\csromanheader{1}{\textbf{มหาสีล}}
\proseitem{205}{ยถา วา ปเนเก โภนฺโต สมณพฺราหฺมณา สทฺธาเทยฺยานิ โภชนานิ ภุญฺชิตฺวา เต เอวรูปาย ติรจฺฉานวิชฺชาย มิจฺฉาชีเวน ชีวิตํ กปฺเปนฺติ.\csromanspacer{} เสยฺยถิทํ, องฺคํ นิมิตฺตํ อุปฺปาตํ สุปินํ ลกฺขณํ มูสิกจฺฉินฺนํ อคฺคิโหมํ ทพฺพิโหมํ ถุสโหมํ กณโหมํ ตณฺฑุลโหมํ สปฺปิโหมํ \csromanfolio{64}เตลโหมํ มุขโหมํ โลหิตโหมํ องฺควิชฺชา วตฺถุวิชฺชา ขตฺตวิชฺชา สิววิชฺชา ภูตวิชฺชา ภูริวิชฺชา อหิวิชฺชา วิสวิชฺชา วิจฺฉิกวิชฺชา มูสิกวิชฺชา สกุณวิชฺชา วายสวิชฺชา ปกฺกชฺฌานํ สรปริตฺตาณํ มิคจกฺกํ อิติ วา, อิติ เอวรูปาย ติรจฺฉานวิชฺชาย มิจฺฉาชีวา ปฏิวิรโต โหติ.\csromanspacer{} อิทมฺปิสฺส โหติ สีลสฺมิํ.}

\proseitem{206}{ยถา วา ปเนเก โภนฺโต สมณพฺราหฺมณา สทฺธาเทยฺยานิ โภชนานิ ภุญฺชิตฺวา เต เอวรูปาย ติรจฺฉานวิชฺชาย มิจฺฉาชีเวน ชีวิตํ กปฺเปนฺติ.\csromanspacer{} เสยฺยถิทํ, มณิลกฺขณํ วตฺถลกฺขณํ ทณฺฑลกฺขณํ สตฺถลกฺขณํ อสิลกฺขณํ อุสุลกฺขณํ ธนุลกฺขณํ อาวุธลกฺขณํ อิตฺถิลกฺขณํ ปุริสลกฺขณํ กุมารลกฺขณํ กุมาริลกฺขณํ ทาสลกฺขณํ ทาสิลกฺขณํ หตฺถิลกฺขณํ อสฺสลกฺขณํ มหิํสลกฺขณํ อุสภลกฺขณํ โคลกฺขณํ อชลกฺขณํ เมณฺฑลกฺขณํ กุกฺกุฏลกฺขณํ วฏฺฏกลกฺขณํ โคธาลกฺขณํ กณฺณิกลกฺขณํ กจฺฉปลกฺขณํ มิคลกฺขณํ อิติ วา, อิติ เอวรูปาย ติรจฺฉานวิชฺชาย มิจฺฉาชีวา ปฏิวิรโต โหติ.\csromanspacer{} อิทมฺปิสฺส โหติ สีลสฺมิํ.}

\proseitem{207}{ยถา วา ปเนเก โภนฺโต สมณพฺราหฺมณา สทฺธาเทยฺยานิ โภชนานิ ภุญฺชิตฺวา เต เอวรูปาย ติรจฺฉานวิชฺชาย มิจฺฉาชีเวน ชีวิตํ กปฺเปนฺติ.\csromanspacer{} เสยฺยถิทํ, รญฺญํ นิยฺยานํ ภวิสฺสติ, รญฺญํ อนิยฺยานํ ภวิสฺสติ, อพฺภนฺตรานํ รญฺญํ อุปยานํ ภวิสฺสติ, พาหิรานํ รญฺญํ อปยานํ ภวิสฺสติ, พาหิรานํ รญฺญํ อุปยานํ ภวิสฺสติ, อพฺภนฺตรานํ รญฺญํ อปยานํ ภวิสฺสติ, อพฺภนฺตรานํ รญฺญํ ชโย ภวิสฺสติ, พาหิรานํ รญฺญํ ปราชโย ภวิสฺสติ, พาหิรานํ รญฺญํ ชโย ภวิสฺสติ, อพฺภนฺตรานํ รญฺญํ ปราชโย ภวิสฺสติ, อิติ อิมสฺส ชโย ภวิสฺสติ, อิมสฺส ปราชโย ภวิสฺสติ อิติ วา, อิติ เอวรูปาย ติรจฺฉานวิชฺชาย มิจฺฉาชีวา ปฏิวิรโต โหติ.\csromanspacer{} อิทมฺปิสฺส โหติ สีลสฺมิํ.}

\proseitem{208}{ยถา วา ปเนเก โภนฺโต สมณพฺราหฺมณา สทฺธาเทยฺยานิ โภชนานิ ภุญฺชิตฺวา เต เอวรูปาย ติรจฺฉานวิชฺชาย มิจฺฉาชีเวน ชีวิตํ กปฺเปนฺติ.\csromanspacer{} เสยฺยถิทํ, จนฺทคฺคาโห ภวิสฺสติ, สูริยคฺคาโห ภวิสฺสติ, นกฺขตฺตคฺคาโห ภวิสฺสติ, จนฺทิมสูริยานํ ปถคมนํ ภวิสฺสติ, จนฺทิมสูริยานํ อุปฺปถคมนํ ภวิสฺสติ, นกฺขตฺตานํ ปถคมนํ ภวิสฺสติ, นกฺขตฺตานํ อุปฺปถคมนํ}

\vspace{\tipitakaparskip}
\csromancenter{สามญฺญผลสุตฺต ภวิสฺสติ, อุกฺกาปาโต ภวิสฺสติ, ทิสาฑาโห ภวิสฺสติ, ภูมิจาโล ภวิสฺสติ, เทวทุทฺรภิ ภวิสฺสติ, จนฺทิมสูริยนกฺขตฺตานํ อุคฺคมนํ โอคมนํ สํกิเลสํ โวทานํ ภวิสฺสติ, เอวํวิปาโก จนฺทคฺคาโห ภวิสฺสติ, เอวํวิปาโก สูริยคฺคาโห ภวิสฺสติ, เอวํวิปาโก นกฺขตฺตคฺคาโห ภวิสฺสติ, เอวํวิปากํ จนฺทิมสูริยานํ ปถคมนํ ภวิสฺสติ, เอวํวิปากํ จนฺทิมสูริยานํ อุปฺปถคมนํ ภวิสฺสติ, เอวํวิปากํ นกฺขตฺตานํ ปถคมนํ ภวิสฺสติ, เอวํวิปากํ นกฺขตฺตานํ อุปฺปถคมนํ ภวิสฺสติ, เอวํวิปาโก อุกฺกาปาโต ภวิสฺสติ, เอวํวิปาโก ทิสาฑาโห ภวิสฺสติ, เอวํวิปาโก ภูมิจาโล ภวิสฺสติ, เอวํวิปาโก เทวทุทฺรภิ ภวิสฺสติ, เอวํวิปากํ จนฺทิมสูริยนกฺขตฺตานํ อุคฺคมนํ โอคมนํ สํกิเลสํ โวทานํ ภวิสฺสติ อิติ วา, อิติ เอวรูปาย ติรจฺฉานวิชฺชาย มิจฺฉาชีวา ปฏิวิรโต โหติ.\csromanspacer{} อิทมฺปิสฺส โหติ สีลสฺมิํ.}
\proseitem{209}{\csromanfolio{65}ยถา วา ปเนเก โภนฺโต สมณพฺราหฺมณา สทฺธาเทยฺยานิ โภชนานิ ภุญฺชิตฺวา เต เอวรูปาย ติรจฺฉานวิชฺชาย มิจฺฉาชีเวน ชีวิตํ กปฺเปนฺติ.\csromanspacer{} เสยฺยถิทํ, สุวุฏฺฐิกา ภวิสฺสติ, ทุพฺพุฏฺฐิกา ภวิสฺสติ, สุภิกฺขํ ภวิสฺสติ, ทุพฺภิกฺขํ ภวิสฺสติ, เขมํ ภวิสฺสติ, ภยํ ภวิสฺสติ, โรโค ภวิสฺสติ, อาโรคฺยํ ภวิสฺสติ, มุทฺทา, คณนา, สงฺขานํ, กาเวยฺยํ, โลกายตํ อิติ วา, อิติ เอวรูปาย ติรจฺฉานวิชฺชาย มิจฺฉาชีวา ปฏิวิรโต โหติ.\csromanspacer{} อิทมฺปิสฺส โหติ สีลสฺมิํ.}

\proseitem{210}{ยถา วา ปเนเก โภนฺโต สมณพฺราหฺมณา สทฺธาเทยฺยานิ โภชนานิ ภุญฺชิตฺวา เต เอวรูปาย ติรจฺฉานวิชฺชาย มิจฺฉาชีเวน ชีวิตํ กปฺเปนฺติ.\csromanspacer{} เสยฺยถิทํ, อาวาหนํ วิวาหนํ สํวรณํ วิวรณํ สํกิรณํ วิกิรณํ สุภคกรณํ ทุพฺภคกรณํ วิรุทฺธคพฺภกรณํ ชิวฺหานิพนฺธนํ หนุสํหนนํ หตฺถาภิชปฺปนํ หนุชปฺปนํ กณฺณชปฺปนํ อาทาสปญฺหํ กุมาริกปญฺหํ เทวปญฺหํ อาทิจฺจุปฏฺฐานํ มหตุปฏฺฐานํ อพฺภุชฺชลนํ สิริวฺหายนํ อิติ วา, อิติ เอวรูปาย ติรจฺฉานวิชฺชาย มิจฺฉาชีวา ปฏิวิรโต โหติ.\csromanspacer{} อิทมฺปิสฺส โหติ สีลสฺมิํ.}

\proseitem{211}{ยถา วา ปเนเก โภนฺโต สมณพฺราหฺมณา สทฺธาเทยฺยานิ โภชนานิ ภุญฺชิตฺวา เต เอวรูปาย ติรจฺฉานวิชฺชาย มิจฺฉาชีเวน ชีวิตํ \csromanfolio{66}กปฺเปนฺติ.\csromanspacer{} เสยฺยถิทํ, สนฺติกมฺมํ ปณิธิกมฺมํ ภูตกมฺมํ ภูริกมฺมํ วสฺสกมฺมํ โวสฺสกมฺมํ วตฺถุกมฺมํ วตฺถุปริกมฺมํ อาจมนํ นฺหาปนํ ชุหนํ วมนํ วิเรจนํ อุทฺธํวิเรจนํ อโธวิเรจนํ สีสวิเรจนํ กณฺณเตลํ เนตฺตตปฺปนํ นตฺถุกมฺมํ อญฺชนํ ปจฺจญฺชนํ สาลากิยํ สลฺลกตฺติยํ ทารกติกิจฺฉา, มูลเภสชฺชานํ อนุปฺปทานํ, โอสธีนํ ปฏิโมกฺโข อิติ วา, อิติ เอวรูปาย ติรจฺฉานวิชฺชาย มิจฺฉาชีวา ปฏิวิรโต โหติ.\csromanspacer{} อิทมฺปิสฺส โหติ สีลสฺมิํ.}

\proseitemwithcloserruled{212}{ส โข โส มหาราช ภิกฺขุ เอวํ สีลสมฺปนฺโน น กุโตจิ ภยํ สมนุปสฺสติ, ยทิทํ สีลสํวรโต.\csromanspacer{} เสยฺยถาปิ มหาราช ราชา ขตฺติโย มุทฺธาภิสิตฺโต นิหตปจฺจามิตฺโต น กุโตจิ ภยํ สมนุปสฺสติ, ยทิทํ ปจฺจตฺถิกโต.\csromanspacer{} เอวเมว โข มหาราช ภิกฺขุ เอวํ สีลสมฺปนฺโน น กุโตจิ ภยํ สมนุปสฺสติ, ยทิทํ สีลสํวรโต, โส อิมินา อริเยน สีลกฺขนฺเธน สมนฺนาคโต อชฺฌตฺตํ อนวชฺชสุขํ ปฏิสํเวเทติ.\csromanspacer{} เอวํ โข มหาราช ภิกฺขุ สีลสมฺปนฺโน โหติ.}{มหาสีลํ นิฏฺฐิตํ.}
\csromanmatikaanchor{mtk.39}
\csromantocmarkref{section}{อินฺทฺริยสํวร}{mtk.39}
\bookmark[dest={mtk.39},level=1]{อินฺทฺริยสํวร}
\csromanheader{1}{\textbf{อินฺทฺริยสํวร}}
\proseitem{213}{กถญฺจ มหาราช ภิกฺขุ อินฺทฺริเยสุ คุตฺตทฺวาโร โหติ.\csromanspacer{} อิธ มหาราช ภิกฺขุ จกฺขุนา รูปํ ทิสฺวา น นิมิตฺตคฺคาหี โหติ นานุพฺยญฺชนคฺคาหี, ยตฺวาธิกรณเมนํ จกฺขุนฺทฺริยํ อสํวุตํ วิหรนฺตํ อภิชฺฌา โทมนสฺสา ปาปกา อกุสลา ธมฺมา อนฺวาสฺสเวยฺยุํ, ตสฺส สํวราย ปฏิปชฺชติ, รกฺขติ จกฺขุนฺทฺริยํ, จกฺขุนฺทฺริเย สํวรํ อาปชฺชติ, โสเตน สทฺทํ สุตฺวา ฯเปฯ ฆาเนน คนฺธํ ฆายิตฺวา ฯเปฯ ชิวฺหาย รสํ สายิตฺวา ฯเปฯ กาเยน โผฏฺฐพฺพํ ผุสิตฺวา ฯเปฯ มนสา ธมฺมํ วิญฺญาย น นิมิตฺตคฺคาหี โหติ นานุพฺยญฺชนคฺคาหี, ยตฺวาธิกรณเมนํ มนินฺทฺริยํ อสํวุตํ วิหรนฺตํ อภิชฺฌา โทมนสฺสา ปาปกา อกุสลา ธมฺมา อนฺวาสฺสเวยฺยุํ, ตสฺส สํวราย ปฏิปชฺชติ, รกฺขติ มนินฺทฺริยํ, มนินฺทฺริเย สํวรํ อาปชฺชติ, โส อิมินา อริเยน อินฺทฺริยสํวเรน สมนฺนาคโต อชฺฌตฺตํ อพฺยาเสกสุขํ ปฏิสํเวเทติ.\csromanspacer{} เอวํ โข มหาราช ภิกฺขุ อินฺทฺริเยสุ คุตฺตทฺวาโร โหติ.}

\csromanmatikaanchor{mtk.40}
\csromantocmarkref{section}{สติสมฺปชญฺญ}{mtk.40}
\bookmark[dest={mtk.40},level=1]{สติสมฺปชญฺญ}
\csromanheader{1}{2. สามญฺญผลสุตฺต \textbf{สติสมฺปชญฺญ}}
\proseitem{214}{\csromanfolio{67}กถญฺจ มหาราช ภิกฺขุ สติสมฺปชญฺเญน สมนฺนาคโต โหติ.\csromanspacer{} อิธ มหาราช ภิกฺขุ อภิกฺกนฺเต ปฏิกฺกนฺเต สมฺปชานการี โหติ, อาโลกิเต วิโลกิเต สมฺปชานการี โหติ, สมิญฺชิเต ปสาริเต สมฺปชานการี โหติ, สํฆาฏิปตฺตจีวรธารเณ สมฺปชานการี โหติ, อสิเต ปีเต ขายิเต สายิเต สมฺปชานการี โหติ, อุจฺจารปสฺสาวกมฺเม สมฺปชานการี โหติ, คเต ฐิเต นิสินฺเน สุตฺเต ชาคริเต ภาสิเต ตุณฺหีภาเว สมฺปชานการี โหติ.\csromanspacer{} เอวํ โข มหาราช ภิกฺขุ สติสมฺปชญฺเญน สมนฺนาคโต โหติ.}

\csromanmatikaanchor{mtk.41}
\csromantocmarkref{section}{สนฺโตส}{mtk.41}
\bookmark[dest={mtk.41},level=1]{สนฺโตส}
\csromanheader{1}{\textbf{สนฺโตส}}
\proseitem{215}{กถญฺจ มหาราช ภิกฺขุ สนฺตุฏฺโฐ โหติ.\csromanspacer{} อิธ มหาราช ภิกฺขุ สนฺตุฏฺโฐ โหติ กายปริหาริเกน จีวเรน กุจฺฉิปริหาริเกน ปิณฺฑปาเตน, โส เยน เยเนว ปกฺกมติ, สมาทาเยว ปกฺกมติ.\csromanspacer{} เสยฺยถาปิ มหาราช ปกฺขี สกุโณ เยน เยเนว เฑติ, สปตฺตภาโรว เฑติ.\csromanspacer{} เอวเมว โข มหาราช ภิกฺขุ สนฺตุฏฺโฐ โหติ กายปริหาริเกน จีวเรน กุจฺฉิปริหาริเกน ปิณฺฑปาเตน, โส เยน เยเนว ปกฺกมติ, สมาทาเยว ปกฺกมติ.\csromanspacer{} เอวํ โข มหาราช ภิกฺขุ สนฺตุฏฺโฐ โหติ.}

\csromanmatikaanchor{mtk.42}
\csromantocmarkref{section}{นีวรณปฺปหาน}{mtk.42}
\bookmark[dest={mtk.42},level=1]{นีวรณปฺปหาน}
\csromanheader{1}{\textbf{นีวรณปฺปหาน}}
\proseitem{216}{โส อิมินา จ อริเยน สีลกฺขนฺเธน สมนฺนาคโต อิมินา จ อริเยน อินฺทฺริยสํวเรน สมนฺนาคโต อิมินา จ อริเยน สติสมฺปชญฺเญน สมนฺนาคโต อิมาย จ อริยาย สนฺตุฏฺฐิยา สมนฺนาคโต วิวิตฺตํ เสนาสนํ ภชติ อรญฺญํ รุกฺขมูลํ ปพฺพตํ กนฺทรํ คิริคุหํ สุสานํ วนปตฺถํ อพฺโภกาสํ ปลาลปุญฺชํ.\csromanspacer{} โส ปจฺฉาภตฺตํ ปิณฺฑปาตปฺปฏิกฺกนฺโต นิสีทติ ปลฺลงฺกํ อาภุชิตฺวา อุชุํ กายํ ปณิธาย ปริมุขํ สติํ อุปฏฺฐเปตฺวา.}

\proseitem{217}{โส อภิชฺฌํ โลเก ปหาย วิคตาภิชฺเฌน เจตสา วิหรติ, อภิชฺฌาย จิตฺตํ ปริโสเธติ.\csromanspacer{} พฺยาปาทปโทสํ ปหาย \csromanfolio{68}อพฺยาปนฺนจิตฺโต วิหรติ สพฺพปาณภูต\-หิตา\-นุกมฺปี, พฺยาปาทปโทสา จิตฺตํ ปริโสเธติ.\csromanspacer{} ถินมิทฺธํ ปหาย วิคตถินมิทฺโธ วิหรติ อาโลกสญฺญี สโต สมฺปชาโน, ถินมิทฺธา จิตฺตํ ปริโสเธติ.\csromanspacer{} อุทฺธจฺจกุกฺกุจฺจํ ปหาย อนุทฺธโต วิหรติ อชฺฌตฺตํ วูปสนฺตจิตฺโต, อุทฺธจฺจกุกฺกุจฺจา จิตฺตํ ปริโสเธติ.\csromanspacer{} วิจิกิจฺฉํ ปหาย ติณฺณวิจิกิจฺโฉ วิหรติ อกถํกถี กุสเลสุ ธมฺเมสุ, วิจิกิจฺฉาย จิตฺตํ ปริโสเธติ.}

\proseitem{218}{เสยฺยถาปิ มหาราช ปุริโส อิณํ อาทาย กมฺมนฺเต ปโยเชยฺย, ตสฺส เต กมฺมนฺตา สมิชฺเฌยฺยุํ, โส ยานิ จ โปราณานิ อิณมูลานิ, ตานิ จ พฺยนฺติํ กเรยฺย\footnote{พฺยนฺตีกเรยฺย (สี, สฺยา, กํ)}, สิยา จสฺส อุตฺตริํ อวสิฏฺฐํ ทารภรณาย.\csromanspacer{} ตสฺส เอวมสฺส “อหํ โข ปุพฺเพ อิณํ อาทาย กมฺมนฺเต ปโยเชสิํ, ตสฺส เม เต กมฺมนฺตา สมิชฺฌิํสุ, โสหํ ยานิ จ โปราณานิ อิณมูลานิ, ตานิ จ พฺยนฺติํ อกาสิํ, อตฺถิ จ เม อุตฺตริํ อวสิฏฺฐํ ทารภรณายา”ติ.\csromanspacer{} โส ตโตนิทานํ ลเภถ ปาโมชฺชํ, อธิคจฺเฉยฺย โสมนสฺสํ.}

\proseitem{219}{เสยฺยถาปิ มหาราช ปุริโส อาพาธิโก อสฺส ทุกฺขิโต พาฬฺหคิลาโน, ภตฺตญฺจสฺส นจฺฉาเทยฺย, น จสฺส กาเย พลมตฺตา, โส อปเรน สมเยน ตมฺหา อาพาธา มุจฺเจยฺย, ภตฺตญฺจสฺส ฉาเทยฺย, สิยา จสฺส กาเย พลมตฺตา.\csromanspacer{} ตสฺส เอวมสฺส “อหํ โข ปุพฺเพ อาพาธิโก อโหสิํ ทุกฺขิโต พาฬฺหคิลาโน, ภตฺตญฺจ เม นจฺฉาเทสิ, น จ เม อาสิ\footnote{น จสฺส เม (ก)} กาเย พลมตฺตา, โสมฺหิ เอตรหิ ตมฺหา อาพาธา มุตฺโต, ภตฺตญฺจ เม ฉาเทติ, อตฺถิ จ เม กาเย พลมตฺตา”ติ.\csromanspacer{} โส ตโตนิทานํ ลเภถ ปาโมชฺชํ, อธิคจฺเฉยฺย โสมนสฺสํ.}

\proseitem{220}{เสยฺยถาปิ มหาราช ปุริโส พนฺธนาคาเร พทฺโธ อสฺส, โส อปเรน สมเยน ตมฺหา พนฺธนาคารา มุจฺเจยฺย โสตฺถินา อพฺภเยน\footnote{อพฺพเยน (สี, ก)}, น จสฺส กิญฺจิ โภคานํ วโย.\csromanspacer{} ตสฺส เอวมสฺส “อหํ โข ปุพฺเพ พนฺธนาคาเร พทฺโธ อโหสิํ, โสมฺหิ เอตรหิ ตมฺหา พนฺธนาคารา}

\vspace{\tipitakaparskip}
\csromancenter{สามญฺญผลสุตฺต มุตฺโต โสตฺถินา อพฺภเยน, นตฺถิ จ เม กิญฺจิ โภคานํ วโย”ติ.\csromanspacer{} โส ตโตนิทานํ ลเภถ ปาโมชฺชํ, อธิคจฺเฉยฺย โสมนสฺสํ.}
\proseitem{221}{\csromanfolio{69}เสยฺยถาปิ มหาราช ปุริโส ทาโส อสฺส อนตฺตาธีโน ปราธีโน น เยนกามํคโม, โส อปเรน สมเยน ตมฺหา ทาสพฺยา มุจฺเจยฺย อตฺตาธีโน อปราธีโน ภุชิสฺโส เยนกามํคโม.\csromanspacer{} ตสฺส เอวมสฺส “อหํ โข ปุพฺเพ ทาโส อโหสิํ อนตฺตาธีโน ปราธีโน น เยนกามํคโม, โสมฺหิ เอตรหิ ตมฺหา ทาสพฺยา มุตฺโต อตฺตาธีโน อปราธีโน ภุชิสฺโส เยนกามํคโม”ติ.\csromanspacer{} โส ตโตนิทานํ ลเภถ ปาโมชฺชํ, อธิคจฺเฉยฺย โสมนสฺสํ.}

\proseitem{222}{เสยฺยถาปิ มหาราช ปุริโส สธโน สโภโค กนฺตารทฺธานมคฺคํ ปฏิปชฺเชยฺย ทุพฺภิกฺขํ สปฺปฏิภยํ, โส อปเรน สมเยน ตํ กนฺตารํ นิตฺถเรยฺย, โสตฺถินา คามนฺตํ อนุปาปุเณยฺย เขมํ อปฺปฏิภยํ.\csromanspacer{} ตสฺส เอวมสฺส “อหํ โข ปุพฺเพ สธโน สโภโค กนฺตารทฺธานมคฺคํ ปฏิปชฺชิํ ทุพฺภิกฺขํ สปฺปฏิภยํ, โสมฺหิ เอตรหิ ตํ กนฺตารํ นิตฺถิณฺโณ, โสตฺถินา คามนฺตํ อนุปฺปตฺโต เขมํ อปฺปฏิภยนฺ”ติ, โส ตโตนิทานํ ลเภถ ปาโมชฺชํ, อธิคจฺเฉยฺย โสมนสฺสํ.}

\proseitem{223}{เอวเมว โข มหาราช ภิกฺขุ ยถา อิณํ ยถา โรคํ ยถา พนฺธนาคารํ ยถา ทาสพฺยํ ยถา กนฺตารทฺธานมคฺคํ.\csromanspacer{} เอวํ อิเม ปญฺจ นีวรเณ อปฺปหีเน อตฺตนิ สมนุปสฺสติ.}

\proseitem{224}{เสยฺยถาปิ มหาราช ยถา อาณณฺยํ ยถา อาโรคฺยํ ยถา พนฺธนาโมกฺขํ ยถา ภุชิสฺสํ ยถา เขมนฺตภูมิํ.\csromanspacer{} เอวเมว โข มหาราช ภิกฺขุ อิเม ปญฺจ นีวรเณ ปหีเน อตฺตนิ สมนุปสฺสติ.}

\proseitem{225}{ตสฺสิเม ปญฺจ นีวรเณ ปหีเน อตฺตนิ สมนุปสฺสโต ปาโมชฺชํ ชายติ, ปมุทิตสฺส ปีติ ชายติ, ปีติมนสฺส กาโย ปสฺสมฺภติ, ปสฺสทฺธกาโย สุขํ เวเทติ, สุขิโน จิตฺตํ สมาธิยติ.}

\csromanmatikaanchor{mtk.43}
\csromantocmarkref{section}{ปฐมฌาน}{mtk.43}
\bookmark[dest={mtk.43},level=1]{ปฐมฌาน}
\csromanheader{1}{\textbf{ปฐมฌาน}}
\proseitem{226}{โส วิวิจฺเจว กาเมหิ วิวิจฺจ อกุสเลหิ ธมฺเมหิ สวิตกฺกํ สวิจารํ วิเวกชํ ปีติสุขํ ปฐมํ ฌานํ อุปสมฺปชฺช วิหรติ, โส \csromanfolio{70}อิมเมว กายํ วิเวกเชน ปีติสุเขน อภิสนฺเทติ ปริสนฺเทติ ปริปูเรติ ปริปฺผรติ, นาสฺส กิญฺจิ สพฺพาวโต กายสฺส วิเวกเชน ปีติสุเขน อปฺผุฏํ โหติ.}

\proseitem{227}{เสยฺยถาปิ มหาราช ทกฺโข นฺหาปโก วา นฺหาปกนฺเตวาสี วา กํสถาเล นฺหานียจุณฺณานิ อากิริตฺวา อุทเกน ปริปฺโผสกํ ปริปฺโผสกํ สนฺเนยฺย, สายํ นฺหานียปิณฺฑิ สฺเนหานุคตา สฺเนหปเรตา สนฺตรพาหิรา ผุฏา สฺเนเหน, น จ ปคฺฆรณี.\csromanspacer{} เอวเมว โข มหาราช ภิกฺขุ อิมเมว กายํ วิเวกเชน ปีติสุเขน อภิสนฺเทติ ปริสนฺเทติ ปริปูเรติ ปริปฺผรติ, นาสฺส กิญฺจิ สพฺพาวโต กายสฺส วิเวกเชน ปีติสุเขน อปฺผุฏํ โหติ.\csromanspacer{} อิทมฺปิ โข มหาราช สนฺทิฏฺฐิกํ สามญฺญผลํ ปุริเมหิ สนฺทิฏฺฐิเกหิ สามญฺญผเลหิ อภิกฺกนฺตตรญฺจ ปณีตตรญฺจ.}

\csromanmatikaanchor{mtk.44}
\csromantocmarkref{section}{ทุติยฌาน}{mtk.44}
\bookmark[dest={mtk.44},level=1]{ทุติยฌาน}
\csromanheader{1}{\textbf{ทุติยฌาน}}
\proseitem{228}{ปุน จปรํ มหาราช ภิกฺขุ วิตกฺกวิจารานํ วูปสมา อชฺฌตฺตํ สมฺปสาทนํ เจตโส เอโกทิภาวํ อวิตกฺกํ อวิจารํ สมาธิชํ ปีติสุขํ ทุติยํ ฌานํ อุปสมฺปชฺช วิหรติ, โส อิมเมว กายํ สมาธิเชน ปีติสุเขน อภิสนฺเทติ ปริสนฺเทติ ปริปูเรติ ปริปฺผรติ, นาสฺส กิญฺจิ สพฺพาวโต กายสฺส สมาธิเชน ปีติสุเขน อปฺผุฏํ โหติ.}

\proseitem{229}{เสยฺยถาปิ มหาราช อุทกรหโท คมฺภีโร อุพฺภิโททโก\footnote{อุพฺภิโตทโก (สฺยา, กํ, ก)} ตสฺส เนวสฺส ปุรตฺถิมาย ทิสาย อุทกสฺส อายมุขํ, น ทกฺขิณาย ทิสาย อุทกสฺส อายมุขํ, น ปจฺฉิมาย ทิสาย อุทกสฺส อายมุขํ, น อุตฺตราย ทิสาย อุทกสฺส อายมุขํ, เทโว จ น กาเลนกาลํ สมฺมาธารํ อนุปฺปเวจฺเฉยฺย.\csromanspacer{} อถ โข ตมฺหาว อุทกรหทา สีตา วาริธารา อุพฺภิชฺชิตฺวา ตเมว อุทกรหทํ สีเตน วารินา อภิสนฺเทยฺย ปริสนฺเทยฺย ปริปูเรยฺย ปริปฺผเรยฺย, นาสฺส กิญฺจิ สพฺพาวโต อุทกรหทสฺส สีเตน วารินา อปฺผุฏํ อสฺส.\csromanspacer{} เอวเมว โข มหาราช ภิกฺขุ อิมเมว กายํ สมาธิเชน ปีติสุเขน อภิสนฺเทติ ปริสนฺเทติ ปริปูเรติ ปริปฺผรติ, นาสฺส กิญฺจิ สพฺพาวโต กายสฺส สมาธิเชน ปีติสุเขน}

\vspace{\tipitakaparskip}
\csromancenter{สามญฺญผลสุตฺต อปฺผุฏํ โหติ.\csromanspacer{} อิทมฺปิ โข มหาราช สนฺทิฏฺฐิกํ สามญฺญผลํ ปุริเมหิ สนฺทิฏฺฐิเกหิ สามญฺญผเลหิ อภิกฺกนฺตรญฺจ ปณีตตรญฺจ.}
\csromanmatikaanchor{mtk.45}
\csromantocmarkref{section}{ตติยฌาน}{mtk.45}
\bookmark[dest={mtk.45},level=1]{ตติยฌาน}
\csromanheader{1}{\textbf{ตติยฌาน}}
\proseitem{230}{\csromanfolio{71}ปุน จปรํ มหาราช ภิกฺขุ ปีติยา จ วิราคา อุเปกฺขโก จ วิหรติ สโต สมฺปชาโน, สุขญฺจ กาเยน ปฏิสํเวเทติ, ยํ ตํ อริยา อาจิกฺขนฺติ “อุเปกฺขโก สติมา สุขวิหารี”ติ, ตติยํ ฌานํ อุปสมฺปชฺช วิหรติ, โส อิมเมว กายํ นิปฺปีติเกน สุเขน อภิสนฺเทติ ปริสนฺเทติ ปริปูเรติ ปริปฺผรติ, นาสฺส กิญฺจิ สพฺพาวโต กายสฺส นิปฺปีติเกน สุเขน อปฺผุฏํ โหติ.}

\proseitem{231}{เสยฺยถาปิ มหาราช อุปฺปลินิยํ วา ปทุมินิยํ วา ปุณฺฑรีกินิยํ วา อปฺเปกจฺจานิ อุปฺปลานิ วา ปทุมานิ วา ปุณฺฑรีกานิ วา อุทเก ชาตานิ อุทเก สํวฑฺฒานิ อุทกานุคฺคตานิ อนฺโตนิมุคฺคโปสีนิ, ตานิ ยาว จคฺคา ยาว จ มูลา สีเตน วารินา อภิสนฺนานิ ปริสนฺนานิ\footnote{อภิสนฺทานิ ปริสนฺทานิ (ก)} ปริปูรานิ ปริปฺผุฏานิ\footnote{ปริปฺผุฏฺฐานิ (อิ)}, นาสฺส กิญฺจิ สพฺพาวตํ อุปฺปลานํ วา ปทุมานํ วา ปุณฺฑรีกานํ วา สีเตน วารินา อปฺผุฏํ อสฺส.\csromanspacer{} เอวเมว โข มหาราช ภิกฺขุ อิมเมว กายํ นิปฺปีติเกน สุเขน อภิสนฺเทติ ปริสนฺเทติ ปริปูเรติ ปริปฺผรติ, นาสฺส กิญฺจิ สพฺพาวโต กายสฺส นิปฺปีติเกน สุเขน อปฺผุฏํ โหติ.\csromanspacer{} อิทมฺปิ โข มหาราช สนฺทิฏฺฐิกํ สามญฺญผลํ ปุริเมหิ สนฺทิฏฺฐิเกหิ สามญฺญผเลหิ อภิกฺกนฺตตรญฺจ ปณีตตรญฺจ.}

\csromanmatikaanchor{mtk.46}
\csromantocmarkref{section}{จตุตฺถฌาน}{mtk.46}
\bookmark[dest={mtk.46},level=1]{จตุตฺถฌาน}
\csromanheader{1}{\textbf{จตุตฺถฌาน}}
\proseitem{232}{ปุน จปรํ มหาราช ภิกฺขุ สุขสฺส จ ปหานา ทุกฺขสฺส จ ปหานา ปุพฺเพว โสมนสฺสโทมนสฺสานํ อตฺถงฺคมา อทุกฺขมสุขํ อุเปกฺขาสติปาริสุทฺธิํ จตุตฺถํ ฌานํ อุปสมฺปชฺช วิหรติ, โส อิมเมว กายํ ปริสุทฺเธน เจตสา ปริโยทาเตน ผริตฺวา นิสินฺโน โหติ, นาสฺส กิญฺจิ สพฺพาวโต กายสฺส ปริสุทฺเธน เจตสา ปริโยทาเตน อปฺผุฏํ โหติ.}

\proseitem{233}{\csromanfolio{72}เสยฺยถาปิ มหาราช ปุริโส โอทาเตน วตฺเถน สสีสํ ปารุปิตฺวา นิสินฺโน อสฺส, นาสฺส กิญฺจิ สพฺพาวโต กายสฺส โอทาเตน วตฺเถน อปฺผุฏํ อสฺส.\csromanspacer{} เอวเมว โข มหาราช ภิกฺขุ อิมเมว กายํ ปริสุทฺเธน เจตสา ปริโยทาเตน ผริตฺวา นิสินฺโน โหติ, นาสฺส กิญฺจิ สพฺพาวโต กายสฺส ปริสุทฺเธน เจตสา ปริโยทาเตน อปฺผุฏํ โหติ.\csromanspacer{} อิทมฺปิ โข มหาราช สนฺทิฏฺฐิกํ สามญฺญผลํ ปุริเมหิ สนฺทิฏฺฐิเกหิ สามญฺญผเลหิ อภิกฺกนฺตตรญฺจ ปณีตตรญฺจ.}

\csromanmatikaanchor{mtk.47}
\csromantocmarkref{section}{วิปสฺสนาญาณ}{mtk.47}
\bookmark[dest={mtk.47},level=1]{วิปสฺสนาญาณ}
\csromanheader{1}{\textbf{วิปสฺสนาญาณ (1)}}
\proseitem{234}{โส\footnote{ปุน จปรํ มหาราช ภิกฺขุ โส (ก)} เอวํ สมาหิเต จิตฺเต ปริสุทฺเธ ปริโยทาเต อนงฺคเณ วิคตูปกฺกิเลเส มุทุภูเต กมฺมนิเย ฐิเต อาเนญฺชปฺปตฺเต ญาณทสฺสนาย จิตฺตํ อภินีหรติ อภินินฺนาเมติ, โส เอวํ ปชานาติ “อยํ โข เม กาโย รูปี จาตุมหาภูติโก มาตาเปตฺติกสมฺภโว โอทนกุมฺมาสูปจโย อนิจฺจุจฺฉาทนปริมทฺทนเภทนวิทฺธํสนธมฺโม, อิทญฺจ ปน เม วิญฺญาณํ เอตฺถ สิตํ เอตฺถ ปฏิพทฺธนฺ”ติ.}

\proseitem{235}{เสยฺยถาปิ มหาราช มณิ เวฬุริโย สุโภ ชาติมา อฏฺฐํโส สุปริกมฺมกโต อจฺโฉ วิปฺปสนฺโน อนาวิโล สพฺพาการสมฺปนฺโน, ตตฺราสฺส สุตฺตํ อาวุตํ นีลํ วา ปีตํ วา โลหิตํ วา\footnote{ปีตกํ วา โลหิตกํ วา (ก)} โอทาตํ วา ปณฺฑุสุตฺตํ วา, ตเมนํ จกฺขุมา ปุริโส หตฺเถ กริตฺวา ปจฺจเวกฺเขยฺย “อยํ โข มณิ เวฬุริโย สุโภ ชาติมา อฏฺฐํโส สุปริกมฺมกโต อจฺโฉ วิปฺปสนฺโน อนาวิโล สพฺพาการสมฺปนฺโน, ตตฺริทํ สุตฺตํ อาวุตํ นีลํ วา ปีตํ วา โลหิตํ วา โอทาตํ วา ปณฺฑุสุตฺตํ วา”ติ.\csromanspacer{} เอวเมว โข มหาราช ภิกฺขุ เอวํ สมาหิเต จิตฺเต ปริสุทฺเธ ปริโยทาเต อนงฺคเณ วิคตูปกฺกิเลเส มุทุภูเต กมฺมนิเย ฐิเต อาเนญฺชปฺปตฺเต ญาณทสฺสนาย จิตฺตํ อภินีหรติ อภินินฺนาเมติ, โส เอวํ ปชานาติ “อยํ โข เม กาโย รูปี จาตุมหาภูติโก มาตาเปตฺติกสมฺภโว โอทนกุมฺมาสูปจโย อนิจฺจุจฺฉาทนปริมทฺทนเภทนวิทฺธํสนธมฺโม, อิทญฺจ ปน เม วิญฺญาณํ เอตฺถ สิตํ เอตฺถ ปฏิพทฺธนฺ”ติ.\csromanspacer{} อิทมฺปิ โข มหาราช สนฺทิฏฺฐิกํ สามญฺญผลํ ปุริเมหิ สนฺทิฏฺฐิเกหิ สามญฺญผเลหิ อภิกฺกนฺตตรญฺจ ปณีตตรญฺจ.}

\csromanmatikaanchor{mtk.48}
\csromantocmarkref{section}{มโนมยิทฺธิญาณ}{mtk.48}
\bookmark[dest={mtk.48},level=1]{มโนมยิทฺธิญาณ}
\csromanheader{1}{2. สามญฺญผลสุตฺต \textbf{มโนมยิทฺธิญาณ (2)}}
\proseitem{236}{\csromanfolio{73}โส เอวํ สมาหิเต จิตฺเต ปริสุทฺเธ ปริโยทาเต อนงฺคเณ วิคตูปกฺกิเลเส มุทุภูเต กมฺมนิเย ฐิเต อาเนญฺชปฺปตฺเต มโนมยํ กายํ อภินิมฺมานาย จิตฺตํ อภินีหรติ อภินินฺนาเมติ, โส อิมมฺหา กายา อญฺญํ กายํ อภินิมฺมินาติ รูปิํ มโนมยํ สพฺพงฺคปจฺจงฺคิํ อหีนินฺทฺริยํ.}

\proseitem{237}{เสยฺยถาปิ มหาราช ปุริโส มุญฺชมฺหา อีสิกํ ปวาเหยฺย\footnote{ปพฺพาเหยฺย (สฺยา, ก)}.\csromanspacer{} ตสฺส เอวมสฺส “อยํ มุญฺโช, อยํ อีสิกา, อญฺโญ มุญฺโช อญฺญา อีสิกา, มุญฺชมฺหา เตฺวว อีสิกา ปวาฬฺหา”ติ\footnote{ปพฺพาฬฺหาติ (สฺยา, ก)}.\csromanspacer{} เสยฺยถา วา ปน มหาราช ปุริโส อสิํ โกสิยา ปวาเหยฺย.\csromanspacer{} ตสฺส เอวมสฺส “อยํ อสิ, อยํ โกสิ, อญฺโญ อสิ, อญฺญา โกสิ, โกสิยา เตฺวว อสิ ปวาโฬฺห”ติ.\csromanspacer{} เสยฺยถา วา ปน มหาราช ปุริโส อหิํ กรณฺฑา อุทฺธเรยฺย.\csromanspacer{} ตสฺส เอวมสฺส “อยํ อหิ, อยํ กรณฺโฑ.\csromanspacer{} อญฺโญ อหิ, อญฺโญ กรณฺโฑ, กรณฺฑา เตฺวว อหิ อุพฺภโต”ติ\footnote{อุทฺธริโต (สฺยา, กํ)}.\csromanspacer{} เอวเมว โข มหาราช ภิกฺขุ เอวํ สมาหิเต จิตฺเต ปริสุทฺเธ ปริโยทาเต อนงฺคเณ วิคตูปกฺกิเลเส มุทุภูเต กมฺมนิเย ฐิเต อาเนญฺชปฺปตฺเต มโนมยํ กายํ อภินิมฺมานาย จิตฺตํ อภินีหรติ อภินินฺนาเมติ, โส อิมมฺหา กายา อญฺญํ กายํ อภินิมฺมินาติ รูปิํ มโนมยํ สพฺพงฺคปจฺจงฺคิํ อหีนินฺทฺริยํ.\csromanspacer{} อิทมฺปิ โข มหาราช สนฺทิฏฺฐิกํ สามญฺญผลํ ปุริเมหิ สนฺทิฏฺฐิเกหิ สามญฺญผเลหิ อภิกฺกนฺตตรญฺจ ปณีตตรญฺจ.}

\csromanmatikaanchor{mtk.49}
\csromantocmarkref{section}{อิทฺธิวิธญาณ}{mtk.49}
\bookmark[dest={mtk.49},level=1]{อิทฺธิวิธญาณ}
\csromanheader{1}{\textbf{อิทฺธิวิธญาณ (3)}}
\proseitem{238}{โส เอวํ สมาหิเต จิตฺเต ปริสุทฺเธ ปริโยทาเต อนงฺคเณ วิคตูปกฺกิเลเส มุทุภูเต กมฺมนิเย ฐิเต อาเนญฺชปฺปตฺเต อิทฺธิวิธาย จิตฺตํ อภินีหรติ อภินินฺนาเมติ, โส อเนกวิหิตํ อิทฺธิวิธํ ปจฺจนุโภติ, เอโกปิ หุตฺวา พหุธา โหติ, พหุธาปิ หุตฺวา เอโก โหติ, อาวิภาวํ ติโรภาวํ ติโรกุฏฺฏํ ติโรปาการํ ติโรปพฺพตํ อสชฺชมาโน คจฺฉติ เสยฺยถาปิ อากาเส.\csromanspacer{} ปถวิยาปิ อุมฺมุชฺชนิมุชฺชํ กโรติ เสยฺยถาปิ อุทเก.\csromanspacer{} อุทเกปิ อภิชฺชมาเน\footnote{อภิชฺชมาโน (สี, ก)} คจฺฉติ เสยฺยถาปิ \csromanfolio{74}ปถวิยา.\csromanspacer{} อากาเสปิ ปลฺลงฺเกน กมติ เสยฺยถาปิ ปกฺขี สกุโณ.\csromanspacer{} อิเมปิ จนฺทิมสูริเย เอวํมหิทฺธิเก เอวํมหานุภาเว ปาณินา ปรามสติ ปริมชฺชติ.\csromanspacer{} ยาว พฺรหฺมโลกาปิ กาเยน วสํ วตฺเตติ.}

\proseitem{239}{เสยฺยถาปิ มหาราช ทกฺโข กุมฺภกาโร วา กุมฺภการนฺเตวาสี วา สุปริกมฺมกตาย มตฺติกาย ยํ ยเทว ภาชนวิกติํ อากงฺเขยฺย, ตํ ตเทว กเรยฺย อภินิปฺผาเทยฺย.\csromanspacer{} เสยฺยถา วา ปน มหาราช ทกฺโข ทนฺตกาโร วา ทนฺตการนฺเตวาสี วา สุปริกมฺมกตสฺมิํ ทนฺตสฺมิํ ยํ ยเทว ทนฺตวิกติํ อากงฺเขยฺย, ตํ ตเทว กเรยฺย อภินิปฺผาเทยฺย.\csromanspacer{} เสยฺยถา วา ปน มหาราช ทกฺโข สุวณฺณกาโร วา สุวณฺณการนฺเตวาสี วา สุปริกมฺมกตสฺมิํ สุวณฺณสฺมิํ ยํ ยเทว สุวณฺณวิกติํ อากงฺเขยฺย, ตํ ตเทว กเรยฺย อภินิปฺผาเทยฺย.\csromanspacer{} เอวเมว โข มหาราช ภิกฺขุ เอวํ สมาหิเต จิตฺเต ปริสุทฺเธ ปริโยทาเต อนงฺคเณ วิคตูปกฺกิเลเส มุทุภูเต กมฺมนิเย ฐิเต อาเนญฺชปฺปตฺเต อิทฺธิวิธาย จิตฺตํ อภินีหรติ อภินินฺนาเมติ, โส อเนกวิหิตํ อิทฺธิวิธํ ปจฺจนุโภติ, เอโกปิ หุตฺวา พหุธา โหติ, พหุธาปิ หุตฺวา เอโก โหติ, อาวิภาวํ ติโรภาวํ ติโรกุฏฺฏํ ติโรปาการํ ติโรปพฺพตํ อสชฺชมาโน คจฺฉติ เสยฺยถาปิ อากาเส.\csromanspacer{} ปถวิยาปิ อุมฺมุชฺชนิมุชฺชํ กโรติ เสยฺยถาปิ อุทเก.\csromanspacer{} อุทเกปิ อภิชฺชมาเน คจฺฉติ เสยฺยถาปิ ปถวิยา.\csromanspacer{} อากาเสปิ ปลฺลงฺเกน กมติ เสยฺยถาปิ ปกฺขี สกุโณ.\csromanspacer{} อิเมปิ จนฺทิมสูริเย เอวํมหิทฺธิเก เอวํมหานุภาเว ปาณินา ปรามสติ ปริมชฺชติ.\csromanspacer{} ยาว พฺรหฺมโลกาปิ กาเยน วสํ วตฺเตติ.\csromanspacer{} อิทมฺปิ โข มหาราช สนฺทิฏฺฐิกํ สามญฺญผลํ ปุริเมหิ สนฺทิฏฺฐิเกหิ สามญฺญผเลหิ อภิกฺกนฺตตรญฺจ ปณีตตรญฺจ.}

\csromanmatikaanchor{mtk.50}
\csromantocmarkref{section}{ทิพฺพโสตญาณ}{mtk.50}
\bookmark[dest={mtk.50},level=1]{ทิพฺพโสตญาณ}
\csromanheader{1}{\textbf{ทิพฺพโสตญาณ (4)}}
\proseitem{240}{โส เอวํ สมาหิเต จิตฺเต ปริสุทฺเธ ปริโยทาเต อนงฺคเณ วิคตูปกฺกิเลเส มุทุภูเต กมฺมนิเย ฐิเต อาเนญฺชปฺปตฺเต ทิพฺพาย โสตธาตุยา จิตฺตํ อภินีหรติ อภินินฺนาเมติ, โส ทิพฺพาย โสตธาตุยา วิสุทฺธาย อติกฺกนฺตมานุสิกาย อุโภ สทฺเท สุณาติ ทิพฺเพ จ มานุเส จ เย ทูเร สนฺติเก จ.}

\csromanheader{2}{2. สามญฺญผลสุตฺต}
\proseitem{241}{\csromanfolio{75}เสยฺยถาปิ มหาราช ปุริโส อทฺธานมคฺคปฺปฏิปนฺโน, โส สุเณยฺย เภริสทฺทมฺปิ มุทิงฺคสทฺทมฺปิ\footnote{มุติงฺคสทฺทมฺปิ (สี, อิ)} สงฺข\-ปณว\-ทินฺทิม\-สทฺทมฺปิ\footnote{สงฺขปณวเทณฺฑิมสทฺทมฺปิ (สี, อิ), สงฺขสทฺทํปิ ปณวสทฺทํปิ เทนฺทิมสทฺทํปิ (สฺยา, กํ)}.\csromanspacer{} ตสฺส เอวมสฺส “เภริสทฺโท” อิติปิ, “มุทิงฺคสทฺโท” อิติปิ, “สงฺขปณวทินฺทิมสทฺโท” อิติปิ\footnote{สงฺขสทฺโท อิติปิ ปณวสทฺโท อิติปิ เทนฺทิมสทฺโท อิติปิ (สฺยา, กํ)}.\csromanspacer{} เอวเมว โข มหาราช ภิกฺขุ เอวํ สมาหิเต จิตฺเต ปริสุทฺเธ ปริโยทาเต อนงฺคเณ วิคตูปกฺกิเลเส มุทุภูเต กมฺมนิเย ฐิเต อาเนญฺชปฺปตฺเต ทิพฺพาย โสตธาตุยา จิตฺตํ อภินีหรติ อภินินฺนาเมติ, โส ทิพฺพาย โสตธาตุยา วิสุทฺธาย อติกฺกนฺตมานุสิกาย อุโภ สทฺเท สุณาติ ทิพฺเพ จ มานุเส จ เย ทูเร สนฺติเก จ.\csromanspacer{} อิทมฺปิ โข มหาราช สนฺทิฏฺฐิกํ สามญฺญผลํ ปุริเมหิ สนฺทิฏฺฐิเกหิ สามญฺญผเลหิ อภิกฺกนฺตตรญฺจ ปณีตตรญฺจ.}

\csromanmatikaanchor{mtk.51}
\csromantocmarkref{section}{เจโตปริยญาณ}{mtk.51}
\bookmark[dest={mtk.51},level=1]{เจโตปริยญาณ}
\csromanheader{1}{\textbf{เจโตปริยญาณ (5)}}
\proseitem{242}{โส เอวํ สมาหิเต จิตฺเต ปริสุทฺเธ ปริโยทาเต อนงฺคเณ วิคตูปกฺกิเลเส มุทุภูเต กมฺมนิเย ฐิเต อาเนญฺชปฺปตฺเต เจโตปริยญาณาย จิตฺตํ อภินีหรติ อภินินฺนาเมติ, โส ปรสตฺตานํ ปรปุคฺคลานํ เจตสา เจโต ปริจฺจ ปชานาติ, สราคํ วา จิตฺตํ “สราคํ จิตฺตนฺ”ติ ปชานาติ, วีตราคํ วา จิตฺตํ “วีตราคํ จิตฺตนฺ”ติ ปชานาติ, สโทสํ วา จิตฺตํ “สโทสํ จิตฺตนฺ”ติ ปชานาติ, วีตโทสํ วา จิตฺตํ “วีตโทสํ จิตฺตนฺ”ติ ปชานาติ, สโมหํ วา จิตฺตํ “สโมหํ จิตฺตนฺ”ติ ปชานาติ, วีตโมหํ วา จิตฺตํ “วีตโมหํ จิตฺตนฺ”ติ ปชานาติ, สํขิตฺตํ วา จิตฺตํ “สํขิตฺตํ จิตฺตนฺ”ติ ปชานาติ, วิกฺขิตฺตํ วา จิตฺตํ “วิกฺขิตฺตํ จิตฺตนฺ”ติ ปชานาติ, มหคฺคตํ วา จิตฺตํ “มหคฺคตํ จิตฺตนฺ”ติ ปชานาติ, อมหคฺคตํ วา จิตฺตํ “อมหคฺคตํ จิตฺตนฺ”ติ ปชานาติ, ส-อุตฺตรํ วา จิตฺตํ “ส-อุตฺตรํ จิตฺตนฺ”ติ ปชานาติ, อนุตฺตรํ วา จิตฺตํ “อนุตฺตรํ จิตฺตนฺ”ติ ปชานาติ, สมาหิตํ วา จิตฺตํ “สมาหิตํ จิตฺตนฺ”ติ ปชานาติ, อสมาหิตํ วา จิตฺตํ “อสมาหิตํ จิตฺตนฺ”ติ ปชานาติ, วิมุตฺตํ วา จิตฺตํ “วิมุตฺตํ จิตฺตนฺ”ติ ปชานาติ, อวิมุตฺตํ วา จิตฺตํ “อวิมุตฺตํ จิตฺตนฺ”ติ ปชานาติ.}

\proseitem{243}{\csromanfolio{76}เสยฺยถาปิ มหาราช อิตฺถี วา ปุริโส วา ทหโร ยุวา มณฺฑนชาติโก อาทาเส วา ปริสุทฺเธ ปริโยทาเต อจฺเฉ วา อุทกปตฺเต สกํ มุขนิมิตฺตํ ปจฺจเวกฺขมาโน สกณิกํ วา “สกณิกนฺ”ติ ชาเนยฺย, อกณิกํ วา “อกณิกนฺ”ติ ชาเนยฺย.\csromanspacer{} เอวเมว โข มหาราช ภิกฺขุ เอวํ สมาหิเต จิตฺเต ปริสุทฺเธ ปริโยทาเต อนงฺคเณ วิคตูปกฺกิเลเส มุทุภูเต กมฺมนิเย ฐิเต อาเนญฺชปฺปตฺเต เจโตปริยญาณาย จิตฺตํ อภินีหรติ อภินินฺนาเมติ, โส ปรสตฺตานํ ปรปุคฺคลานํ เจตสา เจโต ปริจฺจ ปชานาติ, สราคํ วา จิตฺตํ “สราคํ จิตฺตนฺ”ติ ปชานาติ, วีตราคํ วา จิตฺตํ “วีตราคํ จิตฺตนฺ”ติ ปชานาติ, สโทสํ วา จิตฺตํ “สโทสํ จิตฺตนฺ”ติ ปชานาติ, วีตโทสํ วา จิตฺตํ “วีตโทสํ จิตฺตนฺ”ติ ปชานาติ, สโมหํ วา จิตฺตํ “สโมหํ จิตฺตนฺ”ติ ปชานาติ, วีตโมหํ วา จิตฺตํ “วีตโมหํ จิตฺตนฺ”ติ ปชานาติ, สํขิตฺตํ วา จิตฺตํ “สํขิตฺตํ จิตฺตนฺ”ติ ปชานาติ, วิกฺขิตฺตํ วา จิตฺตํ “วิกฺขิตฺตํ จิตฺตนฺ”ติ ปชานาติ, มหคฺคตํ วา จิตฺตํ “มหคฺคตํ จิตฺตนฺ”ติ ปชานาติ, อมหคฺคตํ วา จิตฺตํ “อมหคฺคตํ จิตฺตนฺ”ติ ปชานาติ, ส-อุตฺตรํ วา จิตฺตํ “ส-อุตฺตรํ จิตฺตนฺ”ติ ปชานาติ, อนุตฺตรํ วา จิตฺตํ “อนุตฺตรํ จิตฺตนฺ”ติ ปชานาติ, สมาหิตํ วา จิตฺตํ “สมาหิตํ จิตฺตนฺ”ติ ปชานาติ, อสมาหิตํ วา จิตฺตํ “อสมาหิตํ จิตฺตนฺ”ติ ปชานาติ, วิมุตฺตํ วา จิตฺตํ “วิมุตฺตํ จิตฺตํนฺ”ติ ปชานาติ, อวิมุตฺตํ วา จิตฺตํ “อวิมุตฺตํ จิตฺตนฺ”ติ ปชานาติ.\csromanspacer{} อิทมฺปิ โข มหาราช สนฺทิฏฺฐิกํ สามญฺญผลํ ปุริเมหิ สนฺทิฏฺฐิเกหิ สามญฺญผเลหิ อภิกฺกนฺตตรญฺจ ปณีตตรญฺจ.}

\csromanmatikaanchor{mtk.52}
\csromantocmarkref{section}{ปุพฺเพนิวาสานุสฺสติญาณ}{mtk.52}
\bookmark[dest={mtk.52},level=1]{ปุพฺเพนิวาสานุสฺสติญาณ}
\csromanheader{1}{ปุพฺเพ\-นิวาสา\-นุสฺสติ\-ญาณ (6)}
\proseitem{244}{โส เอวํ สมาหิเต จิตฺเต ปริสุทฺเธ ปริโยทาเต อนงฺคเณ วิคตูปกฺกิเลเส มุทุภูเต กมฺมนิเย ฐิเต อาเนญฺชปฺปตฺเต ปุพฺเพ\-นิวาสา\-นุสฺสติ\-ญาณาย จิตฺตํ อภินีหรติ อภินินฺนาเมติ, โส อเนกวิหิตํ ปุพฺเพนิวาสํ อนุสฺสรติ, เสยฺยถิทํ, เอกมฺปิ ชาติํ เทฺวปิ ชาติโย ติสฺโสปิ ชาติโย จตสฺโสปิ ชาติโย ปญฺจปิ ชาติโย ทสปิ ชาติโย วีสมฺปิ ชาติโย ติํสมฺปิ ชาติโย จตฺตาลีสมฺปิ ชาติโย ปญฺญาสมฺปิ ชาติโย ชาติสตมฺปิ ชาติสหสฺสมฺปิ ชาติสตสหสฺสมฺปิ อเนเกปิ สํวฏฺฏกปฺเป อเนเกปิ วิวฏฺฏกปฺเป อเนเกปิ สํวฏฺฏวิวฏฺฏกปฺเป, “อมุตฺราสิํ เอวํนาโม}

\vspace{\tipitakaparskip}
\csromancenter{สามญฺญผลสุตฺต เอวํโคตฺโต เอวํวณฺโณ เอวมาหาโร เอวํสุขทุกฺขปฺปฏิสํเวที เอวมายุปริยนฺโต, โส ตโต จุโต อมุตฺร อุทปาทิํ, ตตฺราปาสิํ เอวํนาโม เอวํโคตฺโต เอวํวณฺโณ เอวมาหาโร เอวํสุขทุกฺขปฺปฏิสํเวที เอวมายุปริยนฺโต, โส ตโต จุโต อิธูปปนฺโน”ติ, อิติ สาการํ ส-อุทฺเทสํ อเนกวิหิตํ ปุพฺเพนิวาสํ อนุสฺสรติ.}
\proseitem{245}{\csromanfolio{77}เสยฺยถาปิ มหาราช ปุริโส สกมฺหา คามา อญฺญํ คามํ คจฺเฉยฺย, ตมฺหาปิ คามา อญฺญํ คามํ คจฺเฉยฺย, โส ตมฺหา คามา สกํเยว คามํ ปจฺจาคจฺเฉยฺย.\csromanspacer{} ตสฺส เอวมสฺส “อหํ โข สกมฺหา คามา อมุํ คามํ อคจฺฉิํ\footnote{อคญฺฉิํ (สฺยา, กํ)}, ตตฺร เอวํ อฏฺฐาสิํ, เอวํ นิสีทิํ, เอวํ อภาสิํ, เอวํ ตุณฺหี อโหสิํ, ตมฺหาปิ คามา อมุํ คามํ อคจฺฉิํ, ตตฺราปิ เอวํ อฏฺฐาสิํ, เอวํ นิสีทิํ, เอวํ อภาสิํ, เอวํ ตุณฺหี อโหสิํ, โสมฺหิ ตมฺหา คามา สกํเยว คามํ ปจฺจาคโต”ติ.\csromanspacer{} เอวเมว โข มหาราช ภิกฺขุ เอวํ สมาหิเต จิตฺเต ปริสุทฺเธ ปริโยทาเต อนงฺคเณ วิคตูปกฺกิเลเส มุทุภูเต กมฺมนิเย ฐิเต อาเนญฺชปฺปตฺเต ปุพฺเพ\-นิวาสา\-นุสฺสติ\-ญาณาย จิตฺตํ อภินีหรติ อภินินฺนาเมติ, โส อเนกวิหิตํ ปุพฺเพนิวาสํ อนุสฺสรติ, เสยฺยถิทํ, เอกมฺปิ ชาติํ เทฺวปิ ชาติโย ติสฺโสปิ ชาติโย จตสฺโสปิ ชาติโย ปญฺจปิ ชาติโย ทสปิ ชาติโย วีสมฺปิ ชาติโย ติํสมฺปิ ชาติโย จตฺตาลีสมฺปิ ชาติโย ปญฺญาสมฺปิ ชาติโย ชาติสตมฺปิ ชาติสหสฺสมฺปิ ชาติสตสหสฺสมฺปิ อเนเกปิ สํวฏฺฏกปฺเป อเนเกปิ วิวฏฺฏกปฺเป อเนเกปิ สํวฏฺฏวิวฏฺฏกปฺเป, “อมุตฺราสิํ เอวํนาโม เอวํโคตฺโต เอวํวณฺโณ เอวมาหาโร เอวํสุขทุกฺขปฺปฏิสํเวที เอวมายุปริยนฺโต, โส ตโต จุโต อมุตฺร อุทปาทิํ, ตตฺราปาสิํ เอวํนาโม เอวํโคตฺโต เอวํวณฺโณ เอวมาหาโร เอวํสุขทุกฺขปฺปฏิสํเวที เอวมายุปริยนฺโต, โส ตโต จุโต อิธูปปนฺโน”ติ, อิติ สาการํ ส-อุทฺเทสํ อเนกวิหิตํ ปุพฺเพนิวาสํ อนุสฺสรติ.\csromanspacer{} อิทมฺปิ โข มหาราช สนฺทิฏฺฐิกํ สามญฺญผลํ ปุริเมหิ สนฺทิฏฺฐิเกหิ สามญฺญผเลหิ อภิกฺกนฺตตรญฺจ ปณีตตรญฺจ.}

\csromanmatikaanchor{mtk.53}
\csromantocmarkref{section}{ทิพฺพจกฺขุญาณ}{mtk.53}
\bookmark[dest={mtk.53},level=1]{ทิพฺพจกฺขุญาณ}
\csromanheader{1}{\textbf{ทิพฺพจกฺขุญาณ (7)}}
\proseitem{246}{\csromanfolio{78}โส เอวํ สมาหิเต จิตฺเต ปริสุทฺเธ ปริโยทาเต อนงฺคเณ วิคตูปกฺกิเลเส มุทุภูเต กมฺมนิเย ฐิเต อาเนญฺชปฺปตฺเต สตฺตานํ จุตูปปาตญาณาย จิตฺตํ อภินีหรติ อภินินฺนาเมติ, โส ทิพฺเพน จกฺขุนา วิสุทฺเธน อติกฺกนฺตมานุสเกน สตฺเต ปสฺสติ จวมาเน อุปปชฺชมาเน หีเน ปณีเต สุวณฺเณ ทุพฺพณฺเณ สุคเต ทุคฺคเต, ยถากมฺมูปเค สตฺเต ปชานาติ “อิเม วต โภนฺโต สตฺตา กายทุจฺจริเตน สมนฺนาคตา วจีทุจฺจริเตน สมนฺนาคตา มโนทุจฺจริเตน สมนฺนาคตา อริยานํ อุปวาทกา มิจฺฉาทิฏฺฐิกา มิจฺฉาทิฏฺฐิ\-กมฺม\-สมาทานา, เต กายสฺส เภทา ปรํ มรณา อปายํ ทุคฺคติํ วินิปาตํ นิรยํ อุปปนฺนา.\csromanspacer{} อิเม วา ปน โภนฺโต สตฺตา กายสุจริเตน สมนฺนาคตา วจีสุจริเตน สมนฺนาคตา มโนสุจริเตน สมนฺนาคตา อริยานํ อนุปวาทกา สมฺมาทิฏฺฐิกา สมฺมา\-ทิฏฺฐิ\-กมฺม\-สมาทานา, เต กายสฺส เภทา ปรํ มรณา สุคติํ สคฺคํ โลกํ อุปปนฺนา”ติ, อิติ ทิพฺเพน จกฺขุนา วิสุทฺเธน อติกฺกนฺตมานุสเกน สตฺเต ปสฺสติ จวมาเน อุปปชฺชมาเน หีเน ปณีเต สุวณฺเณ ทุพฺพณฺเณ สุคเต ทุคฺคเต, ยถากมฺมูปเค สตฺเต ปชานาติ.}

\proseitem{247}{เสยฺยถาปิ มหาราช มชฺเฌ สิงฺฆาฏเก ปาสาโท, ตตฺถ จกฺขุมา ปุริโส ฐิโต ปสฺเสยฺย มนุสฺเส เคหํ ปวิสนฺเตปิ นิกฺขมนฺเตปิ รถิกายปิ วีถิํ สญฺจรนฺเต\footnote{รถิยาปิ รถิํ สญฺจรนฺเต (สี), รถิยาย วิถิํ สญฺจรนฺเตปิ (สฺยา)} มชฺเฌ สิงฺฆาฏเก นิสินฺเนปิ.\csromanspacer{} ตสฺส เอวมสฺส “เอเต มนุสฺสา เคหํ ปวิสนฺติ, เอเต นิกฺขมนฺติ, เอเต รถิกาย วีถิํ สญฺจรนฺติ, เอเต มชฺเฌ สิงฺฆาฏเก นิสินฺนา”ติ.\csromanspacer{} เอวเมว โข มหาราช ภิกฺขุ เอวํ สมาหิเต จิตฺเต ปริสุทฺเธ ปริโยทาเต อนงฺคเณ วิคตูปกฺกิเลเส มุทุภูเต กมฺมนิเย ฐิเต อาเนญฺชปฺปตฺเต สตฺตานํ จุตูปปาตญาณาย จิตฺตํ อภินีหรติ อภินินฺนาเมติ, โส ทิพฺเพน จกฺขุนา วิสุทฺเธน อติกฺกนฺตมานุสเกน สตฺเต ปสฺสติ จวมาเน อุปปชฺชมาเน หีเน ปณีเต สุวณฺเณ ทุพฺพณฺเณ สุคเต ทุคฺคเต, ยถากมฺมูปเค สตฺเต ปชานาติ “อิเม วต โภนฺโต สตฺตา กายทุจฺจริเตน สมนฺนาคตา วจีทุจฺจริเตน สมนฺนาคตา มโนทุจฺจริเตน สมนฺนาคตา อริยานํ อุปวาทกา มิจฺฉาทิฏฺฐิกา มิจฺฉาทิฏฺฐิ\-กมฺม\-สมาทานา, เต กายสฺส เภทา ปรํ มรณา อปายํ ทุคฺคติํ วินิปาตํ นิรยํ อุปปนฺนา.}

\vspace{\tipitakaparskip}
\csromancenter{สามญฺญผลสุตฺต อิเม วา ปน โภนฺโต สตฺตา กายสุจริเตน สมนฺนาคตา วจีสุจริเตน สมนฺนาคตา มโนสุจริเตน สมนฺนาคตา อริยานํ อนุปวาทกา สมฺมาทิฏฺฐิกา สมฺมา\-ทิฏฺฐิ\-กมฺม\-สมาทานา, เต กายสฺส เภทา ปรํ มรณา สุคติํ สคฺคํ โลกํ อุปปนฺนา”ติ, อิติ ทิพฺเพน จกฺขุนา วิสุทฺเธน อติกฺกนฺตมานุสเกน สตฺเต ปสฺสติ จวมาเน อุปปชฺชมาเน หีเน ปณีเต สุวณฺเณ ทุพฺพณฺเณ สุคเต ทุคฺคเต, ยถากมฺมูปเค สตฺเต ปชานาติ.\csromanspacer{} อิทมฺปิ โข มหาราช สนฺทิฏฺฐิกํ สามญฺญผลํ ปุริเมหิ สนฺทิฏฺฐิเกหิ สามญฺญผเลหิ อภิกฺกนฺตตรญฺจ ปณีตตรญฺจ.}
\csromanmatikaanchor{mtk.54}
\csromantocmarkref{section}{อาสวกฺขยญาณ}{mtk.54}
\bookmark[dest={mtk.54},level=1]{อาสวกฺขยญาณ}
\csromanheader{1}{\textbf{อาสวกฺขยญาณ (8)}}
\proseitem{248}{\csromanfolio{79}โส เอวํ สมาหิเต จิตฺเต ปริสุทฺเธ ปริโยทาเต อนงฺคเณ วิคตูปกฺกิเลเส มุทุภูเต กมฺมนิเย ฐิเต อาเนญฺชปฺปตฺเต อาสวานํ ขยญาณาย จิตฺตํ อภินีหรติ อภินินฺนาเมติ, โส อิทํ ทุกฺขนฺติ ยถาภูตํ ปชานาติ, อยํ ทุกฺขสมุทโยติ ยถาภูตํ ปชานาติ, อยํ ทุกฺขนิโรโธติ ยถาภูตํ ปชานาติ, อยํ ทุกฺขนิโรธคามินี ปฏิปทาติ ยถาภูตํ ปชานาติ.\csromanspacer{} อิเม อาสวาติ ยถาภูตํ ปชานาติ, อยํ อาสวสมุทโยติ ยถาภูตํ ปชานาติ, อยํ อาสวนิโรโธติ ยถาภูตํ ปชานาติ, อยํ อาสวนิโรธคามินี ปฏิปทาติ ยถาภูตํ ปชานาติ.\csromanspacer{} ตสฺส เอวํ ชานโต เอวํ ปสฺสโต กามาสวาปิ จิตฺตํ วิมุจฺจติ, ภวาสวาปิ จิตฺตํ วิมุจฺจติ, อวิชฺชาสวาปิ จิตฺตํ วิมุจฺจติ, วิมุตฺตสฺมิํ วิมุตฺตมิติ ญาณํ โหติ, “ขีณา ชาติ, วุสิตํ พฺรหฺมจริยํ, กตํ กรณียํ, นาปรํ อิตฺถตฺตายา”ติ ปชานาติ.}

\proseitem{249}{เสยฺยถาปิ มหาราช ปพฺพตสงฺเขเป อุทกรหโท อจฺโฉ วิปฺปสนฺโน อนาวิโล, ตตฺถ จกฺขุมา ปุริโส ตีเร ฐิโต ปสฺเสยฺย สิปฺปิกสมฺพุกมฺปิ สกฺขรกถลมฺปิ มจฺฉคุมฺพมฺปิ จรนฺตมฺปิ ติฏฺฐนฺตมฺปิ.\csromanspacer{} ตสฺส เอวมสฺส “อยํ โข อุทกรหโท อจฺโฉ วิปฺปสนฺโน อนาวิโล, ตตฺริเม สิปฺปิกสมฺพุกาปิ สกฺขรกถลาปิ มจฺฉคุมฺพาปิ จรนฺติปิ ติฏฺฐนฺติปี”ติ.\csromanspacer{} เอวเมว โข มหาราช ภิกฺขุ เอวํ สเมหิเต จิตฺเต ปริสุทฺเธ ปริโยทาเต อนงฺคเณ วิคตูปกฺกิเลเส มุทุภูเต กมฺมนิเย ฐิเต อาเนญฺชปฺปตฺเต อาสวานํ ขยญาณาย จิตฺตํ อภินีหรติ อภินินฺนาเมติ, โส อิทํ ทุกฺขนฺติ ยถาภูตํ ปชานาติ, อยํ ทุกฺขสมุทโยติ ยถาภูตํ \csromanfolio{80}ปชานาติ, อยํ ทุกฺขนิโรโธติ ยถาภูตํ ปชานาติ, อยํ ทุกฺขนิโรธคามินี ปฏิปทาติ ยถาภูตํ ปชานาติ.\csromanspacer{} อิเม อาสวาติ ยถาภูตํ ปชานาติ, อยํ อาสวสมุทโยติ ยถาภูตํ ปชานาติ, อยํ อาสวนิโรโธติ ยถาภูตํ ปชานาติ, อยํ อาสวนิโรธคามินี ปฏิปทาติ ยถาภูตํ ปชานาติ.\csromanspacer{} ตสฺส เอวํ ชานโต เอวํ ปสฺสโต กามาสวาปิ จิตฺตํ วิมุจฺจติ, ภวาสวาปิ จิตฺตํ วิมุจฺจติ, อวิชฺชาสวาปิ จิตฺตํ วิมุจฺจติ, วิมุตฺตสฺมิํ วิมุตฺตมิติ ญาณํ โหติ, “ขีณา ชาติ, วุสิตํ พฺรหฺมจริยํ, กตํ กรณียํ, นาปรํ อิตฺถตฺตายา”ติ ปชานาติ.\csromanspacer{} อิทํ โข มหาราช สนฺทิฏฺฐิกํ สามญฺญผลํ ปุริเมหิ สนฺทิฏฺฐิเกหิ สามญฺญผเลหิ อภิกฺกนฺตตรญฺจ ปณีตตรญฺจ.\csromanspacer{} อิมสฺมา จ ปน มหาราช สนฺทิฏฺฐิกา สามญฺญผลา อญฺญํ สนฺทิฏฺฐิกํ สามญฺญผลํ อุตฺตริตรํ วา ปณีตตรํ วา นตฺถี”ติ.}

\csromanmatikaanchor{mtk.55}
\csromantocmarkref{section}{อชาตสตฺตุอุปาสกตฺตปฏิเวทนา}{mtk.55}
\bookmark[dest={mtk.55},level=1]{อชาตสตฺตุอุปาสกตฺตปฏิเวทนา}
\csromanheader{1}{อชาตสตฺตุ\-อุปาสกตฺต\-ปฏิเวทนา}
\proseitem{250}{เอวํ วุตฺเต ราชา มาคโธ อชาตสตฺตุ เวเทหิปุตฺโต ภควนฺตํ เอตทโวจ “อภิกฺกนฺตํ ภนฺเต อภิกฺกนฺตํ ภนฺเต, เสยฺยถาปิ ภนฺเต นิกฺกุชฺชิตํ วา อุกฺกุชฺเชยฺย, ปฏิจฺฉนฺนํ วา วิวเรยฺย, มูฬฺหสฺส วา มคฺคํ อาจิกฺเขยฺย, อนฺธกาเร วา เตลปชฺโชตํ ธาเรยฺย ‘จกฺขุมนฺโต รูปานิ ทกฺขนฺตี’ติ, เอวเมวํ ภนฺเต ภควตา อเนกปริยาเยน ธมฺโม ปกาสิโต.\csromanspacer{} เอสาหํ ภนฺเต ภควนฺตํ สรณํ คจฺฉามิ ธมฺมญฺจ ภิกฺขุสํฆญฺจ, อุปาสกํ มํ ภควา ธาเรตุ อชฺชตคฺเค ปาณุเปตํ สรณํ คตํ, อจฺจโย มํ ภนฺเต อจฺจคมา ยถาพาลํ ยถามูฬฺหํ ยถา-อกุสลํ, โยหํ ปิตรํ ธมฺมิกํ ธมฺมราชานํ อิสฺสริยการณา ชีวิตา โวโรเปสิํ, ตสฺส เม ภนฺเต ภควา อจฺจยํ อจฺจยโต ปฏิคฺคณฺหาตุ อายติํ สํวรายา”ติ.}

\proseitem{251}{ตคฺฆ ตฺวํ มหาราช อจฺจโย อจฺจคมา ยถาพาลํ ยถามูฬฺหํ ยถา-อกุสลํ, ยํ ตฺวํ ปิตรํ ธมฺมิกํ ธมฺมราชานํ ชีวิตา โวโรเปสิ, ยโต จ โข ตฺวํ มหาราช อจฺจยํ อจฺจยโต ทิสฺวา ยถาธมฺมํ ปฏิกโรสิ, ตํ เต มยํ ปฏิคฺคณฺหาม, วุทฺธิเหสา มหาราช อริยสฺส วินเย, โย อจฺจยํ อจฺจยโต ทิสฺวา ยถาธมฺมํ ปฏิกโรติ, อายติํ สํวรํ อาปชฺชตีติ.}

\csromanheader{2}{2. สามญฺญผลสุตฺต}
\proseitem{252}{\csromanfolio{81}เอวํ วุตฺเต ราชา มาคโธ อชาตสตฺตุ เวเทหิปุตฺโต ภควนฺตํ เอตทโวจ “หนฺท จ ทานิ มยํ ภนฺเต คจฺฉาม พหุกิจฺจา มยํ พหุกรณียา”ติ.\csromanspacer{} ยสฺสทานิ ตฺวํ มหาราช กาลํ มญฺญสีติ.\csromanspacer{} อถ โข ราชา มาคโธ อชาตสตฺตุ เวเทหิปุตฺโต ภควโต ภาสิตํ อภินนฺทิตฺวา อนุโมทิตฺวา อุฏฺฐายาสนา ภควนฺตํ อภิวาเทตฺวา ปทกฺขิณํ กตฺวา ปกฺกามิ.}

\proseitemwithcloser{253}{อถ โข ภควา อจิรปกฺกนฺตสฺส รญฺโญ มาคธสฺส อชาตสตฺตุสฺส เวเทหิปุตฺตสฺส ภิกฺขู อามนฺเตสิ “ขตายํ ภิกฺขเว ราชา, อุปหตายํ ภิกฺขเว ราชา, สจายํ ภิกฺขเว ราชา ปิตรํ ธมฺมิกํ ธมฺมราชานํ ชีวิตา น โวโรเปสฺสถ, อิมสฺมิํเยว อาสเน วิรชํ วีตมลํ ธมฺมจกฺขุํ อุปฺปชฺชิสฺสถา”ติ.\csromanspacer{} อิทมโวจ ภควา.\csromanspacer{} อตฺตมนา เต ภิกฺขู ภควโต ภาสิตํ อภินนฺทุนฺติ.}{สามญฺญผลสุตฺตํ นิฏฺฐิตํ ทุติยํ.}
\csromanmatikaanchor{mtk.56}
\csromantocmarknopage{chapter}{3. อมฺพฏฺฐสุตฺต}{mtk.56}
\bookmark[dest={mtk.56},level=0]{3. อมฺพฏฺฐสุตฺต}
\setcsromanheadcha{3. อมฺพฏฺฐสุตฺต}
\setcsromanheadhclear{}
\chapterheadpageread{3. \textbf{อมฺพฏฺฐสุตฺต}}
\proseitem{254}{\csromanfolio{82}เอวํ เม สุตํ\csromandash{}เอกํ สมยํ ภควา โกสเลสุ จาริกํ จรมาโน มหตา ภิกฺขุสํเฆน สทฺธิํ ปญฺจมตฺเตหิ ภิกฺขุสเตหิ เยน อิจฺฉานงฺคลํ นาม โกสลานํ พฺราหฺมณคาโม ตทวสริ, ตตฺร สุทํ ภควา อิจฺฉานงฺคเล วิหรติ อิจฺฉานงฺคลวนสณฺเฑ.}

\csromanmatikaanchor{mtk.57}
\csromantocmarkref{section}{โปกฺขรสาติวตฺถุ}{mtk.57}
\bookmark[dest={mtk.57},level=1]{โปกฺขรสาติวตฺถุ}
\csromanheader{1}{\textbf{โปกฺขรสาติวตฺถุ}}
\proseitem{255}{เตน โข ปน สมเยน พฺราหฺมโณ โปกฺขรสาติ\footnote{โปกฺขรสาตี (สี), โปกฺขรสาทิ (อิ)} อุกฺกฏฺฐํ อชฺฌาวสติ สตฺตุสฺสทํ สติณกฏฺโฐทกํ สธญฺญํ ราชโภคฺคํ รญฺญา ปเสนทินา โกสเลน ทินฺนํ ราชทายํ พฺรหฺมเทยฺยํ.\csromanspacer{} อสฺโสสิ โข พฺราหฺมโณ โปกฺขรสาติ “สมโณ ขลุ โภ โคตโม สกฺยปุตฺโต สกฺยกุลา ปพฺพชิโต โกสเลสุ จาริกํ จรมาโน มหตา ภิกฺขุสํเฆน สทฺธิํ ปญฺจมตฺเตหิ ภิกฺขุสเตหิ อิจฺฉานงฺคลํ อนุปฺปตฺโต อิจฺฉานงฺคเล วิหรติ อิจฺฉานงฺคลวนสณฺเฑ, ตํ โข ปน ภวนฺตํ โคตมํ เอวํ กลฺยาโณ กิตฺติสทฺโท อพฺภุคฺคโต ‘อิติปิ โส ภควา อรหํ สมฺมาสมฺพุทฺโธ วิชฺชาจรณสมฺปนฺโน สุคโต โลกวิทู อนุตฺตโร ปุริสทมฺมสารถิ สตฺถา เทวมนุสฺสานํ พุทฺโธ ภควา’\footnote{ภควาติ (สฺยา, กํ), อุปริโสณทณฺฑสุตฺตาทีสุปิ พุทฺธคุณกถายํ เอวเมว ทิสฺสติ.}, โส อิมํ โลกํ สเทวกํ สมารกํ สพฺรหฺมกํ สสฺสมณพฺราหฺมณิํ ปชํ สเทวมนุสฺสํ สยํ อภิญฺญา สจฺฉิกตฺวา ปเวเทติ, โส ธมฺมํ เทเสติ อาทิกลฺยาณํ มชฺเฌกลฺยาณํ ปริโยสานกลฺยาณํ, สาตฺถํ สพฺยญฺชนํ เกวลปริปุณฺณํ ปริสุทฺธํ พฺรหฺมจริยํ ปกาเสติ, สาธุ โข ปน ตถารูปานํ อรหตํ ทสฺสนํ โหตี”ติ.}

\csromanmatikaanchor{mtk.58}
\csromantocmarkref{section}{อมฺพฏฺฐมาณว}{mtk.58}
\bookmark[dest={mtk.58},level=1]{อมฺพฏฺฐมาณว}
\csromanheader{1}{\textbf{อมฺพฏฺฐมาณว}}
\proseitem{256}{เตน โข ปน สมเยน พฺราหฺมณสฺส โปกฺขรสาติสฺส อมฺพฏฺโฐ นาม มาณโว อนฺเตวาสี โหติ อชฺฌายโก มนฺตธโร ติณฺณํ เวทานํ\footnote{เพทานํ (ก)} ปารคู สนิฆณฺฑุเกฏุภานํ สากฺขรปฺปเภทานํ อิติหาสปญฺจมานํ ปทโก เวยฺยากรโณ โลกายตมหาปุริสลกฺขเณสุ อนวโย อนุญฺญาตปฏิญฺญาโต สเก อาจริยเก เตวิชฺชเก}

\vspace{\tipitakaparskip}
\csromancenter{อมฺพฏฺฐสุตฺต ปาวจเน “ยมหํ ชานามิ, ตํ ตฺวํ ชานาสิ.\csromanspacer{} ยํ ตฺวํ ชานาสิ, ตมหํ ชานามี”ติ.}
\proseitem{257}{\csromanfolio{83}อถ โข พฺราหฺมโณ โปกฺขรสาติ อมฺพฏฺฐํ มาณวํ อามนฺเตสิ “อยํ ตาต อมฺพฏฺฐ สมโณ โคตโม สกฺยปุตฺโต สกฺยกุลา ปพฺพชิโต โกสเลสุ จาริกํ จรมาโน มหตา ภิกฺขุสํเฆน สทฺธิํ ปญฺจมตฺเตหิ ภิกฺขุสเตหิ อิจฺฉานงฺคลํ อนุปฺปตฺโต อิจฺฉานงฺคเล วิหรติ อิจฺฉานงฺคลวนสณฺเฑ, ตํ โข ปน ภวนฺตํ โคตมํ เอวํ กลฺยาโณ กิตฺติสทฺโท อพฺภุคฺคโต ‘อิติปิ โส ภควา อรหํ สมฺมาสมฺพุทฺโธ วิชฺชาจรณสมฺปนฺโน สุคโต โลกวิทู อนุตฺตโร ปุริสทมฺมสารถิ สตฺถา เทวมนุสฺสานํ พุทฺโธ ภควา’, โส อิมํ โลกํ สเทวกํ สมารกํ สพฺรหฺมกํ สสฺสมณพฺราหฺมณิํ ปชํ สเทวมนุสฺสํ สยํ อภิญฺญา สจฺฉิกตฺวา ปเวเทติ, โส ธมฺมํ เทเสติ อาทิกลฺยาณํ มชฺเฌกลฺยาณํ ปริโยสานกลฺยาณํ, สาตฺถํ สพฺยญฺชนํ เกวลปริปุณฺณํ ปริสุทฺธํ พฺรหฺมจริยํ ปกาเสติ, สาธุ โข ปน ตถารูปานํ อรหตํ ทสฺสนํ โหตี”ติ.\csromanspacer{} เอหิ ตฺวํ ตาต อมฺพฏฺฐ เยน สมโณ โคตโม เตนุปสงฺกม, อุปสงฺกมิตฺวา สมณํ โคตมํ ชานาหิ, ยทิ วา ตํ ภวนฺตํ โคตมํ ตถาสนฺตํเยว สทฺโท อพฺภุคฺคโต, ยทิ วา โน ตถา.\csromanspacer{} ยทิ วา โส ภวํ โคตโม ตาทิโส, ยทิ วา น ตาทิโส, ตถา มยํ ตํ ภวนฺตํ โคตมํ เวทิสฺสามา”ติ.}

\proseitem{258}{ยถา กถํ ปนาหํ โภ ตํ ภวนฺตํ โคตมํ ชานิสฺสามิ, “ยทิ วา ตํ ภวนฺตํ โคตมํ ตถาสนฺตํเยว สทฺโท อพฺภุคฺคโต, ยทิ วา โน ตถา.\csromanspacer{} ยทิ วา โส ภวํ โคตโม ตาทิโส, ยทิ วา น ตาทิโส”ติ.}

\prose{อาคตานิ โข ตาต อมฺพฏฺฐ อมฺหากํ มนฺเตสุ ทฺวตฺติํส มหาปุริสลกฺขณานิ, เยหิ สมนฺนาคตสฺส มหาปุริสสฺส เทฺวเยว คติโย ภวนฺติ อนญฺญา.\csromanspacer{} สเจ อคารํ อชฺฌาวสติ, ราชา โหติ จกฺกวตฺตี ธมฺมิโก ธมฺมราชา จาตุรนฺโต วิชิตาวี ชนปทตฺถาวริยปฺปตฺโต สตฺตรตนสมนฺนาคโต.\csromanspacer{} ตสฺสิมานิ สตฺต รตนานิ ภวนฺติ.\csromanspacer{} เสยฺยถิทํ, จกฺกรตนํ, หตฺถิรตนํ, อสฺสรตนํ, มณิรตนํ, อิตฺถิรตนํ, คหปติรตนํ, ปริณายกรตนเมว สตฺตมํ.\csromanspacer{} ปโรสหสฺสํ โข ปนสฺส ปุตฺตา \csromanfolio{84}ภวนฺติ สูรา วีรงฺครูปา ปรเสนปฺปมทฺทนา.\csromanspacer{} โส อิมํ ปถวิํ สาครปริยนฺตํ อทณฺเฑน อสตฺเถน ธมฺเมน อภิวิชิย อชฺฌาวสติ.\csromanspacer{} สเจ โข ปน อคารสฺมา อนคาริยํ ปพฺพชติ, อรหํ โหติ สมฺมาสมฺพุทฺโธ โลเก วิวฏฺฏจฺฉโท.\csromanspacer{} อหํ โข ปน ตาต อมฺพฏฺฐ มนฺตานํ ทาตา, ตฺวํ มนฺตานํ ปฏิคฺคเหตาติ.}

\proseitem{259}{“เอวํ โภ”ติ โข อมฺพฏฺโฐ มาณโว พฺราหฺมณสฺส โปกฺขรสาติสฺส ปฏิสฺสุตฺวา อุฏฺฐายาสนา พฺราหฺมณํ โปกฺขรสาติํ อภิวาเทตฺวา ปทกฺขิณํ กตฺวา วฬวารถมารุยฺห สมฺพหุเลหิ มาณวเกหิ สทฺธิํ เยน อิจฺฉานงฺคลวนสณฺโฑ เตน ปายาสิ, ยาวติกา ยานสฺส ภูมิ ยาเนน คนฺตฺวา ยานา ปจฺโจโรหิตฺวา ปตฺติโกว อารามํ ปาวิสิ.\csromanspacer{} เตน โข ปน สมเยน สมฺพหุลา ภิกฺขู อพฺโภกาเส จงฺกมนฺติ.\csromanspacer{} อถ โข อมฺพฏฺโฐ มาณโว เยน เต ภิกฺขู เตนุปสงฺกมิ, อุปสงฺกมิตฺวา เต ภิกฺขู เอตทโวจ “กหํ นุ โข โภ เอตรหิ โส ภวํ โคตโม วิหรติ, ตํ หิ มยํ ภวนฺตํ โคตมํ ทสฺสนาย อิธูปสงฺกนฺตา”ติ.}

\proseitem{260}{อถ โข เตสํ ภิกฺขูนํ เอตทโหสิ “อยํ โข อมฺพฏฺโฐ มาณโว อภิญฺญาตโกลญฺโญ เจว อภิญฺญาตสฺส จ พฺราหฺมณสฺส โปกฺขรสาติสฺส อนฺเตวาสี, อครุ โข ปน ภควโต เอวรูเปหิ กุลปุตฺเตหิ สทฺธิํ กถาสลฺลาโป โหตี”ติ.\csromanspacer{} เต อมฺพฏฺฐํ มาณวํ เอตทโวจุํ “เอโส อมฺพฏฺฐ วิหาโร สํวุตทฺวาโร, เตน อปฺปสทฺโท อุปสงฺกมิตฺวา อตรมาโน อาฬินฺทํ ปวิสิตฺวา อุกฺกาสิตฺวา อคฺคฬํ อาโกเฏหิ, วิวริสฺสติ เต ภควา ทฺวารนฺ”ติ.}

\proseitem{261}{อถ โข อมฺพฏฺโฐ มาณโว เยน โส วิหาโร สํวุตทฺวาโร, เตน อปฺปสทฺโท อุปสงฺกมิตฺวา อตรมาโน อาฬินฺทํ ปวิสิตฺวา อุกฺกาสิตฺวา อคฺคฬํ อาโกเฏสิ, วิวริ ภควา ทฺวารํ.\csromanspacer{} ปาวิสิ อมฺพฏฺโฐ มาณโว.\csromanspacer{} มาณวกาปิ ปวิสิตฺวา ภควตา สทฺธิํ สมฺโมทิํสุ, สมฺโมทนียํ กถํ สารณียํ วีติสาเรตฺวา เอกมนฺตํ นิสีทิํสุ.\csromanspacer{} อมฺพฏฺโฐ ปน มาณโว จงฺกมนฺโตปิ นิสินฺเนน ภควตา กญฺจิ กญฺจิ\footnote{กิญฺจิ กิญฺจิ (ก)}}

\vspace{\tipitakaparskip}
\csromancenter{อมฺพฏฺฐสุตฺต กถํ สารณียํ วีติสาเรติ, ฐิโตปิ นิสินฺเนน ภควตา กิญฺจิ กิญฺจิ กถํ สารณียํ วีติสาเรติ.}
\proseitem{262}{\csromanfolio{85}อถ โข ภควา อมฺพฏฺฐํ มาณวํ เอตทโวจ “เอวํ นุ เต อมฺพฏฺฐ พฺราหฺมเณหิ วุทฺเธหิ มหลฺลเกหิ อาจริยปาจริเยหิ สทฺธิํ กถาสลฺลาโป โหติ, ยถยิทํ จรํ ติฏฺฐํ นิสินฺเนน มยา กิญฺจิ กิญฺจิ กถํ สารณียํ วีติสาเรตี”ติ.}

\csromanmatikaanchor{mtk.59}
\csromantocmarkref{section}{ปฐมอิพฺภวาท}{mtk.59}
\bookmark[dest={mtk.59},level=1]{ปฐมอิพฺภวาท}
\csromanheader{1}{\textbf{ปฐมอิพฺภวาท}}
\proseitem{263}{โนหิทํ โภ โคตม.\csromanspacer{} คจฺฉนฺโต วา หิ โภ โคตม คจฺฉนฺเตน พฺราหฺมโณ พฺราหฺมเณน สทฺธิํ สลฺลปิตุมรหติ, ฐิโต วา หิ โภ โคตม ฐิเตน พฺราหฺมโณ พฺราหฺมเณน สทฺธิํ สลฺลปิตุมรหติ, นิสินฺโน วา หิ โภ โคตม นิสินฺเนน พฺราหฺมโณ พฺราหฺมเณน สทฺธิํ สลฺลปิตุมรหติ, สยาโน วา หิ โภ โคตม สยาเนน พฺราหฺมโณ พฺราหฺมเณน สทฺธิํ สลฺลปิตุมรหติ, เย จ โข เต โภ โคตม มุณฺฑกา สมณกา อิพฺภา กณฺหา\footnote{กิณฺหา (กสี, อิ)} พนฺธุปาทาปจฺจา, เตหิปิ เม สทฺธิํ เอวํ กถาสลฺลาโป โหติ, ยถริว โภตา โคตเมนาติ.\csromanspacer{} อตฺถิกวโต โข ปน เต อมฺพฏฺฐ อิธาคมนํ อโหสิ, ยาเยว โข ปนตฺถาย อาคจฺเฉยฺยาถ\footnote{อาคจฺเฉยฺยาโถ (สี, อิ)}, ตเมว อตฺถํ สาธุกํ มนสิ กเรยฺยาถ\footnote{มนสิกเรยฺยาโถ (สี, อิ)}, อวุสิตวาเยว โข ปน โภ อยํ อมฺพฏฺโฐ มาณโว วุสิตมานี กิมญฺญตฺร อวุสิตตฺตาติ.}

\proseitem{264}{อถ โข อมฺพฏฺโฐ มาณโว ภควตา อวุสิตวาเทน วุจฺจมาโน กุปิโต อนตฺตมโน ภควนฺตํเยว ขุํเสนฺโต ภควนฺตํเยว วมฺเภนฺโต ภควนฺตํเยว อุปวทมาโน “สมโณ จ เม โภ โคตโม ปาปิโต ภวิสฺสตี”ติ ภควนฺตํ เอตทโวจ “จณฺฑา โภ โคตม สกฺยชาติ, ผรุสา โภ โคตม สกฺยชาติ, ลหุสา โภ โคตม สกฺยชาติ, ภสฺสา โภ โคตม สกฺยชาติ, อิพฺภา สนฺตา อิพฺภา สมานา น พฺราหฺมเณ สกฺกโรนฺติ, น พฺราหฺมเณ ครุํ กโรนฺติ\footnote{ครุกโรนฺติ (สี, สฺยา, กํ, อิ)}, น \csromanfolio{86}พฺราหฺมเณ มาเนนฺติ, น พฺราหฺมเณ ปูเชนฺติ, น พฺราหฺมเณ อปจายนฺติ.\csromanspacer{} ตยิทํ โภ โคตม นจฺฉนฺนํ, ตยิทํ นปฺปติรูปํ, ยทิเม สกฺยา อิพฺภา สนฺตา อิพฺภา สมานา น พฺราหฺมเณ สกฺกโรนฺติ, น พฺราหฺมเณ ครุํ กโรนฺติ, น พฺราหฺมเณ มาเนนฺติ, น พฺราหฺมเณ ปูเชนฺติ, น พฺราหฺมเณ อปจายนฺตีติ.\csromanspacer{} อิติห อมฺพฏฺโฐ มาณโว อิทํ ปฐมํ สเกฺยสุ อิพฺภวาทํ นิปาเตสิ.}

\csromanmatikaanchor{mtk.60}
\csromantocmarkref{section}{ทุติยอิพฺภวาท}{mtk.60}
\bookmark[dest={mtk.60},level=1]{ทุติยอิพฺภวาท}
\csromanheader{1}{\textbf{ทุติยอิพฺภวาท}}
\proseitem{265}{กิํ ปน เต อมฺพฏฺฐ สกฺยา อปรทฺธุนฺติ.\csromanspacer{} เอกมิทาหํ โภ โคตม สมยํ อาจริยสฺส พฺราหฺมณสฺส โปกฺขรสาติสฺส เกนจิเทว กรณีเยน กปิลวตฺถุํ อคมาสิํ, เยน สกฺยานํ สนฺธาคารํ เตนุปสงฺกมิํ, เตน โข ปน สมเยน สมฺพหุลา สกฺยา เจว สกฺยกุมารา จ สนฺธาคาเร\footnote{สนฺถาคาเร (สี, อิ)} อุจฺเจสุ อาสเนสุ นิสินฺนา โหนฺติ อญฺญมญฺญํ องฺคุลิปโตทเกหิ\footnote{องฺคุลิปโตทเกน (อิ)} สญฺชคฺฆนฺตา สํกีฬนฺตา, อญฺญทตฺถุ มมญฺเญว มญฺเญ อนุชคฺฆนฺตา, น มํ โกจิ อาสเนนปิ นิมนฺเตสิ.\csromanspacer{} ตยิทํ โภ โคตม นจฺฉนฺนํ, ตยิทํ นปฺปติรูปํ, ยทิเม สกฺยา อิพฺภา สนฺตา อิพฺภา สมานา น พฺราหฺมเณ สกฺกโรนฺติ, น พฺราหฺมเณ ครุํ กโรนฺติ, น พฺราหฺมเณ มาเนนฺติ, น พฺราหฺมเณ ปูเชนฺติ, น พฺราหฺมเณ อปจายนฺตีติ.\csromanspacer{} อิติห อมฺพฏฺโฐ มาณโว อิทํ ทุติยํ สเกฺยสุ อิพฺภวาทํ นิปาเตสิ.}

\csromanmatikaanchor{mtk.61}
\csromantocmarkref{section}{ตติยอิพฺภวาท}{mtk.61}
\bookmark[dest={mtk.61},level=1]{ตติยอิพฺภวาท}
\csromanheader{1}{\textbf{ตติยอิพฺภวาท}}
\proseitem{266}{ลฏุกิกาปิ โข อมฺพฏฺฐ สกุณิกา สเก กุลาวเก กามลาปินี โหติ.\csromanspacer{} สกํ โข ปเนตํ อมฺพฏฺฐ สกฺยานํ ยทิทํ กปิลวตฺถุํ, นารหตายสฺมา อมฺพฏฺโฐ อิมาย อปฺปมตฺตาย อภิสชฺชิตุนฺติ.\csromanspacer{} จตฺตาโรเม โภ โคตม วณฺณา ขตฺติยา พฺราหฺมณา เวสฺสา สุทฺทา, อิเมสํ หิ โภ โคตม จตุนฺนํ วณฺณานํ ตโย วณฺณา ขตฺติยา จ เวสฺสา จ สุทฺทา จ, อญฺญทตฺถุ พฺราหฺมณสฺเสว ปริจารกา สมฺปชฺชนฺติ.\csromanspacer{} ตยิทํ โภ โคตม นจฺฉนฺนํ, ตยิทํ นปฺปติรูปํ, ยทิเม สกฺยา อิพฺภา สนฺตา อิพฺภา สมานา น พฺราหฺมเณ สกฺกโรนฺติ, น พฺราหฺมเณ ครุํ กโรนฺติ, น พฺราหฺมเณ มาเนนฺติ, น พฺราหฺมเณ ปูเชนฺติ, น พฺราหฺมเณ อปจายนฺตีติ.\csromanspacer{} อิติห อมฺพฏฺโฐ มาณโว อิทํ ตติยํ สเกฺยสุ อิพฺภวาทํ นิปาเตสิ.}

\csromanmatikaanchor{mtk.62}
\csromantocmarkref{section}{ทาสิปุตฺตวาท}{mtk.62}
\bookmark[dest={mtk.62},level=1]{ทาสิปุตฺตวาท}
\csromanheader{1}{3. อมฺพฏฺฐสุตฺต \textbf{ทาสิปุตฺตวาท}}
\proseitem{267}{\csromanfolio{87}อถ โข ภควโต เอตทโหสิ “อติพาฬฺหํ โข อยํ อมฺพฏฺโฐ มาณโว สเกฺยสุ อิพฺภวาเทน นิมฺมาเทติ, ยํนูนาหํ โคตฺตํ ปุจฺเฉยฺยนฺ”ติ.\csromanspacer{} อถ โข ภควา อมฺพฏฺฐํ มาณวํ เอตทโวจ “กถํ โคตฺโตสิ อมฺพฏฺฐา”ติ.\csromanspacer{} กณฺหายโนหมสฺมิ โภ โคตมาติ.\csromanspacer{} โปราณํ โข ปน เต อมฺพฏฺฐ มาตาเปตฺติกํ นามโคตฺตํ อนุสฺสรโต อยฺยปุตฺตา สกฺยา ภวนฺติ, ทาสิปุตฺโต ตฺวมสิ สกฺยานํ, สกฺยา โข ปน อมฺพฏฺฐ ราชานํ โอกฺกากํ ปิตามหํ ทหนฺติ.}

\prose{ภูตปุพฺพํ อมฺพฏฺฐ ราชา โอกฺกาโก ยา สา มเหสี ปิยา มนาปา, ตสฺสา ปุตฺตสฺส รชฺชํ ปริณาเมตุกาโม เชฏฺฐกุมาเร รฏฺฐสฺมา ปพฺพาเชสิ โอกฺกามุขํ กรกณฺฑํ\footnote{อุกฺกามุขํ กรกณฺฑุํ (สี, สฺยา)} หตฺถินิกํ สินิสูรํ\footnote{สินิปุรํ (สี, สฺยา)}, เต รฏฺฐสฺมา ปพฺพาชิตา หิมวนฺตปสฺเส โปกฺขรณิยา ตีเร มหาสากสณฺโฑ, ตตฺถ วาสํ กปฺเปสุํ, เต ชาติสมฺเภทภยา สกาหิ ภคินีหิ สทฺธิํ สํวาสํ กปฺเปสุํ.}

\prose{อถ โข อมฺพฏฺฐ ราชา โอกฺกาโก อมจฺเจ ปาริสชฺเช อามนฺเตสิ “กหํ นุ โข โภ เอตรหิ กุมารา สมฺมนฺตี”ติ.\csromanspacer{} อตฺถิ เทว หิมวนฺตปสฺเส โปกฺขรณิยา ตีเร มหาสากสณฺโฑ, ตตฺเถตรหิ กุมารา สมฺมนฺติ, เต ชาติสมฺเภทภยา สกาหิ ภคินีหิ สทฺธิํ สํวาสํ กปฺเปนฺตีติ.\csromanspacer{} อถ โข อมฺพฏฺฐ ราชา โอกฺกาโก อุทานํ อุทาเนสิ “สกฺยา วต โภ กุมารา ปรมสกฺยา วต โภ กุมารา”ติ, ตทคฺเค โข ปน อมฺพฏฺฐ สกฺยา ปญฺญายนฺติ, โส จ เนสํ ปุพฺพปุริโส.}

\prose{รญฺโญ โข ปน อมฺพฏฺฐ โอกฺกากสฺส ทิสา นาม ทาสี อโหสิ, สา กณฺหํ นาม\footnote{สา กณฺหํ (อิ)} ชเนสิ, ชาโต กโณฺห ปพฺยาหาสิ “โธวถ มํ อมฺม, นหาเปถ มํ อมฺม, อิมสฺมา มํ อสุจิสฺมา ปริโมเจถ, อตฺถาย โว ภวิสฺสามี”ติ.\csromanspacer{} ยถา โข ปน อมฺพฏฺฐ เอตรหิ มนุสฺสา ปิสาเจ ทิสฺวา ปิสาจาติ สญฺชานนฺติ, เอวเมว โข อมฺพฏฺฐ เตน โข ปน สมเยน มนุสฺสา ปิสาเจ กณฺหาติ สญฺชานนฺติ, เต เอวมาหํสุ “อยํ ชาโต ปพฺยาหาสิ, กโณฺห ชาโต ปิสาโจ ชาโต”ติ, ตทคฺเค โข ปน อมฺพฏฺฐ กณฺหายนา ปญฺญายนฺติ, โส จ กณฺหายนานํ ปุพฺพปุริโส.\csromanspacer{} อิติ โข \csromanfolio{88}เต อมฺพฏฺฐ โปราณํ มาตาเปตฺติกํ นามโคตฺตํ อนุสฺสรโต อยฺยปุตฺตา สกฺยา ภวนฺติ, ทาสิปุตฺโต ตฺวมสิ สกฺยานนฺติ.}

\proseitem{268}{เอวํ วุตฺเต เต มาณวกา ภควนฺตํ เอตทโวจุํ “มา ภวํ โคตโม อมฺพฏฺฐํ อติพาฬฺหํ ทาสิปุตฺตวาเทน นิมฺมาเทสิ, สุชาโต จ โภ โคตม อมฺพฏฺโฐ มาณโว, กุลปุตฺโต จ อมฺพฏฺโฐ มาณโว, พหุสฺสุโต จ อมฺพฏฺโฐ มาณโว, กลฺยาณวากฺกรโณ จ อมฺพฏฺโฐ มาณโว, ปณฺฑิโต จ อมฺพฏฺโฐ มาณโว, ปโหติ จ อมฺพฏฺโฐ มาณโว โภตา โคตเมน สทฺธิํ อสฺมิํ วจเน ปฏิมนฺเตตุนฺ”ติ.}

\proseitem{269}{อถ โข ภควา เต มาณวเก เอตทโวจ “สเจ โข ตุมฺหากํ มาณวกานํ เอวํ โหติ ‘ทุชฺชาโต จ อมฺพฏฺโฐ มาณโว, อกุลปุตฺโต จ อมฺพฏฺโฐ มาณโว, อปฺปสฺสุโต จ อมฺพฏฺโฐ มาณโว, อกลฺยาณวากฺกรโณ จ อมฺพฏฺโฐ มาณโว, ทุปฺปญฺโญ จ อมฺพฏฺโฐ มาณโว, น จ ปโหติ อมฺพฏฺโฐ มาณโว สมเณน โคตเมน สทฺธิํ อสฺมิํ วจเน ปฏิมนฺเตตุนฺ’ติ, ติฏฺฐตุ อมฺพฏฺโฐ มาณโว, ตุเมฺห มยา สทฺธิํ มนฺตโวฺห อสฺมิํ วจเน.\csromanspacer{} สเจ ปน ตุมฺหากํ มาณวกานํ เอวํ โหติ ‘สุชาโต จ อมฺพฏฺโฐ มาณโว, กุลปุตฺโต จ อมฺพฏฺโฐ มาณโว, พหุสฺสุโต จ อมฺพฏฺโฐ มาณโว, กลฺยาณวากฺกรโณ จ อมฺพฏฺโฐ มาณโว, ปณฺฑิโต จ อมฺพฏฺโฐ มาณโว, ปโหติ จ อมฺพฏฺโฐ มาณโว สมเณน โคตเมน สทฺธิํ อสฺมิํ วจเน ปฏิมนฺเตตุนฺ’ติ, ติฏฺฐถ ตุเมฺห, อมฺพฏฺโฐ มาณโว มยา สทฺธิํ ปฏิมนฺเตตู”ติ.\csromanspacer{} สุชาโต จ โภ โคตม อมฺพฏฺโฐ มาณโว, กุลปุตฺโต จ อมฺพฏฺโฐ มาณโว, พหุสฺสุโต จ อมฺพฏฺโฐ มาณโว, กลฺยาณวากฺกรโณ จ อมฺพฏฺโฐ มาณโว, ปณฺฑิโต จ อมฺพฏฺโฐ มาณโว, ปโหติ จ อมฺพฏฺโฐ มาณโว โภตา โคตเมน สทฺธิํ อสฺมิํ วจเน ปฏิมนฺเตตุํ, ตุณฺหี มยํ ภวิสฺสาม, อมฺพฏฺโฐ มาณโว โภตา โคตเมน สทฺธิํ อสฺมิํ วจเน ปฏิมนฺเตตูติ.}

\proseitem{270}{อถ โข ภควา อมฺพฏฺฐํ มาณวํ เอตทโวจ “อยํ โข ปน เต อมฺพฏฺฐ สหธมฺมิโก ปโญฺห อาคจฺฉติ, อกามา พฺยากาตพฺโพ, สเจ ตฺวํ น พฺยากริสฺสสิ, อญฺเญน วา อญฺญํ ปฏิจริสฺสสิ, ตุณฺหี วา ภวิสฺสสิ,}

\vspace{\tipitakaparskip}
\csromancenter{อมฺพฏฺฐสุตฺต ปกฺกมิสฺสสิ วา, เอตฺเถว เต สตฺตธา มุทฺธา ผลิสฺสติ.\csromanspacer{} ตํ กิํ มญฺญสิ อมฺพฏฺฐ, กินฺติ เต สุตํ พฺราหฺมณานํ วุทฺธานํ มหลฺลกานํ อาจริยปาจริยานํ ภาสมานานํ กุโตปภุติกา กณฺหายนา, โก จ กณฺหายนานํ ปุพฺพปุริโส”ติ.}
\prose{\csromanfolio{89}เอวํ วุตฺเต อมฺพฏฺโฐ มาณโว ตุณฺหี อโหสิ, ทุติยมฺปิ โข ภควา อมฺพฏฺฐํ มาณวํ เอตทโวจ “ตํ กิํ มญฺญสิ อมฺพฏฺฐ, กินฺติ เต สุตํ พฺราหฺมณานํ วุทฺธานํ มหลฺลกานํ อาจริยปาจริยานํ ภาสมานานํ กุโตปภุติกา กณฺหายนา, โก จ กณฺหายนานํ ปุพฺพปุริโส”ติ.\csromanspacer{} ทุติยมฺปิ โข อมฺพฏฺโฐ มาณโว ตุณฺหี อโหสิ.\csromanspacer{} อถ โข ภควา อมฺพฏฺฐํ มาณวํ เอตทโวจ “พฺยากโรหิ ทานิ อมฺพฏฺฐ, น ทานิ เต ตุณฺหีภาวสฺส กาโล, โย โข อมฺพฏฺฐ ตถาคเตน ยาวตติยกํ สหธมฺมิกํ ปญฺหํ ปุฏฺโฐ น พฺยากโรติ, เอตฺเถวสฺส สตฺตธา มุทฺธา ผลิสฺสตี”ติ.}

\proseitem{271}{เตน โข ปน สมเยน วชิรปาณี ยกฺโข มหนฺตํ อโยกูฏํ อาทาย อาทิตฺตํ สมฺปชฺชลิตํ สโชติภูตํ\footnote{สญฺโชติภูตํ (สฺยา)} อมฺพฏฺฐสฺส มาณวสฺส อุปริ เวหาสํ ฐิโต โหติ “สจายํ อมฺพฏฺโฐ มาณโว ภควตา ยาวตติยกํ สหธมฺมิกํ ปญฺหํ ปุฏฺโฐ น พฺยากริสฺสติ, เอตฺเถวสฺส สตฺตธา มุทฺธํ ผาเลสฺสามี”ติ.\csromanspacer{} ตํ โข ปน วชิรปาณิํ ยกฺขํ ภควา เจว ปสฺสติ อมฺพฏฺโฐ จ มาณโว.}

\proseitem{272}{อถ โข อมฺพฏฺโฐ มาณโว ภีโต สํวิคฺโค โลมหฏฺฐชาโต ภควนฺตํเยว ตาณํ คเวสี ภควนฺตํเยว เลณํ คเวสี ภควนฺตํเยว สรณํ คเวสี อุปนิสีทิตฺวา ภควนฺตํ เอตทโวจ “กิเมตํ\footnote{กิํ เม ตํ (ก)} ภวํ โคตโม อาห, ปุน ภวํ โคตโม พฺรวิตู”ติ\footnote{พฺรูตุ (สฺยา)}.}

\prose{ตํ กิํ มญฺญสิ อมฺพฏฺฐ, กินฺติ เต สุตํ พฺราหฺมณานํ วุทฺธานํ มหลฺลกานํ อาจริยปาจริยานํ ภาสมานานํ กุโตปภุติกา กณฺหายนา, โก จ กณฺหายนานํ ปุพฺพปุริโสติ.\csromanspacer{} เอวเมว เม โภ โคตม สุตํ ยเถว ภวํ โคตโม อาห, ตโตปภุติกา กณฺหายนา, โส จ กณฺหายนานํ ปุพฺพปุริโสติ.}

\csromanmatikaanchor{mtk.63}
\csromantocmarkref{section}{อมฺพฏฺฐวํสกถา}{mtk.63}
\bookmark[dest={mtk.63},level=1]{อมฺพฏฺฐวํสกถา}
\csromanheader{1}{\textbf{อมฺพฏฺฐวํสกถา}}
\proseitem{273}{\csromanfolio{90}เอวํ วุตฺเต เต มาณวกา อุนฺนาทิโน อุจฺจาสทฺทมหาสทฺทา อเหสุํ “ทุชฺชาโต กิร โภ อมฺพฏฺโฐ มาณโว, อกุลปุตฺโต กิร โภ อมฺพฏฺโฐ มาณโว, ทาสิปุตฺโต กิร โภ อมฺพฏฺโฐ มาณโว สกฺยานํ, อยฺยปุตฺตา กิร โภ อมฺพฏฺฐสฺส มาณวสฺส สกฺยา ภวนฺติ, ธมฺมวาทิํเยว กิร มยํ สมณํ โคตมํ อปสาเทตพฺพํ อมญฺญิมฺหา”ติ.}

\proseitem{274}{อถ โข ภควโต เอตทโหสิ “อติพาฬฺหํ โข อิเม มาณวกา อมฺพฏฺฐํ มาณวํ ทาสิปุตฺตวาเทน นิมฺมาเทนฺติ, ยํนูนาหํ ปริโมเจยฺยนฺ”ติ.\csromanspacer{} อถ โข ภควา เต มาณวเก เอตทโวจ “มา โข ตุเมฺห มาณวกา อมฺพฏฺฐํ มาณวํ อติพาฬฺหํ ทาสิปุตฺตวาเทน นิมฺมาเทถ, อุฬาโร โส กโณฺห อิสิ อโหสิ, โส ทกฺขิณชนปทํ คนฺตฺวา พฺรหฺมมนฺเต อธียิตฺวา ราชานํ โอกฺกากํ อุปสงฺกมิตฺวา มทฺทรูปิํ ธีตรํ ยาจิ, ตสฺส ราชา โอกฺกาโก “โก เนวํ เร อยํ มยฺหํ ทาสิปุตฺโต สมาโน มทฺทรูปิํ ธีตรํ ยาจตี”ติ กุปิโต อนตฺตมโน ขุรปฺปํ สนฺนยฺหิ\footnote{สนฺนหิ (ก)}, โส ตํ ขุรปฺปํ เนว อสกฺขิ มุญฺจิตุํ, โน ปฏิสํหริตุํ.}

\prose{อถ โข มาณวกา อมจฺจา ปาริสชฺชา กณฺหํ อิสิํ อุปสงฺกมิตฺวา เอตทโวจุํ “โสตฺถิ ภทฺทนฺเต\footnote{ภทนฺเต (สี, สฺยา)} โหตุ รญฺโญ, โสตฺถิ ภทฺทนฺเต โหตุ รญฺโญ”ติ.\csromanspacer{} โสตฺถิ ภวิสฺสติ รญฺโญ, อปิ จ ราชา ยทิ อโธ ขุรปฺปํ มุญฺจิสฺสติ, ยาวตา รญฺโญ วิชิตํ, เอตฺตาวตา ปถวี อุนฺทฺริยิสฺสตีติ.\csromanspacer{} โสตฺถิ ภทฺทนฺเต โหตุ รญฺโญ, โสตฺถิ ชนปทสฺสาติ.\csromanspacer{} โสตฺถิ ภวิสฺสติ รญฺโญ โสตฺถิ ชนปทสฺส, อปิ จ ราชา ยทิ อุทฺธํ ขุรปฺปํ มุญฺจิสฺสติ, ยาวตา รญฺโญ วิชิตํ, เอตฺตาวตา สตฺตวสฺสานิ เทโว น วสฺสิสฺสตีติ.\csromanspacer{} โสตฺถิ ภทฺทนฺเต โหตุ รญฺโญ โสตฺถิ ชนปทสฺส เทโว จ วสฺสตูติ.\csromanspacer{} โสตฺถิ ภวิสฺสติ รญฺโญ โสตฺถิ ชนปทสฺส เทโว จ วสฺสิสฺสติ, อปิ จ ราชา เชฏฺฐกุมาเร ขุรปฺปํ ปติฏฺฐาเปตุ, โสตฺถิ กุมาโร ปลฺโลโม ภวิสฺสตีติ.\csromanspacer{} อถ โข มาณวกา อมจฺจา โอกฺกากสฺส อาโรเจสุํ “โอกฺกาโก เชฏฺฐกุมาเร ขุรปฺปํ ปติฏฺฐาเปตุ, โสตฺถิ กุมาโร ปลฺโลโม ภวิสฺสตี”ติ.\csromanspacer{} อถ โข ราชา โอกฺกาโก เชฏฺฐกุมาเร ขุรปฺปํ ปติฏฺฐเปสิ, โสตฺถิ กุมาโร}

\vspace{\tipitakaparskip}
\csromancenter{อมฺพฏฺฐสุตฺต ปลฺโลโม สมภวิ.\csromanspacer{} อถ โข ตสฺส ราชา โอกฺกาโก ภีโต สํวิคฺโค โลมหฏฺฐชาโต พฺรหฺมทณฺเฑน ตชฺชิโต มทฺทรูปิํ ธีตรํ อทาสิ.\csromanspacer{} มา โข ตุเมฺห มาณวกา อมฺพฏฺฐํ มาณวํ อติพาฬฺหํ ทาสิปุตฺตวาเทน นิมฺมาเทถ, อุฬาโร โส กโณฺห อิสิ อโหสีติ.}
\csromanmatikaanchor{mtk.64}
\csromantocmarkref{section}{ขตฺติยเสฏฺฐภาว}{mtk.64}
\bookmark[dest={mtk.64},level=1]{ขตฺติยเสฏฺฐภาว}
\csromanheader{1}{\textbf{ขตฺติยเสฏฺฐภาว}}
\proseitem{275}{\csromanfolio{91}อถ โข ภควา อมฺพฏฺฐํ มาณวํ อามนฺเตสิ “ตํ กิํ มญฺญสิ อมฺพฏฺฐ, อิธ ขตฺติยกุมาโร พฺราหฺมณกญฺญาย สทฺธิํ สํวาสํ กปฺเปยฺย, เตสํ สํวาสมนฺวาย ปุตฺโต ชาเยถ, โย โส ขตฺติยกุมาเรน พฺราหฺมณกญฺญาย ปุตฺโต อุปฺปนฺโน, อปิ นุ โส ลเภถ พฺราหฺมเณสุ อาสนํ วา อุทกํ วา”ติ.\csromanspacer{} ลเภถ โภ โคตม.\csromanspacer{} อปินุ นํ พฺราหฺมณา โภเชยฺยุํ สทฺเธ วา ถาลิปาเก วา ยญฺเญ วา ปาหุเน วาติ.\csromanspacer{} โภเชยฺยุํ โภ โคตม.\csromanspacer{} อปินุ นํ พฺราหฺมณา มนฺเต วาเจยฺยุํ วา โน วาติ.\csromanspacer{} วาเจยฺยุํ โภ โคตม.\csromanspacer{} อปินุสฺส อิตฺถีสุ อาวฏํ วา อสฺส อนาวฏํ วาติ.\csromanspacer{} อนาวฏํ หิสฺส โภ โคตม.\csromanspacer{} อปินุ นํ ขตฺติยา ขตฺติยาภิเสเกน อภิสิญฺเจยฺยุนฺติ.\csromanspacer{} โนหิทํ โภ โคตม, ตํ กิสฺส เหตุ, มาติโต หิ โภ โคตม อนุปปนฺโนติ.}

\prose{ตํ กิํ มญฺญสิ อมฺพฏฺฐ, อิธ พฺราหฺมณกุมาโร ขตฺติยกญฺญาย สทฺธิํ สํวาสํ กปฺเปยฺย, เตสํ สํวาสมนฺวาย ปุตฺโต ชาเยถ, โย โส พฺราหฺมณกุมาเรน ขตฺติยกญฺญาย ปุตฺโต อุปฺปนฺโน, อปินุ โส ลเภถ พฺราหฺมเณสุ อาสนํ วา อุทกํ วาติ.\csromanspacer{} ลเภถ โภ โคตม.\csromanspacer{} อปินุ นํ พฺราหฺมณา โภเชยฺยุํ สทฺเธ วา ถาลิปาเก วา ยญฺเญ วา ปาหุเน วาติ.\csromanspacer{} โภเชยฺยุํ โภ โคตม.\csromanspacer{} อปินุ นํ พฺราหฺมณา มนฺเต วาเจยฺยุํ วา โน วาติ.\csromanspacer{} วาเจยฺยุํ โภ โคตม.\csromanspacer{} อปินุสฺส อิตฺถีสุ อาวฏํ วา อสฺส อนาวฏํ วาติ.\csromanspacer{} อนาวฏํ หิสฺส โภ โคตม.\csromanspacer{} อปินุ นํ ขตฺติยา ขตฺติยาภิเสเกน อภิสิญฺเจยฺยุนฺติ.\csromanspacer{} โนหิทํ โภ โคตม, ตํ กิสฺส เหตุ, ปิติโต หิ โภ โคตม อนุปปนฺโนติ.}

\proseitem{276}{อิติ โข อมฺพฏฺฐ อิตฺถิยา วา อิตฺถิํ กริตฺวา ปุริเสน วา ปุริสํ กริตฺวา ขตฺติยาว เสฏฺฐา หีนา พฺราหฺมณา, ตํ กิํ มญฺญสิ อมฺพฏฺฐ, อิธ พฺราหฺมณา พฺราหฺมณํ กิสฺมิญฺจิเทว ปกรเณ ขุรมุณฺฑํ กริตฺวา \csromanfolio{92}ภสฺสปุเฏน วธิตฺวา รฏฺฐา วา นครา วา ปพฺพาเชยฺยุํ, อปินุ โส ลเภถ พฺราหฺมเณสุ อาสนํ วา อุทกํ วาติ.\csromanspacer{} โนหิทํ โภ โคตม.\csromanspacer{} อปินุ นํ พฺราหฺมณา โภเชยฺยุํ สทฺเธ วา ถาลิปาเก วา ยญฺเญ วา ปาหุเน วาติ.\csromanspacer{} โนหิทํ โภ โคตม.\csromanspacer{} อปินุ นํ พฺราหฺมณา มนฺเต วาเจยฺยุํ วา โน วาติ.\csromanspacer{} โนหิทํ โภ โคตม.\csromanspacer{} อปินุสฺส อิตฺถีสุ อาวฏํ วา อสฺส อนาวฏํ วาติ.\csromanspacer{} อาวฏํ หิสฺส โภ โคตม.}

\prose{ตํ กิํ มญฺญสิ อมฺพฏฺฐ, อิธ ขตฺติยา ขตฺติยํ กิสฺมิญฺจิเทว ปกรเณ ขุรมุณฺฑํ กริตฺวา ภสฺสปุเฏน วธิตฺวา รฏฺฐา วา นครา วา ปพฺพาเชยฺยุํ, อปินุ โส ลเภถ พฺราหฺมเณสุ อาสนํ วา อุทกํ วาติ.\csromanspacer{} ลเภถ โภ โคตม.\csromanspacer{} อปินุ นํ พฺราหฺมณา โภเชยฺยุํ สทฺเธ วา ถาลิปาเก วา ยญฺเญ วา ปาหุเน วาติ.\csromanspacer{} โภเชยฺยุํ โภ โคตม.\csromanspacer{} อปินุ นํ พฺราหฺมณา มนฺเต วาเจยฺยุํ วา โน วาติ.\csromanspacer{} วาเจยฺยุํ โภ โคตม.\csromanspacer{} อปินุสฺส อิตฺถีสุ อาวฏํ วา อสฺส อนาวฏํ วาติ.\csromanspacer{} อนาวฏํ หิสฺส โภ โคตม.}

\proseitem{277}{เอตฺตาวตา โข อมฺพฏฺฐ ขตฺติโย ปรมนิหีนตํ ปตฺโต โหติ, ยเทว นํ ขตฺติยา ขุรมุณฺฑํ กริตฺวา ภสฺสปุเฏน วธิตฺวา รฏฺฐา วา นครา วา ปพฺพาเชนฺติ, อิติ โข อมฺพฏฺฐ ยทา ขตฺติโย ปรมนิหีนตํ ปตฺโต โหติ, ตทาปิ ขตฺติยาว เสฏฺฐา, หีนา พฺราหฺมณา.\csromanspacer{} พฺรหฺมุนา เปสา อมฺพฏฺฐ\footnote{พฺรหฺมุนาปิ อมฺพฏฺฐ (ก), พฺรหฺมุนาปิ เอสา อมฺพฏฺฐ (อิ)} สนงฺกุมาเรน คาถา ภาสิตา\csromandash{}}

\setlength{\csromangathapagewidth}{0pt}
\csromangathameasuregroup{}{“ขตฺติโย เสฏฺโฐ ชเนตสฺมิํ, \\ เย โคตฺตปฏิสาริโน. \\ วิชฺชาจรณสมฺปนฺโน, \\ โส เสฏฺโฐ เทวมานุเส”ติ.}
\setlength{\csromangathaleftcol}{0pt}
\csromangathagroup{\csromangathastanza{“ขตฺติโย เสฏฺโฐ ชเนตสฺมิํ, \\ เย โคตฺตปฏิสาริโน. \\ วิชฺชาจรณสมฺปนฺโน, \\ โส เสฏฺโฐ เทวมานุเส”ติ.}}

\prose{สา โข ปเนสา อมฺพฏฺฐ พฺรหฺมุนา สนงฺกุมาเรน คาถา สุคีตา โน ทุคฺคีตา, สุภาสิตา โน ทุพฺภาสิตา, อตฺถสญฺหิตา โน อนตฺถสญฺหิตา, อนุมตา มยา.\csromanspacer{} อหํปิ หิ อมฺพฏฺฐ เอวํ วทามิ\csromandash{}}

\csromanheader{2}{3. อมฺพฏฺฐสุตฺต}
\setlength{\csromangathapagewidth}{0pt}
\csromangathameasuregroup{}{“ขตฺติโย เสฏฺโฐ ชเนตสฺมิํ, \\ เย โคตฺตปฏิสาริโน. \\ วิชฺชาจรณสมฺปนฺโน, \\ โส เสฏฺโฐ เทวมานุเส”ติ.}
\setlength{\csromangathaleftcol}{0pt}
\csromangathagroup{\csromangathastanza{\csromanfolio{93}“ขตฺติโย เสฏฺโฐ ชเนตสฺมิํ, \\ เย โคตฺตปฏิสาริโน. \\ วิชฺชาจรณสมฺปนฺโน, \\ โส เสฏฺโฐ เทวมานุเส”ติ.}}

\csromanheader{2}{ภาณวาโร ปฐโม.}
\csromansectionrule
\csromanmatikaanchor{mtk.65}
\csromantocmarkref{section}{วิชฺชาจรณกถา}{mtk.65}
\bookmark[dest={mtk.65},level=1]{วิชฺชาจรณกถา}
\csromanheader{1}{\textbf{วิชฺชาจรณกถา}}
\proseitem{278}{กตมํ ปน ตํ โภ โคตม จรณํ, กตมา จ ปน สา วิชฺชาติ.\csromanspacer{} น โข อมฺพฏฺฐ อนุตฺตราย วิชฺชาจรณสมฺปทาย ชาติวาโท วา วุจฺจติ, โคตฺตวาโท วา วุจฺจติ, มานวาโท วา วุจฺจติ “อรหสิ วา มํ ตฺวํ, น วา มํ ตฺวํ อรหสี”ติ.\csromanspacer{} ยตฺถ โข อมฺพฏฺฐ อาวาโห วา โหติ, วิวาโห วา โหติ, อาวาหวิวาโห วา โหติ, เอตฺเถตํ วุจฺจติ ชาติวาโท วา อิติปิ, โคตฺตวาโท วา อิติปิ, มานวาโท วา อิติปิ “อรหสิ วา มํ ตฺวํ, น วา มํ ตฺวํ อรหสี”ติ.\csromanspacer{} เย หิ เกจิ อมฺพฏฺฐ ชาติวาทวินิพทฺธา วา โคตฺตวาทวินิพทฺธา วา มานวาทวินิพทฺธา วา อาวาหวิวาห\-วินิพทฺธา วา, อารกา เต อนุตฺตราย วิชฺชาจรณสมฺปทาย.\csromanspacer{} ปหาย โข อมฺพฏฺฐ ชาติวาทวินิพทฺธญฺจ โคตฺตวาทวินิพทฺธญฺจ มานวาทวินิพทฺธญฺจ อาวาหวิวาห\-วินิพทฺธญฺจ อนุตฺตราย วิชฺชาจรณสมฺปทาย สจฺฉิกิริยา โหตีติ.}

\proseitem{279}{กตมํ ปน ตํ โภ โคตม จรณํ, กตมา จ สา วิชฺชาติ.\csromanspacer{} อิธ อมฺพฏฺฐ ตถาคโต โลเก อุปฺปชฺชติ อรหํ สมฺมาสมฺพุทฺโธ วิชฺชาจรณสมฺปนฺโน สุคโต โลกวิทู อนุตฺตโร ปุริสทมฺมสารถิ สตฺถา เทวมนุสฺสานํ พุทฺโธ ภควา, โส อิมํ โลกํ สเทวกํ สมารกํ สพฺรหฺมกํ สสฺสมณพฺราหฺมณิํ ปชํ สเทวมนุสฺสํ สยํ อภิญฺญา สจฺฉิกตฺวา ปเวเทติ, โส ธมฺมํ เทเสติ อาทิกลฺยาณํ มชฺเฌกลฺยาณํ ปริโยสานกลฺยาณํ สาตฺถํ สพฺยญฺชนํ เกวลปริปุณฺณํ ปริสุทฺธํ พฺรหฺมจริยํ ปกาเสติ, ตํ ธมฺมํ สุณาติ คหปติ วา คหปติปุตฺโต วา อญฺญตรสฺมิํ วา กุเล ปจฺจาชาโต, โส ตํ ธมฺมํ สุตฺวา \csromanfolio{94}ตถาคเต สทฺธํ ปฏิลภติ, โส เตน สทฺธาปฏิลาเภน สมนฺนาคโต อิติ ปฏิสญฺจิกฺขติ ฯเปฯ. (ยถา สามญฺญผเล, เอวํ วิตฺถาเรตพฺพํ.)}

\prose{โส วิวิจฺเจว กาเมหิ วิวิจฺจ อกุสเลหิ ธมฺเมหิ สวิตกฺกํ สวิจารํ วิเวกชํ ปีติสุขํ ปฐมํ ฌานํ อุปสมฺปชฺช วิหรติ ฯเปฯ.\csromanspacer{} อิทมฺปิสฺส โหติ จรณสฺมิํ.}

\prose{ปุน จปรํ อมฺพฏฺฐ ภิกฺขุ วิตกฺกวิจารานํ วูปสมา อชฺฌตฺตํ สมฺปสาทนํ เจตโส เอโกทิภาวํ อวิตกฺกํ อวิจารํ สมาธิชํ ปีติสุขํ ทุติยํ ฌานํ อุปสมฺปชฺช วิหรติ ฯเปฯ.\csromanspacer{} อิทมฺปิสฺส โหติ จรณสฺมิํ.}

\prose{ปุน จปรํ อมฺพฏฺฐ ภิกฺขุ ปีติยา จ วิราคา อุเปกฺขโก จ วิหรติ สโต จ สมฺปชาโน, สุขญฺจ กาเยน ปฏิสํเวเทติ, ยํ ตํ อริยา อาจิกฺขนฺติ “อุเปกฺขโก สติมา สุขวิหารี”ติ, ตติยํ ฌานํ อุปสมฺปชฺช วิหรติ ฯเปฯ.\csromanspacer{} อิทมฺปิสฺส โหติ จรณสฺมิํ.}

\prose{ปุน จปรํ อมฺพฏฺฐ ภิกฺขุ สุขสฺส จ ปหานา ทุกฺขสฺส จ ปหานา ปุพฺเพว โสมนสฺสโทมนสฺสานํ อตฺถงฺคมา อทุกฺขมสุขํ อุเปกฺขาสติปาริสุทฺธิํ จตุตฺถํ ฌานํ อุปสมฺปชฺช วิหรติ ฯเปฯ.\csromanspacer{} อิทมฺปิสฺส โหติ จรณสฺมิํ.\csromanspacer{} อิทํ โข ตํ อมฺพฏฺฐ จรณํ.}

\prose{โส เอวํ สมาหิเต จิตฺเต ปริสุทฺเธ ปริโยทาเต อนงฺคเณ วิคตูปกฺกิเลเส มุทุภูเต กมฺมนิเย ฐิเต อาเนญฺชปฺปตฺเต ญาณทสฺสนาย จิตฺตํ อภินีหรติ อภินินฺนาเมติ ฯเปฯ.\csromanspacer{} อิทมฺปิสฺส โหติ วิชฺชาย ฯเปฯ นาปรํ อิตฺถตฺตายาติ ปชานาติ, อิทมฺปิสฺส โหติ วิชฺชาย.\csromanspacer{} อยํ โข สา อมฺพฏฺฐ วิชฺชา.}

\prose{อยํ วุจฺจติ อมฺพฏฺฐ ภิกฺขุ “วิชฺชาสมฺปนฺโน” อิติปิ, “จรณสมฺปนฺโน” อิติปิ, “วิชฺชาจรณสมฺปนฺโน” อิติปิ, อิมาย จ อมฺพฏฺฐ วิชฺชาสมฺปทาย จรณสมฺปทาย จ อญฺญาวิชฺชาสมฺปทา จ จรณสมฺปทา จ อุตฺตริตรา วา ปณีตตรา วา นตฺถิ.}

\csromanmatikaanchor{mtk.66}
\csromantocmarkref{section}{จตุอปายมุข}{mtk.66}
\bookmark[dest={mtk.66},level=1]{จตุอปายมุข}
\csromanheader{1}{\textbf{จตุอปายมุข}}
\proseitem{280}{อิมาย โข อมฺพฏฺฐ อนุตฺตราย วิชฺชาจรณสมฺปทาย จตฺตาริ อปายมุขานิ ภวนฺติ.\csromanspacer{} กตมานิ จตฺตาริ, อิธ อมฺพฏฺฐ เอกจฺโจ สมโณ วา พฺราหฺมโณ วา อิมญฺเญว อนุตฺตรํ วิชฺชาจรณสมฺปทํ อนภิสมฺภุณมาโน}

\vspace{\tipitakaparskip}
\csromancenter{อมฺพฏฺฐสุตฺต ขาริวิธมาทาย\footnote{ขาริวิวิธมาทาย (สี, สฺยา, อิ)} อรญฺญายตนํ อชฺโฌคาหติ “ปวตฺตผลโภชโน ภวิสฺสามี”ติ, โส อญฺญทตฺถุ วิชฺชา\-จรณ\-สมฺปนฺนสฺเสว ปริจารโก สมฺปชฺชติ.\csromanspacer{} อิมาย โข อมฺพฏฺฐ อนุตฺตราย วิชฺชาจรณสมฺปทาย อิทํ ปฐมํ อปายมุขํ ภวติ.}
\prose{\csromanfolio{95}ปุน จปรํ อมฺพฏฺฐ อิเธกจฺโจ สมโณ วา พฺราหฺมโณ วา อิมญฺเจว อนุตฺตรํ วิชฺชาจรณสมฺปทํ อนภิสมฺภุณมาโน ปวตฺตผลโภชนตญฺจ อนภิสมฺภุณมาโน กุทาลปิฏกํ\footnote{กุทฺทาลปิฏกํ (สี, สฺยา, อิ)} อาทาย อรญฺญวนํ อชฺโฌคาหติ “กนฺทมูลผลโภชโน ภวิสฺสามี”ติ, โส อญฺญทตฺถุ วิชฺชา\-จรณ\-สมฺปนฺนสฺเสว ปริจารโก สมฺปชฺชติ.\csromanspacer{} อิมาย โข อมฺพฏฺฐ อนุตฺตราย วิชฺชาจรณสมฺปทาย อิทํ ทุติยํ อปายมุขํ ภวติ.}

\prose{ปุน จปรํ อมฺพฏฺฐ อิเธกจฺโจ สมโณ วา พฺราหฺมโณ วา อิมญฺเจว อนุตฺตรํ วิชฺชาจรณสมฺปทํ อนภิสมฺภุณมาโน ปวตฺตผลโภชนตญฺจ อนภิสมฺภุณมาโน กนฺทมูลผลโภชนตญฺจ อนภิสมฺภุณมาโน คามสามนฺตํ วา นิคมสามนฺตํ วา อคฺยาคารํ กริตฺวา อคฺคิํ ปริจรนฺโต อจฺฉติ, โส อญฺญทตฺถุ วิชฺชา\-จรณ\-สมฺปนฺนสฺเสว ปริจารโก สมฺปชฺชติ.\csromanspacer{} อิมาย โข อมฺพฏฺฐ อนุตฺตราย วิชฺชาจรณสมฺปทาย อิทํ ตติยํ อปายมุขํ ภวติ.}

\prose{ปุน จปรํ อมฺพฏฺฐ อิเธกจฺโจ สมโณ วา พฺราหฺมโณ วา อิมญฺเจว อนุตฺตรํ วิชฺชาจรณสมฺปทํ อนภิสมฺภุณมาโน ปวตฺตผลโภชนตญฺจ อนภิสมฺภุณมาโน กนฺทมูลผลโภชนตญฺจ อนภิสมฺภุณมาโน อคฺคิปาริจริยญฺจ อนภิสมฺภุณมาโน จาตุมหาปเถ จตุทฺวารํ อคารํ กริตฺวา อจฺฉติ “โย อิมาหิ จตูหิ ทิสาหิ อาคมิสฺสติ สมโณ วา, พฺราหฺมโณ วา, ตมหํ ยถาสตฺติ ยถาพลํ ปฏิปูเชสฺสามี”ติ, โส อญฺญทตฺถุ วิชฺชา\-จรณ\-สมฺปนฺนสฺเสว ปริจารโก สมฺปชฺชติ.\csromanspacer{} อิมาย โข อมฺพฏฺฐ อนุตฺตราย วิชฺชาจรณสมฺปทาย อิทํ จตุตฺถํ อปายมุขํ ภวติ.\csromanspacer{} อิมาย โข อมฺพฏฺฐ อนุตฺตราย วิชฺชาจรณสมฺปทาย อิมานิ จตฺตาริ อปายมุขานิ ภวนฺติ.}

\proseitem{281}{ตํ กิํ มญฺญสิ อมฺพฏฺฐ, อปินุ ตฺวํ อิมาย อนุตฺตราย วิชฺชาจรณสมฺปทาย สนฺทิสฺสสิ สาจริยโกติ.\csromanspacer{} โนหิทํ โภ โคตม, โกจาหํ \csromanfolio{96}โภ โคตม สาจริยโก, กา จ อนุตฺตรา วิชฺชาจรณสมฺปทา, อารกาหํ โภ โคตม อนุตฺตราย วิชฺชาจรณสมฺปทาย สาจริยโกติ.}

\prose{ตํ กิํ มญฺญสิ อมฺพฏฺฐ, อปินุ ตฺวํ อิมญฺเจว อนุตฺตรํ วิชฺชาจรณสมฺปทํ อนภิสมฺภุณมาโน ขาริวิธมาทาย อรญฺญวนมชฺโฌคาหสิ สาจริยโก “ปวตฺตผลโภชโน ภวิสฺสามี”ติ.\csromanspacer{} โนหิทํ โภ โคตม.}

\prose{ตํ กิํ มญฺญสิ อมฺพฏฺฐ, อปินุ ตฺวํ อิมญฺเจว อนุตฺตรํ วิชฺชาจรณสมฺปทํ อนภิสมฺภุณมาโน ปวตฺตผลโภชนตญฺจ อนภิสมฺภุณมาโน กุทาลปิฏกํ อาทาย อรญฺญวนมชฺโฌคาหสิ สาจริยโก “กนฺทมูลผลโภชโน ภวิสฺสามี”ติ.\csromanspacer{} โนหิทํ โภ โคตม.}

\prose{ตํ กิํ มญฺญสิ อมฺพฏฺฐ, อปินุ ตฺวํ อิมญฺเจว อนุตฺตรํ วิชฺชาจรณสมฺปทํ อนภิสมฺภุณมาโน ปวตฺตผลโภชนตญฺจ อนภิสมฺภุณมาโน กนฺทมูลผลโภชนตญฺจ อนภิสมฺภุณมาโน คามสามนฺตํ วา นิคมสามนฺตํ วา อคฺยาคารํ กริตฺวา อคฺคิํ ปริจรนฺโต อจฺฉสิ สาจริยโกติ.\csromanspacer{} โนหิทํ โภ โคตม.}

\prose{ตํ กิํ มญฺญสิ อมฺพฏฺฐ, อปินุ ตฺวํ อิมญฺเจว อนุตฺตรํ วิชฺชาจรณสมฺปทํ อนภิสมฺภุณมาโน ปวตฺตผลโภชนตญฺจ อนภิสมฺภุณมาโน กนฺทมูลผลโภชนตญฺจ อนภิสมฺภุณมาโน อคฺคิปาริจริยญฺจ อนภิสมฺภุณมาโน จาตุมหาปเถ จตุทฺวารํ อคารํ กริตฺวา อจฺฉสิ สาจริยโก “โย อิมาหิ จตูหิ ทิสาหิ อาคมิสฺสติ สมโณ วา พฺราหฺมโณ วา, ตํ มยํ ยถาสตฺติ ยถาพลํ ปฏิปูเชสฺสามา”ติ.\csromanspacer{} โนหิทํ โภ โคตม.}

\proseitem{282}{อิติ โข อมฺพฏฺฐ, อิมาย เจว ตฺวํ อนุตฺตราย วิชฺชาจรณสมฺปทาย ปริหีโน สาจริยโก, เย จิเม อนุตฺตราย วิชฺชาจรณสมฺปทาย จตฺตาริ อปายมุขานิ ภวนฺติ, ตโต จ ตฺวํ ปริหีโน สาจริยโก.\csromanspacer{} ภาสิตา โข ปน เต เอสา อมฺพฏฺฐ อาจริเยน พฺราหฺมเณน โปกฺขรสาตินา วาจา “เก จ มุณฺฑกา สมณกา อิพฺภา กณฺหา พนฺธุปาทาปจฺจา, กา จ เตวิชฺชานํ พฺราหฺมณานํ สากจฺฉา”ติ อตฺตนา อาปายิโกปิ อปริปูรมาโน.\csromanspacer{} ปสฺส อมฺพฏฺฐ ยาว อปรทฺธญฺจ เต อิทํ อาจริยสฺส พฺราหฺมณสฺส โปกฺขรสาติสฺส.}

\csromanmatikaanchor{mtk.67}
\csromantocmarkref{section}{ปุพฺพกอิสิภาวานุโยค}{mtk.67}
\bookmark[dest={mtk.67},level=1]{ปุพฺพกอิสิภาวานุโยค}
\csromanheader{1}{3. อมฺพฏฺฐสุตฺต \textbf{ปุพฺพกอิสิภาวานุโยค}}
\proseitem{283}{\csromanfolio{97}พฺราหฺมโณ โข ปน อมฺพฏฺฐ โปกฺขรสาติ รญฺโญ ปเสนทิสฺส โกสลสฺส ทตฺติกํ ภุญฺชติ, ตสฺส ราชา ปเสนทิ โกสโล สมฺมุขีภาวํปิ น ททาติ, ยทาปิ เตน มนฺเตติ, ติโรทุสฺสนฺเตน มนฺเตติ.\csromanspacer{} ยสฺส โข ปน อมฺพฏฺฐ ธมฺมิกํ ปยาตํ ภิกฺขํ ปฏิคฺคเณฺหยฺย, กถํ ตสฺส ราชา ปเสนทิ โกสโล สมฺมุขีภาวํปิ น ทเทยฺย, ปสฺส อมฺพฏฺฐ ยาว อปรทฺธญฺจ เต อิทํ อาจริยสฺส พฺราหฺมณสฺส โปกฺขรสาติสฺส.}

\proseitem{284}{ตํ กิํ มญฺญสิ อมฺพฏฺฐ, อิธ ราชา ปเสนทิ โกสโล หตฺถิคีวาย วา นิสินฺโน อสฺสปิฏฺเฐ วา นิสินฺโน รถูปตฺถเร วา ฐิโต อุคฺเคหิ วา ราชญฺเญหิ วา กิญฺจิเทว มนฺตนํ มนฺเตยฺย, โส ตมฺหา ปเทสา อปกฺกมฺม เอกมนฺตํ ติฏฺเฐยฺย, อถ อาคจฺเฉยฺย สุทฺโท วา สุทฺททาโส วา, ตสฺมิํ ปเทเส ฐิโต ตเทว มนฺตนํ มนฺเตยฺย “เอวํปิ ราชา ปเสนทิ โกสโล อาห, เอวํปิ ราชา ปเสนทิ โกสโล อาหา”ติ.\csromanspacer{} อปินุ โส ราชภณิตํ วา ภณติ ราชมนฺตนํ วา มนฺเตติ, เอตฺตาวตา โส อสฺส ราชา วา ราชมตฺโต วาติ.\csromanspacer{} โนหิทํ โภ โคตม.}

\proseitem{285}{เอวเมว โข ตฺวํ อมฺพฏฺฐ เย เต อเหสุํ พฺราหฺมณานํ ปุพฺพกา อิสโย มนฺตานํ กตฺตาโร มนฺตานํ ปวตฺตาโร, เยสมิทํ เอตรหิ พฺราหฺมณา โปราณํ มนฺตปทํ คีตํ ปวุตฺตํ สมิหิตํ, ตทนุคายนฺติ ตทนุภาสนฺติ ภาสิตมนุภาสนฺติ วาจิตมนุวาเจนฺติ, เสยฺยถิทํ, อฏฺฐโก วามโก วามเทโว เวสฺสามิตฺโต ยมตคฺคิ\footnote{ยมทคฺคิ (ก)} องฺคีรโส ภารทฺวาโช วาเสฏฺโฐ กสฺสโป ภคุ, ตฺยาหํ มนฺเต อธิยามิ สาจริยโกติ, ตาวตา ตฺวํ ภวิสฺสสิ อิสิ วา อิสิตฺถาย วา ปฏิปนฺโนติ เนตํ ฐานํ วิชฺชติ.}

\proseitem{286}{ตํ กิํ มญฺญสิ อมฺพฏฺฐ, กินฺติ เต สุตํ พฺราหฺมณานํ วุทฺธานํ มหลฺลกานํ อาจริยปาจริยานํ ภาสมานานํ เย เต อเหสุํ พฺราหฺมณานํ ปุพฺพากา อิสโย มนฺตานํ กตฺตาโร มนฺตานํ ปวตฺตาโร, \csromanfolio{98}เยสมิทํ เอตรหิ พฺราหฺมณา โปราณํ มนฺตปทํ คีตํ ปวุตฺตํ สมิหิตํ, ตทนุคายนฺติ ตทนุภาสนฺติ ภาสิตมนุภาสนฺติ วาจิตมนุวาเจนฺติ, เสยฺยถิทํ, อฏฺฐโก วามโก วามเทโว เวสฺสามิตฺโต ยมตคฺคิ องฺคีรโส ภารทฺวาโช วาเสฏฺโฐ กสฺสโป ภคุ, เอวํ สุ เต สุนฺหาตา สุวิลิตฺตา กปฺปิตเกสมสฺสู อามุกฺก\-มณิกุณฺฑลา\-ภรณา\footnote{อามุตฺตมาลาภรณา (สี, สฺยา, อิ)} โอทาตวตฺถวสนา ปญฺจหิ กามคุเณหิ สมปฺปิตา สมงฺคีภูตา ปริจาเรนฺติ, เสยฺยถาปิ ตฺวํ เอตรหิ สาจริยโกติ.\csromanspacer{} โนหิทํ โภ โคตม.}

\prose{ฯเปฯ เอวํ สุ เต สาลีนํ โอทนํ สุจิมํสูปเสจนํ วิจิตกาฬกํ อเนกสูปํ อเนกพฺยญฺชนํ ปริภุญฺชนฺติ, เสยฺยถาปิ ตฺวํ เอตรหิ สาจริยโกติ.\csromanspacer{} โนหิทํ โภ โคตม.}

\prose{ฯเปฯ เอวํ สุ เต เวฐกนตปสฺสาหิ นารีหิ ปริจาเรนฺติ, เสยฺยถาปิ ตฺวํ เอตรหิ สาจริยโกติ.\csromanspacer{} โนหิทํ โภ โคตม.}

\prose{ฯเปฯ เอวํ สุ เต กุตฺตวาเลหิ วฬวารเถหิ ทีฆาหิ ปโตทลฏฺฐีหิ วาหเน วิตุเทนฺตา วิปริยายนฺติ, เสยฺยถาปิ ตฺวํ เอตรหิ สาจริยโกติ.\csromanspacer{} โนหิทํ โภ โคตม.}

\prose{ฯเปฯ เอวํ สุ เต อุกฺกิณฺณปริขาสุ โอกฺขิตฺตปลิฆาสุ นครูปการิกาสุ ทีฆาสิวุเธหิ\footnote{ทีฆาสิพทฺเธหิ (สฺยา, อิ)} ปุริเสหิ รกฺขาเปนฺติ, เสยฺยาถาปิ ตฺวํ เอตรหิ สาจริยโกติ.\csromanspacer{} โนหิทํ โภ โคตม.}

\prose{อิติ โข อมฺพฏฺฐ เนว ตฺวํ อิสิ น อิสิตฺถาย ปฏิปนฺโน สาจริยโก, ยสฺส โข ปน อมฺพฏฺฐ มยิ กงฺขา วา วิมติ วา โส มํ ปเญฺหน, อหํ เวยฺยากรเณน โสธิสฺสามีติ.}

\csromanmatikaanchor{mtk.68}
\csromantocmarkref{section}{เทฺวลกฺขณาทสฺสณ}{mtk.68}
\bookmark[dest={mtk.68},level=1]{เทฺวลกฺขณาทสฺสณ}
\csromanheader{1}{\textbf{เทฺวลกฺขณาทสฺสน}}
\proseitem{287}{อถ โข ภควา วิหารา นิกฺขมฺม จงฺกมํ อพฺภุฏฺฐาสิ.\csromanspacer{} อมฺพฏฺโฐปิ มาณโว วิหารา นิกฺขมฺม จงฺกมํ อพฺภุฏฺฐาสิ.\csromanspacer{} อถ โข อมฺพฏฺโฐ มาณโว ภควนฺตํ จงฺกมนฺตํ อนุจงฺกมมาโน ภควโต กาเย ทฺวตฺติํสมหาปุริสลกฺขณานิ สมนฺเนสิ, อทฺทสา โข อมฺพฏฺโฐ มาณโว ภควโต}

\vspace{\tipitakaparskip}
\csromancenter{อมฺพฏฺฐสุตฺต กาเย ทฺวตฺติํสมหาปุริสลกฺขณานิ เยภุยฺเยน ฐเปตฺวา เทฺว, ทฺวีสุ มหาปุริสลกฺขเณสุ กงฺขติ วิจิกิจฺฉติ นาธิมุจฺจติ น สมฺปสีทติ, โกโสหิเต จ วตฺถคุเยฺห ปหูตชิวฺหตาย จ.}
\proseitem{288}{\csromanfolio{99}อถ โข ภควโต เอตทโหสิ “ปสฺสติ โข เม อยํ อมฺพฏฺโฐ มาณโว ทฺวตฺติํสมหาปุริสลกฺขณานิ เยภุยฺเยน ฐเปตฺวา เทฺว, ทฺวีสุ มหาปุริสลกฺขเณสุ กงฺขติ วิจิกิจฺฉติ นาธิมุจฺจติ น สมฺปสีทติ, โกโสหิเต จ วตฺถคุเยฺห ปหูตชิวฺหตาย จา”ติ.\csromanspacer{} อถ โข ภควา ตถารูปํ อิทฺธาภิสงฺขารํ อภิสงฺขาสิ.\csromanspacer{} ยถา อทฺทส อมฺพฏฺโฐ มาณโว ภควโต โกโสหิตํ วตฺถคุยฺหํ.\csromanspacer{} อถ โข ภควา ชิวฺหํ นินฺนาเมตฺวา อุโภปิ กณฺณโสตานิ อนุมสิ ปฏิมสิ, อุโภปิ นาสิกโสตานิ อนุมสิ ปฏิมสิ, เกวลํปิ นลาฏมณฺฑลํ ชิวฺหาย ฉาเทสิ.\csromanspacer{} อถ โข อมฺพฏฺฐสฺส มาณวสฺส เอตทโหสิ “สมนฺนาคโต โข สมโณ โคตโม ทฺวตฺติํสมหาปุริสลกฺขเณหิ ปริปุณฺเณหิ โน อปริปุณฺเณหี”ติ ภควนฺตํ เอตทโวจ “หนฺท จ ทานิ มยํ โภ โคตม คจฺฉาม, พหุกิจฺจา มยํ พหุกรณียา”ติ.\csromanspacer{} ยสฺสทานิ ตฺวํ อมฺพฏฺฐ กาลํ มญฺญสีติ.\csromanspacer{} อถ โข อมฺพฏฺโฐ มาณโว วฬวารถมารุยฺห ปกฺกามิ.}

\proseitem{289}{เตน โข ปน สมเยน พฺราหฺมโณ โปกฺขรสาติ อุกฺกฏฺฐาย นิกฺขมิตฺวา มหตา พฺราหฺมณคเณน สทฺธิํ สเก อาราเม นิสินฺโน โหติ อมฺพฏฺฐํเยว มาณวํ ปฏิมาเนนฺโต.\csromanspacer{} อถ โข อมฺพฏฺโฐ มาณโว เยน สโก อาราโม เตน ปายาสิ, ยาวติกา ยานสฺส ภูมิ ยาเนน คนฺตฺวา ยานา ปจฺโจโรหิตฺวา ปตฺติโกว เยน พฺราหฺมโณ โปกฺขรสาติ เตนุปสงฺกมิ, อุปสงฺกมิตฺวา พฺราหฺมณํ โปกฺขรสาติํ อภิวาเทตฺวา เอกมนฺตํ นิสีทิ.}

\proseitem{290}{เอกมนฺตํ นิสินฺนํ โข อมฺพฏฺฐํ มาณวํ พฺราหฺมโณ โปกฺขรสาติ เอตทโวจ “กจฺจิ ตาต อมฺพฏฺฐ อทฺทส ตํ ภวนฺตํ โคตมนฺ”ติ.\csromanspacer{} อทฺทสาม โข มยํ โภ ตํ ภวนฺตํ โคตมนฺติ.\csromanspacer{} กจฺจิ ตาต อมฺพฏฺฐ ตํ ภวนฺตํ โคตมํ ตถา สนฺตํเยว สทฺโท อพฺภุคฺคโต โน อญฺญถา, กจฺจิ ปน โส ภวํ โคตโม ตาทิโส โน \csromanfolio{100}อญฺญาทิโสติ.\csromanspacer{} ตถา สนฺตํเยว โภ ตํ ภวนฺตํ โคตมํ สทฺโท อพฺภุคฺคโต โน อญฺญถา, ตาทิโสว โส ภวํ โคตโม โน อญฺญาทิโส, สมนฺนาคโต จ โส ภวํ โคตโม ทฺวตฺติํสมหาปุริสลกฺขเณหิ ปริปุณฺเณหิ โน อปริปุณฺเณหีติ.\csromanspacer{} อหุ ปน เต ตาต อมฺพฏฺฐ สมเณน โคตเมน สทฺธิํ โกจิเทว กถาสลฺลาโปติ.\csromanspacer{} อหุ โข เม โภ สมเณน โคตเมน สทฺธิํ โกจิเทว กถาสลฺลาโปติ.\csromanspacer{} ยถา กถํ ปน เต ตาต อมฺพฏฺฐ อหุ สมเณน โคตเมน สทฺธิํ โกจิเทว กถาสลฺลาโปติ.\csromanspacer{} อถ โข อมฺพฏฺโฐ มาณโว ยาวตโก\footnote{ยาวติโก (ก, อิ)} อโหสิ ภควตา สทฺธิํ กถาสลฺลาโป, ตํ สพฺพํ พฺราหฺมณสฺส โปกฺขรสาติสฺส อาโรเจสิ.}

\proseitem{291}{เอวํ วุตฺเต พฺราหฺมโณ โปกฺขรสาติ อมฺพฏฺฐํ มาณวํ เอตทโวจ “อโห วต เร อมฺหากํ ปณฺฑิตก\footnote{ปณฺฑิตกา,}, อโห วต เร อมฺหากํ พหุสฺสุตก\footnote{พหุสฺสุตกา,}, อโห วต เร อมฺหากํ เตวิชฺชก\footnote{เตวิชฺชกา (ก)}, เอวรูเปน กิร โภ ปุริโส อตฺถจรเกน กายสฺส เภทา ปรํ มรณา อปายํ ทุคฺคติํ วินิปาตํ นิรยํ อุปปชฺเชยฺย.\csromanspacer{} ยเทว โข ตฺวํ อมฺพฏฺฐ ตํ ภวนฺตํ โคตมํ เอวํ อาสชฺช อาสชฺช อวจาสิ, อถ โข โส ภวํ โคตโม อเมฺหปิ เอวํ อุปเนยฺย อุปเนยฺย อวจ.\csromanspacer{} อโห วต เร อมฺหากํ ปณฺฑิตก, อโห วต เร อมฺหากํ พหุสฺสุตก, อโห วต เร อมฺหากํ เตวิชฺชก, เอวรูเปน กิร โภ ปุริโส อตฺถจรเกน กายสฺส เภทา ปรํ มรณา อปายํ ทุคฺคติํ วินิปาตํ นิรยํ อุปปชฺเชยฺยา”ติ, กุปิโต\footnote{โส กุปิโต (อิ)} อนตฺตมโน อมฺพฏฺฐํ มาณวํ ปทสาเยว ปวตฺเตสิ.\csromanspacer{} อิจฺฉติ จ ตาวเทว ภควนฺตํ ทสฺสนาย อุปสงฺกมิตุํ.}

\csromanmatikaanchor{mtk.69}
\csromantocmarkref{section}{โปกฺขรสาติพุทฺธูปสงฺกมน}{mtk.69}
\bookmark[dest={mtk.69},level=1]{โปกฺขรสาติพุทฺธูปสงฺกมน}
\csromanheader{1}{\textbf{โปกฺขรสาติพุทฺธูปสงฺกมน}}
\proseitem{292}{อถ โข เต พฺราหฺมณา พฺราหฺมณํ โปกฺขรสาติํ เอตทโวจุํ “อติวิกาโล โข โภ อชฺช สมณํ โคตมํ ทสฺสนาย อุปสงฺกมิตุํ เสฺวทานิ\footnote{ทานิ เสฺว (สี, ก)} ภวํ โปกฺขรสาติ สมณํ โคตมํ ทสฺสานาย อุปสงฺกมิสฺสตี”ติ.\csromanspacer{} อถ โข พฺราหฺมโณ โปกฺขรสาติ สเก นิเวสเน ปณีตํ ขาทนียํ โภชนียํ ปฏิยาทาเปตฺวา ยาเน อาโรเปตฺวา อุกฺกาสุ}

\vspace{\tipitakaparskip}
\csromancenter{อมฺพฏฺฐสุตฺต ธาริยมานาสุ อุกฺกฏฺฐาย นิยฺยาสิ.\csromanspacer{} เยน อิจฺฉานงฺคลวนสณฺโฑ เตน ปายาสิ.\csromanspacer{} ยาวติกา ยานสฺส ภูมิ, ยาเนน คนฺตฺวา ยานา ปจฺโจโรหิตฺวา ปตฺติโกว เยน ภควา เตนุปสงฺกมิ, อุปสงฺกมิตฺวา ภควตา สทฺธิํ สมฺโมทิ, สมฺโมทนียํ กถํ สารณียํ วีติสาเรตฺวา เอกมนฺตํ นิสีทิ.}
\proseitem{293}{\csromanfolio{101}เอกมนฺตํ นิสินฺโน โข พฺราหฺมโณ โปกฺขรสาติ ภควนฺตํ เอตทโวจ “อาคมา นุ ขฺวิธ โภ โคตม อมฺหากํ อนฺเตวาสี อมฺพฏฺโฐ มาณโว”ติ.\csromanspacer{} อาคมา โข เต\footnote{เตธ (สฺยา), เต อิธ (อิ)} พฺราหฺมณ อนฺเตวาสี อมฺพฏฺโฐ มาณโวติ.\csromanspacer{} อหุ ปน เต โภ โคตม อมฺพฏฺเฐน มาณเวน สทฺธิํ โกจิเทว กถาสลฺลาโปติ.\csromanspacer{} อหุ โข เม พฺราหฺมณ อมฺพฏฺเฐน มาณเวน สทฺธิํ โกจิเทว กถาสลฺลาโปติ.\csromanspacer{} ยถากถํ ปน เต โภ โคตม อหุ อมฺพฏฺเฐน มาณเวน สทฺธิํ โกจิเทว กถาสลฺลาโปติ.\csromanspacer{} อถ โข ภควา ยาวตโก อโหสิ อมฺพฏฺเฐน มาณเวน สทฺธิํ กถาสลฺลาโป, ตํ สพฺพํ พฺราหฺมณสฺส โปกฺขรสาติสฺส อาโรเจสิ.\csromanspacer{} เอวํ วุตฺเต พฺราหฺมโณ โปกฺขรสาติ ภควนฺตํ เอตทโวจ “พาโล โภ โคตม อมฺพฏฺโฐ มาณโว, ขมตุ ภวํ โคตโม อมฺพฏฺฐสฺส มาณวสฺสา”ติ.\csromanspacer{} สุขี โหตุ พฺราหฺมณ อมฺพฏฺโฐ มาณโวติ.}

\proseitem{294}{อถ โข พฺราหฺมโณ โปกฺขรสาติ ภควโต กาเย ทฺวตฺติํสมหาปุริสลกฺขณานิ สมนฺเนสิ.\csromanspacer{} อทฺทสา โข พฺราหฺมโณ โปกฺขรสาติ ภควโต กาเย ทฺวตฺติํสมหาปุริสลกฺขณานิ เยภุยฺเยน ฐเปตฺวา เทฺว, ทฺวีสุ มหาปุริสลกฺขเณสุ กงฺขติ วิจิกิจฺฉติ นาธิมุจฺจติ น สมฺปสีทติ โกโสหิเต จ วตฺถคุเยฺห ปหูตชิวฺหตาย จ.}

\proseitem{295}{อถ โข ภควโต เอตทโหสิ “ปสฺสติ โข เม อยํ พฺราหฺมโณ โปกฺขรสาติ ทฺวตฺติํสมหาปุริสลกฺขณานิ เยภุยฺเยน ฐเปตฺวา เทฺว, ทฺวีสุ มหาปุริสลกฺขเณสุ กงฺขติ วิจิกิจฺฉติ นาธิมุจฺจติ น สมฺปสีทติ โกโสหิเต จ วตฺถคุเยฺห ปหูตชิวฺหตาย จา”ติ.\csromanspacer{} อถ โข ภควา ตถารูปํ อิทฺธาภิสงฺขารํ อภิสงฺขาสิ.\csromanspacer{} ยถา อทฺทส พฺราหฺมโณ โปกฺขรสาติ ภควโต โกโสหิตํ วตฺถคุยฺหํ.\csromanspacer{} อถ \csromanfolio{102}โข ภควา ชิวฺหํ นินฺนาเมตฺวา อุโภปิ กณฺณโสตานิ อนุมสิ ปฏิมสิ, อุโภปิ นาสิกโสตานิ อนุมสิ ปฏิมสิ, เกวลํปิ นลาฏมณฺฑลํ ชิวฺหาย ฉาเทสิ.}

\proseitem{296}{อถ โข พฺราหฺมณสฺส โปกฺขรสาติสฺส เอตทโหสิ “สมนฺนาคโต โข สมโณ โคตโม ทฺวตฺติํสมหาปุริสลกฺขเณหิ ปริปุณฺเณหิ โน อปริปุณฺเณหี”ติ ภควนฺตํ เอตทโวจ “อธิวาเสตุ เม ภวํ โคตโม อชฺชตนาย ภตฺตํ สทฺธิํ ภิกฺขุสํเฆนา”ติ.\csromanspacer{} อธิวาเสสิ ภควา ตุณฺหีภาเวน.}

\proseitem{297}{อถ โข พฺราหฺมโณ โปกฺขรสาติ ภควโต อธิวาสนํ วิทิตฺวา ภควโต กาลํ อาโรเจสิ “กาโล โภ โคตม นิฏฺฐิตํ ภตฺตนฺ”ติ.\csromanspacer{} อถ โข ภควา ปุพฺพณฺหสมยํ นิวาเสตฺวา ปตฺตจีวรมาทาย สทฺธิํ ภิกฺขุสํเฆน เยน พฺราหฺมณสฺส โปกฺขรสาติสฺส นิเวสนํ เตนุปสงฺกมิ, อุปสงฺกมิตฺวา ปญฺญตฺเต อาสเน นิสีทิ.\csromanspacer{} อถ โข พฺราหฺมโณ โปกฺขรสาติ ภควนฺตํ ปณีเตน ขาทนีเยน โภชนีเยน สหตฺถา สนฺตปฺเปสิ สมฺปวาเรสิ, มาณวกาปิ ภิกฺขุสํฆํ.\csromanspacer{} อถ โข พฺราหฺมโณ โปกฺขรสาติ ภควนฺตํ ภุตฺตาวิํ โอนีตปตฺตปาณิํ อญฺญตรํ นีจํ อาสนํ คเหตฺวา เอกมนฺตํ นิสีทิ.}

\proseitem{298}{เอกมนฺตํ นิสินฺนสฺส โข พฺราหฺมณสฺส โปกฺขรสาติสฺส ภควา อนุปุพฺพิํ กถํ กเถสิ, เสยฺยถิทํ, ทานกถํ สีลกถํ สคฺคกถํ กามานํ อาทีนวํ โอการํ สํกิเลสํ เนกฺขมฺเม อานิสํสํ ปกาเสสิ.\csromanspacer{} ยทา ภควา อญฺญาสิ พฺราหฺมณํ โปกฺขรสาติํ กลฺลจิตฺตํ มุทุจิตฺตํ วินีวรณจิตฺตํ อุทคฺคจิตฺตํ ปสนฺนจิตฺตํ, อถ ยา พุทฺธานํ สามุกฺกํสิกา ธมฺมเทสนา, ตํ ปกาเสสิ ทุกฺขํ สมุทยํ นิโรธํ มคฺคํ, เสยฺยถาปิ นาม สุทฺธํ วตฺถํ อปคตกาฬกํ สมฺมเทว รชนํ ปฏิคฺคเณฺหยฺย, เอวเมว พฺราหฺมณสฺส โปกฺขรสาติสฺส ตสฺมิํเยว อาสเน วิรชํ วีตมลํ ธมฺมจกฺขุํ อุทปาทิ “ยํ กิญฺจิ สมุทยธมฺมํ, สพฺพํ ตํ นิโรธธมฺมนฺ”ติ.}

\csromanmatikaanchor{mtk.70}
\csromantocmarkref{section}{โปกฺขรสาติอุปาสกตฺตปฏิเวทนา}{mtk.70}
\bookmark[dest={mtk.70},level=1]{โปกฺขรสาติอุปาสกตฺตปฏิเวทนา}
\csromanheader{1}{โปกฺขรสาติ\-อุปาสกตฺต\-ปฏิเวทนา}
\proseitem{299}{อถ โข พฺราหฺมโณ โปกฺขรสาติ ทิฏฺฐธมฺโม ปตฺตธมฺโม วิทิตธมฺโม ปริโยคาฬฺหธมฺโม ติณฺณวิจิกิจฺโฉ วิคตกถํกโถ}

\vspace{\tipitakaparskip}
\csromancenter{อมฺพฏฺฐสุตฺต เวสารชฺชปฺปตฺโต อปรปฺปจฺจโย สตฺถุสาสเน ภควนฺตํ เอตทโวจ “อภิกฺกนฺตํ โภ โคตม, อภิกฺกนฺตํ โภ โคตม, เสยฺยถาปิ โภ โคตม นิกฺกุชฺชิตํ วา อุกฺกุชฺเชยฺย, ปฏิจฺฉนฺนํ วา วิวเรยฺย, มูฬฺหสฺส วา มคฺคํ อาจิกฺเขยฺย, อนฺธกาเร วา เตลปชฺโชตํ ธาเรยฺย ‘จกฺขุมนฺโต รูปานิ ทกฺขนฺตี’ติ, เอวเมวํ โภตา โคตเมน อเนกปริยาเยน ธมฺโม ปกาสิโต.\csromanspacer{} เอสาหํ โภ โคตม สปุตฺโต สภริโย สปริโส สามจฺโจ ภวนฺตํ โคตมํ สรณํ คจฺฉามิ ธมฺมญฺจ ภิกฺขุสํฆญฺจ, อุปาสกํ มํ ภวํ โคตโม ธาเรตุ อชฺชตคฺเค ปาณุเปตํ สรณํ คตํ, ยถา จ ภวํ โคตโม อุกฺกฏฺฐาย อญฺญานิ อุปาสกกุลานิ อุปสงฺกมติ, เอวเมว ภวํ โคตโม โปกฺขรสาติกุลํ อุปสงฺกมตุ.\csromanspacer{} ตตฺถ เย เต มาณวกา วา มาณวิกา วา ภวนฺตํ โคตมํ อภิวาเทสฺสนฺติ วา ปจฺจุฏฺฐิสฺสนฺติ\footnote{ปจฺจุฏฺฐสฺสนฺติ (อิ)} วา อาสนํ วา อุทกํ วา ทสฺสนฺติ จิตฺตํ วา ปสาเทสฺสนฺติ, เตสํ ตํ ภวิสฺสติ ทีฆรตฺตํ หิตาย สุขายา”ติ.\csromanspacer{} กลฺยาณํ วุจฺจติ พฺราหฺมณาติ.}
\nitthitam{อมฺพฏฺฐสุตฺตํ นิฏฺฐิตํ ตติยํ.}
\csromanmatikaanchor{mtk.71}
\csromantocmarknopage{chapter}{4. โสณทณฺฑสุตฺต}{mtk.71}
\bookmark[dest={mtk.71},level=0]{4. โสณทณฺฑสุตฺต}
\setcsromanheadcha{4. โสณทณฺฑสุตฺต}
\setcsromanheadhclear{}
\chapterheadpageread{4. \textbf{โสณทณฺฑสุตฺต}}
\csromanmatikaanchor{mtk.72}
\csromantocmarkref{section}{จมฺเปยฺยกพฺราหฺมณคหปติกา}{mtk.72}
\bookmark[dest={mtk.72},level=1]{จมฺเปยฺยกพฺราหฺมณคหปติกา}
\csromanheader{1}{จมฺเปยฺยก\-พฺราหฺมณคหปติกา}
\proseitem{300}{\csromanfolio{104}เอวํ เม สุตํ\csromandash{}เอกํ สมยํ ภควา องฺเคสุ จาริกํ จรมาโน มหตา ภิกฺขุสํเฆน สทฺธิํ ปญฺจมตฺเตหิ ภิกฺขุสเตหิ เยน จมฺปา ตทวสริ.\csromanspacer{} ตตฺร สุทํ ภควา จมฺปายํ วิหรติ คคฺคราย โปกฺขรณิยา ตีเร.\csromanspacer{} เตน โข ปน สมเยน โสณทณฺโฑ พฺราหฺมโณ จมฺปํ อชฺฌาวสติ สตฺตุสฺสทํ สติณกฏฺโฐทกํ สธญฺญํ ราชโภคฺคํ รญฺญา มาคเธน เสนิเยน พิมฺพิสาเรน ทินฺนํ ราชทายํ พฺรหฺมเทยฺยํ.}

\proseitem{301}{อสฺโสสุํ โข จมฺเปยฺยกา พฺราหฺมณคหปติกา “สมโณ ขลุ โภ โคตโม สกฺยปุตฺโต สกฺยกุลา ปพฺพชิโต องฺเคสุ จาริกํ จรมาโน มหตา ภิกฺขุสํเฆน สทฺธิํ ปญฺจมตฺเตหิ ภิกฺขุสเตหิ จมฺปํ อนุปฺปตฺโต จมฺปายํ วิหรติ คคฺคราย โปกฺขรณิยา ตีเร.\csromanspacer{} ตํ โข ปน ภวนฺตํ โคตมํ เอวํ กลฺยาโณ กิตฺติสทฺโท อพฺภุคฺคโต ‘อิติปิ โส ภควา อรหํ สมฺมาสมฺพุทฺโธ วิชฺชาจรณสมฺปนฺโน สุคโต โลกวิทู อนุตฺตโร ปุริสทมฺมสารถิ สตฺถา เทวมนุสฺสานํ พุทฺโธ ภควา’, โส อิมํ โลกํ สเทวกํ สมารกํ สพฺรหฺมกํ สสฺสมณพฺราหฺมณิํ ปชํ สเทวมนุสฺสํ สยํ อภิญฺญา สจฺฉิกตฺวา ปเวเทติ, โส ธมฺมํ เทเสติ อาทิกลฺยาณํ มชฺเฌกลฺยาณํ ปริโยสานกลฺยาณํ สาตฺถํ สพฺยญฺชนํ เกวลปริปุณฺณํ ปริสุทฺธํ พฺรหฺมจริยํ ปกาเสติ, สาธุ โข ปน ตถารูปานํ อรหตํ ทสฺสนํ โหตี”ติ.\csromanspacer{} อถ โข จมฺเปยฺยกา พฺราหฺมณคหปติกา จมฺปาย นิกฺขมิตฺวา สงฺฆสงฺฆี\footnote{สงฺฆา สงฺฆี (สี, สฺยา, อิ)} คณีภูตา เยน คคฺครา โปกฺขรณี เตนุปสงฺกมนฺติ.}

\proseitem{302}{เตน โข ปน สมเยน โสณทณฺโฑ พฺราหฺมโณ อุปริปาสาเท ทิวาเสยฺยํ อุปคโต โหติ.\csromanspacer{} อทฺทสา โข โสณทณฺโฑ พฺราหฺมโณ จมฺเปยฺยเก พฺราหฺมณคหปติเก จมฺปาย นิกฺขมิตฺวา สงฺฆสงฺฆี\footnote{สงฺเฆ สงฺฆี (สี, อิ) สงฺฆา สงฺฆี (สฺยา)} คณีภูเต เยน คคฺครา โปกฺขรณี เตนุปสงฺกมนฺเต, ทิสฺวา ขตฺตํ อามนฺเตสิ “กิํ นุ โข โภ ขตฺเต จมฺเปยฺยกา พฺราหฺมณคหปติกา}

\vspace{\tipitakaparskip}
\csromancenter{โสณทณฺฑสุตฺต จมฺปาย นิกฺขมิตฺวา สงฺฆสงฺฆี คณีภูตา เยน คคฺครา โปกฺขรณี เตนุปสงฺกมนฺตี”ติ.\csromanspacer{} อตฺถิ โข โภ สมโณ โคตโม สกฺยปุตฺโต สกฺยกุลา ปพฺพชิโต องฺเคสุ จาริกํ จรมาโน มหตา ภิกฺขุสํเฆน สทฺธิํ ปญฺจมตฺเตหิ ภิกฺขุสเตหิ จมฺปํ อนุปฺปตฺโต จมฺปายํ วิหรติ คคฺคราย โปกฺขรณิยา ตีเร.\csromanspacer{} ตํ โข ปน ภวนฺตํ โคตมํ เอวํ กลฺยาโณ กิตฺติสทฺโท อพฺภุคฺคโต “อิติปิ โส ภควา อรหํ สมฺมาสมฺพุทฺโธ วิชฺชาจรณสมฺปนฺโน สุคโต โลกวิทู อนุตฺตโร ปุริสทมฺมสารถิ สตฺถา เทวมนุสฺสานํ พุทฺโธ ภควา”ติ.\csromanspacer{} ตเมเต ภวนฺตํ โคตมํ ทสฺสนาย อุปสงฺกมนฺตีติ.\csromanspacer{} เตน หิ โภ ขตฺเต เยน จมฺเปยฺยกา พฺราหฺมณคหปติกา เตนุปสงฺกม, อุปสงฺกมิตฺวา จมฺเปยฺยเก พฺราหฺมณคหปติเก เอวํ วเทหิ “โสณทณฺโฑ โภ พฺราหฺมโณ เอวมาห อาคเมนฺตุ กิร ภวนฺโต, โสณทณฺโฑปิ พฺราหฺมโณ สมณํ โคตมํ ทสฺสนาย อุปสงฺกมิสฺสตี”ติ. “เอวํ โภ”ติ โข โส ขตฺตา โสณทณฺฑสฺส พฺราหฺมณสฺส ปฏิสฺสุตฺวา เยน จมฺเปยฺยกา พฺราหฺมณคหปติกา เตนุปสงฺกมิ, อุปสงฺกมิตฺวา จมฺเปยฺยเก พฺราหฺมณคหปติเก เอตทโวจ “โสณทณฺโฑ โภ พฺราหฺมโณ เอวมาห อาคเมนฺตุ กิร ภวนฺโต, โสณทณฺโฑปิ พฺราหฺมโณ สมณํ โคตมํ ทสฺสนาย อุปสงฺกมิสฺสตี”ติ.}
\csromanmatikaanchor{mtk.73}
\csromantocmarkref{section}{โสณทณฺฑคุณกถา}{mtk.73}
\bookmark[dest={mtk.73},level=1]{โสณทณฺฑคุณกถา}
\csromanheader{1}{\textbf{โสณทณฺฑคุณกถา}}
\proseitem{303}{\csromanfolio{105}เตน โข ปน สมเยน นานาเวรชฺชกานํ พฺราหฺมณานํ ปญฺจมตฺตานิ พฺราหฺมณสตานิ จมฺปายํ ปฏิวสนฺติ เกนจิเทว กรณีเยน, อสฺโสสุํ โข เต พฺราหฺมณา “โสณทณฺโฑ กิร พฺราหฺมโณ สมณํ โคตมํ ทสฺสนาย อุปสงฺกมิสฺสตี”ติ.\csromanspacer{} อถ โข เต พฺราหฺมณา เยน โสณทณฺโฑ พฺราหฺมโณ เตนุปสงฺกมิํสุ, อุปสงฺกมิตฺวา โสณทณฺฑํ พฺราหฺมณํ เอตทโวจุํ “สจฺจํ กิร ภวํ โสณทณฺโฑ สมณํ โคตมํ ทสฺสนาย อุปสงฺกมิสฺสตี”ติ.\csromanspacer{} เอวํ โข เม โภ โหติ “อหมฺปิ สมณํ โคตมํ ทสฺสนาย อุปสงฺกมิสฺสามี”ติ.}

\prose{มา ภวํ โสณทณฺโฑ สมณํ โคตมํ ทสฺสนาย อุปสงฺกมิ, น อรหติ ภวํ โสณทณฺโฑ สมณํ โคตมํ ทสฺสนาย อุปสงฺกมิตุํ, สเจ ภวํ โสณทณฺโฑ สมณํ โคตมํ ทสฺสนาย อุปสงฺกมิสฺสติ, \csromanfolio{106}โภโต โสณทณฺฑสฺส ยโส หายิสฺสติ, สมณสฺส โคตมสฺส ยโส อภิวฑฺฒิสฺสติ.\csromanspacer{} ยมฺปิ โภโต โสณทณฺฑสฺส ยโส หายิสฺสติ, สมณสฺส โคตมสฺส ยโส อภิวฑฺฒิสฺสติ, อิมินาปงฺเคน น อรหติ ภวํ โสณทณฺโฑ สมณํ โคตมํ ทสฺสนาย อุปสงฺกมิตุํ, สมโณเตฺวว โคตโม อรหติ ภวนฺตํ โสณทณฺฑํ ทสฺสนาย อุปสงฺกมิตุํ.}

\prose{ภวํ หิ โสณทณฺโฑ อุภโต สุชาโต มาติโต จ ปิติโต จ สํสุทฺธคหณิโก ยาว สตฺตมา ปิตามหยุคา อกฺขิตฺโต อนุปกฺกุฏฺโฐ ชาติวาเทน.\csromanspacer{} ยมฺปิ ภวํ โสณทณฺโฑ อุภโต สุชาโต มาติโต จ ปิติโต จ สํสุทฺธคหณิโก ยาว สตฺตมา ปิตามหยุคา อกฺขิตฺโต อนุปกฺกุฏฺโฐ ชาติวาเทน, อิมินาปงฺเคน น อรหติ ภวํ โสณทณฺโฑ สมณํ โคตมํ ทสฺสนาย อุปสงฺกมิตุํ, สมโณเตฺวว โคตโม อรหติ ภวนฺตํ โสณทณฺฑํ ทสฺสนาย อุปสงฺกมิตุํ.}

\prose{ภวํ หิ โสณทณฺโฑ อฑฺโฒ มหทฺธโน มหาโภโค ฯเปฯ.}

\prose{ภวํ หิ โสณทณฺโฑ อชฺฌายโก มนฺตธโร ติณฺณํ เวทานํ ปารคู สนิฆณฺฑุเกฏุภานํ สากฺขรปฺปเภทานํ อิติหาสปญฺจมานํ ปทโก เวยฺยากรโณ โลกายตมหาปุริสลกฺขเณสุ อนวโย ฯเปฯ.}

\prose{ภวํ หิ โสณทณฺโฑ อภิรูโป ทสฺสนีโย ปาสาทิโก ปรมาย วณฺณโปกฺขรตาย สมนฺนาคโต พฺรหฺมวณฺณี พฺรหฺมวจฺฉสี\footnote{พฺรหฺมฑฺฒี (สี), พฺรหฺมวจฺจสี (อิ)} อขุทฺทาวกาโส ทสฺสนาย ฯเปฯ.}

\prose{ภวํ หิ โสณทณฺโฑ สีลวา วุทฺธสีลี วุทฺธสีเลน สมนฺนาคโต ฯเปฯ.}

\prose{ภวํ หิ โสณทณฺโฑ กลฺยาณวาโจ กลฺยาณวากฺกรโณ โปริยา วาจาย สมนฺนาคโต วิสฺสฏฺฐาย อเนลคลาย\footnote{อเนฬคลาย (สี, อิ), อเนลคฬาย (ก)} อตฺถสฺส วิญฺญาปนิยา ฯเปฯ.}

\prose{ภวํ หิ โสณทณฺโฑ พหูนํ อาจริยปาจริโย ตีณิ มาณวกสตานิ มนฺเต วาเจติ, พหู โข ปน นานาทิสา นานาชนปทา มาณวกา}

\vspace{\tipitakaparskip}
\csromancenter{โสณทณฺฑสุตฺต อาคจฺฉนฺติ โภโต โสณทณฺฑสฺส สนฺติเก มนฺตตฺถิกา มนฺเต อธิยิตุกามา ฯเปฯ.}
\prose{\csromanfolio{107}ภวํ หิ โสณทณฺโฑ ชิณฺโณ วุทฺโธ มหลฺลโก อทฺธคโต วโย- อนุปฺปตฺโต, สมโณ โคตโม ตรุโณ เจว ตรุณปพฺพชิโต จ ฯเปฯ.}

\prose{ภวํ หิ โสณทณฺโฑ รญฺโญ มาคธสฺส เสนิยสฺส พิมฺพิสารสฺส สกฺกโต ครุกโต มานิโต ปูชิโต อปจิโต ฯเปฯ.}

\prose{ภวํ หิ โสณทณฺโฑ พฺราหฺมณสฺส โปกฺขรสาติสฺส สกฺกโต ครุกโต มานิโต ปูชิโต อปจิโต ฯเปฯ.}

\prose{ภวํ หิ โสณทณฺโฑ จมฺปํ อชฺฌาวสติ สตฺตุสฺสทํ สติณกฏฺโฐทกํ สธญฺญํ ราชโภคฺคํ รญฺญา มาคเธน เสนิเยน พิมฺพิสาเรน ทินฺนํ ราชทายํ พฺรหฺมเทยฺยํ.\csromanspacer{} ยมฺปิ ภวํ โสณทณฺโฑ จมฺปํ อชฺฌาวสติ สตฺตุสฺสทํ สติณกฏฺโฐทกํ สธญฺญํ ราชโภคฺคํ รญฺญา มาคเธน เสนิเยน พิมฺพิสาเรน ทินฺนํ ราชทายํ พฺรหฺมเทยฺยํ, อิมินาปงฺเคน น อรหติ ภวํ โสณทณฺโฑ สมณํ โคตมํ ทสฺสนาย อุปสงฺกมิตุํ, สมโณเตฺวว โคตโม อรหติ ภวนฺตํ โสณทณฺฑํ ทสฺสนาย อุปสงฺกมิตุนฺติ.}

\csromanmatikaanchor{mtk.74}
\csromantocmarkref{section}{พุทฺธคุณกถา}{mtk.74}
\bookmark[dest={mtk.74},level=1]{พุทฺธคุณกถา}
\csromanheader{1}{\textbf{พุทฺธคุณกถา}}
\csromanheader{2}{304. เอวํ วุตฺเต โสณทณฺโฑ พฺราหฺมโณ เต พฺราหฺมเณ เอตทโวจ\csromandash{}}
\prose{“เตน หิ โภ มมปิ สุณาถ, ยถา มยเมว อรหาม ตํ ภวนฺตํ โคตมํ ทสฺสนาย อุปสงฺกมิตุํ, นเตฺวว อรหติ โส ภวํ โคตโม อมฺหากํ ทสฺสนาย อุปสงฺกมิตุํ.\csromanspacer{} สมโณ ขลุ โภ โคตโม อุภโต สุชาโต มาติโต จ ปิติโต จ สํสุทฺธคหณิโก ยาว สตฺตมา ปิตามหยุคา อกฺขิตฺโต อนุปกฺกุฏฺโฐ ชาติวาเทน.\csromanspacer{} ยมฺปิ โภ สมโณ โคตโม อุภโต สุชาโต มาติโต จ ปิติโต จ สํสุทฺธคหณิโก ยาว สตฺตมา ปิตามหยุคา อกฺขิตฺโต อนุปกฺกุฏฺโฐ ชาติวาเทน, อิมินาปงฺเคน น อรหติ โส ภวํ โคตโม อมฺหากํ ทสฺสนาย}

\prose{\csromanfolio{108}อุปสงฺกมิตุํ, อถ โข มยเมว อรหาม ตํ ภวนฺตํ โคตมํ ทสฺสนาย อุปสงฺกมิตุํ.}

\prose{สมโณ ขลุ โภ โคตโม มหนฺตํ ญาติสํฆํ โอหาย ปพฺพชิโต ฯเปฯ.}

\prose{สมโณ ขลุ โภ โคตโม ปหูตํ หิรญฺญสุวณฺณํ โอหาย ปพฺพชิโต ภูมิคตญฺจ เวหาสฏฺฐญฺจ ฯเปฯ.}

\prose{สมโณ ขลุ โภ โคตโม ทหโรว สมาโน ยุวา สุสุกาฬเกโส ภเทฺรน โยพฺพเนน สมนฺนาคโต ปฐเมน วยสา อคารสฺมา อนคาริยํ ปพฺพชิโต ฯเปฯ.}

\prose{สมโณ ขลุ โภ โคตโม อกามกานํ มาตาปิตูนํ อสฺสุมุขานํ รุทนฺตานํ เกสมสฺสุํ โอหาเรตฺวา กาสายานิ วตฺถานิ อจฺฉาเทตฺวา อคารสฺมา อนคาริยํ ปพฺพชิโต ฯเปฯ.}

\prose{สมโณ ขลุ โภ โคตโม อภิรูโป ทสฺสนีโย ปาสาทิโก ปรมาย วณฺณโปกฺขรตาย สมนฺนาคโต พฺรหฺมวณฺณี พฺรหฺมวจฺฉสี อขุทฺทาวกาโส ทสฺสนาย ฯเปฯ.}

\prose{สมโณ ขลุ โภ โคตโม สีลวา อริยสีลี กุสลสีลี กุสลสีเลน สมนฺนาคโต ฯเปฯ.}

\prose{สมโณ ขลุ โภ โคตโม กลฺยาณวาโจ กลฺยาณวากฺกรโณ โปริยา วาจาย สมนฺนาคโต วิสฺสฏฺฐาย อเนลคลาย อตฺถสฺส วิญฺญาปนิยา ฯเปฯ.}

\prose{สมโณ ขลุ โภ โคตโม พหูนํ อาจริยปาจริโย ฯเปฯ.}

\prose{สมโณ ขลุ โภ โคตโม ขีณกามราโค วิคตจาปลฺโล ฯเปฯ.}

\prose{สมโณ ขลุ โภ โคตโม กมฺมวาที กิริยวาที อปาปปุเรกฺขาโร พฺรหฺมญฺญาย ปชาย ฯเปฯ.}

\prose{สมโณ ขลุ โภ โคตโม อุจฺจา กุลา ปพฺพชิโต อสมฺภินฺนขตฺติยกุลา ฯเปฯ.}

\prose{สมโณ ขลุ โภ โคตโม อฑฺฒา กุลา ปพฺพชิโต มหทฺธนา มหาโภคา ฯเปฯ.}

\csromanheader{2}{4. โสณทณฺฑสุตฺต}
\prose{\csromanfolio{109}สมณํ ขลุ โภ โคตมํ ติโรรฏฺฐา ติโรชนปทา ปญฺหํ ปุจฺฉิตุํ อาคจฺฉนฺติ ฯเปฯ.}

\prose{สมณํ ขลุ โภ โคตมํ อเนกานิ เทวตาสหสฺสานิ ปาเณหิ สรณํ คตานิ ฯเปฯ.}

\prose{สมณํ ขลุ โภ โคตมํ เอวํ กลฺยาโณ กิตฺติสทฺโท อพฺภุคฺคโต ‘อิติปิ โส ภควา อรหํ สมฺมาสมฺพุทฺโธ วิชฺชาจรณสมฺปนฺโน สุคโต โลกวิทู อนุตฺตโร ปุริสทมฺมสารถิ สตฺถา เทวมนุสฺสานํ พุทฺโธ ภควา’ติ ฯเปฯ.}

\prose{สมโณ ขลุ โภ โคตโม ทฺวตฺติํสมหาปุริสลกฺขเณหิ สมนฺนาคโต ฯเปฯ.}

\prose{สมโณ ขลุ โภ โคตโม เอหิสฺวาคตวาที สขิโล สมฺโมทโก อพฺภากุฏิโก อุตฺตานมุโข ปุพฺพภาสี ฯเปฯ.}

\prose{สมโณ ขลุ โภ โคตโม จตุนฺนํ ปริสานํ สกฺกโต ครุกโต มานิโต ปูชิโต อปจิโต ฯเปฯ.}

\prose{สมเณ ขลุ โภ โคตเม พหู เทวา จ มนุสฺสา จ อภิปฺปสนฺนา ฯเปฯ.}

\prose{สมโณ ขลุ โภ โคตโม ยสฺมิํ คาเม วา นิคเม วา ปฏิวสติ, น ตสฺมิํ คาเม วา นิคเม วา อมนุสฺสา มนุสฺเส วิเหเฐนฺติ ฯเปฯ.}

\prose{สมโณ ขลุ โภ โคตโม สงฺฆี คณี คณาจริโย ปุถุติตฺถกรานํ อคฺคมกฺขายติ, ยถา โข ปน โภ เอเตสํ สมณพฺราหฺมณานํ ยถา วา ตถา วา ยโส สมุทาคจฺฉติ, น เหวํ สมณสฺส โคตมสฺส ยโส สมุทาคโต, อถ โข อนุตฺตราย วิชฺชาจรณสมฺปทาย สมณสฺส โคตมสฺส ยโส สมุทาคโต ฯเปฯ.}

\prose{สมณํ ขลุ โภ โคตมํ ราชา มาคโธ เสนิโย พิมฺพิสาโร สปุตฺโต สภริโย สปริโส สามจฺโจ ปาเณหิ สรณํ คโต ฯเปฯ.}

\prose{\csromanfolio{110}สมณํ ขลุ โภ โคตมํ ราชา ปเสนทิ โกสโล สปุตฺโต สภริโย สปริโส สามจฺโจ ปาเณหิ สรณํ คโต ฯเปฯ.}

\prose{สมณํ ขลุ โคตมํ พฺราหฺมโณ โปกฺขรสาติ สปุตฺโต สภริโย สปริโส สามจฺโจ ปาเณหิ สรณํ คโต ฯเปฯ.}

\prose{สมโณ ขลุ โภ โคตโม รญฺโญ มาคธสฺส เสนิยสฺส พิมฺพิสารสฺส สกฺกโต ครุกโต มานิโต ปูชิโต อปจิโต ฯเปฯ.}

\prose{สมโณ ขลุ โภ โคตโม รญฺโญ ปเสนทิสฺส โกสลสฺส สกฺกโต ครุกโต มานิโต ปูชิโต อปจิโต ฯเปฯ.}

\prose{สมโณ ขลุ โภ โคตโม พฺราหฺมณสฺส โปกฺขรสาติสฺส สกฺกโต ครุกโต มานิโต ปูชิโต อปจิโต ฯเปฯ.}

\prose{สมโณ ขลุ โภ โคตโม จมฺปํ อนุปฺปตฺโต จมฺปายํ วิหรติ คคฺคราย โปกฺขรณิยา ตีเร, เย โข ปน โภ เกจิ สมณา วา พฺราหฺมณา วา อมฺหากํ คามเขตฺตํ อาคจฺฉนฺติ, อติถี โน เต โหนฺติ, อติถี โข ปนเมฺหหิ สกฺกาตพฺพา ครุกาตพฺพา มาเนตพฺพา ปูเชตพฺพา อปเจตพฺพา.\csromanspacer{} ยมฺปิ โภ สมโณ โคตโม จมฺปํ อนุปฺปตฺโต จมฺปายํ วิหรติ คคฺคราย โปกฺขรณิยา ตีเร, อติถิมฺหากํ สมโณ โคตโม, อติถิ โข ปนเมฺหหิ สกฺกาตพฺโพ ครุกาตพฺโพ มาเนตพฺโพ ปูเชตพฺโพ อปเจตพฺโพ, อิมินาปงฺเคน น อรหติ โส ภวํ โคตโม อมฺหากํ ทสฺสนาย อุปสงฺกมิตุํ, อถ โข มยเมว อรหาม ตํ ภวนฺตํ โคตมํ ทสฺสนาย อุปสงฺกมิตุํ.\csromanspacer{} เอตฺตเก โข อหํ โภ ตสฺส โภโต โคตมสฺส วณฺเณ ปริยาปุณามิ, โน จ โข โส ภวํ โคตโม เอตฺตกวณฺโณ, อปริมาณวณฺโณ หิ โส ภวํ โคตโม”ติ.}

\proseitem{305}{เอวํ วุตฺเต เต พฺราหฺมณา โสณทณฺฑํ พฺราหฺมณํ เอตทโวจุํ “ยถา โข ภวํ โสณทณฺโฑ สมณสฺส โคตมสฺส วณฺเณ ภาสติ.\csromanspacer{} อิโต เจปิ โส ภวํ โคตโม โยชนสเต วิหรติ, อลเมว สทฺเธน กุลปุตฺเตน ทสฺสนาย อุปสงฺกมิตุํ อปิ ปุโฏเสนา”ติ.\csromanspacer{} เตน หิ โภ สพฺเพว มยํ สมณํ โคตมํ ทสฺสนาย อุปสงฺกมิสฺสามาติ.}

\csromanheader{2}{4. โสณทณฺฑสุตฺต \textbf{โสณทณฺฑ ปริวิตกฺก}}
\csromanmatikaanchor{mtk.75}
\csromantocmarkref{section}{โสณทณฺฑปริวิตกฺก}{mtk.75}
\bookmark[dest={mtk.75},level=1]{โสณทณฺฑปริวิตกฺก}
\proseitem{306}{\csromanfolio{111}อถ โข โสณทณฺโฑ พฺราหฺมโณ มหตา พฺราหฺมณคเณน สทฺธิํ เยน คคฺคารา โปกฺขรณี เตนุปสงฺกมิ.\csromanspacer{} อถ โข โสณทณฺฑสฺส พฺราหฺมณสฺส ติโรวนสณฺฑคตสฺส เอวํ เจตโส ปริวิตกฺโก อุทปาทิ “อหญฺเจว โข ปน สมณํ โคตมํ ปญฺหํ ปุจฺเฉยฺยํ, ตตฺร เจ มํ สมโณ โคตโม เอวํ วเทยฺย ‘น โข เอส พฺราหฺมณ ปโญฺห เอวํ ปุจฺฉิตพฺโพ, เอวํ นาเมส พฺราหฺมณ ปโญฺห ปุจฺฉิตพฺโพ’ติ, เตน มํ อยํ ปริสา ปริภเวยฺย ‘พาโล โสณทณฺโฑ พฺราหฺมโณ อพฺยตฺโต, นาสกฺขิ สมณํ โคตมํ โยนิโส ปญฺหํ ปุจฺฉิตุนฺ’ติ.\csromanspacer{} ยํ โข ปนายํ ปริสา ปริภเวยฺย, ยโสปิ ตสฺส หาเยถ.\csromanspacer{} ยสฺส โข ปน ยโส หาเยถ, โภคาปิ ตสฺส หาเยยฺยุํ, ยโสลทฺธา โข ปน อมฺหากํ โภคา.\csromanspacer{} มมญฺเจว โข ปน สมโณ โคตโม ปญฺหํ ปุจฺเฉยฺย, ตสฺส จาหํ ปญฺหสฺส เวยฺยากรเณน จิตฺตํ น อาราเธยฺยํ, ตตฺร เจ มํ สมโณ โคตโม เอวํ วเทยฺย ‘น โข เอส พฺราหฺมณ ปโญฺห เอวํ พฺยากาตพฺโพ, เอวํ นาเมส พฺราหฺมณ ปโญฺห พฺยากาตพฺโพ’ติ, เตน มํ อยํ ปริสา ปริภเวยฺย ‘พาโล โสณทณฺโฑ พฺราหฺมโณ อพฺยตฺโต, นาสกฺขิ สมณสฺส โคตมสฺส ปญฺหสฺส เวยฺยากรเณน จิตฺตํ อาราเธตุนฺ’ติ.\csromanspacer{} ยํ โข ปนายํ ปริสา ปริภเวยฺย, ยโสปิ ตสฺส หาเยถ.\csromanspacer{} ยสฺส โข ปน ยโส หาเยถ, โภคาปิ ตสฺส หาเยยฺยุํ, ยโสลทฺธา โข ปน อมฺหากํ โภคา.\csromanspacer{} อหญฺเจว โข ปน เอวํ สมีปคโต สมาโน อทิสฺวาว สมณํ โคตมํ นิวตฺเตยฺยํ, เตน มํ อยํ ปริสา ปริภเวยฺย ‘พาโล โสณทณฺโฑ พฺราหฺมโณ อพฺยตฺโต มานถทฺโธ ภีโต จ, โน วิสหติ สมณํ โคตมํ ทสฺสนาย อุปสงฺกมิตุํ, กถํ หิ นาม เอวํ สมีปคโต สมาโน อทิสฺวา สมณํ โคตมํ นิวตฺติสฺสตี’ติ.\csromanspacer{} ยํ โข ปนายํ ปริสา ปริภเวยฺย, ยโสปิ ตสฺส หาเยถ.\csromanspacer{} ยสฺส โข ปน ยโส หาเยถ, โภคาปิ ตสฺส หาเยยฺยุํ, ยโสลทฺธา โข ปนมฺหากํ โภคา”ติ.}

\proseitem{307}{อถ โข โสณทณฺโฑ พฺราหฺมโณ เยน ภควา เตนุปสงฺกมิ, อุปสงฺกมิตฺวา ภควตา สทฺธิํ สมฺโมทิ, สมฺโมทนียํ กถํ สารณียํ วีติสาเรตฺวา เอกมนฺตํ นิสีทิ.\csromanspacer{} จมฺเปยฺยกาปิ โข \csromanfolio{112}พฺราหฺมณคหปติกา อปฺเปกจฺเจ ภควนฺตํ อภิวาเทตฺวา เอกมนฺตํ นิสีทิํสุ, อปฺเปกจฺเจ ภควตา สทฺธิํ สมฺโมทิํสุ, สมฺโมทนียํ กถํ สารณียํ วีติสาเรตฺวา เอกมนฺตํ นิสีทิํสุ, อปฺเปกจฺเจ เยน ภควา เตนญฺชลิํ ปณาเมตฺวา เอกมนฺตํ นิสีทิํสุ, อปฺเปกจฺเจ นามโคตฺตํ สาเวตฺวา เอกมนฺตํ นิสีทิํสุ, อปฺเปกจฺเจ ตุณฺหีภูตา เอกมนฺตํ นิสีทิํสุ.}

\proseitem{308}{ตตฺรปิ สุทํ โสณทณฺโฑ พฺราหฺมโณ เอตเทว พหุลมนุวิตกฺเกนฺโต นิสินฺโน โหติ “อหญฺเจว โข ปน สมณํ โคตมํ ปญฺหํ ปุจฺเฉยฺยํ, ตตฺร เจ มํ สมโณ โคตโม เอวํ วเทยฺย ‘น โข เอส พฺราหฺมณ ปโญฺห เอวํ ปุจฺฉิตพฺโพ, เอวํ นาเมส พฺราหฺมณ ปโญฺห ปุจฺฉิตพฺโพ’ติ.\csromanspacer{} เตน มํ อยํ ปริสา ปริภเวยฺย ‘พาโล โสณทณฺโฑ พฺราหฺมโณ อพฺยตฺโต, นาสกฺขิ สมณํ โคตมํ โยนิโส ปญฺหํ ปุจฺฉิตุนฺ’ติ.\csromanspacer{} ยํ โข ปนายํ ปริสา ปริภเวยฺย, ยโสปิ ตสฺส หาเยถ.\csromanspacer{} ยสฺส โข ปน ยโส หาเยถ, โภคาปิ ตสฺส หาเยยฺยุํ, ยโสลทฺธา โข ปนมฺหากํ โภคา.\csromanspacer{} มมญฺเจว โข ปน สมโณ โคตโม ปญฺหํ ปุจฺเฉยฺย, ตสฺส จาหํ ปญฺหสฺส เวยฺยากรเณน จิตฺตํ น อาโรเธยฺยํ.\csromanspacer{} ตตฺร เจ มํ สมโณ โคตโม เอวํ วเทยฺย ‘น โข เอส พฺราหฺมณ ปโญฺห เอวํ พฺยากาตพฺโพ, เอวํ นาเมส พฺราหฺมณ ปโญฺห พฺยากาตพฺโพ’ติ.\csromanspacer{} เตน มํ อยํ ปริสา ปริภเวยฺย ‘พาโล โสณทณฺโฑ พฺราหฺมโณ อพฺยตฺโต, นาสกฺขิ สมณสฺส โคตมสฺส ปญฺหสฺส เวยฺยากรเณน จิตฺตํ อาราเธตุนฺ’ติ.\csromanspacer{} ยํ โข ปนายํ ปริสา ปริภเวยฺย, ยโสปิ ตสฺส หาเยถ.\csromanspacer{} ยสฺส โข ปน ยโส หาเยถ, โภคาปิ ตสฺส หาเยยฺยุํ, ยโสลทฺธา โข ปนมฺหากํ โภคา.\csromanspacer{} อโห วต มํ สมโณ โคตโม สเก อาจริยเก เตวิชฺชเก ปญฺหํ ปุจฺเฉยฺย, อทฺธา วตสฺสาหํ จิตฺตํ อาราเธยฺยํ ปญฺหสฺส เวยฺยากรเณนา”ติ.}

\csromanmatikaanchor{mtk.76}
\csromantocmarkref{section}{พฺราหฺมณปญฺญตฺติ}{mtk.76}
\bookmark[dest={mtk.76},level=1]{พฺราหฺมณปญฺญตฺติ}
\csromanheader{1}{\textbf{พฺราหฺมณปญฺญตฺติ}}
\proseitem{309}{อถ โข ภควโต โสณทณฺฑสฺส พฺราหฺมณสฺส เจตสา เจโต\-ปริวิตกฺกม\-ญฺญาย เอตทโหสิ “วิหญฺญติ โข อยํ โสณทณฺโฑ พฺราหฺมโณ สเกน จิตฺเตน, ยํนูนาหํ โสณทณฺฑํ พฺราหฺมณํ สเก อาจริยเก เตวิชฺชเก ปญฺหํ ปุจฺเฉยฺยนฺ”ติ.\csromanspacer{} อถ โข ภควา}

\vspace{\tipitakaparskip}
\csromancenter{โสณทณฺฑสุตฺต โสณทณฺฑํ พฺราหฺมณํ เอตทโวจ “กติหิ ปน พฺราหฺมณ องฺเคหิ สมนฺนาคตํ พฺราหฺมณา พฺราหฺมณํ ปญฺญเปนฺติ, ‘พฺราหฺมโณสฺมี’ติ จ วทมาโน สมฺมา วเทยฺย, น จ ปน มุสาวาทํ อาปชฺเชยฺยา”ติ.}
\proseitem{310}{\csromanfolio{113}อถ โข โสณทณฺฑสฺส พฺราหฺมณสฺส เอตทโหสิ “ยํ วต โน อโหสิ อิจฺฉิตํ, ยํ อากงฺขิตํ, ยํ อธิปฺเปตํ, ยํ อภิปตฺถิตํ ‘อโห วต มํ สมโณ โคตโม สเก อาจริยเก เตวิชฺชเก ปญฺหํ ปุจฺเฉยฺย, อทฺธา วตสฺสาหํ จิตฺตํ อาราเธยฺยํ ปญฺหสฺส เวยฺยากรเณนา’ติ.\csromanspacer{} ตตฺร มํ สมโณ โคตโม สเก อาจริยเก เตวิชฺชเก ปญฺหํ ปุจฺฉติ.\csromanspacer{} อทฺธา วตสฺสาหํ จิตฺตํ อาราเธสฺสามิ ปญฺหสฺส เวยฺยากรเณนา”ติ.}

\proseitem{311}{อถ โข โสณทณฺโฑ พฺราหฺมโณ อพฺภุนฺนาเมตฺวา กายํ อนุวิโลเกตฺวา ปริสํ ภควนฺตํ เอตทโวจ “ปญฺจหิ โภ โคตม องฺเคหิ สมนฺนาคตํ พฺราหฺมณา พฺราหฺมณํ ปญฺญเปนฺติ, ‘พฺราหฺมโณสฺมี’ติ จ วทมาโน สมฺมา วเทยฺย, น จ ปน มุสาวาทํ อาปชฺเชยฺย.\csromanspacer{} กตเมหิ ปญฺจหิ.\csromanspacer{} อิธ โภ โคตม พฺราหฺมโณ อุภโต สุชาโต โหติ มาติโต จ ปิติโต จ สํสุทฺธคหณิโก ยาว สตฺตมา ปิตามหยุคา อกฺขิตฺโต อนุปกฺกุฏฺโฐ ชาติวาเทน, อชฺฌายโก โหติ มนฺตธโร ติณฺณํ เวทานํ ปารคู สนิฆณฺฑุเกฏุภานํ สากฺขรปฺปเภทานํ อิติหาสปญฺจมานํ ปทโก เวยฺยากรโณ โลกายตมหาปุริสลกฺขเณสุ อนวโย, อภิรูโป โหติ ทสฺสนีโย ปาสาทิโก ปรมาย วณฺณโปกฺขรตาย สมนฺนาคโต พฺรหฺมวณฺณี พฺรหฺมวจฺฉสี อขุทฺทาวกาโส ทสฺสนาย, สีลวา โหติ วุทฺธสีลี วุทฺธสีเลน สมนฺนาคโต, ปณฺฑิโต จ โหติ เมธาวี ปฐโม วา ทุติโย วา สุชํ ปคฺคณฺหนฺตานํ.\csromanspacer{} อิเมหิ โข โภ โคตม ปญฺจหิ องฺเคหิ สมนฺนาคตํ พฺราหฺมณา พฺราหฺมณํ ปญฺญเปนฺติ, ‘พฺราหฺมโณสฺมี’ติ จ วทมาโน สมฺมา วเทยฺย, น จ ปน มุสาวาทํ อาปชฺเชยฺยา”ติ.}

\prose{อิเมสํ ปน พฺราหฺมณ ปญฺจนฺนํ องฺคานํ สกฺกา เอกํ องฺคํ ฐปยิตฺวา จตูหงฺเคหิ สมนฺนาคตํ พฺราหฺมณา พฺราหฺมณํ ปญฺญเปตุํ, “พฺราหฺมโณสฺมี”ติ จ วทมาโน สมฺมา วเทยฺย, น จ ปน มุสาวาทํ อาปชฺเชยฺยาติ. \csromanfolio{114}สกฺกา โภ โคตม, อิเมสํ หิ โภ โคตม ปญฺจนฺนํ องฺคานํ วณฺณํ ฐปยาม, กิํ หิ วณฺโณ กริสฺสติ.\csromanspacer{} ยโต โข โภ โคตม พฺราหฺมโณ อุภโต สุชาโต โหติ มาติโต จ ปิติโต จ สํสุทฺธคหณิโก ยาว สตฺตมา ปิตามหยุคา อกฺขิตฺโต อนุปกฺกุฏฺโฐ ชาติวาเทน, อชฺฌายโก จ โหติ มนฺตธโร จ ติณฺณํ เวทานํ ปารคู สนิฆณฺฑุเกฏุภานํ สากฺขรปฺปเภทานํ อิติหาสปญฺจมานํ ปทโก เวยฺยากรโณ โลกายตมหาปุริสลกฺขเณสุ อนวโย, สีลวา จ โหติ วุทฺธสีลี วุทฺธสีเลน สมนฺนาคโต, ปณฺฑิโต จ โหติ เมธาวี ปฐโม วา ทุติโย วา สุชํ ปคฺคณฺหนฺตานํ.\csromanspacer{} อิเมหิ โข โภ โคตม จตูหงฺเคหิ สมนฺนาคตํ พฺราหฺมณา พฺราหฺมณํ ปญฺญเปนฺติ, “พฺราหฺมโณสฺมี”ติ จ วทมาโน สมฺมา วเทยฺย, น จ ปน มุสาวาทํ อาปชฺเชยฺยาติ.}

\proseitem{312}{อิเมสํ ปน พฺราหฺมณ จตุนฺนํ องฺคานํ สกฺกา เอกํ องฺคํ ฐปยิตฺวา ตีหงฺเคหิ สมนฺนาคตํ พฺราหฺมณา พฺราหฺมณํ ปญฺญเปตุํ, “พฺราหฺมโณสฺมี”ติ จ วทมาโน สมฺมา วเทยฺย, น จ ปน มุสาวาทํ อาปชฺเชยฺยาติ.\csromanspacer{} สกฺกา โภ โคตม, อิเมสํ หิ โภ โคตม จตุนฺนํ องฺคานํ มนฺเต ฐปยาม, กิํ หิ มนฺตา กริสฺสนฺติ.\csromanspacer{} ยโต โข โภ โคตม พฺราหฺมโณ อุภโต สุชาโต โหติ มาติโต จ ปิติโต จ สํสุทฺธคหณิโก ยาว สตฺตมา ปิตามหยุคา อกฺขิตฺโต อนุปกฺกุฏฺโฐ ชาติวาเทน, สีลวา จ โหติ วุทฺธสีลี วุทฺธสีเลน สมนฺนาคโต, ปณฺฑิโต จ โหติ เมธาวี ปฐโม วา ทุติโย วา สุชํ ปคฺคณฺหนฺตานํ.\csromanspacer{} อิเมหิ โข โภ โคตม ตีหงฺเคหิ สมนฺนาคตํ พฺราหฺมณา พฺราหฺมณํ ปญฺญเปนฺติ, “พฺราหฺมโณสฺมี”ติ จ วทมาโน สมฺมา วเทยฺย, น จ ปน มุสาวาทํ อาปชฺเชยฺยาติ.}

\prose{อิเมสํ ปน พฺราหฺมณ ติณฺณํ องฺคานํ สกฺกา เอกํ องฺคํ ฐปยิตฺวา ทฺวีหงฺเคหิ สมนฺนาคตํ พฺราหฺมณา พฺราหฺมณํ ปญฺญเปตุํ, “พฺราหฺมโณสฺมี”ติ จ วทมาโน สมฺมา วเทยฺย, น จ ปน มุสาวาทํ อาปชฺเชยฺยาติ.\csromanspacer{} สกฺกา โภ โคตม, อิเมสํ หิ โภ โคตม ติณฺณํ องฺคานํ ชาติํ ฐปยาม, กิํ หิ ชาติ กริสฺสติ.\csromanspacer{} ยโต โข โภ โคตม พฺราหฺมโณ สีลวา โหติ วุทฺธสีลี วุทฺธสีเลน สมนฺนาคโต, ปณฺฑิโต จ}

\vspace{\tipitakaparskip}
\csromancenter{โสณทณฺฑสุตฺต โหติ เมธาวี ปฐโม วา ทุติโย วา สุชํ ปคฺคณฺหนฺตานํ.\csromanspacer{} อิเมหิ โข โภ โคตม ทฺวีหงฺเคหิ สมนฺนาคตํ พฺราหฺมณา พฺราหฺมณํ ปญฺญเปนฺติ, “พฺราหฺมโณสฺมี”ติ จ วทมาโน สมฺมา วเทยฺย, น จ ปน มุสาวาทํ อาปชฺเชยฺยาติ.}
\proseitem{313}{\csromanfolio{115}เอวํ วุตฺเต เต พฺราหฺมณา โสณทณฺฑํ พฺราหฺมณํ เอตทโวจุํ “มา ภวํ โสณทณฺโฑ เอวํ อวจ, มา ภวํ โสณทณฺโฑ เอวํ อวจ, อปวทเตว ภวํ โสณทณฺโฑ วณฺณํ, อปวทติ มนฺเต, อปวทติ ชาติํ, เอกํเสน ภวํ โสณทณฺโฑ สมณสฺเสว โคตมสฺส วาทํ อนุปกฺขนฺทตี”ติ.}

\proseitem{314}{อถ โข ภควา เต พฺราหฺมเณ เอตทโวจ “สเจ โข ตุมฺหากํ พฺราหฺมณานํ เอวํ โหติ ‘อปฺปสฺสุโต จ โสณทณฺโฑ พฺราหฺมโณ, อกลฺยาณวากฺกรโณ จ โสณทณฺโฑ พฺราหฺมโณ, ทุปฺปญฺโญ จ โสณทณฺโฑ พฺราหฺมโณ, น จ ปโหติ โสณทณฺโฑ พฺราหฺมโณ สมเณน โคตเมน สทฺธิํ อสฺมิํ วจเน ปฏิมนฺเตตุนฺ’ติ.\csromanspacer{} ติฏฺฐตุ โสณทณฺโฑ พฺราหฺมโณ, ตุเมฺห มยา สทฺธิํ มนฺตโวฺห อสฺมิํ วจเน.\csromanspacer{} สเจ ปน ตุมฺหากํ พฺราหฺมณานํ เอวํ โหติ ‘พหุสฺสุโต จ โสณทณฺโฑ พฺราหฺมโณ, กลฺยาณวากฺกรโณ จ โสณทณฺโฑ พฺราหฺมโณ, ปณฺฑิโต จ โสณทณฺโฑ พฺราหฺมโณ, ปโหติ จ โสณทณฺโฑ พฺราหฺมโณ สมเณน โคตเมน สทฺธิํ อสฺมิํ วจเน ปฏิมนฺเตตุนฺ’ติ.\csromanspacer{} ติฏฺฐถ ตุเมฺห, โสณทณฺโฑ พฺราหฺมโณ มยา สทฺธิํ ปฏิมนฺเตตู”ติ.}

\proseitem{315}{เอวํ วุตฺเต โสณทณฺโฑ พฺราหฺมโณ ภควนฺตํ เอตทโวจ “ติฏฺฐตุ ภวํ โคตโม, ตุณฺหี ภวํ โคตโม โหตุ, อหเมว เตสํ สหธมฺเมน ปฏิวจนํ กริสฺสามี”ติ.\csromanspacer{} อถ โข โสณทณฺโฑ พฺราหฺมโณ เต พฺราหฺมเณ เอตทโวจ “มา ภวนฺโต เอวํ อวจุตฺถ, มา ภวนฺโต เอวํ อวจุตฺถ ‘อปวทเตว ภวํ โสณทณฺโฑ วณฺณํ, อปวทติ มนฺเต, อปวทติ ชาติํ, เอกํเสน ภวํ โสณทณฺโฑ สมณสฺเสว โคตมสฺส วาทํ อนุปกฺขนฺทตี’ติ.\csromanspacer{} นาหํ โภ อปวทามิ วณฺณํ วา มนฺเต วา ชาติํ วา”ติ.}

\proseitem{316}{\csromanfolio{116}เตน โข ปน สมเยน โสณทณฺฑสฺส พฺราหฺมณสฺส ภาคิเนยฺโย องฺคโก นาม มาณวโก ตสฺสํ ปริสายํ นิสินฺโน โหติ.\csromanspacer{} อถ โข โสณทณฺโฑ พฺราหฺมโณ เต พฺราหฺมเณ เอตทโวจ “ปสฺสนฺติ โน โภนฺโต อิมํ องฺคกํ มาณวกํ อมฺหากํ ภาคิเนยฺยนฺ”ติ.\csromanspacer{} เอวํ โภ.\csromanspacer{} องฺคโก โข โภ มาณวโก อภิรูโป ทสฺสนีโย ปาสาทิโก ปรมาย วณฺณโปกฺขรตาย สมนฺนาคโต พฺรหฺมวณฺณี พฺรหฺมวจฺฉสี อขุทฺทาวกาโส ทสฺสนาย, นาสฺส อิมิสฺสํ ปริสายํ สมสโม อตฺถิ วณฺเณน ฐเปตฺวา สมณํ โคตมํ.\csromanspacer{} องฺคโก โข มาณวโก อชฺฌายโก มนฺตธโร ติณฺณํ เวทานํ ปารคู สนิฆณฺฑุเกฏุภานํ สากฺขรปฺปเภทานํ อิติหาสปญฺจมานํ ปทโก เวยฺยากรโณ โลกายตมหาปุริสลกฺขเณสุ อนวโย, อหมสฺส มนฺเต วาเจตา.\csromanspacer{} องฺคโก โข มาณวโก อุภโต สุชาโต มาติโต จ ปิติโต จ สํสุทฺธคหณิโก ยาว สตฺตมา ปิตามหยุคา อกฺขิตฺโต อนุปกฺกุฏฺโฐ ชาติวาเทน, อหมสฺส มาตาปิตโร ชานามิ.\csromanspacer{} องฺคโก โข มาณวโก ปาณมฺปิ หเนยฺย, อทินฺนมฺปิ อาทิเยยฺย, ปรทารมฺปิ คจฺเฉยฺย, มุสาวาทมฺปิ ภเณยฺย, มชฺชมฺปิ ปิเวยฺย, เอตฺถ ทานิ โภ กิํ วณฺโณ กริสฺสติ, กิํ มนฺตา, กิํ ชาติ.\csromanspacer{} ยโต โข โภ พฺราหฺมโณ สีลวา จ โหติ วุทฺธสีลี วุทฺธสีเลน สมนฺนาคโต, ปณฺฑิโต จ โหติ เมธาวี ปฐโม วา ทุติโย วา สุชํ ปคฺคณฺหนฺตานํ.\csromanspacer{} อิเมหิ โข โภ ทฺวีหงฺเคหิ สมนฺนาคตํ พฺราหฺมณา พฺราหฺมณํ ปญฺญเปนฺติ, “พฺราหฺมโณสฺมี”ติ จ วทมาโน สมฺมา วเทยฺย, น จ ปน มุสาวาทํ อาปชฺเชยฺยาติ.}

\csromanmatikaanchor{mtk.77}
\csromantocmarkref{section}{สีลปญฺญากถา}{mtk.77}
\bookmark[dest={mtk.77},level=1]{สีลปญฺญากถา}
\csromanheader{1}{\textbf{สีลปญฺญากถา}}
\proseitem{317}{อิเมสํ ปน พฺราหฺมณ ทฺวินฺนํ องฺคานํ สกฺกา เอกํ องฺคํ ฐปยิตฺวา เอเกน องฺเคน สมนฺนาคตํ พฺราหฺมณา พฺราหฺมณํ ปญฺญเปตุํ, “พฺราหฺมโณสฺมี”ติ จ วทมาโน สมฺมา วเทยฺย, น จ ปน มุสาวาทํ อาปชฺเชยฺยาติ.\csromanspacer{} โน หิทํ โภ โคตม, สีลปริโธตา หิ โภ โคตม ปญฺญา, ปญฺญปริโธตํ สีลํ.\csromanspacer{} ยตฺถ สีลํ ตตฺถ ปญฺญา, ยตฺถ ปญฺญา ตตฺถ สีลํ.\csromanspacer{} สีลวโต ปญฺญา, ปญฺญวโต สีลํ.\csromanspacer{} สีลปญฺญาณญฺจ ปน โลกสฺมิํ อคฺคมกฺขายติ.\csromanspacer{} เสยฺยถาปิ โภ โคตม หตฺเถน วา หตฺถํ โธเวยฺย, ปาเทน วา ปาทํ โธเวยฺย, เอวเมว โข โภ โคตม สีลปริโธตา ปญฺญา,}

\vspace{\tipitakaparskip}
\csromancenter{โสณทณฺฑสุตฺต ปญฺญาปริโธตํ สีลํ.\csromanspacer{} ยตฺถ สีลํ ตตฺถ ปญฺญา, ยตฺถ ปญฺญา ตตฺถ สีลํ.\csromanspacer{} สีลวโต ปญฺญา, ปญฺญวโต สีลํ.\csromanspacer{} สีลปญฺญาณญฺจ ปน โลกสฺมิํ อคฺคมกฺขายตีติ.\csromanspacer{} เอวเมตํ พฺราหฺมณ, เอวเมตํ พฺราหฺมณ, สีลปริโธตา หิ พฺราหฺมณ ปญฺญา, ปญฺญาปริโธตํ สีลํ.\csromanspacer{} ยตฺถ สีลํ ตตฺถ ปญฺญา, ยตฺถ ปญฺญา ตตฺถ สีลํ.\csromanspacer{} สีลวโต ปญฺญา, ปญฺญวโต สีลํ.\csromanspacer{} สีลปญฺญาณญฺจ ปน โลกสฺมิํ อคฺคมกฺขายติ.\csromanspacer{} เสยฺยถาปิ พฺราหฺมณ หตฺเถน วา หตฺถํ โธเวยฺย, ปาเทน วา ปาทํ โธเวยฺย, เอวเมว โข พฺราหฺมณ สีลปริโธตา ปญฺญา, ปญฺญาปริโธตํ สีลํ.\csromanspacer{} ยตฺถ สีลํ ตตฺถ ปญฺญา, ยตฺถ ปญฺญา ตตฺถ สีลํ.\csromanspacer{} สีลวโต ปญฺญา, ปญฺญวโต สีลํ.\csromanspacer{} สีลปญฺญาณญฺจ ปน โลกสฺมิํ อคฺคมกฺขายติ.}
\proseitem{318}{\csromanfolio{117}กตมํ ปน ตํ พฺราหฺมณ สีลํ, กตมา สา ปญฺญาติ.\csromanspacer{} เอตฺตกปรมาว มยํ โภ โคตม เอตสฺมิํ อตฺเถ, สาธุ วต ภวนฺตํเยว โคตมํ ปฏิภาตุ เอตสฺส ภาสิตสฺส อตฺโถติ.\csromanspacer{} เตน หิ พฺราหฺมณ สุโณหิ สาธุกํ มนสิกโรหิ ภาสิสฺสามีติ. “เอวํ โภ”ติ โข โสณทณฺโฑ พฺราหฺมโณ ภควโต ปจฺจสฺโสสิ.\csromanspacer{} ภควา เอตทโวจ “อิธ พฺราหฺมณ ตถาคโต โลเก อุปฺปชฺชติ อรหํ สมฺมาสมฺพุทฺโธ ฯเปฯ. (ยถา สามญฺญผเล, เอวํ วิตฺถาเรตพฺพํ.) เอวํ โข พฺราหฺมณ ภิกฺขุ สีลสมฺปนฺโน โหติ.\csromanspacer{} อิทํ โข ตํ พฺราหฺมณ สีลํ ฯเปฯ ปฐมํ ฌานํ อุปสมฺปชฺช วิหรติ.\csromanspacer{} ทุติยํ ฌานํ.\csromanspacer{} ตติยํ ฌานํ.\csromanspacer{} จตุตฺถํ ฌานํ อุปสมฺปชฺช วิหรติ ฯเปฯ ญาณทสฺสานาย จิตฺตํ อภินีหรติ, อภินินฺนาเมติ ฯเปฯ.\csromanspacer{} อิทมฺปิสฺส โหติ ปญฺญาย ฯเปฯ นาปรํ อิตฺถตฺตายาติ ปชานาติ.\csromanspacer{} อิทมฺปิสฺส โหติ ปญฺญาย.\csromanspacer{} อยํ โข สา พฺราหฺมณ ปญฺญา”ติ.}

\csromanmatikaanchor{mtk.78}
\csromantocmarkref{section}{โสณทณฺฑ อุปาสกตฺตปฏิเวทนา}{mtk.78}
\bookmark[dest={mtk.78},level=1]{โสณทณฺฑ อุปาสกตฺตปฏิเวทนา}
\csromanheader{1}{\textbf{โสณทณฺฑ อุปาสกตฺตปฏิเวทนา}}
\proseitem{319}{เอวํ วุตฺเต โสณทณฺโฑ พฺราหฺมโณ ภควนฺตํ เอตทโวจ “อภิกฺกนฺตํ โภ โคตม, อภิกฺกนฺตํ โภ โคตม, เสยฺยถาปิ โภ โคตม นิกฺกุชฺชิตํ วา อุกฺกุชฺเชยฺย, ปฏิจฺฉนฺนํ วา วิวเรยฺย, มูฬฺหสฺส วา มคฺคํ อาจิกฺเขยฺย, อนฺธกาเร วา เตลปชฺโชตํ ธาเรยฺย ‘จกฺขุมนฺโต รูปานิ ทกฺขนฺตี’ติ.\csromanspacer{} เอวเมวํ โภตา โคตเมน อเนกปริยาเยน ธมฺโม ปกาสิโต, เอสาหํ ภวนฺตํ โคตมํ สรณํ คจฺฉามิ ธมฺมญฺจ ภิกฺขุสํฆญฺจ, อุปาสกํ มํ ภวํ โคตโม ธาเรตุ อชฺชตคฺเค ปาณุเปตํ สรณํ \csromanfolio{118}คตํ, อธิวาเสตุ จ เม ภวํ โคตโม สฺวาตนาย ภตฺตํ สทฺธิํ ภิกฺขุสํเฆนา”ติ.\csromanspacer{} อธิวาเสสิ ภควา ตุณฺหีภาเวน.}

\proseitem{320}{อถ โข โสณทณฺโฑ พฺราหฺมโณ ภควโต อธิวาสนํ วิทิตฺวา อุฏฺฐายาสนา ภควนฺตํ อภิวาเทตฺวา ปทกฺขิณํ กตฺวา ปกฺกามิ.\csromanspacer{} อถ โข โสณทณฺโฑ พฺราหฺมโณ ตสฺสา รตฺติยา อจฺจเยน สเก นิเวสเน ปณีตํ ขาทนียํ โภชนียํ ปฏิยาทาเปตฺวา ภควโต กาลํ อาโรจาเปสิ “กาโล โภ โคตม, นิฏฺฐิตํ ภตฺตนฺ”ติ.\csromanspacer{} อถ โข ภควา ปุพฺพณฺหสมยํ นิวาเสตฺวา ปตฺตจีวรมาทาย สทฺธิํ ภิกฺขุสํเฆน เยน โสณทณฺฑสฺส พฺราหฺมณสฺส นิเวสนํ เตนุปสงฺกมิ, อุปสงฺกมิตฺวา ปญฺญตฺเต อาสเน นิสีทิ.\csromanspacer{} อถ โข โสณทณฺโฑ พฺราหฺมโณ พุทฺธปฺปมุขํ ภิกฺขุสํฆํ ปณีเตน ขาทนีเยน โภชนีเยน สหตฺถา สนฺตปฺเปสิ สมฺปวาเรสิ.}

\proseitem{321}{อถ โข โสณทณฺโฑ พฺราหฺมโณ ภควนฺตํ ภุตฺตาวิํ โอนีตปตฺตปาณิํ อญฺญตรํ นีจํ อาสนํ คเหตฺวา เอกมนฺตํ นิสีทิ, เอกมนฺตํ นิสินฺโน โข โสณทณฺโฑ พฺราหฺมโณ ภควนฺตํ เอตทโวจ “อหญฺเจว โข ปน โภ โคตม ปริสคโต สมาโน อาสนา วุฏฺฐหิตฺวา ภวนฺตํ โคตมํ อภิวาเทยฺยํ, เตน มํ สา ปริสา ปริภเวยฺย.\csromanspacer{} ยํ โข ปน สา ปริสา ปริภเวยฺย, ยโสปิ ตสฺส หาเยถ.\csromanspacer{} ยสฺส โข ปน ยโส หาเยถ, โภคาปิ ตสฺส หาเยยฺยุํ, ยโสลทฺธา โข ปนมฺหากํ โภคา.\csromanspacer{} อหญฺเจว โข ปน โภ โคตม ปริสคโต สมาโน อญฺชลิํ ปคฺคเณฺหยฺยํ, อาสนา เม ตํ ภวํ โคตโม ปจฺจุฏฺฐานํ ธาเรตุ, อหญฺเจว โข ปน โภ โคตม ปริสคโต สมาโน เวฐนํ โอมุญฺเจยฺยํ, สิรสา เม ตํ ภวํ โคตโม อภิวาทนํ ธาเรตุ.\csromanspacer{} อหญฺเจว โข ปน โภ โคตม ยานคโต สมาโน ยานา ปจฺโจโรหิตฺวา ภวนฺตํ โคตมํ อภิวาเทยฺยํ, เตน มํ สา ปริสา ปริภเวยฺย.\csromanspacer{} ยํ โข ปน สา ปริสา ปริภเวยฺย, ยโสปิ ตสฺส หาเยถ.\csromanspacer{} ยสฺส โข ปน ยโส หาเยถ, โภคาปิ ตสฺส หาเยยฺยุํ, ยโสลทฺธา โข ปนมฺหากํ โภคา.\csromanspacer{} อหญฺเจว โข ปน โภ โคตม ยานคโต สมาโน ปโตทลฏฺฐิํ อพฺภุนฺนาเมยฺยํ, ยานา เม ตํ ภวํ โคตโม ปจฺโจโรหนํ ธาเรตุ,}

\vspace{\tipitakaparskip}
\csromancenter{โสณทณฺฑสุตฺต อหญฺเจว โข ปน โภ โคตม ยานคโต สมาโน ฉตฺตํ อปนาเมยฺยํ, สิรสา เม ตํ ภวํ โคตโม อภิวาทนํ ธาเรตู”ติ.}
\proseitemwithcloser{322}{\csromanfolio{119}อถ โข ภควา โสณทณฺฑํ พฺราหฺมณํ ธมฺมิยา กถาย สนฺทสฺเสตฺวา สมาทเปตฺวา สมุตฺเตเชตฺวา สมฺปหํเสตฺวา อุฏฺฐายาสนา ปกฺกามีติ.}{โสณทณฺฑสุตฺตํ นิฏฺฐิตํ จตุตฺถํ.}
\csromanmatikaanchor{mtk.79}
\csromantocmarknopage{chapter}{5. กูฏทนฺตสุตฺต}{mtk.79}
\bookmark[dest={mtk.79},level=0]{5. กูฏทนฺตสุตฺต}
\setcsromanheadcha{5. กูฏทนฺตสุตฺต}
\setcsromanheadhclear{}
\chapterheadpageread{5. \textbf{กูฏทนฺตสุตฺต}}
\csromanmatikaanchor{mtk.80}
\csromantocmarkref{section}{ขาณุมตกพฺราหฺมณคหปติกา}{mtk.80}
\bookmark[dest={mtk.80},level=1]{ขาณุมตกพฺราหฺมณคหปติกา}
\csromanheader{1}{ขาณุมตก\-พฺราหฺมณคหปติกา}
\proseitem{323}{\csromanfolio{120}เอวํ เม สุตํ\csromandash{}เอกํ สมยํ ภควา มคเธสุ จาริกํ จรมาโน มหตา ภิกฺขุสํเฆน สทฺธิํ ปญฺจมตฺเตหิ ภิกฺขุสเตหิ เยน ขาณุมตํ นาม มคธานํ พฺราหฺมณคาโม ตทวสริ.\csromanspacer{} ตตฺร สุทํ ภควา ขาณุมเต วิหรติ อมฺพลฏฺฐิกายํ.\csromanspacer{} เตน โข ปน สมเยน กูฏทนฺโต พฺราหฺมโณ ขาณุมตํ อชฺฌาวสติ สตฺตุสฺสทํ สติณกฏฺโฐทกํ สธญฺญํ ราชโภคฺคํ รญฺญา มาคเธน เสนิเยน พิมฺพิสาเรน ทินฺนํ ราชทายํ พฺรหฺมเทยฺยํ.\csromanspacer{} เตน โข ปน สมเยน กูฏทนฺตสฺส พฺราหฺมณสฺส มหายญฺโญ อุปกฺขโฏ โหติ, สตฺต จ อุสภสตานิ สตฺต จ วจฺฉตรสตานิ สตฺต จ วจฺฉตรีสตานิ สตฺต จ อชสตานิ สตฺต จ อุรพฺภสตานิ ถูณูปนีตานิ โหนฺติ ยญฺญตฺถาย.}

\proseitem{324}{อสฺโสสุํ โข ขาณุมตกา พฺราหฺมณคหปติกา “สมโณ ขลุ โภ โคตโม สกฺยปุตฺโต สกฺยกุลา ปพฺพชิโต มคเธสุ จาริกํ จรมาโน มหตา ภิกฺขุสํเฆน สทฺธิํ ปญฺจมตฺเตหิ ภิกฺขุสเตหิ ขาณุมตํ อนุปฺปตฺโต ขาณุมเต วิหรติ อมฺพลฏฺฐิกายํ.\csromanspacer{} ตํ โข ปน ภวนฺตํ โคตมํ เอวํ กลฺยาโณ กิตฺติสทฺโท อพฺภุคฺคโต ‘อิติปิ โส ภควา อรหํ สมฺมาสมฺพุทฺโธ วิชฺชาจรณสมฺปนฺโน สุคโต โลกวิทู อนุตฺตโร ปุริสทมฺมสารถิ สตฺถา เทวมนุสฺสานํ พุทฺโธ ภควา’ติ, โส อิมํ โลกํ สเทวกํ สมารกํ สพฺรหฺมกํ สสฺสมณพฺราหฺมณิํ ปชํ สเทวมนุสฺสํ สยํ อภิญฺญา สจฺฉิกตฺวา ปเวเทติ, โส ธมฺมํ เทเสติ อาทิกลฺยาณํ มชฺเฌกลฺยาณํ ปริโยสานกลฺยาณํ สาตฺถํ สพฺยญฺชนํ เกวลปริปุณฺณํ ปริสุทฺธํ พฺรหฺมจริยํ ปกาเสติ, สาธุ โข ปน ตถารูปานํ อรหตํ ทสฺสนํ โหตี”ติ.}

\proseitem{325}{อถ โข ขาณุมตกา พฺราหฺมณคหปติกา ขาณุมตา นิกฺขมิตฺวา สงฺฆสงฺฆี คณีภูตา เยน อมฺพลฏฺฐิกา เตนุปสงฺกมนฺติ.}

\csromanheader{2}{5. กูฏทนฺตสุตฺต}
\proseitem{326}{\csromanfolio{121}เตน โข ปน สมเยน กูฏทนฺโต พฺราหฺมโณ อุปริปาสาเท ทิวาเสยฺยํ อุปคโต โหติ.\csromanspacer{} อทฺทสา โข กูฏทนฺโต พฺราหฺมโณ ขาณุมตเก พฺราหฺมณคหปติเก ขาณุมตา นิกฺขมิตฺวา สงฺฆสงฺฆี คณีภูเต เยน อมฺพลฏฺฐิกา เตนุปสงฺกมนฺเต, ทิสฺวา ขตฺตํ อามนฺเตสิ “กิํ นุ โข โภ ขตฺเต ขาณุมตตา พฺราหฺมณคหปติกา ขาณุมตา นิกฺขมิตฺวา สงฺฆสงฺฆี คณีภูตา เยน อมฺพลฏฺฐิกา เตนุปสงฺกมนฺตี”ติ.}

\proseitem{327}{อตฺถิ โข โภ สมโณ โคตโม สกฺยปุตฺโต สกฺยกุลา ปพฺพชิโต มคเธสุ จาริกํ จรมาโน มหตา ภิกฺขุสํเฆน สทฺธิํ ปญฺจมตฺเตหิ ภิกฺขุสเตหิ ขาณุมตํ อนุปฺปตฺโต ขาณุมเต วิหรติ อมฺพลฏฺฐิกายํ.\csromanspacer{} ตํ โข ปน ภวนฺตํ โคตมํ เอวํ กลฺยาโณ กิตฺติสทฺโท อพฺภุคฺคโต “อิติปิ โส ภควา อรหํ สมฺมาสมฺพุทฺโธ วิชฺชาจรณสมฺปนฺโน สุคโต โลกวิทู อนุตฺตโร ปุริสทมฺมสารถิ สตฺถา เทวมนุสฺสานํ พุทฺโธ ภควา”ติ.\csromanspacer{} ตเมเต ภวนฺตํ โคตมํ ทสฺสนาย อุปสงฺกมนฺตีติ.}

\proseitem{328}{อถ โข กูฏทนฺตสฺส พฺราหฺมณสฺส เอตทโหสิ “สุตํ โข ปน เมตํ ‘สมโณ โคตโม ติวิธํ ยญฺญสมฺปทํ โสฬสปริกฺขารํ ชานาตี’ติ.\csromanspacer{} น โข ปนาหํ ชานามิ ติวิธํ ยญฺญสมฺปทํ โสฬสปริกฺขารํ, อิจฺฉามิ จาหํ มหายญฺญํ ยชิตุํ, ยํนูนาหํ สมณํ โคตมํ อุปสงฺกมิตฺวา ติวิธํ ยญฺญสมฺปทํ โสฬสปริกฺขารํ ปุจฺเฉยฺยนฺ”ติ.}

\proseitem{329}{อถ โข กูฏทนฺโต พฺราหฺมโณ ขตฺตํ อามนฺเตสิ “เตน หิ โภ ขตฺเต เยน ขาณุมตกา พฺราหฺมณคหปติกา เตนุปสงฺกม, อุปสงฺกมิตฺวา ขาณุมตเก พฺราหฺมณคหปติเก เอวํ วเทหิ กูฏทนฺโต โภ พฺราหฺมโณ เอวมาห อาคเมนฺตุ กิร ภวนฺโต, กูฏทนฺโตปิ พฺราหฺมโณ สมณํ โคตมํ ทสฺสนาย อุปสงฺกมิสฺสตี”ติ. “เอวํ โภ”ติ โข โส ขตฺตา กูฏทนฺตสฺส พฺราหฺมณสฺส ปฏิสฺสุตฺวา เยน ขาณุมตกา พฺราหฺมณคหปติกา เตนุปสงฺกมิ, อุปสงฺกมิตฺวา ขาณุมตเก พฺราหฺมณคหปติเก เอตทโวจ “กูฏทนฺโต โภ พฺราหฺมโณ เอวมาห อาคเมนฺตุ กิร โภนฺโต, กูฏทนฺโตปิ พฺราหฺมโณ สมณํ โคตมํ ทสฺสนาย อุปสงฺกมิสฺสตี”ติ.}

\csromanmatikaanchor{mtk.81}
\csromantocmarkref{section}{กูฏทนฺตคุณกถา}{mtk.81}
\bookmark[dest={mtk.81},level=1]{กูฏทนฺตคุณกถา}
\csromanheader{1}{\textbf{กูฏทนฺตคุณกถา}}
\proseitem{330}{\csromanfolio{122}เตน โข ปน สมเยน อเนกานิ พฺราหฺมณสตานิ ขาณุมเต ปฏิวสนฺติ “กูฏทนฺตสฺส พฺราหฺมณสฺส มหายญฺญํ อนุภวิสฺสามา”ติ.\csromanspacer{} อสฺโสสุํ โข เต พฺราหฺมณา “กูฏทนฺโต กิร พฺราหฺมโณ สมณํ โคตมํ ทสฺสนาย อุปสงฺกมิสฺสตี”ติ.\csromanspacer{} อถ โข เต พฺราหฺมณา เยน กูฏทนฺโต พฺราหฺมโณ เตนุปสงฺกมิํสุ.}

\proseitem{331}{อุปสงฺกมิตฺวา กูฏทนฺตํ พฺราหฺมณํ เอตทโวจุํ “สจฺจํ กิร ภวํ กูฏทนฺโต สมณํ โคตมํ ทสฺสนาย อุปสงฺกมิสฺสตี”ติ.\csromanspacer{} เอวํ โข เม โภ โหติ “อหมฺปิ สมณํ โคตมํ ทสฺสนาย อุปสงฺกมิสฺสามี”ติ.\csromanspacer{} มา ภวํ กูฏทนฺโต สมณํ โคตมํ ทสฺสนาย อุปสงฺกมิ, น อรหติ ภวํ กูฏทนฺโต สมณํ โคตมํ ทสฺสนาย อุปสงฺกมิตุํ.\csromanspacer{} สเจ ภวํ กูฏทนฺโต สมณํ โคตมํ ทสฺสนาย อุปสงฺกมิสฺสติ, โภโต กูฏทนฺตสฺส ยโส หายิสฺสติ, สมณสฺส โคตมสฺส ยโส อภิวฑฺฒิสฺสติ.\csromanspacer{} ยมฺปิ โภโต กูฏทนฺตสฺส ยโส หายิสฺสติ, สมณสฺส โคตมสฺส ยโส อภิวฑฺฒิสฺสติ, อิมินาปงฺเคน น อรหติ ภวํ กูฏทนฺโต สมณํ โคตมํ ทสฺสนาย อุปสงฺกมิตุํ, สมโณเตฺวว โคตโม อรหติ ภวนฺตํ กูฏทนฺตํ ทสฺสนาย อุปสงฺกมิตุํ.}

\prose{ภวํ หิ กูฏทนฺโต อุภโต สุชาโต มาติโต จ ปิติโต จ สํสุทฺธคหณิโก ยาว สตฺตมา ปิตามหยุคา อกฺขิตฺโต อนุปกฺกุฏฺโฐ ชาติวาเทน.\csromanspacer{} ยมฺปิ ภวํ กูฏทนฺโต อุภโต สุชาโต มาติโต จ ปิติโต จ สํสุทฺธคหณิโก ยาว สตฺตมา ปิตามหยุคา อกฺขิตฺโต อนุปกฺกุฏฺโฐ ชาติวาเทน, อิมินาปงฺเคน น อรหติ ภวํ กูฏทนฺโต สมณํ โคตมํ ทสฺสนาย อุปสงฺกมิตุํ, สมโณเตฺวว โคตโม อรหติ ภวนฺตํ กูฏทนฺตํ ทสฺสนาย อุปสงฺกมิตุํ.}

\prose{ภวํ หิ กูฏทนฺโต อฑฺโฒ มหทฺธโน มหาโภโค ปหูตวิตฺตูปกรโณ ปหูตชาตรูปรชโต ฯเปฯ.}

\prose{ภวํ หิ กูฏทนฺโต อชฺฌายโก มนฺตธโร ติณฺณํ เวทานํ ปารคู สนิฆณฺฑุเกฏุภานํ สากฺขรปฺปเภทานํ อิติหาสปญฺจมานํ ปทโก เวยฺยากรโณ โลกายตมหาปุริสลกฺขเณสุ อนวโย ฯเปฯ.}

\csromanheader{2}{5. กูฏทนฺตสุตฺต}
\prose{\csromanfolio{123}ภวํ หิ กูฏทนฺโต อภิรูโป ทสฺสนีโย ปาสาทิโก ปรมาย วณฺณโปกฺขรตาย สมนฺนาคโต พฺรหฺมวณฺณี พฺรหฺมวจฺฉสี อขุทฺทาวกาโส ทสฺสนาย ฯเปฯ.}

\prose{ภวํ หิ กูฏทนฺโต สีลวา วุทฺธสีลี วุทฺธสีเลน สมนฺนาคโต ฯเปฯ.}

\prose{ภวํ หิ กูฏทนฺโต กลฺยาณวาโจ กลฺยาณวากฺกรโณ โปริยา วาจาย สมนฺนาคโต วิสฺสฏฺฐาย อเนลคลาย อตฺถสฺส วิญฺญาปนิยา ฯเปฯ.}

\prose{ภวํ หิ กูฏทนฺโต พหูนํ อาจริยปาจริโย ตีณิ มาณวกสตานิ มนฺเต วาเจติ, พหู โข ปน นานาทิสา นานาชนปทา มาณวกา อาคจฺฉนฺติ โภโต กูฏทนฺตสฺส สนฺติเก มนฺตตฺถิกา มนฺเต อธิยิตุกามา ฯเปฯ.}

\prose{ภวํ หิ กูฏทนฺโต ชิณฺโณ วุทฺโธ มหลฺลโก อทฺธคโต วโย- อนุปฺปตฺโต, สมโณ โคตโม ตรุโณ เจว ตรุณปพฺพชิโต จ ฯเปฯ.}

\prose{ภวํ หิ กูฏทนฺโต รญฺโญ มาคธสฺส เสนิยสฺส พิมฺพิสารสฺส สกฺกโต ครุกโต มานิโต ปูชิโต อปจิโต ฯเปฯ.}

\prose{ภวํ หิ กูฏทนฺโต พฺราหฺมณสฺส โปกฺขรสาติสฺส สกฺกโต ครุกโต มานิโต ปูชิโต อปจิโต ฯเปฯ.}

\prose{ภวํ หิ กูฏทนฺโต ขาณุมตํ อชฺฌาวสติ สตฺตุสฺสทํ สติณกฏฺโฐทกํ สธญฺญํ ราชโภคฺคํ รญฺญา มาคเธน เสนิเยน พิมฺพิสาเรน ทินฺนํ ราชทายํ พฺรหฺมเทยฺยํ.\csromanspacer{} ยมฺปิ ภวํ กูฏทนฺโต ขาณุมตํ อชฺฌาวสติ สตฺตุสฺสทํ สติณกฏฺโฐทกํ สธญฺญํ ราชโภคฺคํ รญฺญา มาคเธน เสนิเยน พิมฺพิสาเรน ทินฺนํ ราชทายํ พฺรหฺมเทยฺยํ, อิมินาปงฺเคน น อรหติ ภวํ กูฏทนฺโต สมณํ โคตมํ ทสฺสนาย อุปสงฺกมิตุํ, สมโณเตฺวว โคตโม อรหติ ภวนฺตํ กูฏทนฺตํ ทสฺสนาย อุปสงฺกมิตุนฺติ.}

\csromanmatikaanchor{mtk.82}
\csromantocmarkref{section}{พุทฺธคุณกถา}{mtk.82}
\bookmark[dest={mtk.82},level=1]{พุทฺธคุณกถา}
\csromanheader{1}{\textbf{พุทฺธคุณกถา}}
\csromanheader{2}{332. เอวํ วุตฺเต กูฏทนฺโต พฺราหฺมโณ เต พฺราหฺมเณ เอตทโวจ\csromandash{}}
\prose{\csromanfolio{124}“เตน หิ โภ มมปิ สุณาถ, ยถา มยเมว อรหาม ตํ ภวนฺตํ โคตมํ ทสฺสนาย อุปสงฺกมิตุํ, น เตฺวว อรหติ โส ภวํ โคตโม อมฺหากํ ทสฺสนาย อุปสงฺกมิตุํ.\csromanspacer{} สมโณ ขลุ โภ โคตโม อุภโต สุชาโต มาติโต จ ปิติโต จ สํสุทฺธคหณิโก ยาว สตฺตมา ปิตามหยุคา อกฺขิตฺโต อนุปกฺกุฏฺโฐ ชาติวาเทน.\csromanspacer{} ยมฺปิ โภ สมโณ โคตโม อุภโต สุชาโต มาติโต จ ปิติโต จ สํสุทฺธคหณิโก ยาว สตฺตมา ปิตามหยุคา อกฺขิตฺโต อนุปกฺกุฏฺโฐ ชาติวาเทน, อิมินาปงฺเคน น อรหติ โส ภวํ โคตโม อมฺหากํ ทสฺสนาย อุปสงฺกมิตุํ, อถ โข มยเมว อรหาม ตํ ภวนฺตํ โคตมํ ทสฺสนาย อุปสงฺกมิตุํ.}

\prose{สมโณ ขลุ โภ โคตโม มหนฺตํ ญาติสํฆํ โอหาย ปพฺพชิโต ฯเปฯ.}

\prose{สมโณ ขลุ โภ โคตโม ปหูตํ หิรญฺญสุวณฺณํ โอหาย ปพฺพชิโต ภูมิคตญฺจ เวหาสฏฺฐญฺจ ฯเปฯ.}

\prose{สมโณ ขลุ โภ โคตโม ทหโรว สมาโน ยุวา สุสุกาฬเกโส ภเทฺรน โยพฺพเนน สมนฺนาคโต ปฐเมน วยสา อคารสฺมา อนคาริยํ ปพฺพชิโต ฯเปฯ.}

\prose{สมโณ ขลุ โภ โคตโม อกามกานํ มาตาปิตูนํ อสฺสุมุขานํ รุทนฺตานํ เกสมสฺสุํ โอหาเรตฺวา กาสายานิ วตฺถานิ อจฺฉาเทตฺวา อคารสฺมา อนคาริยํ ปพฺพชิโต ฯเปฯ.}

\prose{สมโณ ขลุ โภ โคตโม อภิรูโป ทสฺสนีโย ปาสาทิโก ปรมาย วณฺณโปกฺขรตาย สมนฺนาคโต พฺรหฺมวณฺณี พฺรหฺมวจฺฉสี อขุทฺทาวกาโส ทสฺสนาย ฯเปฯ.}

\prose{สมโณ ขลุ โภ โคตโม สีลวา อริยสีลี กุสลสีลี กุสลสีเลน สมนฺนาคโต ฯเปฯ.}

\prose{สมโณ ขลุ โภ โคตโม กลฺยาณวาโจ กลฺยาณวากฺกรโณ โปริยา วาจาย สมนฺนาคโต วิสฺสฏฺฐาย อเนลคลาย อตฺถสฺส วิญฺญาปนิยา ฯเปฯ.}

\csromanheader{2}{5. กูฏทนฺตสุตฺต}
\prose{\csromanfolio{125}สมโณ ขลุ โภ โคตโม พหูนํ อาจริยปาจริโย ฯเปฯ.}

\prose{สมโณ ขลุ โภ โคตโม ขีณกามราโค วิคตจาปลฺโล ฯเปฯ.}

\prose{สมโณ ขลุ โภ โคตโม กมฺมวาที กิริยวาที อปาปปุเรกฺขาโร พฺรหฺมญฺญาย ปชาย ฯเปฯ.}

\prose{สมโณ ขลุ โภ โคตโม อุจฺจา กุลา ปพฺพชิโต อสมฺภินฺนขตฺติยกุลา ฯเปฯ.}

\prose{สมโณ ขลุ โภ โคตโม อฑฺฒา กุลา ปพฺพชิโต มหทฺธนา มหาโภคา ฯเปฯ.}

\prose{สมณํ ขลุ โภ โคตมํ ติโรรฏฺฐา ติโรชนปทา ปญฺหํ ปุจฺฉิตุํ อาคจฺฉนฺติ ฯเปฯ.}

\prose{สมณํ ขลุ โภ โคตมํ อเนกานิ เทวตาสหสฺสานิ ปาเณหิ สรณํ คตานิ ฯเปฯ.}

\prose{สมณํ ขลุ โภ โคตมํ เอวํ กลฺยาโณ กิตฺติสทฺโท อพฺภุคฺคโต ‘อิติปิ โส ภควา อรหํ สมฺมาสมฺพุทฺโธ วิชฺชาจรณสมฺปนฺโน สุคโต โลกวิทู อนุตฺตโร ปุริสทมฺมสารถิ สตฺถา เทวมนุสฺสานํ พุทฺโธ ภควา’ติ ฯเปฯ.}

\prose{สมโณ ขลุ โภ โคตโม ทฺวตฺติํสมหาปุริสลกฺขเณหิ สมนฺนาคโต ฯเปฯ.}

\prose{สมโณ ขลุ โภ โคตโม เอหิสฺวาคตวาที สขิโล สมฺโมทโก อพฺภากุฏิโก อุตฺตานมุโข ปุพฺพภาสี ฯเปฯ.}

\prose{สมโณ ขลุ โภ โคตโม จตุนฺนํ ปริสานํ สกฺกโต ครุกโต มานิโต ปูชิโต อปจิโต ฯเปฯ.}

\prose{สมเณ ขลุ โภ โคตเม พหู เทวา จ มนุสฺสา จ อภิปฺปสนฺนา ฯเปฯ.}

\prose{สมโณ ขลุ โภ โคตโม ยสฺมิํ คาเม วา นิคเม วา ปฏิวสติ, น ตสฺมิํ คาเม วา นิคเม วา อมนุสฺสา มนุสฺเส วิเหเฐนฺติ ฯเปฯ.}

\prose{สมโณ ขลุ โภ โคตโม สงฺฆี คณี คณาจริโย ปุถุติตฺถกรานํ อคฺคมกฺขายติ, ยถา โข ปน โภ เอเตสํ สมณพฺราหฺมณานํ}

\prose{\csromanfolio{126}ยถา วา ตถา วา ยโส สมุทาคจฺฉติ, น เหวํ สมณสฺส โคตมสฺส ยโส สมุทาคโต, อถ โข อนุตฺตราย วิชฺชาจรณสมฺปทาย สมณสฺส โคตมสฺส ยโส สมุทาคโต ฯเปฯ.}

\prose{สมณํ ขลุ โภ โคตมํ ราชา มาคโธ เสนิโย พิมฺพิสาโร สปุตฺโต สภริโย สปริโส สามจฺโจ ปาเณหิ สรณํ คโต ฯเปฯ.}

\prose{สมณํ ขลุ โภ โคตมํ ราชา ปเสนทิ โกสโล สปุตฺโต สภริโย สปริโส สามจฺโจ ปาเณหิ สรณํ คโต ฯเปฯ.}

\prose{สมณํ ขลุ โภ โคตมํ พฺราหฺมโณ โปกฺขรสาติ สปุตฺโต สภริโย สปริโส สามจฺโจ ปาเณหิ สรณํ คโต ฯเปฯ.}

\prose{สมโณ ขลุ โภ โคตโม รญฺโญ มาคธสฺส เสนิยสฺส พิมฺพิสารสฺส สกฺกโต ครุกโต มานิโต ปูชิโต อปจิโต ฯเปฯ.}

\prose{สมโณ ขลุ โภ โคตโม รญฺโญ ปเสนทิสฺส โกสลสฺส สกฺกโต ครุกโต มานิโต ปูชิโต อปจิโต ฯเปฯ.}

\prose{สมโณ ขลุ โภ โคตโม พฺราหฺมณสฺส โปกฺขรสาติสฺส สกฺกโต ครุกโต มานิโต ปูชิโต อปจิโต ฯเปฯ.}

\prose{สมโณ ขลุ โภ โคตโม ขาณุมตํ อนุปฺปตฺโต ขาณุมเต วิหรติ อมฺพลฏฺฐิกายํ, เย โข ปน โภ เกจิ สมณา วา พฺราหฺมณา วา อมฺหากํ คามเขตฺตํ อาคจฺฉนฺติ, อติถี โน เต โหนฺติ, อติถี โข ปนเมฺหหิ สกฺกาตพฺพา ครุกาตพฺพา มาเนตพฺพา ปูเชตพฺพา อปเจตพฺพา.\csromanspacer{} ยมฺปิ โภ สมโณ โคตโม ขาณุมตํ อนุปฺปตฺโต ขาณุมเต วิหรติ อมฺพลฏฺฐิกายํ, อติถิมฺหากํ สมโณ โคตโม, อติถี โข ปนเมฺหหิ สกฺกาตพฺโพ ครุกาตพฺโพ มาเนตพฺโพ ปูเชตพฺโพ อปเจตพฺโพ, อิมินาปงฺเคน นารหติ โส ภวํ โคตโม อมฺหากํ ทสฺสนาย อุปสงฺกมิตุํ, อถ โข มยเมว อรหาม ตํ ภวนฺตํ โคตมํ ทสฺสนาย อุปสงฺกมิตุํ.\csromanspacer{} เอตฺตเก โข อหํ โภ ตสฺส โภโต โคตมสฺส วณฺเณ ปริยาปุณามิ, โน จ โข โส ภวํ โคตโม เอตฺตกวณฺโณ, อปริมาณวณฺโณ หิ โส ภวํ โคตโม”ติ.}

\csromanheader{2}{5. กูฏทนฺตสุตฺต}
\proseitem{333}{\csromanfolio{127}เอวํ วุตฺเต เต พฺราหฺมณ กูฏทนฺตํ พฺราหฺมณํ เอตทโวจุํ “ยถา โข ภวํ กูฏทนฺโต สมณสฺส โคตมสฺส วณฺเณ ภาสติ, อิโต เจปิ โส ภวํ โคตโม โยชนสเต วิหรติ, อลเมว สทฺเธน กุลปุตฺเตน ทสฺสนาย อุปสงฺกมิตุํ อปิ ปุโฏเสนา”ติ.\csromanspacer{} เตน หิ โภ สพฺเพว มยํ สมณํ โคตมํ ทสฺสนาย อุปสงฺกมิสฺสามาติ.}

\csromanmatikaanchor{mtk.83}
\csromantocmarkref{section}{มหาวิชิตราชยญฺญกถา}{mtk.83}
\bookmark[dest={mtk.83},level=1]{มหาวิชิตราชยญฺญกถา}
\csromanheader{1}{มหาวิชิต\-ราช\-ยญฺญกถา}
\proseitem{334}{อถ โข กูฏทนฺโต พฺราหฺมโณ มหตา พฺราหฺมณคเณน สทฺธิํ เยน อมฺพลฏฺฐิกา เยน ภควา เตนุปสงฺกมิ, อุปสงฺกมิตฺวา ภควตา สทฺธิํ สมฺโมทิ, สมฺโมทนียํ กถํ สารณียํ วีติสาเรตฺวา เอกมนฺตํ นิสีทิ.\csromanspacer{} ขาณุมตกาปิ โข พฺราหฺมณคหปติกา อปฺเปกจฺเจ ภควนฺตํ อภิวาเทตฺวา เอกมนฺตํ นิสีทิํสุ, อปฺเปกจฺเจ ภควตา สทฺธิํ สมฺโมทิํสุ, สมฺโมทนียํ กถํ สารณียํ วีติสาเรตฺวา เอกมนฺตํ นิสีทิํสุ, อปฺเปกจฺเจ เยน ภควา เตนญฺชลิํ ปณาเมตฺวา เอกมนฺตํ นิสีทิํสุ, อปฺเปกจฺเจ นามโคตฺตํ สาเวตฺวา เอกมนฺตํ นิสีทิํสุ, อปฺเปกจฺเจ ตุณฺหีภูตา เอกมนฺตํ นิสีทิํสุ.}

\proseitem{335}{เอกมนฺตํ นิสินฺโน โข กูฏทนฺโต พฺราหฺมโณ ภควนฺตํ เอตทโวจ “สุตํ เมตํ โภ โคตม ‘สมโณ โคตโม ติวิธํ ยญฺญสมฺปทํ โสฬสปริกฺขารํ ชานาตี’ติ, น โข ปนาหํ ชานามิ ติวิธํ ยญฺญสมฺปทํ โสฬสปริกฺขารํ, อิจฺฉามิ จาหํ มหายญฺญํ ยชิตุํ, สาธุ เม ภวํ โคตโม ติวิธํ ยญฺญสมฺปทํ โสฬสปริกฺขารํ เทเสตู”ติ.}

\proseitem{336}{เตน หิ พฺราหฺมณ สุณาหิ สาธุกํ มนสิกโรหิ ภาสิสฺสามีติ. “เอวํ โภ”ติ โข กูฏทนฺโต พฺราหฺมโณ ภควโต ปจฺจสฺโสสิ.\csromanspacer{} ภควา เอตทโวจ, ภูตปุพฺพํ พฺราหฺมณ ราชา มหาวิชิโต นาม อโหสิ อฑฺโฒ มหทฺธโน มหาโภโค ปหูตชาตรูปรชโต ปหูตวิตฺตูปกรโณ ปหูตธนธญฺโญ ปริปุณฺณโกสโกฏฺฐาคาโร.\csromanspacer{} อถ โข พฺราหฺมณ รญฺโญ มหาวิชิตสฺส รโหคตสฺส ปฏิสลฺลีนสฺส เอวํ เจตโส ปริวิตกฺโก อุทปาทิ “อธิคตา โข เม วิปุลา มานุสกา โภคา, มหนฺตํ ปถวิมณฺฑลํ อภิวิชิย \csromanfolio{128}อชฺฌาวสามิ, ยํนูนาหํ มหายญฺญํ ยเชยฺยํ ยํ มม อสฺส ทีฆรตฺตํ หิตาย สุขายา”ติ.}

\proseitem{337}{อถ โข พฺราหฺมณ ราชา มหาวิชิโต ปุโรหิตํ พฺราหฺมณํ อามนฺเตตฺวา เอตทโวจ “อิธ มยฺหํ พฺราหฺมณ รโหคตสฺส ปฏิสลฺลีนสฺส เอวํ เจตโส ปริวิตกฺโก อุทปาทิ, ‘อธิคตา โข เม วิปุลา มานุสกา โภคา, มหนฺตํ ปถวิมณฺฑลํ อภิวิชิย อชฺฌาวสามิ, ยํนูนาหํ มหายญฺญํ ยเชยฺยํ ยํ มม อสฺส ทีฆรตฺตํ หิตาย สุขายา’ติ, อิจฺฉามหํ พฺราหฺมณ มหายญฺญํ ยชิตุํ, อนุสาสตุ มํ ภวํ ยํ มม อสฺส ทีฆรตฺตํ หิตาย สุขายา”ติ.}

\proseitem{338}{เอวํ วุตฺเต พฺราหฺมณ ปุโรหิโต พฺราหฺมโณ ราชานํ มหาวิชิตํ เอตทโวจ “โภโต โข รญฺโญ ชนปโท สกณฺฏโก ส-อุปฺปีโฬ, คามฆาตาปิ ทิสฺสนฺติ, นิคมฆาตาปิ ทิสฺสนฺติ, นครฆาตาปิ ทิสฺสนฺติ, ปนฺถทุหนาปิ ทิสฺสนฺติ.\csromanspacer{} ภวํ โข ปน ราชา เอวํ สกณฺฏเก ชนปเท ส-อุปฺปีเฬ พลิมุทฺธเรยฺย, อกิจฺจการี อสฺส เตน ภวํ ราชา.\csromanspacer{} สิยา โข ปน โภโต รญฺโญ เอวมสฺส ‘อหเมตํ ทสฺสุขีลํ วเธน วา พนฺเธน วา ชานิยา วา ครหาย วา ปพฺพาชนาย วา สมูหนิสฺสามี’ติ, น โข ปเนตสฺส ทสฺสุขีลสฺส เอวํ สมฺมา สมุคฺฆาโต โหติ, เย เต หตาวเสสกา ภวิสฺสนฺติ, เต ปจฺฉา รญฺโญ ชนปทํ วิเหเฐสฺสนฺติ.\csromanspacer{} อปิ จ โข อิทํ สํวิธานํ อาคมฺม เอวเมตสฺส ทสฺสุขีลสฺส สมฺมา สมุคฺฆาโต โหติ, เตน หิ ภวํ ราชา เย โภโต รญฺโญ ชนปเท อุสฺสหนฺติ กสิโครกฺเข, เตสํ ภวํ ราชา พีชภตฺตํ อนุปฺปเทตุ.\csromanspacer{} เย โภโต รญฺโญ ชนปเท อุสฺสหนฺติ วาณิชฺชาย, เตสํ ภวํ ราชา ปาภตํ อนุปฺปเทตุ.\csromanspacer{} เย โภโต รญฺโญ ชนปเท อุสฺสหนฺติ ราชโปริเส, เตสํ ภวํ ราชา ภตฺตเวตนํ ปกปฺเปตุ.\csromanspacer{} เต จ มนุสฺสา สกมฺมปสุตา รญฺโญ ชนปทํ น วิเหเฐสฺสนฺติ, มหา จ รญฺโญ ราสิโก ภวิสฺสติ.\csromanspacer{} เขมฏฺฐิตา ชนปทา อกณฺฏกา อนุปฺปีฬา มนุสฺสา มุทา โมทมานา อุเร ปุตฺเต นจฺเจนฺตา อปารุตฆรา มญฺเญ วิหริสฺสนฺตี”ติ. “เอวํ โภ”ติ โข พฺราหฺมณ ราชา มหาวิชิโต ปุโรหิตสฺส พฺราหฺมณสฺส ปฏิสฺสุตฺวา เย รญฺโญ ชนปเท อุสฺสหิํสุ กสิโครกฺเข, เตสํ ราชา มหาวิชิโต พีชภตฺตํ อนุปฺปทาสิ.}

\vspace{\tipitakaparskip}
\csromancenter{กูฏทนฺตสุตฺต เย จ รญฺโญ ชนปเท อุสฺสหิํสุ วาณิชฺชาย, เตสํ ราชา มหาวิชิโต ปาภตํ อนุปฺปทาสิ.\csromanspacer{} เย จ รญฺโญ ชนปเท อุสฺสหิํสุ ราชโปริเส, เตสํ ราชา มหาวิชิโต ภตฺตเวตนํ ปกปฺเปสิ.\csromanspacer{} เต จ มนุสฺสา สกมฺมปสุตา รญฺโญ ชนปทํ น วิเหฐิํสุ, มหา จ รญฺโญ ราสิโก อโหสิ.\csromanspacer{} เขมฏฺฐิตา ชนปทา อกณฺฏกา อนุปฺปีฬา มนุสฺสา มุทา โมทมานา อุเร ปุตฺเต นจฺเจนฺตา อปารุตฆรา มญฺเญ วิหริํสุ.\csromanspacer{} อถ โข พฺราหฺมณ ราชา มหาวิชิโต ปุโรหิตํ พฺราหฺมณํ อามนฺเตตฺวา เอตทโวจ “สมูหโต โข เม โภโส ทสฺสุขีโล โภโต สํวิธานํ อาคมฺม มหา จ เม ราสิโก เขมฏฺฐิตา ชนปทา อกณฺฏกา อนุปฺปีฬา มนุสฺสา มุทา โมทมานา อุเร ปุตฺเต นจฺเจนฺตา อปารุตฆรา มญฺเญ วิหรนฺติ.\csromanspacer{} อิจฺฉามหํ พฺราหฺมณ มหายญฺญํ ยชิตุํ, อนุสาสตุ มํ ภวํ ยํ มม อสฺส ทีฆรตฺตํ หิตาย สุขายา”ติ.}
\csromanmatikaanchor{mtk.84}
\csromantocmarkref{section}{จตุปริกฺขาร}{mtk.84}
\bookmark[dest={mtk.84},level=1]{จตุปริกฺขาร}
\csromanheader{1}{\textbf{จตุปริกฺขาร}}
\proseitem{339}{\csromanfolio{129}เตน หิ ภวํ ราชา เย โภโต รญฺโญ ชนปเท ขตฺติยา อานุยนฺตา เนคมา เจว ชานปทา จ เต ภวํ ราชา อามนฺตยตํ “อิจฺฉามหํ โภ มหายญฺญํ ยชิตุํ, อนุชานนฺตุ เม ภวนฺโต ยํ มม อสฺส ทีฆรตฺตํ หิตาย สุขายา”ติ.\csromanspacer{} เย โภโต รญฺโญ ชนปเท อมจฺจา ปาริสชฺชา เนคมา เจว ชานปทา จ ฯเปฯ พฺราหฺมณมหาสาลา เนคมา เจว ชานปทา จ ฯเปฯ คหปติเนจยิกา เนคมา เจว ชานปทา จ, เต ภวํ ราชา อามนฺตยตํ “อิจฺฉามหํ โภ มหายญฺญํ ยชิตุํ, อนุชานนฺตุ เม ภวนฺโต ยํ มม อสฺส ทีฆรตฺตํ หิตาย สุขายา”ติ. “เอวํ โภ”ติ โข พฺราหฺมณ ราชา มหาวิชิโต ปุโรหิตสฺส พฺราหฺมณสฺส ปฏิสฺสุตฺวา เย รญฺโญ ชนปเท ขตฺติยา อานุยนฺตา เนคมา เจว ชานปทา จ, เต ราชา มหาวิชิโต อามนฺเตสิ “อิจฺฉามหํ โภ มหายญฺญํ ยชิตุํ, อนุชานนฺตุ เม ภวนฺโต ยํ มม อสฺส ทีฆรตฺตํ หิตาย สุขายา”ติ.\csromanspacer{} ยชตํ ภวํ ราชา ยญฺญํ, ยญฺญกาโล มหาราชาติ.\csromanspacer{} เย รญฺโญ ชนปเท อมจฺจา ปาริสชฺชา เนคมา เจว ชานปทา จ ฯเปฯ พฺราหฺมณมหาสาลา เนคมา เจว ชานปทา จ ฯเปฯ คหปติเนจยิกา เนคมา เจว ชานปทา จ, เต ราชา มหาวิชิโต อามนฺเตสิ “อิจฺฉามหํ \csromanfolio{130}โภ มหายญฺญํ ยชิตุํ, อนุชานนฺตุ เม ภวนฺโต ยํ มม อสฺส ทีฆรตฺตํ หิตาย สุขายา”ติ.\csromanspacer{} ยชตํ ภวํ ราชา ยญฺญํ, ยญฺญกาโล มหาราชาติ.\csromanspacer{} อิติเม จตฺตาโร อนุมติปกฺขา ตสฺเสว ยญฺญสฺส ปริกฺขารา ภวนฺติ.}

\csromanmatikaanchor{mtk.85}
\csromantocmarkref{section}{อฏฺฐปริกฺขาร}{mtk.85}
\bookmark[dest={mtk.85},level=1]{อฏฺฐปริกฺขาร}
\csromanheader{1}{\textbf{อฏฺฐปริกฺขาร}}
\proseitem{340}{ราชา มหาวิชิโต อฏฺฐหงฺเคหิ สมนฺนาคโต, อุภโต สุชาโต มาติโต จ ปิติโต จ สํสุทฺธคหณิโก ยาว สตฺตมา ปิตามหยุคา อกฺขิตฺโต อนุปกฺกุฏฺโฐ ชาติวาเทน, อภิรูโป ทสฺสนีโย ปาสาทิโก ปรมาย วณฺณโปกฺขรตาย สมนฺนาคโต พฺรหฺมวณฺณี พฺรหฺมวจฺฉสี อขุทฺทาวกาโส ทสฺสนาย, อฑฺโฒ มหทฺธโน มหาโภโค ปหูตชาตรูปรชโต ปหูตวิตฺตูปกรโณ ปหูตธนธญฺโญ ปริปุณฺณโกสโกฏฺฐาคาโร, พลวา จตุรงฺคินิยา เสนาย สมนฺนาคโต อสฺสวาย โอวาทปฏิกราย สหติ\footnote{ปตปติ (สี, อิ), ตปติ (สฺยา)} มญฺเญ ปจฺจตฺถิเก ยสสา, สทฺโธ ทายโก ทานปติ อนาวฏทฺวาโร สมณพฺราหฺมณ\-กปณทฺธิก วณิพฺพก ยาจกานํ โอปานภูโต ปุญฺญานิ กโรติ, พหุสฺสุโต ตสฺส ตสฺส สุตชาตสฺส, ตสฺส ตสฺเสว โข ปน ภาสิตสฺส อตฺถํ ชานาติ “อยํ อิมสฺส ภาสิตสฺส อตฺโถ อยํ อิมสฺส ภาสิตสฺส อตฺโถ”ติ, ปณฺฑิโต วิยตฺโต เมธาวี ปฏิพโล อตีตานาคตปจฺจุนฺเน อตฺเถ จินฺเตตุํ.\csromanspacer{} ราชา มหาวิชิโต อิเมหิ อฏฺฐหงฺเคหิ สมนฺนาคโต.\csromanspacer{} อิติ อิมานิปิ อฏฺฐงฺคานิ ตสฺเสว ยญฺญสฺส ปริกฺขารา ภวนฺติ.}

\csromanmatikaanchor{mtk.86}
\csromantocmarkref{section}{จตุปริกฺขาร}{mtk.86}
\bookmark[dest={mtk.86},level=1]{จตุปริกฺขาร}
\csromanheader{1}{\textbf{จตุปริกฺขาร}}
\proseitem{341}{ปุโรหิโต\footnote{ปุโรหิโตปิ (กสี, ก)} พฺราหฺมโณ จตูหงฺเคหิ สมนฺนาคโต, อุภโต สุชาโต มาติโต จ ปิติโต จ สํสุทฺธคหณิโก ยาว สตฺตมา ปิตามหยุคา อกฺขิตฺโต อนุปกฺกุฏฺโฐ ชาติวาเทน, อชฺฌายโก มนฺตธโร ติณฺณํ เวทานํ ปารคู สนิฆณฺฑุเกฏุภานํ สากฺขรปฺปเภทานํ อิติหาสปญฺจมานํ ปทโก เวยฺยากรโณ โลกายตมหาปุริสลกฺขเณสุ อนวโย, สีลวา วุทฺธสีลี}

\vspace{\tipitakaparskip}
\csromancenter{กูฏทนฺตสุตฺต วุทฺธสีเลน สมนฺนาคโต, ปณฺฑิโต วิยตฺโต เมธาวี ปฐโม วา ทุติโย วา สุชํ ปคฺคณฺหนฺตานํ.\csromanspacer{} ปุโรหิโต พฺราหฺมโณ อิเมหิ จตูหงฺเคหิ สมนฺนาคโต.\csromanspacer{} อิติ อิมานิปิ จตฺตาริ องฺคานิ ตสฺเสว ยญฺญสฺส ปริกฺขารา ภวนฺติ.}
\csromanmatikaanchor{mtk.87}
\csromantocmarkref{section}{ติสฺโส วิธา}{mtk.87}
\bookmark[dest={mtk.87},level=1]{ติสฺโส วิธา}
\csromanheader{1}{\textbf{ติสฺโส วิธา}}
\csromanmatikaanchor{mtk.88}
\csromantocmarkref{section}{ทสอาการ}{mtk.88}
\bookmark[dest={mtk.88},level=1]{ทสอาการ}
\proseitem{342}{\csromanfolio{131}อถ โข พฺราหฺมณ ปุโรหิโต พฺราหฺมโณ รญฺโญ มหาวิชิตสฺส ปุพฺเพว ยญฺญา ติสฺโส วิธา เทเสสิ.\csromanspacer{} สิยา โข ปน โภโต รญฺโญ มหายญฺญํ ยิฏฺฐุกามสฺส\footnote{ยิฏฺฐกามสฺส (ก)} โกจิเทว วิปฺปฏิสาโร “มหา วต เม โภคกฺขนฺโธ วิคจฺฉิสฺสตี”ติ, โส โภตา รญฺญา วิปฺปฏิสาโร น กรณีโย.\csromanspacer{} สิยา โข ปน โภโต รญฺโญ มหายญฺญํ ยชมานสฺส โกจิเทว วิปฺปฏิสาโร “มหา วต เม โภคกฺขนฺโธ วิคจฺฉตี”ติ, โส โภตา รญฺญา วิปฺปฏิสาโร น กรณีโย.\csromanspacer{} สิยา โข ปน โภโต รญฺโญ มหายญฺญํ ยิฏฺฐสฺส โกจิเทว วิปฺปฏิสาโร “มหา วต เม โภคกฺขนฺโธ วิคโต”ติ, โส โภตา รญฺญา วิปฺปฏิสาโร น กรณีโยติ.\csromanspacer{} อิมา โข พฺราหฺมณ ปุโรหิโต พฺราหฺมโณ รญฺโญ มหาวิชิตสฺส ปุพฺเพว ยญฺญา ติสฺโส วิธา เทเสสิ.}

\vspace{\tipitakaparskip}
\csromancenter{\textbf{ทสอาการ.}}
\proseitem{343}{อถ โข พฺราหฺมณ ปุโรหิโต พฺราหฺมโณ รญฺโญ มหาวิชิตสฺส ปุพฺเพว ยญฺญา ทสหากาเรหิ ปฏิคฺคาหเกสุ วิปฺปฏิสารํ ปฏิวิเนสิ.\csromanspacer{} อาคมิสฺสนฺติ โข โภโต ยญฺญํ ปาณาติปาติโนปิ ปาณาติปาตา ปฏิวิรตาปิ, เย ตตฺถ ปาณาติปาติโน, เตสํเยว เตน, เย ตตฺถ ปาณาติปาตา ปฏิวิรตา, เต อารพฺภ ยชตํ ภวํ, สชฺชตํ ภวํ, โมทตํ ภวํ, จิตฺตเมว ภวํ อนฺตรํ ปสาเทตุ.\csromanspacer{} อาคมิสฺสนฺติ โข โภโต ยญฺญํ อทินฺนาทายิโนปิ อทินฺนาทานา ปฏิวิรตาปิ ฯเปฯ.\csromanspacer{} กาเมสุ มิจฺฉาจาริโนปิ กาเมสุมิจฺฉาจารา ปฏิวิรตาปิ.\csromanspacer{} มุสาวาทิโนปิ มุสาวาทา ปฏิวิรตาปิ.\csromanspacer{} ปิสุณวาจิโนปิ ปิสุณาย วาจาย ปฏิวิรตาปิ.\csromanspacer{} ผรุสวาจิโนปิ ผรุสาย วาจาย ปฏิวิรตาปิ.\csromanspacer{} สมฺผปฺปลาปิโนปิ \csromanfolio{132}สมฺผปฺปลาปา ปฏิวิรตาปิ.\csromanspacer{} อภิชฺฌาลุโนปิ อนภิชฺฌาลุโนปิ.\csromanspacer{} พฺยาปนฺนจิตฺตาปิ อพฺยาปนฺนจิตฺตาปิ.\csromanspacer{} มิจฺฉาทิฏฺฐิกาปิ สมฺมาทิฏฺฐิกาปิ, เย ตตฺถ มิจฺฉาทิฏฺฐิกา, เตสํเยว เตน, เย ตตฺถ สมฺมาทิฏฺฐิกา, เต อารพฺภ ยชตํ ภวํ, สชฺชตํ ภวํ, โมทตํ ภวํ, จิตฺตเมว ภวํ อนฺตรํ ปสาเทตูติ.\csromanspacer{} อิเมหิ โข พฺราหฺมณ ปุโรหิโต พฺราหฺมโณ รญฺโญ มหาวิชิตสฺส ปุพฺเพว ยญฺญา ทสหากาเรหิ ปฏิคฺคาหเกสุ วิปฺปฏิสารํ ปฏิวิเนสิ.}

\csromanmatikaanchor{mtk.89}
\csromantocmarkref{section}{โสฬสาการ}{mtk.89}
\bookmark[dest={mtk.89},level=1]{โสฬสาการ}
\csromanheader{1}{\textbf{โสฬสาการ}}
\proseitem{344}{อถ โข พฺราหฺมณ ปุโรหิโต พฺราหฺมโณ รญฺโญ มหาวิชิตสฺส มหายญฺญํ ยชมานสฺส โสฬสหากาเรหิ จิตฺตํ สนฺทสฺเสสิ สมาทเปสิ สมุตฺเตเชสิ สมฺปหํเสสิ.\csromanspacer{} สิยา โข ปน โภโต รญฺโญ มหายญฺญํ ยชมานสฺส โกจิเทว วตฺตา “ราชา โข มหาวิชิโต มหายญฺญํ ยชติ, โน จ โข ตสฺส อามนฺติตา ขตฺติยา อานุยนฺตา เนคมา เจว ชานปทา จ, อถ จ ปน ภวํ ราชา เอวรูปํ มหายญฺญํ ยชตี”ติ, เอวมฺปิ โภโต รญฺโญ วตฺตา ธมฺมโต นตฺถิ.\csromanspacer{} โภตา โข ปน รญฺญา อามนฺติตา ขตฺติยา อานุยนฺตา เนคมา เจว ชานปทา จ.\csromanspacer{} อิมินาเปตํ ภวํ ราชา ชานาตุ, ยชตํ ภวํ, สชฺชตํ ภวํ, โมทตํ ภวํ, จิตฺตเมว ภวํ อนฺตรํ ปสาเทตุ. (1)}

\prose{สิยา โข ปน โภโต รญฺโญ มหายญฺญํ ยชมานสฺส โกจิเทว วตฺตา “ราชา โข มหาวิชิโต มหายญฺญํ ยชติ, โน จ โข ตสฺส อามนฺติตา อมจฺจา ปาริสชฺชา เนคมา เจว ชานปทา จ ฯเปฯ พฺราหฺมณมหาสาลา เนคมา เจว ชานปทา จ ฯเปฯ คหปติเนจยิกา เนคมา เจว ชานปทา จ, อถ จ ปน ภวํ ราชา เอวรูปํ มหายญฺญํ ยชตี”ติ, เอวมฺปิ โภโต รญฺโญ วตฺตา ธมฺมโต นตฺถิ.\csromanspacer{} โภตา โข ปน รญฺญา อามนฺติตา คหปติเนจยิกา เนคมา เจว ชานปทา จ.\csromanspacer{} อิมินาเปตํ ภวํ ราชา ชานาตุ, ยชตํ ภวํ, สชฺชตํ ภวํ, โมทตํ ภวํ, จิตฺตเมว ภวํ อนฺตรํ ปสาเทตุ. (4)}

\prose{สิยา โข ปน โภโต รญฺโญ มหายญฺญํ ยชมานสฺส โกจิเทว วตฺตา “ราชา โข มหาวิชิโต มหายญฺญํ ยชติ, โน จ โข}

\vspace{\tipitakaparskip}
\csromancenter{กูฏทนฺตสุตฺต อุภโต สุชาโต มาติโต จ ปิติโต จ สํสุทฺธคหณิโก ยาว สตฺตมา ปิตามหยุคา อกฺขิตฺโต อนุปกฺกุฏฺโฐ ชาติวาเทน, อถ จ ปน ภวํ ราชา เอวรูปํ มหายญฺญํ ยชตี”ติ, เอวมฺปิ โภโต รญฺโญ วตฺตา ธมฺมโต นตฺถิ.\csromanspacer{} ภวํ โข ปน ราชา อุภโต สุชาโต มาติโต จ ปิติโต จ สํสุทฺธคหณิโก ยาว สตฺตมา ปิตามหยุคา อกฺขิตฺโต อนุปกฺกุฏฺโฐ ชาติวาเทน.\csromanspacer{} อิมินาเปตํ ภวํ ราชา ชานาตุ, ยชตํ ภวํ, สชฺชตํ ภวํ, โมทตํ ภวํ, จิตฺตเมว ภวํ อนฺตรํ ปสาเทตุ. (5)}
\prose{\csromanfolio{133}สิยา โข ปน โภโต รญฺโญ มหายญฺญํ ยชมานสฺส โกจิเทว วตฺตา “ราชา โข มหาวิชิโต มหายญฺญํ ยชติ, โน จ โข อภิรูโป ทสฺสนีโย ปาสาทิโก ปรมาย วณฺณโปกฺขรตาย สมนฺนาคโต พฺรหฺมวณฺณี พฺรหฺมวจฺฉสี อขุทฺทาวกาโส ทสฺสนาย ฯเปฯ โน จ โข อฑฺโฒ มหทฺธโน มหาโภโค ปหูตชาตรูปรชโต ปหูตวิตฺตูปกรโณ ปหูตธนธญฺโญ ปริปุณฺณโกสโกฏฺฐาคาโร ฯเปฯ โน จ โข พลวา จตุรงฺคินิยา เสนาย สมนฺนาคโต อสฺสวาย โอวาทปฏิกราย สหติ มญฺเญ ปจฺจตฺถิเก ยสสา ฯเปฯ โน จ โข สทฺโธ ทายโก ทานปติ อนาวฏทฺวาโร สมณพฺราหฺมณกปณทฺธิกวณิพฺพกยาจกานํ โอปานภูโต ปุญฺญานิ กโรติ ฯเปฯ โน จ โข พหุสฺสุโต ตสฺส ตสฺส สุตชาตสฺส ฯเปฯ โน จ โข ตสฺส ตสฺเสว โข ปน ภาสิตสฺส อตฺถํ ชานาติ “อยํ อิมสฺส ภาสิตสฺส อตฺโถ, อยํ อิมสฺส ภาสิตสฺส อตฺโถ”ติ ฯเปฯ โน จ โข ปณฺฑิโต วิยตฺโต เมธาวี ปฏิพโล อตีตานาคตปจฺจุปฺปนฺเน อตฺเถ จินฺเตตุํ, อถ จ ปน ภวํ ราชา เอวรูปํ มหายญฺญํ ยชตี”ติ, เอวมฺปิ โภโต รญฺโญ วตฺตา ธมฺมโต นตฺถิ.\csromanspacer{} ภวํ โข ปน ราชา ปณฺฑิโต วิยตฺโต เมธาวี ปฏิพโล อตีตานาคตปจฺจุปฺปนฺเน อตฺเถ จินฺเตตุํ.\csromanspacer{} อิมินาเปตํ ภวํ ราชา ชานาตุ, ยชตํ ภวํ, สชฺชตํ ภวํ, โมทตํ ภวํ, จิตฺตเมว ภวํ อนฺตรํ ปสาเทตุ. (12)}

\prose{สิยา โข ปน โภโต รญฺโญ มหายญฺญํ ยชมานสฺส โกจิเทว วตฺตา “ราชา โข มหาวิชิโต มหายญฺญํ ยชติ, โน จ ขฺวสฺส ปุโรหิโต พฺราหฺมโณ อุภโต สุชาโต มาติโต จ ปิติโต จ \csromanfolio{134}สํสุทฺธคหณิโก ยาว สตฺตมา ปิตามหยุคา อกฺขิตฺโต อนุปกฺกุฏฺโฐ ชาติวาเทน, อถ จ ปน ภวํ ราชา เอวรูปํ มหายญฺญํ ยชตี”ติ, เอวมฺปิ โภโต รญฺโญ วตฺตา ธมฺมโต นตฺถิ.\csromanspacer{} โภโต โข ปน รญฺโญ ปุโรหิโต พฺราหฺมโณ อุภโต สุชาโต มาติโต จ ปิติโต จ สํสุทฺธคหณิโก ยาว สตฺตมา ปิตามหยุคา อกฺขิตฺโต อนุปกฺกุฏฺโฐ ชาติวาเทน.\csromanspacer{} อิมินาเปตํ ภวํ ราชา ชานาตุ, ยชตํ ภวํ สชฺชตํ ภวํ, โมทตํ ภวํ, จิตฺตเมว ภวํ อนฺตรํ ปสาเทตุ. (13)}

\prose{สิยา โข ปน โภโต รญฺโญ มหายญฺญํ ยชมานสฺส โกจิเทว วตฺตา “ราชา โข มหาวิชิโต มหายญฺญํ ยชติ, โน จ ขฺวสฺส ปุโรหิโต พฺราหฺมโณ อชฺฌายโก มนฺตธโร ติณฺณํ เวทานํ ปารคู สนิฆณฺฑุเกฏุภานํ สากฺขรปฺปเภทานํ อิติหาสปญฺจมานํ ปทโก เวยฺยากรโณ โลกายตมหาปุริสลกฺขเณสุ อนวโย ฯเปฯ โน จ ขฺวสฺส ปุโรหิโต พฺราหฺมโณ สีลวา วุทฺธสีลี วุทฺธสีเลน สมนฺนาคโต ฯเปฯ โน จ ขฺวสฺส ปุโรหิโต พฺราหฺมโณ ปณฺฑิโต วิยตฺโต เมธาวี ปฐโม วา ทุติโย วา สุชํ ปคฺคณฺหนฺตานํ, อถ จ ปน ภวํ ราชา เอวรูปํ มหายญฺญํ ยชตี”ติ, เอวมฺปิ โภโต รญฺโญ วตฺตา ธมฺมโต นตฺถิ.\csromanspacer{} โภโต โข ปน รญฺโญ ปุโรหิโต พฺราหฺมโณ ปณฺฑิโต วิยตฺโต เมธาวี ปฐโม วา ทุติโย วา สุชํ ปคฺคณฺหนฺตานํ.\csromanspacer{} อิมินาเปตํ ภวํ ราชา ชานาตุ, ยชตํ ภวํ, สชฺชตํ ภวํ, โมทตํ ภวํ, จิตฺตเมว ภวํ อนฺตรํ ปสาเทตูติ.\csromanspacer{} อิเมหิ โข พฺราหฺมณ ปุโรหิโต พฺราหฺมโณ รญฺโญ มหาวิชิตสฺส มหายญฺญํ ยชมานสฺส โสฬสหิ อากาเรหิ จิตฺตํ สนฺทสฺเสสิ สมาทเปสิ สมุตฺเตเชสิ สมฺปหํเสสิ. (16)}

\proseitem{345}{ตสฺมิํ โข พฺราหฺมณ ยญฺเญ เนว คาโว หญฺญิํสุ, น อเชฬกา หญฺญิํสุ, น กุกฺกุฏสูกรา หญฺญิํสุ, น วิวิธา ปาณา สํฆาตํ อาปชฺชิํสุ, น รุกฺขา ฉิชฺชิํสุ ยูปตฺถาย, น ทพฺภา ลูยิํสุ พริหิสตฺถาย\footnote{ปริหิํสตฺถาย (สฺยา, กสี, ก), ปรหิํสตฺถาย (ก)}, เยปิสฺส อเหสุํ ทาสาติ วา เปสฺสาติ วา กมฺมกราติ วา, เตปิ น ทณฺฑตชฺชิตา}

\vspace{\tipitakaparskip}
\csromancenter{กูฏทนฺตสุตฺต น ภยตชฺชิตา น อสฺสุมุขา รุทมานา ปริกมฺมานิ อกํสุ.\csromanspacer{} อถ โข เย อิจฺฉิํสุ, เต อกํสุ, เย น อิจฺฉิํสุ, น เต อกํสุ, ยํ อิจฺฉิํสุ, ตํ อกํสุ, ยํ น อิจฺฉิํสุ, น ตํ อกํสุ, สปฺปิเตลนวนีตทธิมธุผาณิเตน เจว โส ยญฺโญ นิฏฺฐานมคมาสิ.}
\proseitem{346}{\csromanfolio{135}อถ โข พฺราหฺมณ ขตฺติยา อานุยนฺตา เนคมา เจว ชานปทา จ.\csromanspacer{} อมจฺจา ปาริสชฺชา เนคมา เจว ชานปทา จ.\csromanspacer{} พฺราหฺมณมหาสาลา เนคมา เจว ชานปทา จ.\csromanspacer{} คหปติเนจยิกา เนคมา เจว ชานปทา จ ปหูตํ สาปเตยฺยํ อาทาย ราชานํ มหาวิชิตํ อุปสงฺกมิตฺวา เอวมาหํสุ “อิทํเทว ปหูตํ สาปเตยฺยํ เทวญฺเญว อุทฺทิสฺสาภตํ, ตํ เทโว ปฏิคฺคณฺหาตู”ติ.\csromanspacer{} อลํ โภ มมาปิทํ ปหูตํ สาปเตยฺยํ ธมฺมิเกน พลินา อภิสงฺขตํ, ตญฺจ โว โหตุ, อิโต จ ภิยฺโย หรถาติ.\csromanspacer{} เต รญฺญา ปฏิกฺขิตฺตา เอกมนฺตํ อปกฺกมฺม เอวํ สมจินฺเตสุํ “น โข เอตํ อมฺหากํ ปติรูปํ, ยํ มยํ อิมานิ สาปเตยฺยานิ ปุนเทว สกานิ ฆรานิ ปฏิหเรยฺยาม, ราชา โข มหาวิชิโต มหายญฺญํ ยชติ, หนฺทสฺส มยํ อนุยาคิโน โหมา”ติ.}

\proseitem{347}{อถ โข พฺราหฺมณ ปุรตฺถิเมน ยญฺญวาฏสฺส\footnote{ยญฺญาวาฏสฺส (สี, อิ, ก)} ขตฺติยา อานุยนฺตา เนคมา เจว ชานปทา จ ทานานิ ปฏฺฐเปสุํ.\csromanspacer{} ทกฺขิเณน ยญฺญวาฏสฺส อมจฺจา ปาริสชฺชา เนคมา เจว ชานปทา จ ทานานิ ปฏฺฐเปสุํ.\csromanspacer{} ปจฺฉิเมน ยญฺญวาฏสฺส พฺราหฺมณมหาสาลา เนคมา เจว ชานปทา จ ทานานิ ปฏฺฐเปสุํ.\csromanspacer{} อุตฺตเรน ยญฺญวาฏสฺส คหปติเนจยิกา เนคมา เจว ชานปทา จ ทานานิ ปฏฺฐเปสุํ.}

\prose{เตสุปิ โข พฺราหฺมณ ยญฺเญสุ เนว คาโว หญฺญิํสุ, น อเชฬกา หญฺญิํสุ, น กุกฺกุฏสูกรา หญฺญิํสุ, น วิวิธา ปาณา สํฆาตํ อาปชฺชิํสุ, น รุกฺขา ฉิชฺชิํสุ ยูปตฺถาย, น ทพฺภา ลูยิํสุ พริหิสตฺถาย.\csromanspacer{} เยปิ เนสํ อเหสุํ ทาสาติ วา เปสฺสาติ วา กมฺมกราติ วา, เตปิ น ทณฺฑตชฺชิตา น ภยตชฺชิตา น อสฺสุมุขา รุทมานา ปริกมฺมานิ อกํสุ.\csromanspacer{} อถ โข เย อิจฺฉิํสุ, เต อกํสุ, เย น อิจฺฉิํสุ, น เต อกํสุ, ยํ อิจฺฉิํสุ, ตํ อกํสุ, ยํ น อิจฺฉิํสุ, น ตํ อกํสุ, สปฺปิเตลนวนีตทธิมธุผาณิเตน เจว เต ยญฺญา นิฏฺฐานมคมํสุ.}

\prose{\csromanfolio{136}อิติ จตฺตาโร จ อนุมติปกฺขา, ราชา มหาวิชิโต อฏฺฐหงฺเคหิ สมนฺนาคโต, ปุโรหิโต พฺราหฺมโณ จตูหงฺเคหิ สมนฺนาคโต, ติสฺโส จ วิธา.\csromanspacer{} อยํ วุจฺจติ พฺราหฺมณ ติวิธา ยญฺญสมฺปทา โสฬสปริกฺขาราติ.}

\proseitem{348}{เอวํ วุตฺเต เต พฺราหฺมณา อุนฺนาทิโน อุจฺจาสทฺทมหาสทฺทา อเหสุํ “อโห ยญฺโญ, อโห ยญฺญสมฺปทา”ติ.\csromanspacer{} กูฏทนฺโต ปน พฺราหฺมโณ ตุณฺหีภูโตว นิสินฺโน โหติ.\csromanspacer{} อถ โข เต พฺราหฺมณา กูฏทนฺตํ พฺราหฺมณํ เอตทโวจุํ “กสฺมา ปน ภวํ กูฏทนฺโต สมณสฺส โคตมสฺส สุภาสิตํ สุภาสิตโต นาพฺภนุโมทตี”ติ.\csromanspacer{} นาหํ โภ สมณสฺส โคตมสฺส สุภาสิตํ สุภาสิตโต นาพฺภนุโมทามิ.\csromanspacer{} มุทฺธาปิ ตสฺส วิปเตยฺย, โย สมณสฺส โคตมสฺส สุภาสิตํ สุภาสิตโต นาพฺภนุโมเทยฺย.\csromanspacer{} อปิ จ เม โภ เอวํ โหติ, สมโณ โคตโม น เอวมาห “เอวํ เม สุตนฺ”ติ วา “เอวํ อรหติ ภวิตุนฺ”ติ วา.\csromanspacer{} อปิ จ สมโณ โคตโม “เอวํ ตทา อาสิ, อิตฺถํ ตทา อาสิ”เตฺวว ภาสติ.\csromanspacer{} ตสฺส มยฺหํ โภ เอวํ โหติ “อทฺธา สมโณ โคตโม เตน สมเยน ราชา วา อโหสิ มหาวิชิโต ยญฺญสฺสามิ ปุโรหิโต วา พฺราหฺมโณ ตสฺส ยญฺญสฺส ยาเชตา”ติ.\csromanspacer{} อภิชานาติ ปน ภวํ โคตโม เอวรูปํ ยญฺญํ ยชิตฺวา วา ยาเชตฺวา วา กายสฺส เภทา ปรํ มรณา สุคติํ สคฺคํ โลกํ อุปปชฺชิตาติ.\csromanspacer{} อภิชานามหํ พฺราหฺมณ เอวรูปํ ยญฺญํ ยชิตฺวา วา ยาเชตฺวา วา กายสฺส เภทา ปรํ มรณา สุคติํ สคฺคํ โลกํ อุปปชฺชิตา, อหํ เตน สมเยน ปุโรหิโต พฺราหฺมโณ อโหสิํ ตสฺส ยญฺญสฺส ยาเชตาติ.}

\csromanmatikaanchor{mtk.90}
\csromantocmarkref{section}{นิจฺจทาน อนุกุลยญฺญ}{mtk.90}
\bookmark[dest={mtk.90},level=1]{นิจฺจทาน อนุกุลยญฺญ}
\csromanheader{1}{\textbf{นิจฺจทาน อนุกุลยญฺญ}}
\proseitem{349}{อตฺถิ ปน โภ โคตม อญฺโญ ยญฺโญ อิมาย ติวิธาย ยญฺญสมฺปทาย\footnote{ติวิธยญฺญสมฺปทาย (ก)} โสฬสปริกฺขาราย อปฺปฏฺฐตโร\footnote{อปฺปตฺถตโร (สฺยา, กํ)} จ อปฺปสมารมฺภตโร\footnote{อปฺปสมารพฺภตโร (สี, อิ, ก)} จ มหปฺผลตโร จ มหานิสํสตโร จาติ.}

\csromanheader{2}{5. กูฏทนฺตสุตฺต}
\prose{\csromanfolio{137}อตฺถิ โข พฺราหฺมณ อญฺโญ ยญฺโญ อิมาย ติวิธาย ยญฺญสมฺปทาย โสฬสปริกฺขาราย อปฺปฏฺฐตโร จ อปฺปสมารมฺภตโร จ มหปฺผลตโร จ มหานิสํสตโร จาติ.}

\prose{กตโม ปน โส โภ โคตม ยญฺโญ อิมาย ติวิธาย ยญฺญสมฺปทาย โสฬสปริกฺขาราย อปฺปฏฺฐตโร จ อปฺปสมารมฺภตโร จ มหปฺผลตโร จ มหานิสํสตโร จาติ.}

\prose{ยานิ โข ปน ตานิ พฺราหฺมณ นิจฺจทานานิ อนุกุลยญฺญานิ สีลวนฺเต ปพฺพชิเต อุทฺทิสฺส ทิยฺยนฺติ, อยํ โข พฺราหฺมณ ยญฺโญ อิมาย ติวิธาย ยญฺญสมฺปทาย โสฬสปริกฺขาราย อปฺปฏฺฐตโร จ อปฺปสมารมฺภตโร จ มหปฺผลตโร จ มหานิสํสตโร จาติ.}

\prose{โก นุ โข โภ โคตม เหตุ โก ปจฺจโย, เยน ตํ นิจฺจทานํ อนุกุลยญฺญํ อิมาย ติวิธาย ยญฺญสมฺปทาย โสฬสปริกฺขาราย อปฺปฏฺฐตรญฺจ อปฺปสมารมฺภตรญฺจ มหปฺผลตรญฺจ มหานิสํสตรญฺจาติ.}

\prose{น โข พฺราหฺมณ เอวรูปํ ยญฺญํ อุปสงฺกมนฺติ อรหนฺโต วา อรหตฺตมคฺคํ วา สมาปนฺนา.\csromanspacer{} ตํ กิสฺส เหตุ, ทิสฺสนฺติ เหตฺถ พฺราหฺมณ ทณฺฑปฺปหาราปิ คลคฺคหาปิ, ตสฺมา เอวรูปํ ยญฺญํ น อุปสงฺกมนฺติ อรหนฺโต วา อรหตฺตมคฺคํ วา สมาปนฺนา.\csromanspacer{} ยานิ โข ปน ตานิ พฺราหฺมณ นิจฺจทานานิ อนุกุลยญฺญานิ สีลวนฺเต ปพฺพชิเต อุทฺทิสฺส ทิยฺยนฺติ, เอวรูปํ โข พฺราหฺมณ ยญฺญํ อุปสงฺกมนฺติ อรหนฺโต วา อรหตฺตมคฺคํ วา สมาปนฺนา.\csromanspacer{} ตํ กิสฺส เหตุ, น เหตฺถ พฺราหฺมณ ทิสฺสนฺติ ทณฺฑปฺปหาราปิ คลคฺคหาปิ, ตสฺมา เอวรูปํ ยญฺญํ อุปสงฺกมนฺติ อรหนฺโต วา อรหตฺตมคฺคํ วา สมาปนฺนา.\csromanspacer{} อยํ โข พฺราหฺมณ เหตุ อยํ ปจฺจโย, เยน ตํ นิจฺจทานํ อนุกุลยญฺญํ อิมาย ติวิธาย ยญฺญสมฺปทาย โสฬสปริกฺขาราย อปฺปฏฺฐตรญฺจ อปฺปสมารมฺภตรญฺจ มหปฺผลตรญฺจ มหานิสํสตรญฺจาติ.}

\proseitem{350}{อตฺถิ ปน โภ โคตม อญฺโญ ยญฺโญ อิมาย จ ติวิธาย ยญฺญสมฺปทาย โสฬสปริกฺขาราย อิมินา จ นิจฺจทาเนน อนุกุลยญฺเญน อปฺปฏฺฐตโร จ อปฺปสมารมฺภตโร จ มหปฺผลตโร จ มหานิสํสตโร จาติ.}

\prose{\csromanfolio{138}อตฺถิ โข พฺราหฺมณ อญฺโญ ยญฺโญ อิมาย จ ติวิธาย ยญฺญสมฺปทาย โสฬสปริกฺขาราย อิมินา จ นิจฺจทาเนน อนุกุลยญฺเญน อปฺปฏฺฐตโร จ อปฺปสมารมฺภตโร จ มหปฺผลตโร จ มหานิสํสตโร จาติ.}

\prose{กตโม ปน โส โภ โคตม ยญฺโญ อิมาย จ ติวิธาย ยญฺญสมฺปทาย โสฬสปริกฺขาราย อิมินา จ นิจฺจทาเนน อนุกุลยญฺเญน อปฺปฏฺฐตโร จ อปฺปสมารมฺภตโร จ มหปฺผลตโร จ มหานิสํสตโร จาติ.}

\prose{โย โข พฺราหฺมณ จาตุทฺทิสํ สํฆํ อุทฺทิสฺส วิหารํ กโรติ, อยํ โข พฺราหฺมณ ยญฺโญ อิมาย จ ติวิธาย ยญฺญสมฺปทาย โสฬสปริกฺขาราย อิมินา จ นิจฺจทาเนน อนุกุลยญฺเญน อปฺปฏฺฐตโร จ อปฺปสมารมฺภตโร จ มหปฺผลตโร จ มหานิสํสตโร จาติ.}

\proseitem{351}{อตฺถิ ปน โภ โคตม อญฺโญ ยญฺโญ อิมาย จ ติวิธาย ยญฺญสมฺปทาย โสฬสปริกฺขาราย อิมินา จ นิจฺจทาเนน อนุกุลยญฺเญน อิมินา จ วิหารทาเนน อปฺปฏฺฐตโร จ อปฺปสมารมฺภตโร จ มหปฺผลตโร จ มหานิสํสตโร จาติ.}

\prose{อตฺถิ โข พฺราหฺมณ อญฺโญ ยญฺโญ อิมาย จ ติวิธาย ยญฺญสมฺปทาย โสฬสปริกฺขาราย อิมินา จ นิจฺจทาเนน อนุกุลยญฺเญน อิมินา จ วิหารทาเนน อปฺปฏฺฐตโร จ อปฺปสมารมฺภตโร จ มหปฺผลตโร จ มหานิสํสตโร จาติ.}

\prose{กตโม ปน โส โภ โคตม ยญฺโญ อิมาย จ ติวิธาย ยญฺญสมฺปทาย โสฬสปริกฺขาราย อิมินา จ นิจฺจทาเนน อนุกุลยญฺเญน อิมินา จ วิหารทาเนน อปฺปฏฺฐตโร จ อปฺปสมารมฺภตโร จ มหปฺผลตโร จ มหานิสํสตโร จาติ.}

\prose{โย โข พฺราหฺมณ ปสนฺนจิตฺโต พุทฺธํ สรณํ คจฺฉติ, ธมฺมํ สรณํ คจฺฉติ, สํฆํ สรณํ คจฺฉติ, อยํ โข พฺราหฺมณ ยญฺโญ อิมาย จ ติวิธาย ยญฺญสมฺปทาย โสฬสปริกฺขาราย อิมินา จ นิจฺจทาเนน อนุกุลยญฺเญน อิมินา จ วิหารทาเนน อปฺปฏฺฐตโร จ อปฺปสมารมฺภตโร จ มหปฺผลตโร จ มหานิสํสตโร จาติ.}

\csromanheader{2}{5. กูฏทนฺตสุตฺต}
\proseitem{352}{\csromanfolio{139}อตฺถิ ปน โภ โคตม อญฺโญ ยญฺโญ อิมาย จ ติวิธาย ยญฺญสมฺปทาย โสฬสปริกฺขาราย อิมินา จ นิจฺจทาเนน อนุกุลยญฺเญน อิมินา จ วิหารทาเนน อิเมหิ จ สรณคมเนหิ อปฺปฏฺฐตโร จ อปฺปสมารมฺภตโร จ มหปฺผลตโร จ มหานิสํสตโร จาติ.}

\prose{อตฺถิ โข พฺราหฺมณ อญฺโญ ยญฺโญ อิมาย จ ติวิธาย ยญฺญสมฺปทาย โสฬสปริกฺขาราย อิมินา จ นิจฺจทาเนน อนุกุลยญฺเญน อิมินา จ วิหารทาเนน อิเมหิ จ สรณคมเนหิ อปฺปฏฺฐตโร จ อปฺปสมารมฺภตโร จ มหปฺผลตโร จ มหานิสํสตโร จาติ.}

\prose{กตโม ปน โส โภ โคตม ยญฺโญ อิมาย จ ติวิธาย ยญฺญสมฺปทาย โสฬสปริกฺขาราย อิมินา จ นิจฺจทาเนน อนุกุลยญฺเญน อิมินา จ วิหารทาเนน อิเมหิ จ สรณคมเนหิ อปฺปฏฺฐตโร จ อปฺปสมารมฺภตโร จ มหปฺผลตโร จ มหานิสํสตโร จาติ.}

\prose{โย โข พฺราหฺมณ ปสนฺนจิตฺโต สิกฺขาปทานิ สมาทิยติ ปาณติปาตา เวรมณิํ, อทินฺนาทานา เวรมณิํ, กาเมสุมิจฺฉาจารา เวรมณิํ, มุสาวาทา เวรมณิํ, สุรา\-เมรย\-มชฺช\-ปมาทฏฺฐานา เวรมณิํ.\csromanspacer{} อยํ โข พฺราหฺมณ ยญฺโญ อิมาย จ ติวิธาย ยญฺญสมฺปทาย โสฬสปริกฺขาราย อิมินา จ นิจฺจทาเนน อนุกุลยญฺเญน อิมินา จ วิหารทาเนน อิเมหิ จ สรณคมเนหิ อปฺปฏฺฐตโร จ อปฺปสมารมฺภตโร จ มหปฺผลตโร จ มหานิสํสตโร จาติ.}

\proseitem{353}{อตฺถิ ปน โภ โคตม อญฺโญ ยญฺโญ อิมาย จ ติวิธาย ยญฺญสมฺปทาย โสฬสปริกฺขาราย อิมินา จ นิจฺจทาเนน อนุกุลยญฺเญน อิมินา จ วิหารทาเนน อิเมหิ จ สรณคมเนหิ อิเมหิ จ สิกฺขาปเทหิ อปฺปฏฺฐตโร จ อปฺปสมารมฺภตโร จ มหปฺผลตโร จ มหานิสํสตโร จาติ.}

\prose{อตฺถิ โข พฺราหฺมณ อญฺโญ ยญฺโญ อิมาย จ ติวิธาย ยญฺญสมฺปทาย โสฬสปริกฺขาราย อิมินา จ นิจฺจทาเนน อนุกุลยญฺเญน อิมินา จ วิหารทาเนน อิเมหิ จ สรณคมเนหิ อิเมหิ จ สิกฺขาปเทหิ \csromanfolio{140}อปฺปฏฺฐตโร จ อปฺปสมารมฺภตโร จ มหปฺผลตโร จ มหานิสํสตโร จาติ.}

\prose{กตโม ปน โส โภ โคตม ยญฺโญ อิมาย จ ติวิธาย ยญฺญสมฺปทาย โสฬสปริกฺขาราย อิมินา จ นิจฺจทาเนน อนุกุลยญฺเญน อิมินา จ วิหารทาเนน อิเมหิ จ สรณคมเนหิ อิเมหิ จ สิกฺขาปเทหิ อปฺปฏฺฐตโร จ อปฺปสมารมฺภตโร จ มหปฺผลตโร จ มหานิสํสตโร จาติ.}

\prose{อิธ พฺราหฺมณ ตถาคโต โลเก อุปฺปชฺชติ อรหํ สมฺมาสมฺพุทฺโธ ฯเปฯ. (ยถา สามญฺญผเล, เอวํ วิตฺถาเรตพฺพํ.) เอวํ โข พฺราหฺมณ ภิกฺขุ สีลสมฺปนฺโน โหติ ฯเปฯ ปฐมํ ฌานํ อุปสมฺปชฺช วิหรติ.\csromanspacer{} อยํ โข พฺราหฺมณ ยญฺโญ ปุริเมหิ ยญฺเญหิ อปฺปฏฺฐตโร จ อปฺปสมารมฺภตโร จ มหปฺผลตโร จ มหานิสํสตโร จ.}

\prose{ฯเปฯ ทุติยํ ฌานํ.\csromanspacer{} ตติยํ ฌานํ.\csromanspacer{} จตุตฺถํ ฌานํ อุปสมฺปชฺช วิหรติ.\csromanspacer{} อยํปิ โข พฺราหฺมณ ยญฺโญ ปุริเมหิ ยญฺเญหิ อปฺปฏฺฐตโร จ อปฺปสมารมฺภตโร จ มหปฺผลตโร จ มหานิสํสตโร จาติ.}

\prose{ฯเปฯ ญาณทสฺสนาย จิตฺตํ อภินีหรติ อภินินฺนาเมหิ.\csromanspacer{} อยํปิ โข พฺราหฺมณ ยญฺโญ ปุริเมหิ ยญฺเญหิ อปฺปฏฺฐตโร จ อปฺปสมารมฺภตโร จ มหปฺผลตโร จ มหานิสํสตโร จ ฯเปฯ นาปรํ อิตฺถตฺตายาติ ปชานาติ.\csromanspacer{} อยํปิ โข พฺราหฺมณ ยญฺโญ ปุริเมหิ ยญฺเญหิ อปฺปฏฺฐตโร จ อปฺปสมารมฺภตโร จ มหปฺผลตโร จ มหานิสํสตโร จ.\csromanspacer{} อิมาย จ พฺราหฺมณ ยญฺญสมฺปทาย อญฺญา ยญฺญสมฺปทา อุตฺตริตรา วา ปณีตตรา วา นตฺถีติ.}

\csromanmatikaanchor{mtk.91}
\csromantocmarkref{section}{กูฏทนฺต อุปาสกตฺตปฏิเวทนา}{mtk.91}
\bookmark[dest={mtk.91},level=1]{กูฏทนฺต อุปาสกตฺตปฏิเวทนา}
\csromanheader{1}{\textbf{กูฏทนฺต อุปาสกตฺตปฏิเวทนา}}
\proseitem{354}{เอวํ วุตฺเต กูฏทนฺโต พฺราหฺมโณ ภควนฺตํ เอตทโวจ “อภิกฺกนฺตํ โภ โคตม, อภิกฺกนฺตํ โภ โคตม, เสยฺยถาปิ โภ โคตม นิกฺกุชฺชิตํ วา อุกฺกุชฺเชยฺย, ปฏิจฺฉนฺนํ วา วิวเรยฺย, มูฬฺหสฺส วา มคฺคํ อาจิกฺเขยฺย, อนฺธกาเร วา เตลปชฺโชตํ ธาเรยฺย ‘จกฺขุมนฺโต รูปานิ ทกฺขนฺตี’ติ, เอวเมวํ โภตา โคตเมน อเนกปริยาเยน}

\vspace{\tipitakaparskip}
\csromancenter{กูฏทนฺตสุตฺต ธมฺโม ปกาสิโต, เอสาหํ ภวนฺตํ โคตมํ สรณํ คจฺฉามิ ธมฺมญฺจ ภิกฺขุสํฆญฺจ, อุปาสกํ มํ ภวํ โคตโม ธาเรตุ อชฺชตคฺเค ปาณุเปตํ สรณํ คตํ.\csromanspacer{} เอสาหํ โภ โคตม สตฺต จ อุสภสตานิ สตฺต จ วจฺฉตรสตานิ สตฺต จ วจฺฉตรีสตานิ สตฺต จ อชสตานิ สตฺต จ อุรพฺภสตานิ มุญฺจามิ, ชีวิตํ เทมิ, หริตานิ เจว ติณานิ ขาทนฺตุ, สีตานิ จ ปานียานิ ปิวนฺตุ, สีโต จ เนสํ วาโต อุปวายตู”ติ.}
\csromanmatikaanchor{mtk.92}
\csromantocmarkref{section}{โสตาปตฺติผลสจฺฉิกิริยา}{mtk.92}
\bookmark[dest={mtk.92},level=1]{โสตาปตฺติผลสจฺฉิกิริยา}
\csromanheader{1}{โสตาปตฺติผล\-สจฺฉิกิริยา}
\proseitem{355}{\csromanfolio{141}อถ โข ภควา กูฏทนฺตสฺส พฺราหฺมณสฺส อนุปุพฺพิํ กถํ กเถสิ.\csromanspacer{} เสยฺยถิทํ, ทานกถํ สีลกถํ สคฺคกถํ กามานํ อาทีนวํ โอการํ สํกิเลสํ เนกฺขมฺเม อานิสํสํ ปกาเสสิ.\csromanspacer{} ยทา ภควา อญฺญาสิ กูฏทนฺตํ พฺราหฺมณํ กลฺลจิตฺตํ มุทุจิตฺตํ วินีวรณจิตฺตํ อุทคฺคจิตฺตํ ปสนฺนจิตฺตํ, อถ ยา พุทฺธานํ สามุกฺกํสิกา ธมฺมเทสนา, ตํ ปกาเสสิ ทุกฺขํ สมุทยํ นิโรธํ มคฺคํ.\csromanspacer{} เสยฺยถาปิ นาม สุทฺธํ วตฺถํ อปคตกาฬกํ สมฺมเทว รชนํ ปฏิคฺคเณฺหยฺย, เอวเมว กูฏทนฺตสฺส พฺราหฺมณสฺส ตสฺมิํเยว อาสเน วิรชํ วีตมลํ ธมฺมจกฺขุํ อุทปาทิ “ยํ กิญฺจิ สมุทยธมฺมํ, สพฺพํ ตํ นิโรธธมฺมนฺ”ติ.}

\proseitem{356}{อถ โข กูฏทนฺโต พฺราหฺมโณ ทิฏฺฐธมฺโม ปตฺตธมฺโม วิทิตธมฺโม ปริโยคาฬฺหธมฺโม ติณฺณวิจิกิจฺโฉ วิคตกถํกโถ เวสารชฺชปฺปตฺโต อปรปฺปจฺจโย สตฺถุสาสเน ภควนฺตํ เอตทโวจ “อธิวาเสตุ เม ภวํ โคตโม สฺวาตนาย ภตฺตํ สทฺธิํ ภิกฺขุสํเฆนา”ติ.\csromanspacer{} อธิวาเสสิ ภควา ตุณฺหีภาเวน.}

\proseitem{357}{อถ โข กูฏทนฺโต พฺราหฺมโณ ภควโต อธิวาสนํ วิทิตฺวา อุฏฺฐายาสนา ภควนฺตํ อภิวาเทตฺวา ปทกฺขิณํ กตฺวา ปกฺกามิ.\csromanspacer{} อถ โข กูฏทนฺโต พฺราหฺมโณ ตสฺสา รตฺติยา อจฺจเยน สเก ยญฺญวาเฏ ปณีตํ ขาทนียํ โภชนียํ ปฏิยาทาเปตฺวา ภควโต กาลํ อาโรจาเปสิ “กาโล โภ โคตม นิฏฺฐิตํ ภตฺตนฺ”ติ.}

\proseitem{358}{\csromanfolio{142}อถ โข ภควา ปุพฺพณฺหสมยํ นิวาเสตฺวา ปตฺตจีวรมาทาย สทฺธิํ ภิกฺขุสํเฆน เยน กูฏทนฺตสฺส พฺราหฺมณสฺส ยญฺญวาโฏ เตนุปสงฺกมิ, อุปสงฺกมิตฺวา ปญฺญตฺเต อาสเน นิสีทิ.}

\prosewithcloser{อถ โข กูฏทนฺโต พฺราหฺมโณ พุทฺธปฺปมุขํ ภิกฺขุสํฆํ ปณีเตน ขาทนีเยน โภชนีเยน สหตฺถา สนฺตปฺเปสิ สมฺปวาเรสิ.\csromanspacer{} อถ โข กูฏทนฺโต พฺราหฺมโณ ภควนฺตํ ภุตฺตาวิํ โอนีตปตฺตปาณิํ อญฺญตรํ นีจํ อาสนํ คเหตฺวา เอกมนฺตํ นิสีทิ, เอกมนฺตํ นิสินฺนํ โข กูฏทนฺตํ พฺราหฺมณํ ภควา ธมฺมิยา กถาย สนฺทสฺเสตฺวา สมาทเปตฺวา สมุตฺเตเชตฺวา สมฺปหํเสตฺวา อุฏฺฐายาสนา ปกฺกามีติ.}{\textbf{กูฏทนฺตสุตฺตํ นิฏฺฐิตํ ปญฺจมํ.}}
\csromanmatikaanchor{mtk.93}
\csromantocmarknopage{chapter}{6. มหาลิสุตฺต}{mtk.93}
\bookmark[dest={mtk.93},level=0]{6. มหาลิสุตฺต}
\setcsromanheadcha{6. มหาลิสุตฺต}
\setcsromanheadhclear{}
\chapterheadpageread{6. \textbf{มหาลิสุตฺต}}
\csromanmatikaanchor{mtk.94}
\csromantocmarkref{section}{พฺราหฺมณทูตวตฺถุ}{mtk.94}
\bookmark[dest={mtk.94},level=1]{พฺราหฺมณทูตวตฺถุ}
\csromanheader{1}{\textbf{พฺราหฺมณทูตวตฺถุ}}
\proseitem{359}{\csromanfolio{143}เอวํ เม สุตํ\csromandash{}เอกํ สมยํ ภควา เวสาลิยํ วิหรติ มหาวเน กูฏาคารสาลายํ.\csromanspacer{} เตน โข ปน สมเยน สมฺพหุลา โกสลกา จ พฺราหฺมณทูตา มาคธกา จ พฺราหฺมณทูตา เวสาลิยํ ปฏิวสนฺติ เกนจิเทว กรณีเยน.\csromanspacer{} อสฺโสสุํ โข เต โกสลกา จ พฺราหฺมณทูตา มาคธกา จ พฺราหฺมณทูตา “สมโณ ขลุ โภ โคตโม สกฺยปุตฺโต สกฺยกุลา ปพฺพชิโต เวสาลิยํ วิหรติ มหาวเน กูฏาคารสาลายํ.\csromanspacer{} ตํ โข ปน ภวนฺตํ โคตมํ เอวํ กลฺยาโณ กิตฺติสทฺโท อพฺภุคฺคโต ‘อิติปิ โส ภควา อรหํ สมฺมาสมฺพุทฺโธ วิชฺชาจรณสมฺปนฺโน สุคโต โลกวิทู อนุตฺตโร ปุริสทมฺมสารถิ สตฺถา เทวมนุสฺสานํ พุทฺโธ ภควา’.\csromanspacer{} โส อิมํ โลกํ สเทวกํ สมารกํ สพฺรหฺมกํ สสฺสมณพฺราหฺมณิํ ปชํ สเทวมนุสฺสํ สยํ อภิญฺญา สจฺฉิกตฺวา ปเวเทติ.\csromanspacer{} โส ธมฺมํ เทเสติ อาทิกลฺยาณํ มชฺเฌกลฺยาณํ ปริโยสานกลฺยาณํ สาตฺถํ สพฺยญฺชนํ เกวลปริปุณฺณํ ปริสุทฺธํ พฺรหฺมจริยํ ปกาเสติ.\csromanspacer{} สาธุ โข ปน ตถารูปานํ อรหตํ ทสฺสนํ โหตี”ติ.}

\proseitem{360}{อถ โข เต โกสลกา จ พฺราหฺมณทูตา มาคธกา จ พฺราหฺมณทูตา เยน มหาวนํ กูฏาคารสาลา เตนุปสงฺกมิํสุ.\csromanspacer{} เตน โข ปน สมเยน อายสฺมา นาคิโต ภควโต อุปฏฺฐาโก โหติ.\csromanspacer{} อถ โข เต โกสลกา จ พฺราหฺมณทูตา มาคธกา จ พฺราหฺมณทูตา เยนายสฺมา นาคิโต เตนุปสงฺกมิํสุ, อุปสงฺกมิตฺวา อายสฺมนฺตํ นาคิตํ เอตทโวจุํ “กหํ นุ โข โภ นาคิต เอตรหิ โส ภวํ โคตโม วิหรติ, ทสฺสนกามา หิ มยํ ตํ ภวนฺตํ โคตมนฺ”ติ.\csromanspacer{} อกาโล โข อาวุโส ภควนฺตํ ทสฺสนาย, ปฏิสลฺลีโน ภควาติ.\csromanspacer{} อถ โข เต โกสลกา จ พฺราหฺมณทูตา มาคธกา จ พฺราหฺมณทูตา ตตฺเถว เอกมนฺตํ นิสีทิํสุ “ทิสฺวาว มยํ ตํ ภวนฺตํ โคตมํ คมิสฺสามา”ติ.}

\csromanmatikaanchor{mtk.95}
\csromantocmarkref{section}{โอฏฺฐทฺธลิจฺฉวีวตฺถุ}{mtk.95}
\bookmark[dest={mtk.95},level=1]{โอฏฺฐทฺธลิจฺฉวีวตฺถุ}
\csromanheader{1}{\textbf{โอฏฺฐทฺธลิจฺฉวีวตฺถุ}}
\proseitem{361}{\csromanfolio{144}โอฏฺฐทฺโธปิ ลิจฺฉวี มหติยา ลิจฺฉวีปริสาย สทฺธิํ เยน มหาวนํ กูฏาคารสาลา เยนายสฺมา นาคิโต เตนุปสงฺกมิ, อุปสงฺกมิตฺวา อายสฺมนฺตํ นาคิตํ อภิวาเทตฺวา เอกมนฺตํ อฏฺฐาสิ, เอกมนฺตํ ฐิโต โข โอฏฺฐทฺโธปิ ลิจฺฉวี อายสฺมนฺตํ นาคิตํ เอตทโวจ “กหํ นุ โข ภนฺเต นาคิต เอตรหิ โส ภควา วิหรติ อรหํ สมฺมาสมฺพุทฺโธ, ทสฺสนกามา หิ มยํ ตํ ภควนฺตํ อรหนฺตํ สมฺมาสมฺพุทฺธนฺ”ติ.\csromanspacer{} อกาโล โข มหาลิ ภควนฺตํ ทสฺสนาย, ปฏิสลฺลีโน ภควาติ.\csromanspacer{} โอฏฺฐทฺโธปิ ลิจฺฉวี ตตฺเถว เอกมนฺตํ นิสีทิ “ทิสฺวาว อหํ ตํ ภควนฺตํ คมิสฺสามิ อรหนฺตํ สมฺมาสมฺพุทฺธนฺ”ติ.}

\proseitem{362}{อถ โข สีโห สมณุทฺเทโส เยนายสฺมา นาคิโต เตนุปสงฺกมิ, อุปสงฺกมิตฺวา อายสฺมนฺตํ นาคิตํ อภิวาเทตฺวา เอกมนฺตํ อฏฺฐาสิ, เอกมนฺตํ ฐิโต โข สีโห สมณุทฺเทโส อายสฺมนฺตํ นาคิตํ เอตทโวจ “เอเต ภนฺเต กสฺสป สมฺพหุลา โกสลกา จ พฺราหฺมณทูตา มาคธกา จ พฺราหฺมณทูตา อิธูปสงฺกนฺตา ภควนฺตํ ทสฺสนาย, โอฏฺฐทฺโธปิ ลิจฺฉวี มหติยา ลิจฺฉวีปริสาย สทฺธิํ อิธูปสงฺกนฺโต ภควนฺตํ ทสฺสนาย, สาธุ ภนฺเต กสฺสป ลภตํ เอสา ชนตา ภควนฺตํ ทสฺสนายา”ติ.}

\prose{เตน หิ สีห ตฺวญฺเญว ภควโต อาโรเจหีติ. “เอวํ ภนฺเต”ติ โข สีโห สมณุทฺเทโส อายสฺมโต นาคิตสฺส ปฏิสฺสุตฺวา เยน ภควา เตนุปสงฺกมิ, อุปสงฺกมิตฺวา ภควนฺตํ อภิวาเทตฺวา เอกมนฺตํ อฏฺฐาสิ, เอกมนฺตํ ฐิโต โข สีโห สมณุทฺเทโส ภควนฺตํ เอตทโวจ “เอเต ภนฺเต สมฺพหุลา โกสลกา จ พฺราหฺมณทูตา มาคธกา จ พฺราหฺมณทูตา อิธูปสงฺกนฺตา ภควนฺตํ ทสฺสนาย, โอฏฺฐทฺโธปิ ลิจฺฉวี มหติยา ลิจฺฉวีปริสาย สทฺธิํ อิธูปสงฺกนฺโต ภควนฺตํ ทสฺสนาย, สาธุ ภนฺเต ลภตํ เอสา ชนตา ภควนฺตํ ทสฺสนายา”ติ.\csromanspacer{} เตน หิ สีห วิหารปจฺฉายายํ อาสนํ ปญฺญเปหีติ. “เอวํ ภนฺเต”ติ โข สีโห สมณุทฺเทโส ภควโต ปฏิสฺสุตฺวา วิหารปจฺฉายายํ อาสนํ ปญฺญเปสิ.}

\csromanheader{2}{6. มหาลิสุตฺต}
\proseitem{363}{\csromanfolio{145}อถ โข ภควา วิหารา นิกฺขมฺม วิหารปจฺฉายายํ ปญฺญตฺเต อาสเน นิสีทิ.\csromanspacer{} อถ โข เต โกสลกา จ พฺราหฺมณทูตา มาคธกา จ พฺราหฺมณทูตา เยน ภควา เตนุปสงฺกมิํสุ, อุปสงฺกมิตฺวา ภควตา สทฺธิํ สมฺโมทิํสุ, สมฺโมทนียํ กถํ สารณียํ วีติสาเรตฺวา เอกมนฺตํ นิสีทิํสุ.\csromanspacer{} โอฏฺฐทฺโธปิ ลิจฺฉวี มหติยา ลิจฺฉวีปริสาย สทฺธิํ เยน ภควา เตนุปสงฺกมิ, อุปสงฺกมิตฺวา ภควนฺตํ อภิวาเทตฺวา เอกมนฺตํ นิสีทิ.}

\proseitem{364}{เอกมนฺตํ นิสินฺโน โข โอฏฺฐทฺโธ ลิจฺฉวี ภควนฺตํ เอตทโวจ “ปุริมานิ ภนฺเต ทิวสานิ ปุริมตรานิ สุนกฺขตฺโต ลิจฺฉวิปุตฺโต เยนาหํ เตนุปสงฺกมิ, อุปสงฺกมิตฺวา มํ เอตทโวจ ‘ยทคฺเค อหํ มหาลิ ภควนฺตํ อุปนิสฺสาย วิหรามิ, น จิรํ ตีณิ วสฺสานิ, ทิพฺพานิ หิ โข รูปานิ ปสฺสามิ ปิยรูปานิ กามูปสํหิตานิ รชนียานิ, โน จ โข ทิพฺพานิ สทฺทานิ สุณามิ ปิยรูปานิ กามูปสํหิตานิ รชนียานี’ติ, สนฺตาเนว นุ โข ภนฺเต สุนกฺขตฺโต ลิจฺฉวิปุตฺโต ทิพฺพานิ สทฺทานิ นาสฺโสสิ ปิยรูปานิ กามูปสํหิตานิ รชนียานิ, อุทาหุ อสนฺตานี”ติ.}

\csromanmatikaanchor{mtk.96}
\csromantocmarkref{section}{เอกํสภาวิตสมาธิ}{mtk.96}
\bookmark[dest={mtk.96},level=1]{เอกํสภาวิตสมาธิ}
\csromanheader{1}{\textbf{เอกํสภาวิตสมาธิ}}
\proseitem{365}{สนฺตาเนว โข มหาลิ สุนกฺขตฺโต ลิจฺฉวิปุตฺโต ทิพฺพานิ สทฺทานิ นาสฺโสสิ ปิยรูปานิ กามูปสํหิตานิ รชนียานิ, โน อสนฺตานีติ.\csromanspacer{} โก นุ โข ภนฺเต เหตุ, โก ปจฺจโย, เยน สนฺตาเนว สุนกฺขตฺโต ลิจฺฉวิปุตฺโต ทิพฺพานิ สทฺทานิ นาสฺโสสิ ปิยรูปานิ กามูปสํหิตานิ รชนียานิ, โน อสนฺตานีติ.}

\proseitem{366}{อิธ มหาลิ ภิกฺขุโน ปุรตฺถิมาย ทิสาย เอกํสภาวิโต สมาธิ โหติ ทิพฺพานํ รูปานํ ทสฺสนาย ปิยรูปานํ กามูปสํหิตานํ รชนียานํ, โน จ โข ทิพฺพานํ สทฺทานํ สวนาย ปิยรูปานํ กามูปสํหิตานํ รชนียานํ.\csromanspacer{} โส ปุรตฺถิมาย ทิสาย เอกํสภาวิเต สมาธิมฺหิ ทิพฺพานํ รูปานํ ทสฺสนาย ปิยรูปานํ กามูปสํหิตานํ รชนียานํ, โน จ โข ทิพฺพานํ สทฺทานํ สวนาย ปิยรูปานํ กามูปสํหิตานํ รชนียานํ.\csromanspacer{} ปุรตฺถิมาย ทิสาย ทิพฺพานิ รูปานิ ปสฺสติ ปิยรูปานิ กามูปสํหิตานิ รชนียานิ, โน จ โข ทิพฺพานิ สทฺทานิ สุณาติ ปิยรูปานิ กามูปสํหิตานิ รชนียานิ.\csromanspacer{} ตํ กิสฺส เหตุ, เอวํ เหตํ มหาลิ โหติ \csromanfolio{146}ภิกฺขุโน ปุรตฺถิมาย ทิสาย เอกํสภาวิเต สมาธิมฺหิ ทิพฺพานํ รูปานํ ทสฺสนาย ปิยรูปานํ กามูปสํหิตานํ รชนียานํ, โน จ โข ทิพฺพานํ สทฺทานํ สวนาย ปิยรูปานํ กามูปสํหิตานํ รชนียานํ.}

\proseitem{367}{ปุน จปรํ มหาลิ ภิกฺขุโน ทกฺขิณาย ทิสาย ฯเปฯ.\csromanspacer{} ปจฺฉิมาย ทิสาย.\csromanspacer{} อุตฺตราย ทิสาย.\csromanspacer{} อุทฺธมโธ ติริยํ เอกํสภาวิโต สมาธิ โหติ ทิพฺพานํ รูปานํ ทสฺสนาย ปิยรูปานํ กามูปสํหิตานํ รชนียานํ, โน จ โข ทิพฺพานํ สทฺทานํ สวนาย ปิยรูปานํ กามูปสํหิตานํ รชนียานํ.\csromanspacer{} โส อุทฺธมโธ ติริยํ เอกํสภาวิเต สมาธิมฺหิ ทิพฺพานํ รูปานํ ทสฺสนาย ปิยรูปานํ กามูปสํหิตานํ รชนียานํ, โน จ โข ทิพฺพานํ สทฺทานํ สวนาย ปิยรูปานํ กามูปสํหิตานํ รชนียานํ.\csromanspacer{} อุทฺธมโธ ติริยํ ทิพฺพานิ รูปานิ ปสฺสติ ปิยรูปานิ กามูปสํหิตานิ รชนียานิ, โน จ โข ทิพฺพานิ สทฺทานิ สุณาติ ปิยรูปานิ กามูปสํหิตานิ รชนียานิ.\csromanspacer{} ตํ กิสฺส เหตุ, เอวํ เหตํ มหาลิ โหติ ภิกฺขุโน อุทฺธมโธ ติริยํ เอกํสภาวิเต สมาธิมฺหิ ทิพฺพานํ รูปานํ ทสฺสนาย ปิยรูปานํ กามูปสํหิตานํ รชนียานํ, โน จ โข ทิพฺพานํ สทฺทานํ สวนาย ปิยรูปานํ กามูปสํหิตานํ รชนียานํ.}

\proseitem{368}{อิธ มหาลิ ภิกฺขุโน ปุรตฺถิมาย ทิสาย เอกํสภาวิโต สมาธิ โหติ ทิพฺพานํ สทฺทานํ สวนาย ปิยรูปานํ กามูปสํหิตานํ รชนียานํ, โน จ โข ทิพฺพานํ รูปานํ ทสฺสนาย ปิยรูปานํ กามูปสํหิตานํ รชนียานํ.\csromanspacer{} โส ปุรตฺถิมาย ทิสาย เอกํสภาวิเต สมาธิมฺหิ ทิพฺพานํ สวนาย ปิยรูปานํ กามูปสํหิตานํ รชนียานํ, โน จ โข ทิพฺพานํ รูปานํ ทสฺสนาย ปิยรูปานํ กามูปสํหิตานํ รชนียานํ.\csromanspacer{} ปุรตฺถิมาย ทิสาย ทิพฺพานิ สทฺทานิ สุณาติ ปิยรูปานิ กามูปสํหิตานิ รชนียานิ, โน จ โข ทิพฺพานิ รูปานิ ปสฺสติ ปิยรูปานิ กามูปสํหิตานิ รชนียานิ.\csromanspacer{} ตํ กิสฺส เหตุ, เอวํ เหตํ มหาลิ โหติ ภิกฺขุโน ปุรตฺถิมาย ทิสาย เอกํสภาวิเต สมาธิมฺหิ ทิพฺพานํ สทฺทานํ สวนาย ปิยรูปานํ กามูปสํหิตานํ รชนียานํ, โน จ โข ทิพฺพานํ รูปานํ ทสฺสนาย ปิยรูปานํ กามูปสํหิตานํ รชนียานํ.}

\proseitem{369}{ปุน จปรํ มหาลิ ภิกฺขุโน ทกฺขิณาย ทิสาย ฯเปฯ.\csromanspacer{} ปจฺฉิมาย ทิสาย.\csromanspacer{} อุตฺตราย ทิสาย.\csromanspacer{} อุทฺธมโธ ติริยํ เอกํสภาวิโต}

\vspace{\tipitakaparskip}
\csromancenter{มหาลิสุตฺต สมาธิ โหติ ทิพฺพานํ สทฺทานํ สวนาย ปิยรูปานํ กามูปสํหิตานํ รชนียานํ, โน จ โข ทิพฺพานํ รูปานํ ทสฺสนาย ปิยรูปานํ กามูปสํหิตานํ รชนียานํ.\csromanspacer{} โส อุทฺธมโธ ติริยํ เอกํสภาวิเต สมาธิมฺหิ ทิพฺพานํ สทฺทานํ สวนาย ปิยรูปานํ กามูปสํหิตานํ รชนียานํ, โน จ โข ทิพฺพานํ รูปานํ ทสฺสนาย ปิยรูปานํ กามูปสํหิตานํ รชนียานํ.\csromanspacer{} อุทฺธมโธ ติริยํ ทิพฺพานิ สทฺทานิ สุณาติ ปิยรูปานิ กามูปสํหิตานิ รชนียานิ, โน จ โข ทิพฺพานิ รูปานิ ปสฺสติ ปิยรูปานิ กามูปสํหิตานิ รชนียานิ.\csromanspacer{} ตํ กิสฺส เหตุ, เอวํ เหตํ มหาลิ โหติ ภิกฺขุโน อุทฺธมโธ ติริยํ เอกํสภาวิเต สมาธิมฺหิ ทิพฺพานํ สทฺทานํ สวนาย ปิยรูปานํ กามูปสํหิตานํ รชนียานํ, โน จ โข ทิพฺพานํ รูปานํ ทสฺสนาย ปิยรูปานํ กามูปสํหิตานํ รชนียานํ.}
\proseitem{370}{\csromanfolio{147}อิธ มหาลิ ภิกฺขุโน ปุรตฺถิมาย ทิสาย อุภยํสภาวิโต สมาธิ โหติ ทิพฺพานญฺจ รูปานํ ทสฺสนาย ปิยรูปานํ กามูปสํหิตานํ รชนียานํ, ทิพฺพานญฺจ สทฺทานํ สวนาย ปิยรูปานํ กามูปสํหิตานํ รชนียานํ.\csromanspacer{} โส ปุรตฺถิมาย ทิสาย อุภยํสภาวิเต สมาธิมฺหิ ทิพฺพานญฺจ รูปานํ ทสฺสนาย ปิยรูปานํ กามูปสํหิตานํ รชนียานํ, ทิพฺพานญฺจ สทฺทานํ สวนาย ปิยรูปานํ กามูปสํหิตานํ รชนียานํ.\csromanspacer{} ปุรตฺถิมาย ทิสาย ทิพฺพานิ จ รูปานิ ปสฺสติ ปิยรูปานิ กามูปสํหิตานิ รชนียานิ, ทิพฺพานิ จ สทฺทานิ สุณาติ ปิยรูปานิ กามูปสํหิตานิ รชนียานิ.\csromanspacer{} ตํ กิสฺส เหตุ, เอวํ เหตํ มหาลิ โหติ ภิกฺขุโน ปุรตฺถิมาย ทิสาย อุภยํสภาวิเต สมาธิมฺหิ ทิพฺพานญฺจ รูปานํ ทสฺสนาย ปิยรูปานํ กามูปสํหิตานํ รชนียานํ ทิพฺพานญฺจ สทฺทานํ สวนาย ปิยรูปานํ กามูปสํหิตานํ รชนียานํ.}

\proseitem{371}{ปุน จปรํ มหาลิ ภิกฺขุโน ทกฺขิณาย ทิสาย ฯเปฯ.\csromanspacer{} ปจฺฉิมาย ทิสาย.\csromanspacer{} อุตฺตราย ทิสาย.\csromanspacer{} อุทฺธมโธ ติริยํ อุภยํสภาวิโต สมาธิ โหติ ทิพฺพานญฺจ รูปานํ ทสฺสนาย ปิยรูปานํ กามูปสํหิตานํ รชนียานํ, ทิพฺพานญฺจ สทฺทานํ สวนาย ปิยรูปานํ กามูปสํหิตานํ รชนียานํ.\csromanspacer{} โส อุทฺธมโธ ติริยํ อุภยํสภาวิเต สมาธิมฺหิ ทิพฺพานญฺจ รูปานํ ทสฺสนาย ปิยรูปานํ กามูปสํหิตานํ รชนียานํ ทิพฺพานญฺจ สทฺทานํ สวนาย ปิยรูปานํ กามูปสํหิตานํ รชนียานํ.\csromanspacer{} อุทฺธมโธ ติริยํ ทิพฺพานิ จ รูปานิ ปสฺสติ \csromanfolio{148}ปิยรูปานิ กามูปสํหิตานิ รชนียานิ, ทิพฺพานิ จ สทฺทานิ สุณาติ ปิยรูปานิ กามูปสํหิตานิ รชนียานิ.\csromanspacer{} ตํ กิสฺส เหตุ, เอวํ เหตํ มหาลิ โหติ ภิกฺขุโน อุทฺธมโธ ติริยํ อุภยํสภาวิเต สมาธิมฺหิ ทิพฺพานญฺจ รูปานํ ทสฺสนาย ปิยรูปานํ กามูปสํหิตานํ รชนียานํ, ทิพฺพานญฺจ สทฺทานํ สวนาย ปิยรูปานํ กามูปสํหิตานํ รชนียานํ.\csromanspacer{} อยํ โข มหาลิ เหตุ อยํ ปจฺจโย.\csromanspacer{} เยน สนฺตาเนว สุนกฺขตฺโต ลิจฺฉวิปุตฺโต ทิพฺพานิ สทฺทานิ นาสฺโสสิ ปิยรูปานิ กามูปสํหิตานิ รชนียานิ, โน อสนฺตานีติ.}

\proseitem{372}{เอตาสํ นูน ภนฺเต สมาธิภาวนานํ สจฺฉิกิริยาเหตุ ภิกฺขู ภควติ พฺราหฺมจริยํ จรนฺตีติ.\csromanspacer{} น โข มหาลิ เอตาสํ สมาธิภาวนานํ สจฺฉิกิริยาเหตุ ภิกฺขู มยิ พฺรหฺมจริยํ จรนฺติ, อตฺถิ โข มหาลิ อญฺเญว ธมฺมา อุตฺตริตรา จ ปณีตตรา จ, เยสํ สจฺฉิกิริยาเหตุ ภิกฺขู มยิ พฺรหฺมจริยํ จรนฺตีติ.}

\csromanmatikaanchor{mtk.97}
\csromantocmarkref{section}{จตุอริยผล}{mtk.97}
\bookmark[dest={mtk.97},level=1]{จตุอริยผล}
\csromanheader{1}{\textbf{จตุอริยผล}}
\proseitem{373}{กตเม ปน เต ภนฺเต ธมฺมา อุตฺตริตรา จ ปณีตตรา จ, เยสํ สจฺฉิกิริยาเหตุ ภิกฺขู ภควติ พฺรหฺมจริยํ จรนฺตีติ.\csromanspacer{} อิธ มหาลิ ภิกฺขุ ติณฺณํ สญฺโญชนานํ ปริกฺขยา โสตาปนฺโน โหติ อวินิปาตธมฺโม นิยโต สมฺโพธิปรายโณ.\csromanspacer{} อยมฺปิ โข มหาลิ ธมฺโม อุตฺตริตโร จ ปณีตตโร จ ยสฺส สจฺฉิกิริยาเหตุ ภิกฺขู มยิ พฺรหฺมจริยํ จรนฺติ.}

\prose{ปุน จปรํ มหาลิ ภิกฺขุ ติณฺณํ สํโยชนานํ ปริกฺขยา ราคโทสโมหานํ ตนุตฺตา สกทาคามี โหติ, สกิเทว\footnote{สกิํ เทว (ก)} อิมํ โลกํ อาคนฺตฺวา ทุกฺขสฺสนฺตํ กโรติ.\csromanspacer{} อยมฺปิ โข มหาลิ ธมฺโม อุตฺตริตโร จ ปณีตตโร จ ยสฺส สจฺฉิกิริยาเหตุ ภิกฺขู มยิ พฺรหฺมจริยํ จรนฺติ.}

\prose{ปุน จปรํ มหาลิ ภิกฺขุ ปญฺจนฺนํ โอรมฺภาคิยานํ สํโยชนานํ ปริกฺขยา โอปปาติโก โหติ ตตฺถ ปรินิพฺพายี อนาวตฺติธมฺโม ตสฺมา โลกา.\csromanspacer{} อยมฺปิ โข มหาลิ ธมฺโม อุตฺตริตโร จ ปณีตตโร จ ยสฺส สจฺฉิกิริยาเหตุ ภิกฺขู มยิ พฺรหฺมจริยํ จรนฺติ.}

\csromanheader{2}{6. มหาลิสุตฺต}
\prose{\csromanfolio{149}ปุน จปรํ มหาลิ ภิกฺขุ อาสวานํ ขยา อนาสวํ เจโตวิมุตฺติํ ปญฺญาวิมุตฺติํ ทิฏฺเฐว ธมฺเม สยํ อภิญฺญา สจฺฉิกตฺวา อุปสมฺปชฺช วิหรติ.\csromanspacer{} อยมฺปิ โข มหาลิ ธมฺโม อุตฺตริตโร จ ปณีตตโร จ ยสฺส สจฺฉิกิริยาเหตุ ภิกฺขู มยิ พฺรหฺมจริยํ จรนฺติ.\csromanspacer{} อิเม โข เต มหาลิ ธมฺมา อุตฺตริตรา จ ปณีตตรา จ เยสํ สจฺฉิกิริยาเหตุ ภิกฺขู มยิ พฺรหฺมจริยํ จรนฺตีติ.}

\csromanmatikaanchor{mtk.98}
\csromantocmarkref{section}{อริยอฏฺฐงฺคิกมคฺค}{mtk.98}
\bookmark[dest={mtk.98},level=1]{อริยอฏฺฐงฺคิกมคฺค}
\csromanheader{1}{\textbf{อริยอฏฺฐงฺคิกมคฺค}}
\proseitem{374}{อตฺถิ ปน ภนฺเต มคฺโค อตฺถิ ปฏิปทา เอเตสํ ธมฺมานํ สจฺฉิกิริยายาติ.\csromanspacer{} อตฺถิ โข มหาลิ มคฺโค อตฺถิ ปฏิปทา เอเตสํ ธมฺมานํ สจฺฉิกิริยายาติ.}

\proseitem{375}{กตโม ปน ภนฺเต มคฺโค กตมา ปฏิปทา เอเตสํ ธมฺมานํ สจฺฉิกิริยายาติ.\csromanspacer{} อยเมว อริโย อฏฺฐงฺคิโก มคฺโค.\csromanspacer{} เสยฺยถิทํ, สมฺมาทิฏฺฐิ สมฺมาสงฺกปฺโป สมฺมาวาจา สมฺมากมฺมนฺโต สมฺมา-อาชีโว สมฺมาวายาโม สมฺมาสติ สมฺมาสมาธิ.\csromanspacer{} อยํ โข มหาลิ มคฺโค อยํ ปฏิปทา เอเตสํ ธมฺมานํ สจฺฉิกิริยาย.}

\csromanmatikaanchor{mtk.99}
\csromantocmarkref{section}{เทฺวปพฺพชิตวตฺถุ}{mtk.99}
\bookmark[dest={mtk.99},level=1]{เทฺวปพฺพชิตวตฺถุ}
\csromanheader{1}{\textbf{เทฺวปพฺพชิตวตฺถุ}}
\proseitem{376}{เอกมิทาหํ มหาลิ สมยํ โกสมฺพิยํ วิหรามิ โฆสิตาราเม.\csromanspacer{} อถ โข เทฺว ปพฺพชิตา มุณฺฑิโย จ ปริพฺพาชโก ชาลิโย จ ทารุปตฺติกนฺเตวาสี เยนาหํ เตนุปสงฺกมิํสุ, อุปสงฺกมิตฺวา มยา สทฺธิํ สมฺโมทิํสุ, สมฺโมทนียํ กถํ สารณียํ วีติสาเรตฺวา เอกมนฺตํ อฏฺฐํสุ, เอกมนฺตํ ฐิตา โข เต เทฺว ปพฺพชิตา มํ เอตทโวจุํ “กิํ นุ โข อาวุโส โคตม ตํ ชีวํ ตํ สรีรํ, อุทาหุ อญฺญํ ชีวํ อญฺญํ สรีรนฺ”ติ.}

\proseitem{377}{เตน หาวุโส สุณาถ สาธุกํ มนสิ กโรถ ภาสิสฺสามีติ. “เอวมาวุโส”ติ โข เต เทฺว ปพฺพชิตา มม ปจฺจสฺโสสุํ.\csromanspacer{} อหํ เอตทโวจํ.\csromanspacer{} อิธาวุโส ตถาคโต โลเก อุปฺปชฺชติ อรหํ สมฺมาสมฺพุทฺโธ ฯเปฯ. (ยถา สามญฺญผเล, เอวํ วิตฺถาเรตพฺพํ.) เอวํ โข อาวุโส ภิกฺขุ สีลสมฺปนฺโน โหติ.}

\prose{ฯเปฯ ปฐมํ ฌานํ อุปสมฺปชฺช วิหรติ.\csromanspacer{} โย โข อาวุโส ภิกฺขุ เอวํ ชานาติ เอวํ ปสฺสติ, กลฺลํ นุ โข ตสฺเสตํ วจนาย “ตํ ชีวํ ตํ สรีรนฺ”ติ \csromanfolio{150}วา “อญฺญํ ชีวํ อญฺญํ สรีรนฺ”ติ วาติ.\csromanspacer{} โย โส อาวุโส ภิกฺขุ เอวํ ชานาติ เอวํ ปสฺสติ, กลฺลํ ตสฺเสตํ วจนาย “ตํ ชีวํ ตํ สรีรนฺ”ติ วา “อญฺญํ ชีวํ อญฺญํ สรีรนฺ”ติ วาติ.\csromanspacer{} อหํ โข ปเนตํ อาวุโส เอวํ ชานามิ เอวํ ปสฺสามิ.\csromanspacer{} อถ จ ปนาหํ น วทามิ “ตํ ชีวํ ตํ สรีรนฺ”ติ วา “อญฺญํ ชีวํ อญฺญํ สรีรนฺ”ติ วา ฯเปฯ ทุติยํ ฌานํ.\csromanspacer{} ตติยํ ฌานํ.\csromanspacer{} จตุตฺถํ ฌานํ อุปสมฺปชฺช วิหรติ.\csromanspacer{} โย โข อาวุโส ภิกฺขุ เอวํ ชานาติ เอวํ ปสฺสติ, กลฺลํ นุ โข ตสฺเสตํ วจนาย “ตํ ชีวํ ตํ สรีรนฺ”ติ วา “อญฺญํ ชีวํ อญฺญํ สรีรนฺ”ติ วาติ.\csromanspacer{} โย โส อาวุโส ภิกฺขุ เอวํ ชานาติ เอวํ ปสฺสติ, กลฺลํ ตสฺเสตํ วจนาย “ตํ ชีวํ ตํ สรีรนฺ”ติ วา “อญฺญํ ชีวํ อญฺญํ สรีรนฺ”ติ วาติ.\csromanspacer{} อหํ โข ปเนตํ อาวุโส เอวํ ชานามิ เอวํ ปสฺสามิ.\csromanspacer{} อถ จ ปนาหํ น วทามิ “ตํ ชีวํ ตํ สรีรนฺ”ติ วา “อญฺญํ ชีวํ อญฺญํ สรีรนฺ”ติ วา ฯเปฯ ญาณทสฺสนาย จิตฺตํ อภินีหรติ อภินินฺนาเมติ.\csromanspacer{} โย โข อาวุโส ภิกฺขุ เอวํ ชานาติ เอวํ ปสฺสติ, กลฺลํ นุ โข ตสฺเสตํ วจนาย “ตํ ชีวํ ตํ สรีรนฺ”ติ วา “อญฺญํ ชีวํ อญฺญํ สรีรนฺ”ติ วาติ.\csromanspacer{} โย โส อาวุโส ภิกฺขุ เอวํ ชานาติ เอวํ ปสฺสติ, กลฺลํ\footnote{น กลฺลํ (สี, สฺยา, กํ, ก)} ตสฺเสตํ วจนาย “ตํ ชีวํ ตํ สรีรนฺ”ติ วา “อญฺญํ ชีวํ อญฺญํ สรีรนฺ”ติ วาติ.\csromanspacer{} อหํ โข ปเนตํ อาวุโส เอวํ ชานามิ เอวํ ปสฺสามิ.\csromanspacer{} อถ จ ปนาหํ น วทามิ “ตํ ชีวํ ตํ สรีรนฺ”ติ วา “อญฺญํ ชีวํ อญฺญํ สรีรนฺ”ติ วา.}

\prosewithcloser{ฯเปฯ นาปรํ อิตฺถตฺตายาติ ปชานาติ.\csromanspacer{} โย โข อาวุโส ภิกฺขุ เอวํ ชานาติ เอวํ ปสฺสติ, กลฺลํ นุ โข ตสฺเสตํ วจนาย “ตํ ชีวํ ตํ สรีรนฺ”ติ วา “อญฺญํ ชีวํ อญฺญํ สรีรนฺ”ติ วาติ.\csromanspacer{} โย โส อาวุโส ภิกฺขุ เอวํ ชานาติ เอวํ ปสฺสติ, น กลฺลํ ตสฺเสตํ วจนาย “ตํ ชีวํ ตํ สรีรนฺ”ติ วา “อญฺญํ ชีวํ อญฺญํ สรีรนฺ”ติ วาติ.\csromanspacer{} อหํ โข ปเนตํ อาวุโส เอวํ ชานามิ เอวํ ปสฺสามิ.\csromanspacer{} อถ จ ปนาหํ น วทามิ “ตํ ชีวํ ตํ สรีรนฺ”ติ วา “อญฺญํ ชีวํ อญฺญํ สรีรนฺ”ติ วาติ.\csromanspacer{} อิทมโวจ ภควา.\csromanspacer{} อตฺตมโน โอฏฺฐทฺโธ ลิจฺฉวี ภควโต ภาสิตํ อภินนฺทีติ.}{\textbf{มหาลิสุตฺตํ นิฏฺฐิตํ ฉฏฺฐํ.}}
\csromanmatikaanchor{mtk.100}
\csromantocmarknopage{chapter}{7. ชาลิยสุตฺต}{mtk.100}
\bookmark[dest={mtk.100},level=0]{7. ชาลิยสุตฺต}
\setcsromanheadcha{7. ชาลิยสุตฺต}
\setcsromanheadhclear{}
\chapterheadpageread{7. \textbf{ชาลิยสุตฺต}}
\csromanmatikaanchor{mtk.101}
\csromantocmarkref{section}{เทฺวปพฺพชิตวตฺถุ}{mtk.101}
\bookmark[dest={mtk.101},level=1]{เทฺวปพฺพชิตวตฺถุ}
\csromanheader{1}{\textbf{เทฺวปพฺพชิตวตฺถุ}}
\proseitem{378}{\csromanfolio{151}เอวํ เม สุตํ\csromandash{}เอกํ สมยํ ภควา โกสมฺพิยํ วิหรติ โฆสิตาราเม.\csromanspacer{} เตน โข ปน สมเยน เทฺว ปพฺพชิตา มุณฺฑิโย จ ปริพฺพาชโก ชาลิโย จ ทารุปตฺติกนฺเตวาสี เยน ภควา เตนุปสงฺกมิํสุ, อุปสงฺกมิตฺวา ภควตา สทฺธิํ สมฺโมทิํสุ, สมฺโมทนียํ กถํ สารณียํ วีติสาเรตฺวา เอกมนฺตํ อฏฺฐํสุ, เอกมนฺตํ ฐิตา โข เต เทฺว ปพฺพชิตา ภควนฺตํ เอตทโวจุํ “กิํ นุ โข อาวุโส โคตม ตํ ชีวํ ตํ สรีรํ, อุทาหุ อญฺญํ ชีวํ อญฺญํ สรีรนฺ”ติ.}

\proseitem{379}{เตน หาวุโส สุณาถ สาธุกํ มนสิ กโรถ ภาสิสฺสามีติ. “เอวมาวุโส”ติ โข เต เทฺว ปพฺพชิตา ภควโต ปจฺจสฺโสสุํ.\csromanspacer{} ภควา เอตทโวจ.\csromanspacer{} อิธาวุโส ตถาคโต โลเก อุปฺปชฺชติ อรหํ สมฺมาสมฺพุทฺโธ ฯเปฯ. (ยถา สามญฺญผเล, เอวํ วิตฺถาเรตพฺพํ.) เอวํ โข อาวุโส ภิกฺขุ สีลสมฺปนฺโน โหติ.}

\prose{ฯเปฯ ปฐมํ ฌานํ อุปสมฺปชฺช วิหรติ.\csromanspacer{} โย โข อาวุโส ภิกฺขุ เอวํ ชานาติ เอวํ ปสฺสติ, กลฺลํ นุ โข ตสฺเสตํ วจนาย “ตํ ชีวํ ตํ สรีรนฺ”ติ วา “อญฺญํ ชีวํ อญฺญํ สรีรนฺ”ติ วาติ.\csromanspacer{} โย โส อาวุโส ภิกฺขุ เอวํ ชานาติ เอวํ ปสฺสติ, กลฺลํ ตสฺเสตํ วจนาย “ตํ ชีวํ ตํ สรีรนฺ”ติ วา “อญฺญํ ชีวํ อญฺญํ สรีรนฺ”ติ วาติ.\csromanspacer{} อหํ โข ปเนตํ อาวุโส เอวํ ชานามิ เอวํ ปสฺสามิ.\csromanspacer{} อถ จ ปนาหํ น วทามิ “ตํ ชีวํ ตํ สรีรนฺ”ติ วา “อญฺญํ ชีวํ อญฺญํ สรีรนฺ”ติ วา ฯเปฯ ทุติยํ ฌานํ.\csromanspacer{} ตติยํ ฌานํ.\csromanspacer{} จตุตฺถํ ฌานํ อุปสมฺปชฺช วิหรติ.\csromanspacer{} โย โข อาวุโส ภิกฺขุ เอวํ ชานาติ เอวํ ปสฺสติ, กลฺลํ นุ โข ตสฺเสตํ วจนาย “ตํ ชีวํ ตํ สรีรนฺ”ติ วา “อญฺญํ ชีวํ อญฺญํ สรีรนฺ”ติ วาติ.\csromanspacer{} โย โส อาวุโส ภิกฺขุ เอวํ ชานาติ เอวํ ปสฺสติ, กลฺลํ ตสฺเสตํ วจนาย “ตํ ชีวํ ตํ สรีรนฺ”ติ วา “อญฺญํ ชีวํ อญฺญํ สรีรนฺ”ติ วาติ.\csromanspacer{} อหํ โข ปเนตํ อาวุโส เอวํ ชานามิ เอวํ ปสฺสามิ.\csromanspacer{} อถ จ ปนาหํ น วทามิ “ตํ ชีวํ ตํ สรีรนฺ”ติ วา “อญฺญํ ชีวํ อญฺญํ สรีรนฺ”ติ วา ฯเปฯ ญาณทสฺสนาย \csromanfolio{152}จิตฺตํ อภินีหรติ อภินินฺนาเมติ.\csromanspacer{} โย โข อาวุโส ภิกฺขุ เอวํ ชานาติ เอวํ ปสฺสติ, กลฺลํ นุ โข ตสฺเสตํ วจนาย “ตํ ชีวํ ตํ สรีรนฺ”ติ วา “อญฺญํ ชีวํ อญฺญํ สรีรนฺ”ติ วาติ.\csromanspacer{} โย โส อาวุโส ภิกฺขุ เอวํ ชานาติ เอวํ ปสฺสติ, กลฺลํ ตสฺเสตํ วจนาย “ตํ ชีวํ ตํ สรีรนฺ”ติ วา “อญฺญํ ชีวํ อญฺญํ สรีรนฺ”ติ วาติ.\csromanspacer{} อหํ โข ปเนตํ อาวุโส เอวํ ชานามิ เอวํ ปสฺสามิ.\csromanspacer{} อถ จ ปนาหํ น วทามิ “ตํ ชีวํ ตํ สรีรนฺ”ติ วา “อญฺญํ ชีวํ อญฺญํ สรีรนฺ”ติ วา.}

\proseitemwithcloser{380}{ฯเปฯ นาปรํ อิตฺถตฺตายาติ ปชานาติ.\csromanspacer{} โย โข อาวุโส ภิกฺขุ เอวํ ชานาติ เอวํ ปสฺสติ, กลฺลํ นุ โข ตสฺเสตํ วจนาย “ตํ ชีวํ ตํ สรีรนฺ”ติ วา “อญฺญํ ชีวํ อญฺญํ สรีรนฺ”ติ วาติ.\csromanspacer{} โย โส อาวุโส ภิกฺขุ เอวํ ชานาติ เอวํ ปสฺสติ, น กลฺลํ ตสฺเสตํ วจนาย “ตํ ชีวํ ตํ สรีรนฺ”ติ วา “อญฺญํ ชีวํ อญฺญํ สรีรนฺ”ติ วาติ.\csromanspacer{} อหํ โข ปเนตํ อาวุโส เอวํ ชานามิ เอวํ ปสฺสามิ.\csromanspacer{} อถ จ ปนาหํ น วทามิ “ตํ ชีวํ ตํ สรีรนฺ”ติ วา “อญฺญํ ชีวํ อญฺญํ สรีรนฺ”ติ วาติ.\csromanspacer{} อิทมโวจ ภควา.\csromanspacer{} อตฺตมนา เต เทฺว ปพฺพชิตา ภควโต ภาสิตํ อภินนฺทุนฺติ.}{\textbf{ชาลิยสุตฺตํ นิฏฺฐิตํ สตฺตมํ.}}
\csromanmatikaanchor{mtk.102}
\csromantocmarknopage{chapter}{8. มหาสีหนาทสุตฺต}{mtk.102}
\bookmark[dest={mtk.102},level=0]{8. มหาสีหนาทสุตฺต}
\setcsromanheadcha{8. มหาสีหนาทสุตฺต}
\setcsromanheadhclear{}
\chapterheadpageread{8. \textbf{มหาสีหนาทสุตฺต}}
\csromanmatikaanchor{mtk.103}
\csromantocmarkref{section}{อเจลกสฺสปวตฺถุ}{mtk.103}
\bookmark[dest={mtk.103},level=1]{อเจลกสฺสปวตฺถุ}
\csromanheader{1}{\textbf{อเจลกสฺสปวตฺถุ}}
\proseitem{381}{\csromanfolio{153}เอวํ เม สุตํ\csromandash{}เอกํ สมยํ ภควา อุรุญฺญายํ\footnote{อุชุญฺญายํ (สี, สฺยา, กํ, อิ)} วิหรติ กณฺณกตฺถเล มิคทาเย.\csromanspacer{} อถ โข อเจโล กสฺสโป เยน ภควา เตนุปสงฺกมิ, อุปสงฺกมิตฺวา ภควตา สทฺธิํ สมฺโมทิ, สมฺโมทนียํ กถํ สารณียํ วีติสาเรตฺวา เอกมนฺตํ อฏฺฐาสิ, เอกมนฺตํ ฐิโต โข อเจโล กสฺสโป ภควนฺตํ เอตทโวจ “สุตํ เมตํ โภ โคตม ‘สมโณ โคตโม สพฺพํ ตปํ ครหติ, สพฺพํ ตปสฺสิํ ลูขาชีวิํ เอกํเสน อุปกฺโกสติ อุปวทตี’ติ.\csromanspacer{} เย เต โภ โคตม เอวมาหํสุ ‘สมโณ โคตโม สพฺพํ ตปํ ครหติ, สพฺพํ ตปสฺสิํ ลูขาชีวิํ เอกํเสน อุปกฺโกสติ อุปวทตี’ติ, กจฺจิ เต โภโต โคตมสฺส วุตฺตวาทิโน, น จ ภวนฺตํ โคตมํ อภูเตน อพฺภาจิกฺขนฺติ, ธมฺมสฺส จานุธมฺมํ พฺยากโรนฺติ, น จ โกจิ สหธมฺมิโก วาทานุวาโท คารยฺหํ ฐานํ อาคจฺฉติ.\csromanspacer{} อนพฺภกฺขาตุกามา หิ มยํ ภวนฺตํ โคตมนฺ”ติ.}

\proseitem{382}{เย เต กสฺสป เอวมาหํสุ “สมโณ โคตโม สพฺพํ ตปํ ครหติ, สพฺพํ ตปสฺสิํ ลูขาชีวิํ เอกํเสน อุปกฺโกสติ อุปวทตี”ติ, น เมเต วุตฺตวาทิโน, อพฺภาจิกฺขนฺติ จ ปน มํ เต อสตา อภูเตน.\csromanspacer{} อิธาหํ กสฺสป เอกจฺจํ ตปสฺสิํ ลูขาชีวิํ ปสฺสามิ ทิพฺเพน จกฺขุนา วิสุทฺเธน อติกฺกนฺตมานุสเกน กายสฺส เภทา ปรํ มรณา อปายํ ทุคฺคติํ วินิปาตํ นิรยํ อุปปนฺนํ.\csromanspacer{} อิธ ปนาหํ กสฺสป เอกจฺจํ ตปสฺสิํ ลูขาชีวิํ ปสฺสามิ ทิพฺเพน จกฺขุนา วิสุทฺเธน อติกฺกนฺตมานุสเกน กายสฺส เภทา ปรํ มรณา สุคติํ สคฺคํ โลกํ อุปปนฺนํ.}

\proseitem{383}{อิธาหํ กสฺสป เอกจฺจํ ตปสฺสิํ อปฺปทุกฺขวิหาริํ ปสฺสามิ ทิพฺเพน จกฺขุนา วิสุทฺเธน อติกฺกนฺตมานุสเกน กายสฺส เภทา ปรํ มรณา อปายํ ทุคฺคติํ วินิปาตํ นิรยํ อุปปนฺนํ.\csromanspacer{} อิธ ปนาหํ กสฺสป เอกจฺจํ ตปสฺสิํ อปฺปทุกฺขวิหาริํ ปสฺสามิ ทิพฺเพน จกฺขุนา วิสุทฺเธน อติกฺกนฺตมานุสเกน กายสฺส เภทา ปรํ มรณา สุคติํ สคฺคํ โลกํ อุปปนฺนํ.\csromanspacer{} โยหํ กสฺสป อิเมสํ ตปสฺสีนํ เอวํ อาคติญฺจ คติญฺจ จุติญฺจ อุปปตฺติญฺจ ยถาภูตํ \csromanfolio{154}ปชานามิ.\csromanspacer{} โสหํ กิํ สพฺพํ ตปํ ครหิสฺสามิ, สพฺพํ วา ตปสฺสิํ ลูขาชีวิํ เอกํเสน อุปกฺโกสิสฺสามิ อุปวทิสฺสามิ.}

\proseitem{384}{สนฺติ กสฺสป เอเก สมณพฺราหฺมณา ปณฺฑิตา นิปุณา กตปรปฺปวาทา วาลเวธิรูปา, เต ภินฺทนฺตา มญฺเญ จรนฺติ ปญฺญาคเตน ทิฏฺฐิคตานิ.\csromanspacer{} เตหิปิ เม สทฺธิํ เอกจฺเจสุ ฐาเนสุ สเมติ, เอกจฺเจสุ ฐาเนสุ น สเมติ.\csromanspacer{} ยํ เต เอกจฺจํ วทนฺติ “สาธู”ติ, มยมฺปิ ตํ เอกจฺจํ วเทม “สาธู”ติ.\csromanspacer{} ยํ เต เอกจฺจํ วทนฺติ “น สาธู”ติ, มยมฺปิ ตํ เอกจฺจํ วเทม “น สาธู”ติ.\csromanspacer{} ยํ เต เอกจฺจํ วทนฺติ “สาธู”ติ, มยํ ตํ เอกจฺจํ วเทม “น สาธู”ติ.\csromanspacer{} ยํ เต เอกจฺจํ วทนฺติ “น สาธู”ติ, มยํ ตํ เอกจฺจํ วเทม “สาธู”ติ.}

\prose{ยํ มยํ เอกจฺจํ วเทม “สาธู”ติ, ปเรปิ ตํ เอกจฺจํ วทนฺติ “สาธู”ติ.\csromanspacer{} ยํ มยํ เอกจฺจํ วเทม “น สาธู”ติ, ปเรปิ ตํ เอกจฺจํ วทนฺติ “น สาธู”ติ.\csromanspacer{} ยํ มยํ เอกจฺจํ วเทม “น สาธู”ติ, ปเร ตํ เอกจฺจํ วทนฺติ “สาธู”ติ.\csromanspacer{} ยํ มยํ เอกจฺจํ วเทม “สาธู”ติ, ปเร ตํ เอกจฺจํ วทนฺติ “น สาธู”ติ.}

\csromanmatikaanchor{mtk.104}
\csromantocmarkref{section}{สมนุยุญฺชาปนกถา}{mtk.104}
\bookmark[dest={mtk.104},level=1]{สมนุยุญฺชาปนกถา}
\csromanheader{1}{\textbf{สมนุยุญฺชาปนกถา}}
\proseitem{385}{ตฺยาหํ อุปสงฺกมิตฺวา เอวํ วทามิ\csromandash{}เยสุ โน อาวุโส ฐาเนสุ น สเมติ, ติฏฺฐนฺตุ ตานิ ฐานานิ.\csromanspacer{} เยสุ ฐาเนสุ สเมติ, ตตฺถ วิญฺญู สมนุยุญฺชนฺตํ สมนุคาหนฺตํ สมนุภาสนฺตํ สตฺถารา วา สตฺถารํ สํเฆน วา สํฆํ “เย อิเมสํ ภวตํ ธมฺมา อกุสลา อกุสลสงฺขาตา, สาวชฺชา สาวชฺชสงฺขาตา, อเสวิตพฺพา อเสวิตพฺพสงฺขาตา, น อลมริยา น อลมริยสงฺขาตา, กณฺหา กณฺหาสงฺขาตา.\csromanspacer{} โก อิเม ธมฺเม อนวเสสํ ปหาย วตฺตติ, สมโณ วา โคตโม, ปเร วา ปน โภนฺโต คณาจริยา”ติ.}

\proseitem{386}{ฐานํ โข ปเนตํ กสฺสป วิชฺชติ, ยํ วิญฺญู สมนุยุญฺชนฺตา สมนุคาหนฺตา สมนุภาสนฺตา เอวํ วเทยฺยุํ “เย อิเมสํ ภวตํ ธมฺมา อกุสลา อกุสลสงฺขาตา, สาวชฺชา สาวชฺชสงฺขาตา, อเสวิตพฺพา อเสวิตพฺพสงฺขาตา, น อลมริยา น อลมริยสงฺขาตา, กณฺหา กณฺหสงฺขาตา.\csromanspacer{} สมโณ โคตโม อิเม ธมฺเม อนวเสสํ ปหาย วตฺตติ, ยํ วา ปน โภนฺโต ปเร คณาจริยา”ติ.\csromanspacer{} อิติห กสฺสป}

\vspace{\tipitakaparskip}
\csromancenter{มหาสีหนาทสุตฺต วิญฺญู สมนุยุญฺชนฺตา สมนุคาหนฺตา สมนุภาสนฺตา อเมฺหว ตตฺถ เยภุยฺเยน ปสํเสยฺยุํ.}
\proseitem{387}{\csromanfolio{155}อปรมฺปิ โน กสฺสป วิญฺญู สมนุยุญฺชนฺตํ สมนุคาหนฺตํ สมนุภาสนฺตํ สตฺถารา วา สตฺถารํ สํเฆน วา สํฆํ “เย อิเมสํ ภวตํ ธมฺมา กุสลา กุสลสงฺขาตา, อนวชฺชา อนวชฺชสงฺขาตา, เสวิตพฺพา เสวิตพฺพสงฺขาตา, อลมริยา อลมริยสงฺขาตา, สุกฺกา สุกฺกสงฺขาตา.\csromanspacer{} โก อิเม ธมฺเม อนวเสสํ สมาทาย วตฺตติ, สมโณ วา โคตโม, ปเร วา ปน โภนฺโต คณาจริยา”ติ.}

\proseitem{388}{ฐานํ โข ปเนตํ กสฺสป วิชฺชติ, ยํ วิญฺญู สมนุยุญฺชนฺตา สมนุคาหนฺตา สมนุภาสนฺตา เอวํ วเทยฺยุํ “เย อิเมสํ ภวตํ ธมฺมา กุสลา กุสลสงฺขาตา, อนวชฺชา อนวชฺชสงฺขาตา, เสวิตพฺพา เสวิตพฺพสงฺขาตา, อลมริยา อลมริยสงฺขาตา, สุกฺกา สุกฺกสงฺขาตา.\csromanspacer{} สมโณ โคตโม อิเม ธมฺเม อนวเสสํ สมาทาย วตฺตติ, ยํ วา ปน โภนฺโต ปเร คณาจริยา”ติ.\csromanspacer{} อิติห กสฺสป วิญฺญู สมนุยุญฺชนฺตา สมนุคาหนฺตา สมนุภาสนฺตา อเมฺหว ตตฺถ เยภุยฺเยน ปสํเสยฺยุํ.}

\proseitem{389}{อปรมฺปิ โน กสฺสป วิญฺญู สมนุยุญฺชนฺตํ สมนุคาหนฺตํ สมนุภาสนฺตํ สตฺถารา วา สตฺถารํ สํเฆน วา สํฆํ “เย อิเมสํ ภวตํ ธมฺมา อกุสลา อกุสลสงฺขาตา, สาวชฺชา สาวชฺชสงฺขาตา, อเสวิตพฺพา อเสวิตพฺพสงฺขาตา, น อลมริยา น อลมริยสงฺขาตา, กณฺหา กณฺหสงฺขาตา.\csromanspacer{} โก อิเม ธมฺเม อนวเสสํ ปหาย วตฺตติ, โคตมสาวกสํโฆ วา, ปเร วา ปน โภนฺโต คณาจริยสาวกสํฆา”ติ.}

\proseitem{390}{ฐานํ โข ปเนตํ กสฺสป วิชฺชติ, ยํ วิญฺญู สมนุยุญฺชนฺตา สมนุคาหนฺตา สมนุภาสนฺตา เอวํ วเทยฺยุํ “เย อิเมสํ ภวตํ ธมฺมา อกุสลา อกุสลสงฺขาตา, สาวชฺชา สาวชฺชสงฺขาตา, อเสวิตพฺพา อเสวิตพฺพสงฺขาตา, น อลมริยา น อลมริยสงฺขาตา, กณฺหา กณฺหสงฺขาตา.\csromanspacer{} โคตมสาวกสํโฆ อิเม ธมฺเม อนวเสสํ ปหาย วตฺตติ, ยํ วา ปน โภนฺโต ปเร คณาจริยสาวกสํฆา”ติ.\csromanspacer{} อิติห กสฺสป วิญฺญู สมนุยุญฺชนฺตา สมนุคาหนฺตา สมนุภาสนฺตา อเมฺหว ตตฺถ เยภุยฺเยน ปสํเสยฺยุํ.}

\proseitem{391}{\csromanfolio{156}อปรมฺปิ โน กสฺสป วิญฺญู สมนุยุญฺชนฺตํ สมนุคาหนฺตํ สมนุภาสนฺตํ สตฺถารา วา สตฺถารํ สํเฆน วา สํฆํ “เย อิเมสํ ภวตํ ธมฺมา กุสลา กุสลสงฺขาตา, อนวชฺชา อนวชฺชสงฺขาตา, เสวิตพฺพา เสวิตพฺพสงฺขาตา, อลมริยา อลมริยสงฺขาตา, สุกฺกา สุกฺกสงฺขาตา.\csromanspacer{} โก อิเม ธมฺเม อนวเสสํ สมาทาย วตฺตติ, โคตมสาวกสํโฆ วา, ปเร วา ปน โภนฺโต คณาจริยสาวกสํฆา”ติ.}

\proseitem{392}{ฐานํ โข ปเนตํ กสฺสป วิชฺชติ, ยํ วิญฺญู สมนุยุญฺชนฺตา สมนุคาหนฺตา สมนุภาสนฺตา เอวํ วเทยฺยุํ “เย อิเมสํ ภวตํ ธมฺมา กุสลา กุสลสงฺขาตา, อนวชฺชา อนวชฺชสงฺขาตา, เสวิตพฺพา เสวิตพฺพสงฺขาตา, อลมริยา อลมริยสงฺขาตา, สุกฺกา สุกฺกสงฺขาตา.\csromanspacer{} โคตมสาวกสํโฆ อิเม ธมฺเม อนวเสสํ สมาทาย วตฺตติ, ยํ วา ปน โภนฺโต ปเร คณาจริยสาวกสํฆา”ติ.\csromanspacer{} อิติห กสฺสป วิญฺญู สมนุยุญฺชนฺตา สมนุคาหนฺตา สมนุภาสนฺตา อเมฺหว ตตฺถ เยภุยฺเยน ปสํเสยฺยุํ.}

\csromanmatikaanchor{mtk.105}
\csromantocmarkref{section}{อริยอฏฺฐงฺคิกมคฺค}{mtk.105}
\bookmark[dest={mtk.105},level=1]{อริยอฏฺฐงฺคิกมคฺค}
\csromanheader{1}{\textbf{อริยอฏฺฐงฺคิกมคฺค}}
\proseitem{393}{อตฺถิ กสฺสป มคฺโค อตฺถิ ปฏิปทา, ยถาปฏิปนฺโน สามํเยว ญสฺสติ สามํ ทกฺขติ\footnote{ทกฺขิติ (สี)} “สมโณว โคตโม กาลวาที ภูตวาที อตฺถวาที ธมฺมวาที วินยวาที”ติ.\csromanspacer{} กตโม จ กสฺสป มคฺโค, กตมา จ ปฏิปทา, ยถาปฏิปนฺโน สามญฺเญว ญสฺสติ สามํ ทกฺขติ “สมโณว โคตโม กาลวาที ภูตวาที อตฺถวาที ธมฺมวาที วินยวาที”ติ.\csromanspacer{} อยเมว อริโย อฏฺฐงฺคิโก มคฺโค.}

\prose{เสยฺยถิทํ, สมฺมาทิฏฺฐิ สมฺมาสงฺกปฺโป สมฺมาวาจา สมฺมากมฺมนฺโต สมฺมา-อาชีโว สมฺมาวายาโม สมฺมาสติ สมฺมาสมาธิ.\csromanspacer{} อยํ โข กสฺสป มคฺโค, อยํ ปฏิปทา, ยถาปฏิปนฺโน สามญฺเญว ญสฺสติ สามํ ทกฺขติ “สมโณว โคตโม กาลวาที ภูตวาที อตฺถวาที ธมฺมวาที วินยวาที”ติ.}

\csromanmatikaanchor{mtk.106}
\csromantocmarkref{section}{ตโปปกฺกมกถา}{mtk.106}
\bookmark[dest={mtk.106},level=1]{ตโปปกฺกมกถา}
\csromanheader{1}{8. มหาสีหนาทสุตฺต \textbf{ตโปปกฺกมกถา}}
\proseitem{394}{\csromanfolio{157}เอวํ วุตฺเต อเจโล กสฺสโป ภควนฺตํ เอตทโวจ “อิเมปิ โข อาวุโส โคตม ตโปปกฺกมา เอเตสํ สมณพฺราหฺมณานํ สามญฺญสงฺขาตา จ พฺรหฺมญฺญสงฺขาตา จ.\csromanspacer{} อเจลโก โหติ, มุตฺตาจาโร, หตฺถาปเลขโน, น- เอหิภทฺทนฺติโก, นติฏฺฐภทฺทนฺติโก, นาภิหฏํ, น อุทฺทิสฺสกตํ, น นิมนฺตนํ สาทิยติ, โส น กุมฺภิมุขา ปฏิคฺคณฺหาติ, น กโฬปิมุขา ปฏิคฺคณฺหาติ, น เอฬกมนฺตรํ, น ทณฺฑมนฺตรํ, น มุสลมนฺตรํ, น ทฺวินฺนํ ภุญฺชมานานํ, น คพฺภินิยา, น ปายมานาย, น ปุริสนฺตรคตาย, น สํกิตฺตีสุ, น ยตฺถ สา อุปฏฺฐิโต โหติ, น ยตฺถ มกฺขิกา สณฺฑสณฺฑจารินี, น มจฺฉํ, น มํสํ, น สุรํ, น เมรยํ, น ถุโสทกํ ปิวติ, โส เอกาคาริโก วา โหติ เอกาโลปิโก, ทฺวาคาริโก วา โหติ ทฺวาโลปิโก, สตฺตาคาริโก วา โหติ สตฺตาโลปิโก, เอกิสฺสาปิ ทตฺติยา ยาเปติ, ทฺวีหิปิ ทตฺตีหิ ยาเปติ, สตฺตหิปิ ทตฺตีหิ ยาเปติ, เอกาหิกมฺปิ อาหารํ อาหาเรติ, ทฺวีหิกมฺปิ อาหารํ อาหาเรติ, สตฺตาหิกมฺปิ อาหารํ อาหาเรติ, อิติ เอวรูปํ อทฺธมาสิกมฺปิ ปริยายภตฺตโภชนานุโยคมนุยุตฺโต วิหรติ.}

\proseitem{395}{อิเมปิ โข อาวุโส โคตม ตโปปกฺกมา เอเตสํ สมณพฺราหฺมณานํ สามญฺญสงฺขาตา จ พฺรหฺมญฺญสงฺขาตา จ.\csromanspacer{} สากภกฺโข วา โหติ, สามากภกฺโข วา โหติ, นีวารภกฺโข วา โหติ, ททฺทุลภกฺโข วา โหติ, หฏภกฺโข วา โหติ, กณภกฺโข วา โหติ, อาจามภกฺโข วา โหติ, ปิญฺญากภกฺโข วา โหติ, ติณภกฺโข วา โหติ, โคมยภกฺโข วา โหติ, วนมูลผลาหาโร ยาเปติ ปวตฺตผลโภชี.}

\proseitem{396}{อิเมปิ โข อาวุโส โคตม ตโปปกฺกมา เอเตสํ สมณพฺราหฺมณานํ สามญฺญสงฺขาตา จ พฺรหฺมญฺญสงฺขาตา จ.\csromanspacer{} สาณานิปิ ธาเรติ, มสาณานิปิ ธาเรติ, ฉวทุสฺสานิปิ ธาเรติ, ปํสุกูลานิปิ ธาเรติ, ติรีฏานิปิ ธาเรติ, อชินมฺปิ ธาเรติ, อชินกฺขิปมฺปิ ธาเรติ, กุสจีรมฺปิ ธาเรติ, วากจีรมฺปิ ธาเรติ, ผลกจีรมฺปิ ธาเรติ, เกสกมฺพลมฺปิ ธาเรติ, วาฬกมฺพลมฺปิ ธาเรติ, อุลูกปกฺขิกมฺปิ ธาเรติ, เกสมสฺสุโลจโกปิ โหติ เกสมสฺสุโลจนานุโยคมนุยุตฺโต, \csromanfolio{158}อุพฺภฏฺฐโกปิ\footnote{อุพฺภฏฺฐิโกปิ (ก)} โหติ อาสนปฏิกฺขิตฺโต, อุกฺกุฏิโกปิ โหติ อุกฺกุฏิกปฺปธานมนุยุตฺโต, กณฺฏกาปสฺสยิโกปิ โหติ กณฺฏกาปสฺสเย เสยฺยํ กปฺเปติ, ผลกเสยฺยมฺปิ กปฺเปติ, ถณฺฑิลเสยฺยมฺปิ กปฺเปติ, เอกปสฺสยิโกปิ โหติ รโชชลฺลธโร, อพฺโภกาสิโกปิ โหติ ยถาสนฺถติโก, เวกฏิโกปิ โหติ วิกฏโภชนานุโยคมนุยุตฺโต, อปานโกปิ โหติ อปานกตฺตมนุยุตฺโต, สายตติยกมฺปิ อุทโกโรหนานุโยคมนุยุตฺโต วิหรตี”ติ.}

\csromanmatikaanchor{mtk.107}
\csromantocmarkref{section}{ตโปปกฺกมนิรตฺถกตา}{mtk.107}
\bookmark[dest={mtk.107},level=1]{ตโปปกฺกมนิรตฺถกตา}
\csromanheader{1}{\textbf{ตโปปกฺกมนิรตฺถกตา}}
\proseitem{397}{อเจลโก เจปิ กสฺสป โหติ, มุตฺตาจาโร, หตฺถาปเลขโน ฯเปฯ อิติ เอวรูปํ อทฺธมาสิกมฺปิ ปริยายภตฺตโภชนานุโยคมนุยุตฺโต วิหรติ.\csromanspacer{} ตสฺส จายํ สีลสมฺปทา จิตฺตสมฺปทา ปญฺญาสมฺปทา อภาวิตา โหติ อสจฺฉิกตา.\csromanspacer{} อถ โข โส อารกาว สามญฺญา อารกาว พฺรหฺมญฺญา.\csromanspacer{} ยโต โข กสฺสป ภิกฺขุ อเวรํ อพฺยาปชฺชํ เมตฺตจิตฺตํ ภาเวติ, อาสวานญฺจ ขยา อนาสวํ เจโตวิมุตฺติํ ปญฺญาวิมุตฺติํ ทิฏฺเฐว ธมฺเม สยํ อภิญฺญา สจฺฉิกตฺวา อุปสมฺปชฺช วิหรติ.\csromanspacer{} อยํ วุจฺจติ กสฺสป ภิกฺขุ สมโณ อิติปิ พฺราหฺมโณ อิติปิ.}

\prose{สากภกฺโข เจปิ กสฺสป โหติ, สามากภกฺโข ฯเปฯ วนมูลผลาหาโร ยาเปติ ปวตฺตผลโภชี.\csromanspacer{} ตสฺส จายํ สีลสมฺปทา จิตฺตสมฺปทา ปญฺญาสมฺปทา อภาวิตา โหติ อสจฺฉิกตา.\csromanspacer{} อถ โข โส อารกาว สามญฺญา อารกาว พฺราหฺมญฺญา.\csromanspacer{} ยโต โข กสฺสป ภิกฺขุ อเวรํ อพฺยาปชฺชํ เมตฺตจิตฺตํ ภาเวติ, อาสวานญฺจ ขยา อนาสวํ เจโตวิมุตฺติํ ปญฺญาวิมุตฺติํ ทิฏฺเฐว ธมฺเม สยํ อภิญฺญา สจฺฉิกตฺวา อุปสมฺปชฺช วิหรติ.\csromanspacer{} อยํ วุจฺจติ กสฺสป ภิกฺขุ สมโณ อิติปิ พฺราหฺมโณ อิติปิ.}

\prose{สาณานิ เจปิ กสฺสป ธาเรติ, มสาณานิปิ ธาเรติ ฯเปฯ สายตติยกมฺปิ อุทโกโรหนานุโยคมนุยุตฺโต วิหรติ.\csromanspacer{} ตสฺส จายํ สีลสมฺปทา จิตฺตสมฺปทา ปญฺญาสมฺปทา อภาวิตา โหติ อสจฺฉิกตา.\csromanspacer{} อถ โข โส อารกาว สามญฺญา อารกาว}

\vspace{\tipitakaparskip}
\csromancenter{มหาสีหนาทสุตฺต พฺรหฺมญฺญา.\csromanspacer{} ยโต โข กสฺสป ภิกฺขุ อเวรํ อพฺยาปชฺชํ เมตฺตจิตฺตํ ภาเวติ, อาสวานญฺจ ขยา อนาสวํ เจโตวิมุตฺติํ ปญฺญาวิมุตฺติํ ทิฏฺเฐว ธมฺเม สยํ อภิญฺญา สจฺฉิกตฺวา อุปสมฺปชฺช วิหรติ.\csromanspacer{} อยํ วุจฺจติ กสฺสป ภิกฺขุ สมโณ อิติปิ พฺราหฺมโณ อิติปีติ.}
\proseitem{398}{\csromanfolio{159}เอวํ วุตฺเต อเจโล กสฺสโป ภควนฺตํ เอตทโวจ “ทุกฺกรํ โภ โคตม สามญฺญํ ทุกฺกรํ พฺรหฺมญฺญนฺ”ติ.\csromanspacer{} ปกติ โข เอสา กสฺสป โลกสฺมิํ ทุกฺกรํ สามญฺญํ ทุกฺกรํ พฺรหฺมญฺญนฺติ.\csromanspacer{} อเจลโก เจปิ กสฺสป โหติ, มุตฺตาจาโร, หตฺถาปเลขโน ฯเปฯ อิติ เอวรูปํ อทฺธมาสิกมฺปิ ปริยายภตฺตโภชนานุโยคมนุยุตฺโต วิหรติ.\csromanspacer{} อิมาย จ กสฺสป มตฺตาย อิมินา ตโปปกฺกเมน สามญฺญํ วา อภวิสฺส พฺรหฺมญฺญํ วา ทุกฺกรํ สุทุกฺกรํ, เนตํ อภวิสฺส กลฺลํ วจนาย “ทุกฺกรํ สามญฺญํ ทุกฺกรํ พฺรหฺมญฺญนฺ”ติ.}

\prose{สกฺกา จ ปเนตํ อภวิสฺส กาตุํ คหปตินา วา คหปติปุตฺเตน วา อนฺตมโส กุมฺภทาสิยาปิ “หนฺทาหํ อเจลโก โหมิ, มุตฺตาจาโร, หตฺถาปเลขโน ฯเปฯ อิติ เอวรูปํ อทฺธมาสิกมฺปิ ปริยายภตฺตโภชนานุโยคมนุยุตฺโต วิหรามี”ติ.}

\prose{ยสฺมา จ โข กสฺสป อญฺญเตฺรว อิมาย มตฺตาย อญฺญตฺร อิมินา ตโปปกฺกเมน สามญฺญํ วา โหติ พฺรหฺมญฺญํ วา ทุกฺกรํ สุทุกฺกรํ, ตสฺมา เอตํ กลฺลํ วจนาย “ทุกฺกรํ สามญฺญํ ทุกฺกรํ พฺรหฺมญฺญนฺ”ติ.\csromanspacer{} ยโต โข กสฺสป ภิกฺขุ อเวรํ อพฺยาปชฺชํ เมตฺตจิตฺตํ ภาเวติ, อาสวานญฺจ ขยา อนาสวํ เจโตวิมุตฺติํ ปญฺญาวิมุตฺติํ ทิฏฺเฐว ธมฺเม สยํ อภิญฺญา สจฺฉิกตฺวา อุปสมฺปชฺช วิหรติ.\csromanspacer{} อยํ วุจฺจติ กสฺสป ภิกฺขุ สมโณ อิติปิ พฺราหฺมโณ อิติปิ.}

\prose{สากภกฺโข เจปิ กสฺสป โหติ, สามากภกฺโข ฯเปฯ วนมูลผลาหาโร ยาเปติ ปวตฺตผลโภชี.\csromanspacer{} อิมาย จ กสฺสป มตฺตาย อิมินา ตโปปกฺกเมน สามญฺญํ วา อภวิสฺส พฺรหฺมญฺญํ วา ทุกฺกรํ สุทุกฺกรํ, เนตํ อภวิสฺส กลฺลํ วจนาย “ทุกฺกรํ สามญฺญํ ทุกฺกรํ พฺรหฺมญฺญนฺ”ติ.}

\prose{สกฺกา จ ปเนตํ อภวิสฺส กาตุํ คหปตินา วา คหปติปุตฺเตน วา อนฺตมโส กุมฺภทาสิยาปิ “หนฺทาหํ สากภกฺโข วา โหมิ, สามากภกฺโข วา ฯเปฯ วนมูลผลาหาโร ยาเปมิ ปวตฺตผลโภชี”ติ.}

\prose{\csromanfolio{160}ยสฺมา จ โข กสฺสป อญฺญเตฺรว อิมาย มตฺตาย อญฺญตฺร อิมินา ตโปปกฺกเมน สามญฺญํ วา โหติ พฺรหฺมญฺญํ วา ทุกฺกรํ สุทุกฺกรํ, ตสฺมา เอตํ กลฺลํ วจนาย “ทุกฺกรํ สามญฺญํ ทุกฺกรํ พฺรหฺมญฺญนฺ”ติ.\csromanspacer{} ยโต โข กสฺสป ภิกฺขุ อเวรํ อพฺยาปชฺชํ เมตฺตจิตฺตํ ภาเวติ, อาสวานญฺจ ขยา อนาสวํ เจโตวิมุตฺติํ ปญฺญาวิมุตฺติํ ทิฏฺเฐว ธมฺเม สยํ อภิญฺญา สจฺฉิกตฺวา อุปสมฺปชฺช วิหรติ.\csromanspacer{} อยํ วุจฺจติ กสฺสป ภิกฺขุ สมโณ อิติปิ พฺราหฺมโณ อิติปิ.}

\prose{สาณานิ เจปิ กสฺสป ธาเรติ, มสาณานิปิ ธาเรติ ฯเปฯ สายตติยกมฺปิ อุทโกโรหนานุโยคมนุยุตฺโต วิหรติ.\csromanspacer{} อิมาย จ กสฺสป มตฺตาย อิมินา ตโปปกฺกเมน สามญฺญํ วา อภวิสฺส พฺรหฺมญฺญํ วา ทุกฺกรํ สุทุกฺกรํ, เนตํ อภวิสฺส กลฺลํ วจนาย “ทุกฺกรํ สามญฺญํ ทุกฺกรํ พฺรหฺมญฺญนฺ”ติ.}

\prose{สกฺกา จ ปเนตํ อภวิสฺส กาตุํ คหปตินา วา คหปติปุตฺเตน วา อนฺตมโส กุมฺภทาสิยาปิ “หนฺทาหํ สาณานิปิ ธาเรมิ, มสาณานิปิ ธาเรมิ ฯเปฯ สายตติยกมฺปิ อุทโกโรหนานุโยคมนุยุตฺโต วิหรามี”ติ.}

\prose{ยสฺมา จ โข กสฺสป อญฺญเตฺรว อิมาย มตฺตาย อญฺญตฺร อิมินา ตโปปกฺกเมน สามญฺญํ วา โหติ พฺรหฺมญฺญํ วา ทุกฺกรํ สุทุกฺกรํ, ตสฺมา เอตํ กลฺลํ วจนาย “ทุกฺกรํ สามญฺญํ ทุกฺกรํ พฺรหฺมญฺญนฺ”ติ.\csromanspacer{} ยโต โข กสฺสป ภิกฺขุ อเวรํ อพฺยาปชฺชํ เมตฺตจิตฺตํ ภาเวติ, อาสวานญฺจ ขยา อนาสวํ เจโตวิมุตฺติํ ปญฺญาวิมุตฺติํ ทิฏฺเฐว ธมฺเม สยํ อภิญฺญา สจฺฉิกตฺวา อุปสมฺปชฺช วิหรติ.\csromanspacer{} อยํ วุจฺจติ กสฺสป ภิกฺขุ สมโณ อิติปิ พฺราหฺมโณ อิติปีติ.}

\proseitem{399}{เอวํ วุตฺเต อเจโล กสฺสโป ภควนฺตํ เอตทโวจ “ทุชฺชาโน โภ โคตม สมโณ ทุชฺชาโน พฺราหฺมโณ”ติ.\csromanspacer{} ปกติ โข เอสา กสฺสป โลกสฺมิํ ทุชฺชาโน สมโณ ทุชฺชาโน พฺราหฺมโณติ.\csromanspacer{} อเจลโก เจปิ กสฺสป โหติ, มุตฺตาจาโร, หตฺถาปเลขโน ฯเปฯ อิติ เอวรูปํ อทฺธมาสิกมฺปิ ปริยายภตฺตโภชนานุโยคมนุยุตฺโต วิหรติ.\csromanspacer{} อิมาย จ กสฺสป มตฺตาย อิมินา ตโปปกฺกเมน สมโณ วา}

\vspace{\tipitakaparskip}
\csromancenter{มหาสีหนาทสุตฺต อภวิสฺส พฺราหฺมโณ วา ทุชฺชาโน สุทุชฺชาโน, เนตํ อภวิสฺส กลฺลํ วจนาย “ทุชฺชาโน สมโณ ทุชฺชาโน พฺราหฺมโณ”ติ.}
\prose{\csromanfolio{161}สกฺกา จ ปเนโส อภวิสฺส ญาตุํ คหปตินา วา คหปติปุตฺเตน วา อนฺตมโส กุมฺภทาสิยาปิ “อยํ อเจลโก โหติ, มุตฺตาจาโร, หตฺถาปเลขโน ฯเปฯ อิติ เอวรูปํ อทฺธมาสิกมฺปิ ปริยายภตฺตโภชานานุโยคมนุยุตฺโต วิหรตี”ติ.}

\prose{ยสฺมา จ โข กสฺสป อญฺญเตฺรว อิมาย มตฺตาย อญฺญตฺร อิมินา ตโปปกฺกเมน สมโณ วา โหติ พฺราหฺมโณ วา ทุชฺชาโน สุทุชฺชาโน, ตสฺมา เอตํ กลฺลํ วจนาย “ทุชฺชาโน สมโณ ทุชฺชาโน พฺราหฺมโณ”ติ.\csromanspacer{} ยโต โข\footnote{ยโต จ โข (ก)} กสฺสป ภิกฺขุ อเวรํ อพฺยาปชฺชํ เมตฺตจิตฺตํ ภาเวติ, อาสวานญฺจ ขยา อนาสวํ เจโตวิมุตฺติํ ปญฺญาวิมุตฺติํ ทิฏฺเฐว ธมฺเม สยํ อภิญฺญา สจฺฉิกตฺวา อุปสมฺปชฺช วิหรติ.\csromanspacer{} อยํ วุจฺจติ กสฺสป ภิกฺขุ สมโณ อิติปิ พฺราหฺมโณ อิติปิ.}

\prose{สากภกฺโข เจปิ กสฺสป โหติ สามากภกฺโข ฯเปฯ วนมูลผลาหาโร ยาเปติ ปวตฺตผลโภชี.\csromanspacer{} อิมาย จ กสฺสป มตฺตาย อิมินา ตโปปกฺกเมน สมโณ วา อภวิสฺส พฺราหฺมโณ วา ทุชฺชาโน สุทุชฺชาโน, เนตํ อภวิสฺส กลฺลํ วจนาย “ทุชฺชาโน สมโณ ทุชฺชาโน พฺราหฺมโณ”ติ.}

\prose{สกฺกา จ ปเนโส อภวิสฺส ญาตุํ คหปตินา วา คหปติปุตฺเตน วา อนฺตมโส กุมฺภทาสิยาปิ “อยํ สากภกฺโข วา โหติ สามากภกฺโข ฯเปฯ วนมูลผลาหาโร ยาเปติ ปวตฺตผลโภชี”ติ.}

\prose{ยสฺมา จ โข กสฺสป อญฺญเตฺรว อิมาย มตฺตาย อญฺญตฺร อิมินา ตโปปกฺกเมน สมโณ วา โหติ พฺราหฺมโณ วา ทุชฺชาโน สุทุชฺชาโน, ตสฺมา เอตํ กลฺลํ วจนาย “ทุชฺชาโน สมโณ ทุชฺชาโน พฺราหฺมโณ”ติ.\csromanspacer{} ยโต โข กสฺสป ภิกฺขุ อเวรํ อพฺยาปชฺชํ เมตฺตจิตฺตํ ภาเวติ, อาสวานญฺจ ขยา อนาสวํ เจโตวิมุตฺติํ ปญฺญาวิมุตฺติํ ทิฏฺเฐว ธมฺเม สยํ อภิญฺญา สจฺฉิกตฺวา อุปสมฺปชฺช วิหรติ.\csromanspacer{} อยํ วุจฺจติ กสฺสป ภิกฺขุ สมโณ อิติปิ พฺราหฺมโณ อิติปิ.}

\csromanmatikaanchor{mtk.108}
\csromantocmarkref{section}{สีลสมาทิปญฺญาสมฺปทา}{mtk.108}
\bookmark[dest={mtk.108},level=1]{สีลสมาทิปญฺญาสมฺปทา}
\prose{\csromanfolio{162}สาณานิ เจปิ กสฺสป ธาเรติ, มสาณานิปิ ธาเรติ ฯเปฯ สายตติยกมฺปิ อุทโกโรหนานุโยคมนุยุตฺโต วิหรติ.\csromanspacer{} อิมาย จ กสฺสป มตฺตาย อิมินา ตโปปกฺกเมน สมโณ วา อภวิสฺส พฺราหฺมโณ วา ทุชฺชาโน สุทุชฺชาโน, เนตํ อภวิสฺส กลฺลํ วจนาย “ทุชฺชาโน สมโณ ทุชฺชาโน พฺราหฺมโณ”ติ.}

\prose{สกฺกา จ ปเนโส อภวิสฺส ญาตุํ คหปตินา วา คหปติปุตฺเตน วา อนฺตมโส กุมฺภทาสิยาปิ “อยํ สาณานิปิ ธาเรติ, มสาณานิปิ ธาเรติ ฯเปฯ สายตติยกมฺปิ อุทโกโรหนานุโยคมนุยุตฺโต วิหรตี”ติ.}

\prose{ยสฺมา จ โข กสฺสป อญฺญเตฺรว อิมาย มตฺตาย อญฺญตฺร อิมินา ตโปปกฺกเมน สมโณ วา โหติ พฺราหฺมโณ วา ทุชฺชาโน สุทุชฺชาโน, ตสฺมา เอตํ กลฺลํ วจนาย “ทุชฺชาโน สมโณ ทุชฺชาโน พฺราหฺมโณ”ติ.\csromanspacer{} ยโต โข กสฺสป ภิกฺขุ อเวรํ อพฺยาปชฺชํ เมตฺตจิตฺตํ ภาเวติ, อาสวานญฺจ ขยา อนาสวํ เจโตวิมุตฺติํ ปญฺญาวิมุตฺติํ ทิฏฺเฐว ธมฺเม สยํ อภิญฺญา สจฺฉิกตฺวา อุปสมฺปชฺช วิหรติ.\csromanspacer{} อยํ วุจฺจติ กสฺสป ภิกฺขุ สมโณ อิติปิ พฺราหฺมโณ อิติปีติ.}

\csromanheader{2}{สีลสมาธิปญฺญา\-สมฺปทา}
\proseitem{400}{เอวํ วุตฺเต อเจโล กสฺสโป ภควนฺตํ เอตทโวจ “กตมา ปน สา โภ โคตม สีลสมฺปทา, กตมา จิตฺตสมฺปทา, กตมา ปญฺญาสมฺปทา”ติ.\csromanspacer{} อิธ กสฺสป ตถาคโต โลเก อุปฺปชฺชติ อรหํ สมฺมาสมฺพุทฺโธ ฯเปฯ ภยทสฺสาวี สมาทาย สิกฺขติ สิกฺขาปเทสุ, กายกมฺมวจีกมฺเมน สมนฺนาคโต กุสเลน ปริสุทฺธาชีโว สีลสมฺปนฺโน อินฺทฺริเยสุ คุตฺตทฺวาโร สติสมฺปชญฺเญน สมนฺนาคโต สนฺตุฏฺโฐ.}

\proseitem{401}{กถญฺจ กสฺสป ภิกฺขุ สีลสมฺปนฺโน โหติ.\csromanspacer{} อิธ กสฺสป ภิกฺขุ ปาณาติปาตํ ปหาย ปาณาติปาตา ปฏิวิรโต โหติ นิหิตทณฺโฑ นิหิตสตฺโถ ลชฺชี ทยาปนฺโน, สพฺพปาณภูต\-หิตา\-นุกมฺปี วิหรติ.\csromanspacer{} อิทมฺปิสฺส โหติ สีลสมฺปทาย ฯเปฯ.}

\prose{ยถา วา ปเนเก โภนฺโต สมณพฺราหฺมณา สทฺธาเทยฺยานิ โภชนานิ ภุญฺชิตฺวา เต เอวรูปาย ติรจฺฉานวิชฺชาย มิจฺฉาชีเวน ชีวิตํ}

\vspace{\tipitakaparskip}
\csromancenter{มหาสีหนาทสุตฺต กปฺเปนฺติ.\csromanspacer{} เสยฺยถิทํ, สนฺติกมฺมํ ปณิธิกมฺมํ ฯเปฯ โอสธีนํ ปติโมกฺโข อิติ วา, อิติ เอวรูปาย ติรจฺฉานวิชฺชาย มิจฺฉาชีวา ปฏิวิรโต โหติ.\csromanspacer{} อิทมฺปิสฺส โหติ สีลสมฺปทาย.}
\prose{\csromanfolio{163}ส โขโส\footnote{อยํ โข (ก)} กสฺสป ภิกฺขุ เอวํ สีลสมฺปนฺโน น กุโตจิ ภยํ สมนุปสฺสติ, ยทิทํ สีลสํวรโต.\csromanspacer{} เสยฺยถาปิ กสฺสป ราชา ขตฺติโย มุทฺธาวสิตฺโต นิหตปจฺจมิตฺโต น กุโตจิ ภยํ สมนุปสฺสติ, ยทิทํ ปจฺจตฺถิกโต.\csromanspacer{} เอวเมว โข กสฺสป ภิกฺขุ เอวํ สีลสมฺปนฺโน น กุโตจิ ภยํ สมนุปสฺสติ, ยทิทํ สีลสํวรโต, โส อิมินา อริเยน สีลกฺขนฺเธน สมนฺนาคโต อชฺฌตฺตํ อนวชฺชสุขํ ปฏิสํเวเทติ.\csromanspacer{} เอวํ โข กสฺสป ภิกฺขุ สีลสมฺปนฺโน โหติ.\csromanspacer{} อยํ โข กสฺสป สีลสมฺปทา.}

\prose{ฯเปฯ ปฐมํ ฌานํ อุปสมฺปชฺช วิหรติ.\csromanspacer{} อิทมฺปิสฺส โหติ จิตฺตสมฺปทาย ฯเปฯ ทุติยํ ฌานํ.\csromanspacer{} ตติยํ ฌานํ.\csromanspacer{} จตุตฺถํ ฌานํ อุปสมฺปชฺช วิหรติ.\csromanspacer{} อิทมฺปิสฺส โหติ จิตฺตสมฺปทาย.\csromanspacer{} อยํ โข กสฺสป จิตฺตสมฺปทา.}

\prose{โส เอวํ สมาหิเต จิตฺเต ฯเปฯ ญาณทสฺสนาย จิตฺตํ อภินีหรติ อภินินฺนาเมติ.\csromanspacer{} อิทมฺปิสฺส โหติ ปญฺญาสมฺปทาย ฯเปฯ นาปรํ อิตฺถตฺตายาติ ปชานาติ.\csromanspacer{} อิทมฺปิสฺส โหติ ปญฺญาสมฺปทาย.\csromanspacer{} อยํ โข กสฺสป ปญฺญาสมฺปทา.}

\prose{อิมาย จ กสฺสป สีลสมฺปทาย จิตฺตสมฺปทาย ปญฺญาสมฺปทาย อญฺญา สีลสมฺปทา จิตฺตสมฺปทา ปญฺญาสมฺปทา อุตฺตริตรา วา ปณีตตรา วา นตฺถิ.}

\csromanmatikaanchor{mtk.109}
\csromantocmarkref{section}{สีหนาทกถา}{mtk.109}
\bookmark[dest={mtk.109},level=1]{สีหนาทกถา}
\csromanheader{1}{\textbf{สีหนาทกถา}}
\proseitem{402}{สนฺติ กสฺสป เอเก สมณพฺราหฺมณา สีลวาทา, เต อเนกปริยาเยน สีลสฺส วณฺณํ ภาสนฺติ, ยาวตา กสฺสป อริยํ ปรมํ สีลํ, นาหํ ตตฺถ อตฺตโน สมสมํ สมนุปสฺสามิ, กุโต ภิยฺโย, อถ โข อหเมว ตตฺถ ภิยฺโย, ยทิทํ อธิสีลํ.}

\prose{สนฺติ กสฺสป เอเก สมณพฺราหฺมณา ตโปชิคุจฺฉาวาทา, เต อเนกปริยาเยน ตโปชิคุจฺฉาย วณฺณํ ภาสนฺติ, ยาวตา กสฺสป อริยา \csromanfolio{164}ปรมา ตโปชิคุจฺฉา, นาหํ ตตฺถ อตฺตโน สมสมํ สมนุปสฺสามิ, กุโต ภิยฺโย, อถ โข อหเมว ตตฺถ ภิยฺโย, ยทิทํ อธิเชคุจฺฉํ.}

\prose{สนฺติ กสฺสป เอเก สมณพฺราหฺมณา ปญฺญาวาทา, เต อเนกปริยาเยน ปญฺญาย วณฺณํ ภาสนฺติ, ยาวตา กสฺสป อริยา ปรมา ปญฺญา, นาหํ ตตฺถ อตฺตโน สมสมํ สมนุปสฺสามิ, กุโต ภิยฺโย, อถ โข อหเมว ตตฺถ ภิยฺโย, ยทิทํ อธิปญฺญํ.}

\prose{สนฺติ กสฺสป เอเก สมณพฺราหฺมณา วิมุตฺติวาทา, เต อเนกปริยาเยน วิมุตฺติยา วณฺณํ ภาสนฺติ, ยาวตา กสฺสป อริยา ปรมา วิมุตฺติ, นาหํ ตตฺถ อตฺตโน สมสมํ สมนุปสฺสามิ, กุโต ภิยฺโย, อถ โข อหเมว ตตฺถ ภิยฺโย, ยทิทํ อธิวิมุตฺติ.}

\proseitem{403}{ฐานํ โข ปเนตํ กสฺสป วิชฺชติ, ยํ อญฺญติตฺถิยา ปริพฺพาชกา เอวํ วเทยฺยุํ “สีหนาทํ โข สมโณ โคตโม นทติ, ตญฺจ โข สุญฺญาคาเร นทติ, โน ปริสาสู”ติ.\csromanspacer{} เต “มา เหวนฺ”ติสฺสุ วจนียา, “สีหนาทญฺจ สมโณ โคตโม นทติ, ปริสาสุ จ นทตี”ติ เอวมสฺสุ กสฺสป วจนียา.}

\prose{ฐานํ โข ปเนตํ กสฺสป วิชฺชติ, ยํ อญฺญติตฺถิยา ปริพฺพาชกา เอวํ วเทยฺยุํ “สีหนาทญฺจ สมโณ โคตโม นทติ, ปริสาสุ จ นทติ, โน จ โข วิสารโท นทตี”ติ.\csromanspacer{} เต “มา เหวนฺ”ติสฺสุ วจนียา, “สีหนาทญฺจ สมโณ โคตโม นทติ, ปริสาสุ จ นทติ, วิสารโท จ นทตี”ติ เอวมสฺสุ กสฺสป วจนียา.}

\prose{ฐานํ โข ปเนตํ กสฺสป วิชฺชติ, ยํ อญฺญติตฺถิยา ปริพฺพาชกา เอวํ วเทยฺยุํ “สีหนาทญฺจ สมโณ โคตโม นทติ, ปริสาสุ จ นทติ, วิสารโท จ นทติ, โน จ โข นํ ปญฺหํ ปุจฺฉนฺติ ฯเปฯ ปญฺหญฺจ นํ ปุจฺฉนฺติ, โน จ โข เนสํ ปญฺหํ ปุฏฺโฐ พฺยากโรติ ฯเปฯ ปญฺหญฺจ เนสํ ปุฏฺโฐ พฺยากโรติ, โน จ โข ปญฺหสฺส เวยฺยากรเณน จิตฺตํ อาราเธติ ฯเปฯ ปญฺหสฺส จ เวยฺยากรเณน จิตฺตํ อาราเธติ, โน จ โข โสตพฺพํ มญฺญนฺติ ฯเปฯ โสตพฺพํ จสฺส มญฺญนฺติ, โน จ โข สุตฺวา ปสีทนฺติ ฯเปฯ สุตฺวา จสฺส}

\vspace{\tipitakaparskip}
\csromancenter{มหาสีหนาทสุตฺต ปสีทนฺติ, โน จ โข ปสนฺนาการํ กโรนฺติ ฯเปฯ ปสนฺนาการญฺจ กโรนฺติ, โน จ โข ตถตฺตาย ปฏิปชฺชนฺติ ฯเปฯ ตถตฺตาย จ ปฏิปชฺชนฺติ, โน จ โข ปฏิปนฺนา อาราเธนฺตี”ติ.\csromanspacer{} เต “มา เหวนฺ”ติสฺสุ วจนียา, “สีหนาทญฺจ สมโณ โคตโม นทติ, ปริสาสุ จ นทติ, วิสารโท จ นทติ, ปญฺหญฺจ นํ ปุจฺฉนฺติ, ปญฺหญฺจ เนสํ ปุฏฺโฐ พฺยากโรติ, ปญฺหสฺส จ เวยฺยากรเณน จิตฺตํ อาราเธติ, โสตพฺพญฺจสฺส มญฺญนฺติ, สุตฺวา จสฺส ปสีทนฺติ, ปสนฺนาการญฺจ กโรนฺติ, ตถตฺตาย จ ปฏิปชฺชนฺติ, ปฏิปนฺนา จ อาราเธนฺตี”ติ เอวมสฺสุ กสฺสป วจนียา.}
\csromanheader{2}{\textbf{ตตฺถิยปริวาสกถา}}
\csromanmatikaanchor{mtk.110}
\csromantocmarkref{section}{ติตฺถิยปริวาสกถา}{mtk.110}
\bookmark[dest={mtk.110},level=1]{ติตฺถิยปริวาสกถา}
\proseitem{404}{\csromanfolio{165}เอกมิทาหํ กสฺสป สมยํ ราชคเห วิหรามิ คิชฺฌกูเฏ ปพฺพเต, ตตฺร มํ อญฺญตโร ตปพฺรหฺมจารี นิโคฺรโธ นาม อธิเชคุจฺเฉ ปญฺหํ อปุจฺฉิ, ตสฺสาหํ อธิเชคุจฺเฉ ปญฺหํ ปุฏฺโฐ พฺยากาสิํ, พฺยากเต จ ปน เม อตฺตมโน อโหสิ ปรํ วิย มตฺตายาติ.\csromanspacer{} โก หิ ภนฺเต ภควโต ธมฺมํ สุตฺวา น อตฺตมโน อสฺส ปรํ วิย มตฺตาย.\csromanspacer{} อหมฺปิ หิ ภนฺเต ภควโต ธมฺมํ สุตฺวา อตฺตมโน ปรํ วิย มตฺตาย.\csromanspacer{} อภิกฺกนฺตํ ภนฺเต, อภิกฺกนฺตํ ภนฺเต.\csromanspacer{} เสยฺยถาปิ ภนฺเต นิกฺกุชฺชิตํ วา อุกฺกุชฺเชยฺย, ปฏิจฺฉนฺนํ วา วิวเรยฺย, มูฬฺหสฺส วา มคฺคํ อาจิกฺเขยฺย, อนฺธกาเร วา เตลปชฺโชตํ ธาเรยฺย “จกฺขุมนฺโต รูปานิ ทกฺขนฺตี”ติ.\csromanspacer{} เอวเมวํ ภควตา อเนกปริยาเยน ธมฺโม ปกาสิโต, เอสาหํ ภนฺเต ภควนฺตํ สรณํ คจฺฉามิ ธมฺมญฺจ ภิกฺขุสํฆญฺจ, ลเภยฺยาหํ ภนฺเต ภควโต สนฺติเก ปพฺพชฺชํ, ลเภยฺยํ อุปสมฺปทนฺติ.}

\proseitem{405}{โย โข กสฺสป อญฺญติตฺถิยปุพฺโพ อิมสฺมิํ ธมฺมวินเย อากงฺขติ ปพฺพชฺชํ, อากงฺขติ อุปสมฺปทํ, โส จตฺตาโร มาเส ปริวสติ, จตุนฺนํ มาสานํ อจฺจเยน อารทฺธจิตฺตา ภิกฺขู ปพฺพาเชนฺติ, อุปสมฺปาเทนฺติ ภิกฺขุภาวาย.\csromanspacer{} อปิ จ เมตฺถ ปุคฺคลเวมตฺตตา วิทิตาติ.\csromanspacer{} สเจ ภนฺเต อญฺญติตฺถิยปุพฺพา อิมสฺมิํ ธมฺมวินเย อากงฺขนฺติ ปพฺพชฺชํ, อากงฺขนฺติ อุปสมฺปทํ, จตฺตาโร มาเส ปริวสนฺติ, จตุนฺนํ มาสานํ อจฺจเยน อารทฺธจิตฺตา ภิกฺขู ปพฺพาเชนฺติ, อุปสมฺปาเทนฺติ ภิกฺขุภาวาย.\csromanspacer{} อหํ จตฺตาริ วสฺสานิ ปริวสิสฺสามิ, จตุนฺนํ วสฺสานํ อจฺจเยน อารทฺธจิตฺตา ภิกฺขู ปพฺพาเชนฺตุ, อุปสมฺปาเทนฺตุ ภิกฺขุภาวายาติ.}

\prosewithcloser{\csromanfolio{166}อลตฺถ โข อเจโล กสฺสโป ภควโต สนฺติเก ปพฺพชฺชํ, อลตฺถ อุปสมฺปทํ, อจิรูปสมฺปนฺโน โข ปนายสฺมา กสฺสโป เอโก วูปกฏฺโฐ อปฺปมตฺโต อาตาปี ปหิตตฺโต วิหรนฺโต น จิรสฺเสว ยสฺสตฺถาย กุลปุตฺตา สมฺมเทว อคารสฺมา อนคาริยํ ปพฺพชนฺติ, ตทนุตฺตรํ พฺรหฺมจริยปริโยสานํ ทิฏฺเฐว ธมฺเม สยํ อภิญฺญา สจฺฉิกตฺวา อุปสมฺปชฺช วิหาสิ, “ขีณา ชาติ, วุสิตํ พฺรหฺมจริยํ, กตํ กรณียํ, นาปรํ อิตฺถตฺตายา”ติ อพฺภญฺญาสิ, อญฺญตโร โข ปนายสฺมา กสฺสโป อรหตํ อโหสีติ.}{\textbf{มหาสีหนาทสุตฺตํ นิฏฺฐิตํ อฏฺฐมํ.}}
\csromanmatikaanchor{mtk.111}
\csromantocmarknopage{chapter}{9. โปฏฺฐปาทสุตฺต}{mtk.111}
\bookmark[dest={mtk.111},level=0]{9. โปฏฺฐปาทสุตฺต}
\setcsromanheadcha{9. โปฏฺฐปาทสุตฺต}
\setcsromanheadhclear{}
\chapterheadpageread{9. \textbf{โปฏฺฐปาทสุตฺต}}
\csromanmatikaanchor{mtk.112}
\csromantocmarkref{section}{โปฏฺฐปาทปริพฺพาชกวตฺถุ}{mtk.112}
\bookmark[dest={mtk.112},level=1]{โปฏฺฐปาทปริพฺพาชกวตฺถุ}
\csromanheader{1}{โปฏฺฐปาท\-ปริพฺพาชก\-วตฺถุ}
\proseitem{406}{\csromanfolio{167}เอวํ เม สุตํ\csromandash{}เอกํ สมยํ ภควา สาวตฺถิยํ วิหรติ เชตวเน อนาถปิณฺฑิกสฺส อาราเม.\csromanspacer{} เตน โข ปน สมเยน โปฏฺฐปาโท ปริพฺพาชโก สมยปฺปวาทเก ตินฺทุกาจีเร เอกสาลเก มลฺลิกาย อาราเม ปฏิวสติ มหติยา ปริพฺพาชกปริสาย สทฺธิํ ติํสมตฺเตหิ ปริพฺพาชกสเตหิ.\csromanspacer{} อถ โข ภควา ปุพฺพณฺหสมยํ นิวาเสตฺวา ปตฺตจีวรมาทาย สาวตฺถิํ ปิณฺฑาย ปาวิสิ.}

\proseitem{407}{อถ โข ภควโต เอตทโหสิ “อติปฺปโค โข ตาว สาวตฺถิยํ ปิณฺฑาย จริตุํ, ยํนูนาหํ เยน สมยปฺปวาทโก ตินฺทุกาจีโร เอกสาลโก มลฺลิกาย อาราโม, เยน โปฏฺฐปาโท ปริพฺพาชโก เตนุปสงฺกเมยฺยนฺ”ติ.\csromanspacer{} อถ โข ภควา เยน สมยปฺปวาทโก ตินฺทุกาจีโร เอกสาลโก มลฺลิกาย อาราโม เตนุปสงฺกมิ.}

\proseitem{408}{เตน โข ปน สมเยน โปฏฺฐปาโท ปริพฺพาชโก มหติยา ปริพฺพาชกปริสาย สทฺธิํ นิสินฺโน โหติ อุนฺนาทินิยา อุจฺจาสทฺท\-มหาสทฺทาย อเนกวิหิตํ ติรจฺฉานกถํ กเถนฺติยา.\csromanspacer{} เสยฺยถิทํ, ราชกถํ โจรกถํ มหามตฺตกถํ เสนากถํ ภยกถํ ยุทฺธกถํ อนฺนกถํ ปานกถํ วตฺถกถํ สยนกถํ มาลากถํ คนฺธกถํ ญาติกถํ ยานกถํ คามกถํ นิคมกถํ นครกถํ ชนปทกถํ อิตฺถิกถํ สูรกถํ วิสิขากถํ กุมฺภฏฺฐานกถํ ปุพฺพเปตกถํ นานตฺตกถํ โลกกฺขายิกํ สมุทฺทกฺขายิกํ อิติภวาภวกถํ อิติ วา.}

\proseitem{409}{อทฺทสา โข โปฏฺฐปาโท ปริพฺพาชโก ภควนฺตํ ทูรโตว อาคจฺฉนฺตํ, ทิสฺวาน สกํ ปริสํ สณฺฐเปสิ “อปฺปสทฺทา โภนฺโต โหนฺตุ, มา โภนฺโต สทฺทมกตฺถ, อยํ สมโณ โคตโม อาคจฺฉติ, อปฺปสทฺทกาโม โข โส อายสฺมา อปฺปสทฺทสฺส วณฺณวาที, อปฺเปว นาม อปฺปสทฺทํ ปริสํ วิทิตฺวา อุปสงฺกมิตพฺพํ มญฺเญยฺยา”ติ.\csromanspacer{} เอวํ วุตฺเต เต ปริพฺพาชกา ตุณฺหี อเหสุํ.}

\proseitem{410}{\csromanfolio{168}อถ โข ภควา เยน โปฏฺฐปาโท ปริพฺพาชโก เตนุปสงฺกมิ.\csromanspacer{} อถ โข โปฏฺฐปาโท ปริพฺพาชโก ภควนฺตํ เอตทโวจ “เอตุ โข ภนฺเต ภควา, สฺวาคตํ ภนฺเต ภควโต, จิรสฺสํ โข ภนฺเต ภควา อิมํ ปริยายมกาสิ, ยทิทํ อิธาคมนาย, นิสีทตุ ภนฺเต ภควา, อิทํ อาสนํ ปญฺญตฺตนฺ”ติ.}

\prose{นิสีทิ ภควา ปญฺญตฺเต อาสเน.\csromanspacer{} โปฏฺฐปาโทปิ โข ปริพฺพาชโก อญฺญตรํ นีจํ อาสนํ คเหตฺวา เอกมนฺตํ นิสีทิ.\csromanspacer{} เอกมนฺตํ นิสินฺนํ โข โปฏฺฐปาทํ ปริพฺพาชกํ ภควา เอตทโวจ “กาย นุตฺถ\footnote{กาย โนตฺถ (สฺยา, ก)} โปฏฺฐปาท เอตรหิ กถาย สนฺนิสินฺนา, กา จ ปน โว อนฺตรากถา วิปฺปกตา”ติ.}

\csromanmatikaanchor{mtk.113}
\csromantocmarkref{section}{อภิสญฺญานิโรธกถา}{mtk.113}
\bookmark[dest={mtk.113},level=1]{อภิสญฺญานิโรธกถา}
\csromanheader{1}{\textbf{อภิสญฺญานิโรธกถา}}
\proseitem{411}{เอวํ วุตฺเต โปฏฺฐปาโท ปริพฺพาชโก ภควนฺตํ เอตทโวจ “ติฏฺฐเตสา ภนฺเต กถา, ยาย มยํ เอตรหิ กถาย สนฺนิสินฺนา, เนสา ภนฺเต กถา ภควโต ทุลฺลภา ภวิสฺสติ ปจฺฉาปิ สวนาย, ปุริมานิ ภนฺเต ทิวสานิ ปุริมตรานิ, นานาติตฺถิยานํ สมณพฺราหฺมณานํ โกตูหลสาลาย สนฺนิสินฺนานํ สนฺนิปติตานํ อภิสญฺญานิโรเธ กถา อุทปาทิ ‘กถํ นุ โข โภ อภิสญฺญานิโรโธ โหตี’ติ.\csromanspacer{} ตเตฺรกจฺเจ เอวมาหํสุ ‘อเหตู อปฺปจฺจยา ปุริสสฺส สญฺญา อุปฺปชฺชนฺติปิ นิรุชฺฌนฺติปิ.\csromanspacer{} ยสฺมิํ สมเย อุปฺปชฺชนฺติ, สญฺญี ตสฺมิํ สมเย โหติ.\csromanspacer{} ยสฺมิํ สมเย นิรุชฺฌนฺติ, อสญฺญี ตสฺมิํ สมเย โหตี’ติ.\csromanspacer{} อิตฺเถเก อภิสญฺญานิโรธํ ปญฺญเปนฺติ.}

\prose{ตมญฺโญ เอวมาห ‘น โข ปน เมตํ\footnote{น โข นาเมตํ (สี, อิ)} โภ เอวํ ภวิสฺสติ, สญฺญา หิ โภ ปุริสสฺส อตฺตา, สา จ โข อุเปติปิ อเปติปิ.\csromanspacer{} ยสฺมิํ สมเย อุเปติ, สญฺญี ตสฺมิํ สมเย โหติ.\csromanspacer{} ยสฺมิํ สมเย อเปติ, อสญฺญี ตสฺมิํ สมเย โหตี’ติ.\csromanspacer{} อิตฺเถเก อภิสญฺญานิโรธํ ปญฺญเปนฺติ.}

\prose{ตมญฺโญ เอวมาห ‘น โข ปน เมตํ โภ เอวํ ภวิสฺสติ, สนฺติ หิ โภ สมณพฺราหฺมณา มหิทฺธิกา มหานุภาวา, เต อิมสฺส ปุริสสฺส สญฺญํ อุปกฑฺฒนฺติปิ อปกฑฺฒนฺติปิ.\csromanspacer{} ยสฺมิํ สมเย อุปกฑฺฒนฺติ, สญฺญี ตสฺมิํ}

\vspace{\tipitakaparskip}
\csromancenter{โปฏฺฐปาทสุตฺต สมเย โหติ.\csromanspacer{} ยสฺมิํ สมเย อปกฑฺฒนฺติ, อสญฺญี ตสฺมิํ สมเย โหตี’ติ.\csromanspacer{} อิตฺเถเก อภิสญฺญานิโรธํ ปญฺญเปนฺติ.}
\prose{\csromanfolio{169}ตมญฺโญ เอวมาห ‘น โข ปน เมตํ โภ เอวํ ภวิสฺสติ, สนฺติ หิ โภ เทวตา มหิทฺธิกา มหานุภาวา, ตา อิมสฺส ปุริสสฺส สญฺญํ อุปกฑฺฒนฺติปิ อปกฑฺฒนฺติปิ.\csromanspacer{} ยสฺมิํ สมเย อุปกฑฺฒนฺติ, สญฺญี ตสฺมิํ สมเย โหติ.\csromanspacer{} ยสฺมิํ สมเย อปกฑฺฒนฺติ, อสญฺญี ตสฺมิํ สมเย โหตี’ติ.\csromanspacer{} อิตฺเถเก อภิสญฺญานิโรธํ ปญฺญเปนฺติ.}

\prose{ตสฺส มยฺหํ ภนฺเต ภควนฺตํเยว อารพฺภ สติ อุทปาทิ ‘อโห นูน ภควา อโห นูน สุคโต, โย อิเมสํ ธมฺมานํ สุกุสโล’ติ.\csromanspacer{} ภควา ภนฺเต กุสโล ภควา ปกตญฺญู อภิสญฺญานิโรธสฺส, กถํ นุ โข ภนฺเต อภิสญฺญานิโรโธ โหตี”ติ.}

\csromanmatikaanchor{mtk.114}
\csromantocmarkref{section}{สเหตุกสญฺญุปฺปาทนิโรธกถา}{mtk.114}
\bookmark[dest={mtk.114},level=1]{สเหตุกสญฺญุปฺปาทนิโรธกถา}
\csromanheader{1}{สเหตุก\-สญฺญุปฺปาท\-นิโรธ\-กถา}
\proseitem{412}{ตตฺร โปฏฺฐปาท เย เต สมณพฺราหฺมณา เอวมาหํสุ “อเหตู อปฺปจฺจยา ปุริสสฺส สญฺญา อุปฺปชฺชนฺติปิ นิรุชฺฌนฺติปี”ติ, อาทิโตว เตสํ อปรทฺธํ.\csromanspacer{} ตํ กิสฺส เหตุ, สเหตู หิ โปฏฺฐปาท สปฺปจฺจยา ปุริสสฺส สญฺญา อุปฺปชฺชนฺติปิ นิรุชฺฌนฺติปิ.\csromanspacer{} สิกฺขา เอกา สญฺญา อุปฺปชฺชติ, สิกฺขา เอกา สญฺญา นิรุชฺฌติ.}

\proseitem{413}{กา จ สิกฺขาติ ภควา อโวจ.\csromanspacer{} อิธ โปฏฺฐปาท ตถาคโต โลเก อุปฺปชฺชติ อรหํ สมฺมาสมฺพุทฺโธ ฯเปฯ. (ยถา สามญฺญผเล, เอวํ วิตฺถาเรตพฺพํ.) เอวํ โข โปฏฺฐปาท ภิกฺขุ สีลสมฺปนฺโน โหติ ฯเปฯ ตสฺสิเม ปญฺจนีวรเณ ปหีเน อตฺตนิ สมนุปสฺสโต ปาโมชฺชํ ชายติ, ปมุทิตสฺส ปีติ ชายติ, ปีติมนสฺส กาโย ปสฺสมฺภติ, ปสฺสทฺธกาโย สุขํ เวเทติ, สุขิโน จิตฺตํ สมาธิยติ, โส วิวิจฺเจว กาเมหิ วิวิจฺจ อกุสเลหิ ธมฺเมหิ สวิตกฺกํ สวิจารํ วิเวกชํ ปีติสุขํ ปฐมํ ฌานํ อุปสมฺปชฺช วิหรติ, ตสฺส ยา ปุริมา กามสญฺญา, สา นิรุชฺฌติ.\csromanspacer{} วิเวกช\-ปีติสุข\-สุขุม\-สจฺจ\-สญฺญา ตสฺมิํ สมเย โหติ, วิเวกช\-ปีติสุข\-สุขุมสจฺจ\-สญฺญีเยว ตสฺมิํ สมเย โหติ.\csromanspacer{} เอวมฺปิ สิกฺขา เอกา สญฺญา อุปฺปชฺชติ, สิกฺขา เอกา สญฺญา นิรุชฺฌติ.\csromanspacer{} อยํ สิกฺขาติ ภควา อโวจ.}

\prose{\csromanfolio{170}ปุน จปรํ โปฏฺฐปาท ภิกฺขุ วิตกฺกวิจารานํ วูปสมา อชฺฌตฺตํ สมฺปสาทนํ เจตโส เอโกทิภาวํ อวิตกฺกํ อวิจารํ สมาธิชํ ปีติสุขํ ทุติยํ ฌานํ อุปสมฺปชฺช วิหรติ.\csromanspacer{} ตสฺส ยา ปุริมา วิเวกช\-ปีติสุข\-สุขุม\-สจฺจ\-สญฺญา, สา นิรุชฺฌติ.\csromanspacer{} สมาธิช\-ปีติสุข\-สุขุมสจฺจ\-สญฺญา ตสฺมิํ สมเย โหติ, สมาธิช\-ปีติสุข\-สุขุมสจฺจ\-สญฺญีเยว ตสฺมิํ สมเย โหติ.\csromanspacer{} เอวมฺปิ สิกฺขา เอกา สญฺญา อุปฺปชฺชติ, สิกฺขา เอกา สญฺญา นิรุชฺฌติ.\csromanspacer{} อยมฺปิ สิกฺขาติ ภควา อโวจ.}

\prose{ปุน จปรํ โปฏฺฐปาท ภิกฺขุ ปีติยา จ วิราคา อุเปกฺขโก จ วิหรติ สโต จ สมฺปชาโน, สุขญฺจ กาเยน ปฏิสํเวเทติ, ยํ ตํ อริยา อาจิกฺขนฺติ “อุเปกฺขโก สติมา สุขวิหารี”ติ, ตติยํ ฌานํ อุปสมฺปชฺช วิหรติ.\csromanspacer{} ตสฺส ยา ปุริมา สมาธิช\-ปีติสุข\-สุขุมสจฺจ\-สญฺญา, สา นิรุชฺฌติ.\csromanspacer{} อุเปกฺขาสุข\-สุขุม\-สจฺจ\-สญฺญา ตสฺมิํ สมเย โหติ, อุเปกฺขาสุข\-สุขุมสจฺจ\-สญฺญีเยว ตสฺมิํ สมเย โหติ.\csromanspacer{} เอวมฺปิ สิกฺขา เอกา สญฺญา อุปฺปชฺชติ, สิกฺขา เอกา สญฺญา นิรุชฺฌติ.\csromanspacer{} อยมฺปิ สิกฺขาติ ภควา อโวจ.}

\prose{ปุน จปรํ โปฏฺฐปาท ภิกฺขุ สุขสฺส จ ปหานา ทุกฺขสฺส จ ปหานา ปุพฺเพว โสมนสฺสโทมนสฺสานํ อตฺถงฺคมา อทุกฺขมสุขํ อุเปกฺขาสติปาริสุทฺธิํ จตุตฺถํ ฌานํ อุปสมฺปชฺช วิหรติ.\csromanspacer{} ตสฺส ยา ปุริมา อุเปกฺขาสุข\-สุขุม\-สจฺจ\-สญฺญา, สา นิรุชฺฌติ.\csromanspacer{} อทุกฺขมสุข\-สุขุม\-สจฺจ\-สญฺญา ตสฺมิํ สมเย โหติ, อทุกฺขมสุข\-สุขุม\-สจฺจ\-สญฺญีเยว ตสฺมิํ สมเย โหติ.\csromanspacer{} เอวมฺปิ สิกฺขา เอกา สญฺญา อุปฺปชฺชติ, สิกฺขา เอกา สญฺญา นิรุชฺฌติ.\csromanspacer{} อยมฺปิ สิกฺขาติ ภควา อโวจ.}

\prose{ปุน จปรํ โปฏฺฐปาท ภิกฺขุ สพฺพโส รูปสญฺญานํ สมติกฺกมา ปฏิฆสญฺญานํ อตฺถงฺคมา นานตฺตสญฺญานํ อมนสิการา “อนนฺโต อากาโส”ติ อากาสานญฺจายตนํ อุปสมฺปชฺช วิหรติ.\csromanspacer{} ตสฺส ยา ปุริมา รูปสญฺญา\footnote{ปุริมสญฺญา (ก)}, สา นิรุชฺฌติ.\csromanspacer{} อากาสานญฺจายตน\-สุขุม\-สจฺจ\-สญฺญา ตสฺมิํ สมเย โหติ, อากาสานญฺจายตน\-สุขุม\-สจฺจ\-สญฺญีเยว ตสฺมิํ สมเย โหติ.\csromanspacer{} เอวมฺปิ สิกฺขา เอกา สญฺญา อุปฺปชฺชติ, สิกฺขา เอกา สญฺญา นิรุชฺฌติ.\csromanspacer{} อยมฺปิ สิกฺขาติ ภควา อโวจ.}

\csromanheader{2}{9. โปฏฺฐปาทสุตฺต}
\prose{\csromanfolio{171}ปุน จปรํ โปฏฺฐปาท ภิกฺขุ สพฺพโส อากาสานญฺจายตนํ สมติกฺกมฺม “อนนฺตํ วิญฺญาณนฺ”ติ วิญฺญาณญฺจายตนํ อุปสมฺปชฺช วิหรติ, ตสฺส ยา ปุริมา อากาสานญฺจายตน\-สุขุม\-สจฺจ\-สญฺญา, สา นิรุชฺฌติ.\csromanspacer{} วิญฺญาณญฺจายตน\-สุขุม\-สจฺจ\-สญฺญา ตสฺมิํ สมเย โหติ, วิญฺญาณญฺจายตน\-สุขุมสจฺจ\-สญฺญีเยว ตสฺมิํ สมเย โหติ.\csromanspacer{} เอวมฺปิ สิกฺขา เอกา สญฺญา อุปฺปชฺชติ, สิกฺขา เอกา สญฺญา นิรุชฺฌติ.\csromanspacer{} อยมฺปิ สิกฺขาติ ภควา อโวจ.}

\prose{ปุน จปรํ โปฏฺฐปาท ภิกฺขุ สพฺพโส วิญฺญาณญฺจายตนํ สมติกฺกมฺม “นตฺถิ กิญฺจี”ติ อากิญฺจญฺญายตนํ อุปสมฺปชฺช วิหรติ.\csromanspacer{} ตสฺส ยา ปุริมา วิญฺญาณญฺจายตน\-สุขุม\-สจฺจ\-สญฺญา, สา นิรุชฺฌติ.\csromanspacer{} อากิญฺจญฺญายตน\-สุขุม\-สจฺจ\-สญฺญา ตสฺมิํ สมเย โหติ, อากิญฺจญฺญายตน\-สุขุม\-สจฺจ\-สญฺญีเยว ตสฺมิํ สมเย โหติ.\csromanspacer{} เอวมฺปิ สิกฺขา เอกา สญฺญา อุปฺปชฺชติ, สิกฺขา เอกา สญฺญา นิรุชฺฌติ.\csromanspacer{} อยมฺปิ สิกฺขาติ ภควา อโวจ.}

\proseitem{414}{ยโต โข โปฏฺฐปาท ภิกฺขุ อิธ สกสญฺญี โหติ, โส ตโต อมุตฺร ตโต อมุตฺร อนุปุพฺเพน สญฺญคฺคํ ผุสติ, ตสฺส สญฺญคฺเค ฐิตสฺส เอวํ โหติ “เจตยมานสฺส เม ปาปิโย, อเจตยมานสฺส เม เสยฺโย.\csromanspacer{} อหญฺเจว โข ปน เจเตยฺยํ, อภิสงฺขเรยฺยํ, อิมา จ เม สญฺญา นิรุชฺเฌยฺยุํ, อญฺญา จ โอฬาริกา สญฺญา อุปฺปชฺเชยฺยุํ, ยํนูนาหํ น เจว เจเตยฺยํ, น จ อภิสงฺขเรยฺยนฺ”ติ.\csromanspacer{} โส น เจว เจเตหิ, น จ อภิสงฺขโรติ.\csromanspacer{} ตสฺส อเจตยโต อนภิสงฺขโรโต ตา เจว สญฺญา นิรุชฺฌนฺติ, อญฺญา จ โอฬาริกา สญฺญา น อุปฺปชฺชนฺติ, โส นิโรธํ ผุสติ.\csromanspacer{} เอวํ โข โปฏฺฐปาท อนุปุพฺพา\-ภิสญฺญานิโรธ\-สมฺปชาน\-สมาปตฺติ โหติ.}

\prose{ตํ กิํ มญฺญสิ โปฏฺฐปาท, อปิ นุ เต อิโต ปุพฺเพ เอวรูปา อนุปุพฺพา\-ภิสญฺญานิโรธ\-สมฺปชาน\-สมาปตฺติ สุตปุพฺพาติ.\csromanspacer{} โน เหตํ ภนฺเต, เอวํ โข อหํ ภนฺเต ภควโต ภาสิตํ อาชานามิ “ยโต โข โปฏฺฐปาท ภิกฺขุ อิธ สกสญฺญี โหติ, โส ตโต อมุตฺร ตโต อมุตฺร อนุปุพฺเพน สญฺญคฺคํ ผุสติ, ตสฺส สญฺญคฺเคฐิตสฺส เอวํ โหติ ‘เจตยมานสฺส เม ปาปิโย, อเจตยมานสฺส เม เสยฺโย, อหญฺเจว โข ปน \csromanfolio{172}เจเตยฺยํ, อภิสงฺขเรยฺยํ, อิมา จ เม สญฺญา นิรุชฺเฌยฺยุํ, อญฺญา จ โอฬาริกา สญฺญา อุปฺปชฺเชยฺยุํ, ยํนูนาหํ น เจว เจเตยฺยํ, น จ อภิสงฺขเรยฺยนฺ’ติ.\csromanspacer{} โส น เจว เจเตติ, น จาภิสงฺขโรติ, ตสฺส อเจตยโต อนภิสงฺขโรโต ตา เจว สญฺญา นิรุชฺฌนฺติ, อญฺญา จ โอฬาริกา สญฺญา น อุปฺปชฺชนฺติ, โส นิโรธํ ผุสติ.\csromanspacer{} เอวํ โข โปฏฺฐปาท อนุปุพฺพา\-ภิสญฺญานิโรธ\-สมฺปชาน\-สมาปตฺติ โหตี”ติ.\csromanspacer{} เอวํ โปฏฺฐปาทาติ.}

\proseitem{415}{เอกญฺเญว นุ โข ภนฺเต ภควา สญฺญคฺคํ ปญฺญเปติ, อุทาหุ ปุถูปิ สญฺญคฺเค ปญฺญเปตีติ.\csromanspacer{} เอกมฺปิ โข อหํ โปฏฺฐปาท สญฺญคฺคํ ปญฺญเปมิ, ปุถูปิ สญฺญคฺเค ปญฺญเปมีติ.\csromanspacer{} ยถา กถํ ปน ภนฺเต ภควา เอกมฺปิ สญฺญคฺคํ ปญฺญเปติ, ปุถูปิ สญฺญคฺเค ปญฺญเปตีติ.\csromanspacer{} ยถา ยถา โข โปฏฺฐปาท นิโรธํ ผุสติ, ตถา ตถาหํ สญฺญคฺคํ ปญฺญเปมิ.\csromanspacer{} เอวํ โข อหํ โปฏฺฐปาท เอกมฺปิ สญฺญคฺคํ ปญฺญเปมิ, ปุถูปิ สญฺญคฺเค ปญฺญเปมีติ.}

\proseitem{416}{สญฺญา นุ โข ภนฺเต ปฐมํ อุปฺปชฺชติ, ปจฺฉา ญาณํ, อุทาหุ ญาณํ ปฐมํ อุปฺปชฺชติ, ปจฺฉา สญฺญา, อุทาหุ สญฺญา จ ญาณญฺจ อปุพฺพํ อจริมํ อุปฺปชฺชนฺตีติ.\csromanspacer{} สญฺญา โข โปฏฺฐปาท ปฐมํ อุปฺปชฺชติ, ปจฺฉา ญาณํ, สญฺญุปฺปาทา จ ปน ญาณุปฺปาโท โหติ.\csromanspacer{} โส เอวํ ปชานาติ “อิทปฺปจฺจยา กิร เม ญาณํ อุทปาที”ติ, อิมินา โข เอตํ โปฏฺฐปาท ปริยาเยน เวทิตพฺพํ “ยถา สญฺญา ปฐมํ อุปฺปชฺชติ, ปจฺฉา ญาณํ, สญฺญุปฺปาทา จ ปน ญาณุปฺปาโท โหตี”ติ.}

\csromanmatikaanchor{mtk.115}
\csromantocmarkref{section}{สญฺญาอตฺตกถา}{mtk.115}
\bookmark[dest={mtk.115},level=1]{สญฺญาอตฺตกถา}
\csromanheader{1}{\textbf{สญฺญาอตฺตกถา}}
\proseitem{417}{สญฺญา นุ โข ภนฺเต ปุริสสฺส อตฺตา, อุทาหุ อญฺญา สญฺญา อญฺโญ อตฺตาติ.\csromanspacer{} กํ ปน ตฺวํ โปฏฺฐปาท อตฺตานํ ปจฺเจสีติ.\csromanspacer{} โอฬาริกํ โข อหํ ภนฺเต อตฺตานํ ปจฺเจมิ รูปิํ จาตุมหาภูติกํ กพฬีการา\-หาร\-ภกฺขนฺติ\footnote{กพฬีการภกฺขนฺติ (สฺยา, ก)}.\csromanspacer{} โอฬาริโก จ หิ เต โปฏฺฐปาท อตฺตา อภวิสฺส รูปี จาตุมหาภูติโก กพฬีการาหารภกฺโข, เอวํ สนฺตํ โข เต โปฏฺฐปาท อญฺญาว สญฺญา อภวิสฺส อญฺโญ อตฺตาติ.\csromanspacer{} ตทมินาเปตํ}

\vspace{\tipitakaparskip}
\csromancenter{โปฏฺฐปาทสุตฺต โปฏฺฐปาท ปริยาเยน เวทิตพฺพํ ยถา อญฺญาว สญฺญา ภวิสฺสติ อญฺโญ อตฺตา.\csromanspacer{} ติฏฺฐเตว สายํ\footnote{ติฏฺฐเตวายํ (สี, อิ)} โปฏฺฐปาท โอฬาริโก อตฺตา รุปี จาตุมหาภูติโก กพฬีการาหารภกฺโข, อถ อิมสฺส ปุริสสฺส อญฺญา จ สญฺญา อุปฺปชฺชนฺติ, อญฺญา จ สญฺญา นิรุชฺฌนฺติ.\csromanspacer{} อิมินา โข เอตํ โปฏฺฐปาท ปริยาเยน เวทิตพฺพํ “ยถา อญฺญาว สญฺญา ภวิสฺสติ อญฺโญ อตฺตา”ติ.}
\proseitem{418}{\csromanfolio{173}มโนมยํ โข อหํ ภนฺเต อตฺตานํ ปจฺเจมิ “สพฺพงฺคปจฺจงฺคิํ อหีนินฺทฺริยนฺ”ติ.\csromanspacer{} มโนมโย จ หิ เต โปฏฺฐปาท อตฺตา อภวิสฺส สพฺพงฺคปจฺจงฺคี อหีนินฺทฺริโย, เอวํ สนฺตมฺปิ โข เต โปฏฺฐปาท อญฺญาว สญฺญา ภวิสฺสติ อญฺโญ อตฺตา.\csromanspacer{} ตทมินาเปตํ โปฏฺฐปาท ปริยาเยน เวทิตพฺพํ ยถา อญฺญาว สญฺญา ภวิสฺสติ อญฺโญ อตฺตา.\csromanspacer{} ติฏฺฐเตว สายํ โปฏฺฐปาท มโนมโย อตฺตา สพฺพงฺคปจฺจงฺคี อหีนินฺทฺริโย, อถ อิมสฺส ปุริสสฺส อญฺญา จ สญฺญา อุปฺปชฺชนฺติ, อญฺญา จ สญฺญา นิรุชฺฌนฺติ.\csromanspacer{} อิมินาปิ โข เอตํ โปฏฺฐปาท ปริยาเยน เวทิตพฺพํ “ยถา อญฺญาว สญฺญา ภวิสฺสติ อญฺโญ อตฺตา”ติ.}

\proseitem{419}{อรูปิํ โข อหํ ภนฺเต อตฺตานํ ปจฺเจมิ สญฺญามยนฺติ.\csromanspacer{} อรูปี จ หิ เต โปฏฺฐปาท อตฺตา อภวิสฺส สญฺญามโย, เอวํ สนฺตมฺปิ โข เต โปฏฺฐปาท อญฺญาว สญฺญา ภวิสฺสติ อญฺโญ อตฺตา.\csromanspacer{} ตทมินาเปตํ โปฏฺฐปาท ปริยาเยน เวทิตพฺพํ ยถา อญฺญาว สญฺญา ภวิสฺสติ อญฺโญ อตฺตา.\csromanspacer{} ติฏฺฐเตว สายํ โปฏฺฐปาท อรูปี อตฺตา สญฺญามโย, อถ อิมสฺส ปุริสสฺส อญฺญา จ สญฺญา อุปฺปชฺชนฺติ, อญฺญา จ สญฺญา นิรุชฺฌนฺติ.\csromanspacer{} อิมินาปิ โข เอตํ โปฏฺฐปาท ปริยาเยน เวทิตพฺพํ “ยถา อญฺญาว สญฺญา ภวิสฺสติ อญฺโญ อตฺตา”ติ.}

\proseitem{420}{สกฺกา ปเนตํ ภนฺเต มยา ญาตุํ “สญฺญา ปุริสสฺส อตฺตา”ติ วา, “อญฺญาว สญฺญา อญฺโญ อตฺตา”ติ วาติ.\csromanspacer{} ทุชฺชานํ โข เอตํ\footnote{เอวํ (ก)} โปฏฺฐปาท ตยา อญฺญทิฏฺฐิเกน อญฺญขนฺติเกน อญฺญรุจิเกน อญฺญตฺราโยเคน อญฺญตฺราจริยเกน “สญฺญา ปุริสสฺส อตฺตา”ติ วา, “อญฺญาว สญฺญา อญฺโญ อตฺตา”ติ วาติ.}

\prose{\csromanfolio{174}สเจตํ ภนฺเต มยา ทุชฺชานํ อญฺญทิฏฺฐิเกน อญฺญขนฺติเกน อญฺญรุจิเกน อญฺญตฺราโยเคน อญฺญตฺราจริยเกน “สญฺญา ปุริสสฺส อตฺตา”ติ วา, “อญฺญาว สญฺญา อญฺโญ อตฺตา”ติ วา.\csromanspacer{} กิํ ปน ภนฺเต สสฺสโต โลโก, อิทเมว สจฺจํ โมฆมญฺญนฺติ.\csromanspacer{} อพฺยากตํ โข เอตํ โปฏฺฐปาท มยา “สสฺสโต โลโก, อิทเมว สจฺจํ โมฆมญฺญนฺ”ติ.}

\prose{กิํ ปน ภนฺเต อสสฺสโต โลโก, อิทเมว สจฺจํ โมฆมญฺญนฺติ.\csromanspacer{} เอตมฺปิ โข โปฏฺฐปาท มยา อพฺยากตํ “อสสฺสโต โลโก, อิทเมว สจฺจํ โมฆมญฺญนฺ”ติ.}

\prose{กิํ ปน ภนฺเต อนฺตวา โลโก ฯเปฯ อนนฺตวา โลโก.\csromanspacer{} ตํ ชีวํ ตํ สรีรํ.\csromanspacer{} อญฺญํ ชีวํ อญฺญํ สรีรํ.\csromanspacer{} โหติ ตถาคโต ปรํ มรณา.\csromanspacer{} น โหติ ตถาคโต ปรํ มรณา.\csromanspacer{} โหติ จ น จ โหติ ตถาคโต ปรํ มรณา.\csromanspacer{} เนว โหติ น น โหติ ตถาคโต ปรํ มรณา, อิทเมว สจฺจํ โมฆมญฺญนฺติ.\csromanspacer{} เอตมฺปิ โข โปฏฺฐปาท มยา อพฺยากตํ “เนว โหติ น น โหติ ตถาคโต ปรํ มรณา, อิทเมว สจฺจํ โมฆมญฺญนฺ”ติ.}

\prose{กสฺมา ปเนตํ ภนฺเต ภควตา อพฺยากตนฺติ.\csromanspacer{} น เหตํ โปฏฺฐปาท อตฺถสํหิตํ น ธมฺมสํหิตํ นาทิพฺรหฺมจริยกํ, น นิพฺพิทาย น วิราคาย น นิโรธาย น อุปสมาย น อภิญฺญาย น สมฺโพธาย น นิพฺพานาย สํวตฺตติ, ตสฺมา เอตํ มยา อพฺยากตนฺติ.}

\prose{กิํ ปน ภนฺเต ภควตา พฺยากตนฺติ.\csromanspacer{} อิทํ ทุกฺขนฺติ โข โปฏฺฐปาท มยา พฺยากตํ.\csromanspacer{} อยํ ทุกฺขสมุทโยติ โข โปฏฺฐปาท มยา พฺยากตํ.\csromanspacer{} อยํ ทุกฺขนิโรโธติ โข โปฏฺฐปาท มยา พฺยากตํ.\csromanspacer{} อยํ ทุกฺขนิโรธคามินี ปฏิปทาติ โข โปฏฺฐปาท มยา พฺยากตนฺติ.}

\prose{กสฺมา ปเนตํ ภนฺเต ภควตา พฺยากตนฺติ.\csromanspacer{} เอตํ หิ โปฏฺฐปาท อตฺถสํหิตํ, เอตํ ธมฺมสํหิตํ, เอตํ อาทิพฺรหฺมจริยกํ, เอตํ นิพฺพิทาย วิราคาย นิโรธาย อุปสมาย อภิญฺญาย สมฺโพธาย นิพฺพานาย สํวตฺตติ, ตสฺมา เอตํ มยา พฺยากตนฺติ.\csromanspacer{} เอวเมตํ ภควา, เอวเมตํ สุคต, ยสฺสทานิ ภนฺเต ภควา กาลํ มญฺญตีติ.\csromanspacer{} อถ โข ภควา อุฏฺฐายาสนา ปกฺกามิ.}

\csromanheader{2}{9. โปฏฺฐปาทสุตฺต}
\proseitem{421}{\csromanfolio{175}อถ โข เต ปริพฺพาชกา อจิรปกฺกนฺตสฺส ภควโต โปฏฺฐปาทํ ปริพฺพาชกํ สมนฺตโต วาจา\footnote{วาจาย (สฺยา, ก)} สนฺนิโตทเกน สญฺฌพฺภริมกํสุ “เอวเมว ปนายํ ภวํ โปฏฺฐปาโท ยญฺญเทว สมโณ โคตโม ภาสติ, ตํ ตเทวสฺส อพฺภนุโมทติ ‘เอวเมตํ ภควา เอวเมตํ สุคตา’ติ.\csromanspacer{} น โข ปน มยํ กิญฺจิ\footnote{กญฺจิ (อิ)} สมณสฺส โคตมสฺส เอกํสิกํ ธมฺมํ เทสิตํ อาชานาม ‘สสฺสโต โลโก’ติ วา, ‘อสสฺสโต โลโก’ติ วา, ‘อนฺตวา โลโก’ติ วา, ‘อนนฺตวา โลโก’ติ วา, ‘ตํ ชีวํ ตํ สรีรนฺ’ติ วา, ‘อญฺญํ ชีวํ อญฺญํ สรีรนฺ’ติ วา, ‘โหติ ตถาคโต ปรํ มรณา’ติ วา, ‘น โหติ ตถาคโต ปรํ มรณา’ติ วา, ‘โหติ จ น จ โหติ ตถาคโต ปรํ มรณา’ติ วา, ‘เนว โหติ น น โหติ ตถาคโต ปรํ มรณา’ติ วา”ติ.}

\prose{เอวํ วุตฺเต โปฏฺฐปาโท ปริพฺพาชโก เต ปริพฺพาชเก เอตทโวจ “อหมฺปิ โข โภ น กิญฺจิ สมณสฺส โคตมสฺส เอกํสิกํ ธมฺมํ เทสิตํ อาชานามิ ‘สสฺสโต โลโก’ติ วา, ‘อสสฺสโต โลโก’ติ วา ฯเปฯ ‘เนว โหติ น น โหติ ตถาคโต ปรํ มรณา’ติ วา, อปิ จ สมโณ โคตโม ภูตํ ตจฺฉํ ตถํ ปฏิปทํ ปญฺญเปติ ธมฺมฏฺฐิตตํ ธมฺมนิยามตํ, ภูตํ โข ปน ตจฺฉํ ตถํ ปฏิปทํ ปญฺญเปนฺตสฺส ธมฺมฏฺฐิตตํ ธมฺมนิยามตํ, กถํ หิ นาม มาทิโส วิญฺญู สมณสฺส โคตมสฺส สุภาสิตํ สุภาสิตโต นาพฺภานุโมเทยฺยา”ติ.}

\csromanmatikaanchor{mtk.116}
\csromantocmarkref{section}{จิตฺตหตฺถิสาริปุตฺตโปฏฺฐปาทวตฺถุ}{mtk.116}
\bookmark[dest={mtk.116},level=1]{จิตฺตหตฺถิสาริปุตฺตโปฏฺฐปาทวตฺถุ}
\csromanheader{1}{จิตฺต\-หตฺถิสาริปุตฺต\-โปฏฺฐปาท\-วตฺถุ}
\proseitem{422}{อถ โข ทฺวีหตีหสฺส อจฺจเยน จิตฺโต จ หตฺถิสาริปุตฺโต โปฏฺฐปาโท จ ปริพฺพาชโก เยน ภควา เตนุปสงฺกมิํสุ, อุปสงฺกมิตฺวา จิตฺโต หตฺถิสาริปุตฺโต ภควนฺตํ อภิวาเทตฺวา เอกมนฺตํ นิสีทิ.\csromanspacer{} โปฏฺฐปาโท ปน ปริพฺพาชโก ภควตา สทฺธิํ สมฺโมทิ, สมฺโมทนียํ กถํ สารณียํ วีติสาเรตฺวา เอกมนฺตํ นิสีทิ, เอกมนฺตํ นิสินฺโน โข โปฏฺฐปาโท ปริพฺพาชโก ภควนฺตํ เอตทโวจ\csromandash{}ตทา มํ ภนฺเต เต ปริพฺพาชกา อจิรปกฺกนฺตสฺส ภควโต สมนฺตโต วาจาสนฺนิโตทเกน สญฺฌพฺภริมกํสุ “เอวเมว ปนายํ ภวํ โปฏฺฐปาโท ยญฺญเทว \csromanfolio{176}สมโณ โคตโม ภาสติ, ตํ ตเทวสฺส อพฺภนุโมทติ ‘เอวเมตํ ภควา เอวเมตํ สุคตา’ติ.\csromanspacer{} น โข ปน มยํ กิญฺจิ สมณสฺส โคตมสฺส เอกํสิกํ ธมฺมํ เทสิตํ อาชานาม ‘สสฺสโต โลโก’ติ วา, ‘อสสฺสโต โลโก’ติ วา, ‘อนฺตวา โลโก’ติ วา, ‘อนนฺตวา โลโก’ติ วา, ‘ตํ ชีวํ ตํ สรีรนฺ’ติ วา, ‘อญฺญํ ชีวํ อญฺญํ สรีรนฺ’ติ วา, ‘โหติ ตถาคโต ปรํ มรณา’ติ วา, ‘น โหติ ตถาคโต ปรํ มรณา’ติ วา, ‘โหติ จ น จ โหติ ตถาคโต ปรํ มรณา’ติ วา, ‘เนว โหติ น น โหติ ตถาคโต ปรํ มรณา’ติ วา”ติ.\csromanspacer{} เอวํ วุตฺตาหํ ภนฺเต เต ปริพฺพาชเก เอตทโวจํ “อหมฺปิ โข โภ น กิญฺจิ สมณสฺส โคตมสฺส เอกํสิกํ ธมฺมํ เทสิตํ อาชานามิ ‘สสฺสโต โลโก’ติ วา, ‘อสสฺสโต โลโก’ติ วา ฯเปฯ ‘เนว โหติ น น โหติ ตถาคโต ปรํ มรณา’ติ วา, อปิ จ สมโณ โคตโม ภูตํ ตจฺฉํ ตถํ ปฏิปทํ ปญฺญเปติ ธมฺมฏฺฐิตตํ ธมฺมนิยามตํ, ภูตํ โข ปน ตจฺฉํ ตถํ ปฏิปทํ ปญฺญเปนฺตสฺส ธมฺมฏฺฐิตตํ ธมฺมนิยามตํ, กถํ หิ นาม มาทิโส วิญฺญู สมณสฺส โคตมสฺส สุภาสิตํ สุภาสิตโต นาพฺภนุโมเทยฺยา”ติ.}

\proseitem{423}{สพฺเพว โข เอเต โปฏฺฐปาท ปริพฺพาชกา อนฺธา อจกฺขุกา, ตฺวํเยว เนสํ เอโก จกฺขุมา.\csromanspacer{} เอกํสิกาปิ หิ โข โปฏฺฐปาท มยา ธมฺมา เทสิตา ปญฺญตฺตา, อเนกํสิกาปิ หิ โข โปฏฺฐปาท มยา ธมฺมา เทสิตา ปญฺญตฺตา.}

\prose{กตเม จ เต โปฏฺฐปาท มยา อเนกํสิกา ธมฺมา เทสิตา ปญฺญตฺตา.\csromanspacer{} สสฺสโต โลโกติ\footnote{โลโกติ วา (สี, ก)} โข โปฏฺฐปาท มยา อเนกํสิโก ธมฺโม เทสิโต ปญฺญตฺโต, อสสฺสโต โลโกติ1 โข โปฏฺฐปาท มยา อเนกํสิโก ธมฺโม เทสิโต ปญฺญตฺโต, อนฺตวา โลโกติ1 โข โปฏฺฐปาท ฯเปฯ อนนฺตวา โลโกติ1 โข โปฏฺฐปาท.\csromanspacer{} ตํ ชีวํ ตํ สรีรนฺติ โข โปฏฺฐปาท.\csromanspacer{} อญฺญํ ชีวํ อญฺญํ สรีรนฺติ โข โปฏฺฐปาท.\csromanspacer{} โหติ ตถาคโต ปรํ มรณาติ โข โปฏฺฐปาท.\csromanspacer{} น โหติ ตถาคโต ปรํ มรณาติ โข โปฏฺฐปาท.\csromanspacer{} โหติ จ น จ โหติ ตถาคโต ปรํ มรณาติ โข โปฏฺฐปาท.}

\vspace{\tipitakaparskip}
\csromancenter{โปฏฺฐปาทสุตฺต เนว โหติ น น โหติ ตถาคโต ปรํ มรณาติ โข โปฏฺฐปาท มยา อเนกํสิโก ธมฺโม เทสิโต ปญฺญตฺโต.}
\prose{\csromanfolio{177}กสฺมา จ เต โปฏฺฐปาท มยา อเนกํสิกา ธมฺมา เทสิตา ปญฺญตฺตา.\csromanspacer{} นเหเต โปฏฺฐปาท อตฺถสํหิตา น ธมฺมสํหิตา น อาทิพฺรหฺมจริยกา น นิพฺพิทาย น วิราคาย น นิโรธาย น อุปสมาย น อภิญฺญาย น สมฺโพธาย น นิพฺพานาย สํวตฺตนฺติ.\csromanspacer{} ตสฺมา เต มยา อเนกํสิกา ธมฺมา เทสิตา ปญฺญตฺตา.}

\csromanmatikaanchor{mtk.117}
\csromantocmarkref{section}{เอกํสิกธมฺมา}{mtk.117}
\bookmark[dest={mtk.117},level=1]{เอกํสิกธมฺมา}
\csromanheader{1}{\textbf{เอกํสิกธมฺมา}}
\proseitem{424}{กตเม จ เต โปฏฺฐปาท มยา เอกํสิกา ธมฺมา เทสิตา ปญฺญตฺตา.\csromanspacer{} อิทํ ทุกฺขนฺติ โข โปฏฺฐปาท มยา เอกํสิโก ธมฺโม เทสิโต ปญฺญตฺโต, อยํ ทุกฺขสมุทโยติ โข โปฏฺฐปาท มยา เอกํสิโก ธมฺโม เทสิโต ปญฺญตฺโต, อยํ ทุกฺขนิโรโธติ โข โปฏฺฐปาท มยา เอกํสิโก ธมฺโม เทสิโต ปญฺญตฺโต, อยํ ทุกฺขนิโรธคามินี ปฏิปทาติ โข โปฏฺฐปาท มยา เอกํสิโก ธมฺโม เทสิโต ปญฺญตฺโต.}

\prose{กสฺมา จ เต โปฏฺฐปาท มยา เอกํสิกา ธมฺมา เทสิตา ปญฺญตฺตา.\csromanspacer{} เอเต หิ โปฏฺฐปาท อตฺถสํหิตา เอเต ธมฺมสํหิตา เอเต อาทิพฺรหฺมจริยกา เอเต นิพฺพิทาย วิราคาย นิโรธาย อุปสมาย อภิญฺญาย สมฺโพธาย นิพฺพานาย สํวตฺตนฺติ.\csromanspacer{} ตสฺมา เต มยา เอกํสิกา ธมฺมา เทสิตา ปญฺญตฺตา.}

\proseitem{425}{สนฺติ โปฏฺฐปาท เอเก สมณพฺราหฺมณา เอวํวาทิโน เอวํทิฏฺฐิโน “เอกนฺตสุขี อตฺตา โหติ อโรโค ปรํ มรณา”ติ.\csromanspacer{} ตฺยาหํ อุปสงฺกมิตฺวา เอวํ วทามิ “สจฺจํ กิร ตุเมฺห อายสฺมนฺโต เอวํวาทิโน เอวํทิฏฺฐิโน เอกนฺตสุขี อตฺตา โหติ อโรโค ปรํ มรณา”ติ.\csromanspacer{} เต เจ เม เอวํ ปุฏฺฐา “อามา”ติ ปฏิชานนฺติ.\csromanspacer{} ตฺยาหํ เอวํ วทามิ “อปิ ปน ตุเมฺห อายสฺมนฺโต เอกนฺตสุขํ โลกํ ชานํ ปสฺสํ วิหรถา”ติ, อิติ ปุฏฺฐา “โน”ติ วทนฺติ.}

\prose{ตฺยาหํ เอวํ วทามิ “อปิ ปน ตุเมฺห อายสฺมนฺโต เอกํ วา รตฺติํ เอกํ วา ทิวสํ อุปฑฺฒํ วา รตฺติํ อุปฑฺฒํ วา ทิวสํ เอกนฺตสุขิํ อตฺตานํ \csromanfolio{178}สญฺชานาถา”ติ\footnote{สมฺปชานาถาติ (สี, สฺยา, ก)}, อิติ ปุฏฺฐา “โน”ติ วทนฺติ.\csromanspacer{} ตฺยาหํ เอวํ วทามิ “อปิ ปน ตุเมฺห อายสฺมนฺโต ชานาถ อยํ มคฺโค อยํ ปฏิปทา เอกนฺตสุขสฺส โลกสฺส สจฺฉิกิริยายา”ติ, อิติ ปุฏฺฐา “โน”ติ วทนฺติ.}

\prose{ตฺยาหํ เอวํ วทามิ “อปิ ปน ตุเมฺห อายสฺมนฺโต ยา ตา เทวตา เอกนฺตสุขํ โลกํ อุปปนฺนา, ตาสํ ภาสมานานํ สทฺทํ สุณาถ, สุปฺปฏิปนฺนาตฺถ มาริสา อุชุปฺปฏิปนฺนาตฺถ มาริสา เอกนฺตสุขสฺส โลกสฺส สจฺฉิกิริยาย, มยมฺปิ หิ มาริสา เอวํปฏิปนฺนา เอกนฺตสุขํ โลกํ อุปปนฺนา”ติ, อิติ ปุฏฺฐา “โน”ติ วทนฺติ.}

\prose{ตํ กิํ มญฺญสิ โปฏฺฐปาท “นนุ เอวํ สนฺเต เตสํ สมณพฺราหฺมณานํ อปฺปาฏิหีรกตํ ภาสิตํ สมฺปชฺชตี”ติ.\csromanspacer{} อทฺธา โข ภนฺเต เอวํ สนฺเต เตสํ สมณพฺราหฺมณานํ อปฺปาฏิหีรกตํ ภาสิตํ สมฺปชฺชตีติ.}

\proseitem{426}{เสยฺยถาปิ โปฏฺฐปาท ปุริโส เอวํ วเทยฺย “อหํ ยา อิมสฺมิํ ชนปเท ชนปทกลฺยาณี, ตํ อิจฺฉามิ ตํ กาเมมี”ติ.\csromanspacer{} ตเมนํ เอวํ วเทยฺยุํ “อมฺโภ ปุริส ยํ ตฺวํ ชนปทกลฺยาณิํ อิจฺฉสิ กาเมสิ, ชานาสิ ตํ ชนปทกลฺยาณิํ ขตฺติยี วา พฺราหฺมณี วา เวสฺสี วา สุทฺที วา”ติ, อิติ ปุฏฺโฐ “โน”ติ วเทยฺย.\csromanspacer{} ตเมนํ เอวํ วเทยฺยุํ “อมฺโภ ปุริส ยํ ตฺวํ ชนปทกลฺยาณิํ อิจฺฉสิ กาเมสิ, ชนาสิ ตํ ชนปทกลฺยาณิํ เอวํนามา เอวํโคตฺตาติ วา, ทีฆา วา รสฺสา วา มชฺฌิมา วา กาฬี วา สามา วา มงฺคุรจฺฉวี วาติ, อมุกสฺมิํ คาเม วา นิคเม วา นคเร วา”ติ, อิติ ปุฏฺโฐ “โน”ติ วเทยฺย.\csromanspacer{} ตเมนํ เอวํ วเทยฺยุํ “อมฺโภ ปุริส ยํ ตฺวํ น ชานาสิ น ปสฺสติ, ตํ ตฺวํ อิจฺฉสิ กาเมสี”ติ, อิติ ปุฏฺโฐ “อามา”ติ วเทยฺย.}

\prose{ตํ กิํ มญฺญสิ โปฏฺฐปาท “นนุ เอวํ สนฺเต ตสฺส ปุริสสฺส อปฺปาฏิหีรกตํ ภาสิตํ สมฺปชฺชตี”ติ.\csromanspacer{} อทฺธา โข ภนฺเต เอวํ สนฺเต ตสฺส ปุริสสฺส อปฺปาฏิหีรกตํ ภาสิตํ สมฺปชฺชตีติ.}

\prose{เอวเมว โข โปฏฺฐปาท เย เต สมณพฺราหฺมณา เอวํวาทิโน เอวํทิฏฺฐิโน “เอกนฺตสุขี อตฺตา โหติ อโรโค ปรํ มรณา”ติ.\csromanspacer{} ตฺยาหํ อุปสงฺกมิตฺวา เอวํ วทามิ “สจฺจํ กิร ตุเมฺห อายสฺมนฺโต เอวํวาทิโน}

\vspace{\tipitakaparskip}
\csromancenter{โปฏฺฐปาทสุตฺต เอวํทิฏฺฐิโน เอกนฺตสุขี อตฺตา โหติ อโรโค ปรํ มรณา”ติ.\csromanspacer{} เต เจ เม เอวํ ปุฏฺฐา “อามา”ติ ปฏิชานนฺติ.\csromanspacer{} ตฺยาหํ เอวํ วทามิ “อปิ ปน ตุเมฺห อายสฺมนฺโต เอกนฺตสุขํ โลกํ ชานํ ปสฺสํ วิหรถา”ติ, อิติ ปุฏฺฐา “โน”ติ วทนฺติ.}
\prose{\csromanfolio{179}ตฺยาหํ เอวํ วทามิ “อปิ ปน ตุเมฺห อายสฺมนฺโต เอกํ วา รตฺติํ เอกํ วา ทิวสํ อุปฑฺฒํ วา รตฺติํ อุปฑฺฒํ วา ทิวสํ เอกนฺตสุขิํ อตฺตานํ สญฺชานาถา”ติ, อิติ ปุฏฺฐา “โน”ติ วทนฺติ.\csromanspacer{} ตฺยาหํ เอวํ วทามิ “อปิ ปน ตุเมฺห อายสฺมนฺโต ชานาถ อยํ มคฺโค อยํ ปฏิปทา เอกนฺตสุขสฺส โลกสฺส สจฺฉิกิริยายา”ติ, อิติ ปุฏฺฐา “โน”ติ วทนฺติ.}

\prose{ตฺยาหํ เอวํ วทามิ “อปิ ปน ตุเมฺห อายสฺมนฺโต ยา ตา เทวตา เอกนฺตสุขํ โลกํ อุปปนฺนา, ตาสํ ภาสมานานํ สทฺทํ สุณาถ, สุปฺปฏิปนฺนาตฺถ มาริสา อุชุปฺปฏิปนฺนาตฺถ มาริสา เอกนฺตสุขสฺส โลกสฺส สจฺฉิกิริยาย, มยมฺปิ หิ มาริสา เอวํปฏิปนฺนา เอกนฺตสุขํ โลกํ อุปปนฺนา”ติ, อิติ ปุฏฺฐา “โน”ติ วทนฺติ.}

\prose{ตํ กิํ มญฺญสิ โปฏฺฐปาท “นนุ เอวํ สนฺเต เตสํ สมณพฺราหฺมณานํ อปฺปาฏิหีรกตํ ภาสิตํ สมฺปชฺชตี”ติ.\csromanspacer{} อทฺธา โข ภนฺเต เอวํ สนฺเต เตสํ สมณพฺราหฺมณานํ อปฺปาฏิหีรกตํ ภาสิตํ สมฺปชฺชตีติ.}

\proseitem{427}{เสยฺยถาปิ โปฏฺฐปาท ปุริโส จาตุมหาปเถ นิสฺเสณิํ กเรยฺย ปาสาทสฺส อาโรหณาย.\csromanspacer{} ตเมนํ เอวํ วเทยฺยุํ “อมฺโภ ปุริส ยสฺส ตฺวํ\footnote{ยํ ตฺวํ (สี, ก)} ปาสาทสฺส อาโรหณาย นิสฺเสณิํ กโรสิ, ชานาสิ ตํ ปาสาทํ ปุรตฺถิมาย วา ทิสาย ทกฺขิณาย วา ทิสาย ปจฺฉิมาย วา ทิสาย อุตฺตราย วา ทิสาย อุจฺโจ วา นีโจ วา มชฺฌิโม วา”ติ, อิติ ปุฏฺโฐ “โน”ติ วเทยฺย.\csromanspacer{} ตเมนํ เอวํ วเทยฺยุํ “อมฺโภ ปุริส ยํ ตฺวํ น ชานาสิ น ปสฺสติ, ตสฺส ตฺวํ ปาสาทสฺส อาโรหณาย นิสฺเสณิํ กโรสี”ติ, อิติ ปุฏฺโฐ “อามา”ติ}

\prose{ตํ กิํ มญฺญสิ โปฏฺฐปาท “นนุ เอวํ สนฺเต ตสฺส ปุริสสฺส อปฺปาฏิหีรกตํ ภาสิตํ สมฺปชฺชตี”ติ.\csromanspacer{} อทฺธา โข ภนฺเต เอวํ สนฺเต ตสฺส ปุริสสฺส อปฺปาฏิหีรกตํ ภาสิตํ สมฺปชฺชตีติ.}

\prose{\csromanfolio{180}เอวเมว โข โปฏฺฐปาท เย เต สมณพฺราหฺมณา เอวํวาทิโน เอวํทิฏฺฐิโน “เอกนฺตสุขี อตฺตา โหติ อโรโค ปรํ มรณา”ติ.\csromanspacer{} ตฺยาหํ อุปสงฺกมิตฺวา เอวํ วทามิ “สจฺจํ กิร ตุเมฺห อายสฺมานฺโต เอวํวาทิโน เอวํทิฏฺฐิโน เอกนฺตสุขี อตฺตา โหติ อโรโค ปรํ มรณา”ติ.\csromanspacer{} เต เจ เม เอวํ ปุฏฺฐา “อามา”ติ ปฏิชานนฺติ.\csromanspacer{} ตฺยาหํ เอวํ วทามิ “อปิ ปน ตุเมฺห อายสฺมนฺโต เอกนฺตสุขํ โลกํ ชานํ ปสฺสํ วิหรถา”ติ, อิติ ปุฏฺฐา “โน”ติ วทนฺติ.}

\prose{ตฺยาหํ เอวํ วทามิ “อปิ ปน ตุเมฺห อายสฺมนฺโต เอกํ วา รตฺติํ เอกํ วา ทิวสํ อุปฑฺฒํ วา รตฺติํ อุปฑฺฒํ วา ทิวสํ เอกนฺตสุขิํ อตฺตานํ สญฺชานาถา”ติ, อิติ ปุฏฺฐา “โน”ติ วทนฺติ.\csromanspacer{} ตฺยาหํ เอวํ วทามิ “อปิ ปน ตุเมฺห อายสฺมนฺโต ชานาถ อยํ มคฺโค อยํ ปฏิปทา เอกนฺตสุขสฺส โลกสฺส สจฺฉิกิริยายา”ติ, อิติ ปุฏฺฐา “โน”ติ วทนฺติ.}

\prose{ตฺยาหํ เอวํ วทามิ “อปิ ปน ตุเมฺห อายสฺมนฺโต ยา ตา เทวตา เอกนฺตสุขํ โลกํ อุปปนฺนา, ตาสํ เทวตานํ ภาสมานานํ สทฺทํ สุณาถ, สุปฺปฏิปนฺนาตฺถ มาริสา อุชุปฺปฏิปนฺนาตฺถ มาริสา เอกนฺตสุขสฺส โลกสฺส สจฺฉิกิริยาย, มยมฺปิ หิ มาริสา เอวํ ปฏิปนฺนา เอกนฺตสุขํ โลกํ อุปปนฺนา”ติ, อิติ ปุฏฺฐา “โน”ติ วทนฺติ.}

\prose{ตํ กิํ มญฺญสิ โปฏฺฐปาท “นนุ เอวํ สนฺเต เตสํ สมณพฺราหฺมณานํ อปฺปาฏิหีรกตํ ภาสิตํ สมฺปชฺชตี”ติ.\csromanspacer{} อทฺธา โข ภนฺเต เอวํ สนฺเต เตสํ สมณพฺราหฺมณานํ อปฺปาฏิหีรกตํ ภาสิตํ สมฺปชฺชตีติ.}

\csromanmatikaanchor{mtk.118}
\csromantocmarkref{section}{ตโย อตฺตปฏิลาภ}{mtk.118}
\bookmark[dest={mtk.118},level=1]{ตโย อตฺตปฏิลาภ}
\csromanheader{1}{\textbf{ตโย อตฺตปฏิลาภ}}
\proseitem{428}{ตโย โข เม โปฏฺฐปาท อตฺตปฏิลาภา โอฬาริโก อตฺตปฏิลาโภ, มโนมโย อตฺตปฏิลาโภ, อรูโป อตฺตปฏิลาโภ.\csromanspacer{} กตโม จ โปฏฺฐปาท โอฬาริโก อตฺตปฏิลาโภ.\csromanspacer{} รูปี จาตุมหาภูติโก กพฬีการาหารภกฺโข\footnote{กพฬีการภกฺโข (สฺยา, ก)}, อยํ โอฬาริโก อตฺตปฏิลาโภ.\csromanspacer{} กตโม มโนมโย อตฺตปฏิลาโภ.\csromanspacer{} รูปี มโนมโย สพฺพงฺคปจฺจงฺคี อหีนินฺทฺริโย, อยํ มโนมโย อตฺตปฏิลาโภ.\csromanspacer{} กตโม อรูโป อตฺตปฏิลาโภ.\csromanspacer{} อรูปี สญฺญามโย, อยํ อรูโป อตฺตปฏิลาโภ.}

\csromanheader{2}{9. โปฏฺฐปาทสุตฺต}
\proseitem{429}{\csromanfolio{181}โอฬาริกสฺสปิ โข อหํ โปฏฺฐปาท อตฺตปฏิลาภสฺส ปหานาย ธมฺมํ เทเสมิ “ยถาปฏิปนฺนานํ โว สํกิเลสิกา ธมฺมา ปหียิสฺสนฺติ, โวทานิยา ธมฺมา อภิวฑฺฒิสฺสนฺติ, ปญฺญาปาริปูริํ เวปุลฺลตฺตญฺจ ทิฏฺเฐว ธมฺเม สยํ อภิญฺญา สจฺฉิกตฺวา อุปสมฺปชฺช วิหริสฺสถา”ติ.\csromanspacer{} สิยา โข ปน เต โปฏฺฐปาท เอวมสฺส “สํกิเลสิกา ธมฺมา ปหียิสฺสนฺติ, โวทานิยา ธมฺมา อภิวฑฺฒิสฺสนฺติ, ปญฺญาปาริปูริํ เวปุลฺลตฺตญฺจ ทิฏฺเฐว ธมฺเม สยํ อภิญฺญา สจฺฉิกตฺวา อุปสมฺปชฺช วิหริสฺสติ, ทุกฺโข จ โข วิหาโร”ติ.\csromanspacer{} น โข ปเนตํ โปฏฺฐปาท เอวํ ทฏฺฐพฺพํ.\csromanspacer{} สํกิเลสิกา เจว ธมฺมา ปหียิสฺสนฺติ, โวทานิยา จ ธมฺมา อภิวฑฺฒิสฺสนฺติ, ปญฺญาปาริปูริํ เวปุลฺลตฺตญฺจ ทิฏฺเฐว ธมฺเม สยํ อภิญฺญา สจฺฉิกตฺวา อุปสมฺปชฺช วิหริสฺสติ, ปามุชฺชญฺเจว ภวิสฺสติ ปีติ จ ปสฺสทฺธิ จ สติ จ สมฺปชญฺญญฺจ สุโข จ วิหาโร.}

\proseitem{430}{มโนมยสฺสปิ โข อหํ โปฏฺฐปาท อตฺตปฏิลาภสฺส ปหานาย ธมฺมํ เทเสมิ “ยถาปฏิปนฺนานํ โว สํกิเลสิกา ธมฺมา ปหียิสฺสนฺติ, โวทานิยา ธมฺมา อภิวฑฺฒิสฺสนฺติ, ปญฺญาปาริปูริํ เวปุลฺลตฺตญฺจ ทิฏฺเฐว ธมฺเม สยํ อภิญฺญา สจฺฉิกตฺวา อุปสมฺปชฺช วิหริสฺสถา”ติ.\csromanspacer{} สิยา โข ปน เต โปฏฺฐปาท เอวมสฺส “สํกิเลสิกา ธมฺมา ปหียิสฺสนฺติ, โวทานิยา ธมฺมา อภิวฑฺฒิสฺสนฺติ, ปญฺญาปาริปูริํ เวปุลฺลตฺตญฺจ ทิฏฺเฐว ธมฺเม สยํ อภิญฺญา สจฺฉิกตฺวา อุปสมฺปชฺช วิหริสฺสติ, ทุกฺโข จ โข วิหาโร”ติ.\csromanspacer{} น โข ปเนตํ โปฏฺฐปาท เอวํ ทฏฺฐพฺพํ.\csromanspacer{} สํกิเลสิกา เจว ธมฺมา ปหียิสฺสนฺติ, โวทานิยา จ ธมฺมา อภิวฑฺฒิสฺสนฺติ, ปญฺญาปาริปูริํ เวปุลฺลตฺตญฺจ ทิฏฺเฐว ธมฺเม สยํ อภิญฺญา สจฺฉิกตฺวา อุปสมฺปชฺช วิหริสฺสติ, ปามุชฺชญฺเจว ภวิสฺสติ ปีติ จ ปสฺสทฺธิ จ สติ จ สมฺปชญฺญญฺจ สุโข จ วิหาโร.}

\proseitem{431}{อรูปสฺสปิ โข อหํ โปฏฺฐปาท อตฺตปฏิลาภสฺส ปหานาย ธมฺมํ เทเสมิ “ยถาปฏิปนฺนานํ โว สํกิเลสิกา ธมฺมา ปหียิสฺสนฺติ, โวทานิยา ธมฺมา อภิวฑฺฒิสฺสนฺติ, ปญฺญาปาริปูริํ เวปุลฺลตฺตญฺจ ทิฏฺเฐว ธมฺเม สยํ อภิญฺญา สจฺฉิกตฺวา อุปสมฺปชฺช วิหริสฺสถา”ติ.\csromanspacer{} สิยา โข ปน เต โปฏฺฐปาท เอวมสฺส “สํกิเลสิกา ธมฺมา ปหียิสฺสนฺติ, โวทานิยา ธมฺมา อภิวฑฺฒิสฺสนฺติ, ปญฺญาปาริปูริํ เวปุลฺลตฺตญฺจ ทิฏฺเฐว ธมฺเม สยํ อภิญฺญา สจฺฉิกตฺวา อุปสมฺปชฺช วิหริสฺสติ, ทุกฺโข จ โข วิหาโร”ติ.\csromanspacer{} น โข \csromanfolio{182}ปเนตํ โปฏฺฐปาท เอวํ ทฏฺฐพฺพํ.\csromanspacer{} สํกิเลสิกา เจว ธมฺมา ปหียิสฺสนฺติ, โวทานิยา จ ธมฺมา อภิวฑฺฒิสฺสนฺติ, ปญฺญาปาริปูริํ เวปุลฺลตฺตญฺจ ทิฏฺเฐว ธมฺเม สยํ อภิญฺญา สจฺฉิกตฺวา อุปสมฺปชฺช วิหริสฺสติ, ปามุชฺชญฺเจว ภวิสฺสติ ปีติ จ ปสฺสทฺธิ จ สติ จ สมฺปชญฺญญฺจ สุโข จ วิหาโร.}

\proseitem{432}{ปเร เจ โปฏฺฐปาท อเมฺห เอวํ ปุจฺเฉยฺยุํ “กตโม ปน โส อาวุโส โอฬาริโก อตฺตปฏิลาโภ, ยสฺส ตุเมฺห ปหานาย ธมฺมํ เทเสถ, ยถาปฏิปนฺนานํ โว สํกิเลสิกา ธมฺมา ปหียิสฺสนฺติ, โวทานิยา ธมฺมา อภิวฑฺฒิสฺสนฺติ, ปญฺญาปริปูริํ เวปุลฺลตฺตญฺจ ทิฏฺเฐว ธมฺเม สยํ อภิญฺญา สจฺฉิกตฺวา อุปสมฺปชฺช วิหริสฺสถา”ติ.\csromanspacer{} เตสํ มยํ เอวํ ปุฏฺฐา เอวํ พฺยากเรยฺยาม “อยํ วา โส อาวุโส โอฬาริโก อตฺตปฏิลาโภ, ยสฺส มยํ ปหานาย ธมฺมํ เทเสม, ยถาปฏิปนฺนานํ โว สํกิเลสิกา ธมฺมา ปหียิสฺสนฺติ, โวทานิยา ธมฺมา อภิวฑฺฒิสฺสนฺติ, ปญฺญาปาริปูริํ เวปุลฺลตฺตญฺจ ทิฏฺเฐว ธมฺเม สยํ อภิญฺญา สจฺฉิกตฺวา อุปสมฺปชฺช วิหริสฺสถา”ติ.}

\proseitem{433}{ปเร เจ โปฏฺฐปาท อเมฺห เอวํ ปุจฺเฉยฺยุํ “กตโม ปน โส อาวุโส มโนมโย อตฺตปฏิลาโภ, ยสฺส ตุเมฺห ปหานาย ธมฺมํ เทเสถ, ยถาปฏิปนฺนานํ โว สํกิเลสิกา ธมฺมา ปหียิสฺสนฺติ, โวทานิยา ธมฺมา อภิวฑฺฒิสฺสนฺติ, ปญฺญาปาริปูริํ เวปุลฺลตฺตญฺจ ทิฏฺเฐว ธมฺเม สยํ อภิญฺญา สจฺฉิกตฺวา อุปสมฺปชฺช วิหริสฺสถา”ติ.\csromanspacer{} เตสํ มยํ เอวํ ปุฏฺฐา เอวํ พฺยากเรยฺยาม “อยํ วา โส อาวุโส มโนมโย อตฺตปฏิลาโภ ยสฺส มยํ ปหานาย ธมฺมํ เทเสม, ยถาปฏิปนฺนานํ โว สํกิเลสิกา ธมฺมา ปหียิสฺสนฺติ, โวทานิยา ธมฺมา อภิวฑฺฒิสฺสนฺติ, ปญฺญาปาริปูริํ เวปุลฺลตฺตญฺจ ทิฏฺเฐว ธมฺเม สยํ อภิญฺญา สจฺฉิกตฺวา อุปสมฺปชฺช วิหริสฺสถา”ติ.}

\proseitem{434}{ปเร เจ โปฏฺฐปาท อเมฺห เอวํ ปุจฺเฉยฺยุํ “กตโม ปน โส อาวุโส อรูโป อตฺตปฏิลาโภ, ยสฺส ตุเมฺห ปหานาย ธมฺมํ เทเสถ, ยถาปฏิปนฺนานํ โว สํกิเลสิกา ธมฺมา ปหียิสฺสนฺติ, โวทานิยา ธมฺมา อภิวฑฺฒิสฺสนฺติ, ปญฺญาปาริปูริํ เวปุลฺลตฺตญฺจ ทิฏฺเฐว ธมฺเม สยํ อภิญฺญา สจฺฉิกตฺวา อุปสมฺปชฺช วิหริสฺสถา”ติ.\csromanspacer{} เตสํ มยํ เอวํ}

\vspace{\tipitakaparskip}
\csromancenter{โปฏฺฐปาทสุตฺต ปุฏฺฐา เอวํ พฺยากเรยฺยาม “อยํ วา โส อาวุโส อรูโป อตฺตปฏิลาโภ ยสฺส มยํ ปหานาย ธมฺมํ เทเสม, ยถาปฏิปนฺนานํ โว สํกิเลสิกา ธมฺมา ปหียิสฺสนฺติ, โวทานิยา ธมฺมา อภิวฑฺฒิสฺสนฺติ, ปญฺญาปาริปูริํ เวปุลฺลตฺตญฺจ ทิฏฺเฐว ธมฺเม สยํ อภิญฺญา สจฺฉิกตฺวา อุปสมฺปชฺช วิหริสฺสถา”ติ.}
\prose{\csromanfolio{183}ตํ กิํ มญฺญสิ โปฏฺฐปาท “นนุ เอวํ สนฺเต สปฺปาฏิหีรกตํ ภาสิตํ สมฺปชฺชตี”ติ.\csromanspacer{} อทฺธา โข ภนฺเต เอวํ สนฺเต สปฺปาฏิหีรกตํ ภาสิตํ สมฺปชฺชตีติ.}

\proseitem{435}{เสยฺยถาปิ โปฏฺฐปาท ปุริโส นิสฺเสณิํ กเรยฺย ปาสาทสฺส อาโรหณาย ตสฺเสว ปาสาทสฺส เหฏฺฐา.\csromanspacer{} ตเมนํ เอวํ วเทยฺยุํ “อมฺโภ ปุริส ยสฺส ตฺวํ ปาสาทสฺส อาโรหณาย นิสฺเสณิํ กโรสิ, ชานาสิ ตํ ปาสาทํ ปุรตฺถิมาย วา ทิสาย ทกฺขิณาย วา ทิสาย ปจฺฉิมาย วา ทิสาย อุตฺตราย วา ทิสาย อุจฺโจ วา นีโจ วา มชฺฌิโม วา”ติ.\csromanspacer{} โส เอวํ วเทยฺย “อยํ วา โส อาวุโส ปาสาโท, ยสฺสาหํ อาโรหณาย นิสฺเสณิํ กโรมิ ตสฺเสว ปาสาทสฺส เหฏฺฐา”ติ.}

\prose{ตํ กิํ มญฺญสิ โปฏฺฐปาท “นนุ เอวํ สนฺเต ตสฺส ปุริสสฺส สปฺปาฏิหีรกตํ ภาสิตํ สมฺปชฺชตี”ติ.\csromanspacer{} อทฺธา โข ภนฺเต เอวํ สนฺเต ตสฺส ปุริสสฺส สปฺปาฏิหีรกตํ ภาสิตํ สมฺปชฺชตีติ.}

\proseitem{436}{เอวเมว โข โปฏฺฐปาท ปเร เจ อเมฺห เอวํ ปุจฺเฉยฺยุํ “กตโม ปน โส อาวุโส โอฬาริโก อตฺตปฏิลาโภ ฯเปฯ.\csromanspacer{} กตโม ปน โส อาวุโส มโนมโย อตฺตปฏิลาโภ ฯเปฯ.\csromanspacer{} กตโม ปน โส อาวุโส อรูโป อตฺตปฏิลาโภ, ยสฺส ตุเมฺห ปหานาย ธมฺมํ เทเสถ, ยถาปฏิปนฺนานํ โว สํกิเลสิกา ธมฺมา ปหียิสฺสนฺติ, โวทานิยา ธมฺมา อภิวฑฺฒิสฺสนฺติ, ปญฺญาปาริปูริํ เวปุลฺลตฺตญฺจ ทิฏฺเฐว ธมฺเม สยํ อภิญฺญา สจฺฉิกตฺวา อุปสมฺปชฺช วิหริสฺสถา”ติ.\csromanspacer{} เตสํ มยํ เอวํ ปุฏฺฐา เอวํ พฺยากเรยฺยาม “อยํ วา โส อาวุโส อรูโป อตฺตปฏิลาโภ, ยสฺส มยํ ปหานาย ธมฺมํ เทเสม, ยถาปฏิปนฺนานํ โว สํกิเลสิกา ธมฺมา ปหียิสฺสนฺติ, โวทานิยา ธมฺมา อภิวฑฺฒิสฺสนฺติ, \csromanfolio{184}ปญฺญาปาริปูริํ เวปุลฺลตฺตญฺจ ทิฏฺเฐว ธมฺเม สยํ อภิญฺญา สจฺฉิกตฺวา อุปสมฺปชฺช วิหริสฺสถา”ติ.}

\prose{ตํ กิํ มญฺญสิ โปฏฺฐปาท “นนุ เอวํ สนฺเต สปฺปาฏิหีรกตํ ภาสิตํ สมฺปชฺชตี”ติ.\csromanspacer{} อทฺธา โข ภนฺเต เอวํ สนฺเต สปฺปาฏิหีรกตํ ภาสิตํ สมฺปชฺชตีติ.}

\proseitem{437}{เอวํ วุตฺเต จิตฺโต หตฺถิสาริปุตฺโต ภควนฺตํ เอตทโวจ “ยสฺมิํ ภนฺเต สมเย โอฬาริโก อตฺตปฏิลาโภ โหติ, โมฆสฺส ตสฺมิํ สมเย มโนมโย อตฺตปฏิลาโภ โหติ, โมโฆ อรูโป อตฺตปฏิลาโภ โหติ, โอฬาริโก วาสฺส อตฺตปฏิลาโภ ตสฺมิํ สมเย สจฺโจ โหติ.\csromanspacer{} ยสฺมิํ ภนฺเต สมเย มโนมโย อตฺตปฏิลาโภ โหติ, โมฆสฺส ตสฺมิํ สมเย โอฬาริโก อตฺตปฏิลาโภ โหติ, โมโฆ อรูโป อตฺตปฏิลาโภ โหติ, มโนมโย วาสฺส อตฺตปฏิลาโภ ตสฺมิํ สมเย สจฺโจ โหติ.\csromanspacer{} ยสฺมิํ ภนฺเต สมเย อรูโป อตฺตปฏิลาโภ โหติ, โมฆสฺส ตสฺมิํ สมเย โอฬาริโก อตฺตปฏิลาโภ โหติ, โมโฆ มโนมโย อตฺตปฏิลาโภ โหติ, อรูโป วาสฺส อตฺตปฏิลาโภ ตสฺมิํ สมเย สจฺโจ โหตี”ติ.}

\prose{ยสฺมิํ จิตฺต สมเย โอฬาริโก อตฺตปฏิลาโภ โหติ, เนว ตสฺมิํ สมเย มโนมโย อตฺตปฏิลาโภติ สงฺขํ คจฺฉติ, น อรูโป อตฺตปฏิลาโภติ สงฺขํ คจฺฉติ, โอฬาริโก อตฺตปฏิลาโภเตฺวว ตสฺมิํ สมเย สงฺขํ คจฺฉติ.\csromanspacer{} ยสฺมิํ จิตฺต สมเย มโนมโย อตฺตปฏิลาโภ โหติ, เนว ตสฺมิํ สมเย โอฬาริโก อตฺตปฏิลาโภติ สงฺขํ คจฺฉติ, น อรูโป อตฺตปฏิลาโภติ สงฺขํ คจฺฉติ, มโนมโย อตฺตปฏิลาโภเตฺวว ตสฺมิํ สมเย สงฺขํ คจฺฉติ.\csromanspacer{} ยสฺมิํ จิตฺต สมเย อรูโป อตฺตปฏิลาโภ โหติ, เนว ตสฺมิํ สมเย โอฬาริโก อตฺตปฏิลาโภติ สงฺขํ คจฺฉติ, น มโนมโย อตฺตปฏิลาโภติ สงฺขํ คจฺฉติ, อรูโป อตฺตปฏิลาโภเตฺวว ตสฺมิํ สมเย สงฺขํ คจฺฉติ.}

\proseitem{438}{สเจ ตํ จิตฺต เอวํ ปุจฺเฉยฺยุํ “อโหสิ ตฺวํ อตีตมทฺธานํ, น ตฺวํ นาโหสิ, ภวิสฺสสิ ตฺวํ อนาคตมทฺธานํ, น ตฺวํ น ภวิสฺสสิ, อตฺถิ}

\vspace{\tipitakaparskip}
\csromancenter{โปฏฺฐปาทสุตฺต ตฺวํ เอตรหิ, น ตฺวํ นตฺถี”ติ, เอวํ ปุฏฺโฐ ตฺวํ จิตฺต กินฺติ พฺยากเรยฺยาสีติ.\csromanspacer{} สเจ มํ ภนฺเต เอวํ ปุจฺเฉยฺยุํ “อโหสิ ตฺวํ อตีตมทฺธานํ, น ตฺวํ น อโหสิ, ภวิสฺสสิ ตฺวํ อนาคตมทฺธานํ, น ตฺวํ น ภวิสฺสสิ, อตฺถิ ตฺวํ เอตรหิ, น ตฺวํ นตฺถี”ติ.\csromanspacer{} เอวํ ปุฏฺโฐ อหํ ภนฺเต เอวํ พฺยากเรยฺยํ “อโหสาหํ อตีตมทฺธานํ, นาหํ น อโหสิํ, ภวิสฺสามหํ อนาคตมทฺธานํ, นาหํ น ภวิสฺสามิ, อตฺถาหํ เอตรหิ, นาหํ นตฺถี”ติ.\csromanspacer{} เอวํ ปุฏฺโฐ อหํ ภนฺเต เอวํ พฺยากเรยฺยนฺติ.}
\prose{\csromanfolio{185}สเจ ปน ตํ จิตฺเต เอวํ ปุจฺเฉยฺยุํ “โย เต อโหสิ อตีโต อตฺตปฏิลาโภ, โสว\footnote{เสฺวว (สี, อิ), โสเยว (สฺยา)} เต อตฺตปฏิลาโภ สจฺโจ, โมโฆ อนาคโต, โมโฆ ปจฺจุปฺปนฺโน.\csromanspacer{} โย\footnote{โย วา (อิ)} เต ภวิสฺสติ อนาคโต อตฺตปฏิลาโภ, โสว เต อตฺตปฏิลาโภ สจฺโจ, โมโฆ อตีโต, โมโฆ ปจฺจุปฺปนฺโน.\csromanspacer{} โย2 เต เอตรหิ ปจฺจุปฺปนฺโน อตฺตปฏิลาโภ, โสว\footnote{โส จ (ก)} เต อตฺตปฏิลาโภ สจฺโจ, โมโฆ อตีโต, โมโฆ อนาคโต”ติ.\csromanspacer{} เอวํ ปุฏฺโฐ ตฺวํ จิตฺต กินฺติ พฺยากเรยฺยาสีติ.\csromanspacer{} สเจ ปน มํ ภนฺเต เอวํ ปุจฺเฉยฺยุํ “โย เต อโหสิ อตีโต อตฺตปฏิลาโภ, โสว เต อตฺตปฏิลาโภ สจฺโจ, โมโฆ อนาคโต, โมโฆ ปจฺจุปฺปนฺโน.\csromanspacer{} โย เต ภวิสฺสติ อนาคโต อตฺตปฏิลาโภ, โสว เต อตฺตปฏิลาโภ สจฺโจ, โมโฆ อตีโต, โมโฆ ปจฺจุปฺปนฺโน.\csromanspacer{} โย เต เอตรหิ ปจฺจุปฺปนฺโน อตฺตปฏิลาโภ, โสว เต อตฺตปฏิลาโภ สจฺโจ, โมโฆ อตีโต, โมโฆ อนาคโต”ติ.\csromanspacer{} เอวํ ปุฏฺโฐ อหํ ภนฺเต เอวํ พฺยากเรยฺยํ “โย เม อโหสิ อตีโต อตฺตปฏิลาโภ, โสว เม อตฺตปฏิลาโภ ตสฺมิํ สมเย สจฺโจ อโหสิ, โมโฆ อนาคโต, โมโฆ ปจฺจุปฺปนฺโน.\csromanspacer{} โย เม ภวิสฺสติ อนาคโต อตฺตปฏิลาโภ, โสว เม อตฺตปฏิลาโภ ตสฺมิํ สมเย สจฺโจ ภวิสฺสติ, โมโฆ อตีโต, โมโฆ ปจฺจุปฺปนฺโน.\csromanspacer{} โย เม เอตรหิ ปจฺจุปฺปนฺโน อตฺตปฏิลาโภ, โสว เม อตฺตปฏิลาโภ สจฺโจ, โมโฆ อตีโต, โมโฆ อนาคโต”ติ.\csromanspacer{} เอวํ ปุฏฺโฐ อหํ ภนฺเต เอวํ พฺยากเรยฺยนฺติ.}

\proseitem{439}{\csromanfolio{186}เอวเมว โข จิตฺต ยสฺมิํ สมเย โอฬาริโก อตฺตปฏิลาโภ โหติ, เนว ตสฺมิํ สมเย มโนมโย อตฺตปฏิลาโภติ สงฺขํ คจฺฉติ, น อรูโป อตฺตปฏิลาโภติ สงฺขํ คจฺฉติ, โอฬาริโก อตฺตปฏิลาโภเตฺวว ตสฺมิํ สมเย สงฺขํ คจฺฉติ.\csromanspacer{} ยสฺมิํ จิตฺต สมเย มโนมโย อตฺตปฏิลาโภ โหติ ฯเปฯ.\csromanspacer{} ยสฺมิํ จิตฺต สมเย อรูโป อตฺตปฏิลาโภ โหติ, เนว ตสฺมิํ สมเย โอฬาริโก อตฺตปฏิลาโภติ สงฺขํ คจฺฉติ, น มโนมโย อตฺตปฏิลาโภติ สงฺขํ คจฺฉติ, อรูโป อตฺตปฏิลาโภเตฺวว ตสฺมิํ สมเย สงฺขํ คจฺฉติ.}

\proseitem{440}{เสยฺยถาปิ จิตฺต ควา ขีรํ, ขีรมฺหา ทธิ, ทธิมฺหา นวนีตํ, นวนีตมฺหา สปฺปิ, สปฺปิมฺหา สปฺปิมณฺโฑ.\csromanspacer{} ยสฺมิํ สมเย ขีรํ โหติ, เนว ตสฺมิํ สมเย ทธีติ สงฺขํ คจฺฉติ, น นวนีตนฺติ สงฺขํ คจฺฉติ, น สปฺปีติ สงฺขํ คจฺฉติ, น สปฺปิมณฺโฑติ สงฺขํ คจฺฉติ, ขีรํเตฺวว ตสฺมิํ สมเย สงฺขํ คจฺฉติ.\csromanspacer{} ยสฺมิํ สมเย ทธิ โหติ ฯเปฯ นวนีตํ โหติ.\csromanspacer{} สปฺปิ โหติ.\csromanspacer{} สปฺปิมณฺโฑ โหติ, เนว ตสฺมิํ สมเย ขีรนฺติ สงฺขํ คจฺฉติ, น ทธีติ สงฺขํ คจฺฉติ, น นวนีตนฺติ สงฺขํ คจฺฉติ, น สปฺปีติ สงฺขํ คจฺฉติ, สปฺปิมณฺโฑเตฺวว ตสฺมิํ สมเย สงฺขํ คจฺฉติ.\csromanspacer{} เอวเมว โข จิตฺต ยสฺมิํ สมเย โอฬาริโก อตฺตปฏิลาโภ โหติ ฯเปฯ.\csromanspacer{} ยสฺมิํ จิตฺต สมเย มโนมโย อตฺตปฏิลาโภ โหติ ฯเปฯ.\csromanspacer{} ยสฺมิํ จิตฺต สมเย อรูโป อตฺตปฏิลาโภ โหติ, เนว ตสฺมิํ สมเย โอฬาริโก อตฺตปฏิลาโภติ สงฺขํ คจฺฉติ, น มโนมโย อตฺตปฏิลาโภติ สงฺขํ คจฺฉติ, อรูโป อตฺตปฏิลาโภเตฺวว ตสฺมิํ สมเย สงฺขํ คจฺฉติ.\csromanspacer{} อิมา โข จิตฺต โลกสมญฺญา โลกนิรุตฺติโย โลกโวหารา โลกปญฺญตฺติโย, ยาหิ ตถาคโต โวหรติ อปรามสนฺติ.}

\proseitem{441}{เอวํ วุตฺเต โปฏฺฐปาโท ปริพฺพาชโก ภควนฺตํ เอตทโวจ “อภิกฺกนฺตํ ภนฺเต อภิกฺกนฺตํ ภนฺเต, เสยฺยถาปิ ภนฺเต นิกฺกุชฺชิตํ วา อุกฺกุชฺเชยฺย, ปฏิจฺฉนฺนํ วา วิวเรยฺย, มูฬฺหสฺส วา มคฺคํ อาจิกฺเขยฺย, อนฺธกาเร วา เตลปชฺโชตํ ธาเรยฺย ‘จกฺขุมนฺโต รูปานิ ทกฺขนฺตี’ติ, เอวเมวํ ภควตา อเนกปริยาเยน ธมฺโม ปกาสิโต, เอสาหํ ภนฺเต ภควนฺตํ สรณํ คจฺฉามิ ธมฺมญฺจ ภิกฺขุสํฆญฺจ, อุปาสกํ มํ ภควา ธาเรตุ อชฺชตคฺเค ปาณุเปตํ สรณํ คตนฺ”ติ.}

\csromanmatikaanchor{mtk.119}
\csromantocmarkref{section}{จิตฺตหตฺถิสาริปุตฺตอุปสมฺปทา}{mtk.119}
\bookmark[dest={mtk.119},level=1]{จิตฺตหตฺถิสาริปุตฺตอุปสมฺปทา}
\csromanheader{1}{9. โปฏฺฐปาทสุตฺต \textbf{จิตฺตหตฺถิสาริปุตฺตอุปสมฺปทา}}
\proseitem{442}{\csromanfolio{187}จิตฺโต ปน หตฺถิสาริปุตฺโต ภควนฺตํ เอตทโวจ “อภิกฺกนฺตํ ภนฺเต อภิกฺกนฺตํ ภนฺเต, เสยฺยถาปิ ภนฺเต นิกฺกุชฺชิตํ วา อุกฺกุชฺเชยฺย, ปฏิจฺฉนฺนํ วา วิวเรยฺย, มูฬฺหสฺส วา มคฺคํ อาจิกฺเขยฺย, อนฺธกาเร วา เตลปชฺโชตํ ธาเรยฺย ‘จกฺขุมนฺโต รูปานิ ทกฺขนฺตี’ติ, เอวเมวํ ภควตา อเนกปริยาเยน ธมฺโม ปกาสิโต, เอสาหํ ภนฺเต ภควนฺตํ สรณํ คจฺฉามิ ธมฺมญฺจ ภิกฺขุสํฆญฺจ, ลเภยฺยาหํ ภนฺเต ภควโต สนฺติเก ปพฺพชฺชํ ลเภยฺยํ อุปสมฺปทนฺ”ติ.}

\proseitemwithcloser{443}{อลตฺถ โข จิตฺโต หตฺถิสาริปุตฺโต ภควโต สนฺติเก ปพฺพชฺชํ อลตฺถ อุปสมฺปทํ.\csromanspacer{} อจิรูปสมฺปนฺโน โข ปนายสฺมา จิตฺโต หตฺถิสาริปุตฺโต เอโก วูปกฏฺโฐ อปฺปมตฺโต อาตาปี ปหิตตฺโต วิหรนฺโต น จิรสฺเสว ยสฺสตฺถาย กุลปุตฺตา สมฺมเทว อคารสฺมา อนคาริยํ ปพฺพชนฺติ, ตทนุตฺตรํ พฺรหฺมจริยปริโยสานํ ทิฏฺเฐว ธมฺเม สยํ อภิญฺญา สจฺฉิกตฺวา อุปสมฺปชฺช วิหาสิ, “ขีณา ชาติ, วุสิตํ พฺรหฺมจริยํ, กตํ กรณียํ, นาปรํ อิตฺถตฺตายา”ติ อพฺภญฺญาสิ.\csromanspacer{} อญฺญตโร โข ปนายสฺมา จิตฺโต หตฺถิสาริปุตฺโต อรหตํ อโหสีติ.}{โปฏฺฐปาทสุตฺตํ นิฏฺฐิตํ นวมํ.}
\csromanmatikaanchor{mtk.120}
\csromantocmarknopage{chapter}{10. สุภสุตฺต}{mtk.120}
\bookmark[dest={mtk.120},level=0]{10. สุภสุตฺต}
\setcsromanheadcha{10. สุภสุตฺต}
\setcsromanheadhclear{}
\chapterheadpageread{10. \textbf{สุภสุตฺต}}
\csromanmatikaanchor{mtk.121}
\csromantocmarkref{section}{สุภมาณววตฺถุ}{mtk.121}
\bookmark[dest={mtk.121},level=1]{สุภมาณววตฺถุ}
\csromanheader{1}{\textbf{สุภมาณววตฺถุ}}
\proseitem{444}{\csromanfolio{188}เอวํ เม สุตํ\csromandash{}เอกํ สมยํ อายสฺมา อานนฺโท สาวตฺถิยํ วิหรติ เชตวเน อนาถปิณฺฑิกสฺส อาราเม อจิรปรินิพฺพุเต ภควติ.\csromanspacer{} เตน โข ปน สมเยน สุโภ มาณโว โตเทยฺยปุตฺโต สาวตฺถิยํ ปฏิวสติ เกนจิเทว กรณีเยน.}

\proseitem{445}{อถ โข สุโภ มาณโว โตเทยฺยปุตฺโต อญฺญตรํ มาณวกํ อามนฺเตสิ\csromandash{}เอหิ ตฺวํ มาณวก เยน สมโณ อานนฺโท เตนุปสงฺกม, อุปสงฺกมิตฺวา มม วจเนน สมณํ อานนฺทํ อปฺปาพาธํ อปฺปาตงฺกํ ลหุฏฺฐานํ พลํ ผาสุวิหารํ ปุจฺฉ “สุโภ มาณโว โตเทยฺยปุตฺโต ภวนฺตํ อานนฺทํ อปฺปาพาธํ อปฺปาตงฺกํ ลหุฏฺฐานํ พลํ ผาสุวิหารํ ปุจฺฉตี”ติ, เอวญฺจ วเทหิ “สาธุ กิร ภวํ อานนฺโท เยน สุภสฺส มาณวสฺส โตเทยฺยปุตฺตสฺส นิเวสนํ เตนุปสงฺกมตุ อนุกมฺปํ อุปาทายา”ติ.}

\proseitem{446}{“เอวํ โภ”ติ โข โส มาณวโก สุภสฺส มาณวสฺส โตเทยฺยปุตฺตสฺส ปฏิสฺสุตฺวา เยนายสฺมา อานนฺโท เตนุปสงฺกมิ, อุปสงฺกมิตฺวา อายสฺมตา อานนฺเทน สทฺธิํ สมฺโมทิ, สมฺโมทนียํ กถํ สารณียํ วีติสาเรตฺวา เอกมนฺตํ นิสีทิ, เอกมนฺตํ นิสินฺโน โข โส มาณวโก อายสฺมนฺตํ อานนฺทํ เอตทโวจ “สุโภ มาณโว โตเทยฺยปุตฺโต ภวนฺตํ อานนฺทํ อปฺปาพาธํ อปฺปาตงฺกํ ลหุฏฺฐานํ พลํ ผาสุวิหารํ ปุจฺฉติ, เอวญฺจ วเทติ สาธุ กิร ภวํ อานนฺโท เยน สุภสฺส มาณวสฺส โตเทยฺยปุตฺตสฺส นิเวสนํ เตนุปสงฺกมตุ อนุกมฺปํ อุปาทายา”ติ.}

\proseitem{447}{เอวํ วุตฺเต อายสฺมา อานนฺโท ตํ มาณวกํ เอตทโวจ “อกาโล โข มาณวก, อตฺถิ เม อชฺช เภสชฺชมตฺตา ปีตา, อปฺเปวนาม เสฺวปิ อุปสงฺกเมยฺยาม กาลญฺจ สมยญฺจ อุปาทายา”ติ. “เอวํ โภ”ติ โข โส มาณวโก อายสฺมโต อานนฺทสฺส ปฏิสฺสุตฺวา อุฏฺฐายาสนา เยน สุโภ มาณโว โตเทยฺยปุตฺโต}

\vspace{\tipitakaparskip}
\csromancenter{สุภสุตฺต เตนุปสงฺกมิ, อุปสงฺกมิตฺวา สุภํ มาณวํ โตเทยฺยปุตฺตํ เอตทโวจ\csromandash{}อโวจุมฺหา โข มยํ โภโต วจเนน ตํ ภวนฺตํ อานนฺทํ “สุโภ มาณโว โตเทยฺยปุตฺโต ภวนฺตํ อานนฺทํ อปฺปาพาธํ อปฺปาตงฺกํ ลหุฏฺฐานํ พลํ ผาสุวิหารํ ปุจฺฉติ, เอวญฺจ วเทติ ‘สาธุ กิร ภวํ อานนฺโท เยน สุภสฺส มาณวสฺส โตเทยฺยปุตฺตสฺส นิเวสนํ เตนุปสงฺกมตุ อนุกมฺปํ อุปาทายา’ติ.\csromanspacer{} เอวํ วุตฺเต โภ สมโณ อานนฺโท มํ เอตทโวจ “อกาโล โข มาณวก, อตฺถิ เม อชฺช เภสชฺชมตฺตา ปีตา, อปฺเปวนาม เสฺวปิ อุปสงฺกเมยฺยาม กาลญฺจ สมยญฺจ อุปาทายา”ติ.\csromanspacer{} เอตฺตาวตาปิ โข โภ กตเมว เอตํ, ยโต โข โส ภวํ อานนฺโท โอกาสมกาสิ สฺวาตนายปิ อุปสงฺกมนายาติ.}
\proseitem{448}{\csromanfolio{189}อถ โข อายสฺมา อานนฺโท ตสฺสา รตฺติยา อจฺจเยน ปุพฺพณฺหสมยํ นิวาเสตฺวา ปตฺตจีวรมาทาย เจตเกน ภิกฺขุนา ปจฺฉาสมเณน เยน สุภสฺส มาณวสฺส โตเทยฺยปุตฺตสฺส นิเวสนํ เตนุปสงฺกมิ, อุปสงฺกมิตฺวา ปญฺญตฺเต อาสเน นิสีทิ.}

\prose{อถ โข สุโภ มาณโว โตเทยฺยปุตฺโต เยนายสฺมา อานนฺโท เตนุปสงฺกมิ, อุปสงฺกมิตฺวา อายสฺมตา อานนฺเทน สทฺธิํ สมฺโมทิ, สมฺโมทนียํ กถํ สารณียํ วีติสาเรตฺวา เอกมนฺตํ นิสีทิ, เอกมนฺตํ นิสินฺโน โข สุโภ มาณโว โตเทยฺยปุตฺโต อายสฺมนฺตํ อานนฺทํ เอตทโวจ “ภวํ หิ อานนฺโท ตสฺส โภโต โคตมสฺส ทีฆรตฺตํ อุปฏฺฐาโก สนฺติกาวจโร สมีปจารี, ภวเมตํ อานนฺโท ชาเนยฺย, เยสํ โส ภวํ โคตโม ธมฺมานํ วณฺณวาที อโหสิ, ยตฺถ จ อิมํ ชนตํ สมาทเปสิ นิเวเสสิ ปติฏฺฐาเปสิ, กตเมสานํ โข โภ อานนฺท ธมฺมานํ โส ภวํ โคตโม วณฺณวาที อโหสิ, กตฺถ จ อิมํ ชนตํ สมาทเปสิ นิเวเสสิ ปติฏฺฐาเปสี”ติ.}

\proseitem{449}{ติณฺณํ โข มาณว ขนฺธานํ โส ภควา วณฺณวาที อโหสิ, เอตฺถ จ อิมํ ชนตํ สมาทเปสิ นิเวเสสิ ปติฏฺฐาเปสิ.\csromanspacer{} กตเมสํ ติณฺณํ, อริยสฺส สีลกฺขนฺธสฺส อริยสฺส สมาธิกฺขนฺธสฺส อริยสฺส ปญฺญากฺขนฺธสฺส, อิเมสํ โข มาณว ติณฺณํ ขนฺธานํ โส ภควา วณฺณวาที อโหสิ, เอตฺถ จ อิมํ ชนตํ สมาทเปสิ นิเวเสสิ ปติฏฺฐาเปสีติ.}

\csromanheader{2}{สีลกฺขนฺธ}
\csromanmatikaanchor{mtk.122}
\csromantocmarkref{section}{สิลกฺขนฺธ}{mtk.122}
\bookmark[dest={mtk.122},level=1]{สิลกฺขนฺธ}
\proseitem{450}{\csromanfolio{190}กตโม ปน โส โภ อานนฺท อริโย สีลกฺขนฺโธ, ยสฺส โส ภวํ โคตโม วณฺณวาที อโหสิ, ยตฺถ จ อิมํ ชนตํ สมาทเปสิ นิเวเสสิ ปติฏฺฐาเปสีติ.}

\prose{อิธ มาณว ตถาคโต โลเก อุปฺปชฺชติ อรหํ สมฺมาสมฺพุทฺโธ วิชฺชาจรณสมฺปนฺโน สุคโต โลกวิทู อนุตฺตโร ปุริสทมฺมสารถิ สตฺถา เทวมนุสฺสานํ พุทฺโธ ภควา, โส อิมํ โลกํ สเทวกํ สมารกํ สพฺรหฺมกํ สสฺสมณพฺราหฺมณิํ ปชํ สเทวมนุสฺสํ สยํ อภิญฺญา สจฺฉิกตฺวา ปเวเทติ, โส ธมฺมํ เทเสติ อาทิกลฺยาณํ มชฺเฌกลฺยาณํ ปริโยสานกลฺยาณํ สาตฺถํ สพฺยญฺชนํ เกวลปริปุณฺณํ ปริสุทฺธํ พฺรหฺมจริยํ ปกาเสติ, ตํ ธมฺมํ สุณาติ คหปติ วา คหปติปุตฺโต วา อญฺญตรสฺมิํ วา กุเล ปจฺจาชาโต, โส ตํ ธมฺมํ สุตฺวา ตถาคเต สทฺธํ ปฏิลภติ, โส เตน สทฺธาปฏิลาเภน สมนฺนาคโต อิติ ปฏิสญฺจิกฺขติ “สมฺพาโธ ฆราวาโส รโชปโถ, อพฺโภกาโส ปพฺพชฺชา, นยิทํ สุกรํ อคารํ อชฺฌาวสตา เอกนฺตปริปุณฺณํ เอกนฺตปริสุทฺธํ สงฺขลิขิตํ พฺรหฺมจริยํ จริตุํ, ยํนูนาหํ เกสมสฺสุํ โอหาเรตฺวา กาสายานิ วตฺถานิ อจฺฉาเทตฺวา อคารสฺมา อนคาริยํ ปพฺพเชยฺยนฺ”ติ, โส อปเรน สมเยน อปฺปํ วา โภคกฺขนฺธํ ปหาย มหนฺตํ วา โภคกฺขนฺธํ ปหาย อปฺปํ วา ญาติปริวฏฺฏํ ปหาย มหนฺตํ วา ญาติปริวฏฺฏํ ปหาย เกสมสฺสุํ โอหาเรตฺวา กาสายานิ วตฺถานิ อจฺฉาเทตฺวา อคารสฺมา อนคาริยํ ปพฺพชติ, โส เอวํ ปพฺพชิโต สมาโน ปาติโมกฺขสํวรสํวุโต วิหรติ อาจารโคจรสมฺปนฺโน อนุมตฺเตสุ วชฺเชสุ ภยทสฺสาวี สมาทาย สิกฺขติ สิกฺขาปเทสุ, กายกมฺมวจีกมฺเมน สมนฺนาคโต กุสเลน ปริสุทฺธาชีโว, สีลสมฺปนฺโน อินฺทฺริเยสุ คุตฺตทฺวาโร สติสมฺปชญฺเญน สมนฺนาคโต สนฺตุฏฺโฐ.}

\proseitem{451}{กถญฺจ มาณว ภิกฺขุ สีลสมฺปนฺโน โหติ.\csromanspacer{} อิธ มาณว ภิกฺขุ ปาณาติปาตํ ปหาย ปาณาติปาตา ปฏิวิรโต โหติ, นิหิตทณฺโฑ นิหิตสตฺโถ ลชฺชี ทยาปนฺโน, สพฺพปาณภูต\-หิตา\-นุกมฺปี วิหรติ.\csromanspacer{} ยมฺปิ มาณว ภิกฺขุ ปาณาติปาตํ ปหาย ปาณาติปาตา ปฏิวิรโต โหติ, นิหิตทณฺโฑ นิหิตสตฺโถ ลชฺชี ทยาปนฺโน,}

\vspace{\tipitakaparskip}
\csromancenter{สุภสุตฺต สพฺพปาณภูต\-หิตา\-นุกมฺปี วิหรติ, อิทมฺปิสฺส โหติ สีลสฺมิํ. (ตโต ปรํ สพฺพํ วิตฺถาเรตพฺพํ.)}
\prose{\csromanfolio{191}ยถา วา ปเนเก โภนฺโต สมณพฺราหฺมณา สทฺธาเทยฺยานิ โภชนานิ ภุญฺชิตฺวา เต เอวรูปาย ติรจฺฉานวิชฺชาย มิจฺฉาชีเวน ชีวิตํ กปฺเปนฺติ, เสยฺยถิทํ, สนฺติกมฺมํ ปณิธิกมฺมํ ภูตกมฺมํ ภูริกมฺมํ วสฺสกมฺมํ โวสฺสกมฺมํ วตฺถุกมฺมํ วตฺถุปริกมฺมํ อาจมนํ นฺหาปนํ ชุหนํ วมนํ วิเรจนํ อุทฺธํวิเรจนํ อโธวิเรจนํ สีสวิเรจนํ กณฺณเตลํ เนตฺตตปฺปนํ นตฺถุกมฺมํ อญฺชนํ ปจฺจญฺชนํ สาลากิยํ สลฺลกตฺติยํ ทารกติกิจฺฉา มูลเภสชฺชานํ อนุปฺปทานํ โอสธีนํ ปฏิโมกฺโข อิติ วา, อิติ เอวรูปาย ติรจฺฉานวิชฺชาย มิจฺฉาชีวา ปฏิวิรโต โหติ.\csromanspacer{} ยมฺปิ มาณว ภิกฺขุ ยถา วา ปเนเก โภนฺโต สมณพฺราหฺมณา สทฺธาเทยฺยานิ โภชนานิ ภุญฺชิตฺวา เต เอวรูปาย ติรจฺฉานวิชฺชาย มิจฺฉาชีเวน ชีวิตํ กปฺเปนฺติ, เสยฺยถิทํ, สนฺติกมฺมํ ปณิธิกมฺมํ ฯเปฯ โอสธีนํ ปฏิโมกฺโข อิติ วา, อิติ เอวรูปาย ติรจฺฉานวิชฺชาย มิจฺฉาชีวา ปฏิวิรโต โหติ, อิทมฺปิสฺส โหติ สีลสฺมิํ.}

\proseitem{452}{ส โขโส\footnote{อยํ โข โส (ก)} มาณว ภิกฺขุ เอวํ สีลสมฺปนฺโน น กุโตจิ ภยํ สมนุปสฺสติ, ยทิทํ สีลสํวรโต.\csromanspacer{} เสยฺยถาวิ มาณว ราชา ขตฺติโย มุทฺธาวสิตฺโต นิหตปจฺจามิตฺโต น กุโตจิ ภยํ สมนุปสฺสติ, ยทิทํ ปจฺจตฺถิกโต.\csromanspacer{} เอวเมว โข มาณว ภิกฺขุ เอวํ สีลสมฺปนฺโน น กุโตจิ ภยํ สมนุปสฺสติ, ยทิทํ สีลสํวรโต.\csromanspacer{} โส อิมินา อริเยน สีลกฺขนฺเธน สมนฺนาคโต อชฺฌตฺตํ อนวชฺชสุขํ ปฏิสํเวเทติ.\csromanspacer{} เอวํ โข มาณว ภิกฺขุ สีลสมฺปนฺโน โหติ.}

\proseitem{453}{อยํ โข โส มาณว อริโย สีลกฺขนฺโธ, ยสฺส โส ภควา วณฺณวาที อโหสิ, ยตฺถ จ อิมํ ชนตํ สมาทเปสิ นิเวเสสิ ปติฏฺฐาเปสิ.\csromanspacer{} อตฺถิ เจเวตฺถ อุตฺตริกรณียนฺติ.}

\prose{อจฺฉริยํ โภ อานนฺท อพฺภุตํ โภ อานนฺท, โส จายํ โภ อานนฺท อริโย สีลกฺขนฺโธ ปริปุณฺโณ, โน อปริปุณฺโณ.\csromanspacer{} เอวํ ปริปุณฺณญฺจาหํ โภ อานนฺท อริยํ สีลกฺขนฺธํ อิโต พหิทฺธา อญฺเญสุ สมณพฺราหฺมเณสุ \csromanfolio{192}น สมนุปสฺสามิ, เอวํ ปริปุณฺณญฺจ โภ อานนฺท อริยํ สีลกฺขนฺธํ อิโต พหิทฺธา อญฺเญ สมณพฺราหฺมณา อตฺตนิ สมนุปสฺเสยฺยุํ, เต ตาวตเกเนว อตฺตมนา อสฺสุ “อลเมตฺตาวตา กตเมตฺตาวตา, อนุปฺปตฺโต โน สามญฺญตฺโถ, นตฺถิ โน กิญฺจิ อุตฺตริกรณียนฺ”ติ.\csromanspacer{} อถ จ ปน ภวํ อานนฺโท เอวมาห อตฺถิ เจเวตฺถ อุตฺตริกรณียนฺติ\footnote{อิมสฺส อนนฺตรํ สีอิโปตฺถเกสุ “ปฐมภาณวารํ”ติ ปาโฐ ทิสฺสติ.}.}

\csromanmatikaanchor{mtk.123}
\csromantocmarkref{section}{สมาธิกฺขนฺธ}{mtk.123}
\bookmark[dest={mtk.123},level=1]{สมาธิกฺขนฺธ}
\csromanheader{1}{\textbf{สมาธิกฺขนฺธ}}
\proseitem{454}{กตโม ปน โส โภ อานนฺท อริโย สมาธิกฺขนฺโธ, ยสฺส โส ภวํ โคตโม วณฺณวาที อโหสิ, ยตฺถ จ อิมํ ชนตํ สมาทเปสิ นิเวเสสิ ปติฏฺฐาเปสีติ.}

\prose{กถญฺจ มาณว ภิกฺขุ อินฺทฺริเยสุ คุตฺตทฺวาโร โหติ.\csromanspacer{} อิธ มาณว ภิกฺขุ จกฺขุนา รูปํ ทิสฺวา น นิมิตฺตคฺคาหี โหติ นานุพฺยญฺชนคฺคาหี, ยตฺวาธิกรณเมนํ จกฺขุนฺทฺริยํ อสํวุตํ วิหรนฺตํ อภิชฺฌาโทมนสฺสา ปาปกา อกุสลา ธมฺมา อนฺวาสฺสเวยฺยุํ.\csromanspacer{} ตสฺส สํวราย ปฏิปชฺชติ, รกฺขติ จกฺขุนฺทฺริยํ, จกฺขุนฺทฺริเย สํวรํ อาปชฺชติ.\csromanspacer{} โสเตน สทฺทํ สุตฺวา ฯเปฯ ฆาเนน คนฺธํ ฆายิตฺวา.\csromanspacer{} ชิวฺหาย รสํ สายิตฺวา.\csromanspacer{} กาเยน โผฏฺฐพฺพํ ผุสิตฺวา.\csromanspacer{} มนสา ธมฺมํ วิญฺญาย น นิมิตฺตคฺคาหี โหติ นานุพฺยญฺชนคฺคาหี, ยตฺวาธิกรณเมนํ มนินฺทฺริยํ อสํวุตํ วิหรนฺตํ อภิชฺฌาโทมนสฺสา ปาปกา อกุสลา ธมฺมา อนฺวาสฺสเวยฺยุํ.\csromanspacer{} ตสฺส สํวราย ปฏิปชฺชติ, รกฺขติ มนินฺทฺริยํ, มนินฺทฺริเย สํวรํ อาปชฺชติ.\csromanspacer{} โส อิมินา อริเยน อินฺทฺริยสํวเรน สมนฺนาคโต อชฺฌตฺตํ อพฺยาเสกสุขํ ปฏิสํเวเทติ.\csromanspacer{} เอวํ โข มาณว ภิกฺขุ อินฺทฺริเยสุ คุตฺตทฺวาโร โหติ.}

\proseitem{455}{กถญฺจ มาณว ภิกฺขุ สติสมฺปชญฺเญน สมนฺนาคโต โหติ.\csromanspacer{} อิธ มาณว ภิกฺขุ อภิกฺกนฺเต ปฏิกฺกนฺเต สมฺปชานการี โหติ, อาโลกิเต วิโลกิเต สมฺปชานการี โหติ, สมิญฺชิเต ปสาริเต สมฺปชานการี โหติ, สํฆาฏิปตฺตจีวรธารเณ สมฺปชานการี โหติ, อสิเต ปีเต ขายิเต สายิเต สมฺปชานการี โหติ, อุจฺจารปสฺสาวกมฺเม สมฺปชานการี โหติ, คเต ฐิเต}

\vspace{\tipitakaparskip}
\csromancenter{สุภสุตฺต นิสินฺเน สุตฺเต ชาคริเต ภาสิเต ตุณฺหีภาเว สมฺปชานการี โหติ.\csromanspacer{} เอวํ โข มาณว ภิกฺขุ สติสมฺปชญฺเญน สมนฺนาคโต โหติ.}
\proseitem{456}{\csromanfolio{193}กถญฺจ มาณว ภิกฺขุ สนฺตุฏฺโฐ โหติ.\csromanspacer{} อิธ มาณว ภิกฺขุ สนฺตุฏฺโฐ โหติ กายปริหาริเกน จีวเรน กุจฺฉิปริหาริเกน ปิณฺฑปาเตน, โส เยน เยเนว ปกฺกมติ, สมาทาเยว ปกฺกมติ.\csromanspacer{} เสยฺยถาปิ มาณว ปกฺขี สกุโณ เยน เยเนว เฑติ, สปตฺตภาโรว เฑติ, เอวเมว โข มาณว ภิกฺขุ สนฺตุฏฺโฐ โหติ กายปริหาริเกน จีวเรน กุจฺฉิปริหาริเกน ปิณฺฑปาเตน, โส เยน เยเนว ปกฺกมติ, สมาทาเยว ปกฺกมติ.\csromanspacer{} เอวํ โข มาณว ภิกฺขุ สนฺตุฏฺโฐ โหติ.}

\proseitem{457}{โส อิมินา จ อริเยน สีลกฺขนฺเธน สมนฺนาคโต อิมินา จ อริเยน อินฺทฺริยสํวเรน สมนฺนาคโต อิมินา จ อริเยน สติสมฺปชญฺเญน สมนฺนาคโต อิมาย จ อริยาย สนฺตุฏฺฐิยา สมนฺนาคโต วิวิตฺตํ เสนาสนํ ภชติ อรญฺญํ รุกฺขมูลํ ปพฺพตํ กนฺทรํ คิริคุหํ สุสานํ วนปตฺถํ อพฺโภกาสํ ปลาลปุญฺชํ.\csromanspacer{} โส ปจฺฉาภตฺตํ ปิณฺฑปาตปฏิกฺกนฺโต นิสีทติ ปลฺลงฺกํ อาภุชิตฺวา อุชุํ กายํ ปณิธาย ปริมุขํ สติํ อุปฏฺฐเปตฺวา.}

\proseitem{458}{โส อภิชฺฌํ โลเก ปหาย วิคตาภิชฺเฌน เจตสา วิหรติ อภิชฺฌาย จิตฺตํ ปริโสเธติ.\csromanspacer{} พฺยาปาทปโทสํ ปหาย อพฺยาปนฺนจิตฺโต วิหรติ สพฺพปาณภูต\-หิตา\-นุกมฺปี พฺยาปาทปโทสา จิตฺตํ ปริโสเธติ.\csromanspacer{} ถินมิทฺธํ ปหาย วิคตถินมิทฺโธ วิหรติ อาโลกสญฺญี สโต สมฺปชาโน, ถินมิทฺธา จิตฺตํ ปริโสเธติ.\csromanspacer{} อุทฺธจฺจกุกฺกุจฺจํ ปหาย อนุทฺธโต วิหรติ อชฺฌตฺตํ วูปสนฺตจิตฺโต อุทฺธจฺจกุกฺกุจฺจา จิตฺตํ ปริโสเธติ.\csromanspacer{} วิจิกิจฺฉํ ปหาย ติณฺณวิจิกิจฺโฉ วิหรติ อกถํกถี กุสเลสุ ธมฺเมสุ, วิจิกิจฺฉาย จิตฺตํ ปริโสเธติ.}

\proseitem{459}{เสยฺยถาปิ มาณว ปุริโส อิณํ อาทาย กมฺมนฺเต ปโยเชยฺย, ตสฺส เต กมฺมนฺตา สมิชฺเฌยฺยุํ, โส ยานิ จ โปราณานิ อิณมูลานิ ตานิ จ พฺยนฺติํ กเรยฺย, สิยา จสฺส อุตฺตริํ อวสิฏฺฐํ ทารภรณาย, ตสฺส เอวมสฺส อหํ โข ปุพฺเพ อิณํ อาทาย กมฺมนฺเต \csromanfolio{194}ปโยเชสิํ, ตสฺส เม เต กมฺมนฺตา สมิชฺฌิํสุ, โสหํ ยานิ จ โปราณานิ อิณมูลานิ ตานิ จ พฺยานฺติํ อกาสิํ, อตฺถิ จ เม อุตฺตริํ อวสิฏฺฐํ ทารภรณายาติ.\csromanspacer{} โส ตโตนิทานํ ลเภถ ปาโมชฺชํ, อธิคจฺเฉยฺย โสมนสฺสํ.}

\proseitem{460}{เสยฺยถาปิ มาณว ปุริโส อาพาธิโก อสฺส ทุกฺขิโต พาฬฺหคิลาโน, ภตฺตญฺจสฺส นจฺฉาเทยฺย, น จสฺส กาเย พลมตฺตา, โส อปเรน สมเยน ตมฺหา อาพาธา มุจฺเจยฺย, ภตฺตญฺจสฺส ฉาเทยฺย, สิยา จสฺส กาเย พลมตฺตา.\csromanspacer{} ตสฺส เอวมสฺส “อหํ โข ปุพฺเพ อาพาธิโก อโหสิํ ทุกฺขิโต พาฬฺหคิลาโน, ภตฺตญฺจ เม นจฺฉาเทสิ, น จ เม อาสิ กาเย พลมตฺตา.\csromanspacer{} โสมฺหิ เอตรหิ ตมฺหา อาพาธา มุตฺโต ภตฺตญฺจ เม ฉาเทติ, อตฺถิ จ เม กาเย พลมตฺตา”ติ.\csromanspacer{} โส ตโตนิทานํ ลเภถ ปาโมชฺชํ, อธิคจฺเฉยฺย โสมนสฺสํ.}

\proseitem{461}{เสยฺยถาปิ มาณว ปุริโส พนฺธนาคาเร พทฺโธ อสฺส, โส อปเรน สมเยน ตมฺหา พนฺธนาคารา มุจฺเจยฺย โสตฺถินา อพฺภเยน, น จสฺส กิญฺจิ โภคานํ วโย.\csromanspacer{} ตสฺส เอวมสฺส “อหํ โข ปุพฺเพ พนฺธนาคาเร พทฺโธ อโหสิํ.\csromanspacer{} โสมฺหิ เอตรหิ ตมฺหา พนฺธนาคารา มุตฺโต โสตฺถินา อพฺภเยน, นตฺถิ จ เม กิญฺจิ โภคานํ วโย”ติ.\csromanspacer{} โส ตโตนิทานํ ลเภถ ปาโมชฺชํ, อธิคจฺเฉยฺย โสมนสฺสํ.}

\proseitem{462}{เสยฺยถาปิ มาณว ปุริโส ทาโส อสฺส อนตฺตาธีโน ปราธีโน น เยน กามํคโม, โส อปเรน สมเยน ตมฺหา ทาสพฺยา มุจฺเจยฺย, อตฺตาธีโน อปราธีโน ภุชิสฺโส เยน กามํคโม.\csromanspacer{} ตสฺส เอวมสฺส “อหํ โข ปุพฺเพ ทาโส อโหสิํ อนตฺตาธีโน ปราธีโน น เยน กามํคโม.\csromanspacer{} โสมฺหิ เอตรหิ ตมฺหา ทาสพฺยา มุตฺโต อตฺตาธีโน อปราธีโน ภุชิสฺโส เยน กามํคโม”ติ.\csromanspacer{} โส ตโตนิทานํ ลเภถ ปาโมชฺชํ, อธิคจฺเฉยฺย โสมนสฺสํ.}

\proseitem{463}{เสยฺยถาปิ มาณว ปุริโส สธโน สโภโค กนฺตารทฺธานมคฺคํ ปฏิปชฺเชยฺย ทุพฺภิกฺขํ สปฺปฏิภยํ, โส อปเรน สมเยน ตํ กนฺตารํ นิตฺถเรยฺย, โสตฺถินา คามนฺตํ อนุปาปุเณยฺย เขมํ อปฺปฏิภยํ.\csromanspacer{} ตสฺส}

\vspace{\tipitakaparskip}
\csromancenter{สุภสุตฺต เอวมสฺส “อหํ โข ปุพฺเพ สธโน สโภโค กนฺตารทฺธานมคฺคํ ปฏิปชฺชิํ ทุพฺภิกฺขํ สปฺปฏิภยํ.\csromanspacer{} โสมฺหิ เอตรหิ ตํ กนฺตารํ นิตฺถิณฺโณ, โสตฺถินา คามนฺตํ อนุปฺปตฺโต เขมํ อปฺปฏิภยนฺ”ติ.\csromanspacer{} โส ตโตนิทานํ ลเภถ ปาโมชฺชํ, อธิคจฺเฉยฺย โสมนสฺสํ.}
\proseitem{464}{\csromanfolio{195}เอวเมว โข มาณว ภิกฺขุ ยถา อิณํ ยถา โรคํ ยถา พนฺธนาคารํ ยถา ทาสพฺยํ ยถา กนฺตารทฺธานมคฺคํ, เอวํ อิเม ปญฺจ นีวรเณ อปฺปหีเน อตฺตนิ สมนุปสฺสติ.}

\proseitem{465}{เสยฺยถาปิ มาณว ยถา อาณณฺยํ ยถา อาโรคฺยํ ยถา พนฺธนาโมกฺขํ ยถา ภุชิสฺสํ ยถา เขมนฺตภูมิํ.\csromanspacer{} เอวเมว โข ภิกฺขุ อิเม ปญฺจ นีวรเณ ปหีเน อตฺตนิ สมนุปสฺสติ.}

\proseitem{466}{ตสฺสิเม ปญฺจ นีวรเณ ปหีเน อตฺตนิ สมนุปสฺสโต ปาโมชฺชํ ชายติ, ปมุทิตสฺส ปีติ ชายติ, ปีติมนสฺส กาโย ปสฺสมฺภติ, ปสฺสทฺธกาโย สุขํ เวเทติ, สุขิโน จิตฺตํ สมาธิยติ.}

\proseitem{467}{โส วิวิจฺเจว กาเมหิ วิวิจฺจ อกุสเลหิ ธมฺเมหิ สวิตกฺกํ สวิจารํ วิเวกชํ ปีติสุขํ ปฐมํ ฌานํ อุปสมฺปชฺช วิหรติ, โส อิมเมว กายํ วิเวกเชน ปีติสุเขน อภิสนฺเทติ ปริสนฺเทติ ปริปูเรติ ปริปฺผรติ, นาสฺส กิญฺจิ สพฺพาวโต กายสฺส วิเวกเชน ปีติสุเขน อปฺผุฏํ โหติ.}

\prose{เสยฺยถาปิ มาณว ทกฺโข นฺหาปโก วา นฺหาปกนฺเตวาสี วา กํสถาเล นฺหานียจุณฺณานิ อากิริตฺวา อุทเกน ปริปฺโผสกํ ปริปฺโผสกํ สนฺเนยฺย, สายํ นฺหานียปิณฺฑิ สฺเนหานุคตา สฺเนหปเรตา สนฺตรพาหิรา ผุฏา สฺเนเหน, น จ ปคฺฆรณี.\csromanspacer{} เอวเมว โข มาณว ภิกฺขุ อิมเมว กายํ วิเวกเชน ปีติสุเขน อภิสนฺเทติ ปริสนฺเทติ ปริปูเรติ ปริปฺผรติ, นาสฺส กิญฺจิ สพฺพาวโต กายสฺส วิเวกเชน ปีติสุเขน อปฺผุฏํ โหติ.\csromanspacer{} ยมฺปิ มาณว ภิกฺขุ วิวิจฺเจว กาเมหิ วิวิจฺจ อกุสเลหิ ธมฺเมหิ สวิตกฺกํ สวิจารํ วิเวกชํ ปีติสุขํ ปฐมํ ฌานํ อุปสมฺปชฺช วิหรติ, โส อิมเมว กายํ วิเวกเชน ปีติสุเขน อภิสนฺเทติ ปริสนฺเทติ ปริปูเรติ ปริปฺผรติ, นาสฺส กิญฺจิ สพฺพาวโต \csromanfolio{196}กายสฺส วิเวกเชน ปีติสุเขน อปฺผุฏํ โหติ.\csromanspacer{} อิทมฺปิสฺส โหติ สมาธิสฺมิํ.}

\proseitem{468}{ปุน จปรํ มาณว ภิกฺขุ วิตกฺกวิจารานํ วูปสมา อชฺฌตฺตํ สมฺปสาทนํ เจตโส เอโกทิภาวํ อวิตกฺกํ อวิจารํ สมาธิชํ ปีติสุขํ ทุติยํ ฌานํ อุปสมฺปชฺช วิหรติ, โส อิมเมว กายํ สมาธิเชน ปีติสุเขน อภิสนฺเทติ ปริสนฺเทติ ปริปูเรติ ปริปฺผรติ, นาสฺส กิญฺจิ สพฺพาวโต กายสฺส สมาธิเชน ปีติสุเขน อปฺผุฏํ โหติ.}

\prose{เสยฺยถาปิ มาณว อุทกรหโท คมฺภีโร อุพฺภิโททโก.\csromanspacer{} ตสฺส เนวสฺส ปุรตฺถิมาย ทิสาย อุทกสฺส อายมุขํ, น ทกฺขิณาย ทิสาย อุทกสฺส อายมุขํ, น ปจฺฉิมาย ทิสาย อุทกสฺส อายมุขํ, น อุตฺตราย ทิสาย อุทกสฺส อายมุขํ, เทโว จ น กาเลน กาลํ สมฺมา ธารํ อนุปเวจฺเฉยฺย.\csromanspacer{} อถ โข ตมฺหาว อุทกรหทา สีตา วาริธารา อุพฺภิชฺชิตฺวา ตเมว อุทกรหทํ สีเตน วารินา อภิสนฺเทยฺย ปริสนฺเทยฺย ปริปูเรยฺย ปริปฺผเรยฺย, นาสฺส กิญฺจิ สพฺพาวโต อุทกรหทสฺส สีเตน วารินา อปฺผุฏํ อสฺส.\csromanspacer{} เอวเมว โข มาณว ภิกฺขุ ฯเปฯ ยมฺปิ มาณว ภิกฺขุ วิตกฺกวิจารานํ วูปสมา ฯเปฯ ทุติยํ ฌานํ อุปสมฺปชฺช วิหรติ, โส อิมเมว กายํ สมาธิเชน ปีติสุเขน อภิสนฺเทติ ปริสนฺเทติ ปริปูเรติ ปริปฺผรติ, นาสฺส กิญฺจิ สพฺพาวโต กายสฺส สมาธิเชน ปีติสุเขน อปฺผุฏํ โหติ.\csromanspacer{} อิทมฺปิสฺส โหติ สมาธิสฺมิํ.}

\proseitem{469}{ปุน จปรํ มาณว ภิกฺขุ ปีติยา จ วิราคา อุเปกฺขโก จ วิหรติ สโต สมฺปชาโน, สุขญฺจ กาเยน ปฏิสํเวเทติ, ยํ ตํ อริยา อาจิกฺขนฺติ “อุเปกฺขโก สติมา สุขวิหารี”ติ, ตติยํ ฌานํ อุปสมฺปชฺช วิหรติ, โส อิมเมว กายํ นิปฺปีติเกน สุเขน อภิสนฺเทติ ปริสนฺเทติ ปริปูเรติ ปริปฺผรติ, นาสฺส กิญฺจิ สพฺพาวโต กายสฺส นิปฺปีติเกน สุเขน อปฺผุฏํ โหติ.}

\prose{เสยฺยถาปิ มาณว อุปฺปลินิยํ วา ปทุมินิยํ วา ปุณฺฑรีกินิยํ วา อปฺเปกจฺจานิ อุปฺปลานิ วา ปทุมานิ วา ปุณฺฑรีกานิ วา อุทเก ชาตานิ อุทเก สํวฑฺฒานิ อุทกานุคฺคตานิ อนฺโตนิมุคฺคโปสีนิ, ตานิ ยาว จคฺคา ยาว จ มูลา สีเตน วารินา อภิสนฺนานิ ปริสนฺนานิ ปริปูรานิ ปริปฺผุฏานิ, นาสฺส กิญฺจิ สพฺพาวตํ อุปฺปลานํ วา ปทุมานํ วา ปุณฺฑรีกานํ}

\vspace{\tipitakaparskip}
\csromancenter{สุภสุตฺต วา สีเตน วารินา อปฺผุฏํ อสฺส.\csromanspacer{} เอวเมว โข มาณว ภิกฺขุ ฯเปฯ ยมฺปิ มาณว ภิกฺขุ ปีติยา จ วิราคา ฯเปฯ.\csromanspacer{} ตติยํ ฌานํ อุปสมฺปชฺช วิหรติ, โส อิมเมว กายํ นิปฺปีติเกน สุเขน อภิสนฺเทติ ปริสนฺเทติ ปริปูเรติ ปริปฺผรติ, นาสฺส กิญฺจิ สพฺพาวโต กายสฺส นิปฺปีติเกน สุเขน อปฺผุฏํ โหติ.\csromanspacer{} อิทมฺปิสฺส โหติ สมาธิสฺมิํ.}
\proseitem{470}{\csromanfolio{197}ปุน จปรํ มาณว ภิกฺขุ สุขสฺส จ ปหานา ทุกฺขสฺส จ ปหานา ปุพฺเพว โสมนสฺสโทมนสฺสานํ อตฺถงฺคมา อทุกฺขมสุขํ อุเปกฺขาสติปาริสุทฺธิํ จตุตฺถํ ฌานํ อุปสมฺปชฺช วิหรติ, โส อิมเมว กายํ ปริสุทฺเธน เจตสา ปริโยทาเตน ผริตฺวา นิสินฺโน โหติ, นาสฺส กิญฺจิ สพฺพาวโต กายสฺส ปริสุทฺเธน เจตสา ปริโยทาเตน อปฺผุฏํ โหติ.}

\prose{เสยฺยถาปิ มาณว ปุริโส โอทาเตน วตฺเถน สสีสํ ปารุปิตฺวา นิสินฺโน อสฺส, นาสฺส กิญฺจิ สพฺพาวโต กายสฺส โอทาเตน วตฺเถน อปฺผุฏํ อสฺส.\csromanspacer{} เอวเมว โข มาณว ภิกฺขุ ฯเปฯ ยมฺปิ มาณว ภิกฺขุ สุขสฺส จ ปหานา ทุกฺขสฺส จ ปหานา ปุพฺเพว โสมนสฺสโทมนสฺสานํ อตฺถงฺคมา อทุกฺขมสุขํ อุเปกฺขาสติปาริสุทฺธิํ จตุตฺถํ ฌานํ อุปสมฺปชฺช วิหรติ, โส อิมเมว กายํ ปริสุทฺเธน เจตสา ปริโยทาเตน ผริตฺวา นิสินฺโน โหติ, นาสฺส กิญฺจิ สพฺพาวโต กายสฺส ปริสุทฺเธน เจตสา ปริโยทาเตน อปฺผุฏํ โหติ.\csromanspacer{} อิทมฺปิสฺส โหติ สมาธิสฺมิํ.}

\proseitem{471}{อยํ โข โส มาณว อริโย สมาธิกฺขนฺโธ, ยสฺส โส ภควา วณฺณวาที อโหสิ, ยตฺถ จ อิมํ ชนตํ สมาทเปสิ นิเวเสสิ ปติฏฺฐาเปสิ.\csromanspacer{} อตฺถิ เจเวตฺถ อุตฺตริกรณียนฺติ.}

\prose{อจฺฉริยํ โภ อานนฺท, อพฺภุตํ โภ อานนฺท, โส จายํ โภ อานนฺท อริโย สมาธิกฺขนฺโธ ปริปุณฺโณ, โน อปริปุณฺโณ.\csromanspacer{} เอวํ ปริปุณฺณญฺจาหํ โภ อานนฺท อริยํ สมาธิกฺขนฺธํ อิโต พหิทฺธา อญฺเญสุ สมณพฺราหฺมเณสุ น สมนุปสฺสามิ.\csromanspacer{} เอวํ ปริปุณฺณญฺจ โภ อานนฺท อริยํ สมาธิกฺขนฺธํ อิโต พหิทฺธา อญฺเญ สมณพฺราหฺมณา อตฺตนิ สมนุปสฺเสยฺยุํ, เต ตาวตเกเนว อตฺตมนา อสฺสุ “อลเมตฺตาวตา, กตเมตฺตาวตา, อนุปฺปตฺโต โน สามญฺญตฺโถ, นตฺถิ โน กิญฺจิ อุตฺตริกรณียนฺ”ติ.\csromanspacer{} อถ จ ปน ภวํ อานนฺโท เอวมาห “อตฺถิ เจเวตฺถ อุตฺตริกรณียนฺ”ติ.}

\csromanmatikaanchor{mtk.124}
\csromantocmarkref{section}{ปญฺญากฺขนฺธ}{mtk.124}
\bookmark[dest={mtk.124},level=1]{ปญฺญากฺขนฺธ}
\csromanheader{1}{\textbf{ปญฺญากฺขนฺธ}}
\proseitem{472}{\csromanfolio{198}กตโม ปน โส โภ อานนฺท อริโย ปญฺญากฺขนฺโธ, ยสฺส โภ ภวํ โคตโม วณฺณวาที อโหสิ, ยตฺถ จ อิมํ ชนตํ สมาทเปสิ นิเวเสสิ ปติฏฺฐาเปสีติ.}

\prose{โส เอวํ สมาหิเต จิตฺเต ปริสุทฺเธ ปริโยทาเต อนงฺคเณ วิคตูปกฺกิเลเส มุทุภูเต กมฺมนิเย ฐิเต อาเนญฺชปฺปตฺเต ญาณทสฺสนาย จิตฺตํ อภินีหรติ อภินินฺนาเมติ.\csromanspacer{} โส เอวํ ปชานาติ “อยํ โข เม กาโย รูปี จาตุมหาภูติโก มาตาเปตฺติกสมฺภโว โอทนกุมฺมาสูปจโย อนิจฺจุจฺฉาทนปริมทฺทนเภทนวิทฺธํสนธมฺโม, อิทญฺจ ปน เม วิญฺญาณํ เอตฺถ สิตํ เอตฺถ ปฏิพทฺธนฺ”ติ.}

\prose{เสยฺยถาปิ มาณว มณิ เวฬุริโย สุโภ ชาติมา อฏฺฐํโส สุปริกมฺมกโต อจฺโฉ วิปฺปสนฺโน อนาวิโล สพฺพาการสมฺปนฺโน.\csromanspacer{} ตตฺราสฺส สุตฺตํ อาวุตํ นีลํ วา ปีตํ วา โลหิตํ วา โอทาตํ วา ปณฺฑุสุตฺตํ วา.\csromanspacer{} ตเมนํ จกฺขุมา ปุริโส หตฺเถ กริตฺวา ปจฺจเวกฺเขยฺย “อยํ โข มณิ เวฬุริโย สุโภ ชาติมา อฏฺฐํโส สุปริกมฺมกโต อจฺโฉ วิปฺปสนฺโน อนาวิโล สพฺพาการสมฺปนฺโน.\csromanspacer{} ตตฺริทํ สุตฺตํ อาวุตํ นีลํ วา ปีตํ วา โลหิตํ วา โอทาตํ วา ปณฺฑุสุตฺตํ วา”ติ.\csromanspacer{} เอวเมว โข มาณว ภิกฺขุ เอวํ สมาหิเต จิตฺเต ปริสุทฺเธ ปริโยทาเต อนงฺคเณ วิคตูปกฺกิเลเส มุทุภูเต กมฺมนิเย ฐิเต อาเนญฺชปฺปตฺเต ญาณทสฺสนาย จิตฺตํ อภินีหรติ อภินินฺนาเมติ.\csromanspacer{} โส เอวํ ปชานาติ “อยํ โข เม กาโย รูปี จาตุมหาภูติโก มาตาเปตฺติกสมฺภโว โอทนกุมฺมาสูปจโย อนิจฺจุจฺฉาทนปริมทฺทนเภทนวิทฺธํสนธมฺโม.\csromanspacer{} อิทญฺจ ปน เม วิญฺญาณํ เอตฺถ สิตํ เอตฺถ ปฏิพทฺธนฺ”ติ.\csromanspacer{} ยมฺปิ มาณว ภิกฺขุ เอวํ สมาหิเต จิตฺเต ฯเปฯ อาเนญฺชปฺปตฺเต ญาณทสฺสนาย จิตฺตํ อภินีหรติ อภินินฺนาเมติ.\csromanspacer{} โส เอวํ ปชานาติ ฯเปฯ.\csromanspacer{} เอตฺถ ปฏิพทฺธนฺ”ติ.\csromanspacer{} อิทมฺปิสฺส โหติ ปญฺญาย.}

\proseitem{473}{โส เอวํ สมาหิเต จิตฺเต ปริสุทฺเธ ปริโยทาเต อนงฺคเณ วิคตูปกฺกิเลเส มุทุภูเต กมฺมนิเย ฐิเต อาเนญฺชปฺปตฺเต มโนมยํ กายํ อภินิมฺมานาย จิตฺตํ อภินีหรติ อภินินฺนาเมติ, โส อิมมฺหา กายา อญฺญํ กายํ อภินิมฺมินาติ รูปิํ มโนมยํ สพฺพงฺคปจฺจงฺคิํ อหีนินฺทฺริยํ.}

\csromanheader{2}{10. สุภสุตฺต}
\prose{\csromanfolio{199}เสยฺยถาปิ มาณว ปุริโส มุญฺชมฺหา อีสิกํ ปวาเหยฺย.\csromanspacer{} ตสฺส เอวมสฺส “อยํ มุญฺโช อยํ อีสิกา, อญฺโญ มุญฺโช อญฺญา อีสิกา, มุญฺชมฺหาเตฺวว อีสิกา ปวาฬฺหา”ติ.\csromanspacer{} เสยฺยถา วา ปน มาณว ปุริโส อสิํ โกสิยา ปวาเหยฺย.\csromanspacer{} ตสฺส เอวมสฺส “อยํ อสิ อยํ โกสิ, อญฺโญ อสิ อญฺญา โกสิ, โกสิยาเตฺวว อสิ ปวาโฬฺห”ติ.\csromanspacer{} เสยฺยถา วา ปน มาณว ปุริโส อหิํ กรณฺฑา อุทฺธเรยฺย.\csromanspacer{} ตสฺส เอวมสฺส “อยํ อหิ อยํ กรณฺโฑ, อญฺโญ อหิ อญฺโญ กรณฺโฑ, กรณฺฑาเตฺวว อหิ อุพฺภโต”ติ.\csromanspacer{} เอวเมว โข มาณว ภิกฺขุ ฯเปฯ ยมฺปิ มาณว ภิกฺขุ เอวํ สมาหิเต จิตฺเต ปริสุทฺเธ ปริโยทาเต อนงฺคเณ วิคตูปกฺกิเลเส มุทุภูเต กมฺมนิเย ฐิเต อาเนญฺชปฺปตฺเต มโนมยํ กายํ อภินิมฺมานาย จิตฺตํ อภินีหรติ อภินินฺนาเมติ ฯเปฯ.\csromanspacer{} อิทมฺปิสฺส โหติ ปญฺญาย.}

\proseitem{474}{โส เอวํ สมาหิเต จิตฺเต ปริสุทฺเธ ปริโยทาเต อนงฺคเณ วิคตูปกฺกิเลเส มุทุภูเต กมฺมนิเย ฐิเต อาเนญฺชปฺปตฺเต อิทฺธิวิธาย จิตฺตํ อภินีหรติ อภินินฺนาเมติ.\csromanspacer{} โส อเนกวิหิตํ อิทฺธิวิธํ ปจฺจนุโภติ, เอโกปิ หุตฺวา พหุธา โหติ, พหุธาปิ หุตฺวา เอโก โหติ.\csromanspacer{} อาวิภาวํ ติโรภาวํ ติโรกุฏฺฏํ ติโรปาการํ ติโรปพฺพตํ อสชฺชมาโน คจฺฉติ เสยฺยถาปิ อากาเส.\csromanspacer{} ปถวิยาปิ อุมฺมุชฺชนิมุชฺชํ กโรติ เสยฺยถาปิ อุทเก.\csromanspacer{} อุทเกปิ อภิชฺชมาเน คจฺฉติ เสยฺยถาปิ ปถวิยํ.\csromanspacer{} อากาเสปิ ปลฺลงฺเกน กมติ เสยฺยถาปิ ปกฺขี สกุโณ.\csromanspacer{} อิเมปิ จนฺทิมสูริเย เอวํ มหิทฺธิเก เอวํ มหานุภาเว ปาณินา ปรามสติ ปริมชฺชติ.\csromanspacer{} ยาว พฺรหฺมโลกาปิ กาเยน วสํ วตฺเตติ.}

\prose{เสยฺยถาปิ มาณว ทกฺโข กุมฺภกาโร วา กุมฺภการนฺเตวาสี วา สุปริกมฺมกตาย มตฺติกาย ยญฺญเทว ภาชนวิกติํ อากงฺเขยฺย.\csromanspacer{} ตํ ตเทว กเรยฺย อภินิปฺผาเทยฺย.\csromanspacer{} เสยฺยถา วา ปน มาณว ทกฺโข ทนฺตกาโร วา ทนฺตการนฺเตวาสี วา สุปริกมฺมกตสฺมิํ ทนฺตสฺมิํ ยญฺญเทว ทนฺตวิกติํ อากงฺเขยฺย.\csromanspacer{} ตํ ตเทว กเรยฺย อภินิปฺผาเทยฺย.\csromanspacer{} เสยฺยถา วา ปน มาณว ทกฺโข สุวณฺณกาโร วา สุวณฺณการนฺเตวาสี วา สุปริกมฺมกตสฺมิํ สุวณฺณสฺมิํ ยญฺญเทว สุวณฺณวิกติํ อากงฺเขยฺย, ตํ ตเทว กเรยฺย อภินิปฺผาเทยฺย.\csromanspacer{} เอวเมว โข มาณว \csromanfolio{200}ภิกฺขุ ฯเปฯ ยมฺปิ มาณว ภิกฺขุ เอวํ สมาหิเต จิตฺเต ปริสุทฺเธ ปริโยทาเต อนงฺคเณ วิคตูปกฺกิเลเส มุทุภูเต กมฺมนิเย ฐิเต อาเนญฺชปฺปตฺเต อิทฺธิวิธาย จิตฺตํ อภินีหรติ อภินินฺนาเมติ.\csromanspacer{} โส อเนกวิหิตํ อิทฺธิวิธํ ปจฺจนุโภติ, เอโกปิ หุตฺวา พหุธา โหติ ฯเปฯ ยาว พฺรหฺมโลกาปิ กาเยน วสํ วตฺเตติ.\csromanspacer{} อิทมฺปิสฺส โหติ ปญฺญาย.}

\proseitem{475}{โส เอวํ สมาหิเต จิตฺเต ฯเปฯ อาเนญฺชปฺปตฺเต ทิพฺพาย โสตธาตุยา จิตฺตํ อภินีหรติ อภินินฺนาเมติ.\csromanspacer{} โส ทิพฺพาย โสตธาตุยา วิสุทฺธาย อติกฺกนฺตมานุสิกาย อุโภ สทฺเท สุณาติ ทิพฺเพ จ มานุเส จ เย ทูเร สนฺติเก จ.\csromanspacer{} เสยฺยถาปิ มาณว ปุริโส อทฺธานมคฺคปฺปฏิปนฺโน, โส สุเณยฺย เภริสทฺทํปิ มุทิงฺคสทฺทํปิ สงฺขปณวทินฺทิมสทฺทํปิ.\csromanspacer{} ตสฺส เอวมสฺส เภริสทฺโท อิติปิ มุทิงฺคสทฺโท อิติปิ สงฺขปณวทินฺทิมสทฺโท อิติปิ\footnote{อิติปีติ (ก)}.\csromanspacer{} เอวเมว โข มาณว ภิกฺขุ ฯเปฯ ยมฺปิ มาณว ภิกฺขุ เอวํ สมาหิเต จิตฺเต ฯเปฯ อาเนญฺชปฺปตฺเต ทิพฺพาย โสตธาตุยา จิตฺตํ อภินีหรติ อภินินฺนาเมติ, โส ทิพฺพาย โสตธาตุยา วิสุทฺธาย อติกฺกนฺตมานุสิกาย อุโภ สทฺเท สุณาติ ทิพฺเพ จ มานุเส จ เย ทูเร สนฺติเก จ.\csromanspacer{} อิทมฺปิสฺส โหติ ปญฺญาย.}

\proseitem{476}{โส เอวํ สมาหิเต จิตฺเต ปริสุทฺเธ ปริโยทาเต อนงฺคเณ วิคตูปกฺกิเลเส มุทุภูเต กมฺมนิเย ฐิเต อาเนญฺชปฺปตฺเต เจโตปริยญาณย จิตฺตํ อภินีหรติ อภินินฺนาเมติ.\csromanspacer{} โส ปรสตฺตานํ ปรปุคฺคลานํ เจตสา เจโต ปริจฺจ ปชานาติ, สราคํ วา จิตฺตํ สราคํ จิตฺตนฺติ ปชานาติ, วีตราคํ วา จิตฺตํ วีตราคํ จิตฺตนฺติ ปชานาติ, สโทสํ วา จิตฺตํ สโทสํ จิตฺตนฺติ ปชานาติ, วีตโทสํ วา จิตฺตํ วีตโทสํ จิตฺตนฺติ ปชานาติ, สโมหํ วา จิตฺตํ สโมหํ จิตฺตนฺติ ปชานาติ, วีตโมหํ วา จิตฺตํ วีตโมหํ จิตฺตนฺติ ปชานาติ, สํขิตฺตํ วา จิตฺตํ สํขิตฺตํ จิตฺตนฺติ ปชานาติ, วิกฺขิตฺตํ วา จิตฺตํ วิกฺขิตฺตํ จิตฺตนฺติ ปชานาติ, มหคฺคตํ วา จิตฺตํ มหคฺคตํ จิตฺตนฺติ ปชานาติ, อมหคฺคตํ วา จิตฺตํ อมหคฺคตํ จิตฺตนฺติ ปชานาติ, ส- อุตฺตรํ วา จิตฺตํ ส-อุตฺตรํ จิตฺตนฺติ ปชานาติ, อนุตฺตรํ วา จิตฺตํ อนุตฺตรํ จิตฺตนฺติ ปชานาติ, สมาหิตํ วา จิตฺตํ สมาหิตํ จิตฺตนฺติ ปชานาติ, อสมาหิตํ วา จิตฺตํ อสมาหิตํ จิตฺตนฺติ ปชานาติ,}

\vspace{\tipitakaparskip}
\csromancenter{สุภสุตฺต วิมุตฺตํ วา จิตฺตํ วิมุตฺตํ จิตฺตนฺติ ปชานาติ, อวิมุตฺตํ วา จิตฺตํ อวิมุตฺตํ จิตฺตนฺติ ปชานาติ.}
\prose{\csromanfolio{201}เสยฺยถาปิ มาณว อิตฺถี วา ปุริโส วา ทหโร ยุวา มณฺฑนชาติโก อาทาเส วา ปริสุทฺเธ ปริโยทาเต อจฺเฉ วา อุทกปตฺเต สกํ มุขนิมิตฺตํ ปจฺจเวกฺขมาโน สกณิกํ วา สกณิกนฺติ ชาเนยฺย, อกณิกํ วา อกณิกนฺติ ชาเนยฺย.\csromanspacer{} เอวเมว โข มาณว ภิกฺขุ ฯเปฯ ยมฺปิ มาณว ภิกฺขุ เอวํ สมาหิเต ฯเปฯ อาเนญฺชปฺปตฺเต เจโตปริยญาณาย จิตฺตํ อภินีหรติ อภินินฺนาเมติ.\csromanspacer{} โส ปรสตฺตานํ ปรปุคฺคลานํ เจตสา เจโต ปริจฺจ ปชานาติ, สราคํ วา จิตฺตํ สราคํ จิตฺตนฺติ ปชานาติ ฯเปฯ อวิมุตฺตํ วา จิตฺตํ อวิมุตฺตํ จิตฺตนฺติ ปชานาติ.\csromanspacer{} อิทมฺปิสฺส โหติ ปญฺญาย.}

\proseitem{477}{โส เอวํ สมาหิเต จิตฺเต ฯเปฯ อาเนญฺชปฺปตฺเต ปุพฺเพ\-นิวาสา\-นุสฺสติ\-ญาณาย จิตฺตํ อภินีหรติ อภินินฺนาเมติ.\csromanspacer{} โส อเนกวิหิตํ ปุพฺเพนิวาสํ อนุสฺสรติ.\csromanspacer{} เสยฺยถิทํ, เอกํปิ ชาติํ เทฺวปิ ชาติโย ติสฺโสปิ ชาติโย จตสฺโสปิ ชาติโย ปญฺจปิ ชาติโย ทสปิ ชาติโย วีสํปิ ชาติโย ติํสํปิ ชาติโย จตฺตาลีสํปิ ชาติโย ปญฺญาสํปิ ชาติโย ชาติสตํปิ ชาติสหสฺสํปิ ชาติสตสหสฺสํปิ อเนเกปิ สํวฏฺฏกปฺเป อเนเกปิ วิวฏฺฏกปฺเป อเนเกปิ สํวฏฺฏวิวฏฺฏกปฺเป, อมุตฺราสิํ เอวํนาโม เอวํโคตฺโต เอวํวณฺโณ เอวมาหาโร เอวํสุขทุกฺขปฺปฏิสํเวที เอวมายุปริยนฺโต, โส ตโต จุโต อมุตฺร อุทปาทิํ, ตตฺราปาสิํ เอวํนาโม เอวํโคตฺโต เอวํวณฺโณ เอวมาหาโร เอวํสุขทุกฺขปฺปฏิสํเวที เอวมายุปริยนฺโต, โส ตโต จุโต อิธูปปนฺโนติ.\csromanspacer{} อิติ สาการํ ส-อุทฺเทสํ อเนกวิหิตํ ปุพฺเพนิวาสํ อนุสฺสรติ.}

\prose{เสยฺยถาปิ มาณว ปุริโส สกมฺหา คามา อญฺญํ คามํ คจฺเฉยฺย, ตมฺหาปิ คามา อญฺญํ คามํ คจฺเฉยฺย, โส ตมฺหา คามา สกํเยว คามํ ปจฺจาคจฺเฉยฺย.\csromanspacer{} ตสฺส เอวมสฺส อหํ โข สกมฺหา คามา อมุํ คามํ อคจฺฉิํ, ตตฺร เอวํ อฏฺฐาสิํ เอวํ นิสีทิํ เอวํ อภาสิํ เอวํ ตุณฺหี อโหสิํ, โส ตมฺหาปิ คามา อมุํ คามํ คจฺฉิํ, ตตฺราปิ เอวํ อฏฺฐาสิํ เอวํ นิสีทิํ เอวํ อภาสิํ เอวํ ตุณฺหี อโหสิํ.\csromanspacer{} โสมฺหิ ตมฺหา คามา สกํเยว คามํ ปจฺจาคโตติ.\csromanspacer{} เอวเมว โข มาณว ภิกฺขุ ฯเปฯ ยมฺปิ มาณว ภิกฺขุ \csromanfolio{202}เอวํ สมาหิเต จิตฺเต ฯเปฯ อาเนญฺชปฺปตฺเต ปุพฺเพ\-นิวาสา\-นุสฺสติ\-ญาณาย จิตฺตํ อภินีหรติ อภินินฺนาเมติ.\csromanspacer{} โส อเนกวิหิตํ ปุพฺเพนิวาสํ อนุสฺสรติ.\csromanspacer{} เสยฺยถิทํ, เอกํปิ ชาติํ ฯเปฯ.\csromanspacer{} อิติ สาการํ ส-อุทฺเทสํ อเนกวิหิตํ ปุพฺเพนิวาสํ อนุสฺสรติ.\csromanspacer{} อิทมฺปิสฺส โหติ ปญฺญาย.}

\proseitem{478}{โส เอวํ สมาหิเต จิตฺเต ฯเปฯ อาเนญฺชปฺปตฺเต สตฺตานํ จุตูปปาตญาณาย จิตฺตํ อภินีหรติ อภินินฺนาเมติ.\csromanspacer{} โส ทิพฺเพน จกฺขุนา วิสุทฺเธน อติกฺกนฺตมานุสเกน สตฺเต ปสฺสติ จวมาเน อุปปชฺชมาเน หีเน ปณีเต สุวณฺเณ ทุพฺพณฺเณ สุคเต ทุคฺคเต, ยถากมฺมูปเค สตฺเต ปชานาติ “อิเม วต โภนฺโต สตฺตา กายทุจฺจริเตน สมนฺนาคตา วจีทุจฺจริเตน สมนฺนาคตา มโนทุจฺจริเตน สมนฺนาคตา อริยานํ อุปวาทกา มิจฺฉาทิฏฺฐิกา มิจฺจาทิฏฺฐิกมฺมสมาทานา, เต กายสฺส เภทา ปรํ มรณา อปายํ ทุคฺคติํ วินิปาตํ นิรยํ อุปปนฺนา.\csromanspacer{} อิเม วา ปน โภนฺโต สตฺตา กายสุจริเตน สมนฺนาคตา วจีสุจริเตน สมนฺนาคตา มโนสุจริเตน สมนฺนาคตา อริยานํ อนุปวาทกา สมฺมาทิฏฺฐิกา สมฺมา\-ทิฏฺฐิ\-กมฺม\-สมาทานา, เต กายสฺส เภทา ปรํ มรณา สุคติํ สคฺคํ โลกํ อุปปนฺนา”ติ.\csromanspacer{} อิติ ทิพฺเพน จกฺขุนา วิสุทฺเธน อติกฺกนฺตมานุสเกน สตฺเต ปสฺสติ จวมาเน อุปปชฺชมาเน หีเน ปณีเต สุวณฺเณ ทุพฺพณฺเณ สุคเต ทุคฺคเต, ยถากมฺมูปเค สตฺเต ปชานาติ.}

\prose{เสยฺยถาปิ มาณว มชฺเฌสิงฺฆาฏเก ปาสาโท, ตตฺถ จกฺขุมา ปุริโส ฐิโต ปสฺเสยฺย มนุสฺเส เคหํ ปวิสนฺเตปิ นิกฺขมนฺเตปิ รถิกายปิ วีถิํ สญฺจรนฺเต มชฺเฌสิงฺฆาฏเก นิสินฺเนปิ, ตสฺส เอวมสฺส “เอเต มนุสฺสา เคหํ ปวิสนฺติ, เอเต นิกฺขมนฺติ, เอเต รถิกาย วีถิํ สญฺจรนฺติ, เอเต มชฺเฌสิงฺฆาฏเก นิสินฺนา”ติ.\csromanspacer{} เอวเมว โข มาณว ภิกฺขุ ฯเปฯ ยมฺปิ มาณว ภิกฺขุ เอวํ สมาหิเต จิตฺเต ฯเปฯ อาเนญฺชปฺปตฺเต สตฺตานํ จุตูปปาตญาณาย จิตฺตํ อภินีหรติ อภินินฺนาเมติ.\csromanspacer{} โส ทิพฺเพน จกฺขุนา วิสุทฺเธน อติกฺกนฺตมานุสเกน สตฺเต ปสฺสติ จวมาเน อุปปชฺชมาเน หีเน ปณีเต สุวณฺเณ ทุพฺพณฺเณ สุคเต ทุคฺคเต, ยถากมฺมูปเค สตฺเต ปชานาติ.\csromanspacer{} อิทมฺปิสฺส โหติ ปญฺญาย.}

\proseitem{479}{โส เอวํ สมาหิเต จิตฺเต ปริสุทฺเธ ปริโยทาเต อนงฺคเณ วิคตูปกฺกิเลเส มุทุภูเต กมฺมนิเย ฐิเต อาเนญฺชปฺปตฺเต อาสวานํ}

\vspace{\tipitakaparskip}
\csromancenter{สุภสุตฺต ขยญาณาย จิตฺตํ อภินีหรติ อภินินฺนาเมติ.\csromanspacer{} โส อิทํ ทุกฺขนฺติ ยถาภูตํ ปชานาติ, อยํ ทุกฺขสมุทโยติ ยถาภูตํ ปชานาติ, อยํ ทุกฺขนิโรโธติ ยถาภูตํ ปชานาติ, อยํ ทุกฺขนิโรธคามินี ปฏิปทาติ ยถาภูตํ ปชานาติ, อิเม อาสวาติ ยถาภูตํ ปชานาติ, อยํ อาสวสมุทโยติ ยถาภูตํ ปชานาติ, อยํ อาสวนิโรโธติ ยถาภูตํ ปชานาติ, อยํ อาสวนิโรธคามินี ปฏิปทาติ ยถาภูตํ ปชานาติ.\csromanspacer{} ตสฺส เอวํ ชานโต เอวํ ปสฺสโต กามาสวาปิ จิตฺตํ วิมุจฺจติ, ภวาสวาปิ จิตฺตํ วิมุจฺจติ, อวิชฺชาสวาปิ จิตฺตํ วิมุจฺจติ, วิมุตฺตสฺมิํ วิมุตฺตมิติ ญาณํ โหติ, “ขีณา ชาติ, วุสิตํ พฺรหฺมจริยํ, กตํ กรณียํ, นาปรํ อิตฺถตฺตายา”ติ ปชานาติ.}
\prose{\csromanfolio{203}เสยฺยถาปิ มาณว ปพฺพตสงฺเขเป อุทกรหโท อจฺโฉ วิปฺปสนฺโน อนาวิโล, ตตฺถ จกฺขุมา ปุริโส ตีเร ฐิโต ปสฺเสยฺย สิปฺปิกสมฺพุกมฺปิสกฺขรกถลมฺปิ มจฺฉคุมฺพมฺปิ จรนฺตมฺปิ ติฏฺฐนฺตมฺปิ.\csromanspacer{} ตสฺส เอวมสฺส “อยํ โข อุทกรหโท อจฺโฉ วิปฺปสนฺโน อนาวิโล, ตตฺริเม สิปฺปิกสมฺพุกาปิ สกฺขรกถลาปิ มจฺฉคุมฺพาปิ จรนฺติปิ ติฏฺฐนฺติปี”ติ.\csromanspacer{} เอวเมว โข มาณว ภิกฺขุ ฯเปฯ ยมฺปิ มาณว ภิกฺขุ เอวํ สมาหิเต จิตฺเต ฯเปฯ อาเนญฺชปฺปตฺเต อาสวานํ ขยญาณาย จิตฺตํ อภินีหรติ อภินินฺนาเมติ.\csromanspacer{} โส อิทํ ทุกฺขนฺติ ยถาภูตํ ปชานาติ ฯเปฯ อาสวนิโรธคามินี ปฏิปทาติ ยถาภูตํ ปชานาติ.\csromanspacer{} ตสฺส เอวํ ชานโต เอวํ ปสฺสโต กามาสวาปิ จิตฺตํ วิมุจฺจติ, ภวาสวาปิ จิตฺตํ วิมุจฺจติ, อวิชฺชาสวาปิ จิตฺตํ วิมุจฺจติ, วิมุตฺตสฺมิํ วิมุตฺตมิติ ญาณํ โหติ, “ขีณา ชาติ, วุสิตํ พฺรหฺมจริยํ, กตํ กรณียํ, นาปรํ อิตฺถตฺตายา”ติ ปชานาติ.\csromanspacer{} อิทมฺปิสฺส โหติ ปญฺญาย.}

\proseitemwithcloser{480}{อยํ โข โส มาณว อริโย ปญฺญากฺขนฺโธ, ยสฺส โส ภควา วณฺณวาที อโหสิ, ยตฺถ จ อิมํ ชนตํ สมาทเปสิ นิเวเสสิ ปติฏฺฐาเปสิ.\csromanspacer{} นตฺถิ เจเวตฺถ อุตฺตริกรณียนฺติ.\csromanspacer{} อจฺฉริยํ โภ อานนฺท, อพฺภุตํ โภ อานนฺท, โส จายํ โภ อานนฺท อริโย ปญฺญากฺขนฺโธ ปริปุณฺโณ, โน อปริปุณฺโณ.\csromanspacer{} เอวํ ปริปุณฺณญฺจาหํ โภ อานนฺท อริยํ ปญฺญากฺขนฺธํ อิโต พหิทฺธา อญฺเญสุ สมณพฺราหฺมเณสุ น สมนุปสฺสามิ, \csromanfolio{204}นตฺถิ เจเวตฺถ\footnote{น สมนุปสฺสามิ ฯเปฯ นตฺถิ โน กิญฺจิ (สฺยา, ก)} อุตฺตริกรณียํ\footnote{อุตฺตริํ กรณียนฺติ (สี, สฺยา, อิ) อุตฺตริกรณียนฺติ (ก)}.\csromanspacer{} อภิกฺกนฺตํ โภ อานนฺท, อภิกฺกนฺตํ โภ อานนฺท, เสยฺยถาปิ โภ อานนฺท นิกฺกุชฺชิตํ วา อุกฺกุชฺเชยฺย, ปฏิจฺฉนฺนํ วา วิวเรยฺย, มูฬฺหสฺส วา มคฺคํ อาจิกฺเขยฺย, อนฺธกาเร วา เตลปชฺโชตํ ธาเรยฺย “จกฺขุมนฺโต รูปานิ ทกฺขนฺตี”ติ.\csromanspacer{} เอวเมวํ โภตา อานนฺเทน อเนกปริยาเยน ธมฺโม ปกาสิโต, เอสาหํ โภ อานนฺท ตํ ภวนฺตํ โคตมํ สรณํ คจฺฉามิ ธมฺมญฺจ ภิกฺขุสํฆญฺจ, อุปาสกํ มํ ภวํ อานนฺโท ธาเรตุ อชฺชตคฺเค ปาณุเปตํ สรณํ คตนฺติ.}{\textbf{สุภสุตฺตํ นิฏฺฐิตํ ทสมํ.}}
\csromanmatikaanchor{mtk.125}
\csromantocmarknopage{chapter}{11. เกวฏฺฏสุตฺต}{mtk.125}
\bookmark[dest={mtk.125},level=0]{11. เกวฏฺฏสุตฺต}
\setcsromanheadcha{11. เกวฏฺฏสุตฺต}
\setcsromanheadhclear{}
\chapterheadpageread{11. \textbf{เกวฏฺฏสุตฺต}}
\csromanmatikaanchor{mtk.126}
\csromantocmarkref{section}{เกวฏฺฏคหปติปุตฺตวตฺถุ}{mtk.126}
\bookmark[dest={mtk.126},level=1]{เกวฏฺฏคหปติปุตฺตวตฺถุ}
\csromanheader{1}{เกวฏฺฏ\-คหปติปุตฺต\-วตฺถุ}
\proseitem{481}{\csromanfolio{205}เอวํ เม สุตํ\csromandash{}เอกํ สมยํ ภควา นาฬนฺทายํ วิหรติ ปาวาริกมฺพวเน.\csromanspacer{} อถ โข เกวฏฺโฏ คหปติปุตฺโต เยน ภควา เตนุปสงฺกมิ, อุปสงฺกมิตฺวา ภควนฺตํ อภิวาเทตฺวา เอกมนฺตํ นิสีทิ, เอกมนฺตํ นิสินฺโน โข เกวฏฺโฏ คหปติปุตฺโต ภควนฺตํ เอตทโวจ “อยํ ภนฺเต นาฬนฺทา อิทฺธา เจว ผีตา จ พหุชนา อากิณฺณมนุสฺสา ภควติ อภิปฺปสนฺนา, สาธุ ภนฺเต ภควา เอกํ ภิกฺขุํ สมาทิสตุ, โย อุตฺตริมนุสฺสธมฺมา อิทฺธิปาฏิหาริยํ กริสฺสติ, เอวายํ นาฬนฺทา ภิยฺโยโส มตฺตาย ภควติ อภิปฺปสีทิสฺสตี”ติ.\csromanspacer{} เอวํ วุตฺเต ภควา เกวฏฺฏํ คหปติปุตฺตํ เอตทโวจ “น โข อหํ เกวฏฺฏ ภิกฺขูนํ เอวํ ธมฺมํ เทเสมิ, เอถ ตุเมฺห ภิกฺขเว คิหีนํ โอทาตวสนานํ อุตฺตริมนุสฺสธมฺมา อิทฺธิปาฏิหาริยํ กโรถา”ติ.}

\proseitem{482}{ทุติยมฺปิ โข เกวฏฺโฏ คหปติปุตฺโต ภควนฺตํ เอตทโวจ “นาหํ ภนฺเต ภควนฺตํ ธํเสมิ, อปิ จ เอวํ วทามิ, อยํ ภนฺเต นาฬนฺทา อิทฺธา เจว ผีตา จ พหุชนา อากิณฺณมนุสฺสา ภควติ อภิปฺปสนฺนา, สาธุ ภนฺเต ภควา เอกํ ภิกฺขุํ สมาทิสตุ, โย อุตฺตริมนุสฺสธมฺมา อิทฺธิปาฏิหาริยํ กริสฺสติ, เอวายํ นาฬนฺทา ภิยฺโยโส มตฺตาย ภควติ อภิปฺปสีทิสฺสตี”ติ.\csromanspacer{} ทุติยมฺปิ โข ภควา เกวฏฺฏํ คหปติปุตฺตํ เอตทโวจ “น โข อหํ เกวฏฺฏ ภิกฺขูนํ เอวํ ธมฺมํ เทเสมิ, เอถ ตุเมฺห ภิกฺขเว คิหีนํ โอทาตวสนานํ อุตฺตริมนุสฺสธมฺมา อิทฺธิปาฏิหาริยํ กโรถา”ติ.}

\prose{ตติยมฺปิ โข เกวฏฺโฏ คหปติปุตฺโต ภควนฺตํ เอตทโวจ “นาหํ ภนฺเต ภควนฺตํ ธํเสมิ, อปิ จ เอวํ วทามิ, อยํ ภนฺเต นาฬนฺทา อิทฺธา เจว ผีตา จ พหุชนา อากิณฺณมนุสฺสา ภควติ อภิปฺปสนฺนา, สาธุ ภนฺเต ภควา เอกํ ภิกฺขุํ สมาทิสตุ, โย อุตฺตริมนุสฺสธมฺมา อิทฺธิปาฏิหาริยํ กริสฺสติ.\csromanspacer{} เอวายํ นาฬนฺทา ภิยฺโยโส มตฺตาย ภควติ อภิปฺปสีทิสฺสตี”ติ.}

\csromanmatikaanchor{mtk.127}
\csromantocmarkref{section}{อิทฺธิปาฏิหาริย}{mtk.127}
\bookmark[dest={mtk.127},level=1]{อิทฺธิปาฏิหาริย}
\csromanheader{1}{\textbf{อิทฺธิปาฏิหาริย}}
\proseitem{483}{\csromanfolio{206}ตีณิ โข อิมานิ เกวฏฺฏ ปาฏิหาริยานิ มยา สยํ อภิญฺญา สจฺฉิกตฺวา ปเวทิตานิ.\csromanspacer{} กตมานิ ตีณิ, อิทฺธิปาฏิหาริยํ, อาเทสนาปาฏิหาริยํ, อนุสาสนีปาฏิหาริยํ.}

\proseitem{484}{กตมญฺจ เกวฏฺฏ อิทฺธิปาฏิหาริยํ.\csromanspacer{} อิธ เกวฏฺฏ ภิกฺขุ อเนกวิหิตํ อิทฺธิวิธํ ปจฺจนุโภติ, เอโกปิ หุตฺวา พหุธา โหติ, พหุธาปิ หุตฺวา เอโก โหติ, อาวิภาวํ ติโรภาวํ ติโรกุฏฺฏํ ติโรปาการํ ติโรปพฺพตํ อสชฺชมาโน คจฺฉติ เสยฺยถาปิ อากาเส, ปถวิยาปิ อุมฺมุชฺชนิมุชฺชํ กโรติ เสยฺยถาปิ อุทเก, อุทเกปิ อภิชฺชมาเน คจฺฉติ เสยฺยถาปิ ปถวิยํ, อากาเสปิ ปลฺลงฺเกน กมติ เสยฺยถาปิ ปกฺขี สกุโณ, อิเมปิ จนฺทิมสูริเย เอวํ มหิทฺธิเก เอวํ มหานุภาเว ปาณินา ปรามสติ ปริมชฺชติ, ยาว พฺรหฺมโลกาปิ กาเยน วสํ วตฺเตติ.}

\prose{ตเมนํ อญฺญตโร สทฺโธ ปสนฺโน ปสฺสติ ตํ ภิกฺขุํ อเนกวิหิตํ อิทฺธิวิธํ ปจฺจนุโภนฺตํ เอโกปิ หุตฺวา พหุธา โหนฺตํ, พหุธาปิ หุตฺวา เอโก โหนฺตํ, อาวิภาวํ ติโรภาวํ ติโรกุฏฺฏํ ติโรปาการํ ติโรปพฺพตํ อสชฺชมานํ คจฺฉนฺตํ เสยฺยถาปิ อากาเส, ปถวิยาปิ อุมฺมุชฺชนิมุชฺชํ กโรนฺตํ เสยฺยถาปิ อุทเก, อุทเกปิ อภิชฺชมาเน คจฺฉนฺตํ เสยฺยถาปิ ปถวิยํ, อากาเสปิ ปลฺลงฺเกน กมนฺตํ เสยฺยถาปิ ปกฺขี สกุโณ, อิเมปิ จนฺทิมสูริเย เอวํ มหิทฺธิเก เอวํ มหานุภาเว ปาณินา ปรามสนฺตํ ปริมชฺชนฺตํ ยาว พฺรหฺมโลกาปิ กาเยน วสํ วตฺเตนฺตํ.}

\prose{ตเมนํ โส สทฺโธ ปสนฺโน อญฺญตรสฺส อสฺสทฺธสฺส อปฺปสนฺนสฺส อาโรเจติ “อจฺฉริยํ วต โภ อพฺภุตํ วต โภ สมณสฺส มหิทฺธิกตา มหานุภาวตา, อมาหํ ภิกฺขุํ อทฺทสํ อเนกวิหิตํ อิทฺธิวิธํ ปจฺจนุโภนฺตํ, เอโกปิ หุตฺวา พหุธา โหนฺตํ, พหุธาปิ หุตฺวา เอโก โหนฺตํ ฯเปฯ ยาว พฺรหฺมโลกาปิ กาเยน วสํ วตฺเตนฺตนฺ”ติ.}

\prose{ตเมนํ โส อสฺสทฺโธ อปฺปสนฺโน ตํ สทฺธํ ปสนฺนํ เอวํ วเทยฺย “อตฺถิ โข โภ คนฺธารี นาม วิชฺชา, ตาย โส ภิกฺขุ อเนกวิหิตํ อิทฺธิวิธํ ปจฺจนุโภติ, เอโกปิ หุตฺวา พหุธา โหติ, พหุธาปิ หุตฺวา เอโก โหติ ฯเปฯ ยาว พฺรหฺมโลกาปิ กาเยน วสํ วตฺเตตี”ติ.}

\csromanheader{2}{11. เกวฏฺฏสุตฺต}
\prose{\csromanfolio{207}ตํ กิํ มญฺญสิ เกวฏฺฏ อปิ นุ โส อสฺสทฺโธ อปฺปสนฺโน ตํ สทฺธํ ปสนฺนํ เอวํ วเทยฺยาติ.\csromanspacer{} วเทยฺย ภนฺเตติ.\csromanspacer{} อิมํ โข อหํ เกวฏฺฏ อิทฺธิปาฏิหาริเย อาทีนวํ สมฺปสฺสมาโน อิทฺธิปาฏิหาริเยน อฏฺฏียามิ หรายามิ ชิคุจฺฉามิ.}

\csromanmatikaanchor{mtk.128}
\csromantocmarkref{section}{อาเทสนาปาฏิหาริย}{mtk.128}
\bookmark[dest={mtk.128},level=1]{อาเทสนาปาฏิหาริย}
\csromanheader{1}{\textbf{อาเทสนาปาฏิหาริย}}
\proseitem{485}{กตมญฺจ เกวฏฺฏ อาเทสนาปาฏิหาริยํ.\csromanspacer{} อิธ เกวฏฺฏ ภิกฺขุ ปรสตฺตานํ ปรปุคฺคลานํ จิตฺตํปิ อาทิสติ, เจตสิกํปิ อาทิสติ, วิตกฺกิตํปิ อาทิสติ, วิจาริตํปิ อาทิสติ “เอวํปิ เต มโน, อิตฺถํปิ เต มโน, อิติปิ เต จิตฺตนฺ”ติ.}

\prose{ตเมนํ อญฺญตโร สทฺโธ ปสนฺโน ปสฺสติ ตํ ภิกฺขุํ ปรสตฺตานํ ปรปุคฺคลานํ จิตฺตํปิ อาทิสนฺตํ เจตสิกํปิ อาทิสนฺตํ วิตกฺกิตํปิ อาทิสนฺตํ วิจาริตํปิ อาทิสนฺตํ “เอวํปิ เต มโน, อิตฺถํปิ เต มโน, อิติปิ เต จิตฺตนฺ”ติ, ตเมนํ โส สทฺโธ ปสนฺโน อญฺญตรสฺส อสฺสทฺธสฺส อปฺปสนฺนสฺส อาโรเจติ\csromandash{}อจฺฉริยํ วต โภ อพฺภุตํ วต โภ สมณสฺส มหิทฺธิกตา มหานุภาวตา, อมาหํ ภิกฺขุํ อทฺทสํ ปรสตฺตานํ ปรปุคฺคลานํ จิตฺตํปิ อาทิสนฺตํ เจตสิกํปิ อาทิสนฺตํ วิตกฺกิตํปิ อาทิสนฺตํ วิจาริตํปิ อาทิสนฺตํ “เอวํปิ เต มโน, อิตฺถํปิ เต มโน, อิติปิ เต จิตฺตนฺ”ติ.}

\prose{ตเมนํ โส อสฺสทฺโธ อปฺปสนฺโน ตํ สทฺธํ ปสนฺนํ เอวํ วเทยฺย “อตฺถิ โข โภ มณิกา นาม วิชฺชา, ตาย โส ภิกฺขุ ปรสตฺตานํ ปรปุคฺคลานํ จิตฺตํปิ อาทิสติ, เจตสิกํปิ อาทิสติ, วิตกฺกิตํปิ อาทิสติ, วิจาริตํปิ อาทิสติ “เอวํปิ เต มโน, อิตฺถํปิ เต มโน, อิติปิ เต จิตฺตนฺ”ติ.}

\prose{ตํ กิํมญฺญสิ เกวฏฺฏ, อปิ นุ โส อสฺสทฺโธ อปฺปสนฺโน ตํ สทฺธํ ปสนฺนํ เอวํ วเทยฺยาติ.\csromanspacer{} วเทยฺย ภนฺเตติ.\csromanspacer{} อิมํ โข อหํ เกวฏฺฏ อาเทสนาปาฏิหาริเย อาทีนวํ สมฺปสฺสมาโน อาเทสนาปาฏิหาริเยน อฏฺฏียามิ หรายามิ ชิคุจฺฉามิ.}

\csromanmatikaanchor{mtk.129}
\csromantocmarkref{section}{อนุสาสนีปาฏิหาริย}{mtk.129}
\bookmark[dest={mtk.129},level=1]{อนุสาสนีปาฏิหาริย}
\csromanheader{1}{\textbf{อนุสาสนีปาฏิหาริย}}
\proseitem{486}{กตมญฺจ เกวฏฺฏ อนุสาสนีปาฏิหาริยํ.\csromanspacer{} อิธ เกวฏฺฏ ภิกฺขุ เอวมนุสาสติ “เอวํ วิตกฺเกถ, มา เอวํ วิตกฺกยิตฺถ, เอวํ มนสิกโรถ \csromanfolio{208}มา เอวํ มนสากตฺถ, อิทํ ปชหถ, อิทํ อุปสมฺปชฺช วิหรถา”ติ.\csromanspacer{} อิทํ วุจฺจติ เกวฏฺฏ อนุสาสนีปาฏิหาริยํ.}

\prose{ปุน จปรํ เกวฏฺฏ อิธ ตถาคโต โลเก อุปฺปชฺชติ อรหํ สมฺมาสมฺพุทฺโธ ฯเปฯ. (ยถา สามญฺญผเล, เอวํ วิตฺถาเรตพฺพํ.) เอวํ โข เกวฏฺฏ ภิกฺขุ สีลสมฺปนฺโน โหติ ฯเปฯ ปฐมํ ฌานํ อุปสมฺปชฺช วิหรติ.\csromanspacer{} อิทํปิ วุจฺจติ เกวฏฺฏ อนุสาสนีปาฏิหาริยํ.}

\prose{ฯเปฯ ทุติยํ ฌานํ.\csromanspacer{} ตติยํ ฌานํ.\csromanspacer{} จตุตฺถํ ฌานํ อุปสมฺปชฺช วิหรติ.\csromanspacer{} อิทํปิ วุจฺจติ เกวฏฺฏ อนุสาสนีปาฏิหาริยํ.\csromanspacer{} ญาณทสฺสานาย จิตฺตํ อภินีหรติ อภินินฺนาเมติ ฯเปฯ.\csromanspacer{} อิทํปิ วุจฺจติ เกวฏฺฏ อนุสาสนีปาฏิหาริยํ.\csromanspacer{} นาปรํ อิตฺถตฺตายาติ ปชานาติ ฯเปฯ.\csromanspacer{} อิทํปิ วุจฺจติ เกวฏฺฏ อนุสาสนีปาฏิหาริยํ.}

\prose{อิมานิ โข เกวฏฺฏ ตีณิ ปาฏิหาริยานิ มยา สยํ อภิญฺญา สจฺฉิกตฺวา ปเวทิตานิ.}

\csromanmatikaanchor{mtk.130}
\csromantocmarkref{section}{ภูตนิโรเธสกภิกฺขุวตฺถุ}{mtk.130}
\bookmark[dest={mtk.130},level=1]{ภูตนิโรเธสกภิกฺขุวตฺถุ}
\csromanheader{1}{\textbf{ภูตนิโรเธสกภิกฺขุวตฺถุ}}
\proseitem{487}{ภูตปุพฺพํ เกวฏฺฏ อิมสฺมิํเยว ภิกฺขุสํเฆ อญฺญตรสฺส ภิกฺขุโน เอวํ เจตโส ปริวิตกฺโก อุทปาทิ “กตฺถ นุ โข อิเม จตฺตาโร มหาภูตา อปริเสสา นิรุชฺฌนฺติ, เสยฺยถิทํ, ปถวีธาตุ อาโปธาตุ เตโชธาตุ วาโยธาตู”ติ.}

\proseitem{488}{อถ โข โส เกวฏฺฏ ภิกฺขุ ตถารูปํ สมาทิํ สมาปชฺชิ, ยถาสเมหิเต จิตฺเต เทวยานิโย มคฺโค ปาตุรโหสิ.\csromanspacer{} อถ โข โส เกวฏฺฏ ภิกฺขุ เยน จาตุมหาราชิกา เทวา เตนุปสงฺกมิ, อุปสงฺกมิตฺวา จาตุมหาราชิเก เทเว เอตทโวจ “กตฺถ นุ โข อาวุโส อิเม จตฺตาโร มหาภูตา อปริเสสา นิรุชฺฌนฺติ, เสยฺยถิทํ, ปถวีธาตุ อาโปธาตุ เตโชธาตุ วาโยธาตู”ติ.}

\prose{เอวํ วุตฺเต เกวฏฺฏ จาตุมหาราชิกา เทวา ตํ ภิกฺขุํ เอตทโวจุํ “มยํปิ โข ภิกฺขุ น ชานาม, ยตฺถิเม จตฺตาโร มหาภูตา อปริเสสา นิรุชฺฌนฺติ, เสยฺยถิทํ, ปถวีธาตุ อาโปธาตุ เตโชธาตุ วาโยธาตูติ\footnote{วาโยธาตุ. อตฺถิ โข (อิ, เอวมุปริปิ)}.\csromanspacer{} อตฺถิ โข1 ภิกฺขุ จตฺตาโร มหาราชาโน อเมฺหหิ}

\vspace{\tipitakaparskip}
\csromancenter{เกวฏฺฏสุตฺต อภิกฺกนฺตตรา จ ปณีตตรา จ, เต โข เอตํ ชาเนยฺยุํ, ยตฺถิเม จตฺตาโร มหาภูตา อปริเสสา นิรุชฺฌนฺติ, เสยฺยถิทํ, ปถวีธาตุ อาโปธาตุ เตโชธาตุ วาโยธาตู”ติ.}
\proseitem{489}{\csromanfolio{209}อถ โข โส เกวฏฺฏ ภิกฺขุ เยน จตฺตาโร มหาราชาโน เตนุปสงฺกมิ, อุปสงฺกมิตฺวา จตฺตาโร มหาราเช เอตทโวจ “กตฺถ นุ โข อาวุโส อิเม จตฺตาโร มหาภูตา อปริเสสา นิรุชฺฌนฺติ, เสยฺยถิทํ, ปถวีธาตุ อาโปธาตุ เตโชธาตุ วาโยธาตู”ติ.\csromanspacer{} เอวํ วุตฺเต เกวฏฺฏ จตฺตาโร มหาราชาโน ตํ ภิกฺขุํ เอตทโวจุํ “มยํปิ โข ภิกฺขุ น ชานาม, ยตฺถิเม จตฺตาโร มหาภูตา อปริเสสา นิรุชฺฌนฺติ, เสยฺยถิทํ, ปถวีธาตุ อาโปธาตุ เตโชธาตุ วาโยธาตูติ.\csromanspacer{} อตฺถิ โข ภิกฺขุ ตาวติํสา นาม เทวา อเมฺหหิ อภิกฺกนฺตตรา จ ปณีตตรา จ, เต โข เอตํ ชาเนยฺยุํ, ยตฺถิเม จตฺตาโร มหาภูตา อปริเสสา นิรุชฺฌนฺติ, เสยฺยถิทํ, ปถวีธาตุ อาโปธาตุ เตโชธาตุ วาโยธาตู”ติ.}

\proseitem{490}{อถ โข โส เกวฏฺฏ ภิกฺขุ เยน ตาวติํสา เทวา เตนุปสงฺกมิ, อุปสงฺกมิตฺวา ตาวติํเส เทเว เอตทโวจ “กตฺถ นุ โข อาวุโส อิเม จตฺตาโร มหาภูตา อปริเสสา นิรุชฺฌนฺติ, เสยฺยถิทํ, ปถวีธาตุ อาโปธาตุ เตโชธาตุ วาโยธาตู”ติ.\csromanspacer{} เอวํ วุตฺเต เกวฏฺฏ ตาวติํสา เทวา ตํ ภิกฺขุํ เอตทโวจุํ “มยํปิ โข ภิกฺขุ น ชานาม, ยตฺถิเม จตฺตาโร มหาภูตา อปริเสสา นิรุชฺฌนฺติ, เสยฺยถิทํ, ปถวีธาตุ อาโปธาตุ เตโชธาตุ วาโยธาตูติ.\csromanspacer{} อตฺถิ โข ภิกฺขุ สกฺโก นาม เทวานมินฺโท อเมฺหหิ อภิกฺกนฺตตโร จ ปณีตตโร จ, โส โข เอตํ ชาเนยฺย, ยตฺถิเม จตฺตาโร มหาภูตา อปริเสสา นิรุชฺฌนฺติ, เสยฺยถิทํ, ปถวีธาตุ อาโปธาตุ เตโชธาตุ วาโยธาตู”ติ.}

\proseitem{491}{อถ โข โส เกวฏฺฏ ภิกฺขุ เยน สกฺโก เทวานมินฺโท เตนุปสงฺกมิ, อุปสงฺกมิตฺวา สกฺกํ เทวานมินฺทํ เอตทโวจ “กตฺถ นุ โข อาวุโส อิเม จตฺตาโร มหาภูตา อปริเสสา นิรุชฺฌนฺติ, เสยฺยถิทํ, ปถวีธาตุ อาโปธาตุ เตโชธาตุ วาโยธาตู”ติ.\csromanspacer{} เอวํ วุตฺเต เกวฏฺฏ สกฺโก เทวานมินฺโท ตํ ภิกฺขุํ เอตทโวจ “อหํปิ โข ภิกฺขุ \csromanfolio{210}น ชานามิ, ยตฺถิเม จตฺตาโร มหาภูตา อปริเสสา นิรุชฺฌนฺติ, เสยฺยถิทํ, ปถวีธาตุ อาโปธาตุ เตโชธาตุ วาโยธาตูติ.\csromanspacer{} อตฺถิ โข ภิกฺขุ ยามา นาม เทวา ฯเปฯ.\csromanspacer{} สุยาโม นาม เทวปุตฺโต.\csromanspacer{} ตุสิตา นาม เทวา.\csromanspacer{} สนฺตุสฺสิโต นาม เทวปุตฺโต.\csromanspacer{} นิมฺมานรตี นาม เทวา.\csromanspacer{} สุนิมฺมิโต นาม เทวปุตฺโต.\csromanspacer{} ปรนิมฺมิตวสวตฺตี นาม เทวา.\csromanspacer{} วสวตฺตี นาม เทวปุตฺโต อเมฺหหิ อภิกฺกนฺตตโร จ ปณีตตโร จ, โส โข เอตํ ชาเนยฺย, ยตฺถิเม จตฺตาโร มหาภูตา อปริเสสา นิรุชฺฌนฺติ, เสยฺยถิทํ, ปถวีธาตุ อาโปธาตุ เตโชธาตุ วาโยธาตู”ติ.}

\proseitem{492}{อถ โข โส เกวฏฺฏ ภิกฺขุ เยน วสวตฺตี เทวปุตฺโต เตนุปสงฺกมิ, อุปสงฺกมิตฺวา วสวตฺติํ เทวปุตฺตํ เอตทโวจ “กตฺถ นุ โข อาวุโส อิเม จตฺตาโร มหาภูตา อปริเสสา นิรุชฺฌนฺติ, เสยฺยถิทํ, ปถวีธาตุ อาโปธาตุ เตโชธาตุ วาโยธาตู”ติ.\csromanspacer{} เอวํ วุตฺเต เกวฏฺฏ วสวตฺตี เทวปุตฺโต ตํ ภิกฺขุํ เอตทโวจ “อหํปิ โข ภิกฺขุ น ชานามิ, ยตฺถิเม จตฺตาโร มหาภูตา อปริเสสา นิรุชฺฌนฺติ, เสยฺยถิทํ, ปถวีธาตุ อาโปธาตุ เตโชธาตุ วาโยธาตูติ.\csromanspacer{} อตฺถิ โข ภิกฺขุ พฺรหฺมกายิกา นาม เทวา อเมฺหหิ อภิกฺกนฺตตรา จ ปณีตตรา จ, เต โข เอตํ ชาเนยฺยุํ, ยตฺถิเม จตฺตาโร มหาภูตา อปริเสสา นิรุชฺฌนฺติ, เสยฺยถิทํ, ปถวีธาตุ อาโปธาตุ เตโชธาตุ วาโยธาตู”ติ.}

\proseitem{493}{อถ โข โส เกวฏฺฏ ภิกฺขุ ตถารูปํ สมาธิํ สมาปชฺชิ, ยถาสมาหิเต จิตฺเต พฺรหฺมยานิโย มคฺโค ปาตุรโหสิ.\csromanspacer{} อถ โข โส เกวฏฺฏ ภิกฺขุ เยน พฺรหฺมกายิกา เทวา เตนุปสงฺกมิ, อุปสงฺกมิตฺวา พฺรหฺมกายิเก เทเว เอตทโวจ “กตฺถ นุ โข อาวุโส อิเม จตฺตาโร มหาภูตา อปริเสสา นิรุชฺฌนฺติ, เสยฺยถิทํ, ปถวีธาตุ อาโปธาตุ เตโชธาตุ วาโยธาตู”ติ.\csromanspacer{} เอวํ วุตฺเต เกวฏฺฏ พฺรหฺมกายิกา เทวา ตํ ภิกฺขุํ เอตทโวจุํ “มยํปิ โข ภิกฺขุ น ชานาม, ยตฺถิเม จตฺตาโร มหาภูตา อปริเสสา นิรุชฺฌนฺติ, เสยฺยถิทํ, ปถวีธาตุ อาโปธาตุ เตโชธาตุ วาโยธาตูติ.\csromanspacer{} อตฺถิ โข ภิกฺขุ พฺรหฺมา มหาพฺรหฺมา อภิภู อนภิภูโต อญฺญทตฺถุทโส วสวตฺตี}

\vspace{\tipitakaparskip}
\csromancenter{เกวฏฺฏสุตฺต อิสฺสโร กตฺตา นิมฺมาตา เสฏฺโฐ สชิตา วสี ปิตา ภูตภพฺยานํ อเมฺหหิ อภิกฺกนฺตตโร จ ปณีตตโร จ, โส โข เอตํ ชาเนยฺย, ยตฺถิเม จตฺตาโร มหาภูตา อปริเสสา นิรุชฺฌนฺติ, เสยฺยถิทํ, ปถวีธาตุ อาโปธาตุ เตโชธาตุ วาโยธาตู”ติ.}
\prose{\csromanfolio{211}กหํ ปนาวุโส เอตรหิ โส มหาพฺรหฺมาติ.\csromanspacer{} มยํปิ โข ภิกฺขุ น ชานาม, ยตฺถ วา พฺรหฺมา เยน วา พฺรหฺมา ยหิํ วา พฺรหฺมา, อปิ จ ภิกฺขุ ยถา นิมิตฺตา ทิสฺสนฺติ, อาโลโก สญฺชายติ, โอภาโส ปาตุภวติ, พฺรหฺมา ปาตุภวิสฺสติ, พฺรหฺมุโน เหตํ ปุพฺพนิมิตฺตํ ปาตุภาวาย, ยทิทํ อาโลโก สญฺชายติ, โอภาโส ปาตุภวตีติ.\csromanspacer{} อถ โข โส เกวฏฺฏ มหาพฺรหฺมา นจิรสฺเสว ปาตุรโหสิ.}

\proseitem{494}{อถ โข โส เกวฏฺฏ ภิกฺขุ เยน โส มหาพฺรหฺมา เตนุปสงฺกมิ, อุปสงฺกมิตฺวา ตํ มหาพฺรหฺมานํ เอตทโวจ “กตฺถ นุ โข อาวุโส อิเม จตฺตาโร มหาภูตา อปริเสสา นิรุชฺฌนฺติ, เสยฺยถิทํ, ปถวีธาตุ อาโปธาตุ เตโชธาตุ วาโยธาตู”ติ.\csromanspacer{} เอวํ วุตฺเต เกวฏฺฏ โส มหาพฺรหฺมา ตํ ภิกฺขุํ เอตทโวจ “อหมสฺมิ ภิกฺขุ พฺรหฺมา มหาพฺรหฺมา อภิภู อนภิภูโต อญฺญทตฺถุทโส วสวตฺตี อิสฺสโร กตฺตา นิมฺมาตา เสฏฺโฐ สชิตา วสี ปิตา ภูตภพฺยานนฺ”ติ.}

\prose{ทุติยมฺปิ โข โส เกวฏฺฏ ภิกฺขุ ตํ มหาพฺรหฺมานํ เอตทโวจ\csromandash{} น โขหํ ตํ อาวุโส เอวํ ปุจฺฉามิ “ตฺวมสิ พฺรหฺมา มหาพฺรหฺมา อภิภู อนภิภูโต อญฺญทตฺถุทโส วสวตฺตี อิสฺสโร กตฺตา นิมฺมาตา เสฏฺโฐ สชิตา วสี ปิตา ภูตภพฺยานนฺ”ติ.\csromanspacer{} เอวญฺจ โข อหํ ตํ อาวุโส ปุจฺฉามิ “กตฺถ นุ โข อาวุโส อิเม จตฺตาโร มหาภูตา อปริเสสา นิรุชฺฌนฺติ, เสยฺยถิทํ, ปถวีธาตุ อาโปธาตุ เตโชธาตุ วาโยธาตู”ติ.}

\prose{ทุติยมฺปิ โข โส เกวฏฺฏ มหาพฺรหฺมา ตํ ภิกฺขุํ เอตทโวจ “อหมสฺมิ ภิกฺขุ พฺรหฺมา มหาพฺรหฺมา อภิภู อนภิภูโต อญฺญทตฺถุทโส วสวตฺตี อิสฺสโร กตฺตา นิมฺมาตา เสฏฺโฐ สชิตา วสี ปิตา ภูตภพฺยานนฺ”ติ.\csromanspacer{} ตติยมฺปิ โข โส เกวฏฺฏ ภิกฺขุ ตํ มหาพฺรหฺมานํ เอตทโวจ\csromandash{}น โขหํ ตํ อาวุโส เอวํ ปุจฺฉามิ “ตฺวมสิ พฺรหฺมา มหาพฺรหฺมา \csromanfolio{212}อภิภู อนภิภูโต อญฺญทตฺถุทโส วสวตฺตี อิสฺสโร กตฺตา นิมฺมาตา เสฏฺโฐ สชิตา วสี ปิตา ภูตภพฺยานนฺ”ติ.\csromanspacer{} เอวญฺจ โข อหํ ตํ อาวุโส ปุจฺฉามิ “กตฺถ นุ โข อาวุโส อิเม จตฺตาโร มหาภูตา อปริเสสา นิรุชฺฌนฺติ, เสยฺยถิทํ, ปถวีธาตุ อาโปธาตุ เตโชธาตุ วาโยธาตู”ติ.}

\proseitem{495}{อถ โข โส เกวฏฺฏ มหาพฺรหฺมา ตํ ภิกฺขุํ พาหายํ คเหตฺวา เอกมนฺตํ อปเนตฺวา ตํ ภิกฺขุํ เอตทโวจ “อิเม โข มํ ภิกฺขุ พฺรหฺมกายิกา เทวา เอวํ ชานนฺติ ‘นตฺถิ กิญฺจิ พฺรหฺมุโน อญฺญาตํ, นตฺถิ กิญฺจิ พฺรหฺมุโน อทิฏฺฐํ, นตฺถิ กิญฺจิ พฺรหฺมุโน อวิทิตํ นตฺถิ กิญฺจิ พฺรหฺมุโน อสจฺฉิกตนฺ’ติ.\csromanspacer{} ตสฺมาหํ เตสํ สมฺมุขา น พฺยากาสิํ.\csromanspacer{} อหํปิ โข ภิกฺขุ น ชานามิ, ยตฺถิเม จตฺตาโร มหาภูตา อปริเสสา นิรุชฺฌนฺติ, เสยฺยถิทํ, ปถวีธาตุ อาโปธาตุ เตโชธาตุ วาโยธาตูติ.\csromanspacer{} ตสฺมาติห ภิกฺขุ ตุเยฺหเวตํ ทุกฺกฏํ, ตุเยฺหเวตํ อปรทฺธํ, ยํ ตฺวํ ตํ ภควนฺตํ อติธาวิตฺวา พหิทฺธา ปริเยฏฺฐิํ อาปชฺชสิ อิมสฺส ปญฺหสฺส เวยฺยากรณาย.\csromanspacer{} คจฺฉ ตฺวํ ภิกฺขุ ตเมว ภควนฺตํ อุปสงฺกมิตฺวา อิมํ ปญฺหํ ปุจฺฉ, ยถา จ เต ภควา พฺยากโรติ, ตถา นํ ธาเรยฺยาสี”ติ.}

\proseitem{496}{อถ โข โส เกวฏฺฏ ภิกฺขุ เสยฺยถาปิ นาม พลวา ปุริโส สมิญฺชิตํ วา พาหํ ปสาเรยฺย, ปสาริตํ วา พาหํ สมิญฺเชยฺย.\csromanspacer{} เอวเมว พฺรหฺมโลเก อนฺตรหิโต มม ปุรโต ปาตุรโหสิ.\csromanspacer{} อถ โข โส เกวฏฺฏ ภิกฺขุ มํ อภิวาเทตฺวา เอกมนฺตํ นิสีทิ, เอกมนฺตํ นิสินฺโน โข เกวฏฺฏ โส ภิกฺขุ มํ เอตทโวจ “กตฺถ นุ โข ภนฺเต อิเม จตฺตาโร มหาภูตา อปริเสสา นิรุชฺฌนฺติ, เสยฺยถิทํ, ปถวีธาตุ อาโปธาตุ เตโชธาตุ วาโยธาตู”ติ.}

\csromanmatikaanchor{mtk.131}
\csromantocmarkref{section}{ตีรทสฺสิสกุณุปมา}{mtk.131}
\bookmark[dest={mtk.131},level=1]{ตีรทสฺสิสกุณุปมา}
\csromanheader{1}{\textbf{ตีรทสฺสิสกุณุปมา}}
\proseitem{497}{เอวํ วุตฺเต อหํ เกวฏฺฏ ตํ ภิกฺขุํ เอตทโวจํ\csromandash{} ภูตปุพฺพํ ภิกฺขุ สามุทฺทิกา วาณิชา ตีรทสฺสิํ สกุณํ คเหตฺวา นาวาย สมุทฺทํ อชฺโฌคาหนฺติ.\csromanspacer{} เต อตีรทกฺขินิยา นาวาย ตีรทสฺสิํ สกุณํ มุญฺจนฺติ, โส คจฺฉเตว ปุรตฺถิมํ ทิสํ, คจฺฉติ ทกฺขิณํ ทิสํ, คจฺฉติ ปจฺฉิมํ ทิสํ, คจฺฉติ อุตฺตรํ ทิสํ, คจฺฉติ อุทฺธํ ทิสํ, คจฺฉติ อนุทิสํ.\csromanspacer{} สเจ โส สมนฺตา}

\vspace{\tipitakaparskip}
\csromancenter{เกวฏฺฏสุตฺต ตีรํ ปสฺสติ, ตถาคตโกว\footnote{ตถาปกฺกนฺโตว (สฺยา)} โหติ.\csromanspacer{} สเจ ปน โส สมนฺตา ตีรํ น ปสฺสติ, ตเมว นาวํ ปจฺจาคจฺฉติ.\csromanspacer{} เอวเมว โข ตฺวํ ภิกฺขุ ยโต ยาว พฺรหฺมโลกา ปริเยสมาโน อิมสฺส ปญฺหสฺส เวยฺยากรณํ นาชฺฌคา, อถ มมญฺเญว สนฺติเก ปจฺจาคโต.\csromanspacer{} น โข เอโส ภิกฺขุ ปโญฺห เอวํ ปุจฺฉิตพฺโพ “กตฺถ นุ โข ภนฺเต อิเม จตฺตาโร มหาภูตา อปริเสสา นิรุชฺฌนฺติ, เสยฺยถิทํ, ปถวีธาตุ อาโปธาตุ เตโชธาตุ วาโยธาตู”ติ.}
\csromanhangingbegin
\hangingheaditem{498}{\csromanfolio{213}เอวญฺจ โข เอโส ภิกฺขุ ปโญฺห ปุจฺฉิตพฺโพ.}
\hangingbody{“กตฺถ อาโป จ ปถวี, เตโช วาโย น คาธติ.}
\hangingbody{กตฺถ ทีฆญฺจ รสฺสญฺจ, อณุํ ถูลํ สุภาสุภํ.}
\hangingbody{กตฺถ นามญฺจ รูปญฺจ, อเสสํ อุปรุชฺฌตี”ติ.}
\csromanhangingend

\csromanhangingbegin
\hangingheaditem{499}{ตตฺร เวยฺยากรณํ ภวติ.}
\hangingbody{“วิญฺญาณํ อนิทสฺสนํ, อนนฺตํ สพฺพโตปภํ.}
\hangingbody{เอตฺถ อาโป จ ปถวี, เตโช วาโย น คาธติ.}
\hangingbody{เอตฺถ ทีฆญฺจ รสฺสญฺจ, อณุํ ถูลํ สุภาสุภํ.}
\hangingbody{เอตฺถ นามญฺจ รูปญฺจ, อเสสํ อุปรุชฺฌติ.}
\hangingbody{วิญฺญาณสฺส นิโรเธน, เอตฺเถตํ อุปรุชฺฌตี”ติ.}
\csromanhangingend

\proseitemwithcloser{500}{อิทมโวจ ภควา.\csromanspacer{} อตฺตมโน เกวฏฺโฏ คหปติปุตฺโต ภควโต ภาสิตํ อภินนฺทีติ.}{เกวฏฺฏสุตฺตํ นิฏฺฐิตํ เอกาทสมํ.}
\csromanmatikaanchor{mtk.132}
\csromantocmarknopage{chapter}{12. โลหิจฺจสุตฺต}{mtk.132}
\bookmark[dest={mtk.132},level=0]{12. โลหิจฺจสุตฺต}
\setcsromanheadcha{12. โลหิจฺจสุตฺต}
\setcsromanheadhclear{}
\chapterheadpageread{12. \textbf{โลหิจฺจสุตฺต}}
\csromanmatikaanchor{mtk.133}
\csromantocmarkref{section}{โลหิจฺจพฺราหฺมณวตฺถุ}{mtk.133}
\bookmark[dest={mtk.133},level=1]{โลหิจฺจพฺราหฺมณวตฺถุ}
\csromanheader{1}{\textbf{โลหิจฺจพฺราหฺมณวตฺถุ}}
\proseitem{501}{\csromanfolio{214}เอวํ เม สุตํ\csromandash{}เอกํ สมยํ ภควา โกสเลสุ จาริกํ จรมาโน มหตา ภิกฺขุสํเฆน สทฺธิํ ปญฺจมตฺเตหิ ภิกฺขุสเตหิ เยน สาลวติกา ตทวสริ.\csromanspacer{} เตน โข ปน สมเยน โลหิจฺโจ พฺราหฺมโณ สาลวติกํ อชฺฌาวสติ สตฺตุสฺสทํ สติณกฏฺโฐทกํ สธญฺญํ ราชโภคฺคํ รญฺญา ปเสนทินา โกสเลน ทินฺนํ ราชทายํ พฺรหฺมเทยฺยํ.}

\proseitem{502}{เตน โข ปน สมเยน โลหิจฺจสฺส พฺราหฺมณสฺส เอวรูปํ ปาปกํ ทิฏฺฐิคตํ อุปฺปนฺนํ โหติ “อิธ สมโณ วา พฺราหฺมโณ วา กุสลํ ธมฺมํ อธิคจฺเฉยฺย, กุสลํ ธมฺมํ อธิคนฺตฺวา น ปรสฺส อาโรเจยฺย, กิญฺหิ ปโร ปรสฺส กริสฺสติ.\csromanspacer{} เสยฺยถาปิ นาม ปุราณํ พนฺธนํ ฉินฺทิตฺวา อญฺญํ นวํ พนฺธนํ กเรยฺย, เอวํ สมฺปทมิทํ ปาปกํ โลภธมฺมํ วทามิ, กิญฺหิ ปโร ปรสฺส กริสฺสตี”ติ.}

\proseitem{503}{อสฺโสสิ โข โลหิจฺโจ พฺราหฺมโณ “สมโณ ขลุ โภ โคตโม สกฺยปุตฺโต สกฺยกุลา ปพฺพชิโต โกสเลสุ จาริกํ จรมาโน มหตา ภิกฺขุสํเฆน สทฺธิํ ปญฺจมตฺเตหิ ภิกฺขุสเตหิ สาลวติกํ อนุปฺปตฺโต.\csromanspacer{} ตํ โข ปน ภวนฺตํ โคตมํ เอวํ กลฺยาโณ กิตฺติสทฺโท อพฺภุคฺคโต ‘อิติปิ โส ภควา อรหํ สมฺมาสมฺพุทฺโธ วิชฺชาจรณสมฺปนฺโน สุคโต โลกวิทู อนุตฺตโร ปุริสทมฺมสารถิ สตฺถา เทวมนุสฺสานํ พุทฺโธ ภควา’.\csromanspacer{} โส อิมํ โลกํ สเทวกํ สมารกํ สพฺรหฺมกํ สสฺสมณพฺราหฺมณิํ ปชํ สเทวมนุสฺสํ สยํ อภิญฺญา สจฺฉิกตฺวา ปเวเทติ, โส ธมฺมํ เทเสติ อาทิกลฺยาณํ มชฺเฌกลฺยาณํ ปริโยสานกลฺยาณํ สาตฺถํ สพฺยญฺชนํ เกวลปริปุณฺณํ ปริสุทฺธํ พฺรหฺมจริยํ ปกาเสติ, สาธุ โข ปน ตถารูปานํ อรหตํ ทสฺสนํ โหตี”ติ.}

\proseitem{504}{อถ โข โลหิจฺโจ พฺราหฺมโณ โรสิกํ\footnote{เภสิกํ (สี, อิ)} นฺหาปิตํ อามนฺเตสิ\csromandash{}เอหิ ตฺวํ สมฺม โรสิเก เยน สมโณ โคตโม}

\vspace{\tipitakaparskip}
\csromancenter{โลหิจฺจสุตฺต เตนุปสงฺกม, อุปสงฺกมิตฺวา มม วจเนน สมณํ โคตมํ อปฺปาพาธํ อปฺปาตงฺกํ ลหุฏฺฐานํ พลํ ผาสุวิหารํ ปุจฺฉ “โลหิจฺโจ โภ โคตม พฺราหฺมโณ ภวนฺตํ โคตมํ อปฺปาพาธํ อปฺปาตงฺกํ ลหุฏฺฐานํ พลํ ผาสุวิหารํ ปุจฺฉตี”ติ.\csromanspacer{} เอวญฺจ วเทหิ “อธิวาเสตุ กิร ภวํ โคตโม โลหิจฺจสฺส พฺราหฺมณสฺส สฺวาตนาย ภตฺตํ สทฺธิํ ภิกฺขุสํเฆนา”ติ.}
\proseitem{505}{\csromanfolio{215}“เอวํ โภ”ติ\footnote{เอวํ ภนฺเตติ (สี, อิ)} โข โรสิกา นฺหาปิโต โลหิจฺจสฺส พฺราหฺมณสฺส ปฏิสฺสุตฺวา เยน ภควา เตนุปสงฺกมิ, อุปสงฺกมิตฺวา ภควนฺตํ อภิวาเทตฺวา เอกมนฺตํ นิสีทิ.\csromanspacer{} เอกมนฺตํ นิสินฺโน โข โรสิกา นฺหาปิโต ภควนฺตํ เอตทโวจ “โลหิจฺโจ ภนฺเต พฺราหฺมโณ ภควนฺตํ อปฺปาพาธํ อปฺปาตงฺกํ ลหุฏฺฐานํ พลํ ผาสุวิหารํ ปุจฺฉติ, เอวญฺจ วเทติ อธิวาเสตุ กิร ภนฺเต ภควา โลหิจฺจสฺส พฺราหฺมณสฺส สฺวาตนาย ภตฺตํ สทฺธิํ ภิกฺขุสํเฆนา”ติ.\csromanspacer{} อธิวาเสสิ ภควา ตุณฺหีภาเวน.}

\proseitem{506}{อถ โข โรสิกา นฺหาปิโต ภควโต อธิวาสนํ วิทิตฺวา อุฏฺฐายาสนา ภควนฺตํ อภิวาเทตฺวา ปทกฺขิณํ กตฺวา เยน โลหิจฺโจ พฺราหฺมโณ เตนุปสงฺกมิ, อุปสงฺกมิตฺวา โลหิจฺจํ พฺราหฺมณํ เอตทโวจ “อโวจุมฺหา โข มยํ โภโต\footnote{มยํ ภนฺเต ตว (สี, อิ)} วจเนน ตํ ภควนฺตํ โลหิจฺโจ ภนฺเต พฺราหฺมโณ ภควนฺตํ อปฺปาพาธํ อปฺปาตงฺกํ ลหุฏฺฐานํ พลํ ผาสุวิหารํ ปุจฺฉติ, เอวญฺจ วเทติ อธิวาเสตุ กิร ภนฺเต ภควา โลหิจฺจสฺส พฺราหฺมณสฺส สฺวาตนาย ภตฺตํ สทฺธิํ ภิกฺขุสํเฆนา’ติ.\csromanspacer{} อธิวุตฺถญฺจ ปน เตน ภควตา”ติ.}

\proseitem{507}{อถ โข โลหิจฺโจ พฺราหฺมโณ ตสฺสา รตฺติยา อจฺจเยน สเก นิเวสเน ปณีตํ ขาทนียํ โภชนียํ ปฏิยาทาเปตฺวา โรสิกํ นฺหาปิตํ อามนฺเตสิ “เอหิ ตฺวํ สมฺม โรสิเก เยน สมโณ โคตโม เตนุปสงฺกม, อุปสงฺกมิตฺวา สมณสฺส โคตมสฺส กาลํ อาโรเจหิ กาโล โภ โคตม นิฏฺฐิตํ ภตฺตนฺ”ติ. “เอวํ โภ”ติ โข โรสิกา นฺหาปิโต โลหิจฺจสฺส พฺราหฺมณสฺส ปฏิสฺสุตฺวา เยน ภควา เตนุปสงฺกมิ, อุปสงฺกมิตฺวา ภควนฺตํ อภิวาเทตฺวา เอกมนฺตํ \csromanfolio{216}อฏฺฐาสิ.\csromanspacer{} เอกมนฺตํ ฐิโต โข โรสิกา นฺหาปิโต ภควโต กาลํ อาโรเจสิ “กาโล ภนฺเต นิฏฺฐิตํ ภตฺตนฺ”ติ.}

\proseitem{508}{อถ โข ภควา ปุพฺพณฺหสมยํ นิวาเสตฺวา ปตฺตจีวรมาทาย สทฺธิํ ภิกฺขุสํเฆน เยน สาลวติกา เตนุปสงฺกมิ.\csromanspacer{} เตน โข ปน สมเยน โรสิกา นฺหาปิโต ภควนฺตํ ปิฏฺฐิโต ปิฏฺฐิโต อนุพนฺโธ โหติ.\csromanspacer{} อถ โข โรสิกา นฺหาปิโต ภควนฺตํ เอตทโวจ “โลหิจฺจสฺส ภนฺเต พฺราหฺมณสฺส เอวรูปํ ปาปกํ ทิฏฺฐิคตํ อุปฺปนฺนํ ‘อิธ สมโณ วา พฺราหฺมโณ วา กุสลํ ธมฺมํ อธิคจฺเฉยฺย, กุสลํ ธมฺมํ อธิคนฺตฺวา น ปรสฺส อาโรเจยฺย, กิญฺหิ ปโร ปรสฺส กริสฺสติ.\csromanspacer{} เสยฺยถาปิ นาม ปุราณํ พนฺธนํ ฉินฺทิตฺวา อญฺญํ นวํ พนฺธนํ กเรยฺย, เอวํ สมฺปทมิทํ ปาปกํ โลภธมฺมํ วทามิ, กิญฺหิ ปโร ปรสฺส กริสฺสตี’ติ.\csromanspacer{} สาธุ ภนฺเต ภควา โลหิจฺจํ พฺราหฺมณํ เอตสฺมา ปาปกา ทิฏฺฐิคตา วิเวเจตู”ติ.\csromanspacer{} อปฺเปว นาม สิยา โรสิเก, อปฺเปว นาม สิยา โรสิเกติ.}

\prose{อถ โข ภควา เยน โลหิจฺจสฺส พฺราหฺมณสฺส นิเวสนํ เตนุปสงฺกมิ, อุปสงฺกมิตฺวา ปญฺญตฺเต อาสเน นิสีทิ.\csromanspacer{} อถ โข โลหิจฺโจ พฺราหฺมโณ พุทฺธปฺปมุขํ ภิกฺขุสํฆํ ปณีเตน ขาทนีเยน โภชนีเยน สหตฺถา สนฺตปฺเปสิ สมฺปวาเรสิ.}

\csromanmatikaanchor{mtk.134}
\csromantocmarkref{section}{โลหิจฺจพฺราหฺมณานุโยค}{mtk.134}
\bookmark[dest={mtk.134},level=1]{โลหิจฺจพฺราหฺมณานุโยค}
\csromanheader{1}{โลหิจฺจ\-พฺราหฺมณา\-นุโยค}
\proseitem{509}{อถ โข โลหิจฺโจ พฺราหฺมโณ ภควนฺตํ ภุตฺตาวิํ โอนีตปตฺตปาณิํ อญฺญตรํ นีจํ อาสนํ คเหตฺวา เอกมนฺตํ นิสีทิ.\csromanspacer{} เอกมนฺตํ นิสินฺนํ โข โลหิจฺจํ พฺราหฺมณํ ภควา เอตทโวจ\csromandash{} สจฺจํ กิร เต โลหิจฺจ เอวรูปํ ปาปกํ ทิฏฺฐิคตํ อุปฺปนฺนํ “อิธ สมโณ วา พฺราหฺมโณ วา กุสลํ ธมฺมํ อธิคจฺเฉยฺย, กุสลํ ธมฺมํ อธิคนฺตฺวา น ปรสฺส อาโรเจยฺย, กิญฺหิ ปโร ปรสฺส กริสฺสติ.\csromanspacer{} เสยฺยถาปิ นาม ปุราณํ พนฺธนํ ฉินฺทิตฺวา อญฺญํ นวํ พนฺธนํ กเรยฺย, เอวํ สมฺปทมิทํ ปาปกํ โลภธมฺมํ วทามิ, กิญฺหิ ปโร ปรสฺส กริสฺสตี”ติ.\csromanspacer{} เอวํ โภ โคตม.\csromanspacer{} ตํ กิํ มญฺญสิ โลหิจฺจ นนุ ตฺวํ สาลวติกํ อชฺฌาวสสีติ.\csromanspacer{} เอวํ โภ โคตม.\csromanspacer{} โย นุ โข โลหิจฺจ เอวํ วเทยฺย “โลหิจฺโจ พฺราหฺมโณ สาลวติกํ อชฺฌาวสติ, ยา สาลวติกาย สมุทยสญฺชาติ, โลหิจฺโจว ตํ}

\vspace{\tipitakaparskip}
\csromancenter{โลหิจฺจสุตฺต พฺราหฺมโณ เอกโก ปริภุญฺเชยฺย, น อญฺเญสํ ทเทยฺยา”ติ.\csromanspacer{} เอวํ วาที โส เย ตํ อุปชีวนฺติ, เตสํ อนฺตรายกโร วา โหติ, โน วาติ.}
\prose{\csromanfolio{217}อนฺตรายกโร โภ โคตม.\csromanspacer{} อนฺตรายกโร สมาโน หิตานุกมฺปี วา เตสํ โหติ อหิตานุกมฺปี วาติ.\csromanspacer{} อหิตานุกมฺปี โภ โคตม.\csromanspacer{} อหิตานุกมฺปิสฺส เมตฺตํ วา เตสุ จิตฺตํ ปจฺจุปฏฺฐิตํ โหติ สปตฺตกํ วาติ.\csromanspacer{} สปตฺตกํ โภ โคตม.\csromanspacer{} สปตฺตเก จิตฺเต ปจฺจุปฏฺฐิเต มิจฺฉาทิฏฺฐิ วา โหติ สมฺมาทิฏฺฐิ วาติ.\csromanspacer{} มิจฺฉาทิฏฺฐิ โภ โคตม.\csromanspacer{} มิจฺฉาทิฏฺฐิสฺส โข อหํ โลหิจฺจ ทฺวินฺนํ คตีนํ อญฺญตรํ คติํ วทามิ นิรยํ วา ติรจฺฉานโยนิํ วา.}

\proseitem{510}{ตํ กิํ มญฺญสิ โลหิจฺจ นนุ ราชา ปเสนทิ โกสโล กาสิโกสลํ อชฺฌาวสตีติ.\csromanspacer{} เอวํ โภ โคตม.\csromanspacer{} โย นุ โข โลหิจฺจ เอวํ วเทยฺย “ราชา ปเสนทิ โกสโล กาสิโกสลํ อชฺฌาวสติ, ยา กาสิโกสเล สมุทยสญฺชาติ, ราชาว ตํ ปเสนทิ โกสโล เอกโก ปริภุญฺเชยฺย, น อญฺเญสํ ทเทยฺยา”ติ.\csromanspacer{} เอวํ วาที โส เย ราชานํ ปเสนทิํ โกสลํ อุปชีวนฺติ ตุเมฺห เจว อญฺเญ จ, เตสํ อนฺตรายกโร วา โหติ, โน วา ติ.}

\prose{อนฺตรายกโร โภ โคตม.\csromanspacer{} อนฺตรายกโร สมาโน หิตานุกมฺปี วา เตสํ โหติ อหิตานุกมฺปี วาติ.\csromanspacer{} อหิตานุกมฺปี โภ โคตม.\csromanspacer{} อหิตานุกมฺปิสฺส เมตฺตํ วา เตสุ จิตฺตํ ปจฺจุปฏฺฐิตํ โหติ สปตฺตกํ วาติ.\csromanspacer{} สปตฺตกํ โภ โคตม.\csromanspacer{} สปตฺตเก จิตฺเต ปจฺจุปฏฺฐิเต มิจฺฉาทิฏฺฐิ วา โหติ สมฺมาทิฏฺฐิ วาติ.\csromanspacer{} มิจฺฉาทิฏฺฐิ โภ โคตม.\csromanspacer{} มิจฺฉาทิฏฺฐิสฺส โข อหํ โลหิจฺจ ทฺวินฺนํ คตีนํ อญฺญตรํ คติํ วทามิ นิรยํ วา ติรจฺฉานโยนิํ วา.}

\proseitem{511}{อิติ กิร โลหิจฺจ โย เอวํ วเทยฺย “โลหิจฺโจ พฺราหฺมโณ สาลวติกํ อชฺฌาวสติ, ยา สาลวติกาย สมุทยสญฺชาติ, โลหิจฺโจว ตํ พฺราหฺมโณ เอกโก ปริภุญฺเชยฺย, น อญฺเญสํ ทเทยฺยา”ติ.\csromanspacer{} เอวํ วาที โส เย ตํ อุปชีวนฺติ, เตสํ อนฺตรายกโร โหติ.\csromanspacer{} อนฺตรายกโร สมาโน อหิตานุกมฺปี โหติ, อหิตานุกมฺปิสฺส สปตฺตกํ จิตฺตํ ปจฺจุปฏฺฐิตํ โหติ, สปตฺตเก จิตฺเต ปจฺจุปฏฺฐิเต มิจฺฉาทิฏฺฐิ โหติ.\csromanspacer{} เอวเมว โข โลหิจฺจ โย เอวํ วเทยฺย “อิธ \csromanfolio{218}สมโณ วา พฺราหฺมโณ วา กุสลํ ธมฺมํ อธิคจฺเฉยฺย, กุสลํ ธมฺมํ อธิคนฺตฺวา น ปรสฺส อาโรเจยฺย, กิญฺหิ ปโร ปรสฺส กริสฺสติ.\csromanspacer{} เสยฺยถาปิ นาม ปุราณํ พนฺธนํ ฉินฺทิตฺวา อญฺญํ นวํ พนฺธนํ กเรยฺย ฯเปฯ กริสฺสตี”ติ.\csromanspacer{} เอวํ วาที โส เย เต กุลปุตฺตา ตถาคตปฺปเวทิตํ ธมฺมวินยํ อาคมฺม เอวรูปํ อุฬารํ วิเสสํ อธิคจฺฉนฺติ, โสตาปตฺติผลํปิ สจฺฉิกโรนฺติ, สกทาคามิผลํปิ สจฺฉิกโรนฺติ, อนาคามิผลํปิ สจฺฉิกโรนฺติ, อรหตฺตมฺปิ สจฺฉิกโรนฺติ, เย จิเม ทิพฺพา คพฺภา ปริปาเจนฺติ ทิพฺพานํ ภวานํ อภินิพฺพตฺติยา, เตสํ อนฺตรายกโร โหติ, อนฺตรายกโร สมาโน อหิตานุกมฺปี โหติ, อหิตานุกมฺปิสฺส สปตฺตกํ จิตฺตํ ปจฺจุปฏฺฐิตํ โหติ, สปตฺตเก จิตฺเต ปจฺจุปฏฺฐิเต มิจฺฉาทิฏฺฐิ โหติ.\csromanspacer{} มิจฺฉาทิฏฺฐิสฺส โข อหํ โลหิจฺจ ทฺวินฺนํ คตีนํ อญฺญตรํ คติํ วทามิ นิรยํ วา ติรจฺฉานโยนิํ วา.}

\proseitem{512}{อิติ กิร โลหิจฺจ โย เอวํ วเทยฺย “ราชา ปเสนทิ โกสโล กาสิโกสลํ อชฺฌาวสติ, ยา กาสิโกสเล สมุทยสญฺชาติ, ราชาว ตํ ปเสนทิ โกสโล เอกโก ปริภุญฺเชยฺย, น อญฺเญสํ ทเทยฺยา”ติ.\csromanspacer{} เอวํ วาที โส เย ราชานํ ปเสนทิํ โกสลํ อุปชีวนฺติ ตุเมฺห เจว อญฺเญ จ, เตสํ อนฺตรายกโร โหติ.\csromanspacer{} อนฺตรายกโร สมาโน อหิตานุกมฺปี โหติ, อหิตานุกมฺปิสฺส สปตฺตกํ จิตฺตํ ปจฺจุปฏฺฐิตํ โหติ, สปตฺตเก จิตฺเต ปจฺจุปฏฺฐิเต มิจฺฉาทิฏฺฐิ โหติ.\csromanspacer{} เอวเมว โข โลหิจฺจ โย เอวํ วเทยฺย, “อิธ สมโณ วา พฺราหฺมโณ วา กุสลํ ธมฺมํ อธิคจฺเฉยฺย, กุสลํ ธมฺมํ อธิคนฺตฺวา น ปรสฺส อาโรเจยฺย, กิญฺหิ ปโร ปรสฺส กริสฺสติ.\csromanspacer{} เสยฺยถาปิ นาม ฯเปฯ กิญฺหิ ปโร ปรสฺส กริสฺสตี”ติ.\csromanspacer{} เอวํ วาที โส เย เต กุลปุตฺตา ตถาคตปฺปเวทิตํ ธมฺมวินยํ อาคมฺม เอวรูปํ อุฬารํ วิเสสํ อธิคจฺฉนฺติ, โสตาปตฺติผลํปิ สจฺฉิกโรนฺติ, สกทาคามิผลํปิ สจฺฉิกโรนฺติ, อนาคามิผลํปิ สจฺฉิกโรนฺติ, อรหตฺตํปิ สจฺฉิกโรนฺติ, เย จิเม ทิพฺพา คพฺภา ปริปาเจนฺติ ทิพฺพานํ ภวานํ อภินิพฺพตฺติยา, เตสํ อนฺตรายกโร โหติ, อนฺตรายกโร สมาโน อหิตานุกมฺปี โหติ, อหิตานุกมฺปิสฺส สปตฺตกํ จิตฺตํ ปจฺจุปฏฺฐิตํ โหติ, สปตฺตเก จิตฺเต ปจฺจุปฏฺฐิเต มิจฺฉาทิฏฺฐิ โหติ.\csromanspacer{} มิจฺฉาทิฏฺฐิสฺส โข อหํ โลหิจฺจ ทฺวินฺนํ คตีนํ อญฺญตรํ คติํ วทามิ นิรยํ วา ติรจฺฉานโยนิํ วา.}

\csromanmatikaanchor{mtk.135}
\csromantocmarkref{section}{ตโยโจทนารหา}{mtk.135}
\bookmark[dest={mtk.135},level=1]{ตโยโจทนารหา}
\csromanheader{1}{12. โลหิจฺจสุตฺต \textbf{ตโยโจทนารหา}}
\proseitem{513}{\csromanfolio{219}ตโย โขเม โลหิจฺจ สตฺถาโร, เย โลเก โจทนารหา, โย จ ปเนวรูเป สตฺถาโร โจเทติ, สา โจทนา ภูตา ตจฺฉา ธมฺมิกา อนวชฺชา.\csromanspacer{} กตเม ตโย.\csromanspacer{} อิธ โลหิจฺจ เอกจฺโจ สตฺถา ยสฺสตฺถาย อคารสฺมา อนคาริยํ ปพฺพชิโต โหติ, สฺวาสฺส สามญฺญตฺโถ อนนุปฺปตฺโต โหติ.\csromanspacer{} โส ตํ สามญฺญตฺถํ อนนุปาปุณิตฺวา สาวกานํ ธมฺมํ เทเสติ “อิทํ โว หิตาย อิทํ โว สุขายา”ติ.\csromanspacer{} ตสฺส สาวกา น สุสฺสูสนฺติ, น โสตํ โอทหนฺติ, น อญฺญา จิตฺตํ อุปฏฺฐเปนฺติ, โวกฺกมฺม จ สตฺถุสาสนา วตฺตนฺติ.\csromanspacer{} โส เอวมสฺส โจเทตพฺโพ “อายสฺมา โข ยสฺสตฺถาย อคารสฺมา อนคาริยํ ปพฺพชิโต, โส เต สามญฺญตฺโถ อนนุปฺปตฺโต, ตํ ตฺวํ สามญฺญตฺถํ อนนุปาปุณิตฺวา สาวกานํ ธมฺมํ เทเสสิ ‘อิทํ โว หิตาย อิทํ โว สุขายา’ติ.\csromanspacer{} ตสฺส เต สาวกา น สุสฺสูสนฺติ, น โสตํ โอทหนฺติ, น อญฺญา จิตฺตํ อุปฏฺฐเปนฺติ, โวกฺกมฺม จ สตฺถุสาสนา วตฺตนฺติ.\csromanspacer{} เสยฺยถาปิ นาม โอสกฺกนฺติยา วา อุสฺสกฺเกยฺย, ปรมฺมุขิํ วา อาลิงฺเคยฺย, เอวํ สมฺปทมิทํ ปาปกํ โลภธมฺมํ วทามิ, กิญฺหิ ปโร ปรสฺส กริสฺสตี”ติ.\csromanspacer{} อยํ โข โลหิจฺจ ปฐโม สตฺถา, โย โลเก โจทนารโห, โย จ ปเนวรูปํ สตฺถารํ โจเทติ, สา โจทนา ภูตา ตจฺฉา ธมฺมิกา อนวชฺชา.}

\proseitem{514}{ปุน จปรํ โลหิจฺจ อิเธกจฺโจ สตฺถา ยสฺสตฺถาย อคารสฺมา อนคาริยํ ปพฺพชิโต โหติ, สฺวาสฺส สามญฺญตฺโถ อนนุปฺปตฺโต โหติ.\csromanspacer{} โส ตํ สามญฺญตฺถํ อนนุปาปุณิตฺวา สาวกานํ ธมฺมํ เทเสติ “อิทํ โว หิตาย อิทํ โว สุขายา”ติ.\csromanspacer{} ตสฺส สาวกา สุสฺสูสนฺติ, โสตํ โอทหนฺติ, อญฺญา จิตฺตํ อุปฏฺฐเปนฺติ, น จ โวกฺกมฺม สตฺถุสาสนา วตฺตนฺติ.\csromanspacer{} โส เอวมสฺส โจเทตพฺโพ “อายสฺมา โข ยสฺสตฺถาย อคารสฺมา อนคาริยํ ปพฺพชิโต, โส เต สามญฺญตฺโถ อนนุปฺปตฺโต, ตํ ตฺวํ สามญฺญตฺถํ อนนุปาปุณิตฺวา สาวกานํ ธมฺมํ เทเสสิ ‘อิทํ โว หิตาย อิทํ โว สุขายา’ติ.\csromanspacer{} ตสฺส เต สาวกา สุสฺสูสนฺติ, โสตํ โอทหนฺติ, อญฺญา จิตฺตํ อุปฏฺฐเปนฺติ, น จ โวกฺกมฺม สตฺถุสาสนา วตฺตนฺติ.\csromanspacer{} เสยฺยถาปิ นาม สกํ เขตฺตํ โอหาย ปรํ เขตฺตํ นิทฺทายิตพฺพํ \csromanfolio{220}มญฺเญยฺย, เอวํ สมฺปทมิทํ ปาปกํ โลภธมฺมํ วทามิ, กิญฺหิ ปโร ปรสฺส กริสฺสตี”ติ.\csromanspacer{} อยํ โข โลหิจฺจ ทุติโย สตฺถา, โย โลเก โจทนารโห, โย จ ปเนวรูปํ สตฺถารํ โจเทติ, สา โจทนา ภูตา ตจฺฉา ธมฺมิกา อนวชฺชา.}

\proseitem{515}{ปุน จปรํ โลหิจฺจ อิเธกจฺโจ สตฺถา ยสฺสตฺถาย อคารสฺมา อนคาริยํ ปพฺพชิโต โหติ, สฺวาสฺส สามญฺญตฺโถ อนุปฺปตฺโต โหติ.\csromanspacer{} โส ตํ สามญฺญตฺถํ อนุปาปุณิตฺวา สาวกานํ ธมฺมํ เทเสติ “อิทํโว หิตาย อิทํ โว สุขายา”ติ.\csromanspacer{} ตสฺส สาวกา น สุสฺสูสนฺติ, น โสตํ โอทหนฺติ, น อญฺญา จิตฺตํ อุปฏฺฐเปนฺติ, โวกฺกมฺม จ สตฺถุสาสนา วตฺตนฺติ.\csromanspacer{} โส เอวมสฺส โจเทตพฺโพ “อายสฺมา โข ยสฺสตฺถาย อคารสฺมา อนคาริยํ ปพฺพชิโต, โส เต สามญฺญตฺโถ อนุปฺปตฺโต, ตํ ตฺวํ สามญฺญตฺถํ อนุปาปุณิตฺวา สาวกานํ ธมฺมํ เทเสติ ‘อิทํ โว หิตาย อิทํ โว สุขายา’ติ.\csromanspacer{} ตสฺส เต สาวกา น สุสฺสูสนฺติ, น โสตํ โอทหนฺติ, น อญฺญา จิตฺตํ อุปฏฺฐเปนฺติ, โวกฺกมฺม จ สตฺถุสาสนา วตฺตนฺติ.\csromanspacer{} เสยฺยถาปิ นาม ปุราณํ พนฺธนํ ฉินฺทิตฺวา อญฺญํ นวํ พนฺธนํ กเรยฺย, เอวํ สมฺปทมิทํ ปาปกํ โลภธมฺมํ วทามิ, กิญฺหิ ปโร ปรสฺส กริสฺสตี”ติ.\csromanspacer{} อยํ โข โลหิจฺจ ตติโย สตฺถา, โย โลเก โจทนารโห, โย จ ปเนวรูปํ สตฺถารํ โจเทติ, สา โจทนา ภูตา ตจฺฉา ธมฺมิกา อนวชฺชา.\csromanspacer{} อิเม โข โลหิจฺจ ตโย สตฺถาโร, เย โลเก โจทนารหา, โย จ ปเนวรูเป สตฺถาโร โจเทติ, สา โจทนา ภูตา ตจฺฉา ธมฺมิกา อนวชฺชาติ.}

\csromanmatikaanchor{mtk.136}
\csromantocmarkref{section}{นโจทนารหสตฺถุ}{mtk.136}
\bookmark[dest={mtk.136},level=1]{นโจทนารหสตฺถุ}
\csromanheader{1}{\textbf{นโจทนารหสตฺถุ}}
\proseitem{516}{เอวํ วุตฺเต โลหิจฺโจ พฺราหฺมโณ ภควนฺตํ เอตทโวจ “อตฺถิ ปน โภ โคตม โกจิ สตฺถา, โย โลเก น โจทนารโห”ติ.\csromanspacer{} อตฺถิ โข โลหิจฺจ สตฺถา, โย โลเก นโจทนารโหติ.\csromanspacer{} กตโม ปน โส โภ โคตม สตฺถา, โย โลเก นโจทนารโหติ.}

\prose{อิธ โลหิจฺจ ตถาคโต โลเก อุปฺปชฺชติ อรหํ สมฺมาสมฺพุทฺโธ ฯเปฯ. (ยถา สามญฺญผเล, เอวํ วิตฺถาเรตพฺพํ.) เอวํ โข โลหิจฺจ ภิกฺขุ}

\vspace{\tipitakaparskip}
\csromancenter{โลหิจฺจสุตฺต สีลสมฺปนฺโน โหติ ฯเปฯ ปฐมํ ฌานํ อุปสมฺปชฺช วิหรติ.\csromanspacer{} ยสฺมิํ โข โลหิจฺจ สตฺถริ สาวโก เอวรูปํ อุฬารํ วิเสสํ อธิคจฺฉติ.\csromanspacer{} อยมฺปิ โข โลหิจฺจ สตฺถา, โย โลเก นโจทนารโห, โย จ ปเนวรูปํ สตฺถารํ โจเทติ, สา โจทนา อภูตา อตจฺฉา อธมฺมิกา สาวชฺชา.}
\prose{\csromanfolio{221}ฯเปฯ ทุติยํ ฌานํ.\csromanspacer{} ตติยํ ฌานํ.\csromanspacer{} จตุตฺถํ ฌานํ อุปสมฺปชฺช วิหรติ.\csromanspacer{} ยสฺมิํ โข โลหิจฺจ สตฺถริ สาวโก เอวรูปํ อุฬารํ วิเสสํ อธิคจฺฉติ.\csromanspacer{} อยมฺปิ โข โลหิจฺจ สตฺถา, โย โลเก นโจทนารโห, โย จ ปเนวรูปํ สตฺถารํ โจเทติ, สา โจทนา อภูตา อตจฺฉา อธมฺมิกา สาวชฺชา.}

\prose{ฯเปฯ ญาณทสฺสนาย จิตฺตํ อภินีหรติ อภินินฺนาเมติ.\csromanspacer{} ยสฺมิํ โข โลหิจฺจ สตฺถริ สาวโก เอวรูปํ อุฬารํ วิเสสํ อธิคจฺฉติ.\csromanspacer{} อยมฺปิ โข โลหิจฺจ สตฺถา, โย โลเก นโจทนารโห, โย จ ปเนวรูปํ สตฺถารํ โจเทติ, สา โจทนา อภูตา อตจฺฉา อธมฺมิกา สาวชฺชา.}

\prose{ฯเปฯ นาปรํ อิตฺถตฺตายาติ ปชานาติ.\csromanspacer{} ยสฺมิํ โข โลหิจฺจ สตฺถริ สาวโก เอวรูปํ อุฬารํ วิเสสํ อธิคจฺฉติ.\csromanspacer{} อยมฺปิ โข โลหิจฺจ สตฺถา, โย โลเก นโจทนารโห, โย จ ปเนวรูปํ สตฺถารํ โจเทติ, สา โจทนา อภูตา อตจฺฉา อธมฺมิกา สาวชฺชาติ.}

\proseitemwithcloser{517}{เอวํ วุตฺเต โลหิจฺโจ พฺราหฺมโณ ภควนฺตํ เอตทโวจ “เสยฺยถาปิ โภ โคตม ปุริโส ปุริสํ นรกปปาตํ ปตนฺตํ เกเสสุ คเหตฺวา อุทฺธริตฺวา ถเล ปติฏฺฐเปยฺย, เอวเมวาหํ โภตา โคตเมน นรกปปาตํ ปปตนฺโต อุทฺธริตฺวา ถเล ปติฏฺฐาปิโต.\csromanspacer{} อภิกฺกนฺตํ โภ โคตม, อภิกฺกนฺตํ โภ โคตม, เสยฺยถาปิ โภ โคตม นิกฺกุชฺชิตํ วา อุกฺกุชฺเชยฺย, ปฏิจฺฉนฺนํ วา วิวเรยฺย, มูฬฺหสฺส วา มคฺคํ อาจิกฺเขยฺย, อนฺธกาเร วา เตลปชฺโชตํ ธาเรยฺย ‘จกฺขุมนฺโต รูปานิ ทกฺขนฺตี’ติ.\csromanspacer{} เอวเมวํ โภตา โคตเมน อเนกปริยาเยน ธมฺโม ปกาสิโต.\csromanspacer{} เอสาหํ ภวนฺตํ โคตมํ สรณํ คจฺฉามิ ธมฺมญฺจ ภิกฺขุสํฆญฺจ, อุปาสกํ มํ ภวํ โคตโม ธาเรตุ อชฺชตคฺเค ปาณุเปตํ สรณํ คตนฺติ.}{โลหิจฺจสุตฺตํ นิฏฺฐิตํ ทฺวาทสมํ.}
\csromanmatikaanchor{mtk.137}
\csromantocmarknopage{chapter}{13. เตวิชฺชสุตฺต}{mtk.137}
\bookmark[dest={mtk.137},level=0]{13. เตวิชฺชสุตฺต}
\setcsromanheadcha{13. เตวิชฺชสุตฺต}
\setcsromanheadhclear{}
\chapterheadpageread{13. \textbf{เตวิชฺชสุตฺต}}
\proseitem{518}{\csromanfolio{222}เอวํ เม สุตํ\csromandash{}เอกํ สมยํ ภควา โกสเลสุ จาริกํ จรมาโน มหตา ภิกฺขุสํเฆน สทฺธิํ ปญฺจมตฺเตหิ ภิกฺขุสเตหิ เยน มนสากฏํ นาม โกสลานํ พฺราหฺมณคาโม ตทวสริ.\csromanspacer{} ตตฺร สุทํ ภควา มนสากเฏ วิหรติ อุตฺตเรน มนสากฏสฺส อจิรวติยา นทิยา ตีเร อมฺพวเน.}

\proseitem{519}{เตน โข ปน สมเยน สมฺพหุลา อภิญฺญาตา อภิญฺญาตา พฺราหฺมณมหาสาลา มนสากเฏ ปฏิวสนฺติ.\csromanspacer{} เสยฺยถิทํ, จงฺกี พฺราหฺมโน ตารุกฺโข พฺราหฺมโณ โปกฺขรสาติ พฺราหฺมโณ ชาณุโสณิ พฺราหฺมโณ โตเทยฺโย พฺราหฺมโณ อญฺเญ จ อภิญฺญาตา อภิญฺญาตา พฺราหฺมณมหาสาลา.}

\proseitem{520}{อถ โข วาเสฏฺฐภารทฺวาชานํ มาณวานํ ชงฺฆวิหารํ อนุจงฺกมนฺตานํ อนุวิจรนฺตานํ มคฺคามคฺเค กถา อุทปาทิ.\csromanspacer{} อถ โข วาเสฏฺโฐ มาณโว เอวมาห “อยเมว อุชุมคฺโค, อยมญฺชสายโน นิยฺยานิโก นิยฺยาติ ตกฺกรสฺส พฺรหฺมสหพฺยตาย, ยฺวายํ อกฺขาโต พฺราหฺมเณน โปกฺขรสาตินา”ติ.\csromanspacer{} ภารทฺวาโชปิ มาณโว เอวมาห “อยเมว อุชุมคฺโค, อยมญฺชสายโน นิยฺยานิโก นิยฺยาติ ตกฺกรสฺส พฺรหฺมสหพฺยตาย, ยฺวายํ อกฺขาโต พฺราหฺมเณน ตารุกฺเขนา”ติ.\csromanspacer{} เนว โข อสกฺขิ วาเสฏฺโฐ มาณโว ภารทฺวาชํ มาณวํ สญฺญาเปตุํ, น ปน อสกฺขิ ภารทฺวาโช มาณโว วาเสฏฺฐํ มาณวํ สญฺญาเปตุํ.}

\proseitem{521}{อถ โข วาเสฏฺโฐ มาณโว ภารทฺวาชํ มาณวํ อามนฺเตสิ “อยํ โข ภารทฺวาช สมโณ โคตโม สกฺยปุตฺโต สกฺยกุลา ปพฺพชิโต มนสากเฏ วิหรติ อุตฺตเรน มนสากฏสฺส อจิรวติยา นทิยา ตีเร อมฺพวเน.\csromanspacer{} ตํ โข ปน ภวนฺตํ โคตมํ เอวํ กลฺยาโณ กิตฺติสทฺโท อพฺภุคฺคโต “อิติปิ โส ภควา อรหํ สมฺมาสมฺพุทฺโธ วิชฺชาจรณสมฺปนฺโน สุคโต โลกวิทู อนุตฺตโร ปุริสทมฺมสารถิ}

\vspace{\tipitakaparskip}
\csromancenter{เตวิชฺชสุตฺต สตฺถา เทวมนุสฺสานํ พุทฺโธ ภควา”ติ.\csromanspacer{} อายาม โภ ภารทฺวาช เยน สมโณ โคตโม เตนุปสงฺกมิสฺสาม, อุปสงฺกมิตฺวา เอตมตฺถํ สมณํ โคตมํ ปุจฺฉิสฺสาม.\csromanspacer{} ยถา โน สมโณ โคตโม พฺยากริสฺสติ, ตถา นํ ธาเรสฺสามาติ. “เอวํ โภ”ติ โข ภารทฺวาโช มาณโว วาเสฏฺฐสฺส มาณวสฺส ปจฺจสฺโสสิ.}
\csromanmatikaanchor{mtk.138}
\csromantocmarkref{section}{มคฺคามคฺคกถา}{mtk.138}
\bookmark[dest={mtk.138},level=1]{มคฺคามคฺคกถา}
\csromanheader{1}{\textbf{มคฺคามคฺคกถา}}
\proseitem{522}{\csromanfolio{223}อถ โข วาเสฏฺฐภารทฺวาชา มาณวา เยน ภควา เตนุปสงฺกมิํสุ, อุปสงฺกมิตฺวา ภควตา สทฺธิํ สมฺโมทิํสุ, สมฺโมทนียํ กถํ สารณียํ วีติสาเรตฺวา เอกมนฺตํ นิสีทิํสุ.\csromanspacer{} เอกมนฺตํ นิสินฺโน โข วาเสฏฺโฐ มาณโว ภควนฺตํ เอตทโวจ “อิธ โภ โคตม อมฺหากํ ชงฺฆวิหารํ อนุจงฺกมนฺตานํ อนุวิจรนฺตานํ มคฺคามคฺเค กถา อุทปาทิ.\csromanspacer{} อหํ เอวํ วทามิ ‘อยเมว อุชุมคฺโค อยมญฺชสายโน นิยฺยานิโก นิยฺยาติ ตกฺกรสฺส พฺรหฺมสหพฺยตาย, ยฺวายํ อกฺขาโต พฺราหฺมเณน โปกฺขรสาตินา’ติ.\csromanspacer{} ภารทฺวาโช มาณโว เอวมาห ‘อยเมว อุชุมคฺโค อยมญฺชสายโน นิยฺยานิโก นิยฺยาติ ตกฺกรสฺส พฺรหฺมสหพฺยตาย, ยฺวายํ อกฺขาโต พฺราหฺมเณน ตารุกฺเขนา’ติ.\csromanspacer{} เอตฺถ โภ โคตม อตฺเถว วิคฺคโห, อตฺถิ วิวาโท, อตฺถิ นานาวาโท”ติ.}

\proseitem{523}{อิติ กิร วาเสฏฺฐ ตฺวํ เอวํ วเทสิ “อยเมว อุชุมคฺโค อยมญฺชสายโน นิยฺยานิโก นิยฺยาติ ตกฺกรสฺส พฺรหฺมสหพฺยตาย, ยฺวายํ อกฺขาโต พฺราหฺมเณน โปกฺขรสาตินา”ติ.\csromanspacer{} ภารทฺวาโช มาณโว เอวมาห “อยเมว อุชุมคฺโค อยมญฺชสายโน นิยฺยานิโก นิยฺยาติ ตกฺกรสฺส พฺรหฺมสหพฺยตาย, ยฺวายํ อกฺขาโต พฺราหฺมเณน ตารุกฺเขนา”ติ.\csromanspacer{} อถ กิสฺมิํ ปน โว วาเสฏฺฐ วิคฺคโห, กิสฺมิํ วิวาโท, กิสฺมิํ นานาวาโทติ.}

\proseitem{524}{มคฺคามคฺเค โภ โคตม, กิญฺจาปิ โภ โคตม พฺราหฺมณา นานามคฺเค ปญฺญเปนฺติ, อทฺธริยา พฺราหฺมณา ติตฺติริยา พฺราหฺมณา ฉนฺโทกา พฺราหฺมณา พวฺหาริชฺฌา พฺราหฺมณา, อถ โข สพฺพานิ ตานิ นิยฺยานิกา นิยฺยนฺติ ตกฺกรสฺส พฺรหฺมสหพฺยตาย.}

\prose{\csromanfolio{224}เสยฺยถาปิ โภ โคตม คามสฺส วา นิคมสฺส วา อวิทูเร พหูนิ เจปิ นานามคฺคานิ ภวนฺติ, อถ โข สพฺพานิ ตานิ คามสโมสรณานิ ภวนฺติ.\csromanspacer{} เอวเมว โข โภ โคตม กิญฺจาปิ พฺราหฺมณา นานามคฺเค ปญฺญาเปนฺติ, อทฺธริยา พฺราหฺมณา ติตฺติริยา พฺราหฺมณา ฉนฺโทกา พฺราหฺมณา พวฺหาริชฺฌา พฺราหฺมณา, อถ โข สพฺพานิ ตานิ นิยฺยานิกา นิยฺยนฺติ ตกฺกรสฺส พฺรหฺมสหพฺยตายาติ.}

\csromanmatikaanchor{mtk.139}
\csromantocmarkref{section}{วาเสฏฺฐมาณวานุโยค}{mtk.139}
\bookmark[dest={mtk.139},level=1]{วาเสฏฺฐมาณวานุโยค}
\csromanheader{1}{วาเสฏฺฐ\-มาณวา\-นุโยค}
\proseitem{525}{“นิยฺยนฺตี”ติ วาเสฏฺฐ วเทสิ. “นิยฺยนฺตี”ติ โภ โคตม วทามิ. “นิยฺยนฺตี”ติ วาเสฏฺฐ วเทสิ. “นิยฺยนฺตี”ติ โภ โคตม วทามิ. “นิยฺยนฺตี”ติ วาเสฏฺฐ วเทสิ. “นิยฺยนฺตี”ติ โภ โคตม วทามิ.}

\prose{กิํ ปน วาเสฏฺฐ อตฺถิ โกจิ เตวิชฺชานํ พฺราหฺมณานํ เอกพฺราหฺมโณปิ, เยน พฺรหฺมา สกฺขิทิฏฺโฐติ.\csromanspacer{} โน หิทํ โภ โคตม.}

\prose{กิํ ปน วาเสฏฺฐ อตฺถิ โกจิ เตวิชฺชานํ พฺราหฺมณานํ เอกาจริโยปิ, เยน พฺรหฺมา สกฺขิทิฏฺโฐติ.\csromanspacer{} โน หิทํ โภ โคตม.}

\prose{กิํ ปน วาเสฏฺฐ อตฺถิ โกจิ เตวิชฺชานํ พฺราหฺมณานํ เอกาจริยปาจริโยปิ, เยน พฺรหฺมา สกฺขิทิฏฺโฐติ.\csromanspacer{} โน หิทํ โภ โคตม.}

\prose{กิํ ปน วาเสฏฺฐ อตฺถิ โกจิ เตวิชฺชานํ พฺราหฺมณานํ ยาว สตฺตมา อาจริยามหยุคา\footnote{สตฺตมาจริยมหยุคา (สฺยา)}, เยน พฺรหฺมา สกฺขิทิฏฺโฐติ.\csromanspacer{} โน หิทํ โภ โคตม.}

\proseitem{526}{กิํ ปน วาเสฏฺฐ เยปิ เตวิชฺชานํ พฺราหฺมณานํ ปุพฺพกา อิสโย มนฺตานํ กตฺตาโร มนฺตานํ ปวตฺตาโร, เยสมิทํ เอตรหิ เตวิชฺชา พฺราหฺมณา โปราณํ มนฺตปทํ คีตํ ปวุตฺตํ สมิหิตํ\footnote{สมีหิตํ (สฺยา)}, ตทนุคายนฺติ, ตทนุภาสนฺติ, ภาสิตมนุภาสนฺติ, วาจิตมนุวาเจนฺติ.\csromanspacer{} เสยฺยถิทํ, อฏฺฐโก วามโก วามเทโว เวสฺสามิตฺโต ยมตคฺคิ องฺคีรโส ภารทฺวาโช วาเสฏฺโฐ กสฺสโป ภคุ.\csromanspacer{} เตปิ เอวมาหํสุ “มยเมตํ ชานาม, มยเมตํ ปสฺสาม, ยตฺถ วา พฺรหฺมา, เยน วา พฺรหฺมา, ยติํ วา พฺรหฺมา”ติ.\csromanspacer{} โน หิทํ โภ โคตม.}

\csromanheader{2}{13. เตวิชฺชสุตฺต}
\proseitem{527}{\csromanfolio{225}อิติ กิร วาเสฏฺฐ นตฺถิ โกจิ เตวิชฺชานํ พฺราหฺมณานํ เอกพฺราหฺมโณปิ, เยน พฺรหฺมา สกฺขิทิฏฺโฐ.\csromanspacer{} นตฺถิ โกจิ เตวิชฺชานํ พฺราหฺมณานํ เอกาจริโยปิ, เยน พฺรหฺมา สกฺขิทิฏฺโฐ.\csromanspacer{} นตฺถิ โกจิ เตวิชฺชานํ พฺราหฺมณานํ เอกาจริยปาจริโยปิ, เยน พฺรหฺมา สกฺขิทิฏฺโฐ.\csromanspacer{} นตฺถิ โกจิ เตวิชฺชานํ พฺราหฺมณานํ ยาว สตฺตมา อาจริยามหยุคา, เยน พฺรหฺมา สกฺขิทิฏฺโฐ.\csromanspacer{} เยปิ กิร เตวิชฺชานํ พฺราหฺมณานํ ปุพฺพกา อิสโย มนฺตานํ กตฺตาโร มนฺตานํ ปวตฺตาโร, เยสมิทํ เอตรหิ เตวิชฺชา พฺราหฺมณา โปราณํ มนฺตปทํ คีตํ ปวุตฺตํ สมิหิตํ, ตทนุคายนฺติ, ตทนุภาสนฺติ, ภาสิตมนุภาสนฺติ, วาจิตมนุวาเจนฺติ.\csromanspacer{} เสยฺยถิทํ,อฏฺฐโก วามโก วามเทโว เวสฺสามิตฺโต ยมตคฺคิ องฺคีรโส ภารทฺวาโช วาเสฏฺโฐ กสฺสโป ภคุ.\csromanspacer{} เตปิ น เอวมาหํสุ “มยเมตํ ชานาม, มยเมตํ ปสฺสาม, ยตฺถ วา พฺรหฺมา, เยน วา พฺรหฺมา, ยหิํ วา พฺรหฺมา”ติ.\csromanspacer{} เตว เตวิชฺชา พฺราหฺมณา เอวมาหํสุ “ยํ น ชานาม, ยํ น ปสฺสาม, ตสฺส สหพฺยตาย มคฺคํ เทเสม, อยเมว อุชุมคฺโค อยมญฺชสายโน นิยฺยานิโก นิยฺยาติ ตกฺกรสฺส พฺรหฺมสหพฺยตายา”ติ.}

\proseitem{528}{ตํ กิํ มญฺญสิ วาเสฏฺฐ, นนุ เอวํ สนฺเต เตวิชฺชานํ พฺราหฺมณานํ อปฺปาฏิหีรกตํ ภาสิตํ สมฺปชฺชตีติ.\csromanspacer{} อทฺธา โข โภ โคตม เอวํ สนฺเต เตวิชฺชานํ พฺราหฺมณานํ อปฺปาฏิหีรกตํ ภาสิตํ สมฺปชฺชตีติ.}

\prose{สาธุ วาเสฏฺฐ, เต วต\footnote{เตว (ก)} วาเสฏฺฐ เตวิชฺชา พฺราหฺมณา ยํ น ชานนฺติ, ยํ น ปสฺสนฺติ, ตสฺส สหพฺยตาย มคฺคํ เทเสสฺสนฺติ “อยเมว อุชุมคฺโค อยมญฺชสายโน นิยฺยานิโก นิยฺยาติ ตกฺกรสฺส พฺรหฺมสหพฺยตายา”ติ, เนตํ ฐานํ วิชฺชติ.}

\proseitem{529}{เสยฺยถาปิ วาเสฏฺฐ อนฺธเวณิ ปรมฺปรสํสตฺตา ปุริโมปิ น ปสฺสติ, มชฺฌิโมปิ น ปสฺสติ, ปจฺฉิโมปิ น ปสฺสติ.\csromanspacer{} เอวเมว โข วาเสฏฺฐ อนฺธเวณูปมํ มญฺเญ เตวิชฺชานํ พฺราหฺมณานํ ภาสิตํ, ปุริโมปิ น ปสฺสติ, มชฺฌิโมปิ น ปสฺสติ, ปจฺฉิโมปิ น ปสฺสติ.\csromanspacer{} เตสมิทํ เตวิชฺชานํ พฺราหฺมณานํ ภาสิตํ หสฺสกญฺเญว สมฺปชฺชติ, นามกญฺเญว สมฺปชฺชติ, ริตฺตกญฺเญว สมฺปชฺชติ, ตุจฺฉกญฺเญว สมฺปชฺชติ.}

\proseitem{530}{\csromanfolio{226}ตํ กิํ มญฺญสิ วาเสฏฺฐ, ปสฺสนฺติ เตวิชฺชา พฺราหฺมณา จนฺทิมสูริเย, อญฺเญ จาปิ พหุชนา.\csromanspacer{} ยโต จ จนฺทิมสูริยา อุคฺคจฺฉนฺติ, ยตฺถ จ โอคจฺฉนฺติ, อายาจนฺติ โถมยนฺติ ปญฺชลิกา นมสฺสมานา อนุปริวตฺตนฺตีติ.}

\prose{เอวํ โภ โคตม, ปสฺสนฺติ เตวิชฺชา พฺราหฺมณา จนฺทิมสูริเย, อญฺเญ จาปิ พหุชนา.\csromanspacer{} ยโต จ จนฺทิมสูริยา อุคฺคจฺฉนฺติ, ยตฺถ จ โอคจฺฉนฺติ, อายาจนฺติ โถมยนฺติ ปญฺชลิกา นมสฺสมานา อนุปริวตฺตนฺตีติ.}

\proseitem{531}{ตํ กิํ มญฺญสิ วาเสฏฺฐ, ยํ ปสฺสนฺติ เตวิชฺชา พฺราหฺมณา จนฺทิมสูริเย, อญฺเญ จาปิ พหุชนา.\csromanspacer{} ยโต จ จนฺทิมสูริยา อุคฺคจฺฉนฺติ, ยตฺถ จ โอคจฺฉนฺติ, อายาจนฺติ โถมยนฺติ ปญฺชลิกา นมสฺสมานา อนุปริวตฺตนฺติ.\csromanspacer{} ปโหนฺติ เตวิชฺชา พฺราหฺมณา จนฺทิมสูริยานํ สหพฺยตาย มคฺคํ เทเสตุํ “อยเมว อุชุมคฺโค อยมญฺชสายโน นิยฺยานิโก นิยฺยาติ ตกฺกรสฺส จนฺทิมสูริยานํ สหพฺยตายา”ติ.\csromanspacer{} โน หิทํ โภ โคตม.}

\prose{อิติ กิร วาเสฏฺฐ ยํ ปสฺสนฺติ เตวิชฺชา พฺราหฺมณา จนฺทิมสูริเย, อญฺเญ จาปิ พหุชนา.\csromanspacer{} ยโต จ จนฺทิมสูริยา อุคฺคจฺฉนฺติ, ยตฺถ จ โอคจฺฉนฺติ, อายาจนฺติ โถมยนฺติ ปญฺชลิกา นมสฺสมานา อนุปริวตฺตนฺติ.\csromanspacer{} เตสมฺปิ นปฺปโหนฺติ จนฺทิมสูริยานํ สหพฺยตาย มคฺคํ เทเสตุํ “อยเมว อุชุมคฺโค อยมญฺชสายโน นิยฺยานิโก นิยฺยาติ ตกฺกรสฺส จนฺทิมสูริยานํ สหพฺยตายา”ติ.}

\proseitem{532}{อิติ ปน\footnote{กิํ ปน (สี, สฺยา, อิ)} น กิร เตวิชฺเชหิ พฺราหฺมเณหิ พฺรหฺมา สกฺขิทิฏฺโฐ.\csromanspacer{} นปิ กิร เตวิชฺชานํ พฺราหฺมณานํ อาจริเยหิ พฺรหฺมา สกฺขิทิฏฺโฐ.\csromanspacer{} นปิ กิร เตวิชฺชานํ พฺราหฺมณานํ อาจริยปาจริเยหิ พฺรหฺมา สกฺขิทิฏฺโฐ.\csromanspacer{} นปิ กิร เตวิชฺชานํ พฺราหฺมณานํ ยาว สตฺตมา\footnote{สตฺตเมหิ (?)} อาจริยามหยุเคหิ พฺรหฺมา สกฺขิทิฏฺโฐ.\csromanspacer{} เยปิ กิร เตวิชฺชานํ พฺราหฺมณานํ ปุพฺพกา อิสโย มนฺตานํ กตฺตาโร มนฺตานํ ปวตฺตาโร, เยสมิทํ เอตรหิ เตวิชฺชา พฺราหฺมณา โปราณํ มนฺตปทํ คีตํ ปวุตฺตํ สมิหิตํ, ตทนุคายนฺติ, ตทนุภาสนฺติ, ภาสิตมนุภาสนฺติ, วาจิตมนุวาเจนฺติ.\csromanspacer{} เสยฺยถิทํ, อฏฺฐโก วามโก}

\vspace{\tipitakaparskip}
\csromancenter{เตวิชฺชสุตฺต วามเทโว เวสฺสามิตฺโต ยมตคฺคิ องฺคีรโส ภารทฺวาโช วาเสฏฺโฐ กสฺสโป ภคุ.\csromanspacer{} เตปิ น เอวมาหํสุ “มยเมตํ ชานาม, มยเมตํ ปสฺสาม, ยตฺถ วา พฺรหฺมา, เยน วา พฺรหฺมา, ยหิํ วา พฺรหฺมา”ติ.\csromanspacer{} เตว เตวิชฺชา พฺราหฺมณา เอวมาหํสุ “ยํ น ชานาม, ยํ น ปสฺสาม, ตสฺส สหพฺยตาย มคฺคํ เทเสม, อยเมว อุชุมคฺโค อยมญฺชสายโน นิยฺยานิโก นิยฺยาติ ตกฺกรสฺส พฺรหฺมสหพฺยตายา”ติ.}
\proseitem{533}{\csromanfolio{227}ตํ กิํ มญฺญสิ วาเสฏฺฐ, นนุ เอวํ สนฺเต เตวิชฺชานํ พฺราหฺมณานํ อปฺปาฏิหีรกตํ ภาสิตํ สมฺปชฺชตีติ.\csromanspacer{} อทฺธา โข โภ โคตม เอวํ สนฺเต เตวิชฺชานํ พฺราหฺมณานํ อปฺปาฏิหีรกตํ ภาสิตํ สมฺปชฺชตีติ.}

\prose{สาธุ วาเสฏฺฐ, เต วต วาเสฏฺฐ เตวิชฺชา พฺราหฺมณา ยํ น ชานนฺติ, ยํ น ปสฺสนฺติ, ตสฺส สหพฺยตาย มคฺคํ เทเสสฺสนฺติ “อยเมว อุชุมคฺโค อยมญฺชสายโน นิยฺยานิโก นิยฺยาติ ตกฺกรสฺส พฺรหฺมสหพฺยตายา”ติ, เนตํ ฐานํ วิชฺชติ.}

\csromanmatikaanchor{mtk.140}
\csromantocmarkref{section}{ชนปทกลฺยาณีอุปมา}{mtk.140}
\bookmark[dest={mtk.140},level=1]{ชนปทกลฺยาณีอุปมา}
\csromanheader{1}{\textbf{ชนปทกลฺยาณีอุปมา}}
\proseitem{534}{เสยฺยถาปิ วาเสฏฺฐ ปุริโส เอวํ วเทยฺย “อหํ ยา อิมสฺมิํ ชนปเท ชนปทกลฺยาณี, ตํ อิจฺฉามิ ตํ กาเมมี”ติ.\csromanspacer{} ตเมนํ เอวํ วเทยฺยุํ “อมฺโภ ปุริส ยํ ตฺวํ ชนปทกลฺยาณิํ อิจฺฉสิ กาเมสิ, ชานาสิ ตํ ชนปทกลฺยาณิํ ขตฺติยี วา พฺราหฺมณี วา เวสฺสี วา สุทฺที วา”ติ, อิติ ปุฏฺโฐ “โน”ติ วเทยฺย.}

\prose{ตเมนํ เอวํ วเทยฺยุํ “อมฺโภ ปุริส ยํ ตฺวํ ชนปทกลฺยาณิํ อิจฺฉาสิ กาเมสิ, ชานาสิ ตํ ชนปทกลฺยาณิํ เอวํนามา เอวํโคตฺตาติ วา, ทีฆา วา รสฺสา วา มชฺฌิมา วา กาฬี วา สามา วา มงฺคุรจฺฉวี วาติ, อมุกสฺมิํ คาเม วา นิคเม วา นคเร วาติ, อิติ ปุฏฺโฐ ‘โน’ติ วเทยฺย.\csromanspacer{} ตเมนํ เอวํ วเทยฺยุํ “อมฺโภ ปุริส ยํ ตฺวํ น ชานาสิ น ปสฺสสิ, ตํ ตฺวํ อิจฺฉสิ กาเมสี”ติ, อิติ ปุฏฺโฐ “อามา”ติ วเทยฺย.}

\proseitem{535}{ตํ กิํ มญฺญสิ วาเสฏฺฐ, นนุ เอวํ สนฺเต ตสฺส ปุริสสฺส อปฺปาฏิหีรกตํ ภาสิตํ สมฺปชฺชตีติ.\csromanspacer{} อทฺธา โข โภ โคตม เอวํ สนฺเต ตสฺส ปุริสสฺส อปฺปาฏิหีรกตํ ภาสิตํ สมฺปชฺชตีติ.}

\proseitem{536}{\csromanfolio{228}เอวเมว โข วาเสฏฺฐ น กิร เตวิชฺเชหิ พฺราหฺมเณหิ พฺรหฺมา สกฺขิทิฏฺโฐ.\csromanspacer{} นปิ กิร เตวิชฺชานํ พฺราหฺมณานํ อาจริเยหิ พฺรหฺมา สกฺขิทิฏฺโฐ.\csromanspacer{} นปิ กิร เตวิชฺชานํ พฺราหฺมณานํ อาจริยปาจริเยหิ พฺรหฺมา สกฺขิทิฏฺโฐ.\csromanspacer{} นปิ กิร เตวิชฺชานํ พฺราหฺมณานํ ยาว สตฺตมา อาจริยามหยุเคหิ พฺรหฺมา สกฺขิทิฏฺโฐ.\csromanspacer{} เยปิ กิร เตวิชฺชานํ พฺราหฺมณานํ ปุพฺพกา อิสโย มนฺตานํ กตฺตาโร มนฺตานํ ปวตฺตาโร, เยสมิทํ เอตรหิ เตวิชฺชา พฺราหฺมณา โปราณํ มนฺตปทํ คีตํ ปวุตฺตํ สมิหิตํ, ตทนุคายนฺติ, ตทนุภาสนฺติ, ภาสิตมนุภาสนฺติ, วาจิตมนุวาเจนฺติ.\csromanspacer{} เสยฺยถิทํ, อฏฺฐโก วามโก วามเทโว เวสฺสามิตฺโต ยมตคฺคิ องฺคีรโส ภารทฺวาโช วาเสฏฺโฐ กสฺสโป ภคุ.\csromanspacer{} เตปิ น เอวมาหํสุ “มยเมตํ ชานาม, มยเมตํ ปสฺสาม, ยตฺถ วา พฺรหฺมา, เยน วา พฺรหฺมา, ยหิํ วา พฺรหฺมา”ติ.\csromanspacer{} เตว เตวิชฺชา พฺราหฺมณา เอวมาหํสุ “ยํ น ชานาม, ยํ น ปสฺสาม, ตสฺส สหพฺยตาย มคฺคํ เทเสม, อยเมว อุชุมคฺโค อยมญฺชสายโน นิยฺยานิโก นิยฺยาติ ตกฺกรสฺส พฺรหฺมสหพฺยตายา”ติ.}

\proseitem{537}{ตํ กิํ มญฺญสิ วาเสฏฺฐ, นนุ เอวํ สนฺเต เตวิชฺชานํ พฺราหฺมณานํ อปฺปาฏิหีรกตํ ภาสิตํ สมฺปชฺชตีติ.\csromanspacer{} อทฺธา โข โภ โคตม เอวํ สนฺเต เตวิชฺชานํ พฺราหฺมณานํ อปฺปาฏิหีรกตํ ภาสิตํ สมฺปชฺชตีติ.}

\prose{สาธุ วาเสฏฺฐ, เต วต วาเสฏฺฐ เตวิชฺชา พฺราหฺมณา ยํ น ชานนฺติ, ยํ น ปสฺสนฺติ, ตสฺส สหพฺยตาย มคฺคํ เทเสสฺสนฺติ “อยเมว อุชุมคฺโค อยมญฺชสายโน นิยฺยานิโก นิยฺยาติ ตกฺกรสฺส พฺรหฺมสหพฺยตายา”ติ, เนตํ ฐานํ วิชฺชติ.}

\csromanmatikaanchor{mtk.141}
\csromantocmarkref{section}{นิสฺเสณีอุปมา}{mtk.141}
\bookmark[dest={mtk.141},level=1]{นิสฺเสณีอุปมา}
\csromanheader{1}{\textbf{นิสฺเสณีอุปมา}}
\proseitem{538}{เสยฺยถาปิ วาเสฏฺฐ ปุริโส จาตุมหาปเถ นิสฺเสณิํ กเรยฺย ปาสาทสฺส อาโรหณาย.\csromanspacer{} ตเมนํ เอวํ วเทยฺยุํ “อมฺโภ ปุริส ยสฺส ตฺวํ\footnote{ยํ ตฺวํ (สฺยา)} ปาสาทสฺส อาโรหณาย นิสฺเสณิํ กโรสิ, ชานาสิ ตํ ปาสาทํ ปุรตฺถิมาย วา ทิสาย ทกฺขิณาย วา ทิสาย}

\vspace{\tipitakaparskip}
\csromancenter{เตวิชฺชสุตฺต ปจฺฉิมาย วา ทิสาย อุตฺตราย วา ทิสาย อุจฺโจ วา นีโจ วา มชฺฌิโม วา”ติ, อิติ ปุฏฺโฐ “โน”ติ วเทยฺย.}
\prose{\csromanfolio{229}ตเมนํ เอวํ วเทยฺยุํ “อมฺโภ ปุริส ยํ ตฺวํ น ชานาสิ, น ปสฺสติ, ตสฺส ตฺวํ ปาสาทสฺส อาโรหณาย นิสฺเสณิํ กโรสี”ติ, อิติ ปุฏฺโฐ “อามา”ติ วเทยฺย.}

\proseitem{539}{ตํ กิํ มญฺญสิ วาเสฏฺฐ, นนุ เอวํ สนฺเต ตสฺส ปุริสสฺส อปฺปาฏิหีรกตํ ภาสิตํ สมฺปชฺชตีติ.\csromanspacer{} อทฺธา โข โภ โคตม เอวํ สนฺเต ตสฺส ปุริสสฺส อปฺปาฏิหีรกตํ ภาสิตํ สมฺปชฺชตีติ.}

\proseitem{540}{เอวเมว โข วาเสฏฺฐ น กิร เตวิชฺเชหิ พฺราหฺมเณหิ พฺรหฺมา สกฺขิทิฏฺโฐ, นปิ กิร เตวิชฺชานํ พฺราหฺมณานํ อาจริเยหิ พฺรหฺมา สกฺขิทิฏฺโฐ, นปิ กิร เตวิชฺชานํ พฺราหฺมณานํ อาจริยปาจริเยหิ พฺรหฺมา สกฺขิทิฏฺโฐ, นปิ กิร เตวิชฺชานํ พฺราหฺมณานํ ยาว สตฺตมา อาจริยามหยุเคหิ พฺรหฺมา สกฺขิทิฏฺโฐ, เยปิ กิร เตวิชฺชานํ พฺราหฺมณานํ ปุพฺพกา อิสโย มนฺตานํ กตฺตาโร มนฺตานํ ปวตฺตาโร, เยสมิทํ เอตรหิ เตวิชฺชา พฺราหฺมณา โปราณํ มนฺตปทํ คีตํ ปวุตฺตํ สมิหิตํ, ตทนุคายนฺติ, ตทนุภาสนฺติ, ภาสิตมนุภาสนฺติ, วาจิตมนุวาเจนฺติ.\csromanspacer{} เสยฺยถิทํ, อฏฺฐโก วามโก วามเทโว เวสฺสามิตฺโต ยมตคฺคิ องฺคีรโส ภารทฺวาโช วาเสฏฺโฐ กสฺสโป ภคุ.\csromanspacer{} เตปิ น เอวมหํสุ “มยเมตํ ชานาม, มยเมตํ ปสฺสาม, ยตฺถ วา พฺรหฺมา, เยน วา พฺรหฺมา, ยหิํ วา พฺรหฺมา”ติ.\csromanspacer{} เตว เตวิชฺชา พฺราหฺมณา เอวมาหํสุ “ยํ น ชานาม, ยํ น ปสฺสาม, ตสฺส สหพฺยตาย มคฺคํ เทเสม, อยเมว อุชุมคฺโค อยมญฺชสายโน นิยฺยานิโก นิยฺยาติ ตกฺกรสฺส พฺรหฺมสหพฺยตายา”ติ.}

\proseitem{541}{ตํ กิํ มญฺญสิ วาเสฏฺฐ, นนุ เอวํ สนฺเต เตวิชฺชานํ พฺราหฺมณานํ อปฺปาฏิหีรกตํ ภาสิตํ สมฺปชฺชตีติ.\csromanspacer{} อทฺธา โข โภ โคตม เอวํ สนฺเต เตวิชฺชานํ พฺราหฺมณานํ อปฺปาฏิหีรกตํ ภาสิตํ สมฺปชฺชตีติ.}

\prose{สาธุ วาเสฏฺฐ, เต วต วาเสฏฺฐ เตวิชฺชา พฺราหฺมณา ยํ น ชานนฺติ, ยํ น ปสฺสนฺติ, ตสฺส สหพฺยตาย มคฺคํ เทเสสฺสนฺติ “อยเมว อุชุมคฺโค \csromanfolio{230}อยมญฺชสายโน นิยฺยานิโก นิยฺยาติ ตกฺกรสฺส พฺรหฺมสหพฺยตายา”ติ, เนตํ ฐานํ วิชฺชติ.}

\csromanmatikaanchor{mtk.142}
\csromantocmarkref{section}{อจิรวตีนทีอุปมา}{mtk.142}
\bookmark[dest={mtk.142},level=1]{อจิรวตีนทีอุปมา}
\csromanheader{1}{\textbf{อจิรวตีนทีอุปมา}}
\proseitem{542}{เสยฺยถาปิ วาเสฏฺฐ อยํ อจิรวตี นที ปูรา อุทกสฺส สมติตฺติกา กากเปยฺยา.\csromanspacer{} อถ ปุริโส อาคจฺเฉยฺย ปารตฺถิโก ปารคเวสี ปารคามี ปารํ ตริตุกาโม, โส โอริเม ตีเร ฐิโต ปาริมํ ตีรํ อเวฺหยฺย “เอหิ ปาราปารํ, เอหิ ปาราปารนฺ”ติ.}

\proseitem{543}{ตํ กิํ มญฺญสิ วาเสฏฺฐ, อปิ นุ ตสฺส ปุริสสฺส อวฺหายนเหตุ วา อายาจนเหตุ วา ปตฺถนเหตุ วา อภินนฺทนเหตุ วา อจิรวติยา นทิยา ปาริมํ ตีรํ โอริมํ ตีรํ อาคจฺเฉยฺยาติ.\csromanspacer{} โน หิทํ โภ โคตม.}

\proseitem{544}{เอวเมว โข วาเสฏฺฐ เตวิชฺชา พฺราหฺมณา เย ธมฺมา พฺราหฺมณการกา เต ธมฺเม ปหาย วตฺตมานา, เย ธมฺมา อพฺราหฺมณการกา เต ธมฺเม สมาทาย วตฺตมานา เอวมาหํสุ “อินฺทมวฺหยาม, โสมมวฺหยาม, วรุณมวฺหยาม, อีสานมวฺหยาม, ปชาปติมวฺหยาม, พฺรหฺมมวฺหยาม, มหิทฺธิมวฺหยาม, ยมมวฺหยามา”ติ.}

\prose{เต วต วาเสฏฺฐ เตวิชฺชา พฺราหฺมณา เย ธมฺมา พฺราหฺมณการกา เต ธมฺเม ปหาย วตฺตมานา, เย ธมฺมา อพฺราหฺมณการกา เต ธมฺเม สมาทาย วตฺตมานา อวฺหายนเหตุ วา อายาจนเหตุ วา ปตฺถนเหตุ วา อภินนฺทนเหตุ วา กายสฺส เภทา ปรํ มรณา พฺรหฺมานํ สหพฺยูปคา ภวิสฺสนฺตีติ เนตํ ฐานํ วิชฺชติ.}

\proseitem{545}{เสยฺยถาปิ วาเสฏฺฐ อยํ อจิรวตี นที ปูรา อุทกสฺส สมติตฺติกา กากเปยฺยา.\csromanspacer{} อถ ปุริโส อาคจฺเฉยฺย ปารตฺถิโก ปารคเวสี ปารคามี ปารํ ตริตุกาโม, โส โอริเม ตีเร ทฬฺหาย อนฺทุยา ปจฺฉาพาหํ คาฬฺหพนฺธนํ พทฺโธ.}

\prose{ตํ กิํ มญฺญสิ วาเสฏฺฐ, อปิ นุ โส ปุริโส อจิรวติยา นทิยา โอริมา ตีรา ปาริมํ ตีรํ คจฺเฉยฺยาติ.\csromanspacer{} โน หิทํ โภ โคตม.}

\csromanheader{2}{13. เตวิชฺชสุตฺต}
\proseitem{546}{\csromanfolio{231}เอวเมว โข วาเสฏฺฐ ปญฺจิเม กามคุณา อริยสฺส วินเย อนฺทูติปิ วุจฺจนฺติ, พนฺธนนฺติปิ วุจฺจนฺติ.\csromanspacer{} กตเม ปญฺจ.\csromanspacer{} จกฺขุวิญฺเญยฺยา รูปา อิฏฺฐา กนฺตา มนาปา ปิยรูปา กามูปสํหิตา รชนียา.\csromanspacer{} โสตวิญฺเญยฺยา สทฺทา ฯเปฯ.\csromanspacer{} ฆานวิญฺเญยฺยา คนฺธา.\csromanspacer{} ชิวฺหาวิญฺเญยฺยา รสา.\csromanspacer{} กายวิญฺเญยฺยา โผฏฺฐพฺพา อิฏฺฐา กนฺตา มนาปา ปิยรูปา กามูปสํหิตา รชนียา.}

\prose{อิเม โข วาเสฏฺฐ ปญฺจ กามคุณา อริยสฺส วินเย อนฺทูติปิ วุจฺจนฺติ.\csromanspacer{} พนฺธนนฺติปิ วุจฺจนฺติ.\csromanspacer{} อิเม โข วาเสฏฺฐ ปญฺจ กามคุเณ เตวิชฺชา พฺราหฺมณา คธิตา มุจฺฉิตา อชฺโฌปนฺนา อนาทีนวทสฺสาวิโน อนิสฺสรณปญฺญา ปริภุญฺชนฺติ.\csromanspacer{} เต วต วาเสฏฺฐ เตวิชฺชา พฺราหฺมณา เย ธมฺมา พฺราหฺมณการกา, เต ธมฺเม ปหาย วตฺตมานา, เย ธมฺมา อพฺราหฺมณการกา, เต ธมฺเม สมาทาย วตฺตมานา ปญฺจ กามคุเณ คธิตา มุจฺฉิตา อชฺโฌปนฺนา อนาทีนวทสฺสาวิโน อนิสฺสรณปญฺญา ปริภุญฺชนฺตา กามนฺทุพนฺธนพทฺธา กายสฺส เภทา ปรํ มรณา พฺรหฺมานํ สหพฺยูปคา ภวิสฺสนฺตีติ เนตํ ฐานํ วิชฺชติ.}

\proseitem{547}{เสยฺยถาปิ วาเสฏฺฐ อยํ อจิรวตี นที ปูรา อุทกสฺส สมติตฺติกา กากเปยฺยา.\csromanspacer{} อถ ปุริโส อาคจฺเฉยฺย ปารตฺถิโก ปารคเวสี ปารคามี ปารํ ตริตุกาโม, โส โอริเม ตีเร สสีสํ ปารุปิตฺวา นิปชฺเชยฺย.}

\prose{ตํ กิํ มญฺญสิ วาเสฏฺฐ, อปิ นุ โส ปุริโส อจิรวติยา นทิยา โอริมา ตีรา ปาริมํ ตีรํ คจฺเฉยฺยาติ.\csromanspacer{} โน หิทํ โภ โคตม.}

\proseitem{548}{เอวเมว โข วาเสฏฺฐ ปญฺจิเม นีวรณา อริยสฺส วินเย อาวรณาติปิ วุจฺจนฺติ, นีวรณาติปิ วุจฺจนฺติ, โอนาหนาติปิ วุจฺจนฺติ, ปริโยนาหนาติปิ วุจฺจนฺติ.\csromanspacer{} กตเม ปญฺจ.\csromanspacer{} กามจฺฉนฺทนีวรณํ, พฺยาปาทนีวรณํ, ถินมิทฺธนีวรณํ, อุทฺธจฺจกุกฺกุจฺจนีวรณํ, วิจิกิจฺฉานีวรณํ.\csromanspacer{} อิเม โข วาเสฏฺฐ ปญฺจ นีวรณา อริยสฺส วินเย อาวรณาติปิ วุจฺจนฺติ, นีวรณาติปิ วุจฺจนฺติ, โอนาหนาติปิ วุจฺจนฺติ, ปริโยนาหนาติปิ วุจฺจนฺติ.}

\proseitem{549}{อิเมหิ โข วาเสฏฺฐ ปญฺจหิ นีวรเณหิ เตวิชฺชา พฺราหฺมณา อาวุฏา นิวุฏา โอนทฺธา\footnote{โอผุฏา (สี, ก), โอผุตา (สฺยา)} ปริโยนทฺธา.\csromanspacer{} เต วต วาเสฏฺฐ เตวิชฺชา พฺราหฺมณา \csromanfolio{232}เย ธมฺมา พฺราหฺมณการกา เต ธมฺเม ปหาย วตฺตมานา, เย ธมฺมา อพฺราหฺมณการกา เต ธมฺเม สมาทาย วตฺตมานา ปญฺจหิ นีวรเณหิ อาวุฏา นิวุฏา โอนทฺธา ปริโยนทฺธา\footnote{ปริโยนทฺธา, เต (สฺยา, ก)} กายสฺส เภทา ปรํ มรณา พฺรหฺมานํ สหพฺยูปคา ภวิสฺสนฺตีติ เนตํ ฐานํ วิชฺชติ.}

\csromanmatikaanchor{mtk.143}
\csromantocmarkref{section}{สํสนฺทนกถา}{mtk.143}
\bookmark[dest={mtk.143},level=1]{สํสนฺทนกถา}
\csromanheader{1}{\textbf{สํสนฺทนกถา}}
\proseitem{550}{ตํ กิํ มญฺญสิ วาเสฏฺฐ, กินฺติ เต สุตํ พฺราหฺมณานํ วุทฺธานํ มหลฺลกานํ อาจริยปาจริยานํ ภาสมานานํ, สปริคฺคโห วา พฺรหฺมา อปริคฺคโห วาติ.\csromanspacer{} อปริคฺคโห โภ โคตม.\csromanspacer{} สเวรจิตฺโต วา อเวรจิตฺโต วาติ.\csromanspacer{} อเวรจิตฺโต โภ โคตม.\csromanspacer{} สพฺยาปชฺชจิตฺโต วา อพฺยาปชฺชจิตฺโต วาติ.\csromanspacer{} อพฺยาปชฺชจิตฺโต โภ โคตม.\csromanspacer{} สํกิลิฏฺฐจิตฺโต วา อสํกิลิฏฺฐจิตฺโต วาติ.\csromanspacer{} อสํกิลิฏฺฐจิตฺโต โภ โคตม.\csromanspacer{} วสวตฺตี วา อวสวตฺตี วาติ.\csromanspacer{} วสวตฺตี โภ โคตม.}

\prose{ตํ กิํ มญฺญสิ วาเสฏฺฐ, สปริคฺคหา วา เตวิชฺชา พฺราหฺมณา อปริคฺคหา วาติ.\csromanspacer{} สปริคฺคหา โภ โคตม.\csromanspacer{} สเวรจิตฺตา วา อเวรจิตฺตา วาติ.\csromanspacer{} สเวรจิตฺตา โภ โคตม.\csromanspacer{} สพฺยาปชฺชจิตฺตา วา อพฺยาปชฺชจิตฺตา วาติ.\csromanspacer{} สพฺยาปชฺชจิตฺตา โภ โคตม.\csromanspacer{} สํกิลิฏฺฐจิตฺตา วา อสํกิลิฏฺฐจิตฺตา วาติ.\csromanspacer{} สํกิลิฏฺฐจิตฺตา โภ โคตม.\csromanspacer{} วสวตฺตี วา อวสวตฺตี วาติ.\csromanspacer{} อวสวตฺตี โภ โคตม.}

\proseitem{551}{อิติ กิร วาเสฏฺฐ สปริคฺคหา เตวิชฺชา พฺราหฺมณา อปริคฺคโห พฺรหฺมา, อปิ นุ โข สปริคฺคหานํ เตวิชฺชานํ พฺราหฺมณานํ อปริคฺคเหน พฺรหฺมุนา สทฺธิํ สํสนฺทติ สเมตีติ.\csromanspacer{} โน หิทํ โภ โคตม.\csromanspacer{} สาธุ วาเสฏฺฐ, เต วต วาเสฏฺฐ สปริคฺคหา เตวิชฺชา พฺราหฺมณา กายสฺส เภทา ปรํ มรณา อปริคฺคหสฺส พฺรหฺมุโน สหพฺยูปคา ภวิสฺสนฺตีติ เนตํ ฐานํ วิชฺชติ.}

\prose{อิติ กิร วาเสฏฺฐ สเวรจิตฺตา เตวิชฺชา พฺราหฺมณา อเวรจิตฺโต พฺรหฺมา ฯเปฯ สพฺยาปชฺชจิตฺตา เตวิชฺชา พฺราหฺมณา อพฺยาปชฺชจิตฺโต พฺรหฺมา.\csromanspacer{} สํกิลิฏฺฐจิตฺตา เตวิชฺชา พฺราหฺมณา อสํกิลิฏฺฐจิตฺโต พฺรหฺมา.\csromanspacer{} อวสวตฺตี}

\vspace{\tipitakaparskip}
\csromancenter{เตวิชฺชสุตฺต เตวิชฺชา พฺราหฺมณา วสวตฺตี พฺรหฺมา, อปิ นุ โข อวสวตฺตีนํ เตวิชฺชานํ พฺราหฺมณานํ วสวตฺตินา พฺรหฺมุนา สทฺธิํ สํสนฺทติ สเมตีติ.\csromanspacer{} โน หิทํ โภ โคตม.\csromanspacer{} สาธุ วาเสฏฺฐ, เต วต วาเสฏฺฐ อวสวตฺตี เตวิชฺชา พฺราหฺมณา กายสฺส เภทา ปรํ มรณา วสวตฺติสฺส พฺรหฺมุโน สหพฺยูปคา ภวิสฺสนฺตีติ เนตํ ฐานํ วิชฺชติ.}
\proseitem{552}{\csromanfolio{233}อิธ โข ปน เต วาเสฏฺฐ เตวิชฺชา พฺราหฺมณา อาสีทิตฺวา\footnote{อาทิสิตฺวา (ก)} สํสีทนฺติ, สํสีทิตฺวา วิสารํ\footnote{วิสาทํ (สี, อิ), วิสตฺตํ (สฺยา)} ปาปุณนฺติ, สุกฺขตรํ\footnote{สุกฺขตรณํ (ก)} มญฺเญ ตรนฺติ.\csromanspacer{} ตสฺมา อิทํ เตวิชฺชานํ พฺราหฺมณานํ เตวิชฺชา-อิริณนฺติปิ วุจฺจติ, เตวิชฺชาวิวนนฺติปิ วุจฺจติ, เตวิชฺชาพฺยสนนฺติปิ วุจฺจตีติ.}

\proseitem{553}{เอวํ วุตฺเต วาเสฏฺโฐ มาณโว ภควนฺตํ เอตทโวจ “สุตํ เมตํ โภ โคตม, สมโณ โคตโม พฺรหฺมานํ สหพฺยตาย มคฺคํ ชานาตี”ติ.\csromanspacer{} ตํ กิํ มญฺญสิ วาเสฏฺฐ.\csromanspacer{} อาสนฺเน อิโต มนสากฏํ, น อิโต ทูเร มนสากฏนฺติ.\csromanspacer{} เอวํ โภ โคตม อาสนฺเน อิโต มนสากฏํ, น อิโต ทูเร มนสากฏนฺติ.}

\proseitem{554}{ตํ กิํ มญฺญสิ วาเสฏฺฐ.\csromanspacer{} อิธสฺส ปุริโส มนสากเฏ ชาตสํวทฺโธ.\csromanspacer{} ตเมนํ มนสากฏโต ตาวเทว อวสฏํ มนสากฏสฺส มคฺคํ ปุจฺเฉยฺยุํ.\csromanspacer{} สิยา นุ โข วาเสฏฺฐ ตสฺส ปุริสสฺส มนสากเฏ ชาตสํวทฺธสฺส มนสากฏสฺส มคฺคํ ปุฏฺฐสฺส ทนฺธายิตตฺตํ วา วิตฺถายิตตฺตํ วาติ.\csromanspacer{} โน หิทํ โภ โคตม.\csromanspacer{} ตํ กิสฺส เหตุ, อมุ หิ โภ โคตม ปุริโส มนสากเฏ ชาตสํวทฺโธ, ตสฺส สพฺพาเนว มนสากฏสฺส มคฺคานิ สุวิทิตานีติ.}

\prose{สิยา โข วาเสฏฺฐ ตสฺส ปุริสสฺส มนสากเฏ ชาตสํวทฺธสฺส มนสากฏสฺส มคฺคํ ปุฏฺฐสฺส ทนฺธายิตตฺตํ วา วิตฺถายิตตฺตํ วา, น เตฺวว ตถาคตสฺส พฺรหฺมโลเก วา พฺรหฺมโลกคามินิยา วา ปฏิปทาย ปุฏฺฐสฺส ทนฺธายิตตฺตํ วา วิตฺถายิตตฺตํ วา, พฺรหฺมานญฺจาหํ วาเสฏฺฐ ปชานามิ พฺรหฺมโลกญฺจ พฺรหฺมโลกคามินิญฺจ ปฏิปทํ, ยถา ปฏิปนฺโน จ พฺรหฺมโลกํ อุปปนฺโน, ตญฺจ ปชานามีติ.}

\proseitem{555}{\csromanfolio{234}เอวํ วุตฺเต วาเสฏฺโฐ มาณโว ภควนฺตํ เอตทโวจ “สุตํ เมตํ โภ โคตม, สมโณ โคตโม พฺรหฺมานํ สหพฺยตาย มคฺคํ เทเสตี”ติ.\csromanspacer{} สาธุ โน ภวํ โคตโม พฺรหฺมานํ สหพฺยตาย มคฺคํ เทเสตุ อุลฺลุมฺปตุ ภวํ โคตโม พฺราหฺมณิํ ปชนฺติ.\csromanspacer{} เตน หิ วาเสฏฺฐ สุณาหิ สาธุกํ มนสิ กโรหิ ภาสิสฺสามีติ. “เอวํ โภ”ติ โข วาเสฏฺโฐ มาณโว ภควโต ปจฺจสฺโสสิ.}

\csromanmatikaanchor{mtk.144}
\csromantocmarkref{section}{พฺรหฺมโลกมคฺคเทสนา}{mtk.144}
\bookmark[dest={mtk.144},level=1]{พฺรหฺมโลกมคฺคเทสนา}
\csromanheader{1}{พฺรหฺมโลกมคฺค\-เทสนา}
\proseitem{556}{ภควา เอตทโวจ.\csromanspacer{} อิธ วาเสฏฺฐ ตถาคโต โลเก อุปฺปชฺชติ อรหํ สมฺมาสมฺพุทฺโธ ฯเปฯ. (ยถา สามญฺญผเล, เอวํ วิตฺถาเรตพฺพํ.) เอวํ โข วาเสฏฺฐ ภิกฺขุ สีลสมฺปนฺโน โหติ ฯเปฯ ตสฺสิเม ปญฺจ นีวรเณ ปหีเน อตฺตนิ สมนุปสฺสโต ปาโมชฺชํ ชายติ, ปมุทิตสฺส ปีติ ชายติ, ปีติมนสฺส กาโย ปสฺสมฺภติ, ปสฺสทฺธกาโย สุขํ เวเทติ, สุขิโน จิตฺตํ สมาธิยติ.}

\prose{โส เมตฺตาสหคเตน เจตสา เอกํ ทิสํ ผริตฺวา วิหรติ.\csromanspacer{} ตถา ทุติยํ.\csromanspacer{} ตถา ตติยํ.\csromanspacer{} ตถา จตุตฺถํ.\csromanspacer{} อิติ อุทฺธมโธ ติริยํ สพฺพธิ สพฺพตฺตตาย สพฺพาวนฺตํ โลกํ เมตฺตาสหคเตน เจตสา วิปุเลน มหคฺคเตน อปฺปมาเณน อเวเรน อพฺยาปชฺเชน ผริตฺวา วิหรติ.}

\prose{เสยฺยถาปิ วาเสฏฺฐ พลวา สงฺขธโม อปฺปกสิเรเนว จตุทฺทิสา วิญฺญาเปยฺย, เอวเมว โข วาเสฏฺฐ เอวํ ภาวิตาย เมตฺตาย เจโตวิมุตฺติยา ยํ ปมาณกตํ กมฺมํ น ตํ ตตฺราวสิสฺสติ, น ตํ ตตฺราวติฏฺฐติ.\csromanspacer{} อยมฺปิ โข วาเสฏฺฐ พฺรหฺมานํ สหพฺยตาย มคฺโค.}

\prose{ปุน จปรํ วาเสฏฺฐ ภิกฺขุ กรุณาสหคเตน เจตสา ฯเปฯ มุทิตาสหคเตน เจตสา ฯเปฯ อุเปกฺขาสหคเตน เจตสา เอกํ ทิสํ ผริตฺวา วิหรติ.\csromanspacer{} ตถา ทุติยํ.\csromanspacer{} ตถา ตติยํ.\csromanspacer{} ตถา จตุตฺถํ.\csromanspacer{} อิติ อุทฺธมโธ ติริยํ สพฺพธิ สพฺพตฺตตาย สพฺพาวนฺตํ โลกํ อุเปกฺขาสหคเตน เจตสา วิปุเลน มหคฺคเตน อปฺปมาเณน อเวเรน อพฺยาปชฺเชน ผริตฺวา วิหรติ.}

\csromanheader{2}{13. เตวิชฺชสุตฺต}
\prose{\csromanfolio{235}เสยฺยถาปิ วาเสฏฺฐ พลวา สงฺขธโม อปฺปกสิเรเนว จตุทฺทิสา วิญฺญาเปยฺย.\csromanspacer{} เอวเมว โข วาเสฏฺฐ เอวํ ภาวิตาย อุเปกฺขาย เจโตวิมุตฺติยา ยํ ปมาณกตํ กมฺมํ น ตํ ตตฺราวสิสฺสติ, น ตํ ตตฺราวติฏฺฐติ.\csromanspacer{} อยมฺปิ โข วาเสฏฺฐ พฺรหฺมานํ สหพฺยตาย มคฺโค.}

\proseitem{557}{ตํ กิํ มญฺญสิ วาเสฏฺฐ, เอวํวิหารี ภิกฺขุ สปริคฺคโห วา อปริคฺคโห วาติ.\csromanspacer{} อปริคฺคโห โภ โคตม.\csromanspacer{} สเวรจิตฺโต วา อเวรจิตฺโต วาติ.\csromanspacer{} อเวรจิตฺโต โภ โคตม.\csromanspacer{} สพฺยาปชฺชจิตฺโต วา อพฺยาปชฺชจิตฺโต วาติ.\csromanspacer{} อพฺยาปชฺชจิตฺโต โภ โคตม.\csromanspacer{} สํกิลิฏฺฐจิตฺโต วา อสํกิลิฏฺฐจิตฺโต วาติ.\csromanspacer{} อสํกิลิฏฺฐจิตฺโต โภ โคตม.\csromanspacer{} วสวตฺตี วา อวสวตฺตี วาติ.\csromanspacer{} วสวตฺตี โภ โคตม.}

\prose{อิติ กิร วาเสฏฺฐ อปริคฺคโห ภิกฺขุ อปริคฺคโห พฺรหฺมา, อปิ นุ โข อปริคฺคหสฺส ภิกฺขุโน อปริคฺคเหน พฺรหฺมุนา สทฺธิํ สํสนฺทติ สเมตีติ.\csromanspacer{} เอวํ โภ โคตม.\csromanspacer{} สาธุ วาเสฏฺฐ, โส วต วาเสฏฺฐ อปริคฺคโห ภิกฺขุ กายสฺส เภทา ปรํ มรณา อปริคฺคหสฺส พฺรหฺมุโน สหพฺยูปโค ภวิสฺสตีติ ฐานเมตํ วิชฺชติ.}

\proseitem{558}{อิติ กิร วาเสฏฺฐ อเวรจิตฺโต ภิกฺขุ อเวรจิตฺโต พฺรหฺมา ฯเปฯ อพฺยาปชฺชจิตฺโต ภิกฺขุ อพฺยาปชฺชจิตฺโต พฺรหฺมา.\csromanspacer{} อสํกิลิฏฺฐจิตฺโต ภิกฺขุ อสํกิลิฏฺฐจิตฺโต พฺรหฺมา.\csromanspacer{} วสวตฺตี ภิกฺขุ วสวตฺตี พฺรหฺมา, อปิ นุ โข วสวตฺติสฺส ภิกฺขุโน วสวตฺตินา พฺรหฺมุนา สทฺธิํ สํสนฺทติ สเมตีติ.\csromanspacer{} เอวํ โภ โคตม.\csromanspacer{} สาธุ วาเสฏฺฐ, โส วต วาเสฏฺฐ วสวตฺตี ภิกฺขุ กายสฺส เภทา ปรํ มรณา วสวตฺติสฺส พฺรหฺมุโน สหพฺยูปโค ภวิสฺสตีติ ฐานเมตํ วิชฺชตีติ.}

\csromanmatikaanchor{mtk.145}
\csromantocmarkref{section}{อุทฺทานคาถา}{mtk.145}
\bookmark[dest={mtk.145},level=1]{อุทฺทานคาถา}
\proseitemwithcloser{559}{เอวํ วุตฺเต วาเสฏฺฐ ภารทฺวาชา มาณวา ภควนฺตํ เอตทโวจุํ “อภิกฺกนฺตํ โภ โคตม, อภิกฺกนฺตํ โภ โคตม, เสยฺยถาปิ โภ โคตม นิกฺกุชฺชิตํ วา อุกฺกุชฺเชยฺย, ปฏิจฺฉนฺนํ วา วิวเรยฺย, มูฬฺหสฺส วา มคฺคํ อาจิกฺเขยฺย, อนฺธกาเร วา เตลปชฺโชตํ ธาเรยฺย ‘จกฺขุมนฺโต รูปานิ ทกฺขนฺตี’ติ.\csromanspacer{} เอวเมวํ โภตา โคตเมน อเนกปริยาเยน ธมฺโม \csromanfolio{236}ปกาสิโต, เอเต มยํ ภวนฺตํ โคตมํ สรณํ คจฺฉาม, ธมฺมญฺจ ภิกฺขุสํฆญฺจ, อุปาสเก โน ภวํ โคตโม ธาเรตุ อชฺชตคฺเค ปาณุเปเต สรณํ คเต”ติ.}{เตวิชฺชสุตฺตํ นิฏฺฐิตํ เตรสมํ.}
\nitthitam{สีลกฺขนฺธวคฺโค นิฏฺฐิโต.}
\titlehead{\textbf{ตสฺสุทฺทานํ}}
\setlength{\csromangathapagewidth}{0pt}
\csromangathameasuregroup{}{พฺรหฺมาสามญฺญ-อมฺพฏฺฐ, \\ โสณกูฏมหาลิชาลินี. \\ สีหโปฏฺฐปาทสุโภ เกวฏฺโฏ, \\ โลหิจฺจเตวิชฺชา เตรสาติ.}
\setlength{\csromangathaleftcol}{0pt}
\csromangathagroup{\csromangathastanza{พฺรหฺมาสามญฺญ-อมฺพฏฺฐ, \\ โสณกูฏมหาลิชาลินี. \\ สีหโปฏฺฐปาทสุโภ เกวฏฺโฏ, \\ โลหิจฺจเตวิชฺชา เตรสาติ.}}

\nitthitam{สีลกฺขนฺธวคฺคปาฬิ นิฏฺฐิตา.}
